% ======================================================================
% preamble.tex -- shared preamble of the main paper and the supplement.
% Owner: coordinator.  Drafters add macros by request only (see DRAFTER_GUIDE.md).
%
% USE.  This file contains NO \documentclass.  The two master files are
%   paper_v1/main.tex        \documentclass[11pt]{amsart} % ======================================================================
% preamble.tex -- shared preamble of the main paper and the supplement.
% Owner: coordinator.  Drafters add macros by request only (see DRAFTER_GUIDE.md).
%
% USE.  This file contains NO \documentclass.  The two master files are
%   paper_v1/main.tex        \documentclass[11pt]{amsart} % ======================================================================
% preamble.tex -- shared preamble of the main paper and the supplement.
% Owner: coordinator.  Drafters add macros by request only (see DRAFTER_GUIDE.md).
%
% USE.  This file contains NO \documentclass.  The two master files are
%   paper_v1/main.tex        \documentclass[11pt]{amsart} % ======================================================================
% preamble.tex -- shared preamble of the main paper and the supplement.
% Owner: coordinator.  Drafters add macros by request only (see DRAFTER_GUIDE.md).
%
% USE.  This file contains NO \documentclass.  The two master files are
%   paper_v1/main.tex        \documentclass[11pt]{amsart} \input{preamble} ...
%   paper_v1/supplement.tex  \documentclass[11pt]{amsart}
%                            \newif\ifsupplement \supplementtrue
%                            \input{preamble} ...
% (the main paper does nothing special: \ifsupplement defaults to false).
% Compile from paper_v1/ so that \input{numbers/numbers_paper.tex} is found.
% ======================================================================

% ---- which document are we in ----------------------------------------
\ifdefined\ifsupplement\else\newif\ifsupplement\fi

% ---- packages ---------------------------------------------------------
\usepackage[T1]{fontenc}
\usepackage[utf8]{inputenc}
\usepackage{lmodern}
\usepackage[margin=1.1in]{geometry}
\usepackage{amsmath,amssymb,amsthm,mathtools,bm}
\usepackage{booktabs,array,longtable,tabularx,multirow}
\usepackage{enumitem}
\usepackage{graphicx}
\usepackage{xcolor}
\usepackage{tikz}
\usetikzlibrary{positioning,arrows.meta,calc}
% xr-hyper must come before hyperref
\usepackage{xr-hyper}
\usepackage[colorlinks=true,linkcolor=blue!50!black,citecolor=green!35!black,urlcolor=blue!50!black]{hyperref}

\raggedbottom
\setlength{\emergencystretch}{2em}
\setlist[enumerate]{leftmargin=*,itemsep=2pt,topsep=3pt}
\setlist[itemize]{leftmargin=*,itemsep=2pt,topsep=3pt}

% ---- numbering --------------------------------------------------------
% Main paper: sections 1,2,...; theorems 2.3 (by section); equations (2.3).
% Supplement: the Parts A..K are the \section's, numbered by letter, so
% theorems are A.1, E.4, ... (do not confuse "Theorem A" of the paper, which
% has no dot, with "Theorem A.1" of Part A).
\ifsupplement
  \renewcommand{\thesection}{\Alph{section}}
\fi
\numberwithin{equation}{section}

% ---- theorem environments (numbered by section, one shared counter) ----
\theoremstyle{plain}
\newtheorem{theorem}{Theorem}[section]
\newtheorem{proposition}[theorem]{Proposition}
\newtheorem{lemma}[theorem]{Lemma}
\newtheorem{corollary}[theorem]{Corollary}
\newtheorem{maintheorem}[theorem]{Main Theorem}
% Theorems A, B, C, D of the paper: own counter, lettered, no section prefix
\newtheorem{thmlet}{Theorem}
\renewcommand{\thethmlet}{\Alph{thmlet}}
\theoremstyle{definition}
\newtheorem{definition}[theorem]{Definition}
\newtheorem{convention}[theorem]{Convention}
\newtheorem{example}[theorem]{Example}
\theoremstyle{remark}
\newtheorem{remark}[theorem]{Remark}

% A list of hypotheses labelled (H1),(H2),... : \begin{hyplist}{H} ... \end{hyplist}
% (also {O}, {L}, {B}, {K}, {W} for the families of section 5.3 of OUTLINE.md).
\newenvironment{hyplist}[1]%
  {\begin{enumerate}[label=\textup{(#1\arabic*)},leftmargin=3em,itemsep=2pt]}%
  {\end{enumerate}}

% ---- trust classes (OUTLINE 4.1, 5.1.3) -------------------------------
% Supplement: explicit tags.  Main paper: certified is the default and
% unmarked; every non-certified number is marked "(floating point)".
\ifsupplement
  \newcommand{\tcert}{\textup{[cert]}}   % arb/outward-rounded, pinned script and key
  \newcommand{\test}{\textup{[est]}}     % floating point or estimate; never in a theorem
  \newcommand{\thand}{\textup{[hand]}}   % proved by hand in the text
  \newcommand{\tlit}{\textup{[lit]}}     % cited from the literature
\else
  \newcommand{\tcert}{}
  \newcommand{\test}{\unskip\ \textup{(floating point)}}
  \newcommand{\thand}{}
  \newcommand{\tlit}{}
\fi

% ---- gaps: visible, counted, an error in a final build -----------------
% Usage: \GAP{one sentence saying what is missing}   (and an entry in ISSUES.md)
\newcounter{gapcount}
\newif\iffinalpaper   % \finalpapertrue makes every \GAP a hard error
\newcommand{\GAP}[1]{\stepcounter{gapcount}\typeout{GAP-MARKER \thegapcount}%
  \iffinalpaper\PackageError{paper}{\string\GAP\space used in a final build}{Remove all gaps.}\fi
  {\color{red}\sffamily\bfseries[GAP~\thegapcount:\ }{\color{red}\sffamily #1}{\color{red}\sffamily\bfseries]}}
\AtEndDocument{\typeout{GAP-MARKER total: \thegapcount}}

% ---- file names and JSON keys (supplement only) -------------------------
% No internal path appears in the main paper; a use there is a warning.
\newcommand{\fn}[1]{%
  \ifsupplement\else\PackageWarning{paper}{\string\fn\space used in the main paper}\fi
  \texttt{\detokenize{#1}}}
\newcommand{\key}[1]{\texttt{\detokenize{#1}}}

% ---- cross references between main paper and supplement ----------------
% Labels carry a unit prefix (section 5.3 of OUTLINE.md and DRAFTER_GUIDE.md):
%   main paper:  intro: obj: s3: s4: s5: s6: s7: s8: s9:   (kind: thm- lem- eq- tab- fig- ...)
%   supplement:  pA: pB: ... pK:
% so all labels are globally unique and the prefixes below are only a second guard.
% \Xref{label}  : reference to the OTHER document (hyperlinked via xr-hyper)
% \Lref{label}  : reference in the SAME document (plain \ref)
\ifsupplement
  \externaldocument[M-]{main}[main.pdf]
  \newcommand{\Xref}[1]{\ref{M-#1}}
\else
  \externaldocument[S-]{supplement}[supplement.pdf]
  \newcommand{\Xref}[1]{\ref{S-#1}}
\fi
\newcommand{\Lref}[1]{\ref{#1}}
% \Supp{pE}{lem-chart} -> "Lemma E.3" style text is produced by the drafter:
%   "Lemma~\Xref{pE:lem-chart} of the supplement".  \SuppPart{E} gives "Part E".
\newcommand{\SuppPart}[1]{Part~#1 of the supplement}

% ---- numbers (generated by unit U0; macros only, no literal certified numbers) ----
\IfFileExists{numbers/numbers_paper.tex}%
  {\input{numbers/numbers_paper.tex}}%
  {\PackageWarning{paper}{numbers/numbers_paper.tex not found (unit U0 not finished)}}

% ======================================================================
% NOTATION MACROS (OUTLINE 5.2 and notation.md).  One macro per symbol
% whose spelling is fixed by the notation table.  Plain letters (a, b, f, ...)
% have no macro.
% ======================================================================
% basic
\newcommand{\R}{\mathbb{R}}
\newcommand{\C}{\mathbb{C}}
\newcommand{\N}{\mathbb{N}}
\newcommand{\Z}{\mathbb{Z}}
\newcommand{\eps}{\varepsilon}
\newcommand{\ip}[2]{\langle #1,#2\rangle}
\newcommand{\abs}[1]{\lvert #1\rvert}
\newcommand{\norm}[1]{\lVert #1\rVert}
\newcommand{\dd}{\mathrm d}
\newcommand{\re}{\operatorname{Re}}
\newcommand{\im}{\operatorname{Im}}
% flow and curves
\newcommand{\arclen}{\mathfrak{s}}                 % arclength of the flow curve
\newcommand{\kss}{\kappa_{\mathfrak{ss}}}          % kappa_{ss} in the flow equation
\newcommand{\Area}{\operatorname{Area}}
\newcommand{\Len}{\mathcal L}                      % length of a curve
\newcommand{\spp}{\operatorname{sp}}                 % sup quantities sp^v_0, sp^v_1, sp^f_1 (were S^v_0, S^v_1, S^f_1)
\newcommand{\rot}{\operatorname{rot}}              % turning number
\newcommand{\Tsing}{T_*}                           % singular time
\newcommand{\Tmax}{T_{\max}}                       % maximal time of existence
\newcommand{\pstar}{p_*}                           % image point of the singularity
\newcommand{\Gz}{\Gamma_0}                         % initial curve
\newcommand{\vth}{\vartheta}                       % parameter of Gamma_0
\newcommand{\dclk}{\delta_t}                       % clock defect of step 4
\newcommand{\fall}{\varpi}                         % Fourier rates varpi_n = n^4-n^2
% rescaling
\newcommand{\lamone}{\lambda_1}
\newcommand{\rate}{\varsigma}                      % 1/4 - a
\newcommand{\xia}{\xi_a}                           % (1-4\bar a)^{-1}
% shrinker and step 1
\newcommand{\Sone}{S_1}
\newcommand{\Gsh}{\Gamma_{\!*}}                    % profile curve of S_1
\newcommand{\cSet}{\mathcal S}                     % certified set of shrinkers
\newcommand{\sigtip}{\sigma_{\rm tip}}
\newcommand{\Xo}{X_0}                              % 2.2490776515
\newcommand{\Dtip}{D}
\newcommand{\ktip}{k_0}
\newcommand{\shootM}{M}                            % shooting function
\newcommand{\Ypar}{\mathcal Y}                     % step-1 parallelepiped
\newcommand{\axis}{\Sigma}                         % symmetry axis of S_1
\newcommand{\sigc}{\sigma_{\rm c}}
\newcommand{\sigrev}{\sigma^{\rm rev}}
% linear theory
\newcommand{\Hrho}{\mathcal H_\rho}
\newcommand{\muslow}{\mu_3}
\newcommand{\rre}{\mathfrak r_{\rm re}}
\newcommand{\rim}{\mathfrak r_{\rm im}}
\newcommand{\Evev}{\mathrm{Ev}_{\rm even}}
\newcommand{\Evod}{\mathrm{Ev}_{\rm odd}}
\newcommand{\gvec}{\mathbf g}                      % translation vector in step 2
\newcommand{\cons}[1]{\mathfrak f_{#1}}            % constraint functions
\newcommand{\phirho}{\varphi_\rho}
% energies
\newcommand{\Lam}{\Lambda}
\newcommand{\En}{E}
\newcommand{\EnI}{E_I}
\newcommand{\muI}{\mu_I}                           % 0.08
\newcommand{\epsE}{\varepsilon_E}                  % 1e-5
\newcommand{\epsI}{\varepsilon_I}                  % 3e-5
\newcommand{\wn}[1]{\mathsf w_{#1}}                % ||w~^{(j)}||_rho
\newcommand{\vI}{v_I}                              % zeta w'
\newcommand{\muerr}{\mu_{\rm err}}                 % step-4 error-energy weight 0.036
% far field
\newcommand{\pU}{p_U}                              % 1.0246
\newcommand{\Uone}{U_1}
\newcommand{\sm}{\operatorname{sm}}                % smoothstep t^3(10-15t+6t^2)
\newcommand{\rc}{r_c}                              % centre radius (Lemma U)
\newcommand{\rchi}{r_\chi}                         % 253
\newcommand{\Iw}{I_w}                              % hand-off windows [y_w,1.5y_w]
\newcommand{\qa}{\mathfrak q_1}                    % quasilinear coefficients of the slope equation
\newcommand{\qb}{\mathfrak q_2}
% cylinders (Lemma O)
\newcommand{\cyl}{\ell}                            % cylinder length unit
\newcommand{\MO}{M_O}
\newcommand{\KO}{K_O}
\newcommand{\Hcyl}{H_{\rm cyl}}
\newcommand{\XiO}{\Xi_O}
\newcommand{\fO}{f_O}
% step 4
\newcommand{\Ybar}{\bar Y}
\newcommand{\Lt}{\widetilde L}                     % linearisation about the step-4 base
\newcommand{\Trans}{\Pi}                           % transport integral
\newcommand{\resid}{\mathrm{res}}
\newcommand{\Resid}[1]{\mathrm{Res}_{#1}}
\newcommand{\boxtip}{\Box_{\rm tip}}
\newcommand{\boxfar}{\Box_{\rm far}}
\newcommand{\Eenv}{\mathcal E_{\rm env}}
\newcommand{\betab}{\beta_{\rm base}}
\newcommand{\tauI}{\tau^{\rm I}}                   % Phase I clock
% names of families of conditions: \Hh{3} = (H3), \Oh{4} = (O4), ...
\newcommand{\Hh}[1]{\textup{(H#1)}}
\newcommand{\Oh}[1]{\textup{(O#1)}}
\newcommand{\Lh}[1]{\textup{(L#1)}}
\newcommand{\Bh}[1]{\textup{(B#1)}}
\newcommand{\Kh}[1]{\textup{(K#1)}}
\newcommand{\Wh}[1]{\textup{(W#1)}}
\newcommand{\PhaseI}{Phase~I}
\newcommand{\PhaseII}{Phase~II}
% margin: threshold/bound (always)
\newcommand{\mgn}[1]{{\times}#1}

% ---- figures: placeholder macro so that drafts compile before F1-F7 exist ----
\newcommand{\figplaceholder}[1]{\fbox{\parbox{0.8\linewidth}{\centering\small[figure #1 to be supplied]}}}
 ...
%   paper_v1/supplement.tex  \documentclass[11pt]{amsart}
%                            \newif\ifsupplement \supplementtrue
%                            % ======================================================================
% preamble.tex -- shared preamble of the main paper and the supplement.
% Owner: coordinator.  Drafters add macros by request only (see DRAFTER_GUIDE.md).
%
% USE.  This file contains NO \documentclass.  The two master files are
%   paper_v1/main.tex        \documentclass[11pt]{amsart} \input{preamble} ...
%   paper_v1/supplement.tex  \documentclass[11pt]{amsart}
%                            \newif\ifsupplement \supplementtrue
%                            \input{preamble} ...
% (the main paper does nothing special: \ifsupplement defaults to false).
% Compile from paper_v1/ so that \input{numbers/numbers_paper.tex} is found.
% ======================================================================

% ---- which document are we in ----------------------------------------
\ifdefined\ifsupplement\else\newif\ifsupplement\fi

% ---- packages ---------------------------------------------------------
\usepackage[T1]{fontenc}
\usepackage[utf8]{inputenc}
\usepackage{lmodern}
\usepackage[margin=1.1in]{geometry}
\usepackage{amsmath,amssymb,amsthm,mathtools,bm}
\usepackage{booktabs,array,longtable,tabularx,multirow}
\usepackage{enumitem}
\usepackage{graphicx}
\usepackage{xcolor}
\usepackage{tikz}
\usetikzlibrary{positioning,arrows.meta,calc}
% xr-hyper must come before hyperref
\usepackage{xr-hyper}
\usepackage[colorlinks=true,linkcolor=blue!50!black,citecolor=green!35!black,urlcolor=blue!50!black]{hyperref}

\raggedbottom
\setlength{\emergencystretch}{2em}
\setlist[enumerate]{leftmargin=*,itemsep=2pt,topsep=3pt}
\setlist[itemize]{leftmargin=*,itemsep=2pt,topsep=3pt}

% ---- numbering --------------------------------------------------------
% Main paper: sections 1,2,...; theorems 2.3 (by section); equations (2.3).
% Supplement: the Parts A..K are the \section's, numbered by letter, so
% theorems are A.1, E.4, ... (do not confuse "Theorem A" of the paper, which
% has no dot, with "Theorem A.1" of Part A).
\ifsupplement
  \renewcommand{\thesection}{\Alph{section}}
\fi
\numberwithin{equation}{section}

% ---- theorem environments (numbered by section, one shared counter) ----
\theoremstyle{plain}
\newtheorem{theorem}{Theorem}[section]
\newtheorem{proposition}[theorem]{Proposition}
\newtheorem{lemma}[theorem]{Lemma}
\newtheorem{corollary}[theorem]{Corollary}
\newtheorem{maintheorem}[theorem]{Main Theorem}
% Theorems A, B, C, D of the paper: own counter, lettered, no section prefix
\newtheorem{thmlet}{Theorem}
\renewcommand{\thethmlet}{\Alph{thmlet}}
\theoremstyle{definition}
\newtheorem{definition}[theorem]{Definition}
\newtheorem{convention}[theorem]{Convention}
\newtheorem{example}[theorem]{Example}
\theoremstyle{remark}
\newtheorem{remark}[theorem]{Remark}

% A list of hypotheses labelled (H1),(H2),... : \begin{hyplist}{H} ... \end{hyplist}
% (also {O}, {L}, {B}, {K}, {W} for the families of section 5.3 of OUTLINE.md).
\newenvironment{hyplist}[1]%
  {\begin{enumerate}[label=\textup{(#1\arabic*)},leftmargin=3em,itemsep=2pt]}%
  {\end{enumerate}}

% ---- trust classes (OUTLINE 4.1, 5.1.3) -------------------------------
% Supplement: explicit tags.  Main paper: certified is the default and
% unmarked; every non-certified number is marked "(floating point)".
\ifsupplement
  \newcommand{\tcert}{\textup{[cert]}}   % arb/outward-rounded, pinned script and key
  \newcommand{\test}{\textup{[est]}}     % floating point or estimate; never in a theorem
  \newcommand{\thand}{\textup{[hand]}}   % proved by hand in the text
  \newcommand{\tlit}{\textup{[lit]}}     % cited from the literature
\else
  \newcommand{\tcert}{}
  \newcommand{\test}{\unskip\ \textup{(floating point)}}
  \newcommand{\thand}{}
  \newcommand{\tlit}{}
\fi

% ---- gaps: visible, counted, an error in a final build -----------------
% Usage: \GAP{one sentence saying what is missing}   (and an entry in ISSUES.md)
\newcounter{gapcount}
\newif\iffinalpaper   % \finalpapertrue makes every \GAP a hard error
\newcommand{\GAP}[1]{\stepcounter{gapcount}\typeout{GAP-MARKER \thegapcount}%
  \iffinalpaper\PackageError{paper}{\string\GAP\space used in a final build}{Remove all gaps.}\fi
  {\color{red}\sffamily\bfseries[GAP~\thegapcount:\ }{\color{red}\sffamily #1}{\color{red}\sffamily\bfseries]}}
\AtEndDocument{\typeout{GAP-MARKER total: \thegapcount}}

% ---- file names and JSON keys (supplement only) -------------------------
% No internal path appears in the main paper; a use there is a warning.
\newcommand{\fn}[1]{%
  \ifsupplement\else\PackageWarning{paper}{\string\fn\space used in the main paper}\fi
  \texttt{\detokenize{#1}}}
\newcommand{\key}[1]{\texttt{\detokenize{#1}}}

% ---- cross references between main paper and supplement ----------------
% Labels carry a unit prefix (section 5.3 of OUTLINE.md and DRAFTER_GUIDE.md):
%   main paper:  intro: obj: s3: s4: s5: s6: s7: s8: s9:   (kind: thm- lem- eq- tab- fig- ...)
%   supplement:  pA: pB: ... pK:
% so all labels are globally unique and the prefixes below are only a second guard.
% \Xref{label}  : reference to the OTHER document (hyperlinked via xr-hyper)
% \Lref{label}  : reference in the SAME document (plain \ref)
\ifsupplement
  \externaldocument[M-]{main}[main.pdf]
  \newcommand{\Xref}[1]{\ref{M-#1}}
\else
  \externaldocument[S-]{supplement}[supplement.pdf]
  \newcommand{\Xref}[1]{\ref{S-#1}}
\fi
\newcommand{\Lref}[1]{\ref{#1}}
% \Supp{pE}{lem-chart} -> "Lemma E.3" style text is produced by the drafter:
%   "Lemma~\Xref{pE:lem-chart} of the supplement".  \SuppPart{E} gives "Part E".
\newcommand{\SuppPart}[1]{Part~#1 of the supplement}

% ---- numbers (generated by unit U0; macros only, no literal certified numbers) ----
\IfFileExists{numbers/numbers_paper.tex}%
  {\input{numbers/numbers_paper.tex}}%
  {\PackageWarning{paper}{numbers/numbers_paper.tex not found (unit U0 not finished)}}

% ======================================================================
% NOTATION MACROS (OUTLINE 5.2 and notation.md).  One macro per symbol
% whose spelling is fixed by the notation table.  Plain letters (a, b, f, ...)
% have no macro.
% ======================================================================
% basic
\newcommand{\R}{\mathbb{R}}
\newcommand{\C}{\mathbb{C}}
\newcommand{\N}{\mathbb{N}}
\newcommand{\Z}{\mathbb{Z}}
\newcommand{\eps}{\varepsilon}
\newcommand{\ip}[2]{\langle #1,#2\rangle}
\newcommand{\abs}[1]{\lvert #1\rvert}
\newcommand{\norm}[1]{\lVert #1\rVert}
\newcommand{\dd}{\mathrm d}
\newcommand{\re}{\operatorname{Re}}
\newcommand{\im}{\operatorname{Im}}
% flow and curves
\newcommand{\arclen}{\mathfrak{s}}                 % arclength of the flow curve
\newcommand{\kss}{\kappa_{\mathfrak{ss}}}          % kappa_{ss} in the flow equation
\newcommand{\Area}{\operatorname{Area}}
\newcommand{\Len}{\mathcal L}                      % length of a curve
\newcommand{\spp}{\operatorname{sp}}                 % sup quantities sp^v_0, sp^v_1, sp^f_1 (were S^v_0, S^v_1, S^f_1)
\newcommand{\rot}{\operatorname{rot}}              % turning number
\newcommand{\Tsing}{T_*}                           % singular time
\newcommand{\Tmax}{T_{\max}}                       % maximal time of existence
\newcommand{\pstar}{p_*}                           % image point of the singularity
\newcommand{\Gz}{\Gamma_0}                         % initial curve
\newcommand{\vth}{\vartheta}                       % parameter of Gamma_0
\newcommand{\dclk}{\delta_t}                       % clock defect of step 4
\newcommand{\fall}{\varpi}                         % Fourier rates varpi_n = n^4-n^2
% rescaling
\newcommand{\lamone}{\lambda_1}
\newcommand{\rate}{\varsigma}                      % 1/4 - a
\newcommand{\xia}{\xi_a}                           % (1-4\bar a)^{-1}
% shrinker and step 1
\newcommand{\Sone}{S_1}
\newcommand{\Gsh}{\Gamma_{\!*}}                    % profile curve of S_1
\newcommand{\cSet}{\mathcal S}                     % certified set of shrinkers
\newcommand{\sigtip}{\sigma_{\rm tip}}
\newcommand{\Xo}{X_0}                              % 2.2490776515
\newcommand{\Dtip}{D}
\newcommand{\ktip}{k_0}
\newcommand{\shootM}{M}                            % shooting function
\newcommand{\Ypar}{\mathcal Y}                     % step-1 parallelepiped
\newcommand{\axis}{\Sigma}                         % symmetry axis of S_1
\newcommand{\sigc}{\sigma_{\rm c}}
\newcommand{\sigrev}{\sigma^{\rm rev}}
% linear theory
\newcommand{\Hrho}{\mathcal H_\rho}
\newcommand{\muslow}{\mu_3}
\newcommand{\rre}{\mathfrak r_{\rm re}}
\newcommand{\rim}{\mathfrak r_{\rm im}}
\newcommand{\Evev}{\mathrm{Ev}_{\rm even}}
\newcommand{\Evod}{\mathrm{Ev}_{\rm odd}}
\newcommand{\gvec}{\mathbf g}                      % translation vector in step 2
\newcommand{\cons}[1]{\mathfrak f_{#1}}            % constraint functions
\newcommand{\phirho}{\varphi_\rho}
% energies
\newcommand{\Lam}{\Lambda}
\newcommand{\En}{E}
\newcommand{\EnI}{E_I}
\newcommand{\muI}{\mu_I}                           % 0.08
\newcommand{\epsE}{\varepsilon_E}                  % 1e-5
\newcommand{\epsI}{\varepsilon_I}                  % 3e-5
\newcommand{\wn}[1]{\mathsf w_{#1}}                % ||w~^{(j)}||_rho
\newcommand{\vI}{v_I}                              % zeta w'
\newcommand{\muerr}{\mu_{\rm err}}                 % step-4 error-energy weight 0.036
% far field
\newcommand{\pU}{p_U}                              % 1.0246
\newcommand{\Uone}{U_1}
\newcommand{\sm}{\operatorname{sm}}                % smoothstep t^3(10-15t+6t^2)
\newcommand{\rc}{r_c}                              % centre radius (Lemma U)
\newcommand{\rchi}{r_\chi}                         % 253
\newcommand{\Iw}{I_w}                              % hand-off windows [y_w,1.5y_w]
\newcommand{\qa}{\mathfrak q_1}                    % quasilinear coefficients of the slope equation
\newcommand{\qb}{\mathfrak q_2}
% cylinders (Lemma O)
\newcommand{\cyl}{\ell}                            % cylinder length unit
\newcommand{\MO}{M_O}
\newcommand{\KO}{K_O}
\newcommand{\Hcyl}{H_{\rm cyl}}
\newcommand{\XiO}{\Xi_O}
\newcommand{\fO}{f_O}
% step 4
\newcommand{\Ybar}{\bar Y}
\newcommand{\Lt}{\widetilde L}                     % linearisation about the step-4 base
\newcommand{\Trans}{\Pi}                           % transport integral
\newcommand{\resid}{\mathrm{res}}
\newcommand{\Resid}[1]{\mathrm{Res}_{#1}}
\newcommand{\boxtip}{\Box_{\rm tip}}
\newcommand{\boxfar}{\Box_{\rm far}}
\newcommand{\Eenv}{\mathcal E_{\rm env}}
\newcommand{\betab}{\beta_{\rm base}}
\newcommand{\tauI}{\tau^{\rm I}}                   % Phase I clock
% names of families of conditions: \Hh{3} = (H3), \Oh{4} = (O4), ...
\newcommand{\Hh}[1]{\textup{(H#1)}}
\newcommand{\Oh}[1]{\textup{(O#1)}}
\newcommand{\Lh}[1]{\textup{(L#1)}}
\newcommand{\Bh}[1]{\textup{(B#1)}}
\newcommand{\Kh}[1]{\textup{(K#1)}}
\newcommand{\Wh}[1]{\textup{(W#1)}}
\newcommand{\PhaseI}{Phase~I}
\newcommand{\PhaseII}{Phase~II}
% margin: threshold/bound (always)
\newcommand{\mgn}[1]{{\times}#1}

% ---- figures: placeholder macro so that drafts compile before F1-F7 exist ----
\newcommand{\figplaceholder}[1]{\fbox{\parbox{0.8\linewidth}{\centering\small[figure #1 to be supplied]}}}
 ...
% (the main paper does nothing special: \ifsupplement defaults to false).
% Compile from paper_v1/ so that % numbers_paper.tex -- GENERATED from ledger.json by build_ledger.py; do not edit.
% Rounding: upper bounds up, lower bounds and margins down at the printed digit; margin = threshold / bound.
% Use as $\times\NName$ for margins.  kind: cert certified, est estimate, hand hypothesis constant, doc only a document quotes it.
% NUEhatZero [cert, upper] U_cert_v35.json :: supersolution.EA.0
\newcommand{\NUEhatZero}{\ensuremath{1.2018\cdot10^{-6}}}
% NUEhatOne [cert, upper] U_cert_v35.json :: supersolution.EA.1
\newcommand{\NUEhatOne}{\ensuremath{1.0441\cdot10^{-4}}}
% NUEhatTwo [cert, upper] U_cert_v35.json :: supersolution.EA.2
\newcommand{\NUEhatTwo}{\ensuremath{0.056694}}
% NUBhatZero [cert, upper] U_cert_v35.json :: supersolution.Bhat.0
\newcommand{\NUBhatZero}{\ensuremath{4.8735\cdot10^{-6}}}
% NUBhatOne [cert, upper] U_cert_v35.json :: supersolution.Bhat.1
\newcommand{\NUBhatOne}{\ensuremath{5.5933\cdot10^{-4}}}
% NUBhatTwo [cert, upper] U_cert_v35.json :: supersolution.Bhat.2
\newcommand{\NUBhatTwo}{\ensuremath{0.36675}}
% NUSupersolutionCZeroZero [cert, upper] U_cert_v35.json :: supersolution.C0.0
\newcommand{\NUSupersolutionCZeroZero}{\ensuremath{1.1340\cdot10^{-6}}}
% NUSupersolutionCZeroOne [cert, upper] U_cert_v35.json :: supersolution.C0.1
\newcommand{\NUSupersolutionCZeroOne}{\ensuremath{1.6623\cdot10^{-4}}}
% NUSupersolutionCOneZero [cert, upper] U_cert_v35.json :: supersolution.C1.0
\newcommand{\NUSupersolutionCOneZero}{\ensuremath{1.5241\cdot10^{-9}}}
% NUSupersolutionCOneOne [cert, upper] U_cert_v35.json :: supersolution.C1.1
\newcommand{\NUSupersolutionCOneOne}{\ensuremath{5.5255\cdot10^{-8}}}
% NUSupersolutionNTwoZero [cert, upper] U_cert_v35.json :: supersolution.N2.0
\newcommand{\NUSupersolutionNTwoZero}{\ensuremath{0.071472}}
% NUSupersolutionNTwoOne [cert, upper] U_cert_v35.json :: supersolution.N2.1
\newcommand{\NUSupersolutionNTwoOne}{\ensuremath{0.33012}}
% NUSupersolutionNTwoTwo [cert, upper] U_cert_v35.json :: supersolution.N2.2
\newcommand{\NUSupersolutionNTwoTwo}{\ensuremath{5.0406}}
% NUSupersolutionNTwoOverFourNTwoRefZero [cert, upper] U_cert_v35.json :: supersolution.N2_over_4N2ref.0
\newcommand{\NUSupersolutionNTwoOverFourNTwoRefZero}{\ensuremath{0.18}}
% NUSupersolutionNTwoOverFourNTwoRefOne [cert, upper] U_cert_v35.json :: supersolution.N2_over_4N2ref.1
\newcommand{\NUSupersolutionNTwoOverFourNTwoRefOne}{\ensuremath{0.42}}
% NUSupersolutionNTwoOverFourNTwoRefTwo [cert, upper] U_cert_v35.json :: supersolution.N2_over_4N2ref.2
\newcommand{\NUSupersolutionNTwoOverFourNTwoRefTwo}{\ensuremath{0.26}}
% NUSupersolutionMOneZHypCheckZonesIIISOneTwoBhat [cert, upper] U_cert_v35.json :: supersolution.M1z_hyp_check.zones_I_II.S1_2Bhat
\newcommand{\NUSupersolutionMOneZHypCheckZonesIIISOneTwoBhat}{\ensuremath{0.28585}}
% NUSupersolutionMOneZHypCheckZoneIIIBOnePlusPUOne [cert, upper] U_cert_v35.json :: supersolution.M1z_hyp_check.zone_III.b1_plus_pU1
\newcommand{\NUSupersolutionMOneZHypCheckZoneIIIBOnePlusPUOne}{\ensuremath{0.50786}}
% NUMarginLateMOne [cert, margin] U_cert_v35.json :: margins.late_M1
\newcommand{\NUMarginLateMOne}{\ensuremath{2.353}}
% NUMarginLateMZero [cert, margin] U_cert_v35.json :: margins.late_M0
\newcommand{\NUMarginLateMZero}{\ensuremath{4.892}}
% NUMarginsLateMZeroZ [cert, margin] U_cert_v35.json :: margins.late_M0z
\newcommand{\NUMarginsLateMZeroZ}{\ensuremath{12.99}}
% NUMarginPresentMOne [cert, margin] U_cert_v35.json :: margins.present_M1_min.value
\newcommand{\NUMarginPresentMOne}{\ensuremath{2.770}}
% NUMarginPresentMZero [cert, margin] U_cert_v35.json :: margins.present_M0_min.value
\newcommand{\NUMarginPresentMZero}{\ensuremath{2.225}}
% NUMarginsPresentMZeroZMinValue [cert, margin] U_cert_v35.json :: margins.present_M0z_min.value
\newcommand{\NUMarginsPresentMZeroZMinValue}{\ensuremath{14.37}}
% NULoopGainZeroTotal [cert, upper] U_cert_v35.json :: loop_gain.0.total
\newcommand{\NULoopGainZeroTotal}{\ensuremath{1.23\cdot10^{-7}}}
% NULoopGainOneTotal [cert, upper] U_cert_v35.json :: loop_gain.1.total
\newcommand{\NULoopGainOneTotal}{\ensuremath{0.037169}}
% NULoopGainTwoNeighbour [cert, upper] U_cert_v35.json :: loop_gain.2.neighbour
\newcommand{\NULoopGainTwoNeighbour}{\ensuremath{0.037018}}
% NULoopGainTwoThroughA [cert, upper] U_cert_v35.json :: loop_gain.2.through_A
\newcommand{\NULoopGainTwoThroughA}{\ensuremath{0.0020172}}
% NUfineBhatZero [cert, upper] U_cert_v35.json :: ufine.Bhat.0
\newcommand{\NUfineBhatZero}{\ensuremath{4.8723\cdot10^{-6}}}
% NUfineBhatOne [cert, upper] U_cert_v35.json :: ufine.Bhat.1
\newcommand{\NUfineBhatOne}{\ensuremath{2.6266\cdot10^{-4}}}
% NUfineBhatTwo [cert, upper] U_cert_v35.json :: ufine.Bhat.2
\newcommand{\NUfineBhatTwo}{\ensuremath{0.025842}}
% NUcdeltaU [cert, upper] U_cert_v35.json :: launch.c_dU
\newcommand{\NUcdeltaU}{\ensuremath{1.7187\cdot10^{-8}}}
% NULaunchWmin [cert, lower] U_cert_v35.json :: launch.wmin
\newcommand{\NULaunchWmin}{\ensuremath{0.73492}}
% NULaunchCAD [cert, upper] U_cert_v35.json :: launch.CA_D
\newcommand{\NULaunchCAD}{\ensuremath{2.2429}}
% NULaunchQD [cert, upper] U_cert_v35.json :: launch.qD
\newcommand{\NULaunchQD}{\ensuremath{0.018509}}
% NULaunchDeltaInZero [cert, upper] U_cert_v35.json :: launch.Delta_in.0
\newcommand{\NULaunchDeltaInZero}{\ensuremath{3.131\cdot10^{-12}}}
% NULaunchDeltaInOne [cert, upper] U_cert_v35.json :: launch.Delta_in.1
\newcommand{\NULaunchDeltaInOne}{\ensuremath{3.287\cdot10^{-12}}}
% NULaunchDeltaInTwo [cert, upper] U_cert_v35.json :: launch.Delta_in.2
\newcommand{\NULaunchDeltaInTwo}{\ensuremath{3.451\cdot10^{-12}}}
% NULaunchPhiZero [cert, upper] U_cert_v35.json :: launch.Phi.0
\newcommand{\NULaunchPhiZero}{\ensuremath{6.434\cdot10^{-9}}}
% NULaunchPhiOne [cert, upper] U_cert_v35.json :: launch.Phi.1
\newcommand{\NULaunchPhiOne}{\ensuremath{1.4217\cdot10^{-8}}}
% NULaunchPhiTwo [cert, upper] U_cert_v35.json :: launch.Phi.2
\newcommand{\NULaunchPhiTwo}{\ensuremath{3.8404\cdot10^{-8}}}
% NUGammaZero [cert, upper] U_cert_v35.json :: Gamma.0
\newcommand{\NUGammaZero}{\ensuremath{4.0000358}}
% NUGammaOne [cert, upper] U_cert_v35.json :: Gamma.1
\newcommand{\NUGammaOne}{\ensuremath{16.000322}}
% NUGammaTwo [cert, upper] U_cert_v35.json :: Gamma.2
\newcommand{\NUGammaTwo}{\ensuremath{64.00218}}
% NUKappa [cert, upper] U_cert_v35.json :: kappa
\newcommand{\NUKappa}{\ensuremath{8.0000358}}
% NUKappaTwo [cert, upper] U_cert_v35.json :: kappa2
\newcommand{\NUKappaTwo}{\ensuremath{64.00061}}
% NUGhatZeroZero [cert, upper] U_cert_v35.json :: Ghat.0.0
\newcommand{\NUGhatZeroZero}{\ensuremath{51322.6}}
% NUGhatOneOne [cert, upper] U_cert_v35.json :: Ghat.1.1
\newcommand{\NUGhatOneOne}{\ensuremath{33962.9}}
% NUGhatTwoTwo [cert, upper] U_cert_v35.json :: Ghat.2.2
\newcommand{\NUGhatTwoTwo}{\ensuremath{348341.1}}
% NUGhatOneZero [cert, upper] U_cert_v35.json :: Ghat.1.0
\newcommand{\NUGhatOneZero}{\ensuremath{4.938}}
% NUGhatTwoZero [cert, upper] U_cert_v35.json :: Ghat.2.0
\newcommand{\NUGhatTwoZero}{\ensuremath{27.984}}
% NUGhatTwoOne [cert, upper] U_cert_v35.json :: Ghat.2.1
\newcommand{\NUGhatTwoOne}{\ensuremath{458.53}}
% NUKZeroZero [cert, upper] U_cert_v35.json :: K.0.0
\newcommand{\NUKZeroZero}{\ensuremath{8034.21}}
% NUKZeroOne [cert, upper] U_cert_v35.json :: K.0.1
\newcommand{\NUKZeroOne}{\ensuremath{724.89}}
% NUKOneZero [cert, upper] U_cert_v35.json :: K.1.0
\newcommand{\NUKOneZero}{\ensuremath{0.08152}}
% NUKOneOne [cert, upper] U_cert_v35.json :: K.1.1
\newcommand{\NUKOneOne}{\ensuremath{6336.94}}
% NUKOneTwo [cert, upper] U_cert_v35.json :: K.1.2
\newcommand{\NUKOneTwo}{\ensuremath{545.59}}
% NUKOneThree [cert, upper] U_cert_v35.json :: K.1.3
\newcommand{\NUKOneThree}{\ensuremath{0.58700}}
% NUKTwoZero [cert, upper] U_cert_v35.json :: K.2.0
\newcommand{\NUKTwoZero}{\ensuremath{0.87186}}
% NUKTwoOne [cert, upper] U_cert_v35.json :: K.2.1
\newcommand{\NUKTwoOne}{\ensuremath{528.62}}
% NUKTwoTwo [cert, upper] U_cert_v35.json :: K.2.2
\newcommand{\NUKTwoTwo}{\ensuremath{23425.98}}
% NUKTwoThree [cert, upper] U_cert_v35.json :: K.2.3
\newcommand{\NUKTwoThree}{\ensuremath{555.62}}
% NUFootwOne [cert, upper] U_cert_v35.json :: footw.1
\newcommand{\NUFootwOne}{\ensuremath{2.6361}}
% NUFootwTwo [cert, upper] U_cert_v35.json :: footw.2
\newcommand{\NUFootwTwo}{\ensuremath{2663.3}}
% NUCThreeThree [cert, upper] U_cert_v35.json :: c33
\newcommand{\NUCThreeThree}{\ensuremath{108.02}}
% NUCTwoTwoThree [cert, upper] U_cert_v35.json :: c223
\newcommand{\NUCTwoTwoThree}{\ensuremath{7.3284}}
% NUCEta [cert, upper] U_cert_v35.json :: C_eta
\newcommand{\NUCEta}{\ensuremath{13.376}}
% NULMax [cert, upper] U_cert_v35.json :: L_max
\newcommand{\NULMax}{\ensuremath{3.5352}}
% NURRatio [cert, upper] U_cert_v35.json :: R_ratio
\newcommand{\NURRatio}{\ensuremath{11.593}}
% NUSweepsSigmaInCoverSup [cert, upper] U_cert_v35.json :: sweeps.Sigma_in.cover_sup
\newcommand{\NUSweepsSigmaInCoverSup}{\ensuremath{5.0892}}
% NUSweepsSigmaECoverSup [cert, upper] U_cert_v35.json :: sweeps.Sigma_E.cover_sup
\newcommand{\NUSweepsSigmaECoverSup}{\ensuremath{0.30846}}
% NUSweepsSXCoverSup [cert, upper] U_cert_v35.json :: sweeps.S_X.cover_sup
\newcommand{\NUSweepsSXCoverSup}{\ensuremath{1.5243}}
% NUSweepsSigmaXCoverSup [cert, upper] U_cert_v35.json :: sweeps.Sigma_X.cover_sup
\newcommand{\NUSweepsSigmaXCoverSup}{\ensuremath{3.0081}}
% NUSweepsSigmaRsTwoCoverSup [cert, upper] U_cert_v35.json :: sweeps.Sigma_rs.2.cover_sup
\newcommand{\NUSweepsSigmaRsTwoCoverSup}{\ensuremath{4.3350\cdot10^{-3}}}
% NUDataTermsLateDZero [cert, upper] U_cert_v35.json :: data_terms_late.D.0
\newcommand{\NUDataTermsLateDZero}{\ensuremath{1.7330\cdot10^{-6}}}
% NUDataTermsLateDOne [cert, upper] U_cert_v35.json :: data_terms_late.D.1
\newcommand{\NUDataTermsLateDOne}{\ensuremath{2.5353\cdot10^{-4}}}
% NUDataTermsLateSlopeShareCOneSigmaXOverDkZero [cert, upper] U_cert_v35.json :: data_terms_late.slope_share_C1SigmaX_over_Dk.0
\newcommand{\NUDataTermsLateSlopeShareCOneSigmaXOverDkZero}{\ensuremath{2.7\cdot10^{-3}}}
% NUDataTermsLateSlopeShareCOneSigmaXOverDkOne [cert, upper] U_cert_v35.json :: data_terms_late.slope_share_C1SigmaX_over_Dk.1
\newcommand{\NUDataTermsLateSlopeShareCOneSigmaXOverDkOne}{\ensuremath{6.6\cdot10^{-4}}}
% NURLtTwoMarginZero [cert, margin] U_cert_v35.json :: R_lt2_margin.0
\newcommand{\NURLtTwoMarginZero}{\ensuremath{4.432}}
% NURLtTwoMarginOne [cert, margin] U_cert_v35.json :: R_lt2_margin.1
\newcommand{\NURLtTwoMarginOne}{\ensuremath{6.697}}
% NUThetaTwoEight [cert, upper] U_cert_v35.json :: theta28
\newcommand{\NUThetaTwoEight}{\ensuremath{0.37011}}
% NUDFiveSix [cert, upper] U_cert_v35.json :: D56
\newcommand{\NUDFiveSix}{\ensuremath{1.0368}}
% NUDTwoEight [cert, upper] U_cert_v35.json :: D28
\newcommand{\NUDTwoEight}{\ensuremath{1.5876}}
% NUSThreeThreeFootOverBhatTwo [cert, upper] U_cert_v35.json :: S33_foot_over_Bhat2
\newcommand{\NUSThreeThreeFootOverBhatTwo}{\ensuremath{1.69\cdot10^{-3}}}
% NSAHJWKappaJ [cert, upper] A_interface.json :: HJW.kappa_J
\newcommand{\NSAHJWKappaJ}{\ensuremath{0.042351}}
% NSAGpZero [cert, lower] A_interface.json :: res.tight.gp.0
\newcommand{\NSAGpZero}{\ensuremath{0.24292}}
% NSAGpOne [cert, lower] A_interface.json :: res.tight.gp.1
\newcommand{\NSAGpOne}{\ensuremath{0.24106}}
% NSAGpTwo [cert, lower] A_interface.json :: res.tight.gp.2
\newcommand{\NSAGpTwo}{\ensuremath{0.11961}}
% NSAAminC [cert, lower] A_interface.json :: res.tight.amin_c
\newcommand{\NSAAminC}{\ensuremath{0.97623}}
% NSARhoJ [cert, upper] A_interface.json :: res.tight.rhoJ
\newcommand{\NSARhoJ}{\ensuremath{5.4995\cdot10^{-7}}}
% NSAKappaOTwoTwo [cert, upper] A_interface.json :: res.tight.kappa_o.2.2
\newcommand{\NSAKappaOTwoTwo}{\ensuremath{1.2416\cdot10^{-5}}}
% NSADsrcZero [cert, upper] A_interface.json :: res.tight.Dsrc.0
\newcommand{\NSADsrcZero}{\ensuremath{1.8092\cdot10^{-7}}}
% NSADsrcOne [cert, upper] A_interface.json :: res.tight.Dsrc.1
\newcommand{\NSADsrcOne}{\ensuremath{9.8276\cdot10^{-5}}}
% NSABZeroDWindowSlopeZero [cert, upper] A_interface.json :: res.tight.B0.D_window_slope.0
\newcommand{\NSABZeroDWindowSlopeZero}{\ensuremath{8.7399\cdot10^{-10}}}
% NSABZeroDWindowSlopeOne [cert, upper] A_interface.json :: res.tight.B0.D_window_slope.1
\newcommand{\NSABZeroDWindowSlopeOne}{\ensuremath{5.2599\cdot10^{-8}}}
% NSAFTwoZero [cert, upper] A_interface.json :: res.tight.F2.0
\newcommand{\NSAFTwoZero}{\ensuremath{5.6506\cdot10^{-13}}}
% NSAFTwoOne [cert, upper] A_interface.json :: res.tight.F2.1
\newcommand{\NSAFTwoOne}{\ensuremath{5.9320\cdot10^{-13}}}
% NSACoupOne [cert, upper] A_interface.json :: res.tight.coup.1
\newcommand{\NSACoupOne}{\ensuremath{1.1568\cdot10^{-12}}}
% NSACOneZeroZero [cert, upper] A_interface.json :: res.tight.C1_0.0
\newcommand{\NSACOneZeroZero}{\ensuremath{7.2060\cdot10^{-10}}}
% NSACOneZeroOne [cert, upper] A_interface.json :: res.tight.C1_0.1
\newcommand{\NSACOneZeroOne}{\ensuremath{4.3143\cdot10^{-8}}}
% NSABZeroSupCoreZero [cert, upper] A_interface.json :: res.tight.B0.sup_core.0
\newcommand{\NSABZeroSupCoreZero}{\ensuremath{4.23\cdot10^{-5}}}
% NSACoreFluxMarginTwoFLUXZeroOverKZero [cert, margin] A_interface.json :: res.tight.core.flux_margin.2.FLUX0_over_K0
\newcommand{\NSACoreFluxMarginTwoFLUXZeroOverKZero}{\ensuremath{3.396}}
% NSACoreFluxMarginTwoFLUXTwoOverKTwo [cert, margin] A_interface.json :: res.tight.core.flux_margin.2.FLUX2_over_K2
\newcommand{\NSACoreFluxMarginTwoFLUXTwoOverKTwo}{\ensuremath{6.012}}
% NSAMarginFluxZero [cert, margin] asm_design_N2ref4.json :: res.tight.flux.beta.2.margin_FLUX0
\newcommand{\NSAMarginFluxZero}{\ensuremath{3.365}}
% NSAMarginFluxTwo [cert, margin] asm_design_N2ref4.json :: res.tight.flux.beta.2.margin_FLUX2
\newcommand{\NSAMarginFluxTwo}{\ensuremath{5.955}}
% NSAfourSZero [cert, upper] asm_design_N2ref4.json :: res.tight.S0
\newcommand{\NSAfourSZero}{\ensuremath{7.1878\cdot10^{-3}}}
% NSAfourSOne [cert, upper] asm_design_N2ref4.json :: res.tight.S1
\newcommand{\NSAfourSOne}{\ensuremath{0.10430}}
% NSAMarginCoreSZeroV [cert, margin] asm_design_N2ref4.json :: res.tight.margins.S0v_vs_0.0153
\newcommand{\NSAMarginCoreSZeroV}{\ensuremath{2.116}}
% NSAMarginCoreSOneV [cert, margin] asm_design_N2ref4.json :: res.tight.margins.S1v_vs_0.6
\newcommand{\NSAMarginCoreSOneV}{\ensuremath{5.748}}
% NSAMarginCoreMOne [cert, margin] asm_design_N2ref4.json :: res.tight.margins.M1core
\newcommand{\NSAMarginCoreMOne}{\ensuremath{4.409}}
% NSAMarginCoreMZero [cert, margin] asm_design_N2ref4.json :: res.tight.margins.M0core
\newcommand{\NSAMarginCoreMZero}{\ensuremath{5.450}}
% NSAMarginGhat [cert, margin] asm_design_N2ref4.json :: res.tight.margins.Ghat22
\newcommand{\NSAMarginGhat}{\ensuremath{1.917}}
% NSAMarginRampMOne [cert, margin] asm_design_N2ref4.json :: res.tight.margins.ramp M1
\newcommand{\NSAMarginRampMOne}{\ensuremath{1.184}}
% NSAfourInfoZeroDWindowConst [cert, upper] asm_design_N2ref4.json :: res.tight.info.0.D_window_const
\newcommand{\NSAfourInfoZeroDWindowConst}{\ensuremath{1.7732\cdot10^{-7}}}
% NSAfourInfoOneDWindowConst [cert, upper] asm_design_N2ref4.json :: res.tight.info.1.D_window_const
\newcommand{\NSAfourInfoOneDWindowConst}{\ensuremath{9.8058\cdot10^{-5}}}
% NZDeltaRZero [cert, upper] zone3_interface.json :: DeltaR.0
\newcommand{\NZDeltaRZero}{\ensuremath{4.4986\cdot10^{-8}}}
% NZDeltaROne [cert, upper] zone3_interface.json :: DeltaR.1
\newcommand{\NZDeltaROne}{\ensuremath{6.7972\cdot10^{-6}}}
% NZDeltaRTwo [cert, upper] zone3_interface.json :: DeltaR.2
\newcommand{\NZDeltaRTwo}{\ensuremath{0.011480}}
% NZDeltaRBZero [cert, upper] zone3_interface.json :: DeltaR_b.0
\newcommand{\NZDeltaRBZero}{\ensuremath{2.66\cdot10^{-12}}}
% NZDeltaRBOne [cert, upper] zone3_interface.json :: DeltaR_b.1
\newcommand{\NZDeltaRBOne}{\ensuremath{3.26\cdot10^{-9}}}
% NZDeltaRBTwo [cert, upper] zone3_interface.json :: DeltaR_b.2
\newcommand{\NZDeltaRBTwo}{\ensuremath{4.35\cdot10^{-8}}}
% NZDeltaRDetailNSqZero [cert, upper] zone3_interface.json :: DeltaR_detail.N_sq.0
\newcommand{\NZDeltaRDetailNSqZero}{\ensuremath{8.28\cdot10^{-5}}}
% NZDeltaRDetailNSqOne [cert, upper] zone3_interface.json :: DeltaR_detail.N_sq.1
\newcommand{\NZDeltaRDetailNSqOne}{\ensuremath{0.10149}}
% NZDeltaRDetailNSqTwo [cert, upper] zone3_interface.json :: DeltaR_detail.N_sq.2
\newcommand{\NZDeltaRDetailNSqTwo}{\ensuremath{1.3567}}
% NZDeltaRDetailNSqThree [cert, upper] zone3_interface.json :: DeltaR_detail.N_sq.3
\newcommand{\NZDeltaRDetailNSqThree}{\ensuremath{25350}}
% NZDeltaRDetailPartsTwoPairsThree [cert, upper] zone3_interface.json :: DeltaR_detail.parts.2.pairs.3
\newcommand{\NZDeltaRDetailPartsTwoPairsThree}{\ensuremath{0.011298}}
% NZDeltaRDetailPartsTwoLFourFoot [cert, upper] zone3_interface.json :: DeltaR_detail.parts.2.L4_foot
\newcommand{\NZDeltaRDetailPartsTwoLFourFoot}{\ensuremath{1.6\cdot10^{-4}}}
% NZDeltaRDetailPartsTwoCubic [cert, upper] zone3_interface.json :: DeltaR_detail.parts.2.cubic
\newcommand{\NZDeltaRDetailPartsTwoCubic}{\ensuremath{9.3\cdot10^{-6}}}
% NZDeltaRDetailPartsTwoSThreeThreeFoot [cert, upper] zone3_interface.json :: DeltaR_detail.parts.2.S33_foot
\newcommand{\NZDeltaRDetailPartsTwoSThreeThreeFoot}{\ensuremath{3.7\cdot10^{-7}}}
% NZNIIIZero [cert, upper] zone3_interface.json :: n_III.0
\newcommand{\NZNIIIZero}{\ensuremath{0.071445}}
% NZNIIIOne [cert, upper] zone3_interface.json :: n_III.1
\newcommand{\NZNIIIOne}{\ensuremath{0.32707}}
% NZNIIITwo [cert, upper] zone3_interface.json :: n_III.2
\newcommand{\NZNIIITwo}{\ensuremath{3.0445}}
% NZTauIII [cert, upper] zone3_interface.json :: tau_III
\newcommand{\NZTauIII}{\ensuremath{4.0708}}
% NZoneIIIBoxZero [cert, upper] zone3_interface.json :: box_III_380.b0
\newcommand{\NZoneIIIBoxZero}{\ensuremath{0.10032}}
% NZoneIIIBoxOne [cert, upper] zone3_interface.json :: box_III_380.b1
\newcommand{\NZoneIIIBoxOne}{\ensuremath{0.41974}}
% NZoneIIIBoxZ [cert, upper] zone3_interface.json :: box_III_380.bz
\newcommand{\NZoneIIIBoxZ}{\ensuremath{0.014291}}
% NZBoxIIIThreeEightZeroMarginsBZero [cert, margin] zone3_interface.json :: box_III_380.margins.b0
\newcommand{\NZBoxIIIThreeEightZeroMarginsBZero}{\ensuremath{1.495}}
% NZoneIIIMarginThinnest [cert, margin] zone3_interface.json :: box_III_380.margins.b1
\newcommand{\NZoneIIIMarginThinnest}{\ensuremath{1.238}}
% NZBoxIIIThreeEightZeroMarginsBz [cert, margin] zone3_interface.json :: box_III_380.margins.bz
\newcommand{\NZBoxIIIThreeEightZeroMarginsBz}{\ensuremath{11.19}}
% NZBoxIIISevenZeroZeroBZero [cert, upper] zone3_interface.json :: box_III_700.b0
\newcommand{\NZBoxIIISevenZeroZeroBZero}{\ensuremath{0.22757}}
% NZBoxIIISevenZeroZeroBOne [cert, upper] zone3_interface.json :: box_III_700.b1
\newcommand{\NZBoxIIISevenZeroZeroBOne}{\ensuremath{0.54234}}
% NZUfineKTwoAtBhatRef [cert, upper] zone3_interface.json :: ufine.K2_at_Bhat_ref
\newcommand{\NZUfineKTwoAtBhatRef}{\ensuremath{0.046956}}
% NZUfineKThreeStripAtBhatRef [cert, upper] zone3_interface.json :: ufine.K3_strip_at_Bhat_ref
\newcommand{\NZUfineKThreeStripAtBhatRef}{\ensuremath{0.054140}}
% NRecoveryMinMargin [cert, margin] zone3_interface.json :: cone_family.H6.at_Bhat_ref.min_margin
\newcommand{\NRecoveryMinMargin}{\ensuremath{1.860}}
% NZConeFamilyHSixAtBhatRefConeSlopeMargin [cert, margin] zone3_interface.json :: cone_family.H6.at_Bhat_ref.cone.slope.margin
\newcommand{\NZConeFamilyHSixAtBhatRefConeSlopeMargin}{\ensuremath{2.46}}
% NZConeFamilyHSixAtBhatRefConeKTwoMargin [cert, margin] zone3_interface.json :: cone_family.H6.at_Bhat_ref.cone.K2.margin
\newcommand{\NZConeFamilyHSixAtBhatRefConeKTwoMargin}{\ensuremath{2.54}}
% NZConeFamilyHSixAtBhatRefConeKThreeMargin [cert, margin] zone3_interface.json :: cone_family.H6.at_Bhat_ref.cone.K3.margin
\newcommand{\NZConeFamilyHSixAtBhatRefConeKThreeMargin}{\ensuremath{2.51}}
% NConeTstarMargin [cert, margin] zone3_interface.json :: acceptance.tstar_margin
\newcommand{\NConeTstarMargin}{\ensuremath{2.30}}
% NUoneTFiveYaddLevelOneTwo [cert, upper] u1_ghat.json :: t5_Yadd.level1.2
\newcommand{\NUoneTFiveYaddLevelOneTwo}{\ensuremath{0.0044}}
% NUoneTFiveYaddLevelTwoTwo [cert, upper] u1_ghat.json :: t5_Yadd.level2.2
\newcommand{\NUoneTFiveYaddLevelTwoTwo}{\ensuremath{0.1081}}
% NUoneTFiveYaddLevelTwoThree [cert, upper] u1_ghat.json :: t5_Yadd.level2.3
\newcommand{\NUoneTFiveYaddLevelTwoThree}{\ensuremath{0.374}}
% NUoneWeightThetaGammaZero [cert, upper] u1_ghat.json :: weight.Theta.Gamma.0
\newcommand{\NUoneWeightThetaGammaZero}{\ensuremath{4.00000000000000000000054}}
% NRKZeroZero [cert, upper] ramp_U.json :: K.0.0
\newcommand{\NRKZeroZero}{\ensuremath{13436.03}}
% NRKOneOne [cert, upper] ramp_U.json :: K.1.1
\newcommand{\NRKOneOne}{\ensuremath{10526.19}}
% NRKTwoTwo [cert, upper] ramp_U.json :: K.2.2
\newcommand{\NRKTwoTwo}{\ensuremath{39482.30}}
% NUoneCasesPrimaryYNextZero [cert, upper] u1_ghat.json :: cases.primary.Y_next.0
\newcommand{\NUoneCasesPrimaryYNextZero}{\ensuremath{457.092}}
% NUoneCasesPrimaryYNextOne [cert, upper] u1_ghat.json :: cases.primary.Y_next.1
\newcommand{\NUoneCasesPrimaryYNextOne}{\ensuremath{356.940}}
% NUoneCasesPrimaryYNextTwo [cert, upper] u1_ghat.json :: cases.primary.Y_next.2
\newcommand{\NUoneCasesPrimaryYNextTwo}{\ensuremath{1320.220}}
% NSfourResultsMNineHTwoTightenedFUOneRowsZeroBound [cert, upper] FINAL_manifest_v33.json :: results.M9_H2_tightened_fU1.rows.0.bound
\newcommand{\NSfourResultsMNineHTwoTightenedFUOneRowsZeroBound}{\ensuremath{3.658\cdot10^{-10}}}
% NSfourResultsMNineHTwoTightenedFUOneRowsZeroMargin [cert, margin] FINAL_manifest_v33.json :: results.M9_H2_tightened_fU1.rows.0.margin
\newcommand{\NSfourResultsMNineHTwoTightenedFUOneRowsZeroMargin}{\ensuremath{27.34}}
% NSfourResultsMNineHTwoTightenedFUOneRowsOneBound [cert, upper] FINAL_manifest_v33.json :: results.M9_H2_tightened_fU1.rows.1.bound
\newcommand{\NSfourResultsMNineHTwoTightenedFUOneRowsOneBound}{\ensuremath{1.777\cdot10^{-9}}}
% NSfourResultsMNineHTwoTightenedFUOneRowsOneMargin [cert, margin] FINAL_manifest_v33.json :: results.M9_H2_tightened_fU1.rows.1.margin
\newcommand{\NSfourResultsMNineHTwoTightenedFUOneRowsOneMargin}{\ensuremath{562.8}}
% NSfourResultsMNineHTwoTightenedFUOneRowsTwoBound [cert, upper] FINAL_manifest_v33.json :: results.M9_H2_tightened_fU1.rows.2.bound
\newcommand{\NSfourResultsMNineHTwoTightenedFUOneRowsTwoBound}{\ensuremath{1.665\cdot10^{-6}}}
% NSfourResultsMNineHTwoTightenedFUOneRowsTwoMargin [cert, margin] FINAL_manifest_v33.json :: results.M9_H2_tightened_fU1.rows.2.margin
\newcommand{\NSfourResultsMNineHTwoTightenedFUOneRowsTwoMargin}{\ensuremath{600.9}}
% NHOneAlphaThr [hand, exact] hand :: 
\newcommand{\NHOneAlphaThr}{\ensuremath{1.5\cdot10^{-4}}}
% NHOneSqrtLamThr [hand, exact] hand :: 
\newcommand{\NHOneSqrtLamThr}{\ensuremath{5\cdot10^{-8}}}
% NYbarStep [hand, exact] hand :: 
\newcommand{\NYbarStep}{\ensuremath{1\cdot10^{-12}}}
% NHTwoThrZero [hand, exact] hand :: 
\newcommand{\NHTwoThrZero}{\ensuremath{1\cdot10^{-8}}}
% NHTwoThrOne [hand, exact] hand :: 
\newcommand{\NHTwoThrOne}{\ensuremath{1\cdot10^{-6}}}
% NHTwoThrTwo [hand, exact] hand :: 
\newcommand{\NHTwoThrTwo}{\ensuremath{1\cdot10^{-3}}}
% NHThreeThrZero [hand, exact] hand :: 
\newcommand{\NHThreeThrZero}{\ensuremath{1\cdot10^{-6}}}
% NHThreeThrOne [hand, exact] hand :: 
\newcommand{\NHThreeThrOne}{\ensuremath{5\cdot10^{-6}}}
% NHThreeThrTwo [hand, exact] hand :: 
\newcommand{\NHThreeThrTwo}{\ensuremath{1\cdot10^{-5}}}
% NHFourThrZero [hand, exact] hand :: 
\newcommand{\NHFourThrZero}{\ensuremath{2\cdot10^{-7}}}
% NHFourThrOne [hand, exact] hand :: 
\newcommand{\NHFourThrOne}{\ensuremath{1\cdot10^{-5}}}
% NHFourThrTwo [hand, exact] hand :: 
\newcommand{\NHFourThrTwo}{\ensuremath{1\cdot10^{-5}}}
% NHFourThrThree [hand, exact] hand :: 
\newcommand{\NHFourThrThree}{\ensuremath{1\cdot10^{-5}}}
% NHFourKfourThr [hand, exact] hand :: 
\newcommand{\NHFourKfourThr}{\ensuremath{0.06}}
% NHSixConeSlopeThr [hand, exact] hand :: 
\newcommand{\NHSixConeSlopeThr}{\ensuremath{0.08}}
% NHSixConeKtwoThr [hand, exact] hand :: 
\newcommand{\NHSixConeKtwoThr}{\ensuremath{0.03}}
% NHSixConeNthreeThr [hand, exact] hand :: 
\newcommand{\NHSixConeNthreeThr}{\ensuremath{0.004}}
% NHSixConeJetThr [hand, exact] hand :: 
\newcommand{\NHSixConeJetThr}{\ensuremath{0.12}}
% NHSixConeHFiveThr [hand, exact] hand :: 
\newcommand{\NHSixConeHFiveThr}{\ensuremath{0.0144}}
% NHSixFarHFiveThr [hand, exact] hand :: 
\newcommand{\NHSixFarHFiveThr}{\ensuremath{0.0225}}
% NLemmaOConeM [hand, exact] hand :: 
\newcommand{\NLemmaOConeM}{\ensuremath{0.12}}
% NLemmaOConeK [hand, exact] hand :: 
\newcommand{\NLemmaOConeK}{\ensuremath{0.18}}
% NLemmaOFarM [hand, exact] hand :: 
\newcommand{\NLemmaOFarM}{\ensuremath{0.15}}
% NLemmaOFarK [hand, exact] hand :: 
\newcommand{\NLemmaOFarK}{\ensuremath{0.3}}
% NCoercivityConst [hand, exact] hand :: 
\newcommand{\NCoercivityConst}{\ensuremath{0.046875}}
% NMOnez [hand, exact] hand :: 
\newcommand{\NMOnez}{\ensuremath{0.534}}
% NBetaA [hand, exact] hand :: 
\newcommand{\NBetaA}{\ensuremath{2}}
% NTOneLower [cert, lower] t1clock.json :: t1.0
\newcommand{\NTOneLower}{\ensuremath{4.5978574775}}
% NTOneUpper [cert, upper] t1clock.json :: t1.1
\newcommand{\NTOneUpper}{\ensuremath{4.5978586144}}
% NLambdaOneLower [cert, lower] t1clock.json :: lnlam.0
\newcommand{\NLambdaOneLower}{\ensuremath{6.49\cdot10^{-5}}}
% NLambdaOneUpper [cert, upper] t1clock.json :: lnlam.1
\newcommand{\NLambdaOneUpper}{\ensuremath{6.75\cdot10^{-5}}}
% NTMinusTOneUpper [cert, upper] t1clock.json :: lnlam.1; 1.5e-4.sup_ab
\newcommand{\NTMinusTOneUpper}{\ensuremath{2.1\cdot10^{-17}}}
% NAbarUpper [cert, upper] final_alpha_points_5e-8.json :: 1.5e-4.sup_ab
\newcommand{\NAbarUpper}{\ensuremath{2.19103\cdot10^{-7}}}
% NBetabarUpper [cert, upper] final_alpha_points_5e-8.json :: 1.5e-4.sup_beta
\newcommand{\NBetabarUpper}{\ensuremath{3.33811\cdot10^{-9}}}
% NBaseKmaxLower [cert, lower] t1clock.json :: base_kmax.0
\newcommand{\NBaseKmaxLower}{\ensuremath{2.279\cdot10^{4}}}
% NBaseKmaxUpper [cert, upper] t1clock.json :: base_kmax.1
\newcommand{\NBaseKmaxUpper}{\ensuremath{2.376\cdot10^{4}}}
% NHOneAlphaBase [cert, upper] summary.json :: alpha_abs_max
\newcommand{\NHOneAlphaBase}{\ensuremath{1.1739\cdot10^{-4}}}
% NHOneAlphaErr [cert, upper] t9_h1_clam_s180.json :: C_alpha; inputs.YSTAR
\newcommand{\NHOneAlphaErr}{\ensuremath{7.21\cdot10^{-13}}}
% NHOneAlphaMargin [cert, margin] summary.json :: alpha_abs_max; C_alpha; inputs.YSTAR
\newcommand{\NHOneAlphaMargin}{\ensuremath{1.27}}
% NHOneSqrtLamBase [cert, upper] t9_h1_clam_s180.json :: sqrtLam_base_Wp
\newcommand{\NHOneSqrtLamBase}{\ensuremath{2.282\cdot10^{-8}}}
% NHOneSqrtLamTotal [cert, upper] t9_h1_clam_s180.json :: sqrtLam_total_Wp
\newcommand{\NHOneSqrtLamTotal}{\ensuremath{2.29\cdot10^{-8}}}
% NHOneSqrtLamMargin [cert, margin] t9_h1_clam_s180.json :: sqrtLam_total_Wp
\newcommand{\NHOneSqrtLamMargin}{\ensuremath{2.18}}
% NHOneYbarLimit [cert, lower] t9_h1_clam_s180.json :: Ybar_limit_H1
\newcommand{\NHOneYbarLimit}{\ensuremath{1.12\cdot10^{-9}}}
% NSwitchYUpper [cert, upper] switch_s2.json :: modulated.Y_raw
\newcommand{\NSwitchYUpper}{\ensuremath{6.21\cdot10^{-13}}}
% NSwitchYMargin [cert, margin] switch_s2.json :: modulated.Y_raw
\newcommand{\NSwitchYMargin}{\ensuremath{1.6}}
% NMinCI [cert, lower] constants.json :: meta.summary.min_c_I
\newcommand{\NMinCI}{\ensuremath{0.05283}}
% NNumCertificates [cert, exact] constants.json :: meta.summary.n_cert
\newcommand{\NNumCertificates}{\ensuremath{26}}
% NMinRhoC [cert, lower] constants.json :: meta.closure.rho_c_min
\newcommand{\NMinRhoC}{\ensuremath{0.04711}}
% NHTwoBaseZero [cert, upper] FINAL_manifest_v33.json :: results.M9_H2_tightened.rows.0.base
\newcommand{\NHTwoBaseZero}{\ensuremath{8.011\cdot10^{-11}}}
% NHTwoBaseOne [cert, upper] FINAL_manifest_v33.json :: results.M9_H2_tightened.rows.1.base
\newcommand{\NHTwoBaseOne}{\ensuremath{1.336\cdot10^{-10}}}
% NHTwoBaseTwo [cert, upper] FINAL_manifest_v33.json :: results.M9_H2_tightened.rows.2.base
\newcommand{\NHTwoBaseTwo}{\ensuremath{2.936\cdot10^{-10}}}
% NHTwoErrWZero [cert, upper] FINAL_manifest_v33.json :: results.M9_H2_tightened.rows.0.with_error
\newcommand{\NHTwoErrWZero}{\ensuremath{8.014\cdot10^{-11}}}
% NHTwoMarginWZero [cert, margin] FINAL_manifest_v33.json :: results.M9_H2_tightened.rows.0.margin
\newcommand{\NHTwoMarginWZero}{\ensuremath{124.78}}
% NHTwoErrWOne [cert, upper] FINAL_manifest_v33.json :: results.M9_H2_tightened.rows.1.with_error
\newcommand{\NHTwoErrWOne}{\ensuremath{4.706\cdot10^{-10}}}
% NHTwoMarginWOne [cert, margin] FINAL_manifest_v33.json :: results.M9_H2_tightened.rows.1.margin
\newcommand{\NHTwoMarginWOne}{\ensuremath{2124.9}}
% NHTwoErrWTwo [cert, upper] FINAL_manifest_v33.json :: results.M9_H2_tightened.rows.2.with_error
\newcommand{\NHTwoErrWTwo}{\ensuremath{1.559\cdot10^{-6}}}
% NHTwoMarginWTwo [cert, margin] FINAL_manifest_v33.json :: results.M9_H2_tightened.rows.2.margin
\newcommand{\NHTwoMarginWTwo}{\ensuremath{641.53}}
% NHThreeBaseZero [cert, upper] FINAL_manifest_v33.json :: results.M9_H2f.rows.0.base
\newcommand{\NHThreeBaseZero}{\ensuremath{3.127\cdot10^{-7}}}
% NHThreeBaseOne [cert, upper] FINAL_manifest_v33.json :: results.M9_H2f.rows.1.base
\newcommand{\NHThreeBaseOne}{\ensuremath{1.004\cdot10^{-6}}}
% NHThreeBaseTwo [cert, upper] FINAL_manifest_v33.json :: results.M9_H2f.rows.2.base
\newcommand{\NHThreeBaseTwo}{\ensuremath{2.431\cdot10^{-6}}}
% NHThreeErrZero [cert, upper] FINAL_manifest_v33.json :: results.M9_H2f.rows.0.with_error
\newcommand{\NHThreeErrZero}{\ensuremath{3.1363\cdot10^{-7}}}
% NHThreeMarginZero [cert, margin] FINAL_manifest_v33.json :: results.M9_H2f.rows.0.margin
\newcommand{\NHThreeMarginZero}{\ensuremath{3.1885}}
% NHThreeErrOne [cert, upper] FINAL_manifest_v33.json :: results.M9_H2f.rows.1.with_error
\newcommand{\NHThreeErrOne}{\ensuremath{1.0034\cdot10^{-6}}}
% NHThreeMarginOne [cert, margin] FINAL_manifest_v33.json :: results.M9_H2f.rows.1.margin
\newcommand{\NHThreeMarginOne}{\ensuremath{4.9834}}
% NHThreeErrTwo [cert, upper] FINAL_manifest_v33.json :: results.M9_H2f.rows.2.with_error
\newcommand{\NHThreeErrTwo}{\ensuremath{2.4351\cdot10^{-6}}}
% NHThreeMarginTwo [cert, margin] FINAL_manifest_v33.json :: results.M9_H2f.rows.2.margin
\newcommand{\NHThreeMarginTwo}{\ensuremath{4.1067}}
% NHFourBaseZero [cert, upper] FINAL_manifest_v33.json :: results.M9_base_sums.combined.windows_max.0
\newcommand{\NHFourBaseZero}{\ensuremath{3.48\cdot10^{-8}}}
% NHFourErrZero [cert, upper] FINAL_manifest_v33.json :: results.M9_H3_H2v32.rows.win.0.with_error_arb
\newcommand{\NHFourErrZero}{\ensuremath{3.5087\cdot10^{-8}}}
% NHFourMarginZero [cert, margin] FINAL_manifest_v33.json :: results.M9_H3_H2v32.rows.win.0.margin_arb
\newcommand{\NHFourMarginZero}{\ensuremath{5.700}}
% NHFourBaseOne [cert, upper] FINAL_manifest_v33.json :: results.M9_base_sums.combined.windows_max.1
\newcommand{\NHFourBaseOne}{\ensuremath{1.27\cdot10^{-7}}}
% NHFourErrOne [cert, upper] FINAL_manifest_v33.json :: results.M9_H3_H2v32.rows.win.1.with_error_arb
\newcommand{\NHFourErrOne}{\ensuremath{1.2616\cdot10^{-7}}}
% NHFourMarginOne [cert, margin] FINAL_manifest_v33.json :: results.M9_H3_H2v32.rows.win.1.margin_arb
\newcommand{\NHFourMarginOne}{\ensuremath{79.26}}
% NHFourBaseTwo [cert, upper] FINAL_manifest_v33.json :: results.M9_base_sums.combined.windows_max.2
\newcommand{\NHFourBaseTwo}{\ensuremath{3.76\cdot10^{-7}}}
% NHFourErrTwo [cert, upper] FINAL_manifest_v33.json :: results.M9_H3_H2v32.rows.win.2.with_error_arb
\newcommand{\NHFourErrTwo}{\ensuremath{3.7724\cdot10^{-7}}}
% NHFourMarginTwo [cert, margin] FINAL_manifest_v33.json :: results.M9_H3_H2v32.rows.win.2.margin_arb
\newcommand{\NHFourMarginTwo}{\ensuremath{26.50}}
% NHFourBaseThree [cert, upper] FINAL_manifest_v33.json :: results.M9_base_sums.combined.windows_max.3
\newcommand{\NHFourBaseThree}{\ensuremath{9.71\cdot10^{-7}}}
% NHFourErrThree [cert, upper] FINAL_manifest_v33.json :: results.M9_H3_H2v32.rows.win.3.with_error_arb
\newcommand{\NHFourErrThree}{\ensuremath{1.2822\cdot10^{-6}}}
% NHFourMarginThree [cert, margin] FINAL_manifest_v33.json :: results.M9_H3_H2v32.rows.win.3.margin_arb
\newcommand{\NHFourMarginThree}{\ensuremath{7.799}}
% NHFourYbarLimit [cert, lower] FINAL_manifest_v33.json :: results.M9_H3_H2v32.binding_Ybar_limit
\newcommand{\NHFourYbarLimit}{\ensuremath{1.575\cdot10^{-11}}}
% NHFourYbarMargin [cert, margin] FINAL_manifest_v33.json :: results.M9_H3_H2v32.binding_Ybar_limit
\newcommand{\NHFourYbarMargin}{\ensuremath{15.75}}
% NHFourKfourBase [cert, upper] amp_A1_w180.json :: inputs.BASE_T1.K4
\newcommand{\NHFourKfourBase}{\ensuremath{3.1633\cdot10^{-4}}}
% NHSixConeSlope [cert, upper] FINAL_manifest_v33.json :: results.cone_fixedL.checks.slope.value
\newcommand{\NHSixConeSlope}{\ensuremath{0.03029}}
% NHSixConeSlopeMargin [cert, margin] FINAL_manifest_v33.json :: results.cone_fixedL.checks.slope.margin
\newcommand{\NHSixConeSlopeMargin}{\ensuremath{2.641}}
% NHSixConeKtwo [cert, upper] FINAL_manifest_v33.json :: results.cone_fixedL.checks.K2.value
\newcommand{\NHSixConeKtwo}{\ensuremath{0.009876}}
% NHSixConeKtwoMargin [cert, margin] FINAL_manifest_v33.json :: results.cone_fixedL.checks.K2.margin
\newcommand{\NHSixConeKtwoMargin}{\ensuremath{3.037}}
% NHSixConeNthree [cert, upper] FINAL_manifest_v33.json :: results.cone_fixedL.checks.N3.value
\newcommand{\NHSixConeNthree}{\ensuremath{1.461\cdot10^{-4}}}
% NHSixConeNthreeMargin [cert, margin] FINAL_manifest_v33.json :: results.cone_fixedL.checks.N3.margin
\newcommand{\NHSixConeNthreeMargin}{\ensuremath{27.38}}
% NHSixConeHFive [cert, upper] FINAL_manifest_v33.json :: results.H4_fixedL.total
\newcommand{\NHSixConeHFive}{\ensuremath{1.2865\cdot10^{-5}}}
% NHSixConeHFiveMargin [cert, margin] FINAL_manifest_v33.json :: results.H4_fixedL.margin
\newcommand{\NHSixConeHFiveMargin}{\ensuremath{1119}}
% NHSixConeTstarNeeded [cert, upper] FINAL_manifest_v33.json :: results.H7_fixedL.ratio_upper
\newcommand{\NHSixConeTstarNeeded}{\ensuremath{1.6000015\cdot10^{-7}}}
% NHSixConeTstarLower [cert, lower] FINAL_manifest_v33.json :: results.H7_fixedL.t_star_lower
\newcommand{\NHSixConeTstarLower}{\ensuremath{3.68334\cdot10^{-7}}}
% NHSixConeTstarMargin [cert, margin] FINAL_manifest_v33.json :: results.H7_fixedL.margin_tstar_over_ratio
\newcommand{\NHSixConeTstarMargin}{\ensuremath{2.302}}
% NHSixConeHFiveRatioPeak [cert, upper] h4loc_A1_w180.json :: windows.constants/amp_A1_w180.json.worst.ratio
\newcommand{\NHSixConeHFiveRatioPeak}{\ensuremath{0.2393}}
% NHSixConeHOneSup [cert, upper] FINAL_manifest_v33.json :: results.H1_fixedL.sup_h0_max
\newcommand{\NHSixConeHOneSup}{\ensuremath{0.2163}}
% NHSixConeHOneH [cert, upper] FINAL_manifest_v33.json :: results.H1_fixedL.H_max
\newcommand{\NHSixConeHOneH}{\ensuremath{0.2185}}
% NHSixFarHFiveRatio [cert, upper] h4loc_v32.json :: windows.constants/amp_A1_w180.json.cylinders.18.ratio
\newcommand{\NHSixFarHFiveRatio}{\ensuremath{0.1761}}
% NHSixFarHFiveMarginOld [cert, margin] h4loc_v32.json :: windows.constants/amp_A1_w180.json.cylinders.18.ratio
\newcommand{\NHSixFarHFiveMarginOld}{\ensuremath{5.68}}
% NHSixFarHFiveMargin [cert, margin] FINAL_manifest_v33.json :: results.M9_H4_far.min_margin
\newcommand{\NHSixFarHFiveMargin}{\ensuremath{15.58}}
% NHSixConeHFiveMarginOld [cert, margin] h4loc_A1_w180.json :: windows.constants/amp_A1_w180.json.worst.ratio
\newcommand{\NHSixConeHFiveMarginOld}{\ensuremath{4.18}}
% NHSixFarKtwo [cert, upper] t9_far_v32.json :: h2
\newcommand{\NHSixFarKtwo}{\ensuremath{1.797\cdot10^{-4}}}
% NHSixFarKtwoThr [hand, exact] t9_far_v32.json :: thr.h2
\newcommand{\NHSixFarKtwoThr}{\ensuremath{3.53\cdot10^{-4}}}
% NHSixFarKtwoMargin [cert, margin] t9_far_v32.json :: margins.h2
\newcommand{\NHSixFarKtwoMargin}{\ensuremath{1.964}}
% NHSixFarKthree [cert, upper] t9_far_v32.json :: h3
\newcommand{\NHSixFarKthree}{\ensuremath{2.595\cdot10^{-8}}}
% NHSixFarKthreeThr [hand, exact] t9_far_v32.json :: thr.h3
\newcommand{\NHSixFarKthreeThr}{\ensuremath{1.17\cdot10^{-5}}}
% NHSixFarKthreeMargin [cert, margin] t9_far_v32.json :: margins.h3
\newcommand{\NHSixFarKthreeMargin}{\ensuremath{450.9}}
% NHSixFarKfour [cert, upper] t9_far_v32.json :: h4
\newcommand{\NHSixFarKfour}{\ensuremath{6.964\cdot10^{-9}}}
% NHSixFarKfourThr [hand, exact] t9_far_v32.json :: thr.h4
\newcommand{\NHSixFarKfourThr}{\ensuremath{1.03\cdot10^{-7}}}
% NHSixFarKfourMargin [cert, margin] t9_far_v32.json :: margins.h4
\newcommand{\NHSixFarKfourMargin}{\ensuremath{14.79}}
% NHSixFarSlope [cert, upper] t9_far_v32.json :: slope
\newcommand{\NHSixFarSlope}{\ensuremath{0.04153}}
% NHSixFarSlopeThr [hand, exact] t9_far_v32.json :: thr.slope
\newcommand{\NHSixFarSlopeThr}{\ensuremath{0.0800}}
% NHSixFarSlopeMargin [cert, margin] t9_far_v32.json :: margins.slope
\newcommand{\NHSixFarSlopeMargin}{\ensuremath{1.926}}
% NHSixFarNthree [cert, upper] t9_far_v32.json :: L2
\newcommand{\NHSixFarNthree}{\ensuremath{5.516\cdot10^{-7}}}
% NHSixFarNthreeThr [hand, exact] t9_far_v32.json :: thr.L2
\newcommand{\NHSixFarNthreeThr}{\ensuremath{3.32\cdot10^{-6}}}
% NHSixFarNthreeMargin [cert, margin] t9_far_v32.json :: margins.L2
\newcommand{\NHSixFarNthreeMargin}{\ensuremath{6.019}}
% NFarTstarMargin [cert, margin] lemmaO_v32_consts.json :: far_v32.case.margin_factor_lower
\newcommand{\NFarTstarMargin}{\ensuremath{15.5}}
% NFarTstarLower [cert, lower] lemmaO_v32_consts.json :: far_v32.case.t_star_lower
\newcommand{\NFarTstarLower}{\ensuremath{9.5108\cdot10^{-8}}}
% NFarTstarNeeded [cert, upper] lemmaO_v32_consts.json :: far_v32.case.needed
\newcommand{\NFarTstarNeeded}{\ensuremath{6.134\cdot10^{-9}}}
% NGammaModes [cert, exact] gamma0_final.json :: M
\newcommand{\NGammaModes}{\ensuremath{3762}}
% NGammaHalfwidth [cert, upper] gamma0_final.json :: b_j_halfwidth
\newcommand{\NGammaHalfwidth}{\ensuremath{8.138\cdot10^{-24}}}
% NGammaContraction [cert, upper] gamma0_final.json :: krawczyk_contraction
\newcommand{\NGammaContraction}{\ensuremath{3.01\cdot10^{-12}}}
% NGammaArcs [cert, exact] gamma0_final.json :: embedded.arcs
\newcommand{\NGammaArcs}{\ensuremath{701}}
% NGammaMaxTurn [cert, upper] gamma0_final.json :: embedded.maxturn
\newcommand{\NGammaMaxTurn}{\ensuremath{0.04997}}
% NGammaMaxPair [cert, upper] gamma0_final.json :: embedded.maxpair
\newcommand{\NGammaMaxPair}{\ensuremath{0.09597}}
% NGammaMinSep [cert, lower] gamma0_final.json :: embedded.minsep
\newcommand{\NGammaMinSep}{\ensuremath{8.89\cdot10^{-3}}}
% NGammaHalfGap [cert, lower] gamma0_final.json :: embedded.gap_cert
\newcommand{\NGammaHalfGap}{\ensuremath{0.0248219}}
% NPhaseISlabs [cert, exact] status.json :: slab
\newcommand{\NPhaseISlabs}{\ensuremath{96}}
% NPhaseIModes [cert, exact] status.json :: last_slab.M
\newcommand{\NPhaseIModes}{\ensuremath{2800}}
% NPhaseIKmaxLower [cert, lower] status.json :: kmax_enclosure.0
\newcommand{\NPhaseIKmaxLower}{\ensuremath{29.352874545}}
% NPhaseIKmaxUpper [cert, upper] status.json :: kmax_enclosure.1
\newcommand{\NPhaseIKmaxUpper}{\ensuremath{29.352874546}}
% NPhaseITLower [cert, lower] status.json :: t
\newcommand{\NPhaseITLower}{\ensuremath{4.597850539485830}}
% NPhaseITUpper [cert, upper] status.json :: t
\newcommand{\NPhaseITUpper}{\ensuremath{4.597850539485831}}
% NPeakSqrtLam [cert, upper] rate_table_fin12f4_a0.00015.json :: sqrtL_max
\newcommand{\NPeakSqrtLam}{\ensuremath{5.308\cdot10^{-8}}}
% NTailRate [cert, lower] rate_table_fin12f4_a0.00015.json :: tail.rho
\newcommand{\NTailRate}{\ensuremath{2.693\cdot10^{-3}}}
% NFloorSqrtLam [cert, upper] rate_table_fin12f4_a0.00015.json :: tail.sqrt_LF
\newcommand{\NFloorSqrtLam}{\ensuremath{9.872\cdot10^{-10}}}
% NAlphaFloor [cert, upper] rate_table_fin12f4_a0.00015.json :: tail.alpha_F
\newcommand{\NAlphaFloor}{\ensuremath{6.867\cdot10^{-7}}}
% NSupBetaVerifier [cert, upper] final_alpha_points_5e-8.json :: 1.5e-4.sup_beta
\newcommand{\NSupBetaVerifier}{\ensuremath{3.33811\cdot10^{-9}}}
% NMOneRamp [cert, upper] final_alpha_points_5e-8.json :: 1.5e-4.M1
\newcommand{\NMOneRamp}{\ensuremath{0.4328}}
% NSITwoPrime [cert, upper] FINAL_manifest_v33.json :: results.F12_sliver.sI2_prime
\newcommand{\NSITwoPrime}{\ensuremath{0.015457}}
% NLoopGainCoarse [cert, upper] U_cert_v35.json :: loop_gain.2.total
\newcommand{\NLoopGainCoarse}{\ensuremath{0.03904}}
% NRateMarginCoarse [cert, margin] U_cert_v35.json :: rate_margin
\newcommand{\NRateMarginCoarse}{\ensuremath{28.2}}
% NSOneEpsLower [cert, lower] bracket2.json :: eps_a
\newcommand{\NSOneEpsLower}{\ensuremath{1.5346198892965\cdot10^{-4}}}
% NSOneEpsUpper [cert, upper] bracket2.json :: eps_b
\newcommand{\NSOneEpsUpper}{\ensuremath{1.5346198899299\cdot10^{-4}}}
% NSOneRowsTip [cert, exact] ext_tip_check.json :: rows
\newcommand{\NSOneRowsTip}{\ensuremath{51246}}
% NSOneTipWidth [cert, upper] ext_tip_check.json :: max_pointwise_width.th
\newcommand{\NSOneTipWidth}{\ensuremath{3.2\cdot10^{-12}}}
% NSOneExtWidth [cert, upper] ext813_check.json :: max_pointwise_width.p
\newcommand{\NSOneExtWidth}{\ensuremath{5.7\cdot10^{-10}}}
% NSOneRowsMC [cert, exact] check_outA5.json :: prof_tip_vs_MC.rows_compared
\newcommand{\NSOneRowsMC}{\ensuremath{1025}}
% NSOneRowsMB [cert, exact] check_outA5.json :: ext813_vs_MB.rows_compared
\newcommand{\NSOneRowsMB}{\ensuremath{501}}
% NSOneMinMargin [cert, margin] check_outA5.json :: min:*.min_relative_margin.*
\newcommand{\NSOneMinMargin}{\ensuremath{1.2\cdot10^{-9}}}
% NTwoSegments [cert, exact] contour_result.json :: len:segments
\newcommand{\NTwoSegments}{\ensuremath{5058}}
% NTwoInteriorBoxes [cert, exact] interior_result.json :: len:
\newcommand{\NTwoInteriorBoxes}{\ensuremath{174}}
% NTwoContourBoxes [cert, exact] eta_cache.json :: len:
\newcommand{\NTwoContourBoxes}{\ensuremath{202}}
% NTwoEtaMax [cert, upper] eta_cache.json :: max:*.eta
\newcommand{\NTwoEtaMax}{\ensuremath{0.1163}}
% NTwoLipMax [cert, upper] eta_cache.json :: max:*.Lip
\newcommand{\NTwoLipMax}{\ensuremath{0.272}}
% NTwoLambdaZero [est, upper] cert_b0.25_th0.0_tau0.03.json :: Lam
\newcommand{\NTwoLambdaZero}{\ensuremath{1.790819}}
% NTwoLambdaOne [est, upper] cert_b0.25_th0.48_tau0.03.json :: Lam
\newcommand{\NTwoLambdaOne}{\ensuremath{1.768282}}
% NSOneDTightLower [doc, lower] doc:README.md :: D ∈ [2.50080627654, 2.50080627668]
\newcommand{\NSOneDTightLower}{\ensuremath{2.50080627654}}
% NSOneDTightUpper [doc, upper] doc:README.md :: D ∈ [2.50080627654, 2.50080627668]
\newcommand{\NSOneDTightUpper}{\ensuremath{2.50080627668}}
% NSOneKTightLower [doc, lower] doc:README.md :: k(0) ∈ [1.53982063762, 1.53982063822]
\newcommand{\NSOneKTightLower}{\ensuremath{1.53982063762}}
% NSOneKTightUpper [doc, upper] doc:README.md :: k(0) ∈ [1.53982063762, 1.53982063822]
\newcommand{\NSOneKTightUpper}{\ensuremath{1.53982063822}}
% NSOneDCondLower [doc, lower] doc:README.md :: D = −p(tip) ∈ [2.500806276619695, 2.500806276620049]
\newcommand{\NSOneDCondLower}{\ensuremath{2.500806276619695}}
% NSOneDCondUpper [doc, upper] doc:README.md :: D = −p(tip) ∈ [2.500806276619695, 2.500806276620049]
\newcommand{\NSOneDCondUpper}{\ensuremath{2.500806276620049}}
% NSOneKCondLower [doc, lower] doc:README.md :: k(0) ∈ [1.539820637871836, 1.5398206378721015]
\newcommand{\NSOneKCondLower}{\ensuremath{1.539820637871836}}
% NSOneKCondUpper [doc, upper] doc:README.md :: k(0) ∈ [1.539820637871836, 1.5398206378721015]
\newcommand{\NSOneKCondUpper}{\ensuremath{1.5398206378721015}}
% NSOneTauTipLower [doc, lower] doc:README.md :: τ_tip ∈ [10.249077651120778, 10.249077651121361]
\newcommand{\NSOneTauTipLower}{\ensuremath{10.249077651120778}}
% NSOneTauTipUpper [doc, upper] doc:README.md :: τ_tip ∈ [10.249077651120778, 10.249077651121361]
\newcommand{\NSOneTauTipUpper}{\ensuremath{10.249077651121361}}
% NSOneStateUncert [doc, upper] doc:README.md :: State uncertainty at σ = 8: 4.4e−13
\newcommand{\NSOneStateUncert}{\ensuremath{4.4\cdot10^{-13}}}
% NFarBoxEprime [doc, margin] doc:STEP4_RESULT.md :: (6.4287, 4.6917, 1.9940, 1.69999995)
\newcommand{\NFarBoxEprime}{\ensuremath{6.42}}
% NFarBoxXEsecond [doc, margin] doc:STEP4_RESULT.md :: (6.4287, 4.6917, 1.9940, 1.69999995)
\newcommand{\NFarBoxXEsecond}{\ensuremath{4.69}}
% NFarBoxXsqEthird [doc, margin] doc:STEP4_RESULT.md :: (6.4287, 4.6917, 1.9940, 1.69999995)
\newcommand{\NFarBoxXsqEthird}{\ensuremath{1.99}}
% NFarBoxE [doc, margin] doc:STEP4_RESULT.md :: (6.4287, 4.6917, 1.9940, 1.69999995)
\newcommand{\NFarBoxE}{\ensuremath{1.69}}
% NHFourKfourErr [doc, upper] doc:STEP4_RESULT.md :: | 5.60e-3 against 0.06 | ×10.7 |
\newcommand{\NHFourKfourErr}{\ensuremath{5.60\cdot10^{-3}}}
% NHFourKfourMargin [doc, margin] doc:STEP4_RESULT.md :: | 5.60e-3 against 0.06 | ×10.7 |
\newcommand{\NHFourKfourMargin}{\ensuremath{10.7}}
% NCoercivityCoefficient [doc, upper] doc:step3_proof_v32.tex :: \norm{e^{\tau\mathcal L}u}_\rho\le14.8\,e^{-0.0213\tau}\norm
\newcommand{\NCoercivityCoefficient}{\ensuremath{14.8}}
% NCoercivityRate [doc, lower] doc:step3_proof_v32.tex :: \norm{e^{\tau\mathcal L}u}_\rho\le14.8\,e^{-0.0213\tau}\norm
\newcommand{\NCoercivityRate}{\ensuremath{0.0213}}
% NRegionRampStart [hand, exact] hand :: 
\newcommand{\NRegionRampStart}{\ensuremath{14}}
% NRegionInterface [hand, exact] hand :: 
\newcommand{\NRegionInterface}{\ensuremath{28}}
% NRegionShellEnd [hand, exact] hand :: 
\newcommand{\NRegionShellEnd}{\ensuremath{60}}
% NRegionSegEnd [hand, exact] hand :: 
\newcommand{\NRegionSegEnd}{\ensuremath{350}}
% NRegionAnnEnd [hand, exact] hand :: 
\newcommand{\NRegionAnnEnd}{\ensuremath{700}}
% NZoneIIILow [hand, exact] hand :: 
\newcommand{\NZoneIIILow}{\ensuremath{253}}
% NZoneIIIHigh [hand, exact] hand :: 
\newcommand{\NZoneIIIHigh}{\ensuremath{379.5}}
% NConeCentreMin [hand, exact] hand :: 
\newcommand{\NConeCentreMin}{\ensuremath{300}}
% NConeCentreMax [hand, exact] hand :: 
\newcommand{\NConeCentreMax}{\ensuremath{380}}
% NFarCentreMin [hand, exact] hand :: 
\newcommand{\NFarCentreMin}{\ensuremath{490}}
% NConeL [hand, exact] hand :: 
\newcommand{\NConeL}{\ensuremath{50}}
% NConeX [hand, exact] hand :: 
\newcommand{\NConeX}{\ensuremath{3}}
% NFarL [hand, exact] hand :: 
\newcommand{\NFarL}{\ensuremath{113}}
% NFarX [hand, exact] hand :: 
\newcommand{\NFarX}{\ensuremath{2}}
 is found.
% ======================================================================

% ---- which document are we in ----------------------------------------
\ifdefined\ifsupplement\else\newif\ifsupplement\fi

% ---- packages ---------------------------------------------------------
\usepackage[T1]{fontenc}
\usepackage[utf8]{inputenc}
\usepackage{lmodern}
\usepackage[margin=1.1in]{geometry}
\usepackage{amsmath,amssymb,amsthm,mathtools,bm}
\usepackage{booktabs,array,longtable,tabularx,multirow}
\usepackage{enumitem}
\usepackage{graphicx}
\usepackage{xcolor}
\usepackage{tikz}
\usetikzlibrary{positioning,arrows.meta,calc}
% xr-hyper must come before hyperref
\usepackage{xr-hyper}
\usepackage[colorlinks=true,linkcolor=blue!50!black,citecolor=green!35!black,urlcolor=blue!50!black]{hyperref}

\raggedbottom
\setlength{\emergencystretch}{2em}
\setlist[enumerate]{leftmargin=*,itemsep=2pt,topsep=3pt}
\setlist[itemize]{leftmargin=*,itemsep=2pt,topsep=3pt}

% ---- numbering --------------------------------------------------------
% Main paper: sections 1,2,...; theorems 2.3 (by section); equations (2.3).
% Supplement: the Parts A..K are the \section's, numbered by letter, so
% theorems are A.1, E.4, ... (do not confuse "Theorem A" of the paper, which
% has no dot, with "Theorem A.1" of Part A).
\ifsupplement
  \renewcommand{\thesection}{\Alph{section}}
\fi
\numberwithin{equation}{section}

% ---- theorem environments (numbered by section, one shared counter) ----
\theoremstyle{plain}
\newtheorem{theorem}{Theorem}[section]
\newtheorem{proposition}[theorem]{Proposition}
\newtheorem{lemma}[theorem]{Lemma}
\newtheorem{corollary}[theorem]{Corollary}
\newtheorem{maintheorem}[theorem]{Main Theorem}
% Theorems A, B, C, D of the paper: own counter, lettered, no section prefix
\newtheorem{thmlet}{Theorem}
\renewcommand{\thethmlet}{\Alph{thmlet}}
\theoremstyle{definition}
\newtheorem{definition}[theorem]{Definition}
\newtheorem{convention}[theorem]{Convention}
\newtheorem{example}[theorem]{Example}
\theoremstyle{remark}
\newtheorem{remark}[theorem]{Remark}

% A list of hypotheses labelled (H1),(H2),... : \begin{hyplist}{H} ... \end{hyplist}
% (also {O}, {L}, {B}, {K}, {W} for the families of section 5.3 of OUTLINE.md).
\newenvironment{hyplist}[1]%
  {\begin{enumerate}[label=\textup{(#1\arabic*)},leftmargin=3em,itemsep=2pt]}%
  {\end{enumerate}}

% ---- trust classes (OUTLINE 4.1, 5.1.3) -------------------------------
% Supplement: explicit tags.  Main paper: certified is the default and
% unmarked; every non-certified number is marked "(floating point)".
\ifsupplement
  \newcommand{\tcert}{\textup{[cert]}}   % arb/outward-rounded, pinned script and key
  \newcommand{\test}{\textup{[est]}}     % floating point or estimate; never in a theorem
  \newcommand{\thand}{\textup{[hand]}}   % proved by hand in the text
  \newcommand{\tlit}{\textup{[lit]}}     % cited from the literature
\else
  \newcommand{\tcert}{}
  \newcommand{\test}{\unskip\ \textup{(floating point)}}
  \newcommand{\thand}{}
  \newcommand{\tlit}{}
\fi

% ---- gaps: visible, counted, an error in a final build -----------------
% Usage: \GAP{one sentence saying what is missing}   (and an entry in ISSUES.md)
\newcounter{gapcount}
\newif\iffinalpaper   % \finalpapertrue makes every \GAP a hard error
\newcommand{\GAP}[1]{\stepcounter{gapcount}\typeout{GAP-MARKER \thegapcount}%
  \iffinalpaper\PackageError{paper}{\string\GAP\space used in a final build}{Remove all gaps.}\fi
  {\color{red}\sffamily\bfseries[GAP~\thegapcount:\ }{\color{red}\sffamily #1}{\color{red}\sffamily\bfseries]}}
\AtEndDocument{\typeout{GAP-MARKER total: \thegapcount}}

% ---- file names and JSON keys (supplement only) -------------------------
% No internal path appears in the main paper; a use there is a warning.
\newcommand{\fn}[1]{%
  \ifsupplement\else\PackageWarning{paper}{\string\fn\space used in the main paper}\fi
  \texttt{\detokenize{#1}}}
\newcommand{\key}[1]{\texttt{\detokenize{#1}}}

% ---- cross references between main paper and supplement ----------------
% Labels carry a unit prefix (section 5.3 of OUTLINE.md and DRAFTER_GUIDE.md):
%   main paper:  intro: obj: s3: s4: s5: s6: s7: s8: s9:   (kind: thm- lem- eq- tab- fig- ...)
%   supplement:  pA: pB: ... pK:
% so all labels are globally unique and the prefixes below are only a second guard.
% \Xref{label}  : reference to the OTHER document (hyperlinked via xr-hyper)
% \Lref{label}  : reference in the SAME document (plain \ref)
\ifsupplement
  \externaldocument[M-]{main}[main.pdf]
  \newcommand{\Xref}[1]{\ref{M-#1}}
\else
  \externaldocument[S-]{supplement}[supplement.pdf]
  \newcommand{\Xref}[1]{\ref{S-#1}}
\fi
\newcommand{\Lref}[1]{\ref{#1}}
% \Supp{pE}{lem-chart} -> "Lemma E.3" style text is produced by the drafter:
%   "Lemma~\Xref{pE:lem-chart} of the supplement".  \SuppPart{E} gives "Part E".
\newcommand{\SuppPart}[1]{Part~#1 of the supplement}

% ---- numbers (generated by unit U0; macros only, no literal certified numbers) ----
\IfFileExists{numbers/numbers_paper.tex}%
  {% numbers_paper.tex -- GENERATED from ledger.json by build_ledger.py; do not edit.
% Rounding: upper bounds up, lower bounds and margins down at the printed digit; margin = threshold / bound.
% Use as $\times\NName$ for margins.  kind: cert certified, est estimate, hand hypothesis constant, doc only a document quotes it.
% NUEhatZero [cert, upper] U_cert_v35.json :: supersolution.EA.0
\newcommand{\NUEhatZero}{\ensuremath{1.2018\cdot10^{-6}}}
% NUEhatOne [cert, upper] U_cert_v35.json :: supersolution.EA.1
\newcommand{\NUEhatOne}{\ensuremath{1.0441\cdot10^{-4}}}
% NUEhatTwo [cert, upper] U_cert_v35.json :: supersolution.EA.2
\newcommand{\NUEhatTwo}{\ensuremath{0.056694}}
% NUBhatZero [cert, upper] U_cert_v35.json :: supersolution.Bhat.0
\newcommand{\NUBhatZero}{\ensuremath{4.8735\cdot10^{-6}}}
% NUBhatOne [cert, upper] U_cert_v35.json :: supersolution.Bhat.1
\newcommand{\NUBhatOne}{\ensuremath{5.5933\cdot10^{-4}}}
% NUBhatTwo [cert, upper] U_cert_v35.json :: supersolution.Bhat.2
\newcommand{\NUBhatTwo}{\ensuremath{0.36675}}
% NUSupersolutionCZeroZero [cert, upper] U_cert_v35.json :: supersolution.C0.0
\newcommand{\NUSupersolutionCZeroZero}{\ensuremath{1.1340\cdot10^{-6}}}
% NUSupersolutionCZeroOne [cert, upper] U_cert_v35.json :: supersolution.C0.1
\newcommand{\NUSupersolutionCZeroOne}{\ensuremath{1.6623\cdot10^{-4}}}
% NUSupersolutionCOneZero [cert, upper] U_cert_v35.json :: supersolution.C1.0
\newcommand{\NUSupersolutionCOneZero}{\ensuremath{1.5241\cdot10^{-9}}}
% NUSupersolutionCOneOne [cert, upper] U_cert_v35.json :: supersolution.C1.1
\newcommand{\NUSupersolutionCOneOne}{\ensuremath{5.5255\cdot10^{-8}}}
% NUSupersolutionNTwoZero [cert, upper] U_cert_v35.json :: supersolution.N2.0
\newcommand{\NUSupersolutionNTwoZero}{\ensuremath{0.071472}}
% NUSupersolutionNTwoOne [cert, upper] U_cert_v35.json :: supersolution.N2.1
\newcommand{\NUSupersolutionNTwoOne}{\ensuremath{0.33012}}
% NUSupersolutionNTwoTwo [cert, upper] U_cert_v35.json :: supersolution.N2.2
\newcommand{\NUSupersolutionNTwoTwo}{\ensuremath{5.0406}}
% NUSupersolutionNTwoOverFourNTwoRefZero [cert, upper] U_cert_v35.json :: supersolution.N2_over_4N2ref.0
\newcommand{\NUSupersolutionNTwoOverFourNTwoRefZero}{\ensuremath{0.18}}
% NUSupersolutionNTwoOverFourNTwoRefOne [cert, upper] U_cert_v35.json :: supersolution.N2_over_4N2ref.1
\newcommand{\NUSupersolutionNTwoOverFourNTwoRefOne}{\ensuremath{0.42}}
% NUSupersolutionNTwoOverFourNTwoRefTwo [cert, upper] U_cert_v35.json :: supersolution.N2_over_4N2ref.2
\newcommand{\NUSupersolutionNTwoOverFourNTwoRefTwo}{\ensuremath{0.26}}
% NUSupersolutionMOneZHypCheckZonesIIISOneTwoBhat [cert, upper] U_cert_v35.json :: supersolution.M1z_hyp_check.zones_I_II.S1_2Bhat
\newcommand{\NUSupersolutionMOneZHypCheckZonesIIISOneTwoBhat}{\ensuremath{0.28585}}
% NUSupersolutionMOneZHypCheckZoneIIIBOnePlusPUOne [cert, upper] U_cert_v35.json :: supersolution.M1z_hyp_check.zone_III.b1_plus_pU1
\newcommand{\NUSupersolutionMOneZHypCheckZoneIIIBOnePlusPUOne}{\ensuremath{0.50786}}
% NUMarginLateMOne [cert, margin] U_cert_v35.json :: margins.late_M1
\newcommand{\NUMarginLateMOne}{\ensuremath{2.353}}
% NUMarginLateMZero [cert, margin] U_cert_v35.json :: margins.late_M0
\newcommand{\NUMarginLateMZero}{\ensuremath{4.892}}
% NUMarginsLateMZeroZ [cert, margin] U_cert_v35.json :: margins.late_M0z
\newcommand{\NUMarginsLateMZeroZ}{\ensuremath{12.99}}
% NUMarginPresentMOne [cert, margin] U_cert_v35.json :: margins.present_M1_min.value
\newcommand{\NUMarginPresentMOne}{\ensuremath{2.770}}
% NUMarginPresentMZero [cert, margin] U_cert_v35.json :: margins.present_M0_min.value
\newcommand{\NUMarginPresentMZero}{\ensuremath{2.225}}
% NUMarginsPresentMZeroZMinValue [cert, margin] U_cert_v35.json :: margins.present_M0z_min.value
\newcommand{\NUMarginsPresentMZeroZMinValue}{\ensuremath{14.37}}
% NULoopGainZeroTotal [cert, upper] U_cert_v35.json :: loop_gain.0.total
\newcommand{\NULoopGainZeroTotal}{\ensuremath{1.23\cdot10^{-7}}}
% NULoopGainOneTotal [cert, upper] U_cert_v35.json :: loop_gain.1.total
\newcommand{\NULoopGainOneTotal}{\ensuremath{0.037169}}
% NULoopGainTwoNeighbour [cert, upper] U_cert_v35.json :: loop_gain.2.neighbour
\newcommand{\NULoopGainTwoNeighbour}{\ensuremath{0.037018}}
% NULoopGainTwoThroughA [cert, upper] U_cert_v35.json :: loop_gain.2.through_A
\newcommand{\NULoopGainTwoThroughA}{\ensuremath{0.0020172}}
% NUfineBhatZero [cert, upper] U_cert_v35.json :: ufine.Bhat.0
\newcommand{\NUfineBhatZero}{\ensuremath{4.8723\cdot10^{-6}}}
% NUfineBhatOne [cert, upper] U_cert_v35.json :: ufine.Bhat.1
\newcommand{\NUfineBhatOne}{\ensuremath{2.6266\cdot10^{-4}}}
% NUfineBhatTwo [cert, upper] U_cert_v35.json :: ufine.Bhat.2
\newcommand{\NUfineBhatTwo}{\ensuremath{0.025842}}
% NUcdeltaU [cert, upper] U_cert_v35.json :: launch.c_dU
\newcommand{\NUcdeltaU}{\ensuremath{1.7187\cdot10^{-8}}}
% NULaunchWmin [cert, lower] U_cert_v35.json :: launch.wmin
\newcommand{\NULaunchWmin}{\ensuremath{0.73492}}
% NULaunchCAD [cert, upper] U_cert_v35.json :: launch.CA_D
\newcommand{\NULaunchCAD}{\ensuremath{2.2429}}
% NULaunchQD [cert, upper] U_cert_v35.json :: launch.qD
\newcommand{\NULaunchQD}{\ensuremath{0.018509}}
% NULaunchDeltaInZero [cert, upper] U_cert_v35.json :: launch.Delta_in.0
\newcommand{\NULaunchDeltaInZero}{\ensuremath{3.131\cdot10^{-12}}}
% NULaunchDeltaInOne [cert, upper] U_cert_v35.json :: launch.Delta_in.1
\newcommand{\NULaunchDeltaInOne}{\ensuremath{3.287\cdot10^{-12}}}
% NULaunchDeltaInTwo [cert, upper] U_cert_v35.json :: launch.Delta_in.2
\newcommand{\NULaunchDeltaInTwo}{\ensuremath{3.451\cdot10^{-12}}}
% NULaunchPhiZero [cert, upper] U_cert_v35.json :: launch.Phi.0
\newcommand{\NULaunchPhiZero}{\ensuremath{6.434\cdot10^{-9}}}
% NULaunchPhiOne [cert, upper] U_cert_v35.json :: launch.Phi.1
\newcommand{\NULaunchPhiOne}{\ensuremath{1.4217\cdot10^{-8}}}
% NULaunchPhiTwo [cert, upper] U_cert_v35.json :: launch.Phi.2
\newcommand{\NULaunchPhiTwo}{\ensuremath{3.8404\cdot10^{-8}}}
% NUGammaZero [cert, upper] U_cert_v35.json :: Gamma.0
\newcommand{\NUGammaZero}{\ensuremath{4.0000358}}
% NUGammaOne [cert, upper] U_cert_v35.json :: Gamma.1
\newcommand{\NUGammaOne}{\ensuremath{16.000322}}
% NUGammaTwo [cert, upper] U_cert_v35.json :: Gamma.2
\newcommand{\NUGammaTwo}{\ensuremath{64.00218}}
% NUKappa [cert, upper] U_cert_v35.json :: kappa
\newcommand{\NUKappa}{\ensuremath{8.0000358}}
% NUKappaTwo [cert, upper] U_cert_v35.json :: kappa2
\newcommand{\NUKappaTwo}{\ensuremath{64.00061}}
% NUGhatZeroZero [cert, upper] U_cert_v35.json :: Ghat.0.0
\newcommand{\NUGhatZeroZero}{\ensuremath{51322.6}}
% NUGhatOneOne [cert, upper] U_cert_v35.json :: Ghat.1.1
\newcommand{\NUGhatOneOne}{\ensuremath{33962.9}}
% NUGhatTwoTwo [cert, upper] U_cert_v35.json :: Ghat.2.2
\newcommand{\NUGhatTwoTwo}{\ensuremath{348341.1}}
% NUGhatOneZero [cert, upper] U_cert_v35.json :: Ghat.1.0
\newcommand{\NUGhatOneZero}{\ensuremath{4.938}}
% NUGhatTwoZero [cert, upper] U_cert_v35.json :: Ghat.2.0
\newcommand{\NUGhatTwoZero}{\ensuremath{27.984}}
% NUGhatTwoOne [cert, upper] U_cert_v35.json :: Ghat.2.1
\newcommand{\NUGhatTwoOne}{\ensuremath{458.53}}
% NUKZeroZero [cert, upper] U_cert_v35.json :: K.0.0
\newcommand{\NUKZeroZero}{\ensuremath{8034.21}}
% NUKZeroOne [cert, upper] U_cert_v35.json :: K.0.1
\newcommand{\NUKZeroOne}{\ensuremath{724.89}}
% NUKOneZero [cert, upper] U_cert_v35.json :: K.1.0
\newcommand{\NUKOneZero}{\ensuremath{0.08152}}
% NUKOneOne [cert, upper] U_cert_v35.json :: K.1.1
\newcommand{\NUKOneOne}{\ensuremath{6336.94}}
% NUKOneTwo [cert, upper] U_cert_v35.json :: K.1.2
\newcommand{\NUKOneTwo}{\ensuremath{545.59}}
% NUKOneThree [cert, upper] U_cert_v35.json :: K.1.3
\newcommand{\NUKOneThree}{\ensuremath{0.58700}}
% NUKTwoZero [cert, upper] U_cert_v35.json :: K.2.0
\newcommand{\NUKTwoZero}{\ensuremath{0.87186}}
% NUKTwoOne [cert, upper] U_cert_v35.json :: K.2.1
\newcommand{\NUKTwoOne}{\ensuremath{528.62}}
% NUKTwoTwo [cert, upper] U_cert_v35.json :: K.2.2
\newcommand{\NUKTwoTwo}{\ensuremath{23425.98}}
% NUKTwoThree [cert, upper] U_cert_v35.json :: K.2.3
\newcommand{\NUKTwoThree}{\ensuremath{555.62}}
% NUFootwOne [cert, upper] U_cert_v35.json :: footw.1
\newcommand{\NUFootwOne}{\ensuremath{2.6361}}
% NUFootwTwo [cert, upper] U_cert_v35.json :: footw.2
\newcommand{\NUFootwTwo}{\ensuremath{2663.3}}
% NUCThreeThree [cert, upper] U_cert_v35.json :: c33
\newcommand{\NUCThreeThree}{\ensuremath{108.02}}
% NUCTwoTwoThree [cert, upper] U_cert_v35.json :: c223
\newcommand{\NUCTwoTwoThree}{\ensuremath{7.3284}}
% NUCEta [cert, upper] U_cert_v35.json :: C_eta
\newcommand{\NUCEta}{\ensuremath{13.376}}
% NULMax [cert, upper] U_cert_v35.json :: L_max
\newcommand{\NULMax}{\ensuremath{3.5352}}
% NURRatio [cert, upper] U_cert_v35.json :: R_ratio
\newcommand{\NURRatio}{\ensuremath{11.593}}
% NUSweepsSigmaInCoverSup [cert, upper] U_cert_v35.json :: sweeps.Sigma_in.cover_sup
\newcommand{\NUSweepsSigmaInCoverSup}{\ensuremath{5.0892}}
% NUSweepsSigmaECoverSup [cert, upper] U_cert_v35.json :: sweeps.Sigma_E.cover_sup
\newcommand{\NUSweepsSigmaECoverSup}{\ensuremath{0.30846}}
% NUSweepsSXCoverSup [cert, upper] U_cert_v35.json :: sweeps.S_X.cover_sup
\newcommand{\NUSweepsSXCoverSup}{\ensuremath{1.5243}}
% NUSweepsSigmaXCoverSup [cert, upper] U_cert_v35.json :: sweeps.Sigma_X.cover_sup
\newcommand{\NUSweepsSigmaXCoverSup}{\ensuremath{3.0081}}
% NUSweepsSigmaRsTwoCoverSup [cert, upper] U_cert_v35.json :: sweeps.Sigma_rs.2.cover_sup
\newcommand{\NUSweepsSigmaRsTwoCoverSup}{\ensuremath{4.3350\cdot10^{-3}}}
% NUDataTermsLateDZero [cert, upper] U_cert_v35.json :: data_terms_late.D.0
\newcommand{\NUDataTermsLateDZero}{\ensuremath{1.7330\cdot10^{-6}}}
% NUDataTermsLateDOne [cert, upper] U_cert_v35.json :: data_terms_late.D.1
\newcommand{\NUDataTermsLateDOne}{\ensuremath{2.5353\cdot10^{-4}}}
% NUDataTermsLateSlopeShareCOneSigmaXOverDkZero [cert, upper] U_cert_v35.json :: data_terms_late.slope_share_C1SigmaX_over_Dk.0
\newcommand{\NUDataTermsLateSlopeShareCOneSigmaXOverDkZero}{\ensuremath{2.7\cdot10^{-3}}}
% NUDataTermsLateSlopeShareCOneSigmaXOverDkOne [cert, upper] U_cert_v35.json :: data_terms_late.slope_share_C1SigmaX_over_Dk.1
\newcommand{\NUDataTermsLateSlopeShareCOneSigmaXOverDkOne}{\ensuremath{6.6\cdot10^{-4}}}
% NURLtTwoMarginZero [cert, margin] U_cert_v35.json :: R_lt2_margin.0
\newcommand{\NURLtTwoMarginZero}{\ensuremath{4.432}}
% NURLtTwoMarginOne [cert, margin] U_cert_v35.json :: R_lt2_margin.1
\newcommand{\NURLtTwoMarginOne}{\ensuremath{6.697}}
% NUThetaTwoEight [cert, upper] U_cert_v35.json :: theta28
\newcommand{\NUThetaTwoEight}{\ensuremath{0.37011}}
% NUDFiveSix [cert, upper] U_cert_v35.json :: D56
\newcommand{\NUDFiveSix}{\ensuremath{1.0368}}
% NUDTwoEight [cert, upper] U_cert_v35.json :: D28
\newcommand{\NUDTwoEight}{\ensuremath{1.5876}}
% NUSThreeThreeFootOverBhatTwo [cert, upper] U_cert_v35.json :: S33_foot_over_Bhat2
\newcommand{\NUSThreeThreeFootOverBhatTwo}{\ensuremath{1.69\cdot10^{-3}}}
% NSAHJWKappaJ [cert, upper] A_interface.json :: HJW.kappa_J
\newcommand{\NSAHJWKappaJ}{\ensuremath{0.042351}}
% NSAGpZero [cert, lower] A_interface.json :: res.tight.gp.0
\newcommand{\NSAGpZero}{\ensuremath{0.24292}}
% NSAGpOne [cert, lower] A_interface.json :: res.tight.gp.1
\newcommand{\NSAGpOne}{\ensuremath{0.24106}}
% NSAGpTwo [cert, lower] A_interface.json :: res.tight.gp.2
\newcommand{\NSAGpTwo}{\ensuremath{0.11961}}
% NSAAminC [cert, lower] A_interface.json :: res.tight.amin_c
\newcommand{\NSAAminC}{\ensuremath{0.97623}}
% NSARhoJ [cert, upper] A_interface.json :: res.tight.rhoJ
\newcommand{\NSARhoJ}{\ensuremath{5.4995\cdot10^{-7}}}
% NSAKappaOTwoTwo [cert, upper] A_interface.json :: res.tight.kappa_o.2.2
\newcommand{\NSAKappaOTwoTwo}{\ensuremath{1.2416\cdot10^{-5}}}
% NSADsrcZero [cert, upper] A_interface.json :: res.tight.Dsrc.0
\newcommand{\NSADsrcZero}{\ensuremath{1.8092\cdot10^{-7}}}
% NSADsrcOne [cert, upper] A_interface.json :: res.tight.Dsrc.1
\newcommand{\NSADsrcOne}{\ensuremath{9.8276\cdot10^{-5}}}
% NSABZeroDWindowSlopeZero [cert, upper] A_interface.json :: res.tight.B0.D_window_slope.0
\newcommand{\NSABZeroDWindowSlopeZero}{\ensuremath{8.7399\cdot10^{-10}}}
% NSABZeroDWindowSlopeOne [cert, upper] A_interface.json :: res.tight.B0.D_window_slope.1
\newcommand{\NSABZeroDWindowSlopeOne}{\ensuremath{5.2599\cdot10^{-8}}}
% NSAFTwoZero [cert, upper] A_interface.json :: res.tight.F2.0
\newcommand{\NSAFTwoZero}{\ensuremath{5.6506\cdot10^{-13}}}
% NSAFTwoOne [cert, upper] A_interface.json :: res.tight.F2.1
\newcommand{\NSAFTwoOne}{\ensuremath{5.9320\cdot10^{-13}}}
% NSACoupOne [cert, upper] A_interface.json :: res.tight.coup.1
\newcommand{\NSACoupOne}{\ensuremath{1.1568\cdot10^{-12}}}
% NSACOneZeroZero [cert, upper] A_interface.json :: res.tight.C1_0.0
\newcommand{\NSACOneZeroZero}{\ensuremath{7.2060\cdot10^{-10}}}
% NSACOneZeroOne [cert, upper] A_interface.json :: res.tight.C1_0.1
\newcommand{\NSACOneZeroOne}{\ensuremath{4.3143\cdot10^{-8}}}
% NSABZeroSupCoreZero [cert, upper] A_interface.json :: res.tight.B0.sup_core.0
\newcommand{\NSABZeroSupCoreZero}{\ensuremath{4.23\cdot10^{-5}}}
% NSACoreFluxMarginTwoFLUXZeroOverKZero [cert, margin] A_interface.json :: res.tight.core.flux_margin.2.FLUX0_over_K0
\newcommand{\NSACoreFluxMarginTwoFLUXZeroOverKZero}{\ensuremath{3.396}}
% NSACoreFluxMarginTwoFLUXTwoOverKTwo [cert, margin] A_interface.json :: res.tight.core.flux_margin.2.FLUX2_over_K2
\newcommand{\NSACoreFluxMarginTwoFLUXTwoOverKTwo}{\ensuremath{6.012}}
% NSAMarginFluxZero [cert, margin] asm_design_N2ref4.json :: res.tight.flux.beta.2.margin_FLUX0
\newcommand{\NSAMarginFluxZero}{\ensuremath{3.365}}
% NSAMarginFluxTwo [cert, margin] asm_design_N2ref4.json :: res.tight.flux.beta.2.margin_FLUX2
\newcommand{\NSAMarginFluxTwo}{\ensuremath{5.955}}
% NSAfourSZero [cert, upper] asm_design_N2ref4.json :: res.tight.S0
\newcommand{\NSAfourSZero}{\ensuremath{7.1878\cdot10^{-3}}}
% NSAfourSOne [cert, upper] asm_design_N2ref4.json :: res.tight.S1
\newcommand{\NSAfourSOne}{\ensuremath{0.10430}}
% NSAMarginCoreSZeroV [cert, margin] asm_design_N2ref4.json :: res.tight.margins.S0v_vs_0.0153
\newcommand{\NSAMarginCoreSZeroV}{\ensuremath{2.116}}
% NSAMarginCoreSOneV [cert, margin] asm_design_N2ref4.json :: res.tight.margins.S1v_vs_0.6
\newcommand{\NSAMarginCoreSOneV}{\ensuremath{5.748}}
% NSAMarginCoreMOne [cert, margin] asm_design_N2ref4.json :: res.tight.margins.M1core
\newcommand{\NSAMarginCoreMOne}{\ensuremath{4.409}}
% NSAMarginCoreMZero [cert, margin] asm_design_N2ref4.json :: res.tight.margins.M0core
\newcommand{\NSAMarginCoreMZero}{\ensuremath{5.450}}
% NSAMarginGhat [cert, margin] asm_design_N2ref4.json :: res.tight.margins.Ghat22
\newcommand{\NSAMarginGhat}{\ensuremath{1.917}}
% NSAMarginRampMOne [cert, margin] asm_design_N2ref4.json :: res.tight.margins.ramp M1
\newcommand{\NSAMarginRampMOne}{\ensuremath{1.184}}
% NSAfourInfoZeroDWindowConst [cert, upper] asm_design_N2ref4.json :: res.tight.info.0.D_window_const
\newcommand{\NSAfourInfoZeroDWindowConst}{\ensuremath{1.7732\cdot10^{-7}}}
% NSAfourInfoOneDWindowConst [cert, upper] asm_design_N2ref4.json :: res.tight.info.1.D_window_const
\newcommand{\NSAfourInfoOneDWindowConst}{\ensuremath{9.8058\cdot10^{-5}}}
% NZDeltaRZero [cert, upper] zone3_interface.json :: DeltaR.0
\newcommand{\NZDeltaRZero}{\ensuremath{4.4986\cdot10^{-8}}}
% NZDeltaROne [cert, upper] zone3_interface.json :: DeltaR.1
\newcommand{\NZDeltaROne}{\ensuremath{6.7972\cdot10^{-6}}}
% NZDeltaRTwo [cert, upper] zone3_interface.json :: DeltaR.2
\newcommand{\NZDeltaRTwo}{\ensuremath{0.011480}}
% NZDeltaRBZero [cert, upper] zone3_interface.json :: DeltaR_b.0
\newcommand{\NZDeltaRBZero}{\ensuremath{2.66\cdot10^{-12}}}
% NZDeltaRBOne [cert, upper] zone3_interface.json :: DeltaR_b.1
\newcommand{\NZDeltaRBOne}{\ensuremath{3.26\cdot10^{-9}}}
% NZDeltaRBTwo [cert, upper] zone3_interface.json :: DeltaR_b.2
\newcommand{\NZDeltaRBTwo}{\ensuremath{4.35\cdot10^{-8}}}
% NZDeltaRDetailNSqZero [cert, upper] zone3_interface.json :: DeltaR_detail.N_sq.0
\newcommand{\NZDeltaRDetailNSqZero}{\ensuremath{8.28\cdot10^{-5}}}
% NZDeltaRDetailNSqOne [cert, upper] zone3_interface.json :: DeltaR_detail.N_sq.1
\newcommand{\NZDeltaRDetailNSqOne}{\ensuremath{0.10149}}
% NZDeltaRDetailNSqTwo [cert, upper] zone3_interface.json :: DeltaR_detail.N_sq.2
\newcommand{\NZDeltaRDetailNSqTwo}{\ensuremath{1.3567}}
% NZDeltaRDetailNSqThree [cert, upper] zone3_interface.json :: DeltaR_detail.N_sq.3
\newcommand{\NZDeltaRDetailNSqThree}{\ensuremath{25350}}
% NZDeltaRDetailPartsTwoPairsThree [cert, upper] zone3_interface.json :: DeltaR_detail.parts.2.pairs.3
\newcommand{\NZDeltaRDetailPartsTwoPairsThree}{\ensuremath{0.011298}}
% NZDeltaRDetailPartsTwoLFourFoot [cert, upper] zone3_interface.json :: DeltaR_detail.parts.2.L4_foot
\newcommand{\NZDeltaRDetailPartsTwoLFourFoot}{\ensuremath{1.6\cdot10^{-4}}}
% NZDeltaRDetailPartsTwoCubic [cert, upper] zone3_interface.json :: DeltaR_detail.parts.2.cubic
\newcommand{\NZDeltaRDetailPartsTwoCubic}{\ensuremath{9.3\cdot10^{-6}}}
% NZDeltaRDetailPartsTwoSThreeThreeFoot [cert, upper] zone3_interface.json :: DeltaR_detail.parts.2.S33_foot
\newcommand{\NZDeltaRDetailPartsTwoSThreeThreeFoot}{\ensuremath{3.7\cdot10^{-7}}}
% NZNIIIZero [cert, upper] zone3_interface.json :: n_III.0
\newcommand{\NZNIIIZero}{\ensuremath{0.071445}}
% NZNIIIOne [cert, upper] zone3_interface.json :: n_III.1
\newcommand{\NZNIIIOne}{\ensuremath{0.32707}}
% NZNIIITwo [cert, upper] zone3_interface.json :: n_III.2
\newcommand{\NZNIIITwo}{\ensuremath{3.0445}}
% NZTauIII [cert, upper] zone3_interface.json :: tau_III
\newcommand{\NZTauIII}{\ensuremath{4.0708}}
% NZoneIIIBoxZero [cert, upper] zone3_interface.json :: box_III_380.b0
\newcommand{\NZoneIIIBoxZero}{\ensuremath{0.10032}}
% NZoneIIIBoxOne [cert, upper] zone3_interface.json :: box_III_380.b1
\newcommand{\NZoneIIIBoxOne}{\ensuremath{0.41974}}
% NZoneIIIBoxZ [cert, upper] zone3_interface.json :: box_III_380.bz
\newcommand{\NZoneIIIBoxZ}{\ensuremath{0.014291}}
% NZBoxIIIThreeEightZeroMarginsBZero [cert, margin] zone3_interface.json :: box_III_380.margins.b0
\newcommand{\NZBoxIIIThreeEightZeroMarginsBZero}{\ensuremath{1.495}}
% NZoneIIIMarginThinnest [cert, margin] zone3_interface.json :: box_III_380.margins.b1
\newcommand{\NZoneIIIMarginThinnest}{\ensuremath{1.238}}
% NZBoxIIIThreeEightZeroMarginsBz [cert, margin] zone3_interface.json :: box_III_380.margins.bz
\newcommand{\NZBoxIIIThreeEightZeroMarginsBz}{\ensuremath{11.19}}
% NZBoxIIISevenZeroZeroBZero [cert, upper] zone3_interface.json :: box_III_700.b0
\newcommand{\NZBoxIIISevenZeroZeroBZero}{\ensuremath{0.22757}}
% NZBoxIIISevenZeroZeroBOne [cert, upper] zone3_interface.json :: box_III_700.b1
\newcommand{\NZBoxIIISevenZeroZeroBOne}{\ensuremath{0.54234}}
% NZUfineKTwoAtBhatRef [cert, upper] zone3_interface.json :: ufine.K2_at_Bhat_ref
\newcommand{\NZUfineKTwoAtBhatRef}{\ensuremath{0.046956}}
% NZUfineKThreeStripAtBhatRef [cert, upper] zone3_interface.json :: ufine.K3_strip_at_Bhat_ref
\newcommand{\NZUfineKThreeStripAtBhatRef}{\ensuremath{0.054140}}
% NRecoveryMinMargin [cert, margin] zone3_interface.json :: cone_family.H6.at_Bhat_ref.min_margin
\newcommand{\NRecoveryMinMargin}{\ensuremath{1.860}}
% NZConeFamilyHSixAtBhatRefConeSlopeMargin [cert, margin] zone3_interface.json :: cone_family.H6.at_Bhat_ref.cone.slope.margin
\newcommand{\NZConeFamilyHSixAtBhatRefConeSlopeMargin}{\ensuremath{2.46}}
% NZConeFamilyHSixAtBhatRefConeKTwoMargin [cert, margin] zone3_interface.json :: cone_family.H6.at_Bhat_ref.cone.K2.margin
\newcommand{\NZConeFamilyHSixAtBhatRefConeKTwoMargin}{\ensuremath{2.54}}
% NZConeFamilyHSixAtBhatRefConeKThreeMargin [cert, margin] zone3_interface.json :: cone_family.H6.at_Bhat_ref.cone.K3.margin
\newcommand{\NZConeFamilyHSixAtBhatRefConeKThreeMargin}{\ensuremath{2.51}}
% NConeTstarMargin [cert, margin] zone3_interface.json :: acceptance.tstar_margin
\newcommand{\NConeTstarMargin}{\ensuremath{2.30}}
% NUoneTFiveYaddLevelOneTwo [cert, upper] u1_ghat.json :: t5_Yadd.level1.2
\newcommand{\NUoneTFiveYaddLevelOneTwo}{\ensuremath{0.0044}}
% NUoneTFiveYaddLevelTwoTwo [cert, upper] u1_ghat.json :: t5_Yadd.level2.2
\newcommand{\NUoneTFiveYaddLevelTwoTwo}{\ensuremath{0.1081}}
% NUoneTFiveYaddLevelTwoThree [cert, upper] u1_ghat.json :: t5_Yadd.level2.3
\newcommand{\NUoneTFiveYaddLevelTwoThree}{\ensuremath{0.374}}
% NUoneWeightThetaGammaZero [cert, upper] u1_ghat.json :: weight.Theta.Gamma.0
\newcommand{\NUoneWeightThetaGammaZero}{\ensuremath{4.00000000000000000000054}}
% NRKZeroZero [cert, upper] ramp_U.json :: K.0.0
\newcommand{\NRKZeroZero}{\ensuremath{13436.03}}
% NRKOneOne [cert, upper] ramp_U.json :: K.1.1
\newcommand{\NRKOneOne}{\ensuremath{10526.19}}
% NRKTwoTwo [cert, upper] ramp_U.json :: K.2.2
\newcommand{\NRKTwoTwo}{\ensuremath{39482.30}}
% NUoneCasesPrimaryYNextZero [cert, upper] u1_ghat.json :: cases.primary.Y_next.0
\newcommand{\NUoneCasesPrimaryYNextZero}{\ensuremath{457.092}}
% NUoneCasesPrimaryYNextOne [cert, upper] u1_ghat.json :: cases.primary.Y_next.1
\newcommand{\NUoneCasesPrimaryYNextOne}{\ensuremath{356.940}}
% NUoneCasesPrimaryYNextTwo [cert, upper] u1_ghat.json :: cases.primary.Y_next.2
\newcommand{\NUoneCasesPrimaryYNextTwo}{\ensuremath{1320.220}}
% NSfourResultsMNineHTwoTightenedFUOneRowsZeroBound [cert, upper] FINAL_manifest_v33.json :: results.M9_H2_tightened_fU1.rows.0.bound
\newcommand{\NSfourResultsMNineHTwoTightenedFUOneRowsZeroBound}{\ensuremath{3.658\cdot10^{-10}}}
% NSfourResultsMNineHTwoTightenedFUOneRowsZeroMargin [cert, margin] FINAL_manifest_v33.json :: results.M9_H2_tightened_fU1.rows.0.margin
\newcommand{\NSfourResultsMNineHTwoTightenedFUOneRowsZeroMargin}{\ensuremath{27.34}}
% NSfourResultsMNineHTwoTightenedFUOneRowsOneBound [cert, upper] FINAL_manifest_v33.json :: results.M9_H2_tightened_fU1.rows.1.bound
\newcommand{\NSfourResultsMNineHTwoTightenedFUOneRowsOneBound}{\ensuremath{1.777\cdot10^{-9}}}
% NSfourResultsMNineHTwoTightenedFUOneRowsOneMargin [cert, margin] FINAL_manifest_v33.json :: results.M9_H2_tightened_fU1.rows.1.margin
\newcommand{\NSfourResultsMNineHTwoTightenedFUOneRowsOneMargin}{\ensuremath{562.8}}
% NSfourResultsMNineHTwoTightenedFUOneRowsTwoBound [cert, upper] FINAL_manifest_v33.json :: results.M9_H2_tightened_fU1.rows.2.bound
\newcommand{\NSfourResultsMNineHTwoTightenedFUOneRowsTwoBound}{\ensuremath{1.665\cdot10^{-6}}}
% NSfourResultsMNineHTwoTightenedFUOneRowsTwoMargin [cert, margin] FINAL_manifest_v33.json :: results.M9_H2_tightened_fU1.rows.2.margin
\newcommand{\NSfourResultsMNineHTwoTightenedFUOneRowsTwoMargin}{\ensuremath{600.9}}
% NHOneAlphaThr [hand, exact] hand :: 
\newcommand{\NHOneAlphaThr}{\ensuremath{1.5\cdot10^{-4}}}
% NHOneSqrtLamThr [hand, exact] hand :: 
\newcommand{\NHOneSqrtLamThr}{\ensuremath{5\cdot10^{-8}}}
% NYbarStep [hand, exact] hand :: 
\newcommand{\NYbarStep}{\ensuremath{1\cdot10^{-12}}}
% NHTwoThrZero [hand, exact] hand :: 
\newcommand{\NHTwoThrZero}{\ensuremath{1\cdot10^{-8}}}
% NHTwoThrOne [hand, exact] hand :: 
\newcommand{\NHTwoThrOne}{\ensuremath{1\cdot10^{-6}}}
% NHTwoThrTwo [hand, exact] hand :: 
\newcommand{\NHTwoThrTwo}{\ensuremath{1\cdot10^{-3}}}
% NHThreeThrZero [hand, exact] hand :: 
\newcommand{\NHThreeThrZero}{\ensuremath{1\cdot10^{-6}}}
% NHThreeThrOne [hand, exact] hand :: 
\newcommand{\NHThreeThrOne}{\ensuremath{5\cdot10^{-6}}}
% NHThreeThrTwo [hand, exact] hand :: 
\newcommand{\NHThreeThrTwo}{\ensuremath{1\cdot10^{-5}}}
% NHFourThrZero [hand, exact] hand :: 
\newcommand{\NHFourThrZero}{\ensuremath{2\cdot10^{-7}}}
% NHFourThrOne [hand, exact] hand :: 
\newcommand{\NHFourThrOne}{\ensuremath{1\cdot10^{-5}}}
% NHFourThrTwo [hand, exact] hand :: 
\newcommand{\NHFourThrTwo}{\ensuremath{1\cdot10^{-5}}}
% NHFourThrThree [hand, exact] hand :: 
\newcommand{\NHFourThrThree}{\ensuremath{1\cdot10^{-5}}}
% NHFourKfourThr [hand, exact] hand :: 
\newcommand{\NHFourKfourThr}{\ensuremath{0.06}}
% NHSixConeSlopeThr [hand, exact] hand :: 
\newcommand{\NHSixConeSlopeThr}{\ensuremath{0.08}}
% NHSixConeKtwoThr [hand, exact] hand :: 
\newcommand{\NHSixConeKtwoThr}{\ensuremath{0.03}}
% NHSixConeNthreeThr [hand, exact] hand :: 
\newcommand{\NHSixConeNthreeThr}{\ensuremath{0.004}}
% NHSixConeJetThr [hand, exact] hand :: 
\newcommand{\NHSixConeJetThr}{\ensuremath{0.12}}
% NHSixConeHFiveThr [hand, exact] hand :: 
\newcommand{\NHSixConeHFiveThr}{\ensuremath{0.0144}}
% NHSixFarHFiveThr [hand, exact] hand :: 
\newcommand{\NHSixFarHFiveThr}{\ensuremath{0.0225}}
% NLemmaOConeM [hand, exact] hand :: 
\newcommand{\NLemmaOConeM}{\ensuremath{0.12}}
% NLemmaOConeK [hand, exact] hand :: 
\newcommand{\NLemmaOConeK}{\ensuremath{0.18}}
% NLemmaOFarM [hand, exact] hand :: 
\newcommand{\NLemmaOFarM}{\ensuremath{0.15}}
% NLemmaOFarK [hand, exact] hand :: 
\newcommand{\NLemmaOFarK}{\ensuremath{0.3}}
% NCoercivityConst [hand, exact] hand :: 
\newcommand{\NCoercivityConst}{\ensuremath{0.046875}}
% NMOnez [hand, exact] hand :: 
\newcommand{\NMOnez}{\ensuremath{0.534}}
% NBetaA [hand, exact] hand :: 
\newcommand{\NBetaA}{\ensuremath{2}}
% NTOneLower [cert, lower] t1clock.json :: t1.0
\newcommand{\NTOneLower}{\ensuremath{4.5978574775}}
% NTOneUpper [cert, upper] t1clock.json :: t1.1
\newcommand{\NTOneUpper}{\ensuremath{4.5978586144}}
% NLambdaOneLower [cert, lower] t1clock.json :: lnlam.0
\newcommand{\NLambdaOneLower}{\ensuremath{6.49\cdot10^{-5}}}
% NLambdaOneUpper [cert, upper] t1clock.json :: lnlam.1
\newcommand{\NLambdaOneUpper}{\ensuremath{6.75\cdot10^{-5}}}
% NTMinusTOneUpper [cert, upper] t1clock.json :: lnlam.1; 1.5e-4.sup_ab
\newcommand{\NTMinusTOneUpper}{\ensuremath{2.1\cdot10^{-17}}}
% NAbarUpper [cert, upper] final_alpha_points_5e-8.json :: 1.5e-4.sup_ab
\newcommand{\NAbarUpper}{\ensuremath{2.19103\cdot10^{-7}}}
% NBetabarUpper [cert, upper] final_alpha_points_5e-8.json :: 1.5e-4.sup_beta
\newcommand{\NBetabarUpper}{\ensuremath{3.33811\cdot10^{-9}}}
% NBaseKmaxLower [cert, lower] t1clock.json :: base_kmax.0
\newcommand{\NBaseKmaxLower}{\ensuremath{2.279\cdot10^{4}}}
% NBaseKmaxUpper [cert, upper] t1clock.json :: base_kmax.1
\newcommand{\NBaseKmaxUpper}{\ensuremath{2.376\cdot10^{4}}}
% NHOneAlphaBase [cert, upper] summary.json :: alpha_abs_max
\newcommand{\NHOneAlphaBase}{\ensuremath{1.1739\cdot10^{-4}}}
% NHOneAlphaErr [cert, upper] t9_h1_clam_s180.json :: C_alpha; inputs.YSTAR
\newcommand{\NHOneAlphaErr}{\ensuremath{7.21\cdot10^{-13}}}
% NHOneAlphaMargin [cert, margin] summary.json :: alpha_abs_max; C_alpha; inputs.YSTAR
\newcommand{\NHOneAlphaMargin}{\ensuremath{1.27}}
% NHOneSqrtLamBase [cert, upper] t9_h1_clam_s180.json :: sqrtLam_base_Wp
\newcommand{\NHOneSqrtLamBase}{\ensuremath{2.282\cdot10^{-8}}}
% NHOneSqrtLamTotal [cert, upper] t9_h1_clam_s180.json :: sqrtLam_total_Wp
\newcommand{\NHOneSqrtLamTotal}{\ensuremath{2.29\cdot10^{-8}}}
% NHOneSqrtLamMargin [cert, margin] t9_h1_clam_s180.json :: sqrtLam_total_Wp
\newcommand{\NHOneSqrtLamMargin}{\ensuremath{2.18}}
% NHOneYbarLimit [cert, lower] t9_h1_clam_s180.json :: Ybar_limit_H1
\newcommand{\NHOneYbarLimit}{\ensuremath{1.12\cdot10^{-9}}}
% NSwitchYUpper [cert, upper] switch_s2.json :: modulated.Y_raw
\newcommand{\NSwitchYUpper}{\ensuremath{6.21\cdot10^{-13}}}
% NSwitchYMargin [cert, margin] switch_s2.json :: modulated.Y_raw
\newcommand{\NSwitchYMargin}{\ensuremath{1.6}}
% NMinCI [cert, lower] constants.json :: meta.summary.min_c_I
\newcommand{\NMinCI}{\ensuremath{0.05283}}
% NNumCertificates [cert, exact] constants.json :: meta.summary.n_cert
\newcommand{\NNumCertificates}{\ensuremath{26}}
% NMinRhoC [cert, lower] constants.json :: meta.closure.rho_c_min
\newcommand{\NMinRhoC}{\ensuremath{0.04711}}
% NHTwoBaseZero [cert, upper] FINAL_manifest_v33.json :: results.M9_H2_tightened.rows.0.base
\newcommand{\NHTwoBaseZero}{\ensuremath{8.011\cdot10^{-11}}}
% NHTwoBaseOne [cert, upper] FINAL_manifest_v33.json :: results.M9_H2_tightened.rows.1.base
\newcommand{\NHTwoBaseOne}{\ensuremath{1.336\cdot10^{-10}}}
% NHTwoBaseTwo [cert, upper] FINAL_manifest_v33.json :: results.M9_H2_tightened.rows.2.base
\newcommand{\NHTwoBaseTwo}{\ensuremath{2.936\cdot10^{-10}}}
% NHTwoErrWZero [cert, upper] FINAL_manifest_v33.json :: results.M9_H2_tightened.rows.0.with_error
\newcommand{\NHTwoErrWZero}{\ensuremath{8.014\cdot10^{-11}}}
% NHTwoMarginWZero [cert, margin] FINAL_manifest_v33.json :: results.M9_H2_tightened.rows.0.margin
\newcommand{\NHTwoMarginWZero}{\ensuremath{124.78}}
% NHTwoErrWOne [cert, upper] FINAL_manifest_v33.json :: results.M9_H2_tightened.rows.1.with_error
\newcommand{\NHTwoErrWOne}{\ensuremath{4.706\cdot10^{-10}}}
% NHTwoMarginWOne [cert, margin] FINAL_manifest_v33.json :: results.M9_H2_tightened.rows.1.margin
\newcommand{\NHTwoMarginWOne}{\ensuremath{2124.9}}
% NHTwoErrWTwo [cert, upper] FINAL_manifest_v33.json :: results.M9_H2_tightened.rows.2.with_error
\newcommand{\NHTwoErrWTwo}{\ensuremath{1.559\cdot10^{-6}}}
% NHTwoMarginWTwo [cert, margin] FINAL_manifest_v33.json :: results.M9_H2_tightened.rows.2.margin
\newcommand{\NHTwoMarginWTwo}{\ensuremath{641.53}}
% NHThreeBaseZero [cert, upper] FINAL_manifest_v33.json :: results.M9_H2f.rows.0.base
\newcommand{\NHThreeBaseZero}{\ensuremath{3.127\cdot10^{-7}}}
% NHThreeBaseOne [cert, upper] FINAL_manifest_v33.json :: results.M9_H2f.rows.1.base
\newcommand{\NHThreeBaseOne}{\ensuremath{1.004\cdot10^{-6}}}
% NHThreeBaseTwo [cert, upper] FINAL_manifest_v33.json :: results.M9_H2f.rows.2.base
\newcommand{\NHThreeBaseTwo}{\ensuremath{2.431\cdot10^{-6}}}
% NHThreeErrZero [cert, upper] FINAL_manifest_v33.json :: results.M9_H2f.rows.0.with_error
\newcommand{\NHThreeErrZero}{\ensuremath{3.1363\cdot10^{-7}}}
% NHThreeMarginZero [cert, margin] FINAL_manifest_v33.json :: results.M9_H2f.rows.0.margin
\newcommand{\NHThreeMarginZero}{\ensuremath{3.1885}}
% NHThreeErrOne [cert, upper] FINAL_manifest_v33.json :: results.M9_H2f.rows.1.with_error
\newcommand{\NHThreeErrOne}{\ensuremath{1.0034\cdot10^{-6}}}
% NHThreeMarginOne [cert, margin] FINAL_manifest_v33.json :: results.M9_H2f.rows.1.margin
\newcommand{\NHThreeMarginOne}{\ensuremath{4.9834}}
% NHThreeErrTwo [cert, upper] FINAL_manifest_v33.json :: results.M9_H2f.rows.2.with_error
\newcommand{\NHThreeErrTwo}{\ensuremath{2.4351\cdot10^{-6}}}
% NHThreeMarginTwo [cert, margin] FINAL_manifest_v33.json :: results.M9_H2f.rows.2.margin
\newcommand{\NHThreeMarginTwo}{\ensuremath{4.1067}}
% NHFourBaseZero [cert, upper] FINAL_manifest_v33.json :: results.M9_base_sums.combined.windows_max.0
\newcommand{\NHFourBaseZero}{\ensuremath{3.48\cdot10^{-8}}}
% NHFourErrZero [cert, upper] FINAL_manifest_v33.json :: results.M9_H3_H2v32.rows.win.0.with_error_arb
\newcommand{\NHFourErrZero}{\ensuremath{3.5087\cdot10^{-8}}}
% NHFourMarginZero [cert, margin] FINAL_manifest_v33.json :: results.M9_H3_H2v32.rows.win.0.margin_arb
\newcommand{\NHFourMarginZero}{\ensuremath{5.700}}
% NHFourBaseOne [cert, upper] FINAL_manifest_v33.json :: results.M9_base_sums.combined.windows_max.1
\newcommand{\NHFourBaseOne}{\ensuremath{1.27\cdot10^{-7}}}
% NHFourErrOne [cert, upper] FINAL_manifest_v33.json :: results.M9_H3_H2v32.rows.win.1.with_error_arb
\newcommand{\NHFourErrOne}{\ensuremath{1.2616\cdot10^{-7}}}
% NHFourMarginOne [cert, margin] FINAL_manifest_v33.json :: results.M9_H3_H2v32.rows.win.1.margin_arb
\newcommand{\NHFourMarginOne}{\ensuremath{79.26}}
% NHFourBaseTwo [cert, upper] FINAL_manifest_v33.json :: results.M9_base_sums.combined.windows_max.2
\newcommand{\NHFourBaseTwo}{\ensuremath{3.76\cdot10^{-7}}}
% NHFourErrTwo [cert, upper] FINAL_manifest_v33.json :: results.M9_H3_H2v32.rows.win.2.with_error_arb
\newcommand{\NHFourErrTwo}{\ensuremath{3.7724\cdot10^{-7}}}
% NHFourMarginTwo [cert, margin] FINAL_manifest_v33.json :: results.M9_H3_H2v32.rows.win.2.margin_arb
\newcommand{\NHFourMarginTwo}{\ensuremath{26.50}}
% NHFourBaseThree [cert, upper] FINAL_manifest_v33.json :: results.M9_base_sums.combined.windows_max.3
\newcommand{\NHFourBaseThree}{\ensuremath{9.71\cdot10^{-7}}}
% NHFourErrThree [cert, upper] FINAL_manifest_v33.json :: results.M9_H3_H2v32.rows.win.3.with_error_arb
\newcommand{\NHFourErrThree}{\ensuremath{1.2822\cdot10^{-6}}}
% NHFourMarginThree [cert, margin] FINAL_manifest_v33.json :: results.M9_H3_H2v32.rows.win.3.margin_arb
\newcommand{\NHFourMarginThree}{\ensuremath{7.799}}
% NHFourYbarLimit [cert, lower] FINAL_manifest_v33.json :: results.M9_H3_H2v32.binding_Ybar_limit
\newcommand{\NHFourYbarLimit}{\ensuremath{1.575\cdot10^{-11}}}
% NHFourYbarMargin [cert, margin] FINAL_manifest_v33.json :: results.M9_H3_H2v32.binding_Ybar_limit
\newcommand{\NHFourYbarMargin}{\ensuremath{15.75}}
% NHFourKfourBase [cert, upper] amp_A1_w180.json :: inputs.BASE_T1.K4
\newcommand{\NHFourKfourBase}{\ensuremath{3.1633\cdot10^{-4}}}
% NHSixConeSlope [cert, upper] FINAL_manifest_v33.json :: results.cone_fixedL.checks.slope.value
\newcommand{\NHSixConeSlope}{\ensuremath{0.03029}}
% NHSixConeSlopeMargin [cert, margin] FINAL_manifest_v33.json :: results.cone_fixedL.checks.slope.margin
\newcommand{\NHSixConeSlopeMargin}{\ensuremath{2.641}}
% NHSixConeKtwo [cert, upper] FINAL_manifest_v33.json :: results.cone_fixedL.checks.K2.value
\newcommand{\NHSixConeKtwo}{\ensuremath{0.009876}}
% NHSixConeKtwoMargin [cert, margin] FINAL_manifest_v33.json :: results.cone_fixedL.checks.K2.margin
\newcommand{\NHSixConeKtwoMargin}{\ensuremath{3.037}}
% NHSixConeNthree [cert, upper] FINAL_manifest_v33.json :: results.cone_fixedL.checks.N3.value
\newcommand{\NHSixConeNthree}{\ensuremath{1.461\cdot10^{-4}}}
% NHSixConeNthreeMargin [cert, margin] FINAL_manifest_v33.json :: results.cone_fixedL.checks.N3.margin
\newcommand{\NHSixConeNthreeMargin}{\ensuremath{27.38}}
% NHSixConeHFive [cert, upper] FINAL_manifest_v33.json :: results.H4_fixedL.total
\newcommand{\NHSixConeHFive}{\ensuremath{1.2865\cdot10^{-5}}}
% NHSixConeHFiveMargin [cert, margin] FINAL_manifest_v33.json :: results.H4_fixedL.margin
\newcommand{\NHSixConeHFiveMargin}{\ensuremath{1119}}
% NHSixConeTstarNeeded [cert, upper] FINAL_manifest_v33.json :: results.H7_fixedL.ratio_upper
\newcommand{\NHSixConeTstarNeeded}{\ensuremath{1.6000015\cdot10^{-7}}}
% NHSixConeTstarLower [cert, lower] FINAL_manifest_v33.json :: results.H7_fixedL.t_star_lower
\newcommand{\NHSixConeTstarLower}{\ensuremath{3.68334\cdot10^{-7}}}
% NHSixConeTstarMargin [cert, margin] FINAL_manifest_v33.json :: results.H7_fixedL.margin_tstar_over_ratio
\newcommand{\NHSixConeTstarMargin}{\ensuremath{2.302}}
% NHSixConeHFiveRatioPeak [cert, upper] h4loc_A1_w180.json :: windows.constants/amp_A1_w180.json.worst.ratio
\newcommand{\NHSixConeHFiveRatioPeak}{\ensuremath{0.2393}}
% NHSixConeHOneSup [cert, upper] FINAL_manifest_v33.json :: results.H1_fixedL.sup_h0_max
\newcommand{\NHSixConeHOneSup}{\ensuremath{0.2163}}
% NHSixConeHOneH [cert, upper] FINAL_manifest_v33.json :: results.H1_fixedL.H_max
\newcommand{\NHSixConeHOneH}{\ensuremath{0.2185}}
% NHSixFarHFiveRatio [cert, upper] h4loc_v32.json :: windows.constants/amp_A1_w180.json.cylinders.18.ratio
\newcommand{\NHSixFarHFiveRatio}{\ensuremath{0.1761}}
% NHSixFarHFiveMarginOld [cert, margin] h4loc_v32.json :: windows.constants/amp_A1_w180.json.cylinders.18.ratio
\newcommand{\NHSixFarHFiveMarginOld}{\ensuremath{5.68}}
% NHSixFarHFiveMargin [cert, margin] FINAL_manifest_v33.json :: results.M9_H4_far.min_margin
\newcommand{\NHSixFarHFiveMargin}{\ensuremath{15.58}}
% NHSixConeHFiveMarginOld [cert, margin] h4loc_A1_w180.json :: windows.constants/amp_A1_w180.json.worst.ratio
\newcommand{\NHSixConeHFiveMarginOld}{\ensuremath{4.18}}
% NHSixFarKtwo [cert, upper] t9_far_v32.json :: h2
\newcommand{\NHSixFarKtwo}{\ensuremath{1.797\cdot10^{-4}}}
% NHSixFarKtwoThr [hand, exact] t9_far_v32.json :: thr.h2
\newcommand{\NHSixFarKtwoThr}{\ensuremath{3.53\cdot10^{-4}}}
% NHSixFarKtwoMargin [cert, margin] t9_far_v32.json :: margins.h2
\newcommand{\NHSixFarKtwoMargin}{\ensuremath{1.964}}
% NHSixFarKthree [cert, upper] t9_far_v32.json :: h3
\newcommand{\NHSixFarKthree}{\ensuremath{2.595\cdot10^{-8}}}
% NHSixFarKthreeThr [hand, exact] t9_far_v32.json :: thr.h3
\newcommand{\NHSixFarKthreeThr}{\ensuremath{1.17\cdot10^{-5}}}
% NHSixFarKthreeMargin [cert, margin] t9_far_v32.json :: margins.h3
\newcommand{\NHSixFarKthreeMargin}{\ensuremath{450.9}}
% NHSixFarKfour [cert, upper] t9_far_v32.json :: h4
\newcommand{\NHSixFarKfour}{\ensuremath{6.964\cdot10^{-9}}}
% NHSixFarKfourThr [hand, exact] t9_far_v32.json :: thr.h4
\newcommand{\NHSixFarKfourThr}{\ensuremath{1.03\cdot10^{-7}}}
% NHSixFarKfourMargin [cert, margin] t9_far_v32.json :: margins.h4
\newcommand{\NHSixFarKfourMargin}{\ensuremath{14.79}}
% NHSixFarSlope [cert, upper] t9_far_v32.json :: slope
\newcommand{\NHSixFarSlope}{\ensuremath{0.04153}}
% NHSixFarSlopeThr [hand, exact] t9_far_v32.json :: thr.slope
\newcommand{\NHSixFarSlopeThr}{\ensuremath{0.0800}}
% NHSixFarSlopeMargin [cert, margin] t9_far_v32.json :: margins.slope
\newcommand{\NHSixFarSlopeMargin}{\ensuremath{1.926}}
% NHSixFarNthree [cert, upper] t9_far_v32.json :: L2
\newcommand{\NHSixFarNthree}{\ensuremath{5.516\cdot10^{-7}}}
% NHSixFarNthreeThr [hand, exact] t9_far_v32.json :: thr.L2
\newcommand{\NHSixFarNthreeThr}{\ensuremath{3.32\cdot10^{-6}}}
% NHSixFarNthreeMargin [cert, margin] t9_far_v32.json :: margins.L2
\newcommand{\NHSixFarNthreeMargin}{\ensuremath{6.019}}
% NFarTstarMargin [cert, margin] lemmaO_v32_consts.json :: far_v32.case.margin_factor_lower
\newcommand{\NFarTstarMargin}{\ensuremath{15.5}}
% NFarTstarLower [cert, lower] lemmaO_v32_consts.json :: far_v32.case.t_star_lower
\newcommand{\NFarTstarLower}{\ensuremath{9.5108\cdot10^{-8}}}
% NFarTstarNeeded [cert, upper] lemmaO_v32_consts.json :: far_v32.case.needed
\newcommand{\NFarTstarNeeded}{\ensuremath{6.134\cdot10^{-9}}}
% NGammaModes [cert, exact] gamma0_final.json :: M
\newcommand{\NGammaModes}{\ensuremath{3762}}
% NGammaHalfwidth [cert, upper] gamma0_final.json :: b_j_halfwidth
\newcommand{\NGammaHalfwidth}{\ensuremath{8.138\cdot10^{-24}}}
% NGammaContraction [cert, upper] gamma0_final.json :: krawczyk_contraction
\newcommand{\NGammaContraction}{\ensuremath{3.01\cdot10^{-12}}}
% NGammaArcs [cert, exact] gamma0_final.json :: embedded.arcs
\newcommand{\NGammaArcs}{\ensuremath{701}}
% NGammaMaxTurn [cert, upper] gamma0_final.json :: embedded.maxturn
\newcommand{\NGammaMaxTurn}{\ensuremath{0.04997}}
% NGammaMaxPair [cert, upper] gamma0_final.json :: embedded.maxpair
\newcommand{\NGammaMaxPair}{\ensuremath{0.09597}}
% NGammaMinSep [cert, lower] gamma0_final.json :: embedded.minsep
\newcommand{\NGammaMinSep}{\ensuremath{8.89\cdot10^{-3}}}
% NGammaHalfGap [cert, lower] gamma0_final.json :: embedded.gap_cert
\newcommand{\NGammaHalfGap}{\ensuremath{0.0248219}}
% NPhaseISlabs [cert, exact] status.json :: slab
\newcommand{\NPhaseISlabs}{\ensuremath{96}}
% NPhaseIModes [cert, exact] status.json :: last_slab.M
\newcommand{\NPhaseIModes}{\ensuremath{2800}}
% NPhaseIKmaxLower [cert, lower] status.json :: kmax_enclosure.0
\newcommand{\NPhaseIKmaxLower}{\ensuremath{29.352874545}}
% NPhaseIKmaxUpper [cert, upper] status.json :: kmax_enclosure.1
\newcommand{\NPhaseIKmaxUpper}{\ensuremath{29.352874546}}
% NPhaseITLower [cert, lower] status.json :: t
\newcommand{\NPhaseITLower}{\ensuremath{4.597850539485830}}
% NPhaseITUpper [cert, upper] status.json :: t
\newcommand{\NPhaseITUpper}{\ensuremath{4.597850539485831}}
% NPeakSqrtLam [cert, upper] rate_table_fin12f4_a0.00015.json :: sqrtL_max
\newcommand{\NPeakSqrtLam}{\ensuremath{5.308\cdot10^{-8}}}
% NTailRate [cert, lower] rate_table_fin12f4_a0.00015.json :: tail.rho
\newcommand{\NTailRate}{\ensuremath{2.693\cdot10^{-3}}}
% NFloorSqrtLam [cert, upper] rate_table_fin12f4_a0.00015.json :: tail.sqrt_LF
\newcommand{\NFloorSqrtLam}{\ensuremath{9.872\cdot10^{-10}}}
% NAlphaFloor [cert, upper] rate_table_fin12f4_a0.00015.json :: tail.alpha_F
\newcommand{\NAlphaFloor}{\ensuremath{6.867\cdot10^{-7}}}
% NSupBetaVerifier [cert, upper] final_alpha_points_5e-8.json :: 1.5e-4.sup_beta
\newcommand{\NSupBetaVerifier}{\ensuremath{3.33811\cdot10^{-9}}}
% NMOneRamp [cert, upper] final_alpha_points_5e-8.json :: 1.5e-4.M1
\newcommand{\NMOneRamp}{\ensuremath{0.4328}}
% NSITwoPrime [cert, upper] FINAL_manifest_v33.json :: results.F12_sliver.sI2_prime
\newcommand{\NSITwoPrime}{\ensuremath{0.015457}}
% NLoopGainCoarse [cert, upper] U_cert_v35.json :: loop_gain.2.total
\newcommand{\NLoopGainCoarse}{\ensuremath{0.03904}}
% NRateMarginCoarse [cert, margin] U_cert_v35.json :: rate_margin
\newcommand{\NRateMarginCoarse}{\ensuremath{28.2}}
% NSOneEpsLower [cert, lower] bracket2.json :: eps_a
\newcommand{\NSOneEpsLower}{\ensuremath{1.5346198892965\cdot10^{-4}}}
% NSOneEpsUpper [cert, upper] bracket2.json :: eps_b
\newcommand{\NSOneEpsUpper}{\ensuremath{1.5346198899299\cdot10^{-4}}}
% NSOneRowsTip [cert, exact] ext_tip_check.json :: rows
\newcommand{\NSOneRowsTip}{\ensuremath{51246}}
% NSOneTipWidth [cert, upper] ext_tip_check.json :: max_pointwise_width.th
\newcommand{\NSOneTipWidth}{\ensuremath{3.2\cdot10^{-12}}}
% NSOneExtWidth [cert, upper] ext813_check.json :: max_pointwise_width.p
\newcommand{\NSOneExtWidth}{\ensuremath{5.7\cdot10^{-10}}}
% NSOneRowsMC [cert, exact] check_outA5.json :: prof_tip_vs_MC.rows_compared
\newcommand{\NSOneRowsMC}{\ensuremath{1025}}
% NSOneRowsMB [cert, exact] check_outA5.json :: ext813_vs_MB.rows_compared
\newcommand{\NSOneRowsMB}{\ensuremath{501}}
% NSOneMinMargin [cert, margin] check_outA5.json :: min:*.min_relative_margin.*
\newcommand{\NSOneMinMargin}{\ensuremath{1.2\cdot10^{-9}}}
% NTwoSegments [cert, exact] contour_result.json :: len:segments
\newcommand{\NTwoSegments}{\ensuremath{5058}}
% NTwoInteriorBoxes [cert, exact] interior_result.json :: len:
\newcommand{\NTwoInteriorBoxes}{\ensuremath{174}}
% NTwoContourBoxes [cert, exact] eta_cache.json :: len:
\newcommand{\NTwoContourBoxes}{\ensuremath{202}}
% NTwoEtaMax [cert, upper] eta_cache.json :: max:*.eta
\newcommand{\NTwoEtaMax}{\ensuremath{0.1163}}
% NTwoLipMax [cert, upper] eta_cache.json :: max:*.Lip
\newcommand{\NTwoLipMax}{\ensuremath{0.272}}
% NTwoLambdaZero [est, upper] cert_b0.25_th0.0_tau0.03.json :: Lam
\newcommand{\NTwoLambdaZero}{\ensuremath{1.790819}}
% NTwoLambdaOne [est, upper] cert_b0.25_th0.48_tau0.03.json :: Lam
\newcommand{\NTwoLambdaOne}{\ensuremath{1.768282}}
% NSOneDTightLower [doc, lower] doc:README.md :: D ∈ [2.50080627654, 2.50080627668]
\newcommand{\NSOneDTightLower}{\ensuremath{2.50080627654}}
% NSOneDTightUpper [doc, upper] doc:README.md :: D ∈ [2.50080627654, 2.50080627668]
\newcommand{\NSOneDTightUpper}{\ensuremath{2.50080627668}}
% NSOneKTightLower [doc, lower] doc:README.md :: k(0) ∈ [1.53982063762, 1.53982063822]
\newcommand{\NSOneKTightLower}{\ensuremath{1.53982063762}}
% NSOneKTightUpper [doc, upper] doc:README.md :: k(0) ∈ [1.53982063762, 1.53982063822]
\newcommand{\NSOneKTightUpper}{\ensuremath{1.53982063822}}
% NSOneDCondLower [doc, lower] doc:README.md :: D = −p(tip) ∈ [2.500806276619695, 2.500806276620049]
\newcommand{\NSOneDCondLower}{\ensuremath{2.500806276619695}}
% NSOneDCondUpper [doc, upper] doc:README.md :: D = −p(tip) ∈ [2.500806276619695, 2.500806276620049]
\newcommand{\NSOneDCondUpper}{\ensuremath{2.500806276620049}}
% NSOneKCondLower [doc, lower] doc:README.md :: k(0) ∈ [1.539820637871836, 1.5398206378721015]
\newcommand{\NSOneKCondLower}{\ensuremath{1.539820637871836}}
% NSOneKCondUpper [doc, upper] doc:README.md :: k(0) ∈ [1.539820637871836, 1.5398206378721015]
\newcommand{\NSOneKCondUpper}{\ensuremath{1.5398206378721015}}
% NSOneTauTipLower [doc, lower] doc:README.md :: τ_tip ∈ [10.249077651120778, 10.249077651121361]
\newcommand{\NSOneTauTipLower}{\ensuremath{10.249077651120778}}
% NSOneTauTipUpper [doc, upper] doc:README.md :: τ_tip ∈ [10.249077651120778, 10.249077651121361]
\newcommand{\NSOneTauTipUpper}{\ensuremath{10.249077651121361}}
% NSOneStateUncert [doc, upper] doc:README.md :: State uncertainty at σ = 8: 4.4e−13
\newcommand{\NSOneStateUncert}{\ensuremath{4.4\cdot10^{-13}}}
% NFarBoxEprime [doc, margin] doc:STEP4_RESULT.md :: (6.4287, 4.6917, 1.9940, 1.69999995)
\newcommand{\NFarBoxEprime}{\ensuremath{6.42}}
% NFarBoxXEsecond [doc, margin] doc:STEP4_RESULT.md :: (6.4287, 4.6917, 1.9940, 1.69999995)
\newcommand{\NFarBoxXEsecond}{\ensuremath{4.69}}
% NFarBoxXsqEthird [doc, margin] doc:STEP4_RESULT.md :: (6.4287, 4.6917, 1.9940, 1.69999995)
\newcommand{\NFarBoxXsqEthird}{\ensuremath{1.99}}
% NFarBoxE [doc, margin] doc:STEP4_RESULT.md :: (6.4287, 4.6917, 1.9940, 1.69999995)
\newcommand{\NFarBoxE}{\ensuremath{1.69}}
% NHFourKfourErr [doc, upper] doc:STEP4_RESULT.md :: | 5.60e-3 against 0.06 | ×10.7 |
\newcommand{\NHFourKfourErr}{\ensuremath{5.60\cdot10^{-3}}}
% NHFourKfourMargin [doc, margin] doc:STEP4_RESULT.md :: | 5.60e-3 against 0.06 | ×10.7 |
\newcommand{\NHFourKfourMargin}{\ensuremath{10.7}}
% NCoercivityCoefficient [doc, upper] doc:step3_proof_v32.tex :: \norm{e^{\tau\mathcal L}u}_\rho\le14.8\,e^{-0.0213\tau}\norm
\newcommand{\NCoercivityCoefficient}{\ensuremath{14.8}}
% NCoercivityRate [doc, lower] doc:step3_proof_v32.tex :: \norm{e^{\tau\mathcal L}u}_\rho\le14.8\,e^{-0.0213\tau}\norm
\newcommand{\NCoercivityRate}{\ensuremath{0.0213}}
% NRegionRampStart [hand, exact] hand :: 
\newcommand{\NRegionRampStart}{\ensuremath{14}}
% NRegionInterface [hand, exact] hand :: 
\newcommand{\NRegionInterface}{\ensuremath{28}}
% NRegionShellEnd [hand, exact] hand :: 
\newcommand{\NRegionShellEnd}{\ensuremath{60}}
% NRegionSegEnd [hand, exact] hand :: 
\newcommand{\NRegionSegEnd}{\ensuremath{350}}
% NRegionAnnEnd [hand, exact] hand :: 
\newcommand{\NRegionAnnEnd}{\ensuremath{700}}
% NZoneIIILow [hand, exact] hand :: 
\newcommand{\NZoneIIILow}{\ensuremath{253}}
% NZoneIIIHigh [hand, exact] hand :: 
\newcommand{\NZoneIIIHigh}{\ensuremath{379.5}}
% NConeCentreMin [hand, exact] hand :: 
\newcommand{\NConeCentreMin}{\ensuremath{300}}
% NConeCentreMax [hand, exact] hand :: 
\newcommand{\NConeCentreMax}{\ensuremath{380}}
% NFarCentreMin [hand, exact] hand :: 
\newcommand{\NFarCentreMin}{\ensuremath{490}}
% NConeL [hand, exact] hand :: 
\newcommand{\NConeL}{\ensuremath{50}}
% NConeX [hand, exact] hand :: 
\newcommand{\NConeX}{\ensuremath{3}}
% NFarL [hand, exact] hand :: 
\newcommand{\NFarL}{\ensuremath{113}}
% NFarX [hand, exact] hand :: 
\newcommand{\NFarX}{\ensuremath{2}}
}%
  {\PackageWarning{paper}{numbers/numbers_paper.tex not found (unit U0 not finished)}}

% ======================================================================
% NOTATION MACROS (OUTLINE 5.2 and notation.md).  One macro per symbol
% whose spelling is fixed by the notation table.  Plain letters (a, b, f, ...)
% have no macro.
% ======================================================================
% basic
\newcommand{\R}{\mathbb{R}}
\newcommand{\C}{\mathbb{C}}
\newcommand{\N}{\mathbb{N}}
\newcommand{\Z}{\mathbb{Z}}
\newcommand{\eps}{\varepsilon}
\newcommand{\ip}[2]{\langle #1,#2\rangle}
\newcommand{\abs}[1]{\lvert #1\rvert}
\newcommand{\norm}[1]{\lVert #1\rVert}
\newcommand{\dd}{\mathrm d}
\newcommand{\re}{\operatorname{Re}}
\newcommand{\im}{\operatorname{Im}}
% flow and curves
\newcommand{\arclen}{\mathfrak{s}}                 % arclength of the flow curve
\newcommand{\kss}{\kappa_{\mathfrak{ss}}}          % kappa_{ss} in the flow equation
\newcommand{\Area}{\operatorname{Area}}
\newcommand{\Len}{\mathcal L}                      % length of a curve
\newcommand{\spp}{\operatorname{sp}}                 % sup quantities sp^v_0, sp^v_1, sp^f_1 (were S^v_0, S^v_1, S^f_1)
\newcommand{\rot}{\operatorname{rot}}              % turning number
\newcommand{\Tsing}{T_*}                           % singular time
\newcommand{\Tmax}{T_{\max}}                       % maximal time of existence
\newcommand{\pstar}{p_*}                           % image point of the singularity
\newcommand{\Gz}{\Gamma_0}                         % initial curve
\newcommand{\vth}{\vartheta}                       % parameter of Gamma_0
\newcommand{\dclk}{\delta_t}                       % clock defect of step 4
\newcommand{\fall}{\varpi}                         % Fourier rates varpi_n = n^4-n^2
% rescaling
\newcommand{\lamone}{\lambda_1}
\newcommand{\rate}{\varsigma}                      % 1/4 - a
\newcommand{\xia}{\xi_a}                           % (1-4\bar a)^{-1}
% shrinker and step 1
\newcommand{\Sone}{S_1}
\newcommand{\Gsh}{\Gamma_{\!*}}                    % profile curve of S_1
\newcommand{\cSet}{\mathcal S}                     % certified set of shrinkers
\newcommand{\sigtip}{\sigma_{\rm tip}}
\newcommand{\Xo}{X_0}                              % 2.2490776515
\newcommand{\Dtip}{D}
\newcommand{\ktip}{k_0}
\newcommand{\shootM}{M}                            % shooting function
\newcommand{\Ypar}{\mathcal Y}                     % step-1 parallelepiped
\newcommand{\axis}{\Sigma}                         % symmetry axis of S_1
\newcommand{\sigc}{\sigma_{\rm c}}
\newcommand{\sigrev}{\sigma^{\rm rev}}
% linear theory
\newcommand{\Hrho}{\mathcal H_\rho}
\newcommand{\muslow}{\mu_3}
\newcommand{\rre}{\mathfrak r_{\rm re}}
\newcommand{\rim}{\mathfrak r_{\rm im}}
\newcommand{\Evev}{\mathrm{Ev}_{\rm even}}
\newcommand{\Evod}{\mathrm{Ev}_{\rm odd}}
\newcommand{\gvec}{\mathbf g}                      % translation vector in step 2
\newcommand{\cons}[1]{\mathfrak f_{#1}}            % constraint functions
\newcommand{\phirho}{\varphi_\rho}
% energies
\newcommand{\Lam}{\Lambda}
\newcommand{\En}{E}
\newcommand{\EnI}{E_I}
\newcommand{\muI}{\mu_I}                           % 0.08
\newcommand{\epsE}{\varepsilon_E}                  % 1e-5
\newcommand{\epsI}{\varepsilon_I}                  % 3e-5
\newcommand{\wn}[1]{\mathsf w_{#1}}                % ||w~^{(j)}||_rho
\newcommand{\vI}{v_I}                              % zeta w'
\newcommand{\muerr}{\mu_{\rm err}}                 % step-4 error-energy weight 0.036
% far field
\newcommand{\pU}{p_U}                              % 1.0246
\newcommand{\Uone}{U_1}
\newcommand{\sm}{\operatorname{sm}}                % smoothstep t^3(10-15t+6t^2)
\newcommand{\rc}{r_c}                              % centre radius (Lemma U)
\newcommand{\rchi}{r_\chi}                         % 253
\newcommand{\Iw}{I_w}                              % hand-off windows [y_w,1.5y_w]
\newcommand{\qa}{\mathfrak q_1}                    % quasilinear coefficients of the slope equation
\newcommand{\qb}{\mathfrak q_2}
% cylinders (Lemma O)
\newcommand{\cyl}{\ell}                            % cylinder length unit
\newcommand{\MO}{M_O}
\newcommand{\KO}{K_O}
\newcommand{\Hcyl}{H_{\rm cyl}}
\newcommand{\XiO}{\Xi_O}
\newcommand{\fO}{f_O}
% step 4
\newcommand{\Ybar}{\bar Y}
\newcommand{\Lt}{\widetilde L}                     % linearisation about the step-4 base
\newcommand{\Trans}{\Pi}                           % transport integral
\newcommand{\resid}{\mathrm{res}}
\newcommand{\Resid}[1]{\mathrm{Res}_{#1}}
\newcommand{\boxtip}{\Box_{\rm tip}}
\newcommand{\boxfar}{\Box_{\rm far}}
\newcommand{\Eenv}{\mathcal E_{\rm env}}
\newcommand{\betab}{\beta_{\rm base}}
\newcommand{\tauI}{\tau^{\rm I}}                   % Phase I clock
% names of families of conditions: \Hh{3} = (H3), \Oh{4} = (O4), ...
\newcommand{\Hh}[1]{\textup{(H#1)}}
\newcommand{\Oh}[1]{\textup{(O#1)}}
\newcommand{\Lh}[1]{\textup{(L#1)}}
\newcommand{\Bh}[1]{\textup{(B#1)}}
\newcommand{\Kh}[1]{\textup{(K#1)}}
\newcommand{\Wh}[1]{\textup{(W#1)}}
\newcommand{\PhaseI}{Phase~I}
\newcommand{\PhaseII}{Phase~II}
% margin: threshold/bound (always)
\newcommand{\mgn}[1]{{\times}#1}

% ---- figures: placeholder macro so that drafts compile before F1-F7 exist ----
\newcommand{\figplaceholder}[1]{\fbox{\parbox{0.8\linewidth}{\centering\small[figure #1 to be supplied]}}}
 ...
%   paper_v1/supplement.tex  \documentclass[11pt]{amsart}
%                            \newif\ifsupplement \supplementtrue
%                            % ======================================================================
% preamble.tex -- shared preamble of the main paper and the supplement.
% Owner: coordinator.  Drafters add macros by request only (see DRAFTER_GUIDE.md).
%
% USE.  This file contains NO \documentclass.  The two master files are
%   paper_v1/main.tex        \documentclass[11pt]{amsart} % ======================================================================
% preamble.tex -- shared preamble of the main paper and the supplement.
% Owner: coordinator.  Drafters add macros by request only (see DRAFTER_GUIDE.md).
%
% USE.  This file contains NO \documentclass.  The two master files are
%   paper_v1/main.tex        \documentclass[11pt]{amsart} \input{preamble} ...
%   paper_v1/supplement.tex  \documentclass[11pt]{amsart}
%                            \newif\ifsupplement \supplementtrue
%                            \input{preamble} ...
% (the main paper does nothing special: \ifsupplement defaults to false).
% Compile from paper_v1/ so that \input{numbers/numbers_paper.tex} is found.
% ======================================================================

% ---- which document are we in ----------------------------------------
\ifdefined\ifsupplement\else\newif\ifsupplement\fi

% ---- packages ---------------------------------------------------------
\usepackage[T1]{fontenc}
\usepackage[utf8]{inputenc}
\usepackage{lmodern}
\usepackage[margin=1.1in]{geometry}
\usepackage{amsmath,amssymb,amsthm,mathtools,bm}
\usepackage{booktabs,array,longtable,tabularx,multirow}
\usepackage{enumitem}
\usepackage{graphicx}
\usepackage{xcolor}
\usepackage{tikz}
\usetikzlibrary{positioning,arrows.meta,calc}
% xr-hyper must come before hyperref
\usepackage{xr-hyper}
\usepackage[colorlinks=true,linkcolor=blue!50!black,citecolor=green!35!black,urlcolor=blue!50!black]{hyperref}

\raggedbottom
\setlength{\emergencystretch}{2em}
\setlist[enumerate]{leftmargin=*,itemsep=2pt,topsep=3pt}
\setlist[itemize]{leftmargin=*,itemsep=2pt,topsep=3pt}

% ---- numbering --------------------------------------------------------
% Main paper: sections 1,2,...; theorems 2.3 (by section); equations (2.3).
% Supplement: the Parts A..K are the \section's, numbered by letter, so
% theorems are A.1, E.4, ... (do not confuse "Theorem A" of the paper, which
% has no dot, with "Theorem A.1" of Part A).
\ifsupplement
  \renewcommand{\thesection}{\Alph{section}}
\fi
\numberwithin{equation}{section}

% ---- theorem environments (numbered by section, one shared counter) ----
\theoremstyle{plain}
\newtheorem{theorem}{Theorem}[section]
\newtheorem{proposition}[theorem]{Proposition}
\newtheorem{lemma}[theorem]{Lemma}
\newtheorem{corollary}[theorem]{Corollary}
\newtheorem{maintheorem}[theorem]{Main Theorem}
% Theorems A, B, C, D of the paper: own counter, lettered, no section prefix
\newtheorem{thmlet}{Theorem}
\renewcommand{\thethmlet}{\Alph{thmlet}}
\theoremstyle{definition}
\newtheorem{definition}[theorem]{Definition}
\newtheorem{convention}[theorem]{Convention}
\newtheorem{example}[theorem]{Example}
\theoremstyle{remark}
\newtheorem{remark}[theorem]{Remark}

% A list of hypotheses labelled (H1),(H2),... : \begin{hyplist}{H} ... \end{hyplist}
% (also {O}, {L}, {B}, {K}, {W} for the families of section 5.3 of OUTLINE.md).
\newenvironment{hyplist}[1]%
  {\begin{enumerate}[label=\textup{(#1\arabic*)},leftmargin=3em,itemsep=2pt]}%
  {\end{enumerate}}

% ---- trust classes (OUTLINE 4.1, 5.1.3) -------------------------------
% Supplement: explicit tags.  Main paper: certified is the default and
% unmarked; every non-certified number is marked "(floating point)".
\ifsupplement
  \newcommand{\tcert}{\textup{[cert]}}   % arb/outward-rounded, pinned script and key
  \newcommand{\test}{\textup{[est]}}     % floating point or estimate; never in a theorem
  \newcommand{\thand}{\textup{[hand]}}   % proved by hand in the text
  \newcommand{\tlit}{\textup{[lit]}}     % cited from the literature
\else
  \newcommand{\tcert}{}
  \newcommand{\test}{\unskip\ \textup{(floating point)}}
  \newcommand{\thand}{}
  \newcommand{\tlit}{}
\fi

% ---- gaps: visible, counted, an error in a final build -----------------
% Usage: \GAP{one sentence saying what is missing}   (and an entry in ISSUES.md)
\newcounter{gapcount}
\newif\iffinalpaper   % \finalpapertrue makes every \GAP a hard error
\newcommand{\GAP}[1]{\stepcounter{gapcount}\typeout{GAP-MARKER \thegapcount}%
  \iffinalpaper\PackageError{paper}{\string\GAP\space used in a final build}{Remove all gaps.}\fi
  {\color{red}\sffamily\bfseries[GAP~\thegapcount:\ }{\color{red}\sffamily #1}{\color{red}\sffamily\bfseries]}}
\AtEndDocument{\typeout{GAP-MARKER total: \thegapcount}}

% ---- file names and JSON keys (supplement only) -------------------------
% No internal path appears in the main paper; a use there is a warning.
\newcommand{\fn}[1]{%
  \ifsupplement\else\PackageWarning{paper}{\string\fn\space used in the main paper}\fi
  \texttt{\detokenize{#1}}}
\newcommand{\key}[1]{\texttt{\detokenize{#1}}}

% ---- cross references between main paper and supplement ----------------
% Labels carry a unit prefix (section 5.3 of OUTLINE.md and DRAFTER_GUIDE.md):
%   main paper:  intro: obj: s3: s4: s5: s6: s7: s8: s9:   (kind: thm- lem- eq- tab- fig- ...)
%   supplement:  pA: pB: ... pK:
% so all labels are globally unique and the prefixes below are only a second guard.
% \Xref{label}  : reference to the OTHER document (hyperlinked via xr-hyper)
% \Lref{label}  : reference in the SAME document (plain \ref)
\ifsupplement
  \externaldocument[M-]{main}[main.pdf]
  \newcommand{\Xref}[1]{\ref{M-#1}}
\else
  \externaldocument[S-]{supplement}[supplement.pdf]
  \newcommand{\Xref}[1]{\ref{S-#1}}
\fi
\newcommand{\Lref}[1]{\ref{#1}}
% \Supp{pE}{lem-chart} -> "Lemma E.3" style text is produced by the drafter:
%   "Lemma~\Xref{pE:lem-chart} of the supplement".  \SuppPart{E} gives "Part E".
\newcommand{\SuppPart}[1]{Part~#1 of the supplement}

% ---- numbers (generated by unit U0; macros only, no literal certified numbers) ----
\IfFileExists{numbers/numbers_paper.tex}%
  {\input{numbers/numbers_paper.tex}}%
  {\PackageWarning{paper}{numbers/numbers_paper.tex not found (unit U0 not finished)}}

% ======================================================================
% NOTATION MACROS (OUTLINE 5.2 and notation.md).  One macro per symbol
% whose spelling is fixed by the notation table.  Plain letters (a, b, f, ...)
% have no macro.
% ======================================================================
% basic
\newcommand{\R}{\mathbb{R}}
\newcommand{\C}{\mathbb{C}}
\newcommand{\N}{\mathbb{N}}
\newcommand{\Z}{\mathbb{Z}}
\newcommand{\eps}{\varepsilon}
\newcommand{\ip}[2]{\langle #1,#2\rangle}
\newcommand{\abs}[1]{\lvert #1\rvert}
\newcommand{\norm}[1]{\lVert #1\rVert}
\newcommand{\dd}{\mathrm d}
\newcommand{\re}{\operatorname{Re}}
\newcommand{\im}{\operatorname{Im}}
% flow and curves
\newcommand{\arclen}{\mathfrak{s}}                 % arclength of the flow curve
\newcommand{\kss}{\kappa_{\mathfrak{ss}}}          % kappa_{ss} in the flow equation
\newcommand{\Area}{\operatorname{Area}}
\newcommand{\Len}{\mathcal L}                      % length of a curve
\newcommand{\spp}{\operatorname{sp}}                 % sup quantities sp^v_0, sp^v_1, sp^f_1 (were S^v_0, S^v_1, S^f_1)
\newcommand{\rot}{\operatorname{rot}}              % turning number
\newcommand{\Tsing}{T_*}                           % singular time
\newcommand{\Tmax}{T_{\max}}                       % maximal time of existence
\newcommand{\pstar}{p_*}                           % image point of the singularity
\newcommand{\Gz}{\Gamma_0}                         % initial curve
\newcommand{\vth}{\vartheta}                       % parameter of Gamma_0
\newcommand{\dclk}{\delta_t}                       % clock defect of step 4
\newcommand{\fall}{\varpi}                         % Fourier rates varpi_n = n^4-n^2
% rescaling
\newcommand{\lamone}{\lambda_1}
\newcommand{\rate}{\varsigma}                      % 1/4 - a
\newcommand{\xia}{\xi_a}                           % (1-4\bar a)^{-1}
% shrinker and step 1
\newcommand{\Sone}{S_1}
\newcommand{\Gsh}{\Gamma_{\!*}}                    % profile curve of S_1
\newcommand{\cSet}{\mathcal S}                     % certified set of shrinkers
\newcommand{\sigtip}{\sigma_{\rm tip}}
\newcommand{\Xo}{X_0}                              % 2.2490776515
\newcommand{\Dtip}{D}
\newcommand{\ktip}{k_0}
\newcommand{\shootM}{M}                            % shooting function
\newcommand{\Ypar}{\mathcal Y}                     % step-1 parallelepiped
\newcommand{\axis}{\Sigma}                         % symmetry axis of S_1
\newcommand{\sigc}{\sigma_{\rm c}}
\newcommand{\sigrev}{\sigma^{\rm rev}}
% linear theory
\newcommand{\Hrho}{\mathcal H_\rho}
\newcommand{\muslow}{\mu_3}
\newcommand{\rre}{\mathfrak r_{\rm re}}
\newcommand{\rim}{\mathfrak r_{\rm im}}
\newcommand{\Evev}{\mathrm{Ev}_{\rm even}}
\newcommand{\Evod}{\mathrm{Ev}_{\rm odd}}
\newcommand{\gvec}{\mathbf g}                      % translation vector in step 2
\newcommand{\cons}[1]{\mathfrak f_{#1}}            % constraint functions
\newcommand{\phirho}{\varphi_\rho}
% energies
\newcommand{\Lam}{\Lambda}
\newcommand{\En}{E}
\newcommand{\EnI}{E_I}
\newcommand{\muI}{\mu_I}                           % 0.08
\newcommand{\epsE}{\varepsilon_E}                  % 1e-5
\newcommand{\epsI}{\varepsilon_I}                  % 3e-5
\newcommand{\wn}[1]{\mathsf w_{#1}}                % ||w~^{(j)}||_rho
\newcommand{\vI}{v_I}                              % zeta w'
\newcommand{\muerr}{\mu_{\rm err}}                 % step-4 error-energy weight 0.036
% far field
\newcommand{\pU}{p_U}                              % 1.0246
\newcommand{\Uone}{U_1}
\newcommand{\sm}{\operatorname{sm}}                % smoothstep t^3(10-15t+6t^2)
\newcommand{\rc}{r_c}                              % centre radius (Lemma U)
\newcommand{\rchi}{r_\chi}                         % 253
\newcommand{\Iw}{I_w}                              % hand-off windows [y_w,1.5y_w]
\newcommand{\qa}{\mathfrak q_1}                    % quasilinear coefficients of the slope equation
\newcommand{\qb}{\mathfrak q_2}
% cylinders (Lemma O)
\newcommand{\cyl}{\ell}                            % cylinder length unit
\newcommand{\MO}{M_O}
\newcommand{\KO}{K_O}
\newcommand{\Hcyl}{H_{\rm cyl}}
\newcommand{\XiO}{\Xi_O}
\newcommand{\fO}{f_O}
% step 4
\newcommand{\Ybar}{\bar Y}
\newcommand{\Lt}{\widetilde L}                     % linearisation about the step-4 base
\newcommand{\Trans}{\Pi}                           % transport integral
\newcommand{\resid}{\mathrm{res}}
\newcommand{\Resid}[1]{\mathrm{Res}_{#1}}
\newcommand{\boxtip}{\Box_{\rm tip}}
\newcommand{\boxfar}{\Box_{\rm far}}
\newcommand{\Eenv}{\mathcal E_{\rm env}}
\newcommand{\betab}{\beta_{\rm base}}
\newcommand{\tauI}{\tau^{\rm I}}                   % Phase I clock
% names of families of conditions: \Hh{3} = (H3), \Oh{4} = (O4), ...
\newcommand{\Hh}[1]{\textup{(H#1)}}
\newcommand{\Oh}[1]{\textup{(O#1)}}
\newcommand{\Lh}[1]{\textup{(L#1)}}
\newcommand{\Bh}[1]{\textup{(B#1)}}
\newcommand{\Kh}[1]{\textup{(K#1)}}
\newcommand{\Wh}[1]{\textup{(W#1)}}
\newcommand{\PhaseI}{Phase~I}
\newcommand{\PhaseII}{Phase~II}
% margin: threshold/bound (always)
\newcommand{\mgn}[1]{{\times}#1}

% ---- figures: placeholder macro so that drafts compile before F1-F7 exist ----
\newcommand{\figplaceholder}[1]{\fbox{\parbox{0.8\linewidth}{\centering\small[figure #1 to be supplied]}}}
 ...
%   paper_v1/supplement.tex  \documentclass[11pt]{amsart}
%                            \newif\ifsupplement \supplementtrue
%                            % ======================================================================
% preamble.tex -- shared preamble of the main paper and the supplement.
% Owner: coordinator.  Drafters add macros by request only (see DRAFTER_GUIDE.md).
%
% USE.  This file contains NO \documentclass.  The two master files are
%   paper_v1/main.tex        \documentclass[11pt]{amsart} \input{preamble} ...
%   paper_v1/supplement.tex  \documentclass[11pt]{amsart}
%                            \newif\ifsupplement \supplementtrue
%                            \input{preamble} ...
% (the main paper does nothing special: \ifsupplement defaults to false).
% Compile from paper_v1/ so that \input{numbers/numbers_paper.tex} is found.
% ======================================================================

% ---- which document are we in ----------------------------------------
\ifdefined\ifsupplement\else\newif\ifsupplement\fi

% ---- packages ---------------------------------------------------------
\usepackage[T1]{fontenc}
\usepackage[utf8]{inputenc}
\usepackage{lmodern}
\usepackage[margin=1.1in]{geometry}
\usepackage{amsmath,amssymb,amsthm,mathtools,bm}
\usepackage{booktabs,array,longtable,tabularx,multirow}
\usepackage{enumitem}
\usepackage{graphicx}
\usepackage{xcolor}
\usepackage{tikz}
\usetikzlibrary{positioning,arrows.meta,calc}
% xr-hyper must come before hyperref
\usepackage{xr-hyper}
\usepackage[colorlinks=true,linkcolor=blue!50!black,citecolor=green!35!black,urlcolor=blue!50!black]{hyperref}

\raggedbottom
\setlength{\emergencystretch}{2em}
\setlist[enumerate]{leftmargin=*,itemsep=2pt,topsep=3pt}
\setlist[itemize]{leftmargin=*,itemsep=2pt,topsep=3pt}

% ---- numbering --------------------------------------------------------
% Main paper: sections 1,2,...; theorems 2.3 (by section); equations (2.3).
% Supplement: the Parts A..K are the \section's, numbered by letter, so
% theorems are A.1, E.4, ... (do not confuse "Theorem A" of the paper, which
% has no dot, with "Theorem A.1" of Part A).
\ifsupplement
  \renewcommand{\thesection}{\Alph{section}}
\fi
\numberwithin{equation}{section}

% ---- theorem environments (numbered by section, one shared counter) ----
\theoremstyle{plain}
\newtheorem{theorem}{Theorem}[section]
\newtheorem{proposition}[theorem]{Proposition}
\newtheorem{lemma}[theorem]{Lemma}
\newtheorem{corollary}[theorem]{Corollary}
\newtheorem{maintheorem}[theorem]{Main Theorem}
% Theorems A, B, C, D of the paper: own counter, lettered, no section prefix
\newtheorem{thmlet}{Theorem}
\renewcommand{\thethmlet}{\Alph{thmlet}}
\theoremstyle{definition}
\newtheorem{definition}[theorem]{Definition}
\newtheorem{convention}[theorem]{Convention}
\newtheorem{example}[theorem]{Example}
\theoremstyle{remark}
\newtheorem{remark}[theorem]{Remark}

% A list of hypotheses labelled (H1),(H2),... : \begin{hyplist}{H} ... \end{hyplist}
% (also {O}, {L}, {B}, {K}, {W} for the families of section 5.3 of OUTLINE.md).
\newenvironment{hyplist}[1]%
  {\begin{enumerate}[label=\textup{(#1\arabic*)},leftmargin=3em,itemsep=2pt]}%
  {\end{enumerate}}

% ---- trust classes (OUTLINE 4.1, 5.1.3) -------------------------------
% Supplement: explicit tags.  Main paper: certified is the default and
% unmarked; every non-certified number is marked "(floating point)".
\ifsupplement
  \newcommand{\tcert}{\textup{[cert]}}   % arb/outward-rounded, pinned script and key
  \newcommand{\test}{\textup{[est]}}     % floating point or estimate; never in a theorem
  \newcommand{\thand}{\textup{[hand]}}   % proved by hand in the text
  \newcommand{\tlit}{\textup{[lit]}}     % cited from the literature
\else
  \newcommand{\tcert}{}
  \newcommand{\test}{\unskip\ \textup{(floating point)}}
  \newcommand{\thand}{}
  \newcommand{\tlit}{}
\fi

% ---- gaps: visible, counted, an error in a final build -----------------
% Usage: \GAP{one sentence saying what is missing}   (and an entry in ISSUES.md)
\newcounter{gapcount}
\newif\iffinalpaper   % \finalpapertrue makes every \GAP a hard error
\newcommand{\GAP}[1]{\stepcounter{gapcount}\typeout{GAP-MARKER \thegapcount}%
  \iffinalpaper\PackageError{paper}{\string\GAP\space used in a final build}{Remove all gaps.}\fi
  {\color{red}\sffamily\bfseries[GAP~\thegapcount:\ }{\color{red}\sffamily #1}{\color{red}\sffamily\bfseries]}}
\AtEndDocument{\typeout{GAP-MARKER total: \thegapcount}}

% ---- file names and JSON keys (supplement only) -------------------------
% No internal path appears in the main paper; a use there is a warning.
\newcommand{\fn}[1]{%
  \ifsupplement\else\PackageWarning{paper}{\string\fn\space used in the main paper}\fi
  \texttt{\detokenize{#1}}}
\newcommand{\key}[1]{\texttt{\detokenize{#1}}}

% ---- cross references between main paper and supplement ----------------
% Labels carry a unit prefix (section 5.3 of OUTLINE.md and DRAFTER_GUIDE.md):
%   main paper:  intro: obj: s3: s4: s5: s6: s7: s8: s9:   (kind: thm- lem- eq- tab- fig- ...)
%   supplement:  pA: pB: ... pK:
% so all labels are globally unique and the prefixes below are only a second guard.
% \Xref{label}  : reference to the OTHER document (hyperlinked via xr-hyper)
% \Lref{label}  : reference in the SAME document (plain \ref)
\ifsupplement
  \externaldocument[M-]{main}[main.pdf]
  \newcommand{\Xref}[1]{\ref{M-#1}}
\else
  \externaldocument[S-]{supplement}[supplement.pdf]
  \newcommand{\Xref}[1]{\ref{S-#1}}
\fi
\newcommand{\Lref}[1]{\ref{#1}}
% \Supp{pE}{lem-chart} -> "Lemma E.3" style text is produced by the drafter:
%   "Lemma~\Xref{pE:lem-chart} of the supplement".  \SuppPart{E} gives "Part E".
\newcommand{\SuppPart}[1]{Part~#1 of the supplement}

% ---- numbers (generated by unit U0; macros only, no literal certified numbers) ----
\IfFileExists{numbers/numbers_paper.tex}%
  {\input{numbers/numbers_paper.tex}}%
  {\PackageWarning{paper}{numbers/numbers_paper.tex not found (unit U0 not finished)}}

% ======================================================================
% NOTATION MACROS (OUTLINE 5.2 and notation.md).  One macro per symbol
% whose spelling is fixed by the notation table.  Plain letters (a, b, f, ...)
% have no macro.
% ======================================================================
% basic
\newcommand{\R}{\mathbb{R}}
\newcommand{\C}{\mathbb{C}}
\newcommand{\N}{\mathbb{N}}
\newcommand{\Z}{\mathbb{Z}}
\newcommand{\eps}{\varepsilon}
\newcommand{\ip}[2]{\langle #1,#2\rangle}
\newcommand{\abs}[1]{\lvert #1\rvert}
\newcommand{\norm}[1]{\lVert #1\rVert}
\newcommand{\dd}{\mathrm d}
\newcommand{\re}{\operatorname{Re}}
\newcommand{\im}{\operatorname{Im}}
% flow and curves
\newcommand{\arclen}{\mathfrak{s}}                 % arclength of the flow curve
\newcommand{\kss}{\kappa_{\mathfrak{ss}}}          % kappa_{ss} in the flow equation
\newcommand{\Area}{\operatorname{Area}}
\newcommand{\Len}{\mathcal L}                      % length of a curve
\newcommand{\spp}{\operatorname{sp}}                 % sup quantities sp^v_0, sp^v_1, sp^f_1 (were S^v_0, S^v_1, S^f_1)
\newcommand{\rot}{\operatorname{rot}}              % turning number
\newcommand{\Tsing}{T_*}                           % singular time
\newcommand{\Tmax}{T_{\max}}                       % maximal time of existence
\newcommand{\pstar}{p_*}                           % image point of the singularity
\newcommand{\Gz}{\Gamma_0}                         % initial curve
\newcommand{\vth}{\vartheta}                       % parameter of Gamma_0
\newcommand{\dclk}{\delta_t}                       % clock defect of step 4
\newcommand{\fall}{\varpi}                         % Fourier rates varpi_n = n^4-n^2
% rescaling
\newcommand{\lamone}{\lambda_1}
\newcommand{\rate}{\varsigma}                      % 1/4 - a
\newcommand{\xia}{\xi_a}                           % (1-4\bar a)^{-1}
% shrinker and step 1
\newcommand{\Sone}{S_1}
\newcommand{\Gsh}{\Gamma_{\!*}}                    % profile curve of S_1
\newcommand{\cSet}{\mathcal S}                     % certified set of shrinkers
\newcommand{\sigtip}{\sigma_{\rm tip}}
\newcommand{\Xo}{X_0}                              % 2.2490776515
\newcommand{\Dtip}{D}
\newcommand{\ktip}{k_0}
\newcommand{\shootM}{M}                            % shooting function
\newcommand{\Ypar}{\mathcal Y}                     % step-1 parallelepiped
\newcommand{\axis}{\Sigma}                         % symmetry axis of S_1
\newcommand{\sigc}{\sigma_{\rm c}}
\newcommand{\sigrev}{\sigma^{\rm rev}}
% linear theory
\newcommand{\Hrho}{\mathcal H_\rho}
\newcommand{\muslow}{\mu_3}
\newcommand{\rre}{\mathfrak r_{\rm re}}
\newcommand{\rim}{\mathfrak r_{\rm im}}
\newcommand{\Evev}{\mathrm{Ev}_{\rm even}}
\newcommand{\Evod}{\mathrm{Ev}_{\rm odd}}
\newcommand{\gvec}{\mathbf g}                      % translation vector in step 2
\newcommand{\cons}[1]{\mathfrak f_{#1}}            % constraint functions
\newcommand{\phirho}{\varphi_\rho}
% energies
\newcommand{\Lam}{\Lambda}
\newcommand{\En}{E}
\newcommand{\EnI}{E_I}
\newcommand{\muI}{\mu_I}                           % 0.08
\newcommand{\epsE}{\varepsilon_E}                  % 1e-5
\newcommand{\epsI}{\varepsilon_I}                  % 3e-5
\newcommand{\wn}[1]{\mathsf w_{#1}}                % ||w~^{(j)}||_rho
\newcommand{\vI}{v_I}                              % zeta w'
\newcommand{\muerr}{\mu_{\rm err}}                 % step-4 error-energy weight 0.036
% far field
\newcommand{\pU}{p_U}                              % 1.0246
\newcommand{\Uone}{U_1}
\newcommand{\sm}{\operatorname{sm}}                % smoothstep t^3(10-15t+6t^2)
\newcommand{\rc}{r_c}                              % centre radius (Lemma U)
\newcommand{\rchi}{r_\chi}                         % 253
\newcommand{\Iw}{I_w}                              % hand-off windows [y_w,1.5y_w]
\newcommand{\qa}{\mathfrak q_1}                    % quasilinear coefficients of the slope equation
\newcommand{\qb}{\mathfrak q_2}
% cylinders (Lemma O)
\newcommand{\cyl}{\ell}                            % cylinder length unit
\newcommand{\MO}{M_O}
\newcommand{\KO}{K_O}
\newcommand{\Hcyl}{H_{\rm cyl}}
\newcommand{\XiO}{\Xi_O}
\newcommand{\fO}{f_O}
% step 4
\newcommand{\Ybar}{\bar Y}
\newcommand{\Lt}{\widetilde L}                     % linearisation about the step-4 base
\newcommand{\Trans}{\Pi}                           % transport integral
\newcommand{\resid}{\mathrm{res}}
\newcommand{\Resid}[1]{\mathrm{Res}_{#1}}
\newcommand{\boxtip}{\Box_{\rm tip}}
\newcommand{\boxfar}{\Box_{\rm far}}
\newcommand{\Eenv}{\mathcal E_{\rm env}}
\newcommand{\betab}{\beta_{\rm base}}
\newcommand{\tauI}{\tau^{\rm I}}                   % Phase I clock
% names of families of conditions: \Hh{3} = (H3), \Oh{4} = (O4), ...
\newcommand{\Hh}[1]{\textup{(H#1)}}
\newcommand{\Oh}[1]{\textup{(O#1)}}
\newcommand{\Lh}[1]{\textup{(L#1)}}
\newcommand{\Bh}[1]{\textup{(B#1)}}
\newcommand{\Kh}[1]{\textup{(K#1)}}
\newcommand{\Wh}[1]{\textup{(W#1)}}
\newcommand{\PhaseI}{Phase~I}
\newcommand{\PhaseII}{Phase~II}
% margin: threshold/bound (always)
\newcommand{\mgn}[1]{{\times}#1}

% ---- figures: placeholder macro so that drafts compile before F1-F7 exist ----
\newcommand{\figplaceholder}[1]{\fbox{\parbox{0.8\linewidth}{\centering\small[figure #1 to be supplied]}}}
 ...
% (the main paper does nothing special: \ifsupplement defaults to false).
% Compile from paper_v1/ so that % numbers_paper.tex -- GENERATED from ledger.json by build_ledger.py; do not edit.
% Rounding: upper bounds up, lower bounds and margins down at the printed digit; margin = threshold / bound.
% Use as $\times\NName$ for margins.  kind: cert certified, est estimate, hand hypothesis constant, doc only a document quotes it.
% NUEhatZero [cert, upper] U_cert_v35.json :: supersolution.EA.0
\newcommand{\NUEhatZero}{\ensuremath{1.2018\cdot10^{-6}}}
% NUEhatOne [cert, upper] U_cert_v35.json :: supersolution.EA.1
\newcommand{\NUEhatOne}{\ensuremath{1.0441\cdot10^{-4}}}
% NUEhatTwo [cert, upper] U_cert_v35.json :: supersolution.EA.2
\newcommand{\NUEhatTwo}{\ensuremath{0.056694}}
% NUBhatZero [cert, upper] U_cert_v35.json :: supersolution.Bhat.0
\newcommand{\NUBhatZero}{\ensuremath{4.8735\cdot10^{-6}}}
% NUBhatOne [cert, upper] U_cert_v35.json :: supersolution.Bhat.1
\newcommand{\NUBhatOne}{\ensuremath{5.5933\cdot10^{-4}}}
% NUBhatTwo [cert, upper] U_cert_v35.json :: supersolution.Bhat.2
\newcommand{\NUBhatTwo}{\ensuremath{0.36675}}
% NUSupersolutionCZeroZero [cert, upper] U_cert_v35.json :: supersolution.C0.0
\newcommand{\NUSupersolutionCZeroZero}{\ensuremath{1.1340\cdot10^{-6}}}
% NUSupersolutionCZeroOne [cert, upper] U_cert_v35.json :: supersolution.C0.1
\newcommand{\NUSupersolutionCZeroOne}{\ensuremath{1.6623\cdot10^{-4}}}
% NUSupersolutionCOneZero [cert, upper] U_cert_v35.json :: supersolution.C1.0
\newcommand{\NUSupersolutionCOneZero}{\ensuremath{1.5241\cdot10^{-9}}}
% NUSupersolutionCOneOne [cert, upper] U_cert_v35.json :: supersolution.C1.1
\newcommand{\NUSupersolutionCOneOne}{\ensuremath{5.5255\cdot10^{-8}}}
% NUSupersolutionNTwoZero [cert, upper] U_cert_v35.json :: supersolution.N2.0
\newcommand{\NUSupersolutionNTwoZero}{\ensuremath{0.071472}}
% NUSupersolutionNTwoOne [cert, upper] U_cert_v35.json :: supersolution.N2.1
\newcommand{\NUSupersolutionNTwoOne}{\ensuremath{0.33012}}
% NUSupersolutionNTwoTwo [cert, upper] U_cert_v35.json :: supersolution.N2.2
\newcommand{\NUSupersolutionNTwoTwo}{\ensuremath{5.0406}}
% NUSupersolutionNTwoOverFourNTwoRefZero [cert, upper] U_cert_v35.json :: supersolution.N2_over_4N2ref.0
\newcommand{\NUSupersolutionNTwoOverFourNTwoRefZero}{\ensuremath{0.18}}
% NUSupersolutionNTwoOverFourNTwoRefOne [cert, upper] U_cert_v35.json :: supersolution.N2_over_4N2ref.1
\newcommand{\NUSupersolutionNTwoOverFourNTwoRefOne}{\ensuremath{0.42}}
% NUSupersolutionNTwoOverFourNTwoRefTwo [cert, upper] U_cert_v35.json :: supersolution.N2_over_4N2ref.2
\newcommand{\NUSupersolutionNTwoOverFourNTwoRefTwo}{\ensuremath{0.26}}
% NUSupersolutionMOneZHypCheckZonesIIISOneTwoBhat [cert, upper] U_cert_v35.json :: supersolution.M1z_hyp_check.zones_I_II.S1_2Bhat
\newcommand{\NUSupersolutionMOneZHypCheckZonesIIISOneTwoBhat}{\ensuremath{0.28585}}
% NUSupersolutionMOneZHypCheckZoneIIIBOnePlusPUOne [cert, upper] U_cert_v35.json :: supersolution.M1z_hyp_check.zone_III.b1_plus_pU1
\newcommand{\NUSupersolutionMOneZHypCheckZoneIIIBOnePlusPUOne}{\ensuremath{0.50786}}
% NUMarginLateMOne [cert, margin] U_cert_v35.json :: margins.late_M1
\newcommand{\NUMarginLateMOne}{\ensuremath{2.353}}
% NUMarginLateMZero [cert, margin] U_cert_v35.json :: margins.late_M0
\newcommand{\NUMarginLateMZero}{\ensuremath{4.892}}
% NUMarginsLateMZeroZ [cert, margin] U_cert_v35.json :: margins.late_M0z
\newcommand{\NUMarginsLateMZeroZ}{\ensuremath{12.99}}
% NUMarginPresentMOne [cert, margin] U_cert_v35.json :: margins.present_M1_min.value
\newcommand{\NUMarginPresentMOne}{\ensuremath{2.770}}
% NUMarginPresentMZero [cert, margin] U_cert_v35.json :: margins.present_M0_min.value
\newcommand{\NUMarginPresentMZero}{\ensuremath{2.225}}
% NUMarginsPresentMZeroZMinValue [cert, margin] U_cert_v35.json :: margins.present_M0z_min.value
\newcommand{\NUMarginsPresentMZeroZMinValue}{\ensuremath{14.37}}
% NULoopGainZeroTotal [cert, upper] U_cert_v35.json :: loop_gain.0.total
\newcommand{\NULoopGainZeroTotal}{\ensuremath{1.23\cdot10^{-7}}}
% NULoopGainOneTotal [cert, upper] U_cert_v35.json :: loop_gain.1.total
\newcommand{\NULoopGainOneTotal}{\ensuremath{0.037169}}
% NULoopGainTwoNeighbour [cert, upper] U_cert_v35.json :: loop_gain.2.neighbour
\newcommand{\NULoopGainTwoNeighbour}{\ensuremath{0.037018}}
% NULoopGainTwoThroughA [cert, upper] U_cert_v35.json :: loop_gain.2.through_A
\newcommand{\NULoopGainTwoThroughA}{\ensuremath{0.0020172}}
% NUfineBhatZero [cert, upper] U_cert_v35.json :: ufine.Bhat.0
\newcommand{\NUfineBhatZero}{\ensuremath{4.8723\cdot10^{-6}}}
% NUfineBhatOne [cert, upper] U_cert_v35.json :: ufine.Bhat.1
\newcommand{\NUfineBhatOne}{\ensuremath{2.6266\cdot10^{-4}}}
% NUfineBhatTwo [cert, upper] U_cert_v35.json :: ufine.Bhat.2
\newcommand{\NUfineBhatTwo}{\ensuremath{0.025842}}
% NUcdeltaU [cert, upper] U_cert_v35.json :: launch.c_dU
\newcommand{\NUcdeltaU}{\ensuremath{1.7187\cdot10^{-8}}}
% NULaunchWmin [cert, lower] U_cert_v35.json :: launch.wmin
\newcommand{\NULaunchWmin}{\ensuremath{0.73492}}
% NULaunchCAD [cert, upper] U_cert_v35.json :: launch.CA_D
\newcommand{\NULaunchCAD}{\ensuremath{2.2429}}
% NULaunchQD [cert, upper] U_cert_v35.json :: launch.qD
\newcommand{\NULaunchQD}{\ensuremath{0.018509}}
% NULaunchDeltaInZero [cert, upper] U_cert_v35.json :: launch.Delta_in.0
\newcommand{\NULaunchDeltaInZero}{\ensuremath{3.131\cdot10^{-12}}}
% NULaunchDeltaInOne [cert, upper] U_cert_v35.json :: launch.Delta_in.1
\newcommand{\NULaunchDeltaInOne}{\ensuremath{3.287\cdot10^{-12}}}
% NULaunchDeltaInTwo [cert, upper] U_cert_v35.json :: launch.Delta_in.2
\newcommand{\NULaunchDeltaInTwo}{\ensuremath{3.451\cdot10^{-12}}}
% NULaunchPhiZero [cert, upper] U_cert_v35.json :: launch.Phi.0
\newcommand{\NULaunchPhiZero}{\ensuremath{6.434\cdot10^{-9}}}
% NULaunchPhiOne [cert, upper] U_cert_v35.json :: launch.Phi.1
\newcommand{\NULaunchPhiOne}{\ensuremath{1.4217\cdot10^{-8}}}
% NULaunchPhiTwo [cert, upper] U_cert_v35.json :: launch.Phi.2
\newcommand{\NULaunchPhiTwo}{\ensuremath{3.8404\cdot10^{-8}}}
% NUGammaZero [cert, upper] U_cert_v35.json :: Gamma.0
\newcommand{\NUGammaZero}{\ensuremath{4.0000358}}
% NUGammaOne [cert, upper] U_cert_v35.json :: Gamma.1
\newcommand{\NUGammaOne}{\ensuremath{16.000322}}
% NUGammaTwo [cert, upper] U_cert_v35.json :: Gamma.2
\newcommand{\NUGammaTwo}{\ensuremath{64.00218}}
% NUKappa [cert, upper] U_cert_v35.json :: kappa
\newcommand{\NUKappa}{\ensuremath{8.0000358}}
% NUKappaTwo [cert, upper] U_cert_v35.json :: kappa2
\newcommand{\NUKappaTwo}{\ensuremath{64.00061}}
% NUGhatZeroZero [cert, upper] U_cert_v35.json :: Ghat.0.0
\newcommand{\NUGhatZeroZero}{\ensuremath{51322.6}}
% NUGhatOneOne [cert, upper] U_cert_v35.json :: Ghat.1.1
\newcommand{\NUGhatOneOne}{\ensuremath{33962.9}}
% NUGhatTwoTwo [cert, upper] U_cert_v35.json :: Ghat.2.2
\newcommand{\NUGhatTwoTwo}{\ensuremath{348341.1}}
% NUGhatOneZero [cert, upper] U_cert_v35.json :: Ghat.1.0
\newcommand{\NUGhatOneZero}{\ensuremath{4.938}}
% NUGhatTwoZero [cert, upper] U_cert_v35.json :: Ghat.2.0
\newcommand{\NUGhatTwoZero}{\ensuremath{27.984}}
% NUGhatTwoOne [cert, upper] U_cert_v35.json :: Ghat.2.1
\newcommand{\NUGhatTwoOne}{\ensuremath{458.53}}
% NUKZeroZero [cert, upper] U_cert_v35.json :: K.0.0
\newcommand{\NUKZeroZero}{\ensuremath{8034.21}}
% NUKZeroOne [cert, upper] U_cert_v35.json :: K.0.1
\newcommand{\NUKZeroOne}{\ensuremath{724.89}}
% NUKOneZero [cert, upper] U_cert_v35.json :: K.1.0
\newcommand{\NUKOneZero}{\ensuremath{0.08152}}
% NUKOneOne [cert, upper] U_cert_v35.json :: K.1.1
\newcommand{\NUKOneOne}{\ensuremath{6336.94}}
% NUKOneTwo [cert, upper] U_cert_v35.json :: K.1.2
\newcommand{\NUKOneTwo}{\ensuremath{545.59}}
% NUKOneThree [cert, upper] U_cert_v35.json :: K.1.3
\newcommand{\NUKOneThree}{\ensuremath{0.58700}}
% NUKTwoZero [cert, upper] U_cert_v35.json :: K.2.0
\newcommand{\NUKTwoZero}{\ensuremath{0.87186}}
% NUKTwoOne [cert, upper] U_cert_v35.json :: K.2.1
\newcommand{\NUKTwoOne}{\ensuremath{528.62}}
% NUKTwoTwo [cert, upper] U_cert_v35.json :: K.2.2
\newcommand{\NUKTwoTwo}{\ensuremath{23425.98}}
% NUKTwoThree [cert, upper] U_cert_v35.json :: K.2.3
\newcommand{\NUKTwoThree}{\ensuremath{555.62}}
% NUFootwOne [cert, upper] U_cert_v35.json :: footw.1
\newcommand{\NUFootwOne}{\ensuremath{2.6361}}
% NUFootwTwo [cert, upper] U_cert_v35.json :: footw.2
\newcommand{\NUFootwTwo}{\ensuremath{2663.3}}
% NUCThreeThree [cert, upper] U_cert_v35.json :: c33
\newcommand{\NUCThreeThree}{\ensuremath{108.02}}
% NUCTwoTwoThree [cert, upper] U_cert_v35.json :: c223
\newcommand{\NUCTwoTwoThree}{\ensuremath{7.3284}}
% NUCEta [cert, upper] U_cert_v35.json :: C_eta
\newcommand{\NUCEta}{\ensuremath{13.376}}
% NULMax [cert, upper] U_cert_v35.json :: L_max
\newcommand{\NULMax}{\ensuremath{3.5352}}
% NURRatio [cert, upper] U_cert_v35.json :: R_ratio
\newcommand{\NURRatio}{\ensuremath{11.593}}
% NUSweepsSigmaInCoverSup [cert, upper] U_cert_v35.json :: sweeps.Sigma_in.cover_sup
\newcommand{\NUSweepsSigmaInCoverSup}{\ensuremath{5.0892}}
% NUSweepsSigmaECoverSup [cert, upper] U_cert_v35.json :: sweeps.Sigma_E.cover_sup
\newcommand{\NUSweepsSigmaECoverSup}{\ensuremath{0.30846}}
% NUSweepsSXCoverSup [cert, upper] U_cert_v35.json :: sweeps.S_X.cover_sup
\newcommand{\NUSweepsSXCoverSup}{\ensuremath{1.5243}}
% NUSweepsSigmaXCoverSup [cert, upper] U_cert_v35.json :: sweeps.Sigma_X.cover_sup
\newcommand{\NUSweepsSigmaXCoverSup}{\ensuremath{3.0081}}
% NUSweepsSigmaRsTwoCoverSup [cert, upper] U_cert_v35.json :: sweeps.Sigma_rs.2.cover_sup
\newcommand{\NUSweepsSigmaRsTwoCoverSup}{\ensuremath{4.3350\cdot10^{-3}}}
% NUDataTermsLateDZero [cert, upper] U_cert_v35.json :: data_terms_late.D.0
\newcommand{\NUDataTermsLateDZero}{\ensuremath{1.7330\cdot10^{-6}}}
% NUDataTermsLateDOne [cert, upper] U_cert_v35.json :: data_terms_late.D.1
\newcommand{\NUDataTermsLateDOne}{\ensuremath{2.5353\cdot10^{-4}}}
% NUDataTermsLateSlopeShareCOneSigmaXOverDkZero [cert, upper] U_cert_v35.json :: data_terms_late.slope_share_C1SigmaX_over_Dk.0
\newcommand{\NUDataTermsLateSlopeShareCOneSigmaXOverDkZero}{\ensuremath{2.7\cdot10^{-3}}}
% NUDataTermsLateSlopeShareCOneSigmaXOverDkOne [cert, upper] U_cert_v35.json :: data_terms_late.slope_share_C1SigmaX_over_Dk.1
\newcommand{\NUDataTermsLateSlopeShareCOneSigmaXOverDkOne}{\ensuremath{6.6\cdot10^{-4}}}
% NURLtTwoMarginZero [cert, margin] U_cert_v35.json :: R_lt2_margin.0
\newcommand{\NURLtTwoMarginZero}{\ensuremath{4.432}}
% NURLtTwoMarginOne [cert, margin] U_cert_v35.json :: R_lt2_margin.1
\newcommand{\NURLtTwoMarginOne}{\ensuremath{6.697}}
% NUThetaTwoEight [cert, upper] U_cert_v35.json :: theta28
\newcommand{\NUThetaTwoEight}{\ensuremath{0.37011}}
% NUDFiveSix [cert, upper] U_cert_v35.json :: D56
\newcommand{\NUDFiveSix}{\ensuremath{1.0368}}
% NUDTwoEight [cert, upper] U_cert_v35.json :: D28
\newcommand{\NUDTwoEight}{\ensuremath{1.5876}}
% NUSThreeThreeFootOverBhatTwo [cert, upper] U_cert_v35.json :: S33_foot_over_Bhat2
\newcommand{\NUSThreeThreeFootOverBhatTwo}{\ensuremath{1.69\cdot10^{-3}}}
% NSAHJWKappaJ [cert, upper] A_interface.json :: HJW.kappa_J
\newcommand{\NSAHJWKappaJ}{\ensuremath{0.042351}}
% NSAGpZero [cert, lower] A_interface.json :: res.tight.gp.0
\newcommand{\NSAGpZero}{\ensuremath{0.24292}}
% NSAGpOne [cert, lower] A_interface.json :: res.tight.gp.1
\newcommand{\NSAGpOne}{\ensuremath{0.24106}}
% NSAGpTwo [cert, lower] A_interface.json :: res.tight.gp.2
\newcommand{\NSAGpTwo}{\ensuremath{0.11961}}
% NSAAminC [cert, lower] A_interface.json :: res.tight.amin_c
\newcommand{\NSAAminC}{\ensuremath{0.97623}}
% NSARhoJ [cert, upper] A_interface.json :: res.tight.rhoJ
\newcommand{\NSARhoJ}{\ensuremath{5.4995\cdot10^{-7}}}
% NSAKappaOTwoTwo [cert, upper] A_interface.json :: res.tight.kappa_o.2.2
\newcommand{\NSAKappaOTwoTwo}{\ensuremath{1.2416\cdot10^{-5}}}
% NSADsrcZero [cert, upper] A_interface.json :: res.tight.Dsrc.0
\newcommand{\NSADsrcZero}{\ensuremath{1.8092\cdot10^{-7}}}
% NSADsrcOne [cert, upper] A_interface.json :: res.tight.Dsrc.1
\newcommand{\NSADsrcOne}{\ensuremath{9.8276\cdot10^{-5}}}
% NSABZeroDWindowSlopeZero [cert, upper] A_interface.json :: res.tight.B0.D_window_slope.0
\newcommand{\NSABZeroDWindowSlopeZero}{\ensuremath{8.7399\cdot10^{-10}}}
% NSABZeroDWindowSlopeOne [cert, upper] A_interface.json :: res.tight.B0.D_window_slope.1
\newcommand{\NSABZeroDWindowSlopeOne}{\ensuremath{5.2599\cdot10^{-8}}}
% NSAFTwoZero [cert, upper] A_interface.json :: res.tight.F2.0
\newcommand{\NSAFTwoZero}{\ensuremath{5.6506\cdot10^{-13}}}
% NSAFTwoOne [cert, upper] A_interface.json :: res.tight.F2.1
\newcommand{\NSAFTwoOne}{\ensuremath{5.9320\cdot10^{-13}}}
% NSACoupOne [cert, upper] A_interface.json :: res.tight.coup.1
\newcommand{\NSACoupOne}{\ensuremath{1.1568\cdot10^{-12}}}
% NSACOneZeroZero [cert, upper] A_interface.json :: res.tight.C1_0.0
\newcommand{\NSACOneZeroZero}{\ensuremath{7.2060\cdot10^{-10}}}
% NSACOneZeroOne [cert, upper] A_interface.json :: res.tight.C1_0.1
\newcommand{\NSACOneZeroOne}{\ensuremath{4.3143\cdot10^{-8}}}
% NSABZeroSupCoreZero [cert, upper] A_interface.json :: res.tight.B0.sup_core.0
\newcommand{\NSABZeroSupCoreZero}{\ensuremath{4.23\cdot10^{-5}}}
% NSACoreFluxMarginTwoFLUXZeroOverKZero [cert, margin] A_interface.json :: res.tight.core.flux_margin.2.FLUX0_over_K0
\newcommand{\NSACoreFluxMarginTwoFLUXZeroOverKZero}{\ensuremath{3.396}}
% NSACoreFluxMarginTwoFLUXTwoOverKTwo [cert, margin] A_interface.json :: res.tight.core.flux_margin.2.FLUX2_over_K2
\newcommand{\NSACoreFluxMarginTwoFLUXTwoOverKTwo}{\ensuremath{6.012}}
% NSAMarginFluxZero [cert, margin] asm_design_N2ref4.json :: res.tight.flux.beta.2.margin_FLUX0
\newcommand{\NSAMarginFluxZero}{\ensuremath{3.365}}
% NSAMarginFluxTwo [cert, margin] asm_design_N2ref4.json :: res.tight.flux.beta.2.margin_FLUX2
\newcommand{\NSAMarginFluxTwo}{\ensuremath{5.955}}
% NSAfourSZero [cert, upper] asm_design_N2ref4.json :: res.tight.S0
\newcommand{\NSAfourSZero}{\ensuremath{7.1878\cdot10^{-3}}}
% NSAfourSOne [cert, upper] asm_design_N2ref4.json :: res.tight.S1
\newcommand{\NSAfourSOne}{\ensuremath{0.10430}}
% NSAMarginCoreSZeroV [cert, margin] asm_design_N2ref4.json :: res.tight.margins.S0v_vs_0.0153
\newcommand{\NSAMarginCoreSZeroV}{\ensuremath{2.116}}
% NSAMarginCoreSOneV [cert, margin] asm_design_N2ref4.json :: res.tight.margins.S1v_vs_0.6
\newcommand{\NSAMarginCoreSOneV}{\ensuremath{5.748}}
% NSAMarginCoreMOne [cert, margin] asm_design_N2ref4.json :: res.tight.margins.M1core
\newcommand{\NSAMarginCoreMOne}{\ensuremath{4.409}}
% NSAMarginCoreMZero [cert, margin] asm_design_N2ref4.json :: res.tight.margins.M0core
\newcommand{\NSAMarginCoreMZero}{\ensuremath{5.450}}
% NSAMarginGhat [cert, margin] asm_design_N2ref4.json :: res.tight.margins.Ghat22
\newcommand{\NSAMarginGhat}{\ensuremath{1.917}}
% NSAMarginRampMOne [cert, margin] asm_design_N2ref4.json :: res.tight.margins.ramp M1
\newcommand{\NSAMarginRampMOne}{\ensuremath{1.184}}
% NSAfourInfoZeroDWindowConst [cert, upper] asm_design_N2ref4.json :: res.tight.info.0.D_window_const
\newcommand{\NSAfourInfoZeroDWindowConst}{\ensuremath{1.7732\cdot10^{-7}}}
% NSAfourInfoOneDWindowConst [cert, upper] asm_design_N2ref4.json :: res.tight.info.1.D_window_const
\newcommand{\NSAfourInfoOneDWindowConst}{\ensuremath{9.8058\cdot10^{-5}}}
% NZDeltaRZero [cert, upper] zone3_interface.json :: DeltaR.0
\newcommand{\NZDeltaRZero}{\ensuremath{4.4986\cdot10^{-8}}}
% NZDeltaROne [cert, upper] zone3_interface.json :: DeltaR.1
\newcommand{\NZDeltaROne}{\ensuremath{6.7972\cdot10^{-6}}}
% NZDeltaRTwo [cert, upper] zone3_interface.json :: DeltaR.2
\newcommand{\NZDeltaRTwo}{\ensuremath{0.011480}}
% NZDeltaRBZero [cert, upper] zone3_interface.json :: DeltaR_b.0
\newcommand{\NZDeltaRBZero}{\ensuremath{2.66\cdot10^{-12}}}
% NZDeltaRBOne [cert, upper] zone3_interface.json :: DeltaR_b.1
\newcommand{\NZDeltaRBOne}{\ensuremath{3.26\cdot10^{-9}}}
% NZDeltaRBTwo [cert, upper] zone3_interface.json :: DeltaR_b.2
\newcommand{\NZDeltaRBTwo}{\ensuremath{4.35\cdot10^{-8}}}
% NZDeltaRDetailNSqZero [cert, upper] zone3_interface.json :: DeltaR_detail.N_sq.0
\newcommand{\NZDeltaRDetailNSqZero}{\ensuremath{8.28\cdot10^{-5}}}
% NZDeltaRDetailNSqOne [cert, upper] zone3_interface.json :: DeltaR_detail.N_sq.1
\newcommand{\NZDeltaRDetailNSqOne}{\ensuremath{0.10149}}
% NZDeltaRDetailNSqTwo [cert, upper] zone3_interface.json :: DeltaR_detail.N_sq.2
\newcommand{\NZDeltaRDetailNSqTwo}{\ensuremath{1.3567}}
% NZDeltaRDetailNSqThree [cert, upper] zone3_interface.json :: DeltaR_detail.N_sq.3
\newcommand{\NZDeltaRDetailNSqThree}{\ensuremath{25350}}
% NZDeltaRDetailPartsTwoPairsThree [cert, upper] zone3_interface.json :: DeltaR_detail.parts.2.pairs.3
\newcommand{\NZDeltaRDetailPartsTwoPairsThree}{\ensuremath{0.011298}}
% NZDeltaRDetailPartsTwoLFourFoot [cert, upper] zone3_interface.json :: DeltaR_detail.parts.2.L4_foot
\newcommand{\NZDeltaRDetailPartsTwoLFourFoot}{\ensuremath{1.6\cdot10^{-4}}}
% NZDeltaRDetailPartsTwoCubic [cert, upper] zone3_interface.json :: DeltaR_detail.parts.2.cubic
\newcommand{\NZDeltaRDetailPartsTwoCubic}{\ensuremath{9.3\cdot10^{-6}}}
% NZDeltaRDetailPartsTwoSThreeThreeFoot [cert, upper] zone3_interface.json :: DeltaR_detail.parts.2.S33_foot
\newcommand{\NZDeltaRDetailPartsTwoSThreeThreeFoot}{\ensuremath{3.7\cdot10^{-7}}}
% NZNIIIZero [cert, upper] zone3_interface.json :: n_III.0
\newcommand{\NZNIIIZero}{\ensuremath{0.071445}}
% NZNIIIOne [cert, upper] zone3_interface.json :: n_III.1
\newcommand{\NZNIIIOne}{\ensuremath{0.32707}}
% NZNIIITwo [cert, upper] zone3_interface.json :: n_III.2
\newcommand{\NZNIIITwo}{\ensuremath{3.0445}}
% NZTauIII [cert, upper] zone3_interface.json :: tau_III
\newcommand{\NZTauIII}{\ensuremath{4.0708}}
% NZoneIIIBoxZero [cert, upper] zone3_interface.json :: box_III_380.b0
\newcommand{\NZoneIIIBoxZero}{\ensuremath{0.10032}}
% NZoneIIIBoxOne [cert, upper] zone3_interface.json :: box_III_380.b1
\newcommand{\NZoneIIIBoxOne}{\ensuremath{0.41974}}
% NZoneIIIBoxZ [cert, upper] zone3_interface.json :: box_III_380.bz
\newcommand{\NZoneIIIBoxZ}{\ensuremath{0.014291}}
% NZBoxIIIThreeEightZeroMarginsBZero [cert, margin] zone3_interface.json :: box_III_380.margins.b0
\newcommand{\NZBoxIIIThreeEightZeroMarginsBZero}{\ensuremath{1.495}}
% NZoneIIIMarginThinnest [cert, margin] zone3_interface.json :: box_III_380.margins.b1
\newcommand{\NZoneIIIMarginThinnest}{\ensuremath{1.238}}
% NZBoxIIIThreeEightZeroMarginsBz [cert, margin] zone3_interface.json :: box_III_380.margins.bz
\newcommand{\NZBoxIIIThreeEightZeroMarginsBz}{\ensuremath{11.19}}
% NZBoxIIISevenZeroZeroBZero [cert, upper] zone3_interface.json :: box_III_700.b0
\newcommand{\NZBoxIIISevenZeroZeroBZero}{\ensuremath{0.22757}}
% NZBoxIIISevenZeroZeroBOne [cert, upper] zone3_interface.json :: box_III_700.b1
\newcommand{\NZBoxIIISevenZeroZeroBOne}{\ensuremath{0.54234}}
% NZUfineKTwoAtBhatRef [cert, upper] zone3_interface.json :: ufine.K2_at_Bhat_ref
\newcommand{\NZUfineKTwoAtBhatRef}{\ensuremath{0.046956}}
% NZUfineKThreeStripAtBhatRef [cert, upper] zone3_interface.json :: ufine.K3_strip_at_Bhat_ref
\newcommand{\NZUfineKThreeStripAtBhatRef}{\ensuremath{0.054140}}
% NRecoveryMinMargin [cert, margin] zone3_interface.json :: cone_family.H6.at_Bhat_ref.min_margin
\newcommand{\NRecoveryMinMargin}{\ensuremath{1.860}}
% NZConeFamilyHSixAtBhatRefConeSlopeMargin [cert, margin] zone3_interface.json :: cone_family.H6.at_Bhat_ref.cone.slope.margin
\newcommand{\NZConeFamilyHSixAtBhatRefConeSlopeMargin}{\ensuremath{2.46}}
% NZConeFamilyHSixAtBhatRefConeKTwoMargin [cert, margin] zone3_interface.json :: cone_family.H6.at_Bhat_ref.cone.K2.margin
\newcommand{\NZConeFamilyHSixAtBhatRefConeKTwoMargin}{\ensuremath{2.54}}
% NZConeFamilyHSixAtBhatRefConeKThreeMargin [cert, margin] zone3_interface.json :: cone_family.H6.at_Bhat_ref.cone.K3.margin
\newcommand{\NZConeFamilyHSixAtBhatRefConeKThreeMargin}{\ensuremath{2.51}}
% NConeTstarMargin [cert, margin] zone3_interface.json :: acceptance.tstar_margin
\newcommand{\NConeTstarMargin}{\ensuremath{2.30}}
% NUoneTFiveYaddLevelOneTwo [cert, upper] u1_ghat.json :: t5_Yadd.level1.2
\newcommand{\NUoneTFiveYaddLevelOneTwo}{\ensuremath{0.0044}}
% NUoneTFiveYaddLevelTwoTwo [cert, upper] u1_ghat.json :: t5_Yadd.level2.2
\newcommand{\NUoneTFiveYaddLevelTwoTwo}{\ensuremath{0.1081}}
% NUoneTFiveYaddLevelTwoThree [cert, upper] u1_ghat.json :: t5_Yadd.level2.3
\newcommand{\NUoneTFiveYaddLevelTwoThree}{\ensuremath{0.374}}
% NUoneWeightThetaGammaZero [cert, upper] u1_ghat.json :: weight.Theta.Gamma.0
\newcommand{\NUoneWeightThetaGammaZero}{\ensuremath{4.00000000000000000000054}}
% NRKZeroZero [cert, upper] ramp_U.json :: K.0.0
\newcommand{\NRKZeroZero}{\ensuremath{13436.03}}
% NRKOneOne [cert, upper] ramp_U.json :: K.1.1
\newcommand{\NRKOneOne}{\ensuremath{10526.19}}
% NRKTwoTwo [cert, upper] ramp_U.json :: K.2.2
\newcommand{\NRKTwoTwo}{\ensuremath{39482.30}}
% NUoneCasesPrimaryYNextZero [cert, upper] u1_ghat.json :: cases.primary.Y_next.0
\newcommand{\NUoneCasesPrimaryYNextZero}{\ensuremath{457.092}}
% NUoneCasesPrimaryYNextOne [cert, upper] u1_ghat.json :: cases.primary.Y_next.1
\newcommand{\NUoneCasesPrimaryYNextOne}{\ensuremath{356.940}}
% NUoneCasesPrimaryYNextTwo [cert, upper] u1_ghat.json :: cases.primary.Y_next.2
\newcommand{\NUoneCasesPrimaryYNextTwo}{\ensuremath{1320.220}}
% NSfourResultsMNineHTwoTightenedFUOneRowsZeroBound [cert, upper] FINAL_manifest_v33.json :: results.M9_H2_tightened_fU1.rows.0.bound
\newcommand{\NSfourResultsMNineHTwoTightenedFUOneRowsZeroBound}{\ensuremath{3.658\cdot10^{-10}}}
% NSfourResultsMNineHTwoTightenedFUOneRowsZeroMargin [cert, margin] FINAL_manifest_v33.json :: results.M9_H2_tightened_fU1.rows.0.margin
\newcommand{\NSfourResultsMNineHTwoTightenedFUOneRowsZeroMargin}{\ensuremath{27.34}}
% NSfourResultsMNineHTwoTightenedFUOneRowsOneBound [cert, upper] FINAL_manifest_v33.json :: results.M9_H2_tightened_fU1.rows.1.bound
\newcommand{\NSfourResultsMNineHTwoTightenedFUOneRowsOneBound}{\ensuremath{1.777\cdot10^{-9}}}
% NSfourResultsMNineHTwoTightenedFUOneRowsOneMargin [cert, margin] FINAL_manifest_v33.json :: results.M9_H2_tightened_fU1.rows.1.margin
\newcommand{\NSfourResultsMNineHTwoTightenedFUOneRowsOneMargin}{\ensuremath{562.8}}
% NSfourResultsMNineHTwoTightenedFUOneRowsTwoBound [cert, upper] FINAL_manifest_v33.json :: results.M9_H2_tightened_fU1.rows.2.bound
\newcommand{\NSfourResultsMNineHTwoTightenedFUOneRowsTwoBound}{\ensuremath{1.665\cdot10^{-6}}}
% NSfourResultsMNineHTwoTightenedFUOneRowsTwoMargin [cert, margin] FINAL_manifest_v33.json :: results.M9_H2_tightened_fU1.rows.2.margin
\newcommand{\NSfourResultsMNineHTwoTightenedFUOneRowsTwoMargin}{\ensuremath{600.9}}
% NHOneAlphaThr [hand, exact] hand :: 
\newcommand{\NHOneAlphaThr}{\ensuremath{1.5\cdot10^{-4}}}
% NHOneSqrtLamThr [hand, exact] hand :: 
\newcommand{\NHOneSqrtLamThr}{\ensuremath{5\cdot10^{-8}}}
% NYbarStep [hand, exact] hand :: 
\newcommand{\NYbarStep}{\ensuremath{1\cdot10^{-12}}}
% NHTwoThrZero [hand, exact] hand :: 
\newcommand{\NHTwoThrZero}{\ensuremath{1\cdot10^{-8}}}
% NHTwoThrOne [hand, exact] hand :: 
\newcommand{\NHTwoThrOne}{\ensuremath{1\cdot10^{-6}}}
% NHTwoThrTwo [hand, exact] hand :: 
\newcommand{\NHTwoThrTwo}{\ensuremath{1\cdot10^{-3}}}
% NHThreeThrZero [hand, exact] hand :: 
\newcommand{\NHThreeThrZero}{\ensuremath{1\cdot10^{-6}}}
% NHThreeThrOne [hand, exact] hand :: 
\newcommand{\NHThreeThrOne}{\ensuremath{5\cdot10^{-6}}}
% NHThreeThrTwo [hand, exact] hand :: 
\newcommand{\NHThreeThrTwo}{\ensuremath{1\cdot10^{-5}}}
% NHFourThrZero [hand, exact] hand :: 
\newcommand{\NHFourThrZero}{\ensuremath{2\cdot10^{-7}}}
% NHFourThrOne [hand, exact] hand :: 
\newcommand{\NHFourThrOne}{\ensuremath{1\cdot10^{-5}}}
% NHFourThrTwo [hand, exact] hand :: 
\newcommand{\NHFourThrTwo}{\ensuremath{1\cdot10^{-5}}}
% NHFourThrThree [hand, exact] hand :: 
\newcommand{\NHFourThrThree}{\ensuremath{1\cdot10^{-5}}}
% NHFourKfourThr [hand, exact] hand :: 
\newcommand{\NHFourKfourThr}{\ensuremath{0.06}}
% NHSixConeSlopeThr [hand, exact] hand :: 
\newcommand{\NHSixConeSlopeThr}{\ensuremath{0.08}}
% NHSixConeKtwoThr [hand, exact] hand :: 
\newcommand{\NHSixConeKtwoThr}{\ensuremath{0.03}}
% NHSixConeNthreeThr [hand, exact] hand :: 
\newcommand{\NHSixConeNthreeThr}{\ensuremath{0.004}}
% NHSixConeJetThr [hand, exact] hand :: 
\newcommand{\NHSixConeJetThr}{\ensuremath{0.12}}
% NHSixConeHFiveThr [hand, exact] hand :: 
\newcommand{\NHSixConeHFiveThr}{\ensuremath{0.0144}}
% NHSixFarHFiveThr [hand, exact] hand :: 
\newcommand{\NHSixFarHFiveThr}{\ensuremath{0.0225}}
% NLemmaOConeM [hand, exact] hand :: 
\newcommand{\NLemmaOConeM}{\ensuremath{0.12}}
% NLemmaOConeK [hand, exact] hand :: 
\newcommand{\NLemmaOConeK}{\ensuremath{0.18}}
% NLemmaOFarM [hand, exact] hand :: 
\newcommand{\NLemmaOFarM}{\ensuremath{0.15}}
% NLemmaOFarK [hand, exact] hand :: 
\newcommand{\NLemmaOFarK}{\ensuremath{0.3}}
% NCoercivityConst [hand, exact] hand :: 
\newcommand{\NCoercivityConst}{\ensuremath{0.046875}}
% NMOnez [hand, exact] hand :: 
\newcommand{\NMOnez}{\ensuremath{0.534}}
% NBetaA [hand, exact] hand :: 
\newcommand{\NBetaA}{\ensuremath{2}}
% NTOneLower [cert, lower] t1clock.json :: t1.0
\newcommand{\NTOneLower}{\ensuremath{4.5978574775}}
% NTOneUpper [cert, upper] t1clock.json :: t1.1
\newcommand{\NTOneUpper}{\ensuremath{4.5978586144}}
% NLambdaOneLower [cert, lower] t1clock.json :: lnlam.0
\newcommand{\NLambdaOneLower}{\ensuremath{6.49\cdot10^{-5}}}
% NLambdaOneUpper [cert, upper] t1clock.json :: lnlam.1
\newcommand{\NLambdaOneUpper}{\ensuremath{6.75\cdot10^{-5}}}
% NTMinusTOneUpper [cert, upper] t1clock.json :: lnlam.1; 1.5e-4.sup_ab
\newcommand{\NTMinusTOneUpper}{\ensuremath{2.1\cdot10^{-17}}}
% NAbarUpper [cert, upper] final_alpha_points_5e-8.json :: 1.5e-4.sup_ab
\newcommand{\NAbarUpper}{\ensuremath{2.19103\cdot10^{-7}}}
% NBetabarUpper [cert, upper] final_alpha_points_5e-8.json :: 1.5e-4.sup_beta
\newcommand{\NBetabarUpper}{\ensuremath{3.33811\cdot10^{-9}}}
% NBaseKmaxLower [cert, lower] t1clock.json :: base_kmax.0
\newcommand{\NBaseKmaxLower}{\ensuremath{2.279\cdot10^{4}}}
% NBaseKmaxUpper [cert, upper] t1clock.json :: base_kmax.1
\newcommand{\NBaseKmaxUpper}{\ensuremath{2.376\cdot10^{4}}}
% NHOneAlphaBase [cert, upper] summary.json :: alpha_abs_max
\newcommand{\NHOneAlphaBase}{\ensuremath{1.1739\cdot10^{-4}}}
% NHOneAlphaErr [cert, upper] t9_h1_clam_s180.json :: C_alpha; inputs.YSTAR
\newcommand{\NHOneAlphaErr}{\ensuremath{7.21\cdot10^{-13}}}
% NHOneAlphaMargin [cert, margin] summary.json :: alpha_abs_max; C_alpha; inputs.YSTAR
\newcommand{\NHOneAlphaMargin}{\ensuremath{1.27}}
% NHOneSqrtLamBase [cert, upper] t9_h1_clam_s180.json :: sqrtLam_base_Wp
\newcommand{\NHOneSqrtLamBase}{\ensuremath{2.282\cdot10^{-8}}}
% NHOneSqrtLamTotal [cert, upper] t9_h1_clam_s180.json :: sqrtLam_total_Wp
\newcommand{\NHOneSqrtLamTotal}{\ensuremath{2.29\cdot10^{-8}}}
% NHOneSqrtLamMargin [cert, margin] t9_h1_clam_s180.json :: sqrtLam_total_Wp
\newcommand{\NHOneSqrtLamMargin}{\ensuremath{2.18}}
% NHOneYbarLimit [cert, lower] t9_h1_clam_s180.json :: Ybar_limit_H1
\newcommand{\NHOneYbarLimit}{\ensuremath{1.12\cdot10^{-9}}}
% NSwitchYUpper [cert, upper] switch_s2.json :: modulated.Y_raw
\newcommand{\NSwitchYUpper}{\ensuremath{6.21\cdot10^{-13}}}
% NSwitchYMargin [cert, margin] switch_s2.json :: modulated.Y_raw
\newcommand{\NSwitchYMargin}{\ensuremath{1.6}}
% NMinCI [cert, lower] constants.json :: meta.summary.min_c_I
\newcommand{\NMinCI}{\ensuremath{0.05283}}
% NNumCertificates [cert, exact] constants.json :: meta.summary.n_cert
\newcommand{\NNumCertificates}{\ensuremath{26}}
% NMinRhoC [cert, lower] constants.json :: meta.closure.rho_c_min
\newcommand{\NMinRhoC}{\ensuremath{0.04711}}
% NHTwoBaseZero [cert, upper] FINAL_manifest_v33.json :: results.M9_H2_tightened.rows.0.base
\newcommand{\NHTwoBaseZero}{\ensuremath{8.011\cdot10^{-11}}}
% NHTwoBaseOne [cert, upper] FINAL_manifest_v33.json :: results.M9_H2_tightened.rows.1.base
\newcommand{\NHTwoBaseOne}{\ensuremath{1.336\cdot10^{-10}}}
% NHTwoBaseTwo [cert, upper] FINAL_manifest_v33.json :: results.M9_H2_tightened.rows.2.base
\newcommand{\NHTwoBaseTwo}{\ensuremath{2.936\cdot10^{-10}}}
% NHTwoErrWZero [cert, upper] FINAL_manifest_v33.json :: results.M9_H2_tightened.rows.0.with_error
\newcommand{\NHTwoErrWZero}{\ensuremath{8.014\cdot10^{-11}}}
% NHTwoMarginWZero [cert, margin] FINAL_manifest_v33.json :: results.M9_H2_tightened.rows.0.margin
\newcommand{\NHTwoMarginWZero}{\ensuremath{124.78}}
% NHTwoErrWOne [cert, upper] FINAL_manifest_v33.json :: results.M9_H2_tightened.rows.1.with_error
\newcommand{\NHTwoErrWOne}{\ensuremath{4.706\cdot10^{-10}}}
% NHTwoMarginWOne [cert, margin] FINAL_manifest_v33.json :: results.M9_H2_tightened.rows.1.margin
\newcommand{\NHTwoMarginWOne}{\ensuremath{2124.9}}
% NHTwoErrWTwo [cert, upper] FINAL_manifest_v33.json :: results.M9_H2_tightened.rows.2.with_error
\newcommand{\NHTwoErrWTwo}{\ensuremath{1.559\cdot10^{-6}}}
% NHTwoMarginWTwo [cert, margin] FINAL_manifest_v33.json :: results.M9_H2_tightened.rows.2.margin
\newcommand{\NHTwoMarginWTwo}{\ensuremath{641.53}}
% NHThreeBaseZero [cert, upper] FINAL_manifest_v33.json :: results.M9_H2f.rows.0.base
\newcommand{\NHThreeBaseZero}{\ensuremath{3.127\cdot10^{-7}}}
% NHThreeBaseOne [cert, upper] FINAL_manifest_v33.json :: results.M9_H2f.rows.1.base
\newcommand{\NHThreeBaseOne}{\ensuremath{1.004\cdot10^{-6}}}
% NHThreeBaseTwo [cert, upper] FINAL_manifest_v33.json :: results.M9_H2f.rows.2.base
\newcommand{\NHThreeBaseTwo}{\ensuremath{2.431\cdot10^{-6}}}
% NHThreeErrZero [cert, upper] FINAL_manifest_v33.json :: results.M9_H2f.rows.0.with_error
\newcommand{\NHThreeErrZero}{\ensuremath{3.1363\cdot10^{-7}}}
% NHThreeMarginZero [cert, margin] FINAL_manifest_v33.json :: results.M9_H2f.rows.0.margin
\newcommand{\NHThreeMarginZero}{\ensuremath{3.1885}}
% NHThreeErrOne [cert, upper] FINAL_manifest_v33.json :: results.M9_H2f.rows.1.with_error
\newcommand{\NHThreeErrOne}{\ensuremath{1.0034\cdot10^{-6}}}
% NHThreeMarginOne [cert, margin] FINAL_manifest_v33.json :: results.M9_H2f.rows.1.margin
\newcommand{\NHThreeMarginOne}{\ensuremath{4.9834}}
% NHThreeErrTwo [cert, upper] FINAL_manifest_v33.json :: results.M9_H2f.rows.2.with_error
\newcommand{\NHThreeErrTwo}{\ensuremath{2.4351\cdot10^{-6}}}
% NHThreeMarginTwo [cert, margin] FINAL_manifest_v33.json :: results.M9_H2f.rows.2.margin
\newcommand{\NHThreeMarginTwo}{\ensuremath{4.1067}}
% NHFourBaseZero [cert, upper] FINAL_manifest_v33.json :: results.M9_base_sums.combined.windows_max.0
\newcommand{\NHFourBaseZero}{\ensuremath{3.48\cdot10^{-8}}}
% NHFourErrZero [cert, upper] FINAL_manifest_v33.json :: results.M9_H3_H2v32.rows.win.0.with_error_arb
\newcommand{\NHFourErrZero}{\ensuremath{3.5087\cdot10^{-8}}}
% NHFourMarginZero [cert, margin] FINAL_manifest_v33.json :: results.M9_H3_H2v32.rows.win.0.margin_arb
\newcommand{\NHFourMarginZero}{\ensuremath{5.700}}
% NHFourBaseOne [cert, upper] FINAL_manifest_v33.json :: results.M9_base_sums.combined.windows_max.1
\newcommand{\NHFourBaseOne}{\ensuremath{1.27\cdot10^{-7}}}
% NHFourErrOne [cert, upper] FINAL_manifest_v33.json :: results.M9_H3_H2v32.rows.win.1.with_error_arb
\newcommand{\NHFourErrOne}{\ensuremath{1.2616\cdot10^{-7}}}
% NHFourMarginOne [cert, margin] FINAL_manifest_v33.json :: results.M9_H3_H2v32.rows.win.1.margin_arb
\newcommand{\NHFourMarginOne}{\ensuremath{79.26}}
% NHFourBaseTwo [cert, upper] FINAL_manifest_v33.json :: results.M9_base_sums.combined.windows_max.2
\newcommand{\NHFourBaseTwo}{\ensuremath{3.76\cdot10^{-7}}}
% NHFourErrTwo [cert, upper] FINAL_manifest_v33.json :: results.M9_H3_H2v32.rows.win.2.with_error_arb
\newcommand{\NHFourErrTwo}{\ensuremath{3.7724\cdot10^{-7}}}
% NHFourMarginTwo [cert, margin] FINAL_manifest_v33.json :: results.M9_H3_H2v32.rows.win.2.margin_arb
\newcommand{\NHFourMarginTwo}{\ensuremath{26.50}}
% NHFourBaseThree [cert, upper] FINAL_manifest_v33.json :: results.M9_base_sums.combined.windows_max.3
\newcommand{\NHFourBaseThree}{\ensuremath{9.71\cdot10^{-7}}}
% NHFourErrThree [cert, upper] FINAL_manifest_v33.json :: results.M9_H3_H2v32.rows.win.3.with_error_arb
\newcommand{\NHFourErrThree}{\ensuremath{1.2822\cdot10^{-6}}}
% NHFourMarginThree [cert, margin] FINAL_manifest_v33.json :: results.M9_H3_H2v32.rows.win.3.margin_arb
\newcommand{\NHFourMarginThree}{\ensuremath{7.799}}
% NHFourYbarLimit [cert, lower] FINAL_manifest_v33.json :: results.M9_H3_H2v32.binding_Ybar_limit
\newcommand{\NHFourYbarLimit}{\ensuremath{1.575\cdot10^{-11}}}
% NHFourYbarMargin [cert, margin] FINAL_manifest_v33.json :: results.M9_H3_H2v32.binding_Ybar_limit
\newcommand{\NHFourYbarMargin}{\ensuremath{15.75}}
% NHFourKfourBase [cert, upper] amp_A1_w180.json :: inputs.BASE_T1.K4
\newcommand{\NHFourKfourBase}{\ensuremath{3.1633\cdot10^{-4}}}
% NHSixConeSlope [cert, upper] FINAL_manifest_v33.json :: results.cone_fixedL.checks.slope.value
\newcommand{\NHSixConeSlope}{\ensuremath{0.03029}}
% NHSixConeSlopeMargin [cert, margin] FINAL_manifest_v33.json :: results.cone_fixedL.checks.slope.margin
\newcommand{\NHSixConeSlopeMargin}{\ensuremath{2.641}}
% NHSixConeKtwo [cert, upper] FINAL_manifest_v33.json :: results.cone_fixedL.checks.K2.value
\newcommand{\NHSixConeKtwo}{\ensuremath{0.009876}}
% NHSixConeKtwoMargin [cert, margin] FINAL_manifest_v33.json :: results.cone_fixedL.checks.K2.margin
\newcommand{\NHSixConeKtwoMargin}{\ensuremath{3.037}}
% NHSixConeNthree [cert, upper] FINAL_manifest_v33.json :: results.cone_fixedL.checks.N3.value
\newcommand{\NHSixConeNthree}{\ensuremath{1.461\cdot10^{-4}}}
% NHSixConeNthreeMargin [cert, margin] FINAL_manifest_v33.json :: results.cone_fixedL.checks.N3.margin
\newcommand{\NHSixConeNthreeMargin}{\ensuremath{27.38}}
% NHSixConeHFive [cert, upper] FINAL_manifest_v33.json :: results.H4_fixedL.total
\newcommand{\NHSixConeHFive}{\ensuremath{1.2865\cdot10^{-5}}}
% NHSixConeHFiveMargin [cert, margin] FINAL_manifest_v33.json :: results.H4_fixedL.margin
\newcommand{\NHSixConeHFiveMargin}{\ensuremath{1119}}
% NHSixConeTstarNeeded [cert, upper] FINAL_manifest_v33.json :: results.H7_fixedL.ratio_upper
\newcommand{\NHSixConeTstarNeeded}{\ensuremath{1.6000015\cdot10^{-7}}}
% NHSixConeTstarLower [cert, lower] FINAL_manifest_v33.json :: results.H7_fixedL.t_star_lower
\newcommand{\NHSixConeTstarLower}{\ensuremath{3.68334\cdot10^{-7}}}
% NHSixConeTstarMargin [cert, margin] FINAL_manifest_v33.json :: results.H7_fixedL.margin_tstar_over_ratio
\newcommand{\NHSixConeTstarMargin}{\ensuremath{2.302}}
% NHSixConeHFiveRatioPeak [cert, upper] h4loc_A1_w180.json :: windows.constants/amp_A1_w180.json.worst.ratio
\newcommand{\NHSixConeHFiveRatioPeak}{\ensuremath{0.2393}}
% NHSixConeHOneSup [cert, upper] FINAL_manifest_v33.json :: results.H1_fixedL.sup_h0_max
\newcommand{\NHSixConeHOneSup}{\ensuremath{0.2163}}
% NHSixConeHOneH [cert, upper] FINAL_manifest_v33.json :: results.H1_fixedL.H_max
\newcommand{\NHSixConeHOneH}{\ensuremath{0.2185}}
% NHSixFarHFiveRatio [cert, upper] h4loc_v32.json :: windows.constants/amp_A1_w180.json.cylinders.18.ratio
\newcommand{\NHSixFarHFiveRatio}{\ensuremath{0.1761}}
% NHSixFarHFiveMarginOld [cert, margin] h4loc_v32.json :: windows.constants/amp_A1_w180.json.cylinders.18.ratio
\newcommand{\NHSixFarHFiveMarginOld}{\ensuremath{5.68}}
% NHSixFarHFiveMargin [cert, margin] FINAL_manifest_v33.json :: results.M9_H4_far.min_margin
\newcommand{\NHSixFarHFiveMargin}{\ensuremath{15.58}}
% NHSixConeHFiveMarginOld [cert, margin] h4loc_A1_w180.json :: windows.constants/amp_A1_w180.json.worst.ratio
\newcommand{\NHSixConeHFiveMarginOld}{\ensuremath{4.18}}
% NHSixFarKtwo [cert, upper] t9_far_v32.json :: h2
\newcommand{\NHSixFarKtwo}{\ensuremath{1.797\cdot10^{-4}}}
% NHSixFarKtwoThr [hand, exact] t9_far_v32.json :: thr.h2
\newcommand{\NHSixFarKtwoThr}{\ensuremath{3.53\cdot10^{-4}}}
% NHSixFarKtwoMargin [cert, margin] t9_far_v32.json :: margins.h2
\newcommand{\NHSixFarKtwoMargin}{\ensuremath{1.964}}
% NHSixFarKthree [cert, upper] t9_far_v32.json :: h3
\newcommand{\NHSixFarKthree}{\ensuremath{2.595\cdot10^{-8}}}
% NHSixFarKthreeThr [hand, exact] t9_far_v32.json :: thr.h3
\newcommand{\NHSixFarKthreeThr}{\ensuremath{1.17\cdot10^{-5}}}
% NHSixFarKthreeMargin [cert, margin] t9_far_v32.json :: margins.h3
\newcommand{\NHSixFarKthreeMargin}{\ensuremath{450.9}}
% NHSixFarKfour [cert, upper] t9_far_v32.json :: h4
\newcommand{\NHSixFarKfour}{\ensuremath{6.964\cdot10^{-9}}}
% NHSixFarKfourThr [hand, exact] t9_far_v32.json :: thr.h4
\newcommand{\NHSixFarKfourThr}{\ensuremath{1.03\cdot10^{-7}}}
% NHSixFarKfourMargin [cert, margin] t9_far_v32.json :: margins.h4
\newcommand{\NHSixFarKfourMargin}{\ensuremath{14.79}}
% NHSixFarSlope [cert, upper] t9_far_v32.json :: slope
\newcommand{\NHSixFarSlope}{\ensuremath{0.04153}}
% NHSixFarSlopeThr [hand, exact] t9_far_v32.json :: thr.slope
\newcommand{\NHSixFarSlopeThr}{\ensuremath{0.0800}}
% NHSixFarSlopeMargin [cert, margin] t9_far_v32.json :: margins.slope
\newcommand{\NHSixFarSlopeMargin}{\ensuremath{1.926}}
% NHSixFarNthree [cert, upper] t9_far_v32.json :: L2
\newcommand{\NHSixFarNthree}{\ensuremath{5.516\cdot10^{-7}}}
% NHSixFarNthreeThr [hand, exact] t9_far_v32.json :: thr.L2
\newcommand{\NHSixFarNthreeThr}{\ensuremath{3.32\cdot10^{-6}}}
% NHSixFarNthreeMargin [cert, margin] t9_far_v32.json :: margins.L2
\newcommand{\NHSixFarNthreeMargin}{\ensuremath{6.019}}
% NFarTstarMargin [cert, margin] lemmaO_v32_consts.json :: far_v32.case.margin_factor_lower
\newcommand{\NFarTstarMargin}{\ensuremath{15.5}}
% NFarTstarLower [cert, lower] lemmaO_v32_consts.json :: far_v32.case.t_star_lower
\newcommand{\NFarTstarLower}{\ensuremath{9.5108\cdot10^{-8}}}
% NFarTstarNeeded [cert, upper] lemmaO_v32_consts.json :: far_v32.case.needed
\newcommand{\NFarTstarNeeded}{\ensuremath{6.134\cdot10^{-9}}}
% NGammaModes [cert, exact] gamma0_final.json :: M
\newcommand{\NGammaModes}{\ensuremath{3762}}
% NGammaHalfwidth [cert, upper] gamma0_final.json :: b_j_halfwidth
\newcommand{\NGammaHalfwidth}{\ensuremath{8.138\cdot10^{-24}}}
% NGammaContraction [cert, upper] gamma0_final.json :: krawczyk_contraction
\newcommand{\NGammaContraction}{\ensuremath{3.01\cdot10^{-12}}}
% NGammaArcs [cert, exact] gamma0_final.json :: embedded.arcs
\newcommand{\NGammaArcs}{\ensuremath{701}}
% NGammaMaxTurn [cert, upper] gamma0_final.json :: embedded.maxturn
\newcommand{\NGammaMaxTurn}{\ensuremath{0.04997}}
% NGammaMaxPair [cert, upper] gamma0_final.json :: embedded.maxpair
\newcommand{\NGammaMaxPair}{\ensuremath{0.09597}}
% NGammaMinSep [cert, lower] gamma0_final.json :: embedded.minsep
\newcommand{\NGammaMinSep}{\ensuremath{8.89\cdot10^{-3}}}
% NGammaHalfGap [cert, lower] gamma0_final.json :: embedded.gap_cert
\newcommand{\NGammaHalfGap}{\ensuremath{0.0248219}}
% NPhaseISlabs [cert, exact] status.json :: slab
\newcommand{\NPhaseISlabs}{\ensuremath{96}}
% NPhaseIModes [cert, exact] status.json :: last_slab.M
\newcommand{\NPhaseIModes}{\ensuremath{2800}}
% NPhaseIKmaxLower [cert, lower] status.json :: kmax_enclosure.0
\newcommand{\NPhaseIKmaxLower}{\ensuremath{29.352874545}}
% NPhaseIKmaxUpper [cert, upper] status.json :: kmax_enclosure.1
\newcommand{\NPhaseIKmaxUpper}{\ensuremath{29.352874546}}
% NPhaseITLower [cert, lower] status.json :: t
\newcommand{\NPhaseITLower}{\ensuremath{4.597850539485830}}
% NPhaseITUpper [cert, upper] status.json :: t
\newcommand{\NPhaseITUpper}{\ensuremath{4.597850539485831}}
% NPeakSqrtLam [cert, upper] rate_table_fin12f4_a0.00015.json :: sqrtL_max
\newcommand{\NPeakSqrtLam}{\ensuremath{5.308\cdot10^{-8}}}
% NTailRate [cert, lower] rate_table_fin12f4_a0.00015.json :: tail.rho
\newcommand{\NTailRate}{\ensuremath{2.693\cdot10^{-3}}}
% NFloorSqrtLam [cert, upper] rate_table_fin12f4_a0.00015.json :: tail.sqrt_LF
\newcommand{\NFloorSqrtLam}{\ensuremath{9.872\cdot10^{-10}}}
% NAlphaFloor [cert, upper] rate_table_fin12f4_a0.00015.json :: tail.alpha_F
\newcommand{\NAlphaFloor}{\ensuremath{6.867\cdot10^{-7}}}
% NSupBetaVerifier [cert, upper] final_alpha_points_5e-8.json :: 1.5e-4.sup_beta
\newcommand{\NSupBetaVerifier}{\ensuremath{3.33811\cdot10^{-9}}}
% NMOneRamp [cert, upper] final_alpha_points_5e-8.json :: 1.5e-4.M1
\newcommand{\NMOneRamp}{\ensuremath{0.4328}}
% NSITwoPrime [cert, upper] FINAL_manifest_v33.json :: results.F12_sliver.sI2_prime
\newcommand{\NSITwoPrime}{\ensuremath{0.015457}}
% NLoopGainCoarse [cert, upper] U_cert_v35.json :: loop_gain.2.total
\newcommand{\NLoopGainCoarse}{\ensuremath{0.03904}}
% NRateMarginCoarse [cert, margin] U_cert_v35.json :: rate_margin
\newcommand{\NRateMarginCoarse}{\ensuremath{28.2}}
% NSOneEpsLower [cert, lower] bracket2.json :: eps_a
\newcommand{\NSOneEpsLower}{\ensuremath{1.5346198892965\cdot10^{-4}}}
% NSOneEpsUpper [cert, upper] bracket2.json :: eps_b
\newcommand{\NSOneEpsUpper}{\ensuremath{1.5346198899299\cdot10^{-4}}}
% NSOneRowsTip [cert, exact] ext_tip_check.json :: rows
\newcommand{\NSOneRowsTip}{\ensuremath{51246}}
% NSOneTipWidth [cert, upper] ext_tip_check.json :: max_pointwise_width.th
\newcommand{\NSOneTipWidth}{\ensuremath{3.2\cdot10^{-12}}}
% NSOneExtWidth [cert, upper] ext813_check.json :: max_pointwise_width.p
\newcommand{\NSOneExtWidth}{\ensuremath{5.7\cdot10^{-10}}}
% NSOneRowsMC [cert, exact] check_outA5.json :: prof_tip_vs_MC.rows_compared
\newcommand{\NSOneRowsMC}{\ensuremath{1025}}
% NSOneRowsMB [cert, exact] check_outA5.json :: ext813_vs_MB.rows_compared
\newcommand{\NSOneRowsMB}{\ensuremath{501}}
% NSOneMinMargin [cert, margin] check_outA5.json :: min:*.min_relative_margin.*
\newcommand{\NSOneMinMargin}{\ensuremath{1.2\cdot10^{-9}}}
% NTwoSegments [cert, exact] contour_result.json :: len:segments
\newcommand{\NTwoSegments}{\ensuremath{5058}}
% NTwoInteriorBoxes [cert, exact] interior_result.json :: len:
\newcommand{\NTwoInteriorBoxes}{\ensuremath{174}}
% NTwoContourBoxes [cert, exact] eta_cache.json :: len:
\newcommand{\NTwoContourBoxes}{\ensuremath{202}}
% NTwoEtaMax [cert, upper] eta_cache.json :: max:*.eta
\newcommand{\NTwoEtaMax}{\ensuremath{0.1163}}
% NTwoLipMax [cert, upper] eta_cache.json :: max:*.Lip
\newcommand{\NTwoLipMax}{\ensuremath{0.272}}
% NTwoLambdaZero [est, upper] cert_b0.25_th0.0_tau0.03.json :: Lam
\newcommand{\NTwoLambdaZero}{\ensuremath{1.790819}}
% NTwoLambdaOne [est, upper] cert_b0.25_th0.48_tau0.03.json :: Lam
\newcommand{\NTwoLambdaOne}{\ensuremath{1.768282}}
% NSOneDTightLower [doc, lower] doc:README.md :: D ∈ [2.50080627654, 2.50080627668]
\newcommand{\NSOneDTightLower}{\ensuremath{2.50080627654}}
% NSOneDTightUpper [doc, upper] doc:README.md :: D ∈ [2.50080627654, 2.50080627668]
\newcommand{\NSOneDTightUpper}{\ensuremath{2.50080627668}}
% NSOneKTightLower [doc, lower] doc:README.md :: k(0) ∈ [1.53982063762, 1.53982063822]
\newcommand{\NSOneKTightLower}{\ensuremath{1.53982063762}}
% NSOneKTightUpper [doc, upper] doc:README.md :: k(0) ∈ [1.53982063762, 1.53982063822]
\newcommand{\NSOneKTightUpper}{\ensuremath{1.53982063822}}
% NSOneDCondLower [doc, lower] doc:README.md :: D = −p(tip) ∈ [2.500806276619695, 2.500806276620049]
\newcommand{\NSOneDCondLower}{\ensuremath{2.500806276619695}}
% NSOneDCondUpper [doc, upper] doc:README.md :: D = −p(tip) ∈ [2.500806276619695, 2.500806276620049]
\newcommand{\NSOneDCondUpper}{\ensuremath{2.500806276620049}}
% NSOneKCondLower [doc, lower] doc:README.md :: k(0) ∈ [1.539820637871836, 1.5398206378721015]
\newcommand{\NSOneKCondLower}{\ensuremath{1.539820637871836}}
% NSOneKCondUpper [doc, upper] doc:README.md :: k(0) ∈ [1.539820637871836, 1.5398206378721015]
\newcommand{\NSOneKCondUpper}{\ensuremath{1.5398206378721015}}
% NSOneTauTipLower [doc, lower] doc:README.md :: τ_tip ∈ [10.249077651120778, 10.249077651121361]
\newcommand{\NSOneTauTipLower}{\ensuremath{10.249077651120778}}
% NSOneTauTipUpper [doc, upper] doc:README.md :: τ_tip ∈ [10.249077651120778, 10.249077651121361]
\newcommand{\NSOneTauTipUpper}{\ensuremath{10.249077651121361}}
% NSOneStateUncert [doc, upper] doc:README.md :: State uncertainty at σ = 8: 4.4e−13
\newcommand{\NSOneStateUncert}{\ensuremath{4.4\cdot10^{-13}}}
% NFarBoxEprime [doc, margin] doc:STEP4_RESULT.md :: (6.4287, 4.6917, 1.9940, 1.69999995)
\newcommand{\NFarBoxEprime}{\ensuremath{6.42}}
% NFarBoxXEsecond [doc, margin] doc:STEP4_RESULT.md :: (6.4287, 4.6917, 1.9940, 1.69999995)
\newcommand{\NFarBoxXEsecond}{\ensuremath{4.69}}
% NFarBoxXsqEthird [doc, margin] doc:STEP4_RESULT.md :: (6.4287, 4.6917, 1.9940, 1.69999995)
\newcommand{\NFarBoxXsqEthird}{\ensuremath{1.99}}
% NFarBoxE [doc, margin] doc:STEP4_RESULT.md :: (6.4287, 4.6917, 1.9940, 1.69999995)
\newcommand{\NFarBoxE}{\ensuremath{1.69}}
% NHFourKfourErr [doc, upper] doc:STEP4_RESULT.md :: | 5.60e-3 against 0.06 | ×10.7 |
\newcommand{\NHFourKfourErr}{\ensuremath{5.60\cdot10^{-3}}}
% NHFourKfourMargin [doc, margin] doc:STEP4_RESULT.md :: | 5.60e-3 against 0.06 | ×10.7 |
\newcommand{\NHFourKfourMargin}{\ensuremath{10.7}}
% NCoercivityCoefficient [doc, upper] doc:step3_proof_v32.tex :: \norm{e^{\tau\mathcal L}u}_\rho\le14.8\,e^{-0.0213\tau}\norm
\newcommand{\NCoercivityCoefficient}{\ensuremath{14.8}}
% NCoercivityRate [doc, lower] doc:step3_proof_v32.tex :: \norm{e^{\tau\mathcal L}u}_\rho\le14.8\,e^{-0.0213\tau}\norm
\newcommand{\NCoercivityRate}{\ensuremath{0.0213}}
% NRegionRampStart [hand, exact] hand :: 
\newcommand{\NRegionRampStart}{\ensuremath{14}}
% NRegionInterface [hand, exact] hand :: 
\newcommand{\NRegionInterface}{\ensuremath{28}}
% NRegionShellEnd [hand, exact] hand :: 
\newcommand{\NRegionShellEnd}{\ensuremath{60}}
% NRegionSegEnd [hand, exact] hand :: 
\newcommand{\NRegionSegEnd}{\ensuremath{350}}
% NRegionAnnEnd [hand, exact] hand :: 
\newcommand{\NRegionAnnEnd}{\ensuremath{700}}
% NZoneIIILow [hand, exact] hand :: 
\newcommand{\NZoneIIILow}{\ensuremath{253}}
% NZoneIIIHigh [hand, exact] hand :: 
\newcommand{\NZoneIIIHigh}{\ensuremath{379.5}}
% NConeCentreMin [hand, exact] hand :: 
\newcommand{\NConeCentreMin}{\ensuremath{300}}
% NConeCentreMax [hand, exact] hand :: 
\newcommand{\NConeCentreMax}{\ensuremath{380}}
% NFarCentreMin [hand, exact] hand :: 
\newcommand{\NFarCentreMin}{\ensuremath{490}}
% NConeL [hand, exact] hand :: 
\newcommand{\NConeL}{\ensuremath{50}}
% NConeX [hand, exact] hand :: 
\newcommand{\NConeX}{\ensuremath{3}}
% NFarL [hand, exact] hand :: 
\newcommand{\NFarL}{\ensuremath{113}}
% NFarX [hand, exact] hand :: 
\newcommand{\NFarX}{\ensuremath{2}}
 is found.
% ======================================================================

% ---- which document are we in ----------------------------------------
\ifdefined\ifsupplement\else\newif\ifsupplement\fi

% ---- packages ---------------------------------------------------------
\usepackage[T1]{fontenc}
\usepackage[utf8]{inputenc}
\usepackage{lmodern}
\usepackage[margin=1.1in]{geometry}
\usepackage{amsmath,amssymb,amsthm,mathtools,bm}
\usepackage{booktabs,array,longtable,tabularx,multirow}
\usepackage{enumitem}
\usepackage{graphicx}
\usepackage{xcolor}
\usepackage{tikz}
\usetikzlibrary{positioning,arrows.meta,calc}
% xr-hyper must come before hyperref
\usepackage{xr-hyper}
\usepackage[colorlinks=true,linkcolor=blue!50!black,citecolor=green!35!black,urlcolor=blue!50!black]{hyperref}

\raggedbottom
\setlength{\emergencystretch}{2em}
\setlist[enumerate]{leftmargin=*,itemsep=2pt,topsep=3pt}
\setlist[itemize]{leftmargin=*,itemsep=2pt,topsep=3pt}

% ---- numbering --------------------------------------------------------
% Main paper: sections 1,2,...; theorems 2.3 (by section); equations (2.3).
% Supplement: the Parts A..K are the \section's, numbered by letter, so
% theorems are A.1, E.4, ... (do not confuse "Theorem A" of the paper, which
% has no dot, with "Theorem A.1" of Part A).
\ifsupplement
  \renewcommand{\thesection}{\Alph{section}}
\fi
\numberwithin{equation}{section}

% ---- theorem environments (numbered by section, one shared counter) ----
\theoremstyle{plain}
\newtheorem{theorem}{Theorem}[section]
\newtheorem{proposition}[theorem]{Proposition}
\newtheorem{lemma}[theorem]{Lemma}
\newtheorem{corollary}[theorem]{Corollary}
\newtheorem{maintheorem}[theorem]{Main Theorem}
% Theorems A, B, C, D of the paper: own counter, lettered, no section prefix
\newtheorem{thmlet}{Theorem}
\renewcommand{\thethmlet}{\Alph{thmlet}}
\theoremstyle{definition}
\newtheorem{definition}[theorem]{Definition}
\newtheorem{convention}[theorem]{Convention}
\newtheorem{example}[theorem]{Example}
\theoremstyle{remark}
\newtheorem{remark}[theorem]{Remark}

% A list of hypotheses labelled (H1),(H2),... : \begin{hyplist}{H} ... \end{hyplist}
% (also {O}, {L}, {B}, {K}, {W} for the families of section 5.3 of OUTLINE.md).
\newenvironment{hyplist}[1]%
  {\begin{enumerate}[label=\textup{(#1\arabic*)},leftmargin=3em,itemsep=2pt]}%
  {\end{enumerate}}

% ---- trust classes (OUTLINE 4.1, 5.1.3) -------------------------------
% Supplement: explicit tags.  Main paper: certified is the default and
% unmarked; every non-certified number is marked "(floating point)".
\ifsupplement
  \newcommand{\tcert}{\textup{[cert]}}   % arb/outward-rounded, pinned script and key
  \newcommand{\test}{\textup{[est]}}     % floating point or estimate; never in a theorem
  \newcommand{\thand}{\textup{[hand]}}   % proved by hand in the text
  \newcommand{\tlit}{\textup{[lit]}}     % cited from the literature
\else
  \newcommand{\tcert}{}
  \newcommand{\test}{\unskip\ \textup{(floating point)}}
  \newcommand{\thand}{}
  \newcommand{\tlit}{}
\fi

% ---- gaps: visible, counted, an error in a final build -----------------
% Usage: \GAP{one sentence saying what is missing}   (and an entry in ISSUES.md)
\newcounter{gapcount}
\newif\iffinalpaper   % \finalpapertrue makes every \GAP a hard error
\newcommand{\GAP}[1]{\stepcounter{gapcount}\typeout{GAP-MARKER \thegapcount}%
  \iffinalpaper\PackageError{paper}{\string\GAP\space used in a final build}{Remove all gaps.}\fi
  {\color{red}\sffamily\bfseries[GAP~\thegapcount:\ }{\color{red}\sffamily #1}{\color{red}\sffamily\bfseries]}}
\AtEndDocument{\typeout{GAP-MARKER total: \thegapcount}}

% ---- file names and JSON keys (supplement only) -------------------------
% No internal path appears in the main paper; a use there is a warning.
\newcommand{\fn}[1]{%
  \ifsupplement\else\PackageWarning{paper}{\string\fn\space used in the main paper}\fi
  \texttt{\detokenize{#1}}}
\newcommand{\key}[1]{\texttt{\detokenize{#1}}}

% ---- cross references between main paper and supplement ----------------
% Labels carry a unit prefix (section 5.3 of OUTLINE.md and DRAFTER_GUIDE.md):
%   main paper:  intro: obj: s3: s4: s5: s6: s7: s8: s9:   (kind: thm- lem- eq- tab- fig- ...)
%   supplement:  pA: pB: ... pK:
% so all labels are globally unique and the prefixes below are only a second guard.
% \Xref{label}  : reference to the OTHER document (hyperlinked via xr-hyper)
% \Lref{label}  : reference in the SAME document (plain \ref)
\ifsupplement
  \externaldocument[M-]{main}[main.pdf]
  \newcommand{\Xref}[1]{\ref{M-#1}}
\else
  \externaldocument[S-]{supplement}[supplement.pdf]
  \newcommand{\Xref}[1]{\ref{S-#1}}
\fi
\newcommand{\Lref}[1]{\ref{#1}}
% \Supp{pE}{lem-chart} -> "Lemma E.3" style text is produced by the drafter:
%   "Lemma~\Xref{pE:lem-chart} of the supplement".  \SuppPart{E} gives "Part E".
\newcommand{\SuppPart}[1]{Part~#1 of the supplement}

% ---- numbers (generated by unit U0; macros only, no literal certified numbers) ----
\IfFileExists{numbers/numbers_paper.tex}%
  {% numbers_paper.tex -- GENERATED from ledger.json by build_ledger.py; do not edit.
% Rounding: upper bounds up, lower bounds and margins down at the printed digit; margin = threshold / bound.
% Use as $\times\NName$ for margins.  kind: cert certified, est estimate, hand hypothesis constant, doc only a document quotes it.
% NUEhatZero [cert, upper] U_cert_v35.json :: supersolution.EA.0
\newcommand{\NUEhatZero}{\ensuremath{1.2018\cdot10^{-6}}}
% NUEhatOne [cert, upper] U_cert_v35.json :: supersolution.EA.1
\newcommand{\NUEhatOne}{\ensuremath{1.0441\cdot10^{-4}}}
% NUEhatTwo [cert, upper] U_cert_v35.json :: supersolution.EA.2
\newcommand{\NUEhatTwo}{\ensuremath{0.056694}}
% NUBhatZero [cert, upper] U_cert_v35.json :: supersolution.Bhat.0
\newcommand{\NUBhatZero}{\ensuremath{4.8735\cdot10^{-6}}}
% NUBhatOne [cert, upper] U_cert_v35.json :: supersolution.Bhat.1
\newcommand{\NUBhatOne}{\ensuremath{5.5933\cdot10^{-4}}}
% NUBhatTwo [cert, upper] U_cert_v35.json :: supersolution.Bhat.2
\newcommand{\NUBhatTwo}{\ensuremath{0.36675}}
% NUSupersolutionCZeroZero [cert, upper] U_cert_v35.json :: supersolution.C0.0
\newcommand{\NUSupersolutionCZeroZero}{\ensuremath{1.1340\cdot10^{-6}}}
% NUSupersolutionCZeroOne [cert, upper] U_cert_v35.json :: supersolution.C0.1
\newcommand{\NUSupersolutionCZeroOne}{\ensuremath{1.6623\cdot10^{-4}}}
% NUSupersolutionCOneZero [cert, upper] U_cert_v35.json :: supersolution.C1.0
\newcommand{\NUSupersolutionCOneZero}{\ensuremath{1.5241\cdot10^{-9}}}
% NUSupersolutionCOneOne [cert, upper] U_cert_v35.json :: supersolution.C1.1
\newcommand{\NUSupersolutionCOneOne}{\ensuremath{5.5255\cdot10^{-8}}}
% NUSupersolutionNTwoZero [cert, upper] U_cert_v35.json :: supersolution.N2.0
\newcommand{\NUSupersolutionNTwoZero}{\ensuremath{0.071472}}
% NUSupersolutionNTwoOne [cert, upper] U_cert_v35.json :: supersolution.N2.1
\newcommand{\NUSupersolutionNTwoOne}{\ensuremath{0.33012}}
% NUSupersolutionNTwoTwo [cert, upper] U_cert_v35.json :: supersolution.N2.2
\newcommand{\NUSupersolutionNTwoTwo}{\ensuremath{5.0406}}
% NUSupersolutionNTwoOverFourNTwoRefZero [cert, upper] U_cert_v35.json :: supersolution.N2_over_4N2ref.0
\newcommand{\NUSupersolutionNTwoOverFourNTwoRefZero}{\ensuremath{0.18}}
% NUSupersolutionNTwoOverFourNTwoRefOne [cert, upper] U_cert_v35.json :: supersolution.N2_over_4N2ref.1
\newcommand{\NUSupersolutionNTwoOverFourNTwoRefOne}{\ensuremath{0.42}}
% NUSupersolutionNTwoOverFourNTwoRefTwo [cert, upper] U_cert_v35.json :: supersolution.N2_over_4N2ref.2
\newcommand{\NUSupersolutionNTwoOverFourNTwoRefTwo}{\ensuremath{0.26}}
% NUSupersolutionMOneZHypCheckZonesIIISOneTwoBhat [cert, upper] U_cert_v35.json :: supersolution.M1z_hyp_check.zones_I_II.S1_2Bhat
\newcommand{\NUSupersolutionMOneZHypCheckZonesIIISOneTwoBhat}{\ensuremath{0.28585}}
% NUSupersolutionMOneZHypCheckZoneIIIBOnePlusPUOne [cert, upper] U_cert_v35.json :: supersolution.M1z_hyp_check.zone_III.b1_plus_pU1
\newcommand{\NUSupersolutionMOneZHypCheckZoneIIIBOnePlusPUOne}{\ensuremath{0.50786}}
% NUMarginLateMOne [cert, margin] U_cert_v35.json :: margins.late_M1
\newcommand{\NUMarginLateMOne}{\ensuremath{2.353}}
% NUMarginLateMZero [cert, margin] U_cert_v35.json :: margins.late_M0
\newcommand{\NUMarginLateMZero}{\ensuremath{4.892}}
% NUMarginsLateMZeroZ [cert, margin] U_cert_v35.json :: margins.late_M0z
\newcommand{\NUMarginsLateMZeroZ}{\ensuremath{12.99}}
% NUMarginPresentMOne [cert, margin] U_cert_v35.json :: margins.present_M1_min.value
\newcommand{\NUMarginPresentMOne}{\ensuremath{2.770}}
% NUMarginPresentMZero [cert, margin] U_cert_v35.json :: margins.present_M0_min.value
\newcommand{\NUMarginPresentMZero}{\ensuremath{2.225}}
% NUMarginsPresentMZeroZMinValue [cert, margin] U_cert_v35.json :: margins.present_M0z_min.value
\newcommand{\NUMarginsPresentMZeroZMinValue}{\ensuremath{14.37}}
% NULoopGainZeroTotal [cert, upper] U_cert_v35.json :: loop_gain.0.total
\newcommand{\NULoopGainZeroTotal}{\ensuremath{1.23\cdot10^{-7}}}
% NULoopGainOneTotal [cert, upper] U_cert_v35.json :: loop_gain.1.total
\newcommand{\NULoopGainOneTotal}{\ensuremath{0.037169}}
% NULoopGainTwoNeighbour [cert, upper] U_cert_v35.json :: loop_gain.2.neighbour
\newcommand{\NULoopGainTwoNeighbour}{\ensuremath{0.037018}}
% NULoopGainTwoThroughA [cert, upper] U_cert_v35.json :: loop_gain.2.through_A
\newcommand{\NULoopGainTwoThroughA}{\ensuremath{0.0020172}}
% NUfineBhatZero [cert, upper] U_cert_v35.json :: ufine.Bhat.0
\newcommand{\NUfineBhatZero}{\ensuremath{4.8723\cdot10^{-6}}}
% NUfineBhatOne [cert, upper] U_cert_v35.json :: ufine.Bhat.1
\newcommand{\NUfineBhatOne}{\ensuremath{2.6266\cdot10^{-4}}}
% NUfineBhatTwo [cert, upper] U_cert_v35.json :: ufine.Bhat.2
\newcommand{\NUfineBhatTwo}{\ensuremath{0.025842}}
% NUcdeltaU [cert, upper] U_cert_v35.json :: launch.c_dU
\newcommand{\NUcdeltaU}{\ensuremath{1.7187\cdot10^{-8}}}
% NULaunchWmin [cert, lower] U_cert_v35.json :: launch.wmin
\newcommand{\NULaunchWmin}{\ensuremath{0.73492}}
% NULaunchCAD [cert, upper] U_cert_v35.json :: launch.CA_D
\newcommand{\NULaunchCAD}{\ensuremath{2.2429}}
% NULaunchQD [cert, upper] U_cert_v35.json :: launch.qD
\newcommand{\NULaunchQD}{\ensuremath{0.018509}}
% NULaunchDeltaInZero [cert, upper] U_cert_v35.json :: launch.Delta_in.0
\newcommand{\NULaunchDeltaInZero}{\ensuremath{3.131\cdot10^{-12}}}
% NULaunchDeltaInOne [cert, upper] U_cert_v35.json :: launch.Delta_in.1
\newcommand{\NULaunchDeltaInOne}{\ensuremath{3.287\cdot10^{-12}}}
% NULaunchDeltaInTwo [cert, upper] U_cert_v35.json :: launch.Delta_in.2
\newcommand{\NULaunchDeltaInTwo}{\ensuremath{3.451\cdot10^{-12}}}
% NULaunchPhiZero [cert, upper] U_cert_v35.json :: launch.Phi.0
\newcommand{\NULaunchPhiZero}{\ensuremath{6.434\cdot10^{-9}}}
% NULaunchPhiOne [cert, upper] U_cert_v35.json :: launch.Phi.1
\newcommand{\NULaunchPhiOne}{\ensuremath{1.4217\cdot10^{-8}}}
% NULaunchPhiTwo [cert, upper] U_cert_v35.json :: launch.Phi.2
\newcommand{\NULaunchPhiTwo}{\ensuremath{3.8404\cdot10^{-8}}}
% NUGammaZero [cert, upper] U_cert_v35.json :: Gamma.0
\newcommand{\NUGammaZero}{\ensuremath{4.0000358}}
% NUGammaOne [cert, upper] U_cert_v35.json :: Gamma.1
\newcommand{\NUGammaOne}{\ensuremath{16.000322}}
% NUGammaTwo [cert, upper] U_cert_v35.json :: Gamma.2
\newcommand{\NUGammaTwo}{\ensuremath{64.00218}}
% NUKappa [cert, upper] U_cert_v35.json :: kappa
\newcommand{\NUKappa}{\ensuremath{8.0000358}}
% NUKappaTwo [cert, upper] U_cert_v35.json :: kappa2
\newcommand{\NUKappaTwo}{\ensuremath{64.00061}}
% NUGhatZeroZero [cert, upper] U_cert_v35.json :: Ghat.0.0
\newcommand{\NUGhatZeroZero}{\ensuremath{51322.6}}
% NUGhatOneOne [cert, upper] U_cert_v35.json :: Ghat.1.1
\newcommand{\NUGhatOneOne}{\ensuremath{33962.9}}
% NUGhatTwoTwo [cert, upper] U_cert_v35.json :: Ghat.2.2
\newcommand{\NUGhatTwoTwo}{\ensuremath{348341.1}}
% NUGhatOneZero [cert, upper] U_cert_v35.json :: Ghat.1.0
\newcommand{\NUGhatOneZero}{\ensuremath{4.938}}
% NUGhatTwoZero [cert, upper] U_cert_v35.json :: Ghat.2.0
\newcommand{\NUGhatTwoZero}{\ensuremath{27.984}}
% NUGhatTwoOne [cert, upper] U_cert_v35.json :: Ghat.2.1
\newcommand{\NUGhatTwoOne}{\ensuremath{458.53}}
% NUKZeroZero [cert, upper] U_cert_v35.json :: K.0.0
\newcommand{\NUKZeroZero}{\ensuremath{8034.21}}
% NUKZeroOne [cert, upper] U_cert_v35.json :: K.0.1
\newcommand{\NUKZeroOne}{\ensuremath{724.89}}
% NUKOneZero [cert, upper] U_cert_v35.json :: K.1.0
\newcommand{\NUKOneZero}{\ensuremath{0.08152}}
% NUKOneOne [cert, upper] U_cert_v35.json :: K.1.1
\newcommand{\NUKOneOne}{\ensuremath{6336.94}}
% NUKOneTwo [cert, upper] U_cert_v35.json :: K.1.2
\newcommand{\NUKOneTwo}{\ensuremath{545.59}}
% NUKOneThree [cert, upper] U_cert_v35.json :: K.1.3
\newcommand{\NUKOneThree}{\ensuremath{0.58700}}
% NUKTwoZero [cert, upper] U_cert_v35.json :: K.2.0
\newcommand{\NUKTwoZero}{\ensuremath{0.87186}}
% NUKTwoOne [cert, upper] U_cert_v35.json :: K.2.1
\newcommand{\NUKTwoOne}{\ensuremath{528.62}}
% NUKTwoTwo [cert, upper] U_cert_v35.json :: K.2.2
\newcommand{\NUKTwoTwo}{\ensuremath{23425.98}}
% NUKTwoThree [cert, upper] U_cert_v35.json :: K.2.3
\newcommand{\NUKTwoThree}{\ensuremath{555.62}}
% NUFootwOne [cert, upper] U_cert_v35.json :: footw.1
\newcommand{\NUFootwOne}{\ensuremath{2.6361}}
% NUFootwTwo [cert, upper] U_cert_v35.json :: footw.2
\newcommand{\NUFootwTwo}{\ensuremath{2663.3}}
% NUCThreeThree [cert, upper] U_cert_v35.json :: c33
\newcommand{\NUCThreeThree}{\ensuremath{108.02}}
% NUCTwoTwoThree [cert, upper] U_cert_v35.json :: c223
\newcommand{\NUCTwoTwoThree}{\ensuremath{7.3284}}
% NUCEta [cert, upper] U_cert_v35.json :: C_eta
\newcommand{\NUCEta}{\ensuremath{13.376}}
% NULMax [cert, upper] U_cert_v35.json :: L_max
\newcommand{\NULMax}{\ensuremath{3.5352}}
% NURRatio [cert, upper] U_cert_v35.json :: R_ratio
\newcommand{\NURRatio}{\ensuremath{11.593}}
% NUSweepsSigmaInCoverSup [cert, upper] U_cert_v35.json :: sweeps.Sigma_in.cover_sup
\newcommand{\NUSweepsSigmaInCoverSup}{\ensuremath{5.0892}}
% NUSweepsSigmaECoverSup [cert, upper] U_cert_v35.json :: sweeps.Sigma_E.cover_sup
\newcommand{\NUSweepsSigmaECoverSup}{\ensuremath{0.30846}}
% NUSweepsSXCoverSup [cert, upper] U_cert_v35.json :: sweeps.S_X.cover_sup
\newcommand{\NUSweepsSXCoverSup}{\ensuremath{1.5243}}
% NUSweepsSigmaXCoverSup [cert, upper] U_cert_v35.json :: sweeps.Sigma_X.cover_sup
\newcommand{\NUSweepsSigmaXCoverSup}{\ensuremath{3.0081}}
% NUSweepsSigmaRsTwoCoverSup [cert, upper] U_cert_v35.json :: sweeps.Sigma_rs.2.cover_sup
\newcommand{\NUSweepsSigmaRsTwoCoverSup}{\ensuremath{4.3350\cdot10^{-3}}}
% NUDataTermsLateDZero [cert, upper] U_cert_v35.json :: data_terms_late.D.0
\newcommand{\NUDataTermsLateDZero}{\ensuremath{1.7330\cdot10^{-6}}}
% NUDataTermsLateDOne [cert, upper] U_cert_v35.json :: data_terms_late.D.1
\newcommand{\NUDataTermsLateDOne}{\ensuremath{2.5353\cdot10^{-4}}}
% NUDataTermsLateSlopeShareCOneSigmaXOverDkZero [cert, upper] U_cert_v35.json :: data_terms_late.slope_share_C1SigmaX_over_Dk.0
\newcommand{\NUDataTermsLateSlopeShareCOneSigmaXOverDkZero}{\ensuremath{2.7\cdot10^{-3}}}
% NUDataTermsLateSlopeShareCOneSigmaXOverDkOne [cert, upper] U_cert_v35.json :: data_terms_late.slope_share_C1SigmaX_over_Dk.1
\newcommand{\NUDataTermsLateSlopeShareCOneSigmaXOverDkOne}{\ensuremath{6.6\cdot10^{-4}}}
% NURLtTwoMarginZero [cert, margin] U_cert_v35.json :: R_lt2_margin.0
\newcommand{\NURLtTwoMarginZero}{\ensuremath{4.432}}
% NURLtTwoMarginOne [cert, margin] U_cert_v35.json :: R_lt2_margin.1
\newcommand{\NURLtTwoMarginOne}{\ensuremath{6.697}}
% NUThetaTwoEight [cert, upper] U_cert_v35.json :: theta28
\newcommand{\NUThetaTwoEight}{\ensuremath{0.37011}}
% NUDFiveSix [cert, upper] U_cert_v35.json :: D56
\newcommand{\NUDFiveSix}{\ensuremath{1.0368}}
% NUDTwoEight [cert, upper] U_cert_v35.json :: D28
\newcommand{\NUDTwoEight}{\ensuremath{1.5876}}
% NUSThreeThreeFootOverBhatTwo [cert, upper] U_cert_v35.json :: S33_foot_over_Bhat2
\newcommand{\NUSThreeThreeFootOverBhatTwo}{\ensuremath{1.69\cdot10^{-3}}}
% NSAHJWKappaJ [cert, upper] A_interface.json :: HJW.kappa_J
\newcommand{\NSAHJWKappaJ}{\ensuremath{0.042351}}
% NSAGpZero [cert, lower] A_interface.json :: res.tight.gp.0
\newcommand{\NSAGpZero}{\ensuremath{0.24292}}
% NSAGpOne [cert, lower] A_interface.json :: res.tight.gp.1
\newcommand{\NSAGpOne}{\ensuremath{0.24106}}
% NSAGpTwo [cert, lower] A_interface.json :: res.tight.gp.2
\newcommand{\NSAGpTwo}{\ensuremath{0.11961}}
% NSAAminC [cert, lower] A_interface.json :: res.tight.amin_c
\newcommand{\NSAAminC}{\ensuremath{0.97623}}
% NSARhoJ [cert, upper] A_interface.json :: res.tight.rhoJ
\newcommand{\NSARhoJ}{\ensuremath{5.4995\cdot10^{-7}}}
% NSAKappaOTwoTwo [cert, upper] A_interface.json :: res.tight.kappa_o.2.2
\newcommand{\NSAKappaOTwoTwo}{\ensuremath{1.2416\cdot10^{-5}}}
% NSADsrcZero [cert, upper] A_interface.json :: res.tight.Dsrc.0
\newcommand{\NSADsrcZero}{\ensuremath{1.8092\cdot10^{-7}}}
% NSADsrcOne [cert, upper] A_interface.json :: res.tight.Dsrc.1
\newcommand{\NSADsrcOne}{\ensuremath{9.8276\cdot10^{-5}}}
% NSABZeroDWindowSlopeZero [cert, upper] A_interface.json :: res.tight.B0.D_window_slope.0
\newcommand{\NSABZeroDWindowSlopeZero}{\ensuremath{8.7399\cdot10^{-10}}}
% NSABZeroDWindowSlopeOne [cert, upper] A_interface.json :: res.tight.B0.D_window_slope.1
\newcommand{\NSABZeroDWindowSlopeOne}{\ensuremath{5.2599\cdot10^{-8}}}
% NSAFTwoZero [cert, upper] A_interface.json :: res.tight.F2.0
\newcommand{\NSAFTwoZero}{\ensuremath{5.6506\cdot10^{-13}}}
% NSAFTwoOne [cert, upper] A_interface.json :: res.tight.F2.1
\newcommand{\NSAFTwoOne}{\ensuremath{5.9320\cdot10^{-13}}}
% NSACoupOne [cert, upper] A_interface.json :: res.tight.coup.1
\newcommand{\NSACoupOne}{\ensuremath{1.1568\cdot10^{-12}}}
% NSACOneZeroZero [cert, upper] A_interface.json :: res.tight.C1_0.0
\newcommand{\NSACOneZeroZero}{\ensuremath{7.2060\cdot10^{-10}}}
% NSACOneZeroOne [cert, upper] A_interface.json :: res.tight.C1_0.1
\newcommand{\NSACOneZeroOne}{\ensuremath{4.3143\cdot10^{-8}}}
% NSABZeroSupCoreZero [cert, upper] A_interface.json :: res.tight.B0.sup_core.0
\newcommand{\NSABZeroSupCoreZero}{\ensuremath{4.23\cdot10^{-5}}}
% NSACoreFluxMarginTwoFLUXZeroOverKZero [cert, margin] A_interface.json :: res.tight.core.flux_margin.2.FLUX0_over_K0
\newcommand{\NSACoreFluxMarginTwoFLUXZeroOverKZero}{\ensuremath{3.396}}
% NSACoreFluxMarginTwoFLUXTwoOverKTwo [cert, margin] A_interface.json :: res.tight.core.flux_margin.2.FLUX2_over_K2
\newcommand{\NSACoreFluxMarginTwoFLUXTwoOverKTwo}{\ensuremath{6.012}}
% NSAMarginFluxZero [cert, margin] asm_design_N2ref4.json :: res.tight.flux.beta.2.margin_FLUX0
\newcommand{\NSAMarginFluxZero}{\ensuremath{3.365}}
% NSAMarginFluxTwo [cert, margin] asm_design_N2ref4.json :: res.tight.flux.beta.2.margin_FLUX2
\newcommand{\NSAMarginFluxTwo}{\ensuremath{5.955}}
% NSAfourSZero [cert, upper] asm_design_N2ref4.json :: res.tight.S0
\newcommand{\NSAfourSZero}{\ensuremath{7.1878\cdot10^{-3}}}
% NSAfourSOne [cert, upper] asm_design_N2ref4.json :: res.tight.S1
\newcommand{\NSAfourSOne}{\ensuremath{0.10430}}
% NSAMarginCoreSZeroV [cert, margin] asm_design_N2ref4.json :: res.tight.margins.S0v_vs_0.0153
\newcommand{\NSAMarginCoreSZeroV}{\ensuremath{2.116}}
% NSAMarginCoreSOneV [cert, margin] asm_design_N2ref4.json :: res.tight.margins.S1v_vs_0.6
\newcommand{\NSAMarginCoreSOneV}{\ensuremath{5.748}}
% NSAMarginCoreMOne [cert, margin] asm_design_N2ref4.json :: res.tight.margins.M1core
\newcommand{\NSAMarginCoreMOne}{\ensuremath{4.409}}
% NSAMarginCoreMZero [cert, margin] asm_design_N2ref4.json :: res.tight.margins.M0core
\newcommand{\NSAMarginCoreMZero}{\ensuremath{5.450}}
% NSAMarginGhat [cert, margin] asm_design_N2ref4.json :: res.tight.margins.Ghat22
\newcommand{\NSAMarginGhat}{\ensuremath{1.917}}
% NSAMarginRampMOne [cert, margin] asm_design_N2ref4.json :: res.tight.margins.ramp M1
\newcommand{\NSAMarginRampMOne}{\ensuremath{1.184}}
% NSAfourInfoZeroDWindowConst [cert, upper] asm_design_N2ref4.json :: res.tight.info.0.D_window_const
\newcommand{\NSAfourInfoZeroDWindowConst}{\ensuremath{1.7732\cdot10^{-7}}}
% NSAfourInfoOneDWindowConst [cert, upper] asm_design_N2ref4.json :: res.tight.info.1.D_window_const
\newcommand{\NSAfourInfoOneDWindowConst}{\ensuremath{9.8058\cdot10^{-5}}}
% NZDeltaRZero [cert, upper] zone3_interface.json :: DeltaR.0
\newcommand{\NZDeltaRZero}{\ensuremath{4.4986\cdot10^{-8}}}
% NZDeltaROne [cert, upper] zone3_interface.json :: DeltaR.1
\newcommand{\NZDeltaROne}{\ensuremath{6.7972\cdot10^{-6}}}
% NZDeltaRTwo [cert, upper] zone3_interface.json :: DeltaR.2
\newcommand{\NZDeltaRTwo}{\ensuremath{0.011480}}
% NZDeltaRBZero [cert, upper] zone3_interface.json :: DeltaR_b.0
\newcommand{\NZDeltaRBZero}{\ensuremath{2.66\cdot10^{-12}}}
% NZDeltaRBOne [cert, upper] zone3_interface.json :: DeltaR_b.1
\newcommand{\NZDeltaRBOne}{\ensuremath{3.26\cdot10^{-9}}}
% NZDeltaRBTwo [cert, upper] zone3_interface.json :: DeltaR_b.2
\newcommand{\NZDeltaRBTwo}{\ensuremath{4.35\cdot10^{-8}}}
% NZDeltaRDetailNSqZero [cert, upper] zone3_interface.json :: DeltaR_detail.N_sq.0
\newcommand{\NZDeltaRDetailNSqZero}{\ensuremath{8.28\cdot10^{-5}}}
% NZDeltaRDetailNSqOne [cert, upper] zone3_interface.json :: DeltaR_detail.N_sq.1
\newcommand{\NZDeltaRDetailNSqOne}{\ensuremath{0.10149}}
% NZDeltaRDetailNSqTwo [cert, upper] zone3_interface.json :: DeltaR_detail.N_sq.2
\newcommand{\NZDeltaRDetailNSqTwo}{\ensuremath{1.3567}}
% NZDeltaRDetailNSqThree [cert, upper] zone3_interface.json :: DeltaR_detail.N_sq.3
\newcommand{\NZDeltaRDetailNSqThree}{\ensuremath{25350}}
% NZDeltaRDetailPartsTwoPairsThree [cert, upper] zone3_interface.json :: DeltaR_detail.parts.2.pairs.3
\newcommand{\NZDeltaRDetailPartsTwoPairsThree}{\ensuremath{0.011298}}
% NZDeltaRDetailPartsTwoLFourFoot [cert, upper] zone3_interface.json :: DeltaR_detail.parts.2.L4_foot
\newcommand{\NZDeltaRDetailPartsTwoLFourFoot}{\ensuremath{1.6\cdot10^{-4}}}
% NZDeltaRDetailPartsTwoCubic [cert, upper] zone3_interface.json :: DeltaR_detail.parts.2.cubic
\newcommand{\NZDeltaRDetailPartsTwoCubic}{\ensuremath{9.3\cdot10^{-6}}}
% NZDeltaRDetailPartsTwoSThreeThreeFoot [cert, upper] zone3_interface.json :: DeltaR_detail.parts.2.S33_foot
\newcommand{\NZDeltaRDetailPartsTwoSThreeThreeFoot}{\ensuremath{3.7\cdot10^{-7}}}
% NZNIIIZero [cert, upper] zone3_interface.json :: n_III.0
\newcommand{\NZNIIIZero}{\ensuremath{0.071445}}
% NZNIIIOne [cert, upper] zone3_interface.json :: n_III.1
\newcommand{\NZNIIIOne}{\ensuremath{0.32707}}
% NZNIIITwo [cert, upper] zone3_interface.json :: n_III.2
\newcommand{\NZNIIITwo}{\ensuremath{3.0445}}
% NZTauIII [cert, upper] zone3_interface.json :: tau_III
\newcommand{\NZTauIII}{\ensuremath{4.0708}}
% NZoneIIIBoxZero [cert, upper] zone3_interface.json :: box_III_380.b0
\newcommand{\NZoneIIIBoxZero}{\ensuremath{0.10032}}
% NZoneIIIBoxOne [cert, upper] zone3_interface.json :: box_III_380.b1
\newcommand{\NZoneIIIBoxOne}{\ensuremath{0.41974}}
% NZoneIIIBoxZ [cert, upper] zone3_interface.json :: box_III_380.bz
\newcommand{\NZoneIIIBoxZ}{\ensuremath{0.014291}}
% NZBoxIIIThreeEightZeroMarginsBZero [cert, margin] zone3_interface.json :: box_III_380.margins.b0
\newcommand{\NZBoxIIIThreeEightZeroMarginsBZero}{\ensuremath{1.495}}
% NZoneIIIMarginThinnest [cert, margin] zone3_interface.json :: box_III_380.margins.b1
\newcommand{\NZoneIIIMarginThinnest}{\ensuremath{1.238}}
% NZBoxIIIThreeEightZeroMarginsBz [cert, margin] zone3_interface.json :: box_III_380.margins.bz
\newcommand{\NZBoxIIIThreeEightZeroMarginsBz}{\ensuremath{11.19}}
% NZBoxIIISevenZeroZeroBZero [cert, upper] zone3_interface.json :: box_III_700.b0
\newcommand{\NZBoxIIISevenZeroZeroBZero}{\ensuremath{0.22757}}
% NZBoxIIISevenZeroZeroBOne [cert, upper] zone3_interface.json :: box_III_700.b1
\newcommand{\NZBoxIIISevenZeroZeroBOne}{\ensuremath{0.54234}}
% NZUfineKTwoAtBhatRef [cert, upper] zone3_interface.json :: ufine.K2_at_Bhat_ref
\newcommand{\NZUfineKTwoAtBhatRef}{\ensuremath{0.046956}}
% NZUfineKThreeStripAtBhatRef [cert, upper] zone3_interface.json :: ufine.K3_strip_at_Bhat_ref
\newcommand{\NZUfineKThreeStripAtBhatRef}{\ensuremath{0.054140}}
% NRecoveryMinMargin [cert, margin] zone3_interface.json :: cone_family.H6.at_Bhat_ref.min_margin
\newcommand{\NRecoveryMinMargin}{\ensuremath{1.860}}
% NZConeFamilyHSixAtBhatRefConeSlopeMargin [cert, margin] zone3_interface.json :: cone_family.H6.at_Bhat_ref.cone.slope.margin
\newcommand{\NZConeFamilyHSixAtBhatRefConeSlopeMargin}{\ensuremath{2.46}}
% NZConeFamilyHSixAtBhatRefConeKTwoMargin [cert, margin] zone3_interface.json :: cone_family.H6.at_Bhat_ref.cone.K2.margin
\newcommand{\NZConeFamilyHSixAtBhatRefConeKTwoMargin}{\ensuremath{2.54}}
% NZConeFamilyHSixAtBhatRefConeKThreeMargin [cert, margin] zone3_interface.json :: cone_family.H6.at_Bhat_ref.cone.K3.margin
\newcommand{\NZConeFamilyHSixAtBhatRefConeKThreeMargin}{\ensuremath{2.51}}
% NConeTstarMargin [cert, margin] zone3_interface.json :: acceptance.tstar_margin
\newcommand{\NConeTstarMargin}{\ensuremath{2.30}}
% NUoneTFiveYaddLevelOneTwo [cert, upper] u1_ghat.json :: t5_Yadd.level1.2
\newcommand{\NUoneTFiveYaddLevelOneTwo}{\ensuremath{0.0044}}
% NUoneTFiveYaddLevelTwoTwo [cert, upper] u1_ghat.json :: t5_Yadd.level2.2
\newcommand{\NUoneTFiveYaddLevelTwoTwo}{\ensuremath{0.1081}}
% NUoneTFiveYaddLevelTwoThree [cert, upper] u1_ghat.json :: t5_Yadd.level2.3
\newcommand{\NUoneTFiveYaddLevelTwoThree}{\ensuremath{0.374}}
% NUoneWeightThetaGammaZero [cert, upper] u1_ghat.json :: weight.Theta.Gamma.0
\newcommand{\NUoneWeightThetaGammaZero}{\ensuremath{4.00000000000000000000054}}
% NRKZeroZero [cert, upper] ramp_U.json :: K.0.0
\newcommand{\NRKZeroZero}{\ensuremath{13436.03}}
% NRKOneOne [cert, upper] ramp_U.json :: K.1.1
\newcommand{\NRKOneOne}{\ensuremath{10526.19}}
% NRKTwoTwo [cert, upper] ramp_U.json :: K.2.2
\newcommand{\NRKTwoTwo}{\ensuremath{39482.30}}
% NUoneCasesPrimaryYNextZero [cert, upper] u1_ghat.json :: cases.primary.Y_next.0
\newcommand{\NUoneCasesPrimaryYNextZero}{\ensuremath{457.092}}
% NUoneCasesPrimaryYNextOne [cert, upper] u1_ghat.json :: cases.primary.Y_next.1
\newcommand{\NUoneCasesPrimaryYNextOne}{\ensuremath{356.940}}
% NUoneCasesPrimaryYNextTwo [cert, upper] u1_ghat.json :: cases.primary.Y_next.2
\newcommand{\NUoneCasesPrimaryYNextTwo}{\ensuremath{1320.220}}
% NSfourResultsMNineHTwoTightenedFUOneRowsZeroBound [cert, upper] FINAL_manifest_v33.json :: results.M9_H2_tightened_fU1.rows.0.bound
\newcommand{\NSfourResultsMNineHTwoTightenedFUOneRowsZeroBound}{\ensuremath{3.658\cdot10^{-10}}}
% NSfourResultsMNineHTwoTightenedFUOneRowsZeroMargin [cert, margin] FINAL_manifest_v33.json :: results.M9_H2_tightened_fU1.rows.0.margin
\newcommand{\NSfourResultsMNineHTwoTightenedFUOneRowsZeroMargin}{\ensuremath{27.34}}
% NSfourResultsMNineHTwoTightenedFUOneRowsOneBound [cert, upper] FINAL_manifest_v33.json :: results.M9_H2_tightened_fU1.rows.1.bound
\newcommand{\NSfourResultsMNineHTwoTightenedFUOneRowsOneBound}{\ensuremath{1.777\cdot10^{-9}}}
% NSfourResultsMNineHTwoTightenedFUOneRowsOneMargin [cert, margin] FINAL_manifest_v33.json :: results.M9_H2_tightened_fU1.rows.1.margin
\newcommand{\NSfourResultsMNineHTwoTightenedFUOneRowsOneMargin}{\ensuremath{562.8}}
% NSfourResultsMNineHTwoTightenedFUOneRowsTwoBound [cert, upper] FINAL_manifest_v33.json :: results.M9_H2_tightened_fU1.rows.2.bound
\newcommand{\NSfourResultsMNineHTwoTightenedFUOneRowsTwoBound}{\ensuremath{1.665\cdot10^{-6}}}
% NSfourResultsMNineHTwoTightenedFUOneRowsTwoMargin [cert, margin] FINAL_manifest_v33.json :: results.M9_H2_tightened_fU1.rows.2.margin
\newcommand{\NSfourResultsMNineHTwoTightenedFUOneRowsTwoMargin}{\ensuremath{600.9}}
% NHOneAlphaThr [hand, exact] hand :: 
\newcommand{\NHOneAlphaThr}{\ensuremath{1.5\cdot10^{-4}}}
% NHOneSqrtLamThr [hand, exact] hand :: 
\newcommand{\NHOneSqrtLamThr}{\ensuremath{5\cdot10^{-8}}}
% NYbarStep [hand, exact] hand :: 
\newcommand{\NYbarStep}{\ensuremath{1\cdot10^{-12}}}
% NHTwoThrZero [hand, exact] hand :: 
\newcommand{\NHTwoThrZero}{\ensuremath{1\cdot10^{-8}}}
% NHTwoThrOne [hand, exact] hand :: 
\newcommand{\NHTwoThrOne}{\ensuremath{1\cdot10^{-6}}}
% NHTwoThrTwo [hand, exact] hand :: 
\newcommand{\NHTwoThrTwo}{\ensuremath{1\cdot10^{-3}}}
% NHThreeThrZero [hand, exact] hand :: 
\newcommand{\NHThreeThrZero}{\ensuremath{1\cdot10^{-6}}}
% NHThreeThrOne [hand, exact] hand :: 
\newcommand{\NHThreeThrOne}{\ensuremath{5\cdot10^{-6}}}
% NHThreeThrTwo [hand, exact] hand :: 
\newcommand{\NHThreeThrTwo}{\ensuremath{1\cdot10^{-5}}}
% NHFourThrZero [hand, exact] hand :: 
\newcommand{\NHFourThrZero}{\ensuremath{2\cdot10^{-7}}}
% NHFourThrOne [hand, exact] hand :: 
\newcommand{\NHFourThrOne}{\ensuremath{1\cdot10^{-5}}}
% NHFourThrTwo [hand, exact] hand :: 
\newcommand{\NHFourThrTwo}{\ensuremath{1\cdot10^{-5}}}
% NHFourThrThree [hand, exact] hand :: 
\newcommand{\NHFourThrThree}{\ensuremath{1\cdot10^{-5}}}
% NHFourKfourThr [hand, exact] hand :: 
\newcommand{\NHFourKfourThr}{\ensuremath{0.06}}
% NHSixConeSlopeThr [hand, exact] hand :: 
\newcommand{\NHSixConeSlopeThr}{\ensuremath{0.08}}
% NHSixConeKtwoThr [hand, exact] hand :: 
\newcommand{\NHSixConeKtwoThr}{\ensuremath{0.03}}
% NHSixConeNthreeThr [hand, exact] hand :: 
\newcommand{\NHSixConeNthreeThr}{\ensuremath{0.004}}
% NHSixConeJetThr [hand, exact] hand :: 
\newcommand{\NHSixConeJetThr}{\ensuremath{0.12}}
% NHSixConeHFiveThr [hand, exact] hand :: 
\newcommand{\NHSixConeHFiveThr}{\ensuremath{0.0144}}
% NHSixFarHFiveThr [hand, exact] hand :: 
\newcommand{\NHSixFarHFiveThr}{\ensuremath{0.0225}}
% NLemmaOConeM [hand, exact] hand :: 
\newcommand{\NLemmaOConeM}{\ensuremath{0.12}}
% NLemmaOConeK [hand, exact] hand :: 
\newcommand{\NLemmaOConeK}{\ensuremath{0.18}}
% NLemmaOFarM [hand, exact] hand :: 
\newcommand{\NLemmaOFarM}{\ensuremath{0.15}}
% NLemmaOFarK [hand, exact] hand :: 
\newcommand{\NLemmaOFarK}{\ensuremath{0.3}}
% NCoercivityConst [hand, exact] hand :: 
\newcommand{\NCoercivityConst}{\ensuremath{0.046875}}
% NMOnez [hand, exact] hand :: 
\newcommand{\NMOnez}{\ensuremath{0.534}}
% NBetaA [hand, exact] hand :: 
\newcommand{\NBetaA}{\ensuremath{2}}
% NTOneLower [cert, lower] t1clock.json :: t1.0
\newcommand{\NTOneLower}{\ensuremath{4.5978574775}}
% NTOneUpper [cert, upper] t1clock.json :: t1.1
\newcommand{\NTOneUpper}{\ensuremath{4.5978586144}}
% NLambdaOneLower [cert, lower] t1clock.json :: lnlam.0
\newcommand{\NLambdaOneLower}{\ensuremath{6.49\cdot10^{-5}}}
% NLambdaOneUpper [cert, upper] t1clock.json :: lnlam.1
\newcommand{\NLambdaOneUpper}{\ensuremath{6.75\cdot10^{-5}}}
% NTMinusTOneUpper [cert, upper] t1clock.json :: lnlam.1; 1.5e-4.sup_ab
\newcommand{\NTMinusTOneUpper}{\ensuremath{2.1\cdot10^{-17}}}
% NAbarUpper [cert, upper] final_alpha_points_5e-8.json :: 1.5e-4.sup_ab
\newcommand{\NAbarUpper}{\ensuremath{2.19103\cdot10^{-7}}}
% NBetabarUpper [cert, upper] final_alpha_points_5e-8.json :: 1.5e-4.sup_beta
\newcommand{\NBetabarUpper}{\ensuremath{3.33811\cdot10^{-9}}}
% NBaseKmaxLower [cert, lower] t1clock.json :: base_kmax.0
\newcommand{\NBaseKmaxLower}{\ensuremath{2.279\cdot10^{4}}}
% NBaseKmaxUpper [cert, upper] t1clock.json :: base_kmax.1
\newcommand{\NBaseKmaxUpper}{\ensuremath{2.376\cdot10^{4}}}
% NHOneAlphaBase [cert, upper] summary.json :: alpha_abs_max
\newcommand{\NHOneAlphaBase}{\ensuremath{1.1739\cdot10^{-4}}}
% NHOneAlphaErr [cert, upper] t9_h1_clam_s180.json :: C_alpha; inputs.YSTAR
\newcommand{\NHOneAlphaErr}{\ensuremath{7.21\cdot10^{-13}}}
% NHOneAlphaMargin [cert, margin] summary.json :: alpha_abs_max; C_alpha; inputs.YSTAR
\newcommand{\NHOneAlphaMargin}{\ensuremath{1.27}}
% NHOneSqrtLamBase [cert, upper] t9_h1_clam_s180.json :: sqrtLam_base_Wp
\newcommand{\NHOneSqrtLamBase}{\ensuremath{2.282\cdot10^{-8}}}
% NHOneSqrtLamTotal [cert, upper] t9_h1_clam_s180.json :: sqrtLam_total_Wp
\newcommand{\NHOneSqrtLamTotal}{\ensuremath{2.29\cdot10^{-8}}}
% NHOneSqrtLamMargin [cert, margin] t9_h1_clam_s180.json :: sqrtLam_total_Wp
\newcommand{\NHOneSqrtLamMargin}{\ensuremath{2.18}}
% NHOneYbarLimit [cert, lower] t9_h1_clam_s180.json :: Ybar_limit_H1
\newcommand{\NHOneYbarLimit}{\ensuremath{1.12\cdot10^{-9}}}
% NSwitchYUpper [cert, upper] switch_s2.json :: modulated.Y_raw
\newcommand{\NSwitchYUpper}{\ensuremath{6.21\cdot10^{-13}}}
% NSwitchYMargin [cert, margin] switch_s2.json :: modulated.Y_raw
\newcommand{\NSwitchYMargin}{\ensuremath{1.6}}
% NMinCI [cert, lower] constants.json :: meta.summary.min_c_I
\newcommand{\NMinCI}{\ensuremath{0.05283}}
% NNumCertificates [cert, exact] constants.json :: meta.summary.n_cert
\newcommand{\NNumCertificates}{\ensuremath{26}}
% NMinRhoC [cert, lower] constants.json :: meta.closure.rho_c_min
\newcommand{\NMinRhoC}{\ensuremath{0.04711}}
% NHTwoBaseZero [cert, upper] FINAL_manifest_v33.json :: results.M9_H2_tightened.rows.0.base
\newcommand{\NHTwoBaseZero}{\ensuremath{8.011\cdot10^{-11}}}
% NHTwoBaseOne [cert, upper] FINAL_manifest_v33.json :: results.M9_H2_tightened.rows.1.base
\newcommand{\NHTwoBaseOne}{\ensuremath{1.336\cdot10^{-10}}}
% NHTwoBaseTwo [cert, upper] FINAL_manifest_v33.json :: results.M9_H2_tightened.rows.2.base
\newcommand{\NHTwoBaseTwo}{\ensuremath{2.936\cdot10^{-10}}}
% NHTwoErrWZero [cert, upper] FINAL_manifest_v33.json :: results.M9_H2_tightened.rows.0.with_error
\newcommand{\NHTwoErrWZero}{\ensuremath{8.014\cdot10^{-11}}}
% NHTwoMarginWZero [cert, margin] FINAL_manifest_v33.json :: results.M9_H2_tightened.rows.0.margin
\newcommand{\NHTwoMarginWZero}{\ensuremath{124.78}}
% NHTwoErrWOne [cert, upper] FINAL_manifest_v33.json :: results.M9_H2_tightened.rows.1.with_error
\newcommand{\NHTwoErrWOne}{\ensuremath{4.706\cdot10^{-10}}}
% NHTwoMarginWOne [cert, margin] FINAL_manifest_v33.json :: results.M9_H2_tightened.rows.1.margin
\newcommand{\NHTwoMarginWOne}{\ensuremath{2124.9}}
% NHTwoErrWTwo [cert, upper] FINAL_manifest_v33.json :: results.M9_H2_tightened.rows.2.with_error
\newcommand{\NHTwoErrWTwo}{\ensuremath{1.559\cdot10^{-6}}}
% NHTwoMarginWTwo [cert, margin] FINAL_manifest_v33.json :: results.M9_H2_tightened.rows.2.margin
\newcommand{\NHTwoMarginWTwo}{\ensuremath{641.53}}
% NHThreeBaseZero [cert, upper] FINAL_manifest_v33.json :: results.M9_H2f.rows.0.base
\newcommand{\NHThreeBaseZero}{\ensuremath{3.127\cdot10^{-7}}}
% NHThreeBaseOne [cert, upper] FINAL_manifest_v33.json :: results.M9_H2f.rows.1.base
\newcommand{\NHThreeBaseOne}{\ensuremath{1.004\cdot10^{-6}}}
% NHThreeBaseTwo [cert, upper] FINAL_manifest_v33.json :: results.M9_H2f.rows.2.base
\newcommand{\NHThreeBaseTwo}{\ensuremath{2.431\cdot10^{-6}}}
% NHThreeErrZero [cert, upper] FINAL_manifest_v33.json :: results.M9_H2f.rows.0.with_error
\newcommand{\NHThreeErrZero}{\ensuremath{3.1363\cdot10^{-7}}}
% NHThreeMarginZero [cert, margin] FINAL_manifest_v33.json :: results.M9_H2f.rows.0.margin
\newcommand{\NHThreeMarginZero}{\ensuremath{3.1885}}
% NHThreeErrOne [cert, upper] FINAL_manifest_v33.json :: results.M9_H2f.rows.1.with_error
\newcommand{\NHThreeErrOne}{\ensuremath{1.0034\cdot10^{-6}}}
% NHThreeMarginOne [cert, margin] FINAL_manifest_v33.json :: results.M9_H2f.rows.1.margin
\newcommand{\NHThreeMarginOne}{\ensuremath{4.9834}}
% NHThreeErrTwo [cert, upper] FINAL_manifest_v33.json :: results.M9_H2f.rows.2.with_error
\newcommand{\NHThreeErrTwo}{\ensuremath{2.4351\cdot10^{-6}}}
% NHThreeMarginTwo [cert, margin] FINAL_manifest_v33.json :: results.M9_H2f.rows.2.margin
\newcommand{\NHThreeMarginTwo}{\ensuremath{4.1067}}
% NHFourBaseZero [cert, upper] FINAL_manifest_v33.json :: results.M9_base_sums.combined.windows_max.0
\newcommand{\NHFourBaseZero}{\ensuremath{3.48\cdot10^{-8}}}
% NHFourErrZero [cert, upper] FINAL_manifest_v33.json :: results.M9_H3_H2v32.rows.win.0.with_error_arb
\newcommand{\NHFourErrZero}{\ensuremath{3.5087\cdot10^{-8}}}
% NHFourMarginZero [cert, margin] FINAL_manifest_v33.json :: results.M9_H3_H2v32.rows.win.0.margin_arb
\newcommand{\NHFourMarginZero}{\ensuremath{5.700}}
% NHFourBaseOne [cert, upper] FINAL_manifest_v33.json :: results.M9_base_sums.combined.windows_max.1
\newcommand{\NHFourBaseOne}{\ensuremath{1.27\cdot10^{-7}}}
% NHFourErrOne [cert, upper] FINAL_manifest_v33.json :: results.M9_H3_H2v32.rows.win.1.with_error_arb
\newcommand{\NHFourErrOne}{\ensuremath{1.2616\cdot10^{-7}}}
% NHFourMarginOne [cert, margin] FINAL_manifest_v33.json :: results.M9_H3_H2v32.rows.win.1.margin_arb
\newcommand{\NHFourMarginOne}{\ensuremath{79.26}}
% NHFourBaseTwo [cert, upper] FINAL_manifest_v33.json :: results.M9_base_sums.combined.windows_max.2
\newcommand{\NHFourBaseTwo}{\ensuremath{3.76\cdot10^{-7}}}
% NHFourErrTwo [cert, upper] FINAL_manifest_v33.json :: results.M9_H3_H2v32.rows.win.2.with_error_arb
\newcommand{\NHFourErrTwo}{\ensuremath{3.7724\cdot10^{-7}}}
% NHFourMarginTwo [cert, margin] FINAL_manifest_v33.json :: results.M9_H3_H2v32.rows.win.2.margin_arb
\newcommand{\NHFourMarginTwo}{\ensuremath{26.50}}
% NHFourBaseThree [cert, upper] FINAL_manifest_v33.json :: results.M9_base_sums.combined.windows_max.3
\newcommand{\NHFourBaseThree}{\ensuremath{9.71\cdot10^{-7}}}
% NHFourErrThree [cert, upper] FINAL_manifest_v33.json :: results.M9_H3_H2v32.rows.win.3.with_error_arb
\newcommand{\NHFourErrThree}{\ensuremath{1.2822\cdot10^{-6}}}
% NHFourMarginThree [cert, margin] FINAL_manifest_v33.json :: results.M9_H3_H2v32.rows.win.3.margin_arb
\newcommand{\NHFourMarginThree}{\ensuremath{7.799}}
% NHFourYbarLimit [cert, lower] FINAL_manifest_v33.json :: results.M9_H3_H2v32.binding_Ybar_limit
\newcommand{\NHFourYbarLimit}{\ensuremath{1.575\cdot10^{-11}}}
% NHFourYbarMargin [cert, margin] FINAL_manifest_v33.json :: results.M9_H3_H2v32.binding_Ybar_limit
\newcommand{\NHFourYbarMargin}{\ensuremath{15.75}}
% NHFourKfourBase [cert, upper] amp_A1_w180.json :: inputs.BASE_T1.K4
\newcommand{\NHFourKfourBase}{\ensuremath{3.1633\cdot10^{-4}}}
% NHSixConeSlope [cert, upper] FINAL_manifest_v33.json :: results.cone_fixedL.checks.slope.value
\newcommand{\NHSixConeSlope}{\ensuremath{0.03029}}
% NHSixConeSlopeMargin [cert, margin] FINAL_manifest_v33.json :: results.cone_fixedL.checks.slope.margin
\newcommand{\NHSixConeSlopeMargin}{\ensuremath{2.641}}
% NHSixConeKtwo [cert, upper] FINAL_manifest_v33.json :: results.cone_fixedL.checks.K2.value
\newcommand{\NHSixConeKtwo}{\ensuremath{0.009876}}
% NHSixConeKtwoMargin [cert, margin] FINAL_manifest_v33.json :: results.cone_fixedL.checks.K2.margin
\newcommand{\NHSixConeKtwoMargin}{\ensuremath{3.037}}
% NHSixConeNthree [cert, upper] FINAL_manifest_v33.json :: results.cone_fixedL.checks.N3.value
\newcommand{\NHSixConeNthree}{\ensuremath{1.461\cdot10^{-4}}}
% NHSixConeNthreeMargin [cert, margin] FINAL_manifest_v33.json :: results.cone_fixedL.checks.N3.margin
\newcommand{\NHSixConeNthreeMargin}{\ensuremath{27.38}}
% NHSixConeHFive [cert, upper] FINAL_manifest_v33.json :: results.H4_fixedL.total
\newcommand{\NHSixConeHFive}{\ensuremath{1.2865\cdot10^{-5}}}
% NHSixConeHFiveMargin [cert, margin] FINAL_manifest_v33.json :: results.H4_fixedL.margin
\newcommand{\NHSixConeHFiveMargin}{\ensuremath{1119}}
% NHSixConeTstarNeeded [cert, upper] FINAL_manifest_v33.json :: results.H7_fixedL.ratio_upper
\newcommand{\NHSixConeTstarNeeded}{\ensuremath{1.6000015\cdot10^{-7}}}
% NHSixConeTstarLower [cert, lower] FINAL_manifest_v33.json :: results.H7_fixedL.t_star_lower
\newcommand{\NHSixConeTstarLower}{\ensuremath{3.68334\cdot10^{-7}}}
% NHSixConeTstarMargin [cert, margin] FINAL_manifest_v33.json :: results.H7_fixedL.margin_tstar_over_ratio
\newcommand{\NHSixConeTstarMargin}{\ensuremath{2.302}}
% NHSixConeHFiveRatioPeak [cert, upper] h4loc_A1_w180.json :: windows.constants/amp_A1_w180.json.worst.ratio
\newcommand{\NHSixConeHFiveRatioPeak}{\ensuremath{0.2393}}
% NHSixConeHOneSup [cert, upper] FINAL_manifest_v33.json :: results.H1_fixedL.sup_h0_max
\newcommand{\NHSixConeHOneSup}{\ensuremath{0.2163}}
% NHSixConeHOneH [cert, upper] FINAL_manifest_v33.json :: results.H1_fixedL.H_max
\newcommand{\NHSixConeHOneH}{\ensuremath{0.2185}}
% NHSixFarHFiveRatio [cert, upper] h4loc_v32.json :: windows.constants/amp_A1_w180.json.cylinders.18.ratio
\newcommand{\NHSixFarHFiveRatio}{\ensuremath{0.1761}}
% NHSixFarHFiveMarginOld [cert, margin] h4loc_v32.json :: windows.constants/amp_A1_w180.json.cylinders.18.ratio
\newcommand{\NHSixFarHFiveMarginOld}{\ensuremath{5.68}}
% NHSixFarHFiveMargin [cert, margin] FINAL_manifest_v33.json :: results.M9_H4_far.min_margin
\newcommand{\NHSixFarHFiveMargin}{\ensuremath{15.58}}
% NHSixConeHFiveMarginOld [cert, margin] h4loc_A1_w180.json :: windows.constants/amp_A1_w180.json.worst.ratio
\newcommand{\NHSixConeHFiveMarginOld}{\ensuremath{4.18}}
% NHSixFarKtwo [cert, upper] t9_far_v32.json :: h2
\newcommand{\NHSixFarKtwo}{\ensuremath{1.797\cdot10^{-4}}}
% NHSixFarKtwoThr [hand, exact] t9_far_v32.json :: thr.h2
\newcommand{\NHSixFarKtwoThr}{\ensuremath{3.53\cdot10^{-4}}}
% NHSixFarKtwoMargin [cert, margin] t9_far_v32.json :: margins.h2
\newcommand{\NHSixFarKtwoMargin}{\ensuremath{1.964}}
% NHSixFarKthree [cert, upper] t9_far_v32.json :: h3
\newcommand{\NHSixFarKthree}{\ensuremath{2.595\cdot10^{-8}}}
% NHSixFarKthreeThr [hand, exact] t9_far_v32.json :: thr.h3
\newcommand{\NHSixFarKthreeThr}{\ensuremath{1.17\cdot10^{-5}}}
% NHSixFarKthreeMargin [cert, margin] t9_far_v32.json :: margins.h3
\newcommand{\NHSixFarKthreeMargin}{\ensuremath{450.9}}
% NHSixFarKfour [cert, upper] t9_far_v32.json :: h4
\newcommand{\NHSixFarKfour}{\ensuremath{6.964\cdot10^{-9}}}
% NHSixFarKfourThr [hand, exact] t9_far_v32.json :: thr.h4
\newcommand{\NHSixFarKfourThr}{\ensuremath{1.03\cdot10^{-7}}}
% NHSixFarKfourMargin [cert, margin] t9_far_v32.json :: margins.h4
\newcommand{\NHSixFarKfourMargin}{\ensuremath{14.79}}
% NHSixFarSlope [cert, upper] t9_far_v32.json :: slope
\newcommand{\NHSixFarSlope}{\ensuremath{0.04153}}
% NHSixFarSlopeThr [hand, exact] t9_far_v32.json :: thr.slope
\newcommand{\NHSixFarSlopeThr}{\ensuremath{0.0800}}
% NHSixFarSlopeMargin [cert, margin] t9_far_v32.json :: margins.slope
\newcommand{\NHSixFarSlopeMargin}{\ensuremath{1.926}}
% NHSixFarNthree [cert, upper] t9_far_v32.json :: L2
\newcommand{\NHSixFarNthree}{\ensuremath{5.516\cdot10^{-7}}}
% NHSixFarNthreeThr [hand, exact] t9_far_v32.json :: thr.L2
\newcommand{\NHSixFarNthreeThr}{\ensuremath{3.32\cdot10^{-6}}}
% NHSixFarNthreeMargin [cert, margin] t9_far_v32.json :: margins.L2
\newcommand{\NHSixFarNthreeMargin}{\ensuremath{6.019}}
% NFarTstarMargin [cert, margin] lemmaO_v32_consts.json :: far_v32.case.margin_factor_lower
\newcommand{\NFarTstarMargin}{\ensuremath{15.5}}
% NFarTstarLower [cert, lower] lemmaO_v32_consts.json :: far_v32.case.t_star_lower
\newcommand{\NFarTstarLower}{\ensuremath{9.5108\cdot10^{-8}}}
% NFarTstarNeeded [cert, upper] lemmaO_v32_consts.json :: far_v32.case.needed
\newcommand{\NFarTstarNeeded}{\ensuremath{6.134\cdot10^{-9}}}
% NGammaModes [cert, exact] gamma0_final.json :: M
\newcommand{\NGammaModes}{\ensuremath{3762}}
% NGammaHalfwidth [cert, upper] gamma0_final.json :: b_j_halfwidth
\newcommand{\NGammaHalfwidth}{\ensuremath{8.138\cdot10^{-24}}}
% NGammaContraction [cert, upper] gamma0_final.json :: krawczyk_contraction
\newcommand{\NGammaContraction}{\ensuremath{3.01\cdot10^{-12}}}
% NGammaArcs [cert, exact] gamma0_final.json :: embedded.arcs
\newcommand{\NGammaArcs}{\ensuremath{701}}
% NGammaMaxTurn [cert, upper] gamma0_final.json :: embedded.maxturn
\newcommand{\NGammaMaxTurn}{\ensuremath{0.04997}}
% NGammaMaxPair [cert, upper] gamma0_final.json :: embedded.maxpair
\newcommand{\NGammaMaxPair}{\ensuremath{0.09597}}
% NGammaMinSep [cert, lower] gamma0_final.json :: embedded.minsep
\newcommand{\NGammaMinSep}{\ensuremath{8.89\cdot10^{-3}}}
% NGammaHalfGap [cert, lower] gamma0_final.json :: embedded.gap_cert
\newcommand{\NGammaHalfGap}{\ensuremath{0.0248219}}
% NPhaseISlabs [cert, exact] status.json :: slab
\newcommand{\NPhaseISlabs}{\ensuremath{96}}
% NPhaseIModes [cert, exact] status.json :: last_slab.M
\newcommand{\NPhaseIModes}{\ensuremath{2800}}
% NPhaseIKmaxLower [cert, lower] status.json :: kmax_enclosure.0
\newcommand{\NPhaseIKmaxLower}{\ensuremath{29.352874545}}
% NPhaseIKmaxUpper [cert, upper] status.json :: kmax_enclosure.1
\newcommand{\NPhaseIKmaxUpper}{\ensuremath{29.352874546}}
% NPhaseITLower [cert, lower] status.json :: t
\newcommand{\NPhaseITLower}{\ensuremath{4.597850539485830}}
% NPhaseITUpper [cert, upper] status.json :: t
\newcommand{\NPhaseITUpper}{\ensuremath{4.597850539485831}}
% NPeakSqrtLam [cert, upper] rate_table_fin12f4_a0.00015.json :: sqrtL_max
\newcommand{\NPeakSqrtLam}{\ensuremath{5.308\cdot10^{-8}}}
% NTailRate [cert, lower] rate_table_fin12f4_a0.00015.json :: tail.rho
\newcommand{\NTailRate}{\ensuremath{2.693\cdot10^{-3}}}
% NFloorSqrtLam [cert, upper] rate_table_fin12f4_a0.00015.json :: tail.sqrt_LF
\newcommand{\NFloorSqrtLam}{\ensuremath{9.872\cdot10^{-10}}}
% NAlphaFloor [cert, upper] rate_table_fin12f4_a0.00015.json :: tail.alpha_F
\newcommand{\NAlphaFloor}{\ensuremath{6.867\cdot10^{-7}}}
% NSupBetaVerifier [cert, upper] final_alpha_points_5e-8.json :: 1.5e-4.sup_beta
\newcommand{\NSupBetaVerifier}{\ensuremath{3.33811\cdot10^{-9}}}
% NMOneRamp [cert, upper] final_alpha_points_5e-8.json :: 1.5e-4.M1
\newcommand{\NMOneRamp}{\ensuremath{0.4328}}
% NSITwoPrime [cert, upper] FINAL_manifest_v33.json :: results.F12_sliver.sI2_prime
\newcommand{\NSITwoPrime}{\ensuremath{0.015457}}
% NLoopGainCoarse [cert, upper] U_cert_v35.json :: loop_gain.2.total
\newcommand{\NLoopGainCoarse}{\ensuremath{0.03904}}
% NRateMarginCoarse [cert, margin] U_cert_v35.json :: rate_margin
\newcommand{\NRateMarginCoarse}{\ensuremath{28.2}}
% NSOneEpsLower [cert, lower] bracket2.json :: eps_a
\newcommand{\NSOneEpsLower}{\ensuremath{1.5346198892965\cdot10^{-4}}}
% NSOneEpsUpper [cert, upper] bracket2.json :: eps_b
\newcommand{\NSOneEpsUpper}{\ensuremath{1.5346198899299\cdot10^{-4}}}
% NSOneRowsTip [cert, exact] ext_tip_check.json :: rows
\newcommand{\NSOneRowsTip}{\ensuremath{51246}}
% NSOneTipWidth [cert, upper] ext_tip_check.json :: max_pointwise_width.th
\newcommand{\NSOneTipWidth}{\ensuremath{3.2\cdot10^{-12}}}
% NSOneExtWidth [cert, upper] ext813_check.json :: max_pointwise_width.p
\newcommand{\NSOneExtWidth}{\ensuremath{5.7\cdot10^{-10}}}
% NSOneRowsMC [cert, exact] check_outA5.json :: prof_tip_vs_MC.rows_compared
\newcommand{\NSOneRowsMC}{\ensuremath{1025}}
% NSOneRowsMB [cert, exact] check_outA5.json :: ext813_vs_MB.rows_compared
\newcommand{\NSOneRowsMB}{\ensuremath{501}}
% NSOneMinMargin [cert, margin] check_outA5.json :: min:*.min_relative_margin.*
\newcommand{\NSOneMinMargin}{\ensuremath{1.2\cdot10^{-9}}}
% NTwoSegments [cert, exact] contour_result.json :: len:segments
\newcommand{\NTwoSegments}{\ensuremath{5058}}
% NTwoInteriorBoxes [cert, exact] interior_result.json :: len:
\newcommand{\NTwoInteriorBoxes}{\ensuremath{174}}
% NTwoContourBoxes [cert, exact] eta_cache.json :: len:
\newcommand{\NTwoContourBoxes}{\ensuremath{202}}
% NTwoEtaMax [cert, upper] eta_cache.json :: max:*.eta
\newcommand{\NTwoEtaMax}{\ensuremath{0.1163}}
% NTwoLipMax [cert, upper] eta_cache.json :: max:*.Lip
\newcommand{\NTwoLipMax}{\ensuremath{0.272}}
% NTwoLambdaZero [est, upper] cert_b0.25_th0.0_tau0.03.json :: Lam
\newcommand{\NTwoLambdaZero}{\ensuremath{1.790819}}
% NTwoLambdaOne [est, upper] cert_b0.25_th0.48_tau0.03.json :: Lam
\newcommand{\NTwoLambdaOne}{\ensuremath{1.768282}}
% NSOneDTightLower [doc, lower] doc:README.md :: D ∈ [2.50080627654, 2.50080627668]
\newcommand{\NSOneDTightLower}{\ensuremath{2.50080627654}}
% NSOneDTightUpper [doc, upper] doc:README.md :: D ∈ [2.50080627654, 2.50080627668]
\newcommand{\NSOneDTightUpper}{\ensuremath{2.50080627668}}
% NSOneKTightLower [doc, lower] doc:README.md :: k(0) ∈ [1.53982063762, 1.53982063822]
\newcommand{\NSOneKTightLower}{\ensuremath{1.53982063762}}
% NSOneKTightUpper [doc, upper] doc:README.md :: k(0) ∈ [1.53982063762, 1.53982063822]
\newcommand{\NSOneKTightUpper}{\ensuremath{1.53982063822}}
% NSOneDCondLower [doc, lower] doc:README.md :: D = −p(tip) ∈ [2.500806276619695, 2.500806276620049]
\newcommand{\NSOneDCondLower}{\ensuremath{2.500806276619695}}
% NSOneDCondUpper [doc, upper] doc:README.md :: D = −p(tip) ∈ [2.500806276619695, 2.500806276620049]
\newcommand{\NSOneDCondUpper}{\ensuremath{2.500806276620049}}
% NSOneKCondLower [doc, lower] doc:README.md :: k(0) ∈ [1.539820637871836, 1.5398206378721015]
\newcommand{\NSOneKCondLower}{\ensuremath{1.539820637871836}}
% NSOneKCondUpper [doc, upper] doc:README.md :: k(0) ∈ [1.539820637871836, 1.5398206378721015]
\newcommand{\NSOneKCondUpper}{\ensuremath{1.5398206378721015}}
% NSOneTauTipLower [doc, lower] doc:README.md :: τ_tip ∈ [10.249077651120778, 10.249077651121361]
\newcommand{\NSOneTauTipLower}{\ensuremath{10.249077651120778}}
% NSOneTauTipUpper [doc, upper] doc:README.md :: τ_tip ∈ [10.249077651120778, 10.249077651121361]
\newcommand{\NSOneTauTipUpper}{\ensuremath{10.249077651121361}}
% NSOneStateUncert [doc, upper] doc:README.md :: State uncertainty at σ = 8: 4.4e−13
\newcommand{\NSOneStateUncert}{\ensuremath{4.4\cdot10^{-13}}}
% NFarBoxEprime [doc, margin] doc:STEP4_RESULT.md :: (6.4287, 4.6917, 1.9940, 1.69999995)
\newcommand{\NFarBoxEprime}{\ensuremath{6.42}}
% NFarBoxXEsecond [doc, margin] doc:STEP4_RESULT.md :: (6.4287, 4.6917, 1.9940, 1.69999995)
\newcommand{\NFarBoxXEsecond}{\ensuremath{4.69}}
% NFarBoxXsqEthird [doc, margin] doc:STEP4_RESULT.md :: (6.4287, 4.6917, 1.9940, 1.69999995)
\newcommand{\NFarBoxXsqEthird}{\ensuremath{1.99}}
% NFarBoxE [doc, margin] doc:STEP4_RESULT.md :: (6.4287, 4.6917, 1.9940, 1.69999995)
\newcommand{\NFarBoxE}{\ensuremath{1.69}}
% NHFourKfourErr [doc, upper] doc:STEP4_RESULT.md :: | 5.60e-3 against 0.06 | ×10.7 |
\newcommand{\NHFourKfourErr}{\ensuremath{5.60\cdot10^{-3}}}
% NHFourKfourMargin [doc, margin] doc:STEP4_RESULT.md :: | 5.60e-3 against 0.06 | ×10.7 |
\newcommand{\NHFourKfourMargin}{\ensuremath{10.7}}
% NCoercivityCoefficient [doc, upper] doc:step3_proof_v32.tex :: \norm{e^{\tau\mathcal L}u}_\rho\le14.8\,e^{-0.0213\tau}\norm
\newcommand{\NCoercivityCoefficient}{\ensuremath{14.8}}
% NCoercivityRate [doc, lower] doc:step3_proof_v32.tex :: \norm{e^{\tau\mathcal L}u}_\rho\le14.8\,e^{-0.0213\tau}\norm
\newcommand{\NCoercivityRate}{\ensuremath{0.0213}}
% NRegionRampStart [hand, exact] hand :: 
\newcommand{\NRegionRampStart}{\ensuremath{14}}
% NRegionInterface [hand, exact] hand :: 
\newcommand{\NRegionInterface}{\ensuremath{28}}
% NRegionShellEnd [hand, exact] hand :: 
\newcommand{\NRegionShellEnd}{\ensuremath{60}}
% NRegionSegEnd [hand, exact] hand :: 
\newcommand{\NRegionSegEnd}{\ensuremath{350}}
% NRegionAnnEnd [hand, exact] hand :: 
\newcommand{\NRegionAnnEnd}{\ensuremath{700}}
% NZoneIIILow [hand, exact] hand :: 
\newcommand{\NZoneIIILow}{\ensuremath{253}}
% NZoneIIIHigh [hand, exact] hand :: 
\newcommand{\NZoneIIIHigh}{\ensuremath{379.5}}
% NConeCentreMin [hand, exact] hand :: 
\newcommand{\NConeCentreMin}{\ensuremath{300}}
% NConeCentreMax [hand, exact] hand :: 
\newcommand{\NConeCentreMax}{\ensuremath{380}}
% NFarCentreMin [hand, exact] hand :: 
\newcommand{\NFarCentreMin}{\ensuremath{490}}
% NConeL [hand, exact] hand :: 
\newcommand{\NConeL}{\ensuremath{50}}
% NConeX [hand, exact] hand :: 
\newcommand{\NConeX}{\ensuremath{3}}
% NFarL [hand, exact] hand :: 
\newcommand{\NFarL}{\ensuremath{113}}
% NFarX [hand, exact] hand :: 
\newcommand{\NFarX}{\ensuremath{2}}
}%
  {\PackageWarning{paper}{numbers/numbers_paper.tex not found (unit U0 not finished)}}

% ======================================================================
% NOTATION MACROS (OUTLINE 5.2 and notation.md).  One macro per symbol
% whose spelling is fixed by the notation table.  Plain letters (a, b, f, ...)
% have no macro.
% ======================================================================
% basic
\newcommand{\R}{\mathbb{R}}
\newcommand{\C}{\mathbb{C}}
\newcommand{\N}{\mathbb{N}}
\newcommand{\Z}{\mathbb{Z}}
\newcommand{\eps}{\varepsilon}
\newcommand{\ip}[2]{\langle #1,#2\rangle}
\newcommand{\abs}[1]{\lvert #1\rvert}
\newcommand{\norm}[1]{\lVert #1\rVert}
\newcommand{\dd}{\mathrm d}
\newcommand{\re}{\operatorname{Re}}
\newcommand{\im}{\operatorname{Im}}
% flow and curves
\newcommand{\arclen}{\mathfrak{s}}                 % arclength of the flow curve
\newcommand{\kss}{\kappa_{\mathfrak{ss}}}          % kappa_{ss} in the flow equation
\newcommand{\Area}{\operatorname{Area}}
\newcommand{\Len}{\mathcal L}                      % length of a curve
\newcommand{\spp}{\operatorname{sp}}                 % sup quantities sp^v_0, sp^v_1, sp^f_1 (were S^v_0, S^v_1, S^f_1)
\newcommand{\rot}{\operatorname{rot}}              % turning number
\newcommand{\Tsing}{T_*}                           % singular time
\newcommand{\Tmax}{T_{\max}}                       % maximal time of existence
\newcommand{\pstar}{p_*}                           % image point of the singularity
\newcommand{\Gz}{\Gamma_0}                         % initial curve
\newcommand{\vth}{\vartheta}                       % parameter of Gamma_0
\newcommand{\dclk}{\delta_t}                       % clock defect of step 4
\newcommand{\fall}{\varpi}                         % Fourier rates varpi_n = n^4-n^2
% rescaling
\newcommand{\lamone}{\lambda_1}
\newcommand{\rate}{\varsigma}                      % 1/4 - a
\newcommand{\xia}{\xi_a}                           % (1-4\bar a)^{-1}
% shrinker and step 1
\newcommand{\Sone}{S_1}
\newcommand{\Gsh}{\Gamma_{\!*}}                    % profile curve of S_1
\newcommand{\cSet}{\mathcal S}                     % certified set of shrinkers
\newcommand{\sigtip}{\sigma_{\rm tip}}
\newcommand{\Xo}{X_0}                              % 2.2490776515
\newcommand{\Dtip}{D}
\newcommand{\ktip}{k_0}
\newcommand{\shootM}{M}                            % shooting function
\newcommand{\Ypar}{\mathcal Y}                     % step-1 parallelepiped
\newcommand{\axis}{\Sigma}                         % symmetry axis of S_1
\newcommand{\sigc}{\sigma_{\rm c}}
\newcommand{\sigrev}{\sigma^{\rm rev}}
% linear theory
\newcommand{\Hrho}{\mathcal H_\rho}
\newcommand{\muslow}{\mu_3}
\newcommand{\rre}{\mathfrak r_{\rm re}}
\newcommand{\rim}{\mathfrak r_{\rm im}}
\newcommand{\Evev}{\mathrm{Ev}_{\rm even}}
\newcommand{\Evod}{\mathrm{Ev}_{\rm odd}}
\newcommand{\gvec}{\mathbf g}                      % translation vector in step 2
\newcommand{\cons}[1]{\mathfrak f_{#1}}            % constraint functions
\newcommand{\phirho}{\varphi_\rho}
% energies
\newcommand{\Lam}{\Lambda}
\newcommand{\En}{E}
\newcommand{\EnI}{E_I}
\newcommand{\muI}{\mu_I}                           % 0.08
\newcommand{\epsE}{\varepsilon_E}                  % 1e-5
\newcommand{\epsI}{\varepsilon_I}                  % 3e-5
\newcommand{\wn}[1]{\mathsf w_{#1}}                % ||w~^{(j)}||_rho
\newcommand{\vI}{v_I}                              % zeta w'
\newcommand{\muerr}{\mu_{\rm err}}                 % step-4 error-energy weight 0.036
% far field
\newcommand{\pU}{p_U}                              % 1.0246
\newcommand{\Uone}{U_1}
\newcommand{\sm}{\operatorname{sm}}                % smoothstep t^3(10-15t+6t^2)
\newcommand{\rc}{r_c}                              % centre radius (Lemma U)
\newcommand{\rchi}{r_\chi}                         % 253
\newcommand{\Iw}{I_w}                              % hand-off windows [y_w,1.5y_w]
\newcommand{\qa}{\mathfrak q_1}                    % quasilinear coefficients of the slope equation
\newcommand{\qb}{\mathfrak q_2}
% cylinders (Lemma O)
\newcommand{\cyl}{\ell}                            % cylinder length unit
\newcommand{\MO}{M_O}
\newcommand{\KO}{K_O}
\newcommand{\Hcyl}{H_{\rm cyl}}
\newcommand{\XiO}{\Xi_O}
\newcommand{\fO}{f_O}
% step 4
\newcommand{\Ybar}{\bar Y}
\newcommand{\Lt}{\widetilde L}                     % linearisation about the step-4 base
\newcommand{\Trans}{\Pi}                           % transport integral
\newcommand{\resid}{\mathrm{res}}
\newcommand{\Resid}[1]{\mathrm{Res}_{#1}}
\newcommand{\boxtip}{\Box_{\rm tip}}
\newcommand{\boxfar}{\Box_{\rm far}}
\newcommand{\Eenv}{\mathcal E_{\rm env}}
\newcommand{\betab}{\beta_{\rm base}}
\newcommand{\tauI}{\tau^{\rm I}}                   % Phase I clock
% names of families of conditions: \Hh{3} = (H3), \Oh{4} = (O4), ...
\newcommand{\Hh}[1]{\textup{(H#1)}}
\newcommand{\Oh}[1]{\textup{(O#1)}}
\newcommand{\Lh}[1]{\textup{(L#1)}}
\newcommand{\Bh}[1]{\textup{(B#1)}}
\newcommand{\Kh}[1]{\textup{(K#1)}}
\newcommand{\Wh}[1]{\textup{(W#1)}}
\newcommand{\PhaseI}{Phase~I}
\newcommand{\PhaseII}{Phase~II}
% margin: threshold/bound (always)
\newcommand{\mgn}[1]{{\times}#1}

% ---- figures: placeholder macro so that drafts compile before F1-F7 exist ----
\newcommand{\figplaceholder}[1]{\fbox{\parbox{0.8\linewidth}{\centering\small[figure #1 to be supplied]}}}
 ...
% (the main paper does nothing special: \ifsupplement defaults to false).
% Compile from paper_v1/ so that % numbers_paper.tex -- GENERATED from ledger.json by build_ledger.py; do not edit.
% Rounding: upper bounds up, lower bounds and margins down at the printed digit; margin = threshold / bound.
% Use as $\times\NName$ for margins.  kind: cert certified, est estimate, hand hypothesis constant, doc only a document quotes it.
% NUEhatZero [cert, upper] U_cert_v35.json :: supersolution.EA.0
\newcommand{\NUEhatZero}{\ensuremath{1.2018\cdot10^{-6}}}
% NUEhatOne [cert, upper] U_cert_v35.json :: supersolution.EA.1
\newcommand{\NUEhatOne}{\ensuremath{1.0441\cdot10^{-4}}}
% NUEhatTwo [cert, upper] U_cert_v35.json :: supersolution.EA.2
\newcommand{\NUEhatTwo}{\ensuremath{0.056694}}
% NUBhatZero [cert, upper] U_cert_v35.json :: supersolution.Bhat.0
\newcommand{\NUBhatZero}{\ensuremath{4.8735\cdot10^{-6}}}
% NUBhatOne [cert, upper] U_cert_v35.json :: supersolution.Bhat.1
\newcommand{\NUBhatOne}{\ensuremath{5.5933\cdot10^{-4}}}
% NUBhatTwo [cert, upper] U_cert_v35.json :: supersolution.Bhat.2
\newcommand{\NUBhatTwo}{\ensuremath{0.36675}}
% NUSupersolutionCZeroZero [cert, upper] U_cert_v35.json :: supersolution.C0.0
\newcommand{\NUSupersolutionCZeroZero}{\ensuremath{1.1340\cdot10^{-6}}}
% NUSupersolutionCZeroOne [cert, upper] U_cert_v35.json :: supersolution.C0.1
\newcommand{\NUSupersolutionCZeroOne}{\ensuremath{1.6623\cdot10^{-4}}}
% NUSupersolutionCOneZero [cert, upper] U_cert_v35.json :: supersolution.C1.0
\newcommand{\NUSupersolutionCOneZero}{\ensuremath{1.5241\cdot10^{-9}}}
% NUSupersolutionCOneOne [cert, upper] U_cert_v35.json :: supersolution.C1.1
\newcommand{\NUSupersolutionCOneOne}{\ensuremath{5.5255\cdot10^{-8}}}
% NUSupersolutionNTwoZero [cert, upper] U_cert_v35.json :: supersolution.N2.0
\newcommand{\NUSupersolutionNTwoZero}{\ensuremath{0.071472}}
% NUSupersolutionNTwoOne [cert, upper] U_cert_v35.json :: supersolution.N2.1
\newcommand{\NUSupersolutionNTwoOne}{\ensuremath{0.33012}}
% NUSupersolutionNTwoTwo [cert, upper] U_cert_v35.json :: supersolution.N2.2
\newcommand{\NUSupersolutionNTwoTwo}{\ensuremath{5.0406}}
% NUSupersolutionNTwoOverFourNTwoRefZero [cert, upper] U_cert_v35.json :: supersolution.N2_over_4N2ref.0
\newcommand{\NUSupersolutionNTwoOverFourNTwoRefZero}{\ensuremath{0.18}}
% NUSupersolutionNTwoOverFourNTwoRefOne [cert, upper] U_cert_v35.json :: supersolution.N2_over_4N2ref.1
\newcommand{\NUSupersolutionNTwoOverFourNTwoRefOne}{\ensuremath{0.42}}
% NUSupersolutionNTwoOverFourNTwoRefTwo [cert, upper] U_cert_v35.json :: supersolution.N2_over_4N2ref.2
\newcommand{\NUSupersolutionNTwoOverFourNTwoRefTwo}{\ensuremath{0.26}}
% NUSupersolutionMOneZHypCheckZonesIIISOneTwoBhat [cert, upper] U_cert_v35.json :: supersolution.M1z_hyp_check.zones_I_II.S1_2Bhat
\newcommand{\NUSupersolutionMOneZHypCheckZonesIIISOneTwoBhat}{\ensuremath{0.28585}}
% NUSupersolutionMOneZHypCheckZoneIIIBOnePlusPUOne [cert, upper] U_cert_v35.json :: supersolution.M1z_hyp_check.zone_III.b1_plus_pU1
\newcommand{\NUSupersolutionMOneZHypCheckZoneIIIBOnePlusPUOne}{\ensuremath{0.50786}}
% NUMarginLateMOne [cert, margin] U_cert_v35.json :: margins.late_M1
\newcommand{\NUMarginLateMOne}{\ensuremath{2.353}}
% NUMarginLateMZero [cert, margin] U_cert_v35.json :: margins.late_M0
\newcommand{\NUMarginLateMZero}{\ensuremath{4.892}}
% NUMarginsLateMZeroZ [cert, margin] U_cert_v35.json :: margins.late_M0z
\newcommand{\NUMarginsLateMZeroZ}{\ensuremath{12.99}}
% NUMarginPresentMOne [cert, margin] U_cert_v35.json :: margins.present_M1_min.value
\newcommand{\NUMarginPresentMOne}{\ensuremath{2.770}}
% NUMarginPresentMZero [cert, margin] U_cert_v35.json :: margins.present_M0_min.value
\newcommand{\NUMarginPresentMZero}{\ensuremath{2.225}}
% NUMarginsPresentMZeroZMinValue [cert, margin] U_cert_v35.json :: margins.present_M0z_min.value
\newcommand{\NUMarginsPresentMZeroZMinValue}{\ensuremath{14.37}}
% NULoopGainZeroTotal [cert, upper] U_cert_v35.json :: loop_gain.0.total
\newcommand{\NULoopGainZeroTotal}{\ensuremath{1.23\cdot10^{-7}}}
% NULoopGainOneTotal [cert, upper] U_cert_v35.json :: loop_gain.1.total
\newcommand{\NULoopGainOneTotal}{\ensuremath{0.037169}}
% NULoopGainTwoNeighbour [cert, upper] U_cert_v35.json :: loop_gain.2.neighbour
\newcommand{\NULoopGainTwoNeighbour}{\ensuremath{0.037018}}
% NULoopGainTwoThroughA [cert, upper] U_cert_v35.json :: loop_gain.2.through_A
\newcommand{\NULoopGainTwoThroughA}{\ensuremath{0.0020172}}
% NUfineBhatZero [cert, upper] U_cert_v35.json :: ufine.Bhat.0
\newcommand{\NUfineBhatZero}{\ensuremath{4.8723\cdot10^{-6}}}
% NUfineBhatOne [cert, upper] U_cert_v35.json :: ufine.Bhat.1
\newcommand{\NUfineBhatOne}{\ensuremath{2.6266\cdot10^{-4}}}
% NUfineBhatTwo [cert, upper] U_cert_v35.json :: ufine.Bhat.2
\newcommand{\NUfineBhatTwo}{\ensuremath{0.025842}}
% NUcdeltaU [cert, upper] U_cert_v35.json :: launch.c_dU
\newcommand{\NUcdeltaU}{\ensuremath{1.7187\cdot10^{-8}}}
% NULaunchWmin [cert, lower] U_cert_v35.json :: launch.wmin
\newcommand{\NULaunchWmin}{\ensuremath{0.73492}}
% NULaunchCAD [cert, upper] U_cert_v35.json :: launch.CA_D
\newcommand{\NULaunchCAD}{\ensuremath{2.2429}}
% NULaunchQD [cert, upper] U_cert_v35.json :: launch.qD
\newcommand{\NULaunchQD}{\ensuremath{0.018509}}
% NULaunchDeltaInZero [cert, upper] U_cert_v35.json :: launch.Delta_in.0
\newcommand{\NULaunchDeltaInZero}{\ensuremath{3.131\cdot10^{-12}}}
% NULaunchDeltaInOne [cert, upper] U_cert_v35.json :: launch.Delta_in.1
\newcommand{\NULaunchDeltaInOne}{\ensuremath{3.287\cdot10^{-12}}}
% NULaunchDeltaInTwo [cert, upper] U_cert_v35.json :: launch.Delta_in.2
\newcommand{\NULaunchDeltaInTwo}{\ensuremath{3.451\cdot10^{-12}}}
% NULaunchPhiZero [cert, upper] U_cert_v35.json :: launch.Phi.0
\newcommand{\NULaunchPhiZero}{\ensuremath{6.434\cdot10^{-9}}}
% NULaunchPhiOne [cert, upper] U_cert_v35.json :: launch.Phi.1
\newcommand{\NULaunchPhiOne}{\ensuremath{1.4217\cdot10^{-8}}}
% NULaunchPhiTwo [cert, upper] U_cert_v35.json :: launch.Phi.2
\newcommand{\NULaunchPhiTwo}{\ensuremath{3.8404\cdot10^{-8}}}
% NUGammaZero [cert, upper] U_cert_v35.json :: Gamma.0
\newcommand{\NUGammaZero}{\ensuremath{4.0000358}}
% NUGammaOne [cert, upper] U_cert_v35.json :: Gamma.1
\newcommand{\NUGammaOne}{\ensuremath{16.000322}}
% NUGammaTwo [cert, upper] U_cert_v35.json :: Gamma.2
\newcommand{\NUGammaTwo}{\ensuremath{64.00218}}
% NUKappa [cert, upper] U_cert_v35.json :: kappa
\newcommand{\NUKappa}{\ensuremath{8.0000358}}
% NUKappaTwo [cert, upper] U_cert_v35.json :: kappa2
\newcommand{\NUKappaTwo}{\ensuremath{64.00061}}
% NUGhatZeroZero [cert, upper] U_cert_v35.json :: Ghat.0.0
\newcommand{\NUGhatZeroZero}{\ensuremath{51322.6}}
% NUGhatOneOne [cert, upper] U_cert_v35.json :: Ghat.1.1
\newcommand{\NUGhatOneOne}{\ensuremath{33962.9}}
% NUGhatTwoTwo [cert, upper] U_cert_v35.json :: Ghat.2.2
\newcommand{\NUGhatTwoTwo}{\ensuremath{348341.1}}
% NUGhatOneZero [cert, upper] U_cert_v35.json :: Ghat.1.0
\newcommand{\NUGhatOneZero}{\ensuremath{4.938}}
% NUGhatTwoZero [cert, upper] U_cert_v35.json :: Ghat.2.0
\newcommand{\NUGhatTwoZero}{\ensuremath{27.984}}
% NUGhatTwoOne [cert, upper] U_cert_v35.json :: Ghat.2.1
\newcommand{\NUGhatTwoOne}{\ensuremath{458.53}}
% NUKZeroZero [cert, upper] U_cert_v35.json :: K.0.0
\newcommand{\NUKZeroZero}{\ensuremath{8034.21}}
% NUKZeroOne [cert, upper] U_cert_v35.json :: K.0.1
\newcommand{\NUKZeroOne}{\ensuremath{724.89}}
% NUKOneZero [cert, upper] U_cert_v35.json :: K.1.0
\newcommand{\NUKOneZero}{\ensuremath{0.08152}}
% NUKOneOne [cert, upper] U_cert_v35.json :: K.1.1
\newcommand{\NUKOneOne}{\ensuremath{6336.94}}
% NUKOneTwo [cert, upper] U_cert_v35.json :: K.1.2
\newcommand{\NUKOneTwo}{\ensuremath{545.59}}
% NUKOneThree [cert, upper] U_cert_v35.json :: K.1.3
\newcommand{\NUKOneThree}{\ensuremath{0.58700}}
% NUKTwoZero [cert, upper] U_cert_v35.json :: K.2.0
\newcommand{\NUKTwoZero}{\ensuremath{0.87186}}
% NUKTwoOne [cert, upper] U_cert_v35.json :: K.2.1
\newcommand{\NUKTwoOne}{\ensuremath{528.62}}
% NUKTwoTwo [cert, upper] U_cert_v35.json :: K.2.2
\newcommand{\NUKTwoTwo}{\ensuremath{23425.98}}
% NUKTwoThree [cert, upper] U_cert_v35.json :: K.2.3
\newcommand{\NUKTwoThree}{\ensuremath{555.62}}
% NUFootwOne [cert, upper] U_cert_v35.json :: footw.1
\newcommand{\NUFootwOne}{\ensuremath{2.6361}}
% NUFootwTwo [cert, upper] U_cert_v35.json :: footw.2
\newcommand{\NUFootwTwo}{\ensuremath{2663.3}}
% NUCThreeThree [cert, upper] U_cert_v35.json :: c33
\newcommand{\NUCThreeThree}{\ensuremath{108.02}}
% NUCTwoTwoThree [cert, upper] U_cert_v35.json :: c223
\newcommand{\NUCTwoTwoThree}{\ensuremath{7.3284}}
% NUCEta [cert, upper] U_cert_v35.json :: C_eta
\newcommand{\NUCEta}{\ensuremath{13.376}}
% NULMax [cert, upper] U_cert_v35.json :: L_max
\newcommand{\NULMax}{\ensuremath{3.5352}}
% NURRatio [cert, upper] U_cert_v35.json :: R_ratio
\newcommand{\NURRatio}{\ensuremath{11.593}}
% NUSweepsSigmaInCoverSup [cert, upper] U_cert_v35.json :: sweeps.Sigma_in.cover_sup
\newcommand{\NUSweepsSigmaInCoverSup}{\ensuremath{5.0892}}
% NUSweepsSigmaECoverSup [cert, upper] U_cert_v35.json :: sweeps.Sigma_E.cover_sup
\newcommand{\NUSweepsSigmaECoverSup}{\ensuremath{0.30846}}
% NUSweepsSXCoverSup [cert, upper] U_cert_v35.json :: sweeps.S_X.cover_sup
\newcommand{\NUSweepsSXCoverSup}{\ensuremath{1.5243}}
% NUSweepsSigmaXCoverSup [cert, upper] U_cert_v35.json :: sweeps.Sigma_X.cover_sup
\newcommand{\NUSweepsSigmaXCoverSup}{\ensuremath{3.0081}}
% NUSweepsSigmaRsTwoCoverSup [cert, upper] U_cert_v35.json :: sweeps.Sigma_rs.2.cover_sup
\newcommand{\NUSweepsSigmaRsTwoCoverSup}{\ensuremath{4.3350\cdot10^{-3}}}
% NUDataTermsLateDZero [cert, upper] U_cert_v35.json :: data_terms_late.D.0
\newcommand{\NUDataTermsLateDZero}{\ensuremath{1.7330\cdot10^{-6}}}
% NUDataTermsLateDOne [cert, upper] U_cert_v35.json :: data_terms_late.D.1
\newcommand{\NUDataTermsLateDOne}{\ensuremath{2.5353\cdot10^{-4}}}
% NUDataTermsLateSlopeShareCOneSigmaXOverDkZero [cert, upper] U_cert_v35.json :: data_terms_late.slope_share_C1SigmaX_over_Dk.0
\newcommand{\NUDataTermsLateSlopeShareCOneSigmaXOverDkZero}{\ensuremath{2.7\cdot10^{-3}}}
% NUDataTermsLateSlopeShareCOneSigmaXOverDkOne [cert, upper] U_cert_v35.json :: data_terms_late.slope_share_C1SigmaX_over_Dk.1
\newcommand{\NUDataTermsLateSlopeShareCOneSigmaXOverDkOne}{\ensuremath{6.6\cdot10^{-4}}}
% NURLtTwoMarginZero [cert, margin] U_cert_v35.json :: R_lt2_margin.0
\newcommand{\NURLtTwoMarginZero}{\ensuremath{4.432}}
% NURLtTwoMarginOne [cert, margin] U_cert_v35.json :: R_lt2_margin.1
\newcommand{\NURLtTwoMarginOne}{\ensuremath{6.697}}
% NUThetaTwoEight [cert, upper] U_cert_v35.json :: theta28
\newcommand{\NUThetaTwoEight}{\ensuremath{0.37011}}
% NUDFiveSix [cert, upper] U_cert_v35.json :: D56
\newcommand{\NUDFiveSix}{\ensuremath{1.0368}}
% NUDTwoEight [cert, upper] U_cert_v35.json :: D28
\newcommand{\NUDTwoEight}{\ensuremath{1.5876}}
% NUSThreeThreeFootOverBhatTwo [cert, upper] U_cert_v35.json :: S33_foot_over_Bhat2
\newcommand{\NUSThreeThreeFootOverBhatTwo}{\ensuremath{1.69\cdot10^{-3}}}
% NSAHJWKappaJ [cert, upper] A_interface.json :: HJW.kappa_J
\newcommand{\NSAHJWKappaJ}{\ensuremath{0.042351}}
% NSAGpZero [cert, lower] A_interface.json :: res.tight.gp.0
\newcommand{\NSAGpZero}{\ensuremath{0.24292}}
% NSAGpOne [cert, lower] A_interface.json :: res.tight.gp.1
\newcommand{\NSAGpOne}{\ensuremath{0.24106}}
% NSAGpTwo [cert, lower] A_interface.json :: res.tight.gp.2
\newcommand{\NSAGpTwo}{\ensuremath{0.11961}}
% NSAAminC [cert, lower] A_interface.json :: res.tight.amin_c
\newcommand{\NSAAminC}{\ensuremath{0.97623}}
% NSARhoJ [cert, upper] A_interface.json :: res.tight.rhoJ
\newcommand{\NSARhoJ}{\ensuremath{5.4995\cdot10^{-7}}}
% NSAKappaOTwoTwo [cert, upper] A_interface.json :: res.tight.kappa_o.2.2
\newcommand{\NSAKappaOTwoTwo}{\ensuremath{1.2416\cdot10^{-5}}}
% NSADsrcZero [cert, upper] A_interface.json :: res.tight.Dsrc.0
\newcommand{\NSADsrcZero}{\ensuremath{1.8092\cdot10^{-7}}}
% NSADsrcOne [cert, upper] A_interface.json :: res.tight.Dsrc.1
\newcommand{\NSADsrcOne}{\ensuremath{9.8276\cdot10^{-5}}}
% NSABZeroDWindowSlopeZero [cert, upper] A_interface.json :: res.tight.B0.D_window_slope.0
\newcommand{\NSABZeroDWindowSlopeZero}{\ensuremath{8.7399\cdot10^{-10}}}
% NSABZeroDWindowSlopeOne [cert, upper] A_interface.json :: res.tight.B0.D_window_slope.1
\newcommand{\NSABZeroDWindowSlopeOne}{\ensuremath{5.2599\cdot10^{-8}}}
% NSAFTwoZero [cert, upper] A_interface.json :: res.tight.F2.0
\newcommand{\NSAFTwoZero}{\ensuremath{5.6506\cdot10^{-13}}}
% NSAFTwoOne [cert, upper] A_interface.json :: res.tight.F2.1
\newcommand{\NSAFTwoOne}{\ensuremath{5.9320\cdot10^{-13}}}
% NSACoupOne [cert, upper] A_interface.json :: res.tight.coup.1
\newcommand{\NSACoupOne}{\ensuremath{1.1568\cdot10^{-12}}}
% NSACOneZeroZero [cert, upper] A_interface.json :: res.tight.C1_0.0
\newcommand{\NSACOneZeroZero}{\ensuremath{7.2060\cdot10^{-10}}}
% NSACOneZeroOne [cert, upper] A_interface.json :: res.tight.C1_0.1
\newcommand{\NSACOneZeroOne}{\ensuremath{4.3143\cdot10^{-8}}}
% NSABZeroSupCoreZero [cert, upper] A_interface.json :: res.tight.B0.sup_core.0
\newcommand{\NSABZeroSupCoreZero}{\ensuremath{4.23\cdot10^{-5}}}
% NSACoreFluxMarginTwoFLUXZeroOverKZero [cert, margin] A_interface.json :: res.tight.core.flux_margin.2.FLUX0_over_K0
\newcommand{\NSACoreFluxMarginTwoFLUXZeroOverKZero}{\ensuremath{3.396}}
% NSACoreFluxMarginTwoFLUXTwoOverKTwo [cert, margin] A_interface.json :: res.tight.core.flux_margin.2.FLUX2_over_K2
\newcommand{\NSACoreFluxMarginTwoFLUXTwoOverKTwo}{\ensuremath{6.012}}
% NSAMarginFluxZero [cert, margin] asm_design_N2ref4.json :: res.tight.flux.beta.2.margin_FLUX0
\newcommand{\NSAMarginFluxZero}{\ensuremath{3.365}}
% NSAMarginFluxTwo [cert, margin] asm_design_N2ref4.json :: res.tight.flux.beta.2.margin_FLUX2
\newcommand{\NSAMarginFluxTwo}{\ensuremath{5.955}}
% NSAfourSZero [cert, upper] asm_design_N2ref4.json :: res.tight.S0
\newcommand{\NSAfourSZero}{\ensuremath{7.1878\cdot10^{-3}}}
% NSAfourSOne [cert, upper] asm_design_N2ref4.json :: res.tight.S1
\newcommand{\NSAfourSOne}{\ensuremath{0.10430}}
% NSAMarginCoreSZeroV [cert, margin] asm_design_N2ref4.json :: res.tight.margins.S0v_vs_0.0153
\newcommand{\NSAMarginCoreSZeroV}{\ensuremath{2.116}}
% NSAMarginCoreSOneV [cert, margin] asm_design_N2ref4.json :: res.tight.margins.S1v_vs_0.6
\newcommand{\NSAMarginCoreSOneV}{\ensuremath{5.748}}
% NSAMarginCoreMOne [cert, margin] asm_design_N2ref4.json :: res.tight.margins.M1core
\newcommand{\NSAMarginCoreMOne}{\ensuremath{4.409}}
% NSAMarginCoreMZero [cert, margin] asm_design_N2ref4.json :: res.tight.margins.M0core
\newcommand{\NSAMarginCoreMZero}{\ensuremath{5.450}}
% NSAMarginGhat [cert, margin] asm_design_N2ref4.json :: res.tight.margins.Ghat22
\newcommand{\NSAMarginGhat}{\ensuremath{1.917}}
% NSAMarginRampMOne [cert, margin] asm_design_N2ref4.json :: res.tight.margins.ramp M1
\newcommand{\NSAMarginRampMOne}{\ensuremath{1.184}}
% NSAfourInfoZeroDWindowConst [cert, upper] asm_design_N2ref4.json :: res.tight.info.0.D_window_const
\newcommand{\NSAfourInfoZeroDWindowConst}{\ensuremath{1.7732\cdot10^{-7}}}
% NSAfourInfoOneDWindowConst [cert, upper] asm_design_N2ref4.json :: res.tight.info.1.D_window_const
\newcommand{\NSAfourInfoOneDWindowConst}{\ensuremath{9.8058\cdot10^{-5}}}
% NZDeltaRZero [cert, upper] zone3_interface.json :: DeltaR.0
\newcommand{\NZDeltaRZero}{\ensuremath{4.4986\cdot10^{-8}}}
% NZDeltaROne [cert, upper] zone3_interface.json :: DeltaR.1
\newcommand{\NZDeltaROne}{\ensuremath{6.7972\cdot10^{-6}}}
% NZDeltaRTwo [cert, upper] zone3_interface.json :: DeltaR.2
\newcommand{\NZDeltaRTwo}{\ensuremath{0.011480}}
% NZDeltaRBZero [cert, upper] zone3_interface.json :: DeltaR_b.0
\newcommand{\NZDeltaRBZero}{\ensuremath{2.66\cdot10^{-12}}}
% NZDeltaRBOne [cert, upper] zone3_interface.json :: DeltaR_b.1
\newcommand{\NZDeltaRBOne}{\ensuremath{3.26\cdot10^{-9}}}
% NZDeltaRBTwo [cert, upper] zone3_interface.json :: DeltaR_b.2
\newcommand{\NZDeltaRBTwo}{\ensuremath{4.35\cdot10^{-8}}}
% NZDeltaRDetailNSqZero [cert, upper] zone3_interface.json :: DeltaR_detail.N_sq.0
\newcommand{\NZDeltaRDetailNSqZero}{\ensuremath{8.28\cdot10^{-5}}}
% NZDeltaRDetailNSqOne [cert, upper] zone3_interface.json :: DeltaR_detail.N_sq.1
\newcommand{\NZDeltaRDetailNSqOne}{\ensuremath{0.10149}}
% NZDeltaRDetailNSqTwo [cert, upper] zone3_interface.json :: DeltaR_detail.N_sq.2
\newcommand{\NZDeltaRDetailNSqTwo}{\ensuremath{1.3567}}
% NZDeltaRDetailNSqThree [cert, upper] zone3_interface.json :: DeltaR_detail.N_sq.3
\newcommand{\NZDeltaRDetailNSqThree}{\ensuremath{25350}}
% NZDeltaRDetailPartsTwoPairsThree [cert, upper] zone3_interface.json :: DeltaR_detail.parts.2.pairs.3
\newcommand{\NZDeltaRDetailPartsTwoPairsThree}{\ensuremath{0.011298}}
% NZDeltaRDetailPartsTwoLFourFoot [cert, upper] zone3_interface.json :: DeltaR_detail.parts.2.L4_foot
\newcommand{\NZDeltaRDetailPartsTwoLFourFoot}{\ensuremath{1.6\cdot10^{-4}}}
% NZDeltaRDetailPartsTwoCubic [cert, upper] zone3_interface.json :: DeltaR_detail.parts.2.cubic
\newcommand{\NZDeltaRDetailPartsTwoCubic}{\ensuremath{9.3\cdot10^{-6}}}
% NZDeltaRDetailPartsTwoSThreeThreeFoot [cert, upper] zone3_interface.json :: DeltaR_detail.parts.2.S33_foot
\newcommand{\NZDeltaRDetailPartsTwoSThreeThreeFoot}{\ensuremath{3.7\cdot10^{-7}}}
% NZNIIIZero [cert, upper] zone3_interface.json :: n_III.0
\newcommand{\NZNIIIZero}{\ensuremath{0.071445}}
% NZNIIIOne [cert, upper] zone3_interface.json :: n_III.1
\newcommand{\NZNIIIOne}{\ensuremath{0.32707}}
% NZNIIITwo [cert, upper] zone3_interface.json :: n_III.2
\newcommand{\NZNIIITwo}{\ensuremath{3.0445}}
% NZTauIII [cert, upper] zone3_interface.json :: tau_III
\newcommand{\NZTauIII}{\ensuremath{4.0708}}
% NZoneIIIBoxZero [cert, upper] zone3_interface.json :: box_III_380.b0
\newcommand{\NZoneIIIBoxZero}{\ensuremath{0.10032}}
% NZoneIIIBoxOne [cert, upper] zone3_interface.json :: box_III_380.b1
\newcommand{\NZoneIIIBoxOne}{\ensuremath{0.41974}}
% NZoneIIIBoxZ [cert, upper] zone3_interface.json :: box_III_380.bz
\newcommand{\NZoneIIIBoxZ}{\ensuremath{0.014291}}
% NZBoxIIIThreeEightZeroMarginsBZero [cert, margin] zone3_interface.json :: box_III_380.margins.b0
\newcommand{\NZBoxIIIThreeEightZeroMarginsBZero}{\ensuremath{1.495}}
% NZoneIIIMarginThinnest [cert, margin] zone3_interface.json :: box_III_380.margins.b1
\newcommand{\NZoneIIIMarginThinnest}{\ensuremath{1.238}}
% NZBoxIIIThreeEightZeroMarginsBz [cert, margin] zone3_interface.json :: box_III_380.margins.bz
\newcommand{\NZBoxIIIThreeEightZeroMarginsBz}{\ensuremath{11.19}}
% NZBoxIIISevenZeroZeroBZero [cert, upper] zone3_interface.json :: box_III_700.b0
\newcommand{\NZBoxIIISevenZeroZeroBZero}{\ensuremath{0.22757}}
% NZBoxIIISevenZeroZeroBOne [cert, upper] zone3_interface.json :: box_III_700.b1
\newcommand{\NZBoxIIISevenZeroZeroBOne}{\ensuremath{0.54234}}
% NZUfineKTwoAtBhatRef [cert, upper] zone3_interface.json :: ufine.K2_at_Bhat_ref
\newcommand{\NZUfineKTwoAtBhatRef}{\ensuremath{0.046956}}
% NZUfineKThreeStripAtBhatRef [cert, upper] zone3_interface.json :: ufine.K3_strip_at_Bhat_ref
\newcommand{\NZUfineKThreeStripAtBhatRef}{\ensuremath{0.054140}}
% NRecoveryMinMargin [cert, margin] zone3_interface.json :: cone_family.H6.at_Bhat_ref.min_margin
\newcommand{\NRecoveryMinMargin}{\ensuremath{1.860}}
% NZConeFamilyHSixAtBhatRefConeSlopeMargin [cert, margin] zone3_interface.json :: cone_family.H6.at_Bhat_ref.cone.slope.margin
\newcommand{\NZConeFamilyHSixAtBhatRefConeSlopeMargin}{\ensuremath{2.46}}
% NZConeFamilyHSixAtBhatRefConeKTwoMargin [cert, margin] zone3_interface.json :: cone_family.H6.at_Bhat_ref.cone.K2.margin
\newcommand{\NZConeFamilyHSixAtBhatRefConeKTwoMargin}{\ensuremath{2.54}}
% NZConeFamilyHSixAtBhatRefConeKThreeMargin [cert, margin] zone3_interface.json :: cone_family.H6.at_Bhat_ref.cone.K3.margin
\newcommand{\NZConeFamilyHSixAtBhatRefConeKThreeMargin}{\ensuremath{2.51}}
% NConeTstarMargin [cert, margin] zone3_interface.json :: acceptance.tstar_margin
\newcommand{\NConeTstarMargin}{\ensuremath{2.30}}
% NUoneTFiveYaddLevelOneTwo [cert, upper] u1_ghat.json :: t5_Yadd.level1.2
\newcommand{\NUoneTFiveYaddLevelOneTwo}{\ensuremath{0.0044}}
% NUoneTFiveYaddLevelTwoTwo [cert, upper] u1_ghat.json :: t5_Yadd.level2.2
\newcommand{\NUoneTFiveYaddLevelTwoTwo}{\ensuremath{0.1081}}
% NUoneTFiveYaddLevelTwoThree [cert, upper] u1_ghat.json :: t5_Yadd.level2.3
\newcommand{\NUoneTFiveYaddLevelTwoThree}{\ensuremath{0.374}}
% NUoneWeightThetaGammaZero [cert, upper] u1_ghat.json :: weight.Theta.Gamma.0
\newcommand{\NUoneWeightThetaGammaZero}{\ensuremath{4.00000000000000000000054}}
% NRKZeroZero [cert, upper] ramp_U.json :: K.0.0
\newcommand{\NRKZeroZero}{\ensuremath{13436.03}}
% NRKOneOne [cert, upper] ramp_U.json :: K.1.1
\newcommand{\NRKOneOne}{\ensuremath{10526.19}}
% NRKTwoTwo [cert, upper] ramp_U.json :: K.2.2
\newcommand{\NRKTwoTwo}{\ensuremath{39482.30}}
% NUoneCasesPrimaryYNextZero [cert, upper] u1_ghat.json :: cases.primary.Y_next.0
\newcommand{\NUoneCasesPrimaryYNextZero}{\ensuremath{457.092}}
% NUoneCasesPrimaryYNextOne [cert, upper] u1_ghat.json :: cases.primary.Y_next.1
\newcommand{\NUoneCasesPrimaryYNextOne}{\ensuremath{356.940}}
% NUoneCasesPrimaryYNextTwo [cert, upper] u1_ghat.json :: cases.primary.Y_next.2
\newcommand{\NUoneCasesPrimaryYNextTwo}{\ensuremath{1320.220}}
% NSfourResultsMNineHTwoTightenedFUOneRowsZeroBound [cert, upper] FINAL_manifest_v33.json :: results.M9_H2_tightened_fU1.rows.0.bound
\newcommand{\NSfourResultsMNineHTwoTightenedFUOneRowsZeroBound}{\ensuremath{3.658\cdot10^{-10}}}
% NSfourResultsMNineHTwoTightenedFUOneRowsZeroMargin [cert, margin] FINAL_manifest_v33.json :: results.M9_H2_tightened_fU1.rows.0.margin
\newcommand{\NSfourResultsMNineHTwoTightenedFUOneRowsZeroMargin}{\ensuremath{27.34}}
% NSfourResultsMNineHTwoTightenedFUOneRowsOneBound [cert, upper] FINAL_manifest_v33.json :: results.M9_H2_tightened_fU1.rows.1.bound
\newcommand{\NSfourResultsMNineHTwoTightenedFUOneRowsOneBound}{\ensuremath{1.777\cdot10^{-9}}}
% NSfourResultsMNineHTwoTightenedFUOneRowsOneMargin [cert, margin] FINAL_manifest_v33.json :: results.M9_H2_tightened_fU1.rows.1.margin
\newcommand{\NSfourResultsMNineHTwoTightenedFUOneRowsOneMargin}{\ensuremath{562.8}}
% NSfourResultsMNineHTwoTightenedFUOneRowsTwoBound [cert, upper] FINAL_manifest_v33.json :: results.M9_H2_tightened_fU1.rows.2.bound
\newcommand{\NSfourResultsMNineHTwoTightenedFUOneRowsTwoBound}{\ensuremath{1.665\cdot10^{-6}}}
% NSfourResultsMNineHTwoTightenedFUOneRowsTwoMargin [cert, margin] FINAL_manifest_v33.json :: results.M9_H2_tightened_fU1.rows.2.margin
\newcommand{\NSfourResultsMNineHTwoTightenedFUOneRowsTwoMargin}{\ensuremath{600.9}}
% NHOneAlphaThr [hand, exact] hand :: 
\newcommand{\NHOneAlphaThr}{\ensuremath{1.5\cdot10^{-4}}}
% NHOneSqrtLamThr [hand, exact] hand :: 
\newcommand{\NHOneSqrtLamThr}{\ensuremath{5\cdot10^{-8}}}
% NYbarStep [hand, exact] hand :: 
\newcommand{\NYbarStep}{\ensuremath{1\cdot10^{-12}}}
% NHTwoThrZero [hand, exact] hand :: 
\newcommand{\NHTwoThrZero}{\ensuremath{1\cdot10^{-8}}}
% NHTwoThrOne [hand, exact] hand :: 
\newcommand{\NHTwoThrOne}{\ensuremath{1\cdot10^{-6}}}
% NHTwoThrTwo [hand, exact] hand :: 
\newcommand{\NHTwoThrTwo}{\ensuremath{1\cdot10^{-3}}}
% NHThreeThrZero [hand, exact] hand :: 
\newcommand{\NHThreeThrZero}{\ensuremath{1\cdot10^{-6}}}
% NHThreeThrOne [hand, exact] hand :: 
\newcommand{\NHThreeThrOne}{\ensuremath{5\cdot10^{-6}}}
% NHThreeThrTwo [hand, exact] hand :: 
\newcommand{\NHThreeThrTwo}{\ensuremath{1\cdot10^{-5}}}
% NHFourThrZero [hand, exact] hand :: 
\newcommand{\NHFourThrZero}{\ensuremath{2\cdot10^{-7}}}
% NHFourThrOne [hand, exact] hand :: 
\newcommand{\NHFourThrOne}{\ensuremath{1\cdot10^{-5}}}
% NHFourThrTwo [hand, exact] hand :: 
\newcommand{\NHFourThrTwo}{\ensuremath{1\cdot10^{-5}}}
% NHFourThrThree [hand, exact] hand :: 
\newcommand{\NHFourThrThree}{\ensuremath{1\cdot10^{-5}}}
% NHFourKfourThr [hand, exact] hand :: 
\newcommand{\NHFourKfourThr}{\ensuremath{0.06}}
% NHSixConeSlopeThr [hand, exact] hand :: 
\newcommand{\NHSixConeSlopeThr}{\ensuremath{0.08}}
% NHSixConeKtwoThr [hand, exact] hand :: 
\newcommand{\NHSixConeKtwoThr}{\ensuremath{0.03}}
% NHSixConeNthreeThr [hand, exact] hand :: 
\newcommand{\NHSixConeNthreeThr}{\ensuremath{0.004}}
% NHSixConeJetThr [hand, exact] hand :: 
\newcommand{\NHSixConeJetThr}{\ensuremath{0.12}}
% NHSixConeHFiveThr [hand, exact] hand :: 
\newcommand{\NHSixConeHFiveThr}{\ensuremath{0.0144}}
% NHSixFarHFiveThr [hand, exact] hand :: 
\newcommand{\NHSixFarHFiveThr}{\ensuremath{0.0225}}
% NLemmaOConeM [hand, exact] hand :: 
\newcommand{\NLemmaOConeM}{\ensuremath{0.12}}
% NLemmaOConeK [hand, exact] hand :: 
\newcommand{\NLemmaOConeK}{\ensuremath{0.18}}
% NLemmaOFarM [hand, exact] hand :: 
\newcommand{\NLemmaOFarM}{\ensuremath{0.15}}
% NLemmaOFarK [hand, exact] hand :: 
\newcommand{\NLemmaOFarK}{\ensuremath{0.3}}
% NCoercivityConst [hand, exact] hand :: 
\newcommand{\NCoercivityConst}{\ensuremath{0.046875}}
% NMOnez [hand, exact] hand :: 
\newcommand{\NMOnez}{\ensuremath{0.534}}
% NBetaA [hand, exact] hand :: 
\newcommand{\NBetaA}{\ensuremath{2}}
% NTOneLower [cert, lower] t1clock.json :: t1.0
\newcommand{\NTOneLower}{\ensuremath{4.5978574775}}
% NTOneUpper [cert, upper] t1clock.json :: t1.1
\newcommand{\NTOneUpper}{\ensuremath{4.5978586144}}
% NLambdaOneLower [cert, lower] t1clock.json :: lnlam.0
\newcommand{\NLambdaOneLower}{\ensuremath{6.49\cdot10^{-5}}}
% NLambdaOneUpper [cert, upper] t1clock.json :: lnlam.1
\newcommand{\NLambdaOneUpper}{\ensuremath{6.75\cdot10^{-5}}}
% NTMinusTOneUpper [cert, upper] t1clock.json :: lnlam.1; 1.5e-4.sup_ab
\newcommand{\NTMinusTOneUpper}{\ensuremath{2.1\cdot10^{-17}}}
% NAbarUpper [cert, upper] final_alpha_points_5e-8.json :: 1.5e-4.sup_ab
\newcommand{\NAbarUpper}{\ensuremath{2.19103\cdot10^{-7}}}
% NBetabarUpper [cert, upper] final_alpha_points_5e-8.json :: 1.5e-4.sup_beta
\newcommand{\NBetabarUpper}{\ensuremath{3.33811\cdot10^{-9}}}
% NBaseKmaxLower [cert, lower] t1clock.json :: base_kmax.0
\newcommand{\NBaseKmaxLower}{\ensuremath{2.279\cdot10^{4}}}
% NBaseKmaxUpper [cert, upper] t1clock.json :: base_kmax.1
\newcommand{\NBaseKmaxUpper}{\ensuremath{2.376\cdot10^{4}}}
% NHOneAlphaBase [cert, upper] summary.json :: alpha_abs_max
\newcommand{\NHOneAlphaBase}{\ensuremath{1.1739\cdot10^{-4}}}
% NHOneAlphaErr [cert, upper] t9_h1_clam_s180.json :: C_alpha; inputs.YSTAR
\newcommand{\NHOneAlphaErr}{\ensuremath{7.21\cdot10^{-13}}}
% NHOneAlphaMargin [cert, margin] summary.json :: alpha_abs_max; C_alpha; inputs.YSTAR
\newcommand{\NHOneAlphaMargin}{\ensuremath{1.27}}
% NHOneSqrtLamBase [cert, upper] t9_h1_clam_s180.json :: sqrtLam_base_Wp
\newcommand{\NHOneSqrtLamBase}{\ensuremath{2.282\cdot10^{-8}}}
% NHOneSqrtLamTotal [cert, upper] t9_h1_clam_s180.json :: sqrtLam_total_Wp
\newcommand{\NHOneSqrtLamTotal}{\ensuremath{2.29\cdot10^{-8}}}
% NHOneSqrtLamMargin [cert, margin] t9_h1_clam_s180.json :: sqrtLam_total_Wp
\newcommand{\NHOneSqrtLamMargin}{\ensuremath{2.18}}
% NHOneYbarLimit [cert, lower] t9_h1_clam_s180.json :: Ybar_limit_H1
\newcommand{\NHOneYbarLimit}{\ensuremath{1.12\cdot10^{-9}}}
% NSwitchYUpper [cert, upper] switch_s2.json :: modulated.Y_raw
\newcommand{\NSwitchYUpper}{\ensuremath{6.21\cdot10^{-13}}}
% NSwitchYMargin [cert, margin] switch_s2.json :: modulated.Y_raw
\newcommand{\NSwitchYMargin}{\ensuremath{1.6}}
% NMinCI [cert, lower] constants.json :: meta.summary.min_c_I
\newcommand{\NMinCI}{\ensuremath{0.05283}}
% NNumCertificates [cert, exact] constants.json :: meta.summary.n_cert
\newcommand{\NNumCertificates}{\ensuremath{26}}
% NMinRhoC [cert, lower] constants.json :: meta.closure.rho_c_min
\newcommand{\NMinRhoC}{\ensuremath{0.04711}}
% NHTwoBaseZero [cert, upper] FINAL_manifest_v33.json :: results.M9_H2_tightened.rows.0.base
\newcommand{\NHTwoBaseZero}{\ensuremath{8.011\cdot10^{-11}}}
% NHTwoBaseOne [cert, upper] FINAL_manifest_v33.json :: results.M9_H2_tightened.rows.1.base
\newcommand{\NHTwoBaseOne}{\ensuremath{1.336\cdot10^{-10}}}
% NHTwoBaseTwo [cert, upper] FINAL_manifest_v33.json :: results.M9_H2_tightened.rows.2.base
\newcommand{\NHTwoBaseTwo}{\ensuremath{2.936\cdot10^{-10}}}
% NHTwoErrWZero [cert, upper] FINAL_manifest_v33.json :: results.M9_H2_tightened.rows.0.with_error
\newcommand{\NHTwoErrWZero}{\ensuremath{8.014\cdot10^{-11}}}
% NHTwoMarginWZero [cert, margin] FINAL_manifest_v33.json :: results.M9_H2_tightened.rows.0.margin
\newcommand{\NHTwoMarginWZero}{\ensuremath{124.78}}
% NHTwoErrWOne [cert, upper] FINAL_manifest_v33.json :: results.M9_H2_tightened.rows.1.with_error
\newcommand{\NHTwoErrWOne}{\ensuremath{4.706\cdot10^{-10}}}
% NHTwoMarginWOne [cert, margin] FINAL_manifest_v33.json :: results.M9_H2_tightened.rows.1.margin
\newcommand{\NHTwoMarginWOne}{\ensuremath{2124.9}}
% NHTwoErrWTwo [cert, upper] FINAL_manifest_v33.json :: results.M9_H2_tightened.rows.2.with_error
\newcommand{\NHTwoErrWTwo}{\ensuremath{1.559\cdot10^{-6}}}
% NHTwoMarginWTwo [cert, margin] FINAL_manifest_v33.json :: results.M9_H2_tightened.rows.2.margin
\newcommand{\NHTwoMarginWTwo}{\ensuremath{641.53}}
% NHThreeBaseZero [cert, upper] FINAL_manifest_v33.json :: results.M9_H2f.rows.0.base
\newcommand{\NHThreeBaseZero}{\ensuremath{3.127\cdot10^{-7}}}
% NHThreeBaseOne [cert, upper] FINAL_manifest_v33.json :: results.M9_H2f.rows.1.base
\newcommand{\NHThreeBaseOne}{\ensuremath{1.004\cdot10^{-6}}}
% NHThreeBaseTwo [cert, upper] FINAL_manifest_v33.json :: results.M9_H2f.rows.2.base
\newcommand{\NHThreeBaseTwo}{\ensuremath{2.431\cdot10^{-6}}}
% NHThreeErrZero [cert, upper] FINAL_manifest_v33.json :: results.M9_H2f.rows.0.with_error
\newcommand{\NHThreeErrZero}{\ensuremath{3.1363\cdot10^{-7}}}
% NHThreeMarginZero [cert, margin] FINAL_manifest_v33.json :: results.M9_H2f.rows.0.margin
\newcommand{\NHThreeMarginZero}{\ensuremath{3.1885}}
% NHThreeErrOne [cert, upper] FINAL_manifest_v33.json :: results.M9_H2f.rows.1.with_error
\newcommand{\NHThreeErrOne}{\ensuremath{1.0034\cdot10^{-6}}}
% NHThreeMarginOne [cert, margin] FINAL_manifest_v33.json :: results.M9_H2f.rows.1.margin
\newcommand{\NHThreeMarginOne}{\ensuremath{4.9834}}
% NHThreeErrTwo [cert, upper] FINAL_manifest_v33.json :: results.M9_H2f.rows.2.with_error
\newcommand{\NHThreeErrTwo}{\ensuremath{2.4351\cdot10^{-6}}}
% NHThreeMarginTwo [cert, margin] FINAL_manifest_v33.json :: results.M9_H2f.rows.2.margin
\newcommand{\NHThreeMarginTwo}{\ensuremath{4.1067}}
% NHFourBaseZero [cert, upper] FINAL_manifest_v33.json :: results.M9_base_sums.combined.windows_max.0
\newcommand{\NHFourBaseZero}{\ensuremath{3.48\cdot10^{-8}}}
% NHFourErrZero [cert, upper] FINAL_manifest_v33.json :: results.M9_H3_H2v32.rows.win.0.with_error_arb
\newcommand{\NHFourErrZero}{\ensuremath{3.5087\cdot10^{-8}}}
% NHFourMarginZero [cert, margin] FINAL_manifest_v33.json :: results.M9_H3_H2v32.rows.win.0.margin_arb
\newcommand{\NHFourMarginZero}{\ensuremath{5.700}}
% NHFourBaseOne [cert, upper] FINAL_manifest_v33.json :: results.M9_base_sums.combined.windows_max.1
\newcommand{\NHFourBaseOne}{\ensuremath{1.27\cdot10^{-7}}}
% NHFourErrOne [cert, upper] FINAL_manifest_v33.json :: results.M9_H3_H2v32.rows.win.1.with_error_arb
\newcommand{\NHFourErrOne}{\ensuremath{1.2616\cdot10^{-7}}}
% NHFourMarginOne [cert, margin] FINAL_manifest_v33.json :: results.M9_H3_H2v32.rows.win.1.margin_arb
\newcommand{\NHFourMarginOne}{\ensuremath{79.26}}
% NHFourBaseTwo [cert, upper] FINAL_manifest_v33.json :: results.M9_base_sums.combined.windows_max.2
\newcommand{\NHFourBaseTwo}{\ensuremath{3.76\cdot10^{-7}}}
% NHFourErrTwo [cert, upper] FINAL_manifest_v33.json :: results.M9_H3_H2v32.rows.win.2.with_error_arb
\newcommand{\NHFourErrTwo}{\ensuremath{3.7724\cdot10^{-7}}}
% NHFourMarginTwo [cert, margin] FINAL_manifest_v33.json :: results.M9_H3_H2v32.rows.win.2.margin_arb
\newcommand{\NHFourMarginTwo}{\ensuremath{26.50}}
% NHFourBaseThree [cert, upper] FINAL_manifest_v33.json :: results.M9_base_sums.combined.windows_max.3
\newcommand{\NHFourBaseThree}{\ensuremath{9.71\cdot10^{-7}}}
% NHFourErrThree [cert, upper] FINAL_manifest_v33.json :: results.M9_H3_H2v32.rows.win.3.with_error_arb
\newcommand{\NHFourErrThree}{\ensuremath{1.2822\cdot10^{-6}}}
% NHFourMarginThree [cert, margin] FINAL_manifest_v33.json :: results.M9_H3_H2v32.rows.win.3.margin_arb
\newcommand{\NHFourMarginThree}{\ensuremath{7.799}}
% NHFourYbarLimit [cert, lower] FINAL_manifest_v33.json :: results.M9_H3_H2v32.binding_Ybar_limit
\newcommand{\NHFourYbarLimit}{\ensuremath{1.575\cdot10^{-11}}}
% NHFourYbarMargin [cert, margin] FINAL_manifest_v33.json :: results.M9_H3_H2v32.binding_Ybar_limit
\newcommand{\NHFourYbarMargin}{\ensuremath{15.75}}
% NHFourKfourBase [cert, upper] amp_A1_w180.json :: inputs.BASE_T1.K4
\newcommand{\NHFourKfourBase}{\ensuremath{3.1633\cdot10^{-4}}}
% NHSixConeSlope [cert, upper] FINAL_manifest_v33.json :: results.cone_fixedL.checks.slope.value
\newcommand{\NHSixConeSlope}{\ensuremath{0.03029}}
% NHSixConeSlopeMargin [cert, margin] FINAL_manifest_v33.json :: results.cone_fixedL.checks.slope.margin
\newcommand{\NHSixConeSlopeMargin}{\ensuremath{2.641}}
% NHSixConeKtwo [cert, upper] FINAL_manifest_v33.json :: results.cone_fixedL.checks.K2.value
\newcommand{\NHSixConeKtwo}{\ensuremath{0.009876}}
% NHSixConeKtwoMargin [cert, margin] FINAL_manifest_v33.json :: results.cone_fixedL.checks.K2.margin
\newcommand{\NHSixConeKtwoMargin}{\ensuremath{3.037}}
% NHSixConeNthree [cert, upper] FINAL_manifest_v33.json :: results.cone_fixedL.checks.N3.value
\newcommand{\NHSixConeNthree}{\ensuremath{1.461\cdot10^{-4}}}
% NHSixConeNthreeMargin [cert, margin] FINAL_manifest_v33.json :: results.cone_fixedL.checks.N3.margin
\newcommand{\NHSixConeNthreeMargin}{\ensuremath{27.38}}
% NHSixConeHFive [cert, upper] FINAL_manifest_v33.json :: results.H4_fixedL.total
\newcommand{\NHSixConeHFive}{\ensuremath{1.2865\cdot10^{-5}}}
% NHSixConeHFiveMargin [cert, margin] FINAL_manifest_v33.json :: results.H4_fixedL.margin
\newcommand{\NHSixConeHFiveMargin}{\ensuremath{1119}}
% NHSixConeTstarNeeded [cert, upper] FINAL_manifest_v33.json :: results.H7_fixedL.ratio_upper
\newcommand{\NHSixConeTstarNeeded}{\ensuremath{1.6000015\cdot10^{-7}}}
% NHSixConeTstarLower [cert, lower] FINAL_manifest_v33.json :: results.H7_fixedL.t_star_lower
\newcommand{\NHSixConeTstarLower}{\ensuremath{3.68334\cdot10^{-7}}}
% NHSixConeTstarMargin [cert, margin] FINAL_manifest_v33.json :: results.H7_fixedL.margin_tstar_over_ratio
\newcommand{\NHSixConeTstarMargin}{\ensuremath{2.302}}
% NHSixConeHFiveRatioPeak [cert, upper] h4loc_A1_w180.json :: windows.constants/amp_A1_w180.json.worst.ratio
\newcommand{\NHSixConeHFiveRatioPeak}{\ensuremath{0.2393}}
% NHSixConeHOneSup [cert, upper] FINAL_manifest_v33.json :: results.H1_fixedL.sup_h0_max
\newcommand{\NHSixConeHOneSup}{\ensuremath{0.2163}}
% NHSixConeHOneH [cert, upper] FINAL_manifest_v33.json :: results.H1_fixedL.H_max
\newcommand{\NHSixConeHOneH}{\ensuremath{0.2185}}
% NHSixFarHFiveRatio [cert, upper] h4loc_v32.json :: windows.constants/amp_A1_w180.json.cylinders.18.ratio
\newcommand{\NHSixFarHFiveRatio}{\ensuremath{0.1761}}
% NHSixFarHFiveMarginOld [cert, margin] h4loc_v32.json :: windows.constants/amp_A1_w180.json.cylinders.18.ratio
\newcommand{\NHSixFarHFiveMarginOld}{\ensuremath{5.68}}
% NHSixFarHFiveMargin [cert, margin] FINAL_manifest_v33.json :: results.M9_H4_far.min_margin
\newcommand{\NHSixFarHFiveMargin}{\ensuremath{15.58}}
% NHSixConeHFiveMarginOld [cert, margin] h4loc_A1_w180.json :: windows.constants/amp_A1_w180.json.worst.ratio
\newcommand{\NHSixConeHFiveMarginOld}{\ensuremath{4.18}}
% NHSixFarKtwo [cert, upper] t9_far_v32.json :: h2
\newcommand{\NHSixFarKtwo}{\ensuremath{1.797\cdot10^{-4}}}
% NHSixFarKtwoThr [hand, exact] t9_far_v32.json :: thr.h2
\newcommand{\NHSixFarKtwoThr}{\ensuremath{3.53\cdot10^{-4}}}
% NHSixFarKtwoMargin [cert, margin] t9_far_v32.json :: margins.h2
\newcommand{\NHSixFarKtwoMargin}{\ensuremath{1.964}}
% NHSixFarKthree [cert, upper] t9_far_v32.json :: h3
\newcommand{\NHSixFarKthree}{\ensuremath{2.595\cdot10^{-8}}}
% NHSixFarKthreeThr [hand, exact] t9_far_v32.json :: thr.h3
\newcommand{\NHSixFarKthreeThr}{\ensuremath{1.17\cdot10^{-5}}}
% NHSixFarKthreeMargin [cert, margin] t9_far_v32.json :: margins.h3
\newcommand{\NHSixFarKthreeMargin}{\ensuremath{450.9}}
% NHSixFarKfour [cert, upper] t9_far_v32.json :: h4
\newcommand{\NHSixFarKfour}{\ensuremath{6.964\cdot10^{-9}}}
% NHSixFarKfourThr [hand, exact] t9_far_v32.json :: thr.h4
\newcommand{\NHSixFarKfourThr}{\ensuremath{1.03\cdot10^{-7}}}
% NHSixFarKfourMargin [cert, margin] t9_far_v32.json :: margins.h4
\newcommand{\NHSixFarKfourMargin}{\ensuremath{14.79}}
% NHSixFarSlope [cert, upper] t9_far_v32.json :: slope
\newcommand{\NHSixFarSlope}{\ensuremath{0.04153}}
% NHSixFarSlopeThr [hand, exact] t9_far_v32.json :: thr.slope
\newcommand{\NHSixFarSlopeThr}{\ensuremath{0.0800}}
% NHSixFarSlopeMargin [cert, margin] t9_far_v32.json :: margins.slope
\newcommand{\NHSixFarSlopeMargin}{\ensuremath{1.926}}
% NHSixFarNthree [cert, upper] t9_far_v32.json :: L2
\newcommand{\NHSixFarNthree}{\ensuremath{5.516\cdot10^{-7}}}
% NHSixFarNthreeThr [hand, exact] t9_far_v32.json :: thr.L2
\newcommand{\NHSixFarNthreeThr}{\ensuremath{3.32\cdot10^{-6}}}
% NHSixFarNthreeMargin [cert, margin] t9_far_v32.json :: margins.L2
\newcommand{\NHSixFarNthreeMargin}{\ensuremath{6.019}}
% NFarTstarMargin [cert, margin] lemmaO_v32_consts.json :: far_v32.case.margin_factor_lower
\newcommand{\NFarTstarMargin}{\ensuremath{15.5}}
% NFarTstarLower [cert, lower] lemmaO_v32_consts.json :: far_v32.case.t_star_lower
\newcommand{\NFarTstarLower}{\ensuremath{9.5108\cdot10^{-8}}}
% NFarTstarNeeded [cert, upper] lemmaO_v32_consts.json :: far_v32.case.needed
\newcommand{\NFarTstarNeeded}{\ensuremath{6.134\cdot10^{-9}}}
% NGammaModes [cert, exact] gamma0_final.json :: M
\newcommand{\NGammaModes}{\ensuremath{3762}}
% NGammaHalfwidth [cert, upper] gamma0_final.json :: b_j_halfwidth
\newcommand{\NGammaHalfwidth}{\ensuremath{8.138\cdot10^{-24}}}
% NGammaContraction [cert, upper] gamma0_final.json :: krawczyk_contraction
\newcommand{\NGammaContraction}{\ensuremath{3.01\cdot10^{-12}}}
% NGammaArcs [cert, exact] gamma0_final.json :: embedded.arcs
\newcommand{\NGammaArcs}{\ensuremath{701}}
% NGammaMaxTurn [cert, upper] gamma0_final.json :: embedded.maxturn
\newcommand{\NGammaMaxTurn}{\ensuremath{0.04997}}
% NGammaMaxPair [cert, upper] gamma0_final.json :: embedded.maxpair
\newcommand{\NGammaMaxPair}{\ensuremath{0.09597}}
% NGammaMinSep [cert, lower] gamma0_final.json :: embedded.minsep
\newcommand{\NGammaMinSep}{\ensuremath{8.89\cdot10^{-3}}}
% NGammaHalfGap [cert, lower] gamma0_final.json :: embedded.gap_cert
\newcommand{\NGammaHalfGap}{\ensuremath{0.0248219}}
% NPhaseISlabs [cert, exact] status.json :: slab
\newcommand{\NPhaseISlabs}{\ensuremath{96}}
% NPhaseIModes [cert, exact] status.json :: last_slab.M
\newcommand{\NPhaseIModes}{\ensuremath{2800}}
% NPhaseIKmaxLower [cert, lower] status.json :: kmax_enclosure.0
\newcommand{\NPhaseIKmaxLower}{\ensuremath{29.352874545}}
% NPhaseIKmaxUpper [cert, upper] status.json :: kmax_enclosure.1
\newcommand{\NPhaseIKmaxUpper}{\ensuremath{29.352874546}}
% NPhaseITLower [cert, lower] status.json :: t
\newcommand{\NPhaseITLower}{\ensuremath{4.597850539485830}}
% NPhaseITUpper [cert, upper] status.json :: t
\newcommand{\NPhaseITUpper}{\ensuremath{4.597850539485831}}
% NPeakSqrtLam [cert, upper] rate_table_fin12f4_a0.00015.json :: sqrtL_max
\newcommand{\NPeakSqrtLam}{\ensuremath{5.308\cdot10^{-8}}}
% NTailRate [cert, lower] rate_table_fin12f4_a0.00015.json :: tail.rho
\newcommand{\NTailRate}{\ensuremath{2.693\cdot10^{-3}}}
% NFloorSqrtLam [cert, upper] rate_table_fin12f4_a0.00015.json :: tail.sqrt_LF
\newcommand{\NFloorSqrtLam}{\ensuremath{9.872\cdot10^{-10}}}
% NAlphaFloor [cert, upper] rate_table_fin12f4_a0.00015.json :: tail.alpha_F
\newcommand{\NAlphaFloor}{\ensuremath{6.867\cdot10^{-7}}}
% NSupBetaVerifier [cert, upper] final_alpha_points_5e-8.json :: 1.5e-4.sup_beta
\newcommand{\NSupBetaVerifier}{\ensuremath{3.33811\cdot10^{-9}}}
% NMOneRamp [cert, upper] final_alpha_points_5e-8.json :: 1.5e-4.M1
\newcommand{\NMOneRamp}{\ensuremath{0.4328}}
% NSITwoPrime [cert, upper] FINAL_manifest_v33.json :: results.F12_sliver.sI2_prime
\newcommand{\NSITwoPrime}{\ensuremath{0.015457}}
% NLoopGainCoarse [cert, upper] U_cert_v35.json :: loop_gain.2.total
\newcommand{\NLoopGainCoarse}{\ensuremath{0.03904}}
% NRateMarginCoarse [cert, margin] U_cert_v35.json :: rate_margin
\newcommand{\NRateMarginCoarse}{\ensuremath{28.2}}
% NSOneEpsLower [cert, lower] bracket2.json :: eps_a
\newcommand{\NSOneEpsLower}{\ensuremath{1.5346198892965\cdot10^{-4}}}
% NSOneEpsUpper [cert, upper] bracket2.json :: eps_b
\newcommand{\NSOneEpsUpper}{\ensuremath{1.5346198899299\cdot10^{-4}}}
% NSOneRowsTip [cert, exact] ext_tip_check.json :: rows
\newcommand{\NSOneRowsTip}{\ensuremath{51246}}
% NSOneTipWidth [cert, upper] ext_tip_check.json :: max_pointwise_width.th
\newcommand{\NSOneTipWidth}{\ensuremath{3.2\cdot10^{-12}}}
% NSOneExtWidth [cert, upper] ext813_check.json :: max_pointwise_width.p
\newcommand{\NSOneExtWidth}{\ensuremath{5.7\cdot10^{-10}}}
% NSOneRowsMC [cert, exact] check_outA5.json :: prof_tip_vs_MC.rows_compared
\newcommand{\NSOneRowsMC}{\ensuremath{1025}}
% NSOneRowsMB [cert, exact] check_outA5.json :: ext813_vs_MB.rows_compared
\newcommand{\NSOneRowsMB}{\ensuremath{501}}
% NSOneMinMargin [cert, margin] check_outA5.json :: min:*.min_relative_margin.*
\newcommand{\NSOneMinMargin}{\ensuremath{1.2\cdot10^{-9}}}
% NTwoSegments [cert, exact] contour_result.json :: len:segments
\newcommand{\NTwoSegments}{\ensuremath{5058}}
% NTwoInteriorBoxes [cert, exact] interior_result.json :: len:
\newcommand{\NTwoInteriorBoxes}{\ensuremath{174}}
% NTwoContourBoxes [cert, exact] eta_cache.json :: len:
\newcommand{\NTwoContourBoxes}{\ensuremath{202}}
% NTwoEtaMax [cert, upper] eta_cache.json :: max:*.eta
\newcommand{\NTwoEtaMax}{\ensuremath{0.1163}}
% NTwoLipMax [cert, upper] eta_cache.json :: max:*.Lip
\newcommand{\NTwoLipMax}{\ensuremath{0.272}}
% NTwoLambdaZero [est, upper] cert_b0.25_th0.0_tau0.03.json :: Lam
\newcommand{\NTwoLambdaZero}{\ensuremath{1.790819}}
% NTwoLambdaOne [est, upper] cert_b0.25_th0.48_tau0.03.json :: Lam
\newcommand{\NTwoLambdaOne}{\ensuremath{1.768282}}
% NSOneDTightLower [doc, lower] doc:README.md :: D ∈ [2.50080627654, 2.50080627668]
\newcommand{\NSOneDTightLower}{\ensuremath{2.50080627654}}
% NSOneDTightUpper [doc, upper] doc:README.md :: D ∈ [2.50080627654, 2.50080627668]
\newcommand{\NSOneDTightUpper}{\ensuremath{2.50080627668}}
% NSOneKTightLower [doc, lower] doc:README.md :: k(0) ∈ [1.53982063762, 1.53982063822]
\newcommand{\NSOneKTightLower}{\ensuremath{1.53982063762}}
% NSOneKTightUpper [doc, upper] doc:README.md :: k(0) ∈ [1.53982063762, 1.53982063822]
\newcommand{\NSOneKTightUpper}{\ensuremath{1.53982063822}}
% NSOneDCondLower [doc, lower] doc:README.md :: D = −p(tip) ∈ [2.500806276619695, 2.500806276620049]
\newcommand{\NSOneDCondLower}{\ensuremath{2.500806276619695}}
% NSOneDCondUpper [doc, upper] doc:README.md :: D = −p(tip) ∈ [2.500806276619695, 2.500806276620049]
\newcommand{\NSOneDCondUpper}{\ensuremath{2.500806276620049}}
% NSOneKCondLower [doc, lower] doc:README.md :: k(0) ∈ [1.539820637871836, 1.5398206378721015]
\newcommand{\NSOneKCondLower}{\ensuremath{1.539820637871836}}
% NSOneKCondUpper [doc, upper] doc:README.md :: k(0) ∈ [1.539820637871836, 1.5398206378721015]
\newcommand{\NSOneKCondUpper}{\ensuremath{1.5398206378721015}}
% NSOneTauTipLower [doc, lower] doc:README.md :: τ_tip ∈ [10.249077651120778, 10.249077651121361]
\newcommand{\NSOneTauTipLower}{\ensuremath{10.249077651120778}}
% NSOneTauTipUpper [doc, upper] doc:README.md :: τ_tip ∈ [10.249077651120778, 10.249077651121361]
\newcommand{\NSOneTauTipUpper}{\ensuremath{10.249077651121361}}
% NSOneStateUncert [doc, upper] doc:README.md :: State uncertainty at σ = 8: 4.4e−13
\newcommand{\NSOneStateUncert}{\ensuremath{4.4\cdot10^{-13}}}
% NFarBoxEprime [doc, margin] doc:STEP4_RESULT.md :: (6.4287, 4.6917, 1.9940, 1.69999995)
\newcommand{\NFarBoxEprime}{\ensuremath{6.42}}
% NFarBoxXEsecond [doc, margin] doc:STEP4_RESULT.md :: (6.4287, 4.6917, 1.9940, 1.69999995)
\newcommand{\NFarBoxXEsecond}{\ensuremath{4.69}}
% NFarBoxXsqEthird [doc, margin] doc:STEP4_RESULT.md :: (6.4287, 4.6917, 1.9940, 1.69999995)
\newcommand{\NFarBoxXsqEthird}{\ensuremath{1.99}}
% NFarBoxE [doc, margin] doc:STEP4_RESULT.md :: (6.4287, 4.6917, 1.9940, 1.69999995)
\newcommand{\NFarBoxE}{\ensuremath{1.69}}
% NHFourKfourErr [doc, upper] doc:STEP4_RESULT.md :: | 5.60e-3 against 0.06 | ×10.7 |
\newcommand{\NHFourKfourErr}{\ensuremath{5.60\cdot10^{-3}}}
% NHFourKfourMargin [doc, margin] doc:STEP4_RESULT.md :: | 5.60e-3 against 0.06 | ×10.7 |
\newcommand{\NHFourKfourMargin}{\ensuremath{10.7}}
% NCoercivityCoefficient [doc, upper] doc:step3_proof_v32.tex :: \norm{e^{\tau\mathcal L}u}_\rho\le14.8\,e^{-0.0213\tau}\norm
\newcommand{\NCoercivityCoefficient}{\ensuremath{14.8}}
% NCoercivityRate [doc, lower] doc:step3_proof_v32.tex :: \norm{e^{\tau\mathcal L}u}_\rho\le14.8\,e^{-0.0213\tau}\norm
\newcommand{\NCoercivityRate}{\ensuremath{0.0213}}
% NRegionRampStart [hand, exact] hand :: 
\newcommand{\NRegionRampStart}{\ensuremath{14}}
% NRegionInterface [hand, exact] hand :: 
\newcommand{\NRegionInterface}{\ensuremath{28}}
% NRegionShellEnd [hand, exact] hand :: 
\newcommand{\NRegionShellEnd}{\ensuremath{60}}
% NRegionSegEnd [hand, exact] hand :: 
\newcommand{\NRegionSegEnd}{\ensuremath{350}}
% NRegionAnnEnd [hand, exact] hand :: 
\newcommand{\NRegionAnnEnd}{\ensuremath{700}}
% NZoneIIILow [hand, exact] hand :: 
\newcommand{\NZoneIIILow}{\ensuremath{253}}
% NZoneIIIHigh [hand, exact] hand :: 
\newcommand{\NZoneIIIHigh}{\ensuremath{379.5}}
% NConeCentreMin [hand, exact] hand :: 
\newcommand{\NConeCentreMin}{\ensuremath{300}}
% NConeCentreMax [hand, exact] hand :: 
\newcommand{\NConeCentreMax}{\ensuremath{380}}
% NFarCentreMin [hand, exact] hand :: 
\newcommand{\NFarCentreMin}{\ensuremath{490}}
% NConeL [hand, exact] hand :: 
\newcommand{\NConeL}{\ensuremath{50}}
% NConeX [hand, exact] hand :: 
\newcommand{\NConeX}{\ensuremath{3}}
% NFarL [hand, exact] hand :: 
\newcommand{\NFarL}{\ensuremath{113}}
% NFarX [hand, exact] hand :: 
\newcommand{\NFarX}{\ensuremath{2}}
 is found.
% ======================================================================

% ---- which document are we in ----------------------------------------
\ifdefined\ifsupplement\else\newif\ifsupplement\fi

% ---- packages ---------------------------------------------------------
\usepackage[T1]{fontenc}
\usepackage[utf8]{inputenc}
\usepackage{lmodern}
\usepackage[margin=1.1in]{geometry}
\usepackage{amsmath,amssymb,amsthm,mathtools,bm}
\usepackage{booktabs,array,longtable,tabularx,multirow}
\usepackage{enumitem}
\usepackage{graphicx}
\usepackage{xcolor}
\usepackage{tikz}
\usetikzlibrary{positioning,arrows.meta,calc}
% xr-hyper must come before hyperref
\usepackage{xr-hyper}
\usepackage[colorlinks=true,linkcolor=blue!50!black,citecolor=green!35!black,urlcolor=blue!50!black]{hyperref}

\raggedbottom
\setlength{\emergencystretch}{2em}
\setlist[enumerate]{leftmargin=*,itemsep=2pt,topsep=3pt}
\setlist[itemize]{leftmargin=*,itemsep=2pt,topsep=3pt}

% ---- numbering --------------------------------------------------------
% Main paper: sections 1,2,...; theorems 2.3 (by section); equations (2.3).
% Supplement: the Parts A..K are the \section's, numbered by letter, so
% theorems are A.1, E.4, ... (do not confuse "Theorem A" of the paper, which
% has no dot, with "Theorem A.1" of Part A).
\ifsupplement
  \renewcommand{\thesection}{\Alph{section}}
\fi
\numberwithin{equation}{section}

% ---- theorem environments (numbered by section, one shared counter) ----
\theoremstyle{plain}
\newtheorem{theorem}{Theorem}[section]
\newtheorem{proposition}[theorem]{Proposition}
\newtheorem{lemma}[theorem]{Lemma}
\newtheorem{corollary}[theorem]{Corollary}
\newtheorem{maintheorem}[theorem]{Main Theorem}
% Theorems A, B, C, D of the paper: own counter, lettered, no section prefix
\newtheorem{thmlet}{Theorem}
\renewcommand{\thethmlet}{\Alph{thmlet}}
\theoremstyle{definition}
\newtheorem{definition}[theorem]{Definition}
\newtheorem{convention}[theorem]{Convention}
\newtheorem{example}[theorem]{Example}
\theoremstyle{remark}
\newtheorem{remark}[theorem]{Remark}

% A list of hypotheses labelled (H1),(H2),... : \begin{hyplist}{H} ... \end{hyplist}
% (also {O}, {L}, {B}, {K}, {W} for the families of section 5.3 of OUTLINE.md).
\newenvironment{hyplist}[1]%
  {\begin{enumerate}[label=\textup{(#1\arabic*)},leftmargin=3em,itemsep=2pt]}%
  {\end{enumerate}}

% ---- trust classes (OUTLINE 4.1, 5.1.3) -------------------------------
% Supplement: explicit tags.  Main paper: certified is the default and
% unmarked; every non-certified number is marked "(floating point)".
\ifsupplement
  \newcommand{\tcert}{\textup{[cert]}}   % arb/outward-rounded, pinned script and key
  \newcommand{\test}{\textup{[est]}}     % floating point or estimate; never in a theorem
  \newcommand{\thand}{\textup{[hand]}}   % proved by hand in the text
  \newcommand{\tlit}{\textup{[lit]}}     % cited from the literature
\else
  \newcommand{\tcert}{}
  \newcommand{\test}{\unskip\ \textup{(floating point)}}
  \newcommand{\thand}{}
  \newcommand{\tlit}{}
\fi

% ---- gaps: visible, counted, an error in a final build -----------------
% Usage: \GAP{one sentence saying what is missing}   (and an entry in ISSUES.md)
\newcounter{gapcount}
\newif\iffinalpaper   % \finalpapertrue makes every \GAP a hard error
\newcommand{\GAP}[1]{\stepcounter{gapcount}\typeout{GAP-MARKER \thegapcount}%
  \iffinalpaper\PackageError{paper}{\string\GAP\space used in a final build}{Remove all gaps.}\fi
  {\color{red}\sffamily\bfseries[GAP~\thegapcount:\ }{\color{red}\sffamily #1}{\color{red}\sffamily\bfseries]}}
\AtEndDocument{\typeout{GAP-MARKER total: \thegapcount}}

% ---- file names and JSON keys (supplement only) -------------------------
% No internal path appears in the main paper; a use there is a warning.
\newcommand{\fn}[1]{%
  \ifsupplement\else\PackageWarning{paper}{\string\fn\space used in the main paper}\fi
  \texttt{\detokenize{#1}}}
\newcommand{\key}[1]{\texttt{\detokenize{#1}}}

% ---- cross references between main paper and supplement ----------------
% Labels carry a unit prefix (section 5.3 of OUTLINE.md and DRAFTER_GUIDE.md):
%   main paper:  intro: obj: s3: s4: s5: s6: s7: s8: s9:   (kind: thm- lem- eq- tab- fig- ...)
%   supplement:  pA: pB: ... pK:
% so all labels are globally unique and the prefixes below are only a second guard.
% \Xref{label}  : reference to the OTHER document (hyperlinked via xr-hyper)
% \Lref{label}  : reference in the SAME document (plain \ref)
\ifsupplement
  \externaldocument[M-]{main}[main.pdf]
  \newcommand{\Xref}[1]{\ref{M-#1}}
\else
  \externaldocument[S-]{supplement}[supplement.pdf]
  \newcommand{\Xref}[1]{\ref{S-#1}}
\fi
\newcommand{\Lref}[1]{\ref{#1}}
% \Supp{pE}{lem-chart} -> "Lemma E.3" style text is produced by the drafter:
%   "Lemma~\Xref{pE:lem-chart} of the supplement".  \SuppPart{E} gives "Part E".
\newcommand{\SuppPart}[1]{Part~#1 of the supplement}

% ---- numbers (generated by unit U0; macros only, no literal certified numbers) ----
\IfFileExists{numbers/numbers_paper.tex}%
  {% numbers_paper.tex -- GENERATED from ledger.json by build_ledger.py; do not edit.
% Rounding: upper bounds up, lower bounds and margins down at the printed digit; margin = threshold / bound.
% Use as $\times\NName$ for margins.  kind: cert certified, est estimate, hand hypothesis constant, doc only a document quotes it.
% NUEhatZero [cert, upper] U_cert_v35.json :: supersolution.EA.0
\newcommand{\NUEhatZero}{\ensuremath{1.2018\cdot10^{-6}}}
% NUEhatOne [cert, upper] U_cert_v35.json :: supersolution.EA.1
\newcommand{\NUEhatOne}{\ensuremath{1.0441\cdot10^{-4}}}
% NUEhatTwo [cert, upper] U_cert_v35.json :: supersolution.EA.2
\newcommand{\NUEhatTwo}{\ensuremath{0.056694}}
% NUBhatZero [cert, upper] U_cert_v35.json :: supersolution.Bhat.0
\newcommand{\NUBhatZero}{\ensuremath{4.8735\cdot10^{-6}}}
% NUBhatOne [cert, upper] U_cert_v35.json :: supersolution.Bhat.1
\newcommand{\NUBhatOne}{\ensuremath{5.5933\cdot10^{-4}}}
% NUBhatTwo [cert, upper] U_cert_v35.json :: supersolution.Bhat.2
\newcommand{\NUBhatTwo}{\ensuremath{0.36675}}
% NUSupersolutionCZeroZero [cert, upper] U_cert_v35.json :: supersolution.C0.0
\newcommand{\NUSupersolutionCZeroZero}{\ensuremath{1.1340\cdot10^{-6}}}
% NUSupersolutionCZeroOne [cert, upper] U_cert_v35.json :: supersolution.C0.1
\newcommand{\NUSupersolutionCZeroOne}{\ensuremath{1.6623\cdot10^{-4}}}
% NUSupersolutionCOneZero [cert, upper] U_cert_v35.json :: supersolution.C1.0
\newcommand{\NUSupersolutionCOneZero}{\ensuremath{1.5241\cdot10^{-9}}}
% NUSupersolutionCOneOne [cert, upper] U_cert_v35.json :: supersolution.C1.1
\newcommand{\NUSupersolutionCOneOne}{\ensuremath{5.5255\cdot10^{-8}}}
% NUSupersolutionNTwoZero [cert, upper] U_cert_v35.json :: supersolution.N2.0
\newcommand{\NUSupersolutionNTwoZero}{\ensuremath{0.071472}}
% NUSupersolutionNTwoOne [cert, upper] U_cert_v35.json :: supersolution.N2.1
\newcommand{\NUSupersolutionNTwoOne}{\ensuremath{0.33012}}
% NUSupersolutionNTwoTwo [cert, upper] U_cert_v35.json :: supersolution.N2.2
\newcommand{\NUSupersolutionNTwoTwo}{\ensuremath{5.0406}}
% NUSupersolutionNTwoOverFourNTwoRefZero [cert, upper] U_cert_v35.json :: supersolution.N2_over_4N2ref.0
\newcommand{\NUSupersolutionNTwoOverFourNTwoRefZero}{\ensuremath{0.18}}
% NUSupersolutionNTwoOverFourNTwoRefOne [cert, upper] U_cert_v35.json :: supersolution.N2_over_4N2ref.1
\newcommand{\NUSupersolutionNTwoOverFourNTwoRefOne}{\ensuremath{0.42}}
% NUSupersolutionNTwoOverFourNTwoRefTwo [cert, upper] U_cert_v35.json :: supersolution.N2_over_4N2ref.2
\newcommand{\NUSupersolutionNTwoOverFourNTwoRefTwo}{\ensuremath{0.26}}
% NUSupersolutionMOneZHypCheckZonesIIISOneTwoBhat [cert, upper] U_cert_v35.json :: supersolution.M1z_hyp_check.zones_I_II.S1_2Bhat
\newcommand{\NUSupersolutionMOneZHypCheckZonesIIISOneTwoBhat}{\ensuremath{0.28585}}
% NUSupersolutionMOneZHypCheckZoneIIIBOnePlusPUOne [cert, upper] U_cert_v35.json :: supersolution.M1z_hyp_check.zone_III.b1_plus_pU1
\newcommand{\NUSupersolutionMOneZHypCheckZoneIIIBOnePlusPUOne}{\ensuremath{0.50786}}
% NUMarginLateMOne [cert, margin] U_cert_v35.json :: margins.late_M1
\newcommand{\NUMarginLateMOne}{\ensuremath{2.353}}
% NUMarginLateMZero [cert, margin] U_cert_v35.json :: margins.late_M0
\newcommand{\NUMarginLateMZero}{\ensuremath{4.892}}
% NUMarginsLateMZeroZ [cert, margin] U_cert_v35.json :: margins.late_M0z
\newcommand{\NUMarginsLateMZeroZ}{\ensuremath{12.99}}
% NUMarginPresentMOne [cert, margin] U_cert_v35.json :: margins.present_M1_min.value
\newcommand{\NUMarginPresentMOne}{\ensuremath{2.770}}
% NUMarginPresentMZero [cert, margin] U_cert_v35.json :: margins.present_M0_min.value
\newcommand{\NUMarginPresentMZero}{\ensuremath{2.225}}
% NUMarginsPresentMZeroZMinValue [cert, margin] U_cert_v35.json :: margins.present_M0z_min.value
\newcommand{\NUMarginsPresentMZeroZMinValue}{\ensuremath{14.37}}
% NULoopGainZeroTotal [cert, upper] U_cert_v35.json :: loop_gain.0.total
\newcommand{\NULoopGainZeroTotal}{\ensuremath{1.23\cdot10^{-7}}}
% NULoopGainOneTotal [cert, upper] U_cert_v35.json :: loop_gain.1.total
\newcommand{\NULoopGainOneTotal}{\ensuremath{0.037169}}
% NULoopGainTwoNeighbour [cert, upper] U_cert_v35.json :: loop_gain.2.neighbour
\newcommand{\NULoopGainTwoNeighbour}{\ensuremath{0.037018}}
% NULoopGainTwoThroughA [cert, upper] U_cert_v35.json :: loop_gain.2.through_A
\newcommand{\NULoopGainTwoThroughA}{\ensuremath{0.0020172}}
% NUfineBhatZero [cert, upper] U_cert_v35.json :: ufine.Bhat.0
\newcommand{\NUfineBhatZero}{\ensuremath{4.8723\cdot10^{-6}}}
% NUfineBhatOne [cert, upper] U_cert_v35.json :: ufine.Bhat.1
\newcommand{\NUfineBhatOne}{\ensuremath{2.6266\cdot10^{-4}}}
% NUfineBhatTwo [cert, upper] U_cert_v35.json :: ufine.Bhat.2
\newcommand{\NUfineBhatTwo}{\ensuremath{0.025842}}
% NUcdeltaU [cert, upper] U_cert_v35.json :: launch.c_dU
\newcommand{\NUcdeltaU}{\ensuremath{1.7187\cdot10^{-8}}}
% NULaunchWmin [cert, lower] U_cert_v35.json :: launch.wmin
\newcommand{\NULaunchWmin}{\ensuremath{0.73492}}
% NULaunchCAD [cert, upper] U_cert_v35.json :: launch.CA_D
\newcommand{\NULaunchCAD}{\ensuremath{2.2429}}
% NULaunchQD [cert, upper] U_cert_v35.json :: launch.qD
\newcommand{\NULaunchQD}{\ensuremath{0.018509}}
% NULaunchDeltaInZero [cert, upper] U_cert_v35.json :: launch.Delta_in.0
\newcommand{\NULaunchDeltaInZero}{\ensuremath{3.131\cdot10^{-12}}}
% NULaunchDeltaInOne [cert, upper] U_cert_v35.json :: launch.Delta_in.1
\newcommand{\NULaunchDeltaInOne}{\ensuremath{3.287\cdot10^{-12}}}
% NULaunchDeltaInTwo [cert, upper] U_cert_v35.json :: launch.Delta_in.2
\newcommand{\NULaunchDeltaInTwo}{\ensuremath{3.451\cdot10^{-12}}}
% NULaunchPhiZero [cert, upper] U_cert_v35.json :: launch.Phi.0
\newcommand{\NULaunchPhiZero}{\ensuremath{6.434\cdot10^{-9}}}
% NULaunchPhiOne [cert, upper] U_cert_v35.json :: launch.Phi.1
\newcommand{\NULaunchPhiOne}{\ensuremath{1.4217\cdot10^{-8}}}
% NULaunchPhiTwo [cert, upper] U_cert_v35.json :: launch.Phi.2
\newcommand{\NULaunchPhiTwo}{\ensuremath{3.8404\cdot10^{-8}}}
% NUGammaZero [cert, upper] U_cert_v35.json :: Gamma.0
\newcommand{\NUGammaZero}{\ensuremath{4.0000358}}
% NUGammaOne [cert, upper] U_cert_v35.json :: Gamma.1
\newcommand{\NUGammaOne}{\ensuremath{16.000322}}
% NUGammaTwo [cert, upper] U_cert_v35.json :: Gamma.2
\newcommand{\NUGammaTwo}{\ensuremath{64.00218}}
% NUKappa [cert, upper] U_cert_v35.json :: kappa
\newcommand{\NUKappa}{\ensuremath{8.0000358}}
% NUKappaTwo [cert, upper] U_cert_v35.json :: kappa2
\newcommand{\NUKappaTwo}{\ensuremath{64.00061}}
% NUGhatZeroZero [cert, upper] U_cert_v35.json :: Ghat.0.0
\newcommand{\NUGhatZeroZero}{\ensuremath{51322.6}}
% NUGhatOneOne [cert, upper] U_cert_v35.json :: Ghat.1.1
\newcommand{\NUGhatOneOne}{\ensuremath{33962.9}}
% NUGhatTwoTwo [cert, upper] U_cert_v35.json :: Ghat.2.2
\newcommand{\NUGhatTwoTwo}{\ensuremath{348341.1}}
% NUGhatOneZero [cert, upper] U_cert_v35.json :: Ghat.1.0
\newcommand{\NUGhatOneZero}{\ensuremath{4.938}}
% NUGhatTwoZero [cert, upper] U_cert_v35.json :: Ghat.2.0
\newcommand{\NUGhatTwoZero}{\ensuremath{27.984}}
% NUGhatTwoOne [cert, upper] U_cert_v35.json :: Ghat.2.1
\newcommand{\NUGhatTwoOne}{\ensuremath{458.53}}
% NUKZeroZero [cert, upper] U_cert_v35.json :: K.0.0
\newcommand{\NUKZeroZero}{\ensuremath{8034.21}}
% NUKZeroOne [cert, upper] U_cert_v35.json :: K.0.1
\newcommand{\NUKZeroOne}{\ensuremath{724.89}}
% NUKOneZero [cert, upper] U_cert_v35.json :: K.1.0
\newcommand{\NUKOneZero}{\ensuremath{0.08152}}
% NUKOneOne [cert, upper] U_cert_v35.json :: K.1.1
\newcommand{\NUKOneOne}{\ensuremath{6336.94}}
% NUKOneTwo [cert, upper] U_cert_v35.json :: K.1.2
\newcommand{\NUKOneTwo}{\ensuremath{545.59}}
% NUKOneThree [cert, upper] U_cert_v35.json :: K.1.3
\newcommand{\NUKOneThree}{\ensuremath{0.58700}}
% NUKTwoZero [cert, upper] U_cert_v35.json :: K.2.0
\newcommand{\NUKTwoZero}{\ensuremath{0.87186}}
% NUKTwoOne [cert, upper] U_cert_v35.json :: K.2.1
\newcommand{\NUKTwoOne}{\ensuremath{528.62}}
% NUKTwoTwo [cert, upper] U_cert_v35.json :: K.2.2
\newcommand{\NUKTwoTwo}{\ensuremath{23425.98}}
% NUKTwoThree [cert, upper] U_cert_v35.json :: K.2.3
\newcommand{\NUKTwoThree}{\ensuremath{555.62}}
% NUFootwOne [cert, upper] U_cert_v35.json :: footw.1
\newcommand{\NUFootwOne}{\ensuremath{2.6361}}
% NUFootwTwo [cert, upper] U_cert_v35.json :: footw.2
\newcommand{\NUFootwTwo}{\ensuremath{2663.3}}
% NUCThreeThree [cert, upper] U_cert_v35.json :: c33
\newcommand{\NUCThreeThree}{\ensuremath{108.02}}
% NUCTwoTwoThree [cert, upper] U_cert_v35.json :: c223
\newcommand{\NUCTwoTwoThree}{\ensuremath{7.3284}}
% NUCEta [cert, upper] U_cert_v35.json :: C_eta
\newcommand{\NUCEta}{\ensuremath{13.376}}
% NULMax [cert, upper] U_cert_v35.json :: L_max
\newcommand{\NULMax}{\ensuremath{3.5352}}
% NURRatio [cert, upper] U_cert_v35.json :: R_ratio
\newcommand{\NURRatio}{\ensuremath{11.593}}
% NUSweepsSigmaInCoverSup [cert, upper] U_cert_v35.json :: sweeps.Sigma_in.cover_sup
\newcommand{\NUSweepsSigmaInCoverSup}{\ensuremath{5.0892}}
% NUSweepsSigmaECoverSup [cert, upper] U_cert_v35.json :: sweeps.Sigma_E.cover_sup
\newcommand{\NUSweepsSigmaECoverSup}{\ensuremath{0.30846}}
% NUSweepsSXCoverSup [cert, upper] U_cert_v35.json :: sweeps.S_X.cover_sup
\newcommand{\NUSweepsSXCoverSup}{\ensuremath{1.5243}}
% NUSweepsSigmaXCoverSup [cert, upper] U_cert_v35.json :: sweeps.Sigma_X.cover_sup
\newcommand{\NUSweepsSigmaXCoverSup}{\ensuremath{3.0081}}
% NUSweepsSigmaRsTwoCoverSup [cert, upper] U_cert_v35.json :: sweeps.Sigma_rs.2.cover_sup
\newcommand{\NUSweepsSigmaRsTwoCoverSup}{\ensuremath{4.3350\cdot10^{-3}}}
% NUDataTermsLateDZero [cert, upper] U_cert_v35.json :: data_terms_late.D.0
\newcommand{\NUDataTermsLateDZero}{\ensuremath{1.7330\cdot10^{-6}}}
% NUDataTermsLateDOne [cert, upper] U_cert_v35.json :: data_terms_late.D.1
\newcommand{\NUDataTermsLateDOne}{\ensuremath{2.5353\cdot10^{-4}}}
% NUDataTermsLateSlopeShareCOneSigmaXOverDkZero [cert, upper] U_cert_v35.json :: data_terms_late.slope_share_C1SigmaX_over_Dk.0
\newcommand{\NUDataTermsLateSlopeShareCOneSigmaXOverDkZero}{\ensuremath{2.7\cdot10^{-3}}}
% NUDataTermsLateSlopeShareCOneSigmaXOverDkOne [cert, upper] U_cert_v35.json :: data_terms_late.slope_share_C1SigmaX_over_Dk.1
\newcommand{\NUDataTermsLateSlopeShareCOneSigmaXOverDkOne}{\ensuremath{6.6\cdot10^{-4}}}
% NURLtTwoMarginZero [cert, margin] U_cert_v35.json :: R_lt2_margin.0
\newcommand{\NURLtTwoMarginZero}{\ensuremath{4.432}}
% NURLtTwoMarginOne [cert, margin] U_cert_v35.json :: R_lt2_margin.1
\newcommand{\NURLtTwoMarginOne}{\ensuremath{6.697}}
% NUThetaTwoEight [cert, upper] U_cert_v35.json :: theta28
\newcommand{\NUThetaTwoEight}{\ensuremath{0.37011}}
% NUDFiveSix [cert, upper] U_cert_v35.json :: D56
\newcommand{\NUDFiveSix}{\ensuremath{1.0368}}
% NUDTwoEight [cert, upper] U_cert_v35.json :: D28
\newcommand{\NUDTwoEight}{\ensuremath{1.5876}}
% NUSThreeThreeFootOverBhatTwo [cert, upper] U_cert_v35.json :: S33_foot_over_Bhat2
\newcommand{\NUSThreeThreeFootOverBhatTwo}{\ensuremath{1.69\cdot10^{-3}}}
% NSAHJWKappaJ [cert, upper] A_interface.json :: HJW.kappa_J
\newcommand{\NSAHJWKappaJ}{\ensuremath{0.042351}}
% NSAGpZero [cert, lower] A_interface.json :: res.tight.gp.0
\newcommand{\NSAGpZero}{\ensuremath{0.24292}}
% NSAGpOne [cert, lower] A_interface.json :: res.tight.gp.1
\newcommand{\NSAGpOne}{\ensuremath{0.24106}}
% NSAGpTwo [cert, lower] A_interface.json :: res.tight.gp.2
\newcommand{\NSAGpTwo}{\ensuremath{0.11961}}
% NSAAminC [cert, lower] A_interface.json :: res.tight.amin_c
\newcommand{\NSAAminC}{\ensuremath{0.97623}}
% NSARhoJ [cert, upper] A_interface.json :: res.tight.rhoJ
\newcommand{\NSARhoJ}{\ensuremath{5.4995\cdot10^{-7}}}
% NSAKappaOTwoTwo [cert, upper] A_interface.json :: res.tight.kappa_o.2.2
\newcommand{\NSAKappaOTwoTwo}{\ensuremath{1.2416\cdot10^{-5}}}
% NSADsrcZero [cert, upper] A_interface.json :: res.tight.Dsrc.0
\newcommand{\NSADsrcZero}{\ensuremath{1.8092\cdot10^{-7}}}
% NSADsrcOne [cert, upper] A_interface.json :: res.tight.Dsrc.1
\newcommand{\NSADsrcOne}{\ensuremath{9.8276\cdot10^{-5}}}
% NSABZeroDWindowSlopeZero [cert, upper] A_interface.json :: res.tight.B0.D_window_slope.0
\newcommand{\NSABZeroDWindowSlopeZero}{\ensuremath{8.7399\cdot10^{-10}}}
% NSABZeroDWindowSlopeOne [cert, upper] A_interface.json :: res.tight.B0.D_window_slope.1
\newcommand{\NSABZeroDWindowSlopeOne}{\ensuremath{5.2599\cdot10^{-8}}}
% NSAFTwoZero [cert, upper] A_interface.json :: res.tight.F2.0
\newcommand{\NSAFTwoZero}{\ensuremath{5.6506\cdot10^{-13}}}
% NSAFTwoOne [cert, upper] A_interface.json :: res.tight.F2.1
\newcommand{\NSAFTwoOne}{\ensuremath{5.9320\cdot10^{-13}}}
% NSACoupOne [cert, upper] A_interface.json :: res.tight.coup.1
\newcommand{\NSACoupOne}{\ensuremath{1.1568\cdot10^{-12}}}
% NSACOneZeroZero [cert, upper] A_interface.json :: res.tight.C1_0.0
\newcommand{\NSACOneZeroZero}{\ensuremath{7.2060\cdot10^{-10}}}
% NSACOneZeroOne [cert, upper] A_interface.json :: res.tight.C1_0.1
\newcommand{\NSACOneZeroOne}{\ensuremath{4.3143\cdot10^{-8}}}
% NSABZeroSupCoreZero [cert, upper] A_interface.json :: res.tight.B0.sup_core.0
\newcommand{\NSABZeroSupCoreZero}{\ensuremath{4.23\cdot10^{-5}}}
% NSACoreFluxMarginTwoFLUXZeroOverKZero [cert, margin] A_interface.json :: res.tight.core.flux_margin.2.FLUX0_over_K0
\newcommand{\NSACoreFluxMarginTwoFLUXZeroOverKZero}{\ensuremath{3.396}}
% NSACoreFluxMarginTwoFLUXTwoOverKTwo [cert, margin] A_interface.json :: res.tight.core.flux_margin.2.FLUX2_over_K2
\newcommand{\NSACoreFluxMarginTwoFLUXTwoOverKTwo}{\ensuremath{6.012}}
% NSAMarginFluxZero [cert, margin] asm_design_N2ref4.json :: res.tight.flux.beta.2.margin_FLUX0
\newcommand{\NSAMarginFluxZero}{\ensuremath{3.365}}
% NSAMarginFluxTwo [cert, margin] asm_design_N2ref4.json :: res.tight.flux.beta.2.margin_FLUX2
\newcommand{\NSAMarginFluxTwo}{\ensuremath{5.955}}
% NSAfourSZero [cert, upper] asm_design_N2ref4.json :: res.tight.S0
\newcommand{\NSAfourSZero}{\ensuremath{7.1878\cdot10^{-3}}}
% NSAfourSOne [cert, upper] asm_design_N2ref4.json :: res.tight.S1
\newcommand{\NSAfourSOne}{\ensuremath{0.10430}}
% NSAMarginCoreSZeroV [cert, margin] asm_design_N2ref4.json :: res.tight.margins.S0v_vs_0.0153
\newcommand{\NSAMarginCoreSZeroV}{\ensuremath{2.116}}
% NSAMarginCoreSOneV [cert, margin] asm_design_N2ref4.json :: res.tight.margins.S1v_vs_0.6
\newcommand{\NSAMarginCoreSOneV}{\ensuremath{5.748}}
% NSAMarginCoreMOne [cert, margin] asm_design_N2ref4.json :: res.tight.margins.M1core
\newcommand{\NSAMarginCoreMOne}{\ensuremath{4.409}}
% NSAMarginCoreMZero [cert, margin] asm_design_N2ref4.json :: res.tight.margins.M0core
\newcommand{\NSAMarginCoreMZero}{\ensuremath{5.450}}
% NSAMarginGhat [cert, margin] asm_design_N2ref4.json :: res.tight.margins.Ghat22
\newcommand{\NSAMarginGhat}{\ensuremath{1.917}}
% NSAMarginRampMOne [cert, margin] asm_design_N2ref4.json :: res.tight.margins.ramp M1
\newcommand{\NSAMarginRampMOne}{\ensuremath{1.184}}
% NSAfourInfoZeroDWindowConst [cert, upper] asm_design_N2ref4.json :: res.tight.info.0.D_window_const
\newcommand{\NSAfourInfoZeroDWindowConst}{\ensuremath{1.7732\cdot10^{-7}}}
% NSAfourInfoOneDWindowConst [cert, upper] asm_design_N2ref4.json :: res.tight.info.1.D_window_const
\newcommand{\NSAfourInfoOneDWindowConst}{\ensuremath{9.8058\cdot10^{-5}}}
% NZDeltaRZero [cert, upper] zone3_interface.json :: DeltaR.0
\newcommand{\NZDeltaRZero}{\ensuremath{4.4986\cdot10^{-8}}}
% NZDeltaROne [cert, upper] zone3_interface.json :: DeltaR.1
\newcommand{\NZDeltaROne}{\ensuremath{6.7972\cdot10^{-6}}}
% NZDeltaRTwo [cert, upper] zone3_interface.json :: DeltaR.2
\newcommand{\NZDeltaRTwo}{\ensuremath{0.011480}}
% NZDeltaRBZero [cert, upper] zone3_interface.json :: DeltaR_b.0
\newcommand{\NZDeltaRBZero}{\ensuremath{2.66\cdot10^{-12}}}
% NZDeltaRBOne [cert, upper] zone3_interface.json :: DeltaR_b.1
\newcommand{\NZDeltaRBOne}{\ensuremath{3.26\cdot10^{-9}}}
% NZDeltaRBTwo [cert, upper] zone3_interface.json :: DeltaR_b.2
\newcommand{\NZDeltaRBTwo}{\ensuremath{4.35\cdot10^{-8}}}
% NZDeltaRDetailNSqZero [cert, upper] zone3_interface.json :: DeltaR_detail.N_sq.0
\newcommand{\NZDeltaRDetailNSqZero}{\ensuremath{8.28\cdot10^{-5}}}
% NZDeltaRDetailNSqOne [cert, upper] zone3_interface.json :: DeltaR_detail.N_sq.1
\newcommand{\NZDeltaRDetailNSqOne}{\ensuremath{0.10149}}
% NZDeltaRDetailNSqTwo [cert, upper] zone3_interface.json :: DeltaR_detail.N_sq.2
\newcommand{\NZDeltaRDetailNSqTwo}{\ensuremath{1.3567}}
% NZDeltaRDetailNSqThree [cert, upper] zone3_interface.json :: DeltaR_detail.N_sq.3
\newcommand{\NZDeltaRDetailNSqThree}{\ensuremath{25350}}
% NZDeltaRDetailPartsTwoPairsThree [cert, upper] zone3_interface.json :: DeltaR_detail.parts.2.pairs.3
\newcommand{\NZDeltaRDetailPartsTwoPairsThree}{\ensuremath{0.011298}}
% NZDeltaRDetailPartsTwoLFourFoot [cert, upper] zone3_interface.json :: DeltaR_detail.parts.2.L4_foot
\newcommand{\NZDeltaRDetailPartsTwoLFourFoot}{\ensuremath{1.6\cdot10^{-4}}}
% NZDeltaRDetailPartsTwoCubic [cert, upper] zone3_interface.json :: DeltaR_detail.parts.2.cubic
\newcommand{\NZDeltaRDetailPartsTwoCubic}{\ensuremath{9.3\cdot10^{-6}}}
% NZDeltaRDetailPartsTwoSThreeThreeFoot [cert, upper] zone3_interface.json :: DeltaR_detail.parts.2.S33_foot
\newcommand{\NZDeltaRDetailPartsTwoSThreeThreeFoot}{\ensuremath{3.7\cdot10^{-7}}}
% NZNIIIZero [cert, upper] zone3_interface.json :: n_III.0
\newcommand{\NZNIIIZero}{\ensuremath{0.071445}}
% NZNIIIOne [cert, upper] zone3_interface.json :: n_III.1
\newcommand{\NZNIIIOne}{\ensuremath{0.32707}}
% NZNIIITwo [cert, upper] zone3_interface.json :: n_III.2
\newcommand{\NZNIIITwo}{\ensuremath{3.0445}}
% NZTauIII [cert, upper] zone3_interface.json :: tau_III
\newcommand{\NZTauIII}{\ensuremath{4.0708}}
% NZoneIIIBoxZero [cert, upper] zone3_interface.json :: box_III_380.b0
\newcommand{\NZoneIIIBoxZero}{\ensuremath{0.10032}}
% NZoneIIIBoxOne [cert, upper] zone3_interface.json :: box_III_380.b1
\newcommand{\NZoneIIIBoxOne}{\ensuremath{0.41974}}
% NZoneIIIBoxZ [cert, upper] zone3_interface.json :: box_III_380.bz
\newcommand{\NZoneIIIBoxZ}{\ensuremath{0.014291}}
% NZBoxIIIThreeEightZeroMarginsBZero [cert, margin] zone3_interface.json :: box_III_380.margins.b0
\newcommand{\NZBoxIIIThreeEightZeroMarginsBZero}{\ensuremath{1.495}}
% NZoneIIIMarginThinnest [cert, margin] zone3_interface.json :: box_III_380.margins.b1
\newcommand{\NZoneIIIMarginThinnest}{\ensuremath{1.238}}
% NZBoxIIIThreeEightZeroMarginsBz [cert, margin] zone3_interface.json :: box_III_380.margins.bz
\newcommand{\NZBoxIIIThreeEightZeroMarginsBz}{\ensuremath{11.19}}
% NZBoxIIISevenZeroZeroBZero [cert, upper] zone3_interface.json :: box_III_700.b0
\newcommand{\NZBoxIIISevenZeroZeroBZero}{\ensuremath{0.22757}}
% NZBoxIIISevenZeroZeroBOne [cert, upper] zone3_interface.json :: box_III_700.b1
\newcommand{\NZBoxIIISevenZeroZeroBOne}{\ensuremath{0.54234}}
% NZUfineKTwoAtBhatRef [cert, upper] zone3_interface.json :: ufine.K2_at_Bhat_ref
\newcommand{\NZUfineKTwoAtBhatRef}{\ensuremath{0.046956}}
% NZUfineKThreeStripAtBhatRef [cert, upper] zone3_interface.json :: ufine.K3_strip_at_Bhat_ref
\newcommand{\NZUfineKThreeStripAtBhatRef}{\ensuremath{0.054140}}
% NRecoveryMinMargin [cert, margin] zone3_interface.json :: cone_family.H6.at_Bhat_ref.min_margin
\newcommand{\NRecoveryMinMargin}{\ensuremath{1.860}}
% NZConeFamilyHSixAtBhatRefConeSlopeMargin [cert, margin] zone3_interface.json :: cone_family.H6.at_Bhat_ref.cone.slope.margin
\newcommand{\NZConeFamilyHSixAtBhatRefConeSlopeMargin}{\ensuremath{2.46}}
% NZConeFamilyHSixAtBhatRefConeKTwoMargin [cert, margin] zone3_interface.json :: cone_family.H6.at_Bhat_ref.cone.K2.margin
\newcommand{\NZConeFamilyHSixAtBhatRefConeKTwoMargin}{\ensuremath{2.54}}
% NZConeFamilyHSixAtBhatRefConeKThreeMargin [cert, margin] zone3_interface.json :: cone_family.H6.at_Bhat_ref.cone.K3.margin
\newcommand{\NZConeFamilyHSixAtBhatRefConeKThreeMargin}{\ensuremath{2.51}}
% NConeTstarMargin [cert, margin] zone3_interface.json :: acceptance.tstar_margin
\newcommand{\NConeTstarMargin}{\ensuremath{2.30}}
% NUoneTFiveYaddLevelOneTwo [cert, upper] u1_ghat.json :: t5_Yadd.level1.2
\newcommand{\NUoneTFiveYaddLevelOneTwo}{\ensuremath{0.0044}}
% NUoneTFiveYaddLevelTwoTwo [cert, upper] u1_ghat.json :: t5_Yadd.level2.2
\newcommand{\NUoneTFiveYaddLevelTwoTwo}{\ensuremath{0.1081}}
% NUoneTFiveYaddLevelTwoThree [cert, upper] u1_ghat.json :: t5_Yadd.level2.3
\newcommand{\NUoneTFiveYaddLevelTwoThree}{\ensuremath{0.374}}
% NUoneWeightThetaGammaZero [cert, upper] u1_ghat.json :: weight.Theta.Gamma.0
\newcommand{\NUoneWeightThetaGammaZero}{\ensuremath{4.00000000000000000000054}}
% NRKZeroZero [cert, upper] ramp_U.json :: K.0.0
\newcommand{\NRKZeroZero}{\ensuremath{13436.03}}
% NRKOneOne [cert, upper] ramp_U.json :: K.1.1
\newcommand{\NRKOneOne}{\ensuremath{10526.19}}
% NRKTwoTwo [cert, upper] ramp_U.json :: K.2.2
\newcommand{\NRKTwoTwo}{\ensuremath{39482.30}}
% NUoneCasesPrimaryYNextZero [cert, upper] u1_ghat.json :: cases.primary.Y_next.0
\newcommand{\NUoneCasesPrimaryYNextZero}{\ensuremath{457.092}}
% NUoneCasesPrimaryYNextOne [cert, upper] u1_ghat.json :: cases.primary.Y_next.1
\newcommand{\NUoneCasesPrimaryYNextOne}{\ensuremath{356.940}}
% NUoneCasesPrimaryYNextTwo [cert, upper] u1_ghat.json :: cases.primary.Y_next.2
\newcommand{\NUoneCasesPrimaryYNextTwo}{\ensuremath{1320.220}}
% NSfourResultsMNineHTwoTightenedFUOneRowsZeroBound [cert, upper] FINAL_manifest_v33.json :: results.M9_H2_tightened_fU1.rows.0.bound
\newcommand{\NSfourResultsMNineHTwoTightenedFUOneRowsZeroBound}{\ensuremath{3.658\cdot10^{-10}}}
% NSfourResultsMNineHTwoTightenedFUOneRowsZeroMargin [cert, margin] FINAL_manifest_v33.json :: results.M9_H2_tightened_fU1.rows.0.margin
\newcommand{\NSfourResultsMNineHTwoTightenedFUOneRowsZeroMargin}{\ensuremath{27.34}}
% NSfourResultsMNineHTwoTightenedFUOneRowsOneBound [cert, upper] FINAL_manifest_v33.json :: results.M9_H2_tightened_fU1.rows.1.bound
\newcommand{\NSfourResultsMNineHTwoTightenedFUOneRowsOneBound}{\ensuremath{1.777\cdot10^{-9}}}
% NSfourResultsMNineHTwoTightenedFUOneRowsOneMargin [cert, margin] FINAL_manifest_v33.json :: results.M9_H2_tightened_fU1.rows.1.margin
\newcommand{\NSfourResultsMNineHTwoTightenedFUOneRowsOneMargin}{\ensuremath{562.8}}
% NSfourResultsMNineHTwoTightenedFUOneRowsTwoBound [cert, upper] FINAL_manifest_v33.json :: results.M9_H2_tightened_fU1.rows.2.bound
\newcommand{\NSfourResultsMNineHTwoTightenedFUOneRowsTwoBound}{\ensuremath{1.665\cdot10^{-6}}}
% NSfourResultsMNineHTwoTightenedFUOneRowsTwoMargin [cert, margin] FINAL_manifest_v33.json :: results.M9_H2_tightened_fU1.rows.2.margin
\newcommand{\NSfourResultsMNineHTwoTightenedFUOneRowsTwoMargin}{\ensuremath{600.9}}
% NHOneAlphaThr [hand, exact] hand :: 
\newcommand{\NHOneAlphaThr}{\ensuremath{1.5\cdot10^{-4}}}
% NHOneSqrtLamThr [hand, exact] hand :: 
\newcommand{\NHOneSqrtLamThr}{\ensuremath{5\cdot10^{-8}}}
% NYbarStep [hand, exact] hand :: 
\newcommand{\NYbarStep}{\ensuremath{1\cdot10^{-12}}}
% NHTwoThrZero [hand, exact] hand :: 
\newcommand{\NHTwoThrZero}{\ensuremath{1\cdot10^{-8}}}
% NHTwoThrOne [hand, exact] hand :: 
\newcommand{\NHTwoThrOne}{\ensuremath{1\cdot10^{-6}}}
% NHTwoThrTwo [hand, exact] hand :: 
\newcommand{\NHTwoThrTwo}{\ensuremath{1\cdot10^{-3}}}
% NHThreeThrZero [hand, exact] hand :: 
\newcommand{\NHThreeThrZero}{\ensuremath{1\cdot10^{-6}}}
% NHThreeThrOne [hand, exact] hand :: 
\newcommand{\NHThreeThrOne}{\ensuremath{5\cdot10^{-6}}}
% NHThreeThrTwo [hand, exact] hand :: 
\newcommand{\NHThreeThrTwo}{\ensuremath{1\cdot10^{-5}}}
% NHFourThrZero [hand, exact] hand :: 
\newcommand{\NHFourThrZero}{\ensuremath{2\cdot10^{-7}}}
% NHFourThrOne [hand, exact] hand :: 
\newcommand{\NHFourThrOne}{\ensuremath{1\cdot10^{-5}}}
% NHFourThrTwo [hand, exact] hand :: 
\newcommand{\NHFourThrTwo}{\ensuremath{1\cdot10^{-5}}}
% NHFourThrThree [hand, exact] hand :: 
\newcommand{\NHFourThrThree}{\ensuremath{1\cdot10^{-5}}}
% NHFourKfourThr [hand, exact] hand :: 
\newcommand{\NHFourKfourThr}{\ensuremath{0.06}}
% NHSixConeSlopeThr [hand, exact] hand :: 
\newcommand{\NHSixConeSlopeThr}{\ensuremath{0.08}}
% NHSixConeKtwoThr [hand, exact] hand :: 
\newcommand{\NHSixConeKtwoThr}{\ensuremath{0.03}}
% NHSixConeNthreeThr [hand, exact] hand :: 
\newcommand{\NHSixConeNthreeThr}{\ensuremath{0.004}}
% NHSixConeJetThr [hand, exact] hand :: 
\newcommand{\NHSixConeJetThr}{\ensuremath{0.12}}
% NHSixConeHFiveThr [hand, exact] hand :: 
\newcommand{\NHSixConeHFiveThr}{\ensuremath{0.0144}}
% NHSixFarHFiveThr [hand, exact] hand :: 
\newcommand{\NHSixFarHFiveThr}{\ensuremath{0.0225}}
% NLemmaOConeM [hand, exact] hand :: 
\newcommand{\NLemmaOConeM}{\ensuremath{0.12}}
% NLemmaOConeK [hand, exact] hand :: 
\newcommand{\NLemmaOConeK}{\ensuremath{0.18}}
% NLemmaOFarM [hand, exact] hand :: 
\newcommand{\NLemmaOFarM}{\ensuremath{0.15}}
% NLemmaOFarK [hand, exact] hand :: 
\newcommand{\NLemmaOFarK}{\ensuremath{0.3}}
% NCoercivityConst [hand, exact] hand :: 
\newcommand{\NCoercivityConst}{\ensuremath{0.046875}}
% NMOnez [hand, exact] hand :: 
\newcommand{\NMOnez}{\ensuremath{0.534}}
% NBetaA [hand, exact] hand :: 
\newcommand{\NBetaA}{\ensuremath{2}}
% NTOneLower [cert, lower] t1clock.json :: t1.0
\newcommand{\NTOneLower}{\ensuremath{4.5978574775}}
% NTOneUpper [cert, upper] t1clock.json :: t1.1
\newcommand{\NTOneUpper}{\ensuremath{4.5978586144}}
% NLambdaOneLower [cert, lower] t1clock.json :: lnlam.0
\newcommand{\NLambdaOneLower}{\ensuremath{6.49\cdot10^{-5}}}
% NLambdaOneUpper [cert, upper] t1clock.json :: lnlam.1
\newcommand{\NLambdaOneUpper}{\ensuremath{6.75\cdot10^{-5}}}
% NTMinusTOneUpper [cert, upper] t1clock.json :: lnlam.1; 1.5e-4.sup_ab
\newcommand{\NTMinusTOneUpper}{\ensuremath{2.1\cdot10^{-17}}}
% NAbarUpper [cert, upper] final_alpha_points_5e-8.json :: 1.5e-4.sup_ab
\newcommand{\NAbarUpper}{\ensuremath{2.19103\cdot10^{-7}}}
% NBetabarUpper [cert, upper] final_alpha_points_5e-8.json :: 1.5e-4.sup_beta
\newcommand{\NBetabarUpper}{\ensuremath{3.33811\cdot10^{-9}}}
% NBaseKmaxLower [cert, lower] t1clock.json :: base_kmax.0
\newcommand{\NBaseKmaxLower}{\ensuremath{2.279\cdot10^{4}}}
% NBaseKmaxUpper [cert, upper] t1clock.json :: base_kmax.1
\newcommand{\NBaseKmaxUpper}{\ensuremath{2.376\cdot10^{4}}}
% NHOneAlphaBase [cert, upper] summary.json :: alpha_abs_max
\newcommand{\NHOneAlphaBase}{\ensuremath{1.1739\cdot10^{-4}}}
% NHOneAlphaErr [cert, upper] t9_h1_clam_s180.json :: C_alpha; inputs.YSTAR
\newcommand{\NHOneAlphaErr}{\ensuremath{7.21\cdot10^{-13}}}
% NHOneAlphaMargin [cert, margin] summary.json :: alpha_abs_max; C_alpha; inputs.YSTAR
\newcommand{\NHOneAlphaMargin}{\ensuremath{1.27}}
% NHOneSqrtLamBase [cert, upper] t9_h1_clam_s180.json :: sqrtLam_base_Wp
\newcommand{\NHOneSqrtLamBase}{\ensuremath{2.282\cdot10^{-8}}}
% NHOneSqrtLamTotal [cert, upper] t9_h1_clam_s180.json :: sqrtLam_total_Wp
\newcommand{\NHOneSqrtLamTotal}{\ensuremath{2.29\cdot10^{-8}}}
% NHOneSqrtLamMargin [cert, margin] t9_h1_clam_s180.json :: sqrtLam_total_Wp
\newcommand{\NHOneSqrtLamMargin}{\ensuremath{2.18}}
% NHOneYbarLimit [cert, lower] t9_h1_clam_s180.json :: Ybar_limit_H1
\newcommand{\NHOneYbarLimit}{\ensuremath{1.12\cdot10^{-9}}}
% NSwitchYUpper [cert, upper] switch_s2.json :: modulated.Y_raw
\newcommand{\NSwitchYUpper}{\ensuremath{6.21\cdot10^{-13}}}
% NSwitchYMargin [cert, margin] switch_s2.json :: modulated.Y_raw
\newcommand{\NSwitchYMargin}{\ensuremath{1.6}}
% NMinCI [cert, lower] constants.json :: meta.summary.min_c_I
\newcommand{\NMinCI}{\ensuremath{0.05283}}
% NNumCertificates [cert, exact] constants.json :: meta.summary.n_cert
\newcommand{\NNumCertificates}{\ensuremath{26}}
% NMinRhoC [cert, lower] constants.json :: meta.closure.rho_c_min
\newcommand{\NMinRhoC}{\ensuremath{0.04711}}
% NHTwoBaseZero [cert, upper] FINAL_manifest_v33.json :: results.M9_H2_tightened.rows.0.base
\newcommand{\NHTwoBaseZero}{\ensuremath{8.011\cdot10^{-11}}}
% NHTwoBaseOne [cert, upper] FINAL_manifest_v33.json :: results.M9_H2_tightened.rows.1.base
\newcommand{\NHTwoBaseOne}{\ensuremath{1.336\cdot10^{-10}}}
% NHTwoBaseTwo [cert, upper] FINAL_manifest_v33.json :: results.M9_H2_tightened.rows.2.base
\newcommand{\NHTwoBaseTwo}{\ensuremath{2.936\cdot10^{-10}}}
% NHTwoErrWZero [cert, upper] FINAL_manifest_v33.json :: results.M9_H2_tightened.rows.0.with_error
\newcommand{\NHTwoErrWZero}{\ensuremath{8.014\cdot10^{-11}}}
% NHTwoMarginWZero [cert, margin] FINAL_manifest_v33.json :: results.M9_H2_tightened.rows.0.margin
\newcommand{\NHTwoMarginWZero}{\ensuremath{124.78}}
% NHTwoErrWOne [cert, upper] FINAL_manifest_v33.json :: results.M9_H2_tightened.rows.1.with_error
\newcommand{\NHTwoErrWOne}{\ensuremath{4.706\cdot10^{-10}}}
% NHTwoMarginWOne [cert, margin] FINAL_manifest_v33.json :: results.M9_H2_tightened.rows.1.margin
\newcommand{\NHTwoMarginWOne}{\ensuremath{2124.9}}
% NHTwoErrWTwo [cert, upper] FINAL_manifest_v33.json :: results.M9_H2_tightened.rows.2.with_error
\newcommand{\NHTwoErrWTwo}{\ensuremath{1.559\cdot10^{-6}}}
% NHTwoMarginWTwo [cert, margin] FINAL_manifest_v33.json :: results.M9_H2_tightened.rows.2.margin
\newcommand{\NHTwoMarginWTwo}{\ensuremath{641.53}}
% NHThreeBaseZero [cert, upper] FINAL_manifest_v33.json :: results.M9_H2f.rows.0.base
\newcommand{\NHThreeBaseZero}{\ensuremath{3.127\cdot10^{-7}}}
% NHThreeBaseOne [cert, upper] FINAL_manifest_v33.json :: results.M9_H2f.rows.1.base
\newcommand{\NHThreeBaseOne}{\ensuremath{1.004\cdot10^{-6}}}
% NHThreeBaseTwo [cert, upper] FINAL_manifest_v33.json :: results.M9_H2f.rows.2.base
\newcommand{\NHThreeBaseTwo}{\ensuremath{2.431\cdot10^{-6}}}
% NHThreeErrZero [cert, upper] FINAL_manifest_v33.json :: results.M9_H2f.rows.0.with_error
\newcommand{\NHThreeErrZero}{\ensuremath{3.1363\cdot10^{-7}}}
% NHThreeMarginZero [cert, margin] FINAL_manifest_v33.json :: results.M9_H2f.rows.0.margin
\newcommand{\NHThreeMarginZero}{\ensuremath{3.1885}}
% NHThreeErrOne [cert, upper] FINAL_manifest_v33.json :: results.M9_H2f.rows.1.with_error
\newcommand{\NHThreeErrOne}{\ensuremath{1.0034\cdot10^{-6}}}
% NHThreeMarginOne [cert, margin] FINAL_manifest_v33.json :: results.M9_H2f.rows.1.margin
\newcommand{\NHThreeMarginOne}{\ensuremath{4.9834}}
% NHThreeErrTwo [cert, upper] FINAL_manifest_v33.json :: results.M9_H2f.rows.2.with_error
\newcommand{\NHThreeErrTwo}{\ensuremath{2.4351\cdot10^{-6}}}
% NHThreeMarginTwo [cert, margin] FINAL_manifest_v33.json :: results.M9_H2f.rows.2.margin
\newcommand{\NHThreeMarginTwo}{\ensuremath{4.1067}}
% NHFourBaseZero [cert, upper] FINAL_manifest_v33.json :: results.M9_base_sums.combined.windows_max.0
\newcommand{\NHFourBaseZero}{\ensuremath{3.48\cdot10^{-8}}}
% NHFourErrZero [cert, upper] FINAL_manifest_v33.json :: results.M9_H3_H2v32.rows.win.0.with_error_arb
\newcommand{\NHFourErrZero}{\ensuremath{3.5087\cdot10^{-8}}}
% NHFourMarginZero [cert, margin] FINAL_manifest_v33.json :: results.M9_H3_H2v32.rows.win.0.margin_arb
\newcommand{\NHFourMarginZero}{\ensuremath{5.700}}
% NHFourBaseOne [cert, upper] FINAL_manifest_v33.json :: results.M9_base_sums.combined.windows_max.1
\newcommand{\NHFourBaseOne}{\ensuremath{1.27\cdot10^{-7}}}
% NHFourErrOne [cert, upper] FINAL_manifest_v33.json :: results.M9_H3_H2v32.rows.win.1.with_error_arb
\newcommand{\NHFourErrOne}{\ensuremath{1.2616\cdot10^{-7}}}
% NHFourMarginOne [cert, margin] FINAL_manifest_v33.json :: results.M9_H3_H2v32.rows.win.1.margin_arb
\newcommand{\NHFourMarginOne}{\ensuremath{79.26}}
% NHFourBaseTwo [cert, upper] FINAL_manifest_v33.json :: results.M9_base_sums.combined.windows_max.2
\newcommand{\NHFourBaseTwo}{\ensuremath{3.76\cdot10^{-7}}}
% NHFourErrTwo [cert, upper] FINAL_manifest_v33.json :: results.M9_H3_H2v32.rows.win.2.with_error_arb
\newcommand{\NHFourErrTwo}{\ensuremath{3.7724\cdot10^{-7}}}
% NHFourMarginTwo [cert, margin] FINAL_manifest_v33.json :: results.M9_H3_H2v32.rows.win.2.margin_arb
\newcommand{\NHFourMarginTwo}{\ensuremath{26.50}}
% NHFourBaseThree [cert, upper] FINAL_manifest_v33.json :: results.M9_base_sums.combined.windows_max.3
\newcommand{\NHFourBaseThree}{\ensuremath{9.71\cdot10^{-7}}}
% NHFourErrThree [cert, upper] FINAL_manifest_v33.json :: results.M9_H3_H2v32.rows.win.3.with_error_arb
\newcommand{\NHFourErrThree}{\ensuremath{1.2822\cdot10^{-6}}}
% NHFourMarginThree [cert, margin] FINAL_manifest_v33.json :: results.M9_H3_H2v32.rows.win.3.margin_arb
\newcommand{\NHFourMarginThree}{\ensuremath{7.799}}
% NHFourYbarLimit [cert, lower] FINAL_manifest_v33.json :: results.M9_H3_H2v32.binding_Ybar_limit
\newcommand{\NHFourYbarLimit}{\ensuremath{1.575\cdot10^{-11}}}
% NHFourYbarMargin [cert, margin] FINAL_manifest_v33.json :: results.M9_H3_H2v32.binding_Ybar_limit
\newcommand{\NHFourYbarMargin}{\ensuremath{15.75}}
% NHFourKfourBase [cert, upper] amp_A1_w180.json :: inputs.BASE_T1.K4
\newcommand{\NHFourKfourBase}{\ensuremath{3.1633\cdot10^{-4}}}
% NHSixConeSlope [cert, upper] FINAL_manifest_v33.json :: results.cone_fixedL.checks.slope.value
\newcommand{\NHSixConeSlope}{\ensuremath{0.03029}}
% NHSixConeSlopeMargin [cert, margin] FINAL_manifest_v33.json :: results.cone_fixedL.checks.slope.margin
\newcommand{\NHSixConeSlopeMargin}{\ensuremath{2.641}}
% NHSixConeKtwo [cert, upper] FINAL_manifest_v33.json :: results.cone_fixedL.checks.K2.value
\newcommand{\NHSixConeKtwo}{\ensuremath{0.009876}}
% NHSixConeKtwoMargin [cert, margin] FINAL_manifest_v33.json :: results.cone_fixedL.checks.K2.margin
\newcommand{\NHSixConeKtwoMargin}{\ensuremath{3.037}}
% NHSixConeNthree [cert, upper] FINAL_manifest_v33.json :: results.cone_fixedL.checks.N3.value
\newcommand{\NHSixConeNthree}{\ensuremath{1.461\cdot10^{-4}}}
% NHSixConeNthreeMargin [cert, margin] FINAL_manifest_v33.json :: results.cone_fixedL.checks.N3.margin
\newcommand{\NHSixConeNthreeMargin}{\ensuremath{27.38}}
% NHSixConeHFive [cert, upper] FINAL_manifest_v33.json :: results.H4_fixedL.total
\newcommand{\NHSixConeHFive}{\ensuremath{1.2865\cdot10^{-5}}}
% NHSixConeHFiveMargin [cert, margin] FINAL_manifest_v33.json :: results.H4_fixedL.margin
\newcommand{\NHSixConeHFiveMargin}{\ensuremath{1119}}
% NHSixConeTstarNeeded [cert, upper] FINAL_manifest_v33.json :: results.H7_fixedL.ratio_upper
\newcommand{\NHSixConeTstarNeeded}{\ensuremath{1.6000015\cdot10^{-7}}}
% NHSixConeTstarLower [cert, lower] FINAL_manifest_v33.json :: results.H7_fixedL.t_star_lower
\newcommand{\NHSixConeTstarLower}{\ensuremath{3.68334\cdot10^{-7}}}
% NHSixConeTstarMargin [cert, margin] FINAL_manifest_v33.json :: results.H7_fixedL.margin_tstar_over_ratio
\newcommand{\NHSixConeTstarMargin}{\ensuremath{2.302}}
% NHSixConeHFiveRatioPeak [cert, upper] h4loc_A1_w180.json :: windows.constants/amp_A1_w180.json.worst.ratio
\newcommand{\NHSixConeHFiveRatioPeak}{\ensuremath{0.2393}}
% NHSixConeHOneSup [cert, upper] FINAL_manifest_v33.json :: results.H1_fixedL.sup_h0_max
\newcommand{\NHSixConeHOneSup}{\ensuremath{0.2163}}
% NHSixConeHOneH [cert, upper] FINAL_manifest_v33.json :: results.H1_fixedL.H_max
\newcommand{\NHSixConeHOneH}{\ensuremath{0.2185}}
% NHSixFarHFiveRatio [cert, upper] h4loc_v32.json :: windows.constants/amp_A1_w180.json.cylinders.18.ratio
\newcommand{\NHSixFarHFiveRatio}{\ensuremath{0.1761}}
% NHSixFarHFiveMarginOld [cert, margin] h4loc_v32.json :: windows.constants/amp_A1_w180.json.cylinders.18.ratio
\newcommand{\NHSixFarHFiveMarginOld}{\ensuremath{5.68}}
% NHSixFarHFiveMargin [cert, margin] FINAL_manifest_v33.json :: results.M9_H4_far.min_margin
\newcommand{\NHSixFarHFiveMargin}{\ensuremath{15.58}}
% NHSixConeHFiveMarginOld [cert, margin] h4loc_A1_w180.json :: windows.constants/amp_A1_w180.json.worst.ratio
\newcommand{\NHSixConeHFiveMarginOld}{\ensuremath{4.18}}
% NHSixFarKtwo [cert, upper] t9_far_v32.json :: h2
\newcommand{\NHSixFarKtwo}{\ensuremath{1.797\cdot10^{-4}}}
% NHSixFarKtwoThr [hand, exact] t9_far_v32.json :: thr.h2
\newcommand{\NHSixFarKtwoThr}{\ensuremath{3.53\cdot10^{-4}}}
% NHSixFarKtwoMargin [cert, margin] t9_far_v32.json :: margins.h2
\newcommand{\NHSixFarKtwoMargin}{\ensuremath{1.964}}
% NHSixFarKthree [cert, upper] t9_far_v32.json :: h3
\newcommand{\NHSixFarKthree}{\ensuremath{2.595\cdot10^{-8}}}
% NHSixFarKthreeThr [hand, exact] t9_far_v32.json :: thr.h3
\newcommand{\NHSixFarKthreeThr}{\ensuremath{1.17\cdot10^{-5}}}
% NHSixFarKthreeMargin [cert, margin] t9_far_v32.json :: margins.h3
\newcommand{\NHSixFarKthreeMargin}{\ensuremath{450.9}}
% NHSixFarKfour [cert, upper] t9_far_v32.json :: h4
\newcommand{\NHSixFarKfour}{\ensuremath{6.964\cdot10^{-9}}}
% NHSixFarKfourThr [hand, exact] t9_far_v32.json :: thr.h4
\newcommand{\NHSixFarKfourThr}{\ensuremath{1.03\cdot10^{-7}}}
% NHSixFarKfourMargin [cert, margin] t9_far_v32.json :: margins.h4
\newcommand{\NHSixFarKfourMargin}{\ensuremath{14.79}}
% NHSixFarSlope [cert, upper] t9_far_v32.json :: slope
\newcommand{\NHSixFarSlope}{\ensuremath{0.04153}}
% NHSixFarSlopeThr [hand, exact] t9_far_v32.json :: thr.slope
\newcommand{\NHSixFarSlopeThr}{\ensuremath{0.0800}}
% NHSixFarSlopeMargin [cert, margin] t9_far_v32.json :: margins.slope
\newcommand{\NHSixFarSlopeMargin}{\ensuremath{1.926}}
% NHSixFarNthree [cert, upper] t9_far_v32.json :: L2
\newcommand{\NHSixFarNthree}{\ensuremath{5.516\cdot10^{-7}}}
% NHSixFarNthreeThr [hand, exact] t9_far_v32.json :: thr.L2
\newcommand{\NHSixFarNthreeThr}{\ensuremath{3.32\cdot10^{-6}}}
% NHSixFarNthreeMargin [cert, margin] t9_far_v32.json :: margins.L2
\newcommand{\NHSixFarNthreeMargin}{\ensuremath{6.019}}
% NFarTstarMargin [cert, margin] lemmaO_v32_consts.json :: far_v32.case.margin_factor_lower
\newcommand{\NFarTstarMargin}{\ensuremath{15.5}}
% NFarTstarLower [cert, lower] lemmaO_v32_consts.json :: far_v32.case.t_star_lower
\newcommand{\NFarTstarLower}{\ensuremath{9.5108\cdot10^{-8}}}
% NFarTstarNeeded [cert, upper] lemmaO_v32_consts.json :: far_v32.case.needed
\newcommand{\NFarTstarNeeded}{\ensuremath{6.134\cdot10^{-9}}}
% NGammaModes [cert, exact] gamma0_final.json :: M
\newcommand{\NGammaModes}{\ensuremath{3762}}
% NGammaHalfwidth [cert, upper] gamma0_final.json :: b_j_halfwidth
\newcommand{\NGammaHalfwidth}{\ensuremath{8.138\cdot10^{-24}}}
% NGammaContraction [cert, upper] gamma0_final.json :: krawczyk_contraction
\newcommand{\NGammaContraction}{\ensuremath{3.01\cdot10^{-12}}}
% NGammaArcs [cert, exact] gamma0_final.json :: embedded.arcs
\newcommand{\NGammaArcs}{\ensuremath{701}}
% NGammaMaxTurn [cert, upper] gamma0_final.json :: embedded.maxturn
\newcommand{\NGammaMaxTurn}{\ensuremath{0.04997}}
% NGammaMaxPair [cert, upper] gamma0_final.json :: embedded.maxpair
\newcommand{\NGammaMaxPair}{\ensuremath{0.09597}}
% NGammaMinSep [cert, lower] gamma0_final.json :: embedded.minsep
\newcommand{\NGammaMinSep}{\ensuremath{8.89\cdot10^{-3}}}
% NGammaHalfGap [cert, lower] gamma0_final.json :: embedded.gap_cert
\newcommand{\NGammaHalfGap}{\ensuremath{0.0248219}}
% NPhaseISlabs [cert, exact] status.json :: slab
\newcommand{\NPhaseISlabs}{\ensuremath{96}}
% NPhaseIModes [cert, exact] status.json :: last_slab.M
\newcommand{\NPhaseIModes}{\ensuremath{2800}}
% NPhaseIKmaxLower [cert, lower] status.json :: kmax_enclosure.0
\newcommand{\NPhaseIKmaxLower}{\ensuremath{29.352874545}}
% NPhaseIKmaxUpper [cert, upper] status.json :: kmax_enclosure.1
\newcommand{\NPhaseIKmaxUpper}{\ensuremath{29.352874546}}
% NPhaseITLower [cert, lower] status.json :: t
\newcommand{\NPhaseITLower}{\ensuremath{4.597850539485830}}
% NPhaseITUpper [cert, upper] status.json :: t
\newcommand{\NPhaseITUpper}{\ensuremath{4.597850539485831}}
% NPeakSqrtLam [cert, upper] rate_table_fin12f4_a0.00015.json :: sqrtL_max
\newcommand{\NPeakSqrtLam}{\ensuremath{5.308\cdot10^{-8}}}
% NTailRate [cert, lower] rate_table_fin12f4_a0.00015.json :: tail.rho
\newcommand{\NTailRate}{\ensuremath{2.693\cdot10^{-3}}}
% NFloorSqrtLam [cert, upper] rate_table_fin12f4_a0.00015.json :: tail.sqrt_LF
\newcommand{\NFloorSqrtLam}{\ensuremath{9.872\cdot10^{-10}}}
% NAlphaFloor [cert, upper] rate_table_fin12f4_a0.00015.json :: tail.alpha_F
\newcommand{\NAlphaFloor}{\ensuremath{6.867\cdot10^{-7}}}
% NSupBetaVerifier [cert, upper] final_alpha_points_5e-8.json :: 1.5e-4.sup_beta
\newcommand{\NSupBetaVerifier}{\ensuremath{3.33811\cdot10^{-9}}}
% NMOneRamp [cert, upper] final_alpha_points_5e-8.json :: 1.5e-4.M1
\newcommand{\NMOneRamp}{\ensuremath{0.4328}}
% NSITwoPrime [cert, upper] FINAL_manifest_v33.json :: results.F12_sliver.sI2_prime
\newcommand{\NSITwoPrime}{\ensuremath{0.015457}}
% NLoopGainCoarse [cert, upper] U_cert_v35.json :: loop_gain.2.total
\newcommand{\NLoopGainCoarse}{\ensuremath{0.03904}}
% NRateMarginCoarse [cert, margin] U_cert_v35.json :: rate_margin
\newcommand{\NRateMarginCoarse}{\ensuremath{28.2}}
% NSOneEpsLower [cert, lower] bracket2.json :: eps_a
\newcommand{\NSOneEpsLower}{\ensuremath{1.5346198892965\cdot10^{-4}}}
% NSOneEpsUpper [cert, upper] bracket2.json :: eps_b
\newcommand{\NSOneEpsUpper}{\ensuremath{1.5346198899299\cdot10^{-4}}}
% NSOneRowsTip [cert, exact] ext_tip_check.json :: rows
\newcommand{\NSOneRowsTip}{\ensuremath{51246}}
% NSOneTipWidth [cert, upper] ext_tip_check.json :: max_pointwise_width.th
\newcommand{\NSOneTipWidth}{\ensuremath{3.2\cdot10^{-12}}}
% NSOneExtWidth [cert, upper] ext813_check.json :: max_pointwise_width.p
\newcommand{\NSOneExtWidth}{\ensuremath{5.7\cdot10^{-10}}}
% NSOneRowsMC [cert, exact] check_outA5.json :: prof_tip_vs_MC.rows_compared
\newcommand{\NSOneRowsMC}{\ensuremath{1025}}
% NSOneRowsMB [cert, exact] check_outA5.json :: ext813_vs_MB.rows_compared
\newcommand{\NSOneRowsMB}{\ensuremath{501}}
% NSOneMinMargin [cert, margin] check_outA5.json :: min:*.min_relative_margin.*
\newcommand{\NSOneMinMargin}{\ensuremath{1.2\cdot10^{-9}}}
% NTwoSegments [cert, exact] contour_result.json :: len:segments
\newcommand{\NTwoSegments}{\ensuremath{5058}}
% NTwoInteriorBoxes [cert, exact] interior_result.json :: len:
\newcommand{\NTwoInteriorBoxes}{\ensuremath{174}}
% NTwoContourBoxes [cert, exact] eta_cache.json :: len:
\newcommand{\NTwoContourBoxes}{\ensuremath{202}}
% NTwoEtaMax [cert, upper] eta_cache.json :: max:*.eta
\newcommand{\NTwoEtaMax}{\ensuremath{0.1163}}
% NTwoLipMax [cert, upper] eta_cache.json :: max:*.Lip
\newcommand{\NTwoLipMax}{\ensuremath{0.272}}
% NTwoLambdaZero [est, upper] cert_b0.25_th0.0_tau0.03.json :: Lam
\newcommand{\NTwoLambdaZero}{\ensuremath{1.790819}}
% NTwoLambdaOne [est, upper] cert_b0.25_th0.48_tau0.03.json :: Lam
\newcommand{\NTwoLambdaOne}{\ensuremath{1.768282}}
% NSOneDTightLower [doc, lower] doc:README.md :: D ∈ [2.50080627654, 2.50080627668]
\newcommand{\NSOneDTightLower}{\ensuremath{2.50080627654}}
% NSOneDTightUpper [doc, upper] doc:README.md :: D ∈ [2.50080627654, 2.50080627668]
\newcommand{\NSOneDTightUpper}{\ensuremath{2.50080627668}}
% NSOneKTightLower [doc, lower] doc:README.md :: k(0) ∈ [1.53982063762, 1.53982063822]
\newcommand{\NSOneKTightLower}{\ensuremath{1.53982063762}}
% NSOneKTightUpper [doc, upper] doc:README.md :: k(0) ∈ [1.53982063762, 1.53982063822]
\newcommand{\NSOneKTightUpper}{\ensuremath{1.53982063822}}
% NSOneDCondLower [doc, lower] doc:README.md :: D = −p(tip) ∈ [2.500806276619695, 2.500806276620049]
\newcommand{\NSOneDCondLower}{\ensuremath{2.500806276619695}}
% NSOneDCondUpper [doc, upper] doc:README.md :: D = −p(tip) ∈ [2.500806276619695, 2.500806276620049]
\newcommand{\NSOneDCondUpper}{\ensuremath{2.500806276620049}}
% NSOneKCondLower [doc, lower] doc:README.md :: k(0) ∈ [1.539820637871836, 1.5398206378721015]
\newcommand{\NSOneKCondLower}{\ensuremath{1.539820637871836}}
% NSOneKCondUpper [doc, upper] doc:README.md :: k(0) ∈ [1.539820637871836, 1.5398206378721015]
\newcommand{\NSOneKCondUpper}{\ensuremath{1.5398206378721015}}
% NSOneTauTipLower [doc, lower] doc:README.md :: τ_tip ∈ [10.249077651120778, 10.249077651121361]
\newcommand{\NSOneTauTipLower}{\ensuremath{10.249077651120778}}
% NSOneTauTipUpper [doc, upper] doc:README.md :: τ_tip ∈ [10.249077651120778, 10.249077651121361]
\newcommand{\NSOneTauTipUpper}{\ensuremath{10.249077651121361}}
% NSOneStateUncert [doc, upper] doc:README.md :: State uncertainty at σ = 8: 4.4e−13
\newcommand{\NSOneStateUncert}{\ensuremath{4.4\cdot10^{-13}}}
% NFarBoxEprime [doc, margin] doc:STEP4_RESULT.md :: (6.4287, 4.6917, 1.9940, 1.69999995)
\newcommand{\NFarBoxEprime}{\ensuremath{6.42}}
% NFarBoxXEsecond [doc, margin] doc:STEP4_RESULT.md :: (6.4287, 4.6917, 1.9940, 1.69999995)
\newcommand{\NFarBoxXEsecond}{\ensuremath{4.69}}
% NFarBoxXsqEthird [doc, margin] doc:STEP4_RESULT.md :: (6.4287, 4.6917, 1.9940, 1.69999995)
\newcommand{\NFarBoxXsqEthird}{\ensuremath{1.99}}
% NFarBoxE [doc, margin] doc:STEP4_RESULT.md :: (6.4287, 4.6917, 1.9940, 1.69999995)
\newcommand{\NFarBoxE}{\ensuremath{1.69}}
% NHFourKfourErr [doc, upper] doc:STEP4_RESULT.md :: | 5.60e-3 against 0.06 | ×10.7 |
\newcommand{\NHFourKfourErr}{\ensuremath{5.60\cdot10^{-3}}}
% NHFourKfourMargin [doc, margin] doc:STEP4_RESULT.md :: | 5.60e-3 against 0.06 | ×10.7 |
\newcommand{\NHFourKfourMargin}{\ensuremath{10.7}}
% NCoercivityCoefficient [doc, upper] doc:step3_proof_v32.tex :: \norm{e^{\tau\mathcal L}u}_\rho\le14.8\,e^{-0.0213\tau}\norm
\newcommand{\NCoercivityCoefficient}{\ensuremath{14.8}}
% NCoercivityRate [doc, lower] doc:step3_proof_v32.tex :: \norm{e^{\tau\mathcal L}u}_\rho\le14.8\,e^{-0.0213\tau}\norm
\newcommand{\NCoercivityRate}{\ensuremath{0.0213}}
% NRegionRampStart [hand, exact] hand :: 
\newcommand{\NRegionRampStart}{\ensuremath{14}}
% NRegionInterface [hand, exact] hand :: 
\newcommand{\NRegionInterface}{\ensuremath{28}}
% NRegionShellEnd [hand, exact] hand :: 
\newcommand{\NRegionShellEnd}{\ensuremath{60}}
% NRegionSegEnd [hand, exact] hand :: 
\newcommand{\NRegionSegEnd}{\ensuremath{350}}
% NRegionAnnEnd [hand, exact] hand :: 
\newcommand{\NRegionAnnEnd}{\ensuremath{700}}
% NZoneIIILow [hand, exact] hand :: 
\newcommand{\NZoneIIILow}{\ensuremath{253}}
% NZoneIIIHigh [hand, exact] hand :: 
\newcommand{\NZoneIIIHigh}{\ensuremath{379.5}}
% NConeCentreMin [hand, exact] hand :: 
\newcommand{\NConeCentreMin}{\ensuremath{300}}
% NConeCentreMax [hand, exact] hand :: 
\newcommand{\NConeCentreMax}{\ensuremath{380}}
% NFarCentreMin [hand, exact] hand :: 
\newcommand{\NFarCentreMin}{\ensuremath{490}}
% NConeL [hand, exact] hand :: 
\newcommand{\NConeL}{\ensuremath{50}}
% NConeX [hand, exact] hand :: 
\newcommand{\NConeX}{\ensuremath{3}}
% NFarL [hand, exact] hand :: 
\newcommand{\NFarL}{\ensuremath{113}}
% NFarX [hand, exact] hand :: 
\newcommand{\NFarX}{\ensuremath{2}}
}%
  {\PackageWarning{paper}{numbers/numbers_paper.tex not found (unit U0 not finished)}}

% ======================================================================
% NOTATION MACROS (OUTLINE 5.2 and notation.md).  One macro per symbol
% whose spelling is fixed by the notation table.  Plain letters (a, b, f, ...)
% have no macro.
% ======================================================================
% basic
\newcommand{\R}{\mathbb{R}}
\newcommand{\C}{\mathbb{C}}
\newcommand{\N}{\mathbb{N}}
\newcommand{\Z}{\mathbb{Z}}
\newcommand{\eps}{\varepsilon}
\newcommand{\ip}[2]{\langle #1,#2\rangle}
\newcommand{\abs}[1]{\lvert #1\rvert}
\newcommand{\norm}[1]{\lVert #1\rVert}
\newcommand{\dd}{\mathrm d}
\newcommand{\re}{\operatorname{Re}}
\newcommand{\im}{\operatorname{Im}}
% flow and curves
\newcommand{\arclen}{\mathfrak{s}}                 % arclength of the flow curve
\newcommand{\kss}{\kappa_{\mathfrak{ss}}}          % kappa_{ss} in the flow equation
\newcommand{\Area}{\operatorname{Area}}
\newcommand{\Len}{\mathcal L}                      % length of a curve
\newcommand{\spp}{\operatorname{sp}}                 % sup quantities sp^v_0, sp^v_1, sp^f_1 (were S^v_0, S^v_1, S^f_1)
\newcommand{\rot}{\operatorname{rot}}              % turning number
\newcommand{\Tsing}{T_*}                           % singular time
\newcommand{\Tmax}{T_{\max}}                       % maximal time of existence
\newcommand{\pstar}{p_*}                           % image point of the singularity
\newcommand{\Gz}{\Gamma_0}                         % initial curve
\newcommand{\vth}{\vartheta}                       % parameter of Gamma_0
\newcommand{\dclk}{\delta_t}                       % clock defect of step 4
\newcommand{\fall}{\varpi}                         % Fourier rates varpi_n = n^4-n^2
% rescaling
\newcommand{\lamone}{\lambda_1}
\newcommand{\rate}{\varsigma}                      % 1/4 - a
\newcommand{\xia}{\xi_a}                           % (1-4\bar a)^{-1}
% shrinker and step 1
\newcommand{\Sone}{S_1}
\newcommand{\Gsh}{\Gamma_{\!*}}                    % profile curve of S_1
\newcommand{\cSet}{\mathcal S}                     % certified set of shrinkers
\newcommand{\sigtip}{\sigma_{\rm tip}}
\newcommand{\Xo}{X_0}                              % 2.2490776515
\newcommand{\Dtip}{D}
\newcommand{\ktip}{k_0}
\newcommand{\shootM}{M}                            % shooting function
\newcommand{\Ypar}{\mathcal Y}                     % step-1 parallelepiped
\newcommand{\axis}{\Sigma}                         % symmetry axis of S_1
\newcommand{\sigc}{\sigma_{\rm c}}
\newcommand{\sigrev}{\sigma^{\rm rev}}
% linear theory
\newcommand{\Hrho}{\mathcal H_\rho}
\newcommand{\muslow}{\mu_3}
\newcommand{\rre}{\mathfrak r_{\rm re}}
\newcommand{\rim}{\mathfrak r_{\rm im}}
\newcommand{\Evev}{\mathrm{Ev}_{\rm even}}
\newcommand{\Evod}{\mathrm{Ev}_{\rm odd}}
\newcommand{\gvec}{\mathbf g}                      % translation vector in step 2
\newcommand{\cons}[1]{\mathfrak f_{#1}}            % constraint functions
\newcommand{\phirho}{\varphi_\rho}
% energies
\newcommand{\Lam}{\Lambda}
\newcommand{\En}{E}
\newcommand{\EnI}{E_I}
\newcommand{\muI}{\mu_I}                           % 0.08
\newcommand{\epsE}{\varepsilon_E}                  % 1e-5
\newcommand{\epsI}{\varepsilon_I}                  % 3e-5
\newcommand{\wn}[1]{\mathsf w_{#1}}                % ||w~^{(j)}||_rho
\newcommand{\vI}{v_I}                              % zeta w'
\newcommand{\muerr}{\mu_{\rm err}}                 % step-4 error-energy weight 0.036
% far field
\newcommand{\pU}{p_U}                              % 1.0246
\newcommand{\Uone}{U_1}
\newcommand{\sm}{\operatorname{sm}}                % smoothstep t^3(10-15t+6t^2)
\newcommand{\rc}{r_c}                              % centre radius (Lemma U)
\newcommand{\rchi}{r_\chi}                         % 253
\newcommand{\Iw}{I_w}                              % hand-off windows [y_w,1.5y_w]
\newcommand{\qa}{\mathfrak q_1}                    % quasilinear coefficients of the slope equation
\newcommand{\qb}{\mathfrak q_2}
% cylinders (Lemma O)
\newcommand{\cyl}{\ell}                            % cylinder length unit
\newcommand{\MO}{M_O}
\newcommand{\KO}{K_O}
\newcommand{\Hcyl}{H_{\rm cyl}}
\newcommand{\XiO}{\Xi_O}
\newcommand{\fO}{f_O}
% step 4
\newcommand{\Ybar}{\bar Y}
\newcommand{\Lt}{\widetilde L}                     % linearisation about the step-4 base
\newcommand{\Trans}{\Pi}                           % transport integral
\newcommand{\resid}{\mathrm{res}}
\newcommand{\Resid}[1]{\mathrm{Res}_{#1}}
\newcommand{\boxtip}{\Box_{\rm tip}}
\newcommand{\boxfar}{\Box_{\rm far}}
\newcommand{\Eenv}{\mathcal E_{\rm env}}
\newcommand{\betab}{\beta_{\rm base}}
\newcommand{\tauI}{\tau^{\rm I}}                   % Phase I clock
% names of families of conditions: \Hh{3} = (H3), \Oh{4} = (O4), ...
\newcommand{\Hh}[1]{\textup{(H#1)}}
\newcommand{\Oh}[1]{\textup{(O#1)}}
\newcommand{\Lh}[1]{\textup{(L#1)}}
\newcommand{\Bh}[1]{\textup{(B#1)}}
\newcommand{\Kh}[1]{\textup{(K#1)}}
\newcommand{\Wh}[1]{\textup{(W#1)}}
\newcommand{\PhaseI}{Phase~I}
\newcommand{\PhaseII}{Phase~II}
% margin: threshold/bound (always)
\newcommand{\mgn}[1]{{\times}#1}

% ---- figures: placeholder macro so that drafts compile before F1-F7 exist ----
\newcommand{\figplaceholder}[1]{\fbox{\parbox{0.8\linewidth}{\centering\small[figure #1 to be supplied]}}}
 ...
%   paper_v1/supplement.tex  \documentclass[11pt]{amsart}
%                            \newif\ifsupplement \supplementtrue
%                            % ======================================================================
% preamble.tex -- shared preamble of the main paper and the supplement.
% Owner: coordinator.  Drafters add macros by request only (see DRAFTER_GUIDE.md).
%
% USE.  This file contains NO \documentclass.  The two master files are
%   paper_v1/main.tex        \documentclass[11pt]{amsart} % ======================================================================
% preamble.tex -- shared preamble of the main paper and the supplement.
% Owner: coordinator.  Drafters add macros by request only (see DRAFTER_GUIDE.md).
%
% USE.  This file contains NO \documentclass.  The two master files are
%   paper_v1/main.tex        \documentclass[11pt]{amsart} % ======================================================================
% preamble.tex -- shared preamble of the main paper and the supplement.
% Owner: coordinator.  Drafters add macros by request only (see DRAFTER_GUIDE.md).
%
% USE.  This file contains NO \documentclass.  The two master files are
%   paper_v1/main.tex        \documentclass[11pt]{amsart} \input{preamble} ...
%   paper_v1/supplement.tex  \documentclass[11pt]{amsart}
%                            \newif\ifsupplement \supplementtrue
%                            \input{preamble} ...
% (the main paper does nothing special: \ifsupplement defaults to false).
% Compile from paper_v1/ so that \input{numbers/numbers_paper.tex} is found.
% ======================================================================

% ---- which document are we in ----------------------------------------
\ifdefined\ifsupplement\else\newif\ifsupplement\fi

% ---- packages ---------------------------------------------------------
\usepackage[T1]{fontenc}
\usepackage[utf8]{inputenc}
\usepackage{lmodern}
\usepackage[margin=1.1in]{geometry}
\usepackage{amsmath,amssymb,amsthm,mathtools,bm}
\usepackage{booktabs,array,longtable,tabularx,multirow}
\usepackage{enumitem}
\usepackage{graphicx}
\usepackage{xcolor}
\usepackage{tikz}
\usetikzlibrary{positioning,arrows.meta,calc}
% xr-hyper must come before hyperref
\usepackage{xr-hyper}
\usepackage[colorlinks=true,linkcolor=blue!50!black,citecolor=green!35!black,urlcolor=blue!50!black]{hyperref}

\raggedbottom
\setlength{\emergencystretch}{2em}
\setlist[enumerate]{leftmargin=*,itemsep=2pt,topsep=3pt}
\setlist[itemize]{leftmargin=*,itemsep=2pt,topsep=3pt}

% ---- numbering --------------------------------------------------------
% Main paper: sections 1,2,...; theorems 2.3 (by section); equations (2.3).
% Supplement: the Parts A..K are the \section's, numbered by letter, so
% theorems are A.1, E.4, ... (do not confuse "Theorem A" of the paper, which
% has no dot, with "Theorem A.1" of Part A).
\ifsupplement
  \renewcommand{\thesection}{\Alph{section}}
\fi
\numberwithin{equation}{section}

% ---- theorem environments (numbered by section, one shared counter) ----
\theoremstyle{plain}
\newtheorem{theorem}{Theorem}[section]
\newtheorem{proposition}[theorem]{Proposition}
\newtheorem{lemma}[theorem]{Lemma}
\newtheorem{corollary}[theorem]{Corollary}
\newtheorem{maintheorem}[theorem]{Main Theorem}
% Theorems A, B, C, D of the paper: own counter, lettered, no section prefix
\newtheorem{thmlet}{Theorem}
\renewcommand{\thethmlet}{\Alph{thmlet}}
\theoremstyle{definition}
\newtheorem{definition}[theorem]{Definition}
\newtheorem{convention}[theorem]{Convention}
\newtheorem{example}[theorem]{Example}
\theoremstyle{remark}
\newtheorem{remark}[theorem]{Remark}

% A list of hypotheses labelled (H1),(H2),... : \begin{hyplist}{H} ... \end{hyplist}
% (also {O}, {L}, {B}, {K}, {W} for the families of section 5.3 of OUTLINE.md).
\newenvironment{hyplist}[1]%
  {\begin{enumerate}[label=\textup{(#1\arabic*)},leftmargin=3em,itemsep=2pt]}%
  {\end{enumerate}}

% ---- trust classes (OUTLINE 4.1, 5.1.3) -------------------------------
% Supplement: explicit tags.  Main paper: certified is the default and
% unmarked; every non-certified number is marked "(floating point)".
\ifsupplement
  \newcommand{\tcert}{\textup{[cert]}}   % arb/outward-rounded, pinned script and key
  \newcommand{\test}{\textup{[est]}}     % floating point or estimate; never in a theorem
  \newcommand{\thand}{\textup{[hand]}}   % proved by hand in the text
  \newcommand{\tlit}{\textup{[lit]}}     % cited from the literature
\else
  \newcommand{\tcert}{}
  \newcommand{\test}{\unskip\ \textup{(floating point)}}
  \newcommand{\thand}{}
  \newcommand{\tlit}{}
\fi

% ---- gaps: visible, counted, an error in a final build -----------------
% Usage: \GAP{one sentence saying what is missing}   (and an entry in ISSUES.md)
\newcounter{gapcount}
\newif\iffinalpaper   % \finalpapertrue makes every \GAP a hard error
\newcommand{\GAP}[1]{\stepcounter{gapcount}\typeout{GAP-MARKER \thegapcount}%
  \iffinalpaper\PackageError{paper}{\string\GAP\space used in a final build}{Remove all gaps.}\fi
  {\color{red}\sffamily\bfseries[GAP~\thegapcount:\ }{\color{red}\sffamily #1}{\color{red}\sffamily\bfseries]}}
\AtEndDocument{\typeout{GAP-MARKER total: \thegapcount}}

% ---- file names and JSON keys (supplement only) -------------------------
% No internal path appears in the main paper; a use there is a warning.
\newcommand{\fn}[1]{%
  \ifsupplement\else\PackageWarning{paper}{\string\fn\space used in the main paper}\fi
  \texttt{\detokenize{#1}}}
\newcommand{\key}[1]{\texttt{\detokenize{#1}}}

% ---- cross references between main paper and supplement ----------------
% Labels carry a unit prefix (section 5.3 of OUTLINE.md and DRAFTER_GUIDE.md):
%   main paper:  intro: obj: s3: s4: s5: s6: s7: s8: s9:   (kind: thm- lem- eq- tab- fig- ...)
%   supplement:  pA: pB: ... pK:
% so all labels are globally unique and the prefixes below are only a second guard.
% \Xref{label}  : reference to the OTHER document (hyperlinked via xr-hyper)
% \Lref{label}  : reference in the SAME document (plain \ref)
\ifsupplement
  \externaldocument[M-]{main}[main.pdf]
  \newcommand{\Xref}[1]{\ref{M-#1}}
\else
  \externaldocument[S-]{supplement}[supplement.pdf]
  \newcommand{\Xref}[1]{\ref{S-#1}}
\fi
\newcommand{\Lref}[1]{\ref{#1}}
% \Supp{pE}{lem-chart} -> "Lemma E.3" style text is produced by the drafter:
%   "Lemma~\Xref{pE:lem-chart} of the supplement".  \SuppPart{E} gives "Part E".
\newcommand{\SuppPart}[1]{Part~#1 of the supplement}

% ---- numbers (generated by unit U0; macros only, no literal certified numbers) ----
\IfFileExists{numbers/numbers_paper.tex}%
  {\input{numbers/numbers_paper.tex}}%
  {\PackageWarning{paper}{numbers/numbers_paper.tex not found (unit U0 not finished)}}

% ======================================================================
% NOTATION MACROS (OUTLINE 5.2 and notation.md).  One macro per symbol
% whose spelling is fixed by the notation table.  Plain letters (a, b, f, ...)
% have no macro.
% ======================================================================
% basic
\newcommand{\R}{\mathbb{R}}
\newcommand{\C}{\mathbb{C}}
\newcommand{\N}{\mathbb{N}}
\newcommand{\Z}{\mathbb{Z}}
\newcommand{\eps}{\varepsilon}
\newcommand{\ip}[2]{\langle #1,#2\rangle}
\newcommand{\abs}[1]{\lvert #1\rvert}
\newcommand{\norm}[1]{\lVert #1\rVert}
\newcommand{\dd}{\mathrm d}
\newcommand{\re}{\operatorname{Re}}
\newcommand{\im}{\operatorname{Im}}
% flow and curves
\newcommand{\arclen}{\mathfrak{s}}                 % arclength of the flow curve
\newcommand{\kss}{\kappa_{\mathfrak{ss}}}          % kappa_{ss} in the flow equation
\newcommand{\Area}{\operatorname{Area}}
\newcommand{\Len}{\mathcal L}                      % length of a curve
\newcommand{\spp}{\operatorname{sp}}                 % sup quantities sp^v_0, sp^v_1, sp^f_1 (were S^v_0, S^v_1, S^f_1)
\newcommand{\rot}{\operatorname{rot}}              % turning number
\newcommand{\Tsing}{T_*}                           % singular time
\newcommand{\Tmax}{T_{\max}}                       % maximal time of existence
\newcommand{\pstar}{p_*}                           % image point of the singularity
\newcommand{\Gz}{\Gamma_0}                         % initial curve
\newcommand{\vth}{\vartheta}                       % parameter of Gamma_0
\newcommand{\dclk}{\delta_t}                       % clock defect of step 4
\newcommand{\fall}{\varpi}                         % Fourier rates varpi_n = n^4-n^2
% rescaling
\newcommand{\lamone}{\lambda_1}
\newcommand{\rate}{\varsigma}                      % 1/4 - a
\newcommand{\xia}{\xi_a}                           % (1-4\bar a)^{-1}
% shrinker and step 1
\newcommand{\Sone}{S_1}
\newcommand{\Gsh}{\Gamma_{\!*}}                    % profile curve of S_1
\newcommand{\cSet}{\mathcal S}                     % certified set of shrinkers
\newcommand{\sigtip}{\sigma_{\rm tip}}
\newcommand{\Xo}{X_0}                              % 2.2490776515
\newcommand{\Dtip}{D}
\newcommand{\ktip}{k_0}
\newcommand{\shootM}{M}                            % shooting function
\newcommand{\Ypar}{\mathcal Y}                     % step-1 parallelepiped
\newcommand{\axis}{\Sigma}                         % symmetry axis of S_1
\newcommand{\sigc}{\sigma_{\rm c}}
\newcommand{\sigrev}{\sigma^{\rm rev}}
% linear theory
\newcommand{\Hrho}{\mathcal H_\rho}
\newcommand{\muslow}{\mu_3}
\newcommand{\rre}{\mathfrak r_{\rm re}}
\newcommand{\rim}{\mathfrak r_{\rm im}}
\newcommand{\Evev}{\mathrm{Ev}_{\rm even}}
\newcommand{\Evod}{\mathrm{Ev}_{\rm odd}}
\newcommand{\gvec}{\mathbf g}                      % translation vector in step 2
\newcommand{\cons}[1]{\mathfrak f_{#1}}            % constraint functions
\newcommand{\phirho}{\varphi_\rho}
% energies
\newcommand{\Lam}{\Lambda}
\newcommand{\En}{E}
\newcommand{\EnI}{E_I}
\newcommand{\muI}{\mu_I}                           % 0.08
\newcommand{\epsE}{\varepsilon_E}                  % 1e-5
\newcommand{\epsI}{\varepsilon_I}                  % 3e-5
\newcommand{\wn}[1]{\mathsf w_{#1}}                % ||w~^{(j)}||_rho
\newcommand{\vI}{v_I}                              % zeta w'
\newcommand{\muerr}{\mu_{\rm err}}                 % step-4 error-energy weight 0.036
% far field
\newcommand{\pU}{p_U}                              % 1.0246
\newcommand{\Uone}{U_1}
\newcommand{\sm}{\operatorname{sm}}                % smoothstep t^3(10-15t+6t^2)
\newcommand{\rc}{r_c}                              % centre radius (Lemma U)
\newcommand{\rchi}{r_\chi}                         % 253
\newcommand{\Iw}{I_w}                              % hand-off windows [y_w,1.5y_w]
\newcommand{\qa}{\mathfrak q_1}                    % quasilinear coefficients of the slope equation
\newcommand{\qb}{\mathfrak q_2}
% cylinders (Lemma O)
\newcommand{\cyl}{\ell}                            % cylinder length unit
\newcommand{\MO}{M_O}
\newcommand{\KO}{K_O}
\newcommand{\Hcyl}{H_{\rm cyl}}
\newcommand{\XiO}{\Xi_O}
\newcommand{\fO}{f_O}
% step 4
\newcommand{\Ybar}{\bar Y}
\newcommand{\Lt}{\widetilde L}                     % linearisation about the step-4 base
\newcommand{\Trans}{\Pi}                           % transport integral
\newcommand{\resid}{\mathrm{res}}
\newcommand{\Resid}[1]{\mathrm{Res}_{#1}}
\newcommand{\boxtip}{\Box_{\rm tip}}
\newcommand{\boxfar}{\Box_{\rm far}}
\newcommand{\Eenv}{\mathcal E_{\rm env}}
\newcommand{\betab}{\beta_{\rm base}}
\newcommand{\tauI}{\tau^{\rm I}}                   % Phase I clock
% names of families of conditions: \Hh{3} = (H3), \Oh{4} = (O4), ...
\newcommand{\Hh}[1]{\textup{(H#1)}}
\newcommand{\Oh}[1]{\textup{(O#1)}}
\newcommand{\Lh}[1]{\textup{(L#1)}}
\newcommand{\Bh}[1]{\textup{(B#1)}}
\newcommand{\Kh}[1]{\textup{(K#1)}}
\newcommand{\Wh}[1]{\textup{(W#1)}}
\newcommand{\PhaseI}{Phase~I}
\newcommand{\PhaseII}{Phase~II}
% margin: threshold/bound (always)
\newcommand{\mgn}[1]{{\times}#1}

% ---- figures: placeholder macro so that drafts compile before F1-F7 exist ----
\newcommand{\figplaceholder}[1]{\fbox{\parbox{0.8\linewidth}{\centering\small[figure #1 to be supplied]}}}
 ...
%   paper_v1/supplement.tex  \documentclass[11pt]{amsart}
%                            \newif\ifsupplement \supplementtrue
%                            % ======================================================================
% preamble.tex -- shared preamble of the main paper and the supplement.
% Owner: coordinator.  Drafters add macros by request only (see DRAFTER_GUIDE.md).
%
% USE.  This file contains NO \documentclass.  The two master files are
%   paper_v1/main.tex        \documentclass[11pt]{amsart} \input{preamble} ...
%   paper_v1/supplement.tex  \documentclass[11pt]{amsart}
%                            \newif\ifsupplement \supplementtrue
%                            \input{preamble} ...
% (the main paper does nothing special: \ifsupplement defaults to false).
% Compile from paper_v1/ so that \input{numbers/numbers_paper.tex} is found.
% ======================================================================

% ---- which document are we in ----------------------------------------
\ifdefined\ifsupplement\else\newif\ifsupplement\fi

% ---- packages ---------------------------------------------------------
\usepackage[T1]{fontenc}
\usepackage[utf8]{inputenc}
\usepackage{lmodern}
\usepackage[margin=1.1in]{geometry}
\usepackage{amsmath,amssymb,amsthm,mathtools,bm}
\usepackage{booktabs,array,longtable,tabularx,multirow}
\usepackage{enumitem}
\usepackage{graphicx}
\usepackage{xcolor}
\usepackage{tikz}
\usetikzlibrary{positioning,arrows.meta,calc}
% xr-hyper must come before hyperref
\usepackage{xr-hyper}
\usepackage[colorlinks=true,linkcolor=blue!50!black,citecolor=green!35!black,urlcolor=blue!50!black]{hyperref}

\raggedbottom
\setlength{\emergencystretch}{2em}
\setlist[enumerate]{leftmargin=*,itemsep=2pt,topsep=3pt}
\setlist[itemize]{leftmargin=*,itemsep=2pt,topsep=3pt}

% ---- numbering --------------------------------------------------------
% Main paper: sections 1,2,...; theorems 2.3 (by section); equations (2.3).
% Supplement: the Parts A..K are the \section's, numbered by letter, so
% theorems are A.1, E.4, ... (do not confuse "Theorem A" of the paper, which
% has no dot, with "Theorem A.1" of Part A).
\ifsupplement
  \renewcommand{\thesection}{\Alph{section}}
\fi
\numberwithin{equation}{section}

% ---- theorem environments (numbered by section, one shared counter) ----
\theoremstyle{plain}
\newtheorem{theorem}{Theorem}[section]
\newtheorem{proposition}[theorem]{Proposition}
\newtheorem{lemma}[theorem]{Lemma}
\newtheorem{corollary}[theorem]{Corollary}
\newtheorem{maintheorem}[theorem]{Main Theorem}
% Theorems A, B, C, D of the paper: own counter, lettered, no section prefix
\newtheorem{thmlet}{Theorem}
\renewcommand{\thethmlet}{\Alph{thmlet}}
\theoremstyle{definition}
\newtheorem{definition}[theorem]{Definition}
\newtheorem{convention}[theorem]{Convention}
\newtheorem{example}[theorem]{Example}
\theoremstyle{remark}
\newtheorem{remark}[theorem]{Remark}

% A list of hypotheses labelled (H1),(H2),... : \begin{hyplist}{H} ... \end{hyplist}
% (also {O}, {L}, {B}, {K}, {W} for the families of section 5.3 of OUTLINE.md).
\newenvironment{hyplist}[1]%
  {\begin{enumerate}[label=\textup{(#1\arabic*)},leftmargin=3em,itemsep=2pt]}%
  {\end{enumerate}}

% ---- trust classes (OUTLINE 4.1, 5.1.3) -------------------------------
% Supplement: explicit tags.  Main paper: certified is the default and
% unmarked; every non-certified number is marked "(floating point)".
\ifsupplement
  \newcommand{\tcert}{\textup{[cert]}}   % arb/outward-rounded, pinned script and key
  \newcommand{\test}{\textup{[est]}}     % floating point or estimate; never in a theorem
  \newcommand{\thand}{\textup{[hand]}}   % proved by hand in the text
  \newcommand{\tlit}{\textup{[lit]}}     % cited from the literature
\else
  \newcommand{\tcert}{}
  \newcommand{\test}{\unskip\ \textup{(floating point)}}
  \newcommand{\thand}{}
  \newcommand{\tlit}{}
\fi

% ---- gaps: visible, counted, an error in a final build -----------------
% Usage: \GAP{one sentence saying what is missing}   (and an entry in ISSUES.md)
\newcounter{gapcount}
\newif\iffinalpaper   % \finalpapertrue makes every \GAP a hard error
\newcommand{\GAP}[1]{\stepcounter{gapcount}\typeout{GAP-MARKER \thegapcount}%
  \iffinalpaper\PackageError{paper}{\string\GAP\space used in a final build}{Remove all gaps.}\fi
  {\color{red}\sffamily\bfseries[GAP~\thegapcount:\ }{\color{red}\sffamily #1}{\color{red}\sffamily\bfseries]}}
\AtEndDocument{\typeout{GAP-MARKER total: \thegapcount}}

% ---- file names and JSON keys (supplement only) -------------------------
% No internal path appears in the main paper; a use there is a warning.
\newcommand{\fn}[1]{%
  \ifsupplement\else\PackageWarning{paper}{\string\fn\space used in the main paper}\fi
  \texttt{\detokenize{#1}}}
\newcommand{\key}[1]{\texttt{\detokenize{#1}}}

% ---- cross references between main paper and supplement ----------------
% Labels carry a unit prefix (section 5.3 of OUTLINE.md and DRAFTER_GUIDE.md):
%   main paper:  intro: obj: s3: s4: s5: s6: s7: s8: s9:   (kind: thm- lem- eq- tab- fig- ...)
%   supplement:  pA: pB: ... pK:
% so all labels are globally unique and the prefixes below are only a second guard.
% \Xref{label}  : reference to the OTHER document (hyperlinked via xr-hyper)
% \Lref{label}  : reference in the SAME document (plain \ref)
\ifsupplement
  \externaldocument[M-]{main}[main.pdf]
  \newcommand{\Xref}[1]{\ref{M-#1}}
\else
  \externaldocument[S-]{supplement}[supplement.pdf]
  \newcommand{\Xref}[1]{\ref{S-#1}}
\fi
\newcommand{\Lref}[1]{\ref{#1}}
% \Supp{pE}{lem-chart} -> "Lemma E.3" style text is produced by the drafter:
%   "Lemma~\Xref{pE:lem-chart} of the supplement".  \SuppPart{E} gives "Part E".
\newcommand{\SuppPart}[1]{Part~#1 of the supplement}

% ---- numbers (generated by unit U0; macros only, no literal certified numbers) ----
\IfFileExists{numbers/numbers_paper.tex}%
  {\input{numbers/numbers_paper.tex}}%
  {\PackageWarning{paper}{numbers/numbers_paper.tex not found (unit U0 not finished)}}

% ======================================================================
% NOTATION MACROS (OUTLINE 5.2 and notation.md).  One macro per symbol
% whose spelling is fixed by the notation table.  Plain letters (a, b, f, ...)
% have no macro.
% ======================================================================
% basic
\newcommand{\R}{\mathbb{R}}
\newcommand{\C}{\mathbb{C}}
\newcommand{\N}{\mathbb{N}}
\newcommand{\Z}{\mathbb{Z}}
\newcommand{\eps}{\varepsilon}
\newcommand{\ip}[2]{\langle #1,#2\rangle}
\newcommand{\abs}[1]{\lvert #1\rvert}
\newcommand{\norm}[1]{\lVert #1\rVert}
\newcommand{\dd}{\mathrm d}
\newcommand{\re}{\operatorname{Re}}
\newcommand{\im}{\operatorname{Im}}
% flow and curves
\newcommand{\arclen}{\mathfrak{s}}                 % arclength of the flow curve
\newcommand{\kss}{\kappa_{\mathfrak{ss}}}          % kappa_{ss} in the flow equation
\newcommand{\Area}{\operatorname{Area}}
\newcommand{\Len}{\mathcal L}                      % length of a curve
\newcommand{\spp}{\operatorname{sp}}                 % sup quantities sp^v_0, sp^v_1, sp^f_1 (were S^v_0, S^v_1, S^f_1)
\newcommand{\rot}{\operatorname{rot}}              % turning number
\newcommand{\Tsing}{T_*}                           % singular time
\newcommand{\Tmax}{T_{\max}}                       % maximal time of existence
\newcommand{\pstar}{p_*}                           % image point of the singularity
\newcommand{\Gz}{\Gamma_0}                         % initial curve
\newcommand{\vth}{\vartheta}                       % parameter of Gamma_0
\newcommand{\dclk}{\delta_t}                       % clock defect of step 4
\newcommand{\fall}{\varpi}                         % Fourier rates varpi_n = n^4-n^2
% rescaling
\newcommand{\lamone}{\lambda_1}
\newcommand{\rate}{\varsigma}                      % 1/4 - a
\newcommand{\xia}{\xi_a}                           % (1-4\bar a)^{-1}
% shrinker and step 1
\newcommand{\Sone}{S_1}
\newcommand{\Gsh}{\Gamma_{\!*}}                    % profile curve of S_1
\newcommand{\cSet}{\mathcal S}                     % certified set of shrinkers
\newcommand{\sigtip}{\sigma_{\rm tip}}
\newcommand{\Xo}{X_0}                              % 2.2490776515
\newcommand{\Dtip}{D}
\newcommand{\ktip}{k_0}
\newcommand{\shootM}{M}                            % shooting function
\newcommand{\Ypar}{\mathcal Y}                     % step-1 parallelepiped
\newcommand{\axis}{\Sigma}                         % symmetry axis of S_1
\newcommand{\sigc}{\sigma_{\rm c}}
\newcommand{\sigrev}{\sigma^{\rm rev}}
% linear theory
\newcommand{\Hrho}{\mathcal H_\rho}
\newcommand{\muslow}{\mu_3}
\newcommand{\rre}{\mathfrak r_{\rm re}}
\newcommand{\rim}{\mathfrak r_{\rm im}}
\newcommand{\Evev}{\mathrm{Ev}_{\rm even}}
\newcommand{\Evod}{\mathrm{Ev}_{\rm odd}}
\newcommand{\gvec}{\mathbf g}                      % translation vector in step 2
\newcommand{\cons}[1]{\mathfrak f_{#1}}            % constraint functions
\newcommand{\phirho}{\varphi_\rho}
% energies
\newcommand{\Lam}{\Lambda}
\newcommand{\En}{E}
\newcommand{\EnI}{E_I}
\newcommand{\muI}{\mu_I}                           % 0.08
\newcommand{\epsE}{\varepsilon_E}                  % 1e-5
\newcommand{\epsI}{\varepsilon_I}                  % 3e-5
\newcommand{\wn}[1]{\mathsf w_{#1}}                % ||w~^{(j)}||_rho
\newcommand{\vI}{v_I}                              % zeta w'
\newcommand{\muerr}{\mu_{\rm err}}                 % step-4 error-energy weight 0.036
% far field
\newcommand{\pU}{p_U}                              % 1.0246
\newcommand{\Uone}{U_1}
\newcommand{\sm}{\operatorname{sm}}                % smoothstep t^3(10-15t+6t^2)
\newcommand{\rc}{r_c}                              % centre radius (Lemma U)
\newcommand{\rchi}{r_\chi}                         % 253
\newcommand{\Iw}{I_w}                              % hand-off windows [y_w,1.5y_w]
\newcommand{\qa}{\mathfrak q_1}                    % quasilinear coefficients of the slope equation
\newcommand{\qb}{\mathfrak q_2}
% cylinders (Lemma O)
\newcommand{\cyl}{\ell}                            % cylinder length unit
\newcommand{\MO}{M_O}
\newcommand{\KO}{K_O}
\newcommand{\Hcyl}{H_{\rm cyl}}
\newcommand{\XiO}{\Xi_O}
\newcommand{\fO}{f_O}
% step 4
\newcommand{\Ybar}{\bar Y}
\newcommand{\Lt}{\widetilde L}                     % linearisation about the step-4 base
\newcommand{\Trans}{\Pi}                           % transport integral
\newcommand{\resid}{\mathrm{res}}
\newcommand{\Resid}[1]{\mathrm{Res}_{#1}}
\newcommand{\boxtip}{\Box_{\rm tip}}
\newcommand{\boxfar}{\Box_{\rm far}}
\newcommand{\Eenv}{\mathcal E_{\rm env}}
\newcommand{\betab}{\beta_{\rm base}}
\newcommand{\tauI}{\tau^{\rm I}}                   % Phase I clock
% names of families of conditions: \Hh{3} = (H3), \Oh{4} = (O4), ...
\newcommand{\Hh}[1]{\textup{(H#1)}}
\newcommand{\Oh}[1]{\textup{(O#1)}}
\newcommand{\Lh}[1]{\textup{(L#1)}}
\newcommand{\Bh}[1]{\textup{(B#1)}}
\newcommand{\Kh}[1]{\textup{(K#1)}}
\newcommand{\Wh}[1]{\textup{(W#1)}}
\newcommand{\PhaseI}{Phase~I}
\newcommand{\PhaseII}{Phase~II}
% margin: threshold/bound (always)
\newcommand{\mgn}[1]{{\times}#1}

% ---- figures: placeholder macro so that drafts compile before F1-F7 exist ----
\newcommand{\figplaceholder}[1]{\fbox{\parbox{0.8\linewidth}{\centering\small[figure #1 to be supplied]}}}
 ...
% (the main paper does nothing special: \ifsupplement defaults to false).
% Compile from paper_v1/ so that % numbers_paper.tex -- GENERATED from ledger.json by build_ledger.py; do not edit.
% Rounding: upper bounds up, lower bounds and margins down at the printed digit; margin = threshold / bound.
% Use as $\times\NName$ for margins.  kind: cert certified, est estimate, hand hypothesis constant, doc only a document quotes it.
% NUEhatZero [cert, upper] U_cert_v35.json :: supersolution.EA.0
\newcommand{\NUEhatZero}{\ensuremath{1.2018\cdot10^{-6}}}
% NUEhatOne [cert, upper] U_cert_v35.json :: supersolution.EA.1
\newcommand{\NUEhatOne}{\ensuremath{1.0441\cdot10^{-4}}}
% NUEhatTwo [cert, upper] U_cert_v35.json :: supersolution.EA.2
\newcommand{\NUEhatTwo}{\ensuremath{0.056694}}
% NUBhatZero [cert, upper] U_cert_v35.json :: supersolution.Bhat.0
\newcommand{\NUBhatZero}{\ensuremath{4.8735\cdot10^{-6}}}
% NUBhatOne [cert, upper] U_cert_v35.json :: supersolution.Bhat.1
\newcommand{\NUBhatOne}{\ensuremath{5.5933\cdot10^{-4}}}
% NUBhatTwo [cert, upper] U_cert_v35.json :: supersolution.Bhat.2
\newcommand{\NUBhatTwo}{\ensuremath{0.36675}}
% NUSupersolutionCZeroZero [cert, upper] U_cert_v35.json :: supersolution.C0.0
\newcommand{\NUSupersolutionCZeroZero}{\ensuremath{1.1340\cdot10^{-6}}}
% NUSupersolutionCZeroOne [cert, upper] U_cert_v35.json :: supersolution.C0.1
\newcommand{\NUSupersolutionCZeroOne}{\ensuremath{1.6623\cdot10^{-4}}}
% NUSupersolutionCOneZero [cert, upper] U_cert_v35.json :: supersolution.C1.0
\newcommand{\NUSupersolutionCOneZero}{\ensuremath{1.5241\cdot10^{-9}}}
% NUSupersolutionCOneOne [cert, upper] U_cert_v35.json :: supersolution.C1.1
\newcommand{\NUSupersolutionCOneOne}{\ensuremath{5.5255\cdot10^{-8}}}
% NUSupersolutionNTwoZero [cert, upper] U_cert_v35.json :: supersolution.N2.0
\newcommand{\NUSupersolutionNTwoZero}{\ensuremath{0.071472}}
% NUSupersolutionNTwoOne [cert, upper] U_cert_v35.json :: supersolution.N2.1
\newcommand{\NUSupersolutionNTwoOne}{\ensuremath{0.33012}}
% NUSupersolutionNTwoTwo [cert, upper] U_cert_v35.json :: supersolution.N2.2
\newcommand{\NUSupersolutionNTwoTwo}{\ensuremath{5.0406}}
% NUSupersolutionNTwoOverFourNTwoRefZero [cert, upper] U_cert_v35.json :: supersolution.N2_over_4N2ref.0
\newcommand{\NUSupersolutionNTwoOverFourNTwoRefZero}{\ensuremath{0.18}}
% NUSupersolutionNTwoOverFourNTwoRefOne [cert, upper] U_cert_v35.json :: supersolution.N2_over_4N2ref.1
\newcommand{\NUSupersolutionNTwoOverFourNTwoRefOne}{\ensuremath{0.42}}
% NUSupersolutionNTwoOverFourNTwoRefTwo [cert, upper] U_cert_v35.json :: supersolution.N2_over_4N2ref.2
\newcommand{\NUSupersolutionNTwoOverFourNTwoRefTwo}{\ensuremath{0.26}}
% NUSupersolutionMOneZHypCheckZonesIIISOneTwoBhat [cert, upper] U_cert_v35.json :: supersolution.M1z_hyp_check.zones_I_II.S1_2Bhat
\newcommand{\NUSupersolutionMOneZHypCheckZonesIIISOneTwoBhat}{\ensuremath{0.28585}}
% NUSupersolutionMOneZHypCheckZoneIIIBOnePlusPUOne [cert, upper] U_cert_v35.json :: supersolution.M1z_hyp_check.zone_III.b1_plus_pU1
\newcommand{\NUSupersolutionMOneZHypCheckZoneIIIBOnePlusPUOne}{\ensuremath{0.50786}}
% NUMarginLateMOne [cert, margin] U_cert_v35.json :: margins.late_M1
\newcommand{\NUMarginLateMOne}{\ensuremath{2.353}}
% NUMarginLateMZero [cert, margin] U_cert_v35.json :: margins.late_M0
\newcommand{\NUMarginLateMZero}{\ensuremath{4.892}}
% NUMarginsLateMZeroZ [cert, margin] U_cert_v35.json :: margins.late_M0z
\newcommand{\NUMarginsLateMZeroZ}{\ensuremath{12.99}}
% NUMarginPresentMOne [cert, margin] U_cert_v35.json :: margins.present_M1_min.value
\newcommand{\NUMarginPresentMOne}{\ensuremath{2.770}}
% NUMarginPresentMZero [cert, margin] U_cert_v35.json :: margins.present_M0_min.value
\newcommand{\NUMarginPresentMZero}{\ensuremath{2.225}}
% NUMarginsPresentMZeroZMinValue [cert, margin] U_cert_v35.json :: margins.present_M0z_min.value
\newcommand{\NUMarginsPresentMZeroZMinValue}{\ensuremath{14.37}}
% NULoopGainZeroTotal [cert, upper] U_cert_v35.json :: loop_gain.0.total
\newcommand{\NULoopGainZeroTotal}{\ensuremath{1.23\cdot10^{-7}}}
% NULoopGainOneTotal [cert, upper] U_cert_v35.json :: loop_gain.1.total
\newcommand{\NULoopGainOneTotal}{\ensuremath{0.037169}}
% NULoopGainTwoNeighbour [cert, upper] U_cert_v35.json :: loop_gain.2.neighbour
\newcommand{\NULoopGainTwoNeighbour}{\ensuremath{0.037018}}
% NULoopGainTwoThroughA [cert, upper] U_cert_v35.json :: loop_gain.2.through_A
\newcommand{\NULoopGainTwoThroughA}{\ensuremath{0.0020172}}
% NUfineBhatZero [cert, upper] U_cert_v35.json :: ufine.Bhat.0
\newcommand{\NUfineBhatZero}{\ensuremath{4.8723\cdot10^{-6}}}
% NUfineBhatOne [cert, upper] U_cert_v35.json :: ufine.Bhat.1
\newcommand{\NUfineBhatOne}{\ensuremath{2.6266\cdot10^{-4}}}
% NUfineBhatTwo [cert, upper] U_cert_v35.json :: ufine.Bhat.2
\newcommand{\NUfineBhatTwo}{\ensuremath{0.025842}}
% NUcdeltaU [cert, upper] U_cert_v35.json :: launch.c_dU
\newcommand{\NUcdeltaU}{\ensuremath{1.7187\cdot10^{-8}}}
% NULaunchWmin [cert, lower] U_cert_v35.json :: launch.wmin
\newcommand{\NULaunchWmin}{\ensuremath{0.73492}}
% NULaunchCAD [cert, upper] U_cert_v35.json :: launch.CA_D
\newcommand{\NULaunchCAD}{\ensuremath{2.2429}}
% NULaunchQD [cert, upper] U_cert_v35.json :: launch.qD
\newcommand{\NULaunchQD}{\ensuremath{0.018509}}
% NULaunchDeltaInZero [cert, upper] U_cert_v35.json :: launch.Delta_in.0
\newcommand{\NULaunchDeltaInZero}{\ensuremath{3.131\cdot10^{-12}}}
% NULaunchDeltaInOne [cert, upper] U_cert_v35.json :: launch.Delta_in.1
\newcommand{\NULaunchDeltaInOne}{\ensuremath{3.287\cdot10^{-12}}}
% NULaunchDeltaInTwo [cert, upper] U_cert_v35.json :: launch.Delta_in.2
\newcommand{\NULaunchDeltaInTwo}{\ensuremath{3.451\cdot10^{-12}}}
% NULaunchPhiZero [cert, upper] U_cert_v35.json :: launch.Phi.0
\newcommand{\NULaunchPhiZero}{\ensuremath{6.434\cdot10^{-9}}}
% NULaunchPhiOne [cert, upper] U_cert_v35.json :: launch.Phi.1
\newcommand{\NULaunchPhiOne}{\ensuremath{1.4217\cdot10^{-8}}}
% NULaunchPhiTwo [cert, upper] U_cert_v35.json :: launch.Phi.2
\newcommand{\NULaunchPhiTwo}{\ensuremath{3.8404\cdot10^{-8}}}
% NUGammaZero [cert, upper] U_cert_v35.json :: Gamma.0
\newcommand{\NUGammaZero}{\ensuremath{4.0000358}}
% NUGammaOne [cert, upper] U_cert_v35.json :: Gamma.1
\newcommand{\NUGammaOne}{\ensuremath{16.000322}}
% NUGammaTwo [cert, upper] U_cert_v35.json :: Gamma.2
\newcommand{\NUGammaTwo}{\ensuremath{64.00218}}
% NUKappa [cert, upper] U_cert_v35.json :: kappa
\newcommand{\NUKappa}{\ensuremath{8.0000358}}
% NUKappaTwo [cert, upper] U_cert_v35.json :: kappa2
\newcommand{\NUKappaTwo}{\ensuremath{64.00061}}
% NUGhatZeroZero [cert, upper] U_cert_v35.json :: Ghat.0.0
\newcommand{\NUGhatZeroZero}{\ensuremath{51322.6}}
% NUGhatOneOne [cert, upper] U_cert_v35.json :: Ghat.1.1
\newcommand{\NUGhatOneOne}{\ensuremath{33962.9}}
% NUGhatTwoTwo [cert, upper] U_cert_v35.json :: Ghat.2.2
\newcommand{\NUGhatTwoTwo}{\ensuremath{348341.1}}
% NUGhatOneZero [cert, upper] U_cert_v35.json :: Ghat.1.0
\newcommand{\NUGhatOneZero}{\ensuremath{4.938}}
% NUGhatTwoZero [cert, upper] U_cert_v35.json :: Ghat.2.0
\newcommand{\NUGhatTwoZero}{\ensuremath{27.984}}
% NUGhatTwoOne [cert, upper] U_cert_v35.json :: Ghat.2.1
\newcommand{\NUGhatTwoOne}{\ensuremath{458.53}}
% NUKZeroZero [cert, upper] U_cert_v35.json :: K.0.0
\newcommand{\NUKZeroZero}{\ensuremath{8034.21}}
% NUKZeroOne [cert, upper] U_cert_v35.json :: K.0.1
\newcommand{\NUKZeroOne}{\ensuremath{724.89}}
% NUKOneZero [cert, upper] U_cert_v35.json :: K.1.0
\newcommand{\NUKOneZero}{\ensuremath{0.08152}}
% NUKOneOne [cert, upper] U_cert_v35.json :: K.1.1
\newcommand{\NUKOneOne}{\ensuremath{6336.94}}
% NUKOneTwo [cert, upper] U_cert_v35.json :: K.1.2
\newcommand{\NUKOneTwo}{\ensuremath{545.59}}
% NUKOneThree [cert, upper] U_cert_v35.json :: K.1.3
\newcommand{\NUKOneThree}{\ensuremath{0.58700}}
% NUKTwoZero [cert, upper] U_cert_v35.json :: K.2.0
\newcommand{\NUKTwoZero}{\ensuremath{0.87186}}
% NUKTwoOne [cert, upper] U_cert_v35.json :: K.2.1
\newcommand{\NUKTwoOne}{\ensuremath{528.62}}
% NUKTwoTwo [cert, upper] U_cert_v35.json :: K.2.2
\newcommand{\NUKTwoTwo}{\ensuremath{23425.98}}
% NUKTwoThree [cert, upper] U_cert_v35.json :: K.2.3
\newcommand{\NUKTwoThree}{\ensuremath{555.62}}
% NUFootwOne [cert, upper] U_cert_v35.json :: footw.1
\newcommand{\NUFootwOne}{\ensuremath{2.6361}}
% NUFootwTwo [cert, upper] U_cert_v35.json :: footw.2
\newcommand{\NUFootwTwo}{\ensuremath{2663.3}}
% NUCThreeThree [cert, upper] U_cert_v35.json :: c33
\newcommand{\NUCThreeThree}{\ensuremath{108.02}}
% NUCTwoTwoThree [cert, upper] U_cert_v35.json :: c223
\newcommand{\NUCTwoTwoThree}{\ensuremath{7.3284}}
% NUCEta [cert, upper] U_cert_v35.json :: C_eta
\newcommand{\NUCEta}{\ensuremath{13.376}}
% NULMax [cert, upper] U_cert_v35.json :: L_max
\newcommand{\NULMax}{\ensuremath{3.5352}}
% NURRatio [cert, upper] U_cert_v35.json :: R_ratio
\newcommand{\NURRatio}{\ensuremath{11.593}}
% NUSweepsSigmaInCoverSup [cert, upper] U_cert_v35.json :: sweeps.Sigma_in.cover_sup
\newcommand{\NUSweepsSigmaInCoverSup}{\ensuremath{5.0892}}
% NUSweepsSigmaECoverSup [cert, upper] U_cert_v35.json :: sweeps.Sigma_E.cover_sup
\newcommand{\NUSweepsSigmaECoverSup}{\ensuremath{0.30846}}
% NUSweepsSXCoverSup [cert, upper] U_cert_v35.json :: sweeps.S_X.cover_sup
\newcommand{\NUSweepsSXCoverSup}{\ensuremath{1.5243}}
% NUSweepsSigmaXCoverSup [cert, upper] U_cert_v35.json :: sweeps.Sigma_X.cover_sup
\newcommand{\NUSweepsSigmaXCoverSup}{\ensuremath{3.0081}}
% NUSweepsSigmaRsTwoCoverSup [cert, upper] U_cert_v35.json :: sweeps.Sigma_rs.2.cover_sup
\newcommand{\NUSweepsSigmaRsTwoCoverSup}{\ensuremath{4.3350\cdot10^{-3}}}
% NUDataTermsLateDZero [cert, upper] U_cert_v35.json :: data_terms_late.D.0
\newcommand{\NUDataTermsLateDZero}{\ensuremath{1.7330\cdot10^{-6}}}
% NUDataTermsLateDOne [cert, upper] U_cert_v35.json :: data_terms_late.D.1
\newcommand{\NUDataTermsLateDOne}{\ensuremath{2.5353\cdot10^{-4}}}
% NUDataTermsLateSlopeShareCOneSigmaXOverDkZero [cert, upper] U_cert_v35.json :: data_terms_late.slope_share_C1SigmaX_over_Dk.0
\newcommand{\NUDataTermsLateSlopeShareCOneSigmaXOverDkZero}{\ensuremath{2.7\cdot10^{-3}}}
% NUDataTermsLateSlopeShareCOneSigmaXOverDkOne [cert, upper] U_cert_v35.json :: data_terms_late.slope_share_C1SigmaX_over_Dk.1
\newcommand{\NUDataTermsLateSlopeShareCOneSigmaXOverDkOne}{\ensuremath{6.6\cdot10^{-4}}}
% NURLtTwoMarginZero [cert, margin] U_cert_v35.json :: R_lt2_margin.0
\newcommand{\NURLtTwoMarginZero}{\ensuremath{4.432}}
% NURLtTwoMarginOne [cert, margin] U_cert_v35.json :: R_lt2_margin.1
\newcommand{\NURLtTwoMarginOne}{\ensuremath{6.697}}
% NUThetaTwoEight [cert, upper] U_cert_v35.json :: theta28
\newcommand{\NUThetaTwoEight}{\ensuremath{0.37011}}
% NUDFiveSix [cert, upper] U_cert_v35.json :: D56
\newcommand{\NUDFiveSix}{\ensuremath{1.0368}}
% NUDTwoEight [cert, upper] U_cert_v35.json :: D28
\newcommand{\NUDTwoEight}{\ensuremath{1.5876}}
% NUSThreeThreeFootOverBhatTwo [cert, upper] U_cert_v35.json :: S33_foot_over_Bhat2
\newcommand{\NUSThreeThreeFootOverBhatTwo}{\ensuremath{1.69\cdot10^{-3}}}
% NSAHJWKappaJ [cert, upper] A_interface.json :: HJW.kappa_J
\newcommand{\NSAHJWKappaJ}{\ensuremath{0.042351}}
% NSAGpZero [cert, lower] A_interface.json :: res.tight.gp.0
\newcommand{\NSAGpZero}{\ensuremath{0.24292}}
% NSAGpOne [cert, lower] A_interface.json :: res.tight.gp.1
\newcommand{\NSAGpOne}{\ensuremath{0.24106}}
% NSAGpTwo [cert, lower] A_interface.json :: res.tight.gp.2
\newcommand{\NSAGpTwo}{\ensuremath{0.11961}}
% NSAAminC [cert, lower] A_interface.json :: res.tight.amin_c
\newcommand{\NSAAminC}{\ensuremath{0.97623}}
% NSARhoJ [cert, upper] A_interface.json :: res.tight.rhoJ
\newcommand{\NSARhoJ}{\ensuremath{5.4995\cdot10^{-7}}}
% NSAKappaOTwoTwo [cert, upper] A_interface.json :: res.tight.kappa_o.2.2
\newcommand{\NSAKappaOTwoTwo}{\ensuremath{1.2416\cdot10^{-5}}}
% NSADsrcZero [cert, upper] A_interface.json :: res.tight.Dsrc.0
\newcommand{\NSADsrcZero}{\ensuremath{1.8092\cdot10^{-7}}}
% NSADsrcOne [cert, upper] A_interface.json :: res.tight.Dsrc.1
\newcommand{\NSADsrcOne}{\ensuremath{9.8276\cdot10^{-5}}}
% NSABZeroDWindowSlopeZero [cert, upper] A_interface.json :: res.tight.B0.D_window_slope.0
\newcommand{\NSABZeroDWindowSlopeZero}{\ensuremath{8.7399\cdot10^{-10}}}
% NSABZeroDWindowSlopeOne [cert, upper] A_interface.json :: res.tight.B0.D_window_slope.1
\newcommand{\NSABZeroDWindowSlopeOne}{\ensuremath{5.2599\cdot10^{-8}}}
% NSAFTwoZero [cert, upper] A_interface.json :: res.tight.F2.0
\newcommand{\NSAFTwoZero}{\ensuremath{5.6506\cdot10^{-13}}}
% NSAFTwoOne [cert, upper] A_interface.json :: res.tight.F2.1
\newcommand{\NSAFTwoOne}{\ensuremath{5.9320\cdot10^{-13}}}
% NSACoupOne [cert, upper] A_interface.json :: res.tight.coup.1
\newcommand{\NSACoupOne}{\ensuremath{1.1568\cdot10^{-12}}}
% NSACOneZeroZero [cert, upper] A_interface.json :: res.tight.C1_0.0
\newcommand{\NSACOneZeroZero}{\ensuremath{7.2060\cdot10^{-10}}}
% NSACOneZeroOne [cert, upper] A_interface.json :: res.tight.C1_0.1
\newcommand{\NSACOneZeroOne}{\ensuremath{4.3143\cdot10^{-8}}}
% NSABZeroSupCoreZero [cert, upper] A_interface.json :: res.tight.B0.sup_core.0
\newcommand{\NSABZeroSupCoreZero}{\ensuremath{4.23\cdot10^{-5}}}
% NSACoreFluxMarginTwoFLUXZeroOverKZero [cert, margin] A_interface.json :: res.tight.core.flux_margin.2.FLUX0_over_K0
\newcommand{\NSACoreFluxMarginTwoFLUXZeroOverKZero}{\ensuremath{3.396}}
% NSACoreFluxMarginTwoFLUXTwoOverKTwo [cert, margin] A_interface.json :: res.tight.core.flux_margin.2.FLUX2_over_K2
\newcommand{\NSACoreFluxMarginTwoFLUXTwoOverKTwo}{\ensuremath{6.012}}
% NSAMarginFluxZero [cert, margin] asm_design_N2ref4.json :: res.tight.flux.beta.2.margin_FLUX0
\newcommand{\NSAMarginFluxZero}{\ensuremath{3.365}}
% NSAMarginFluxTwo [cert, margin] asm_design_N2ref4.json :: res.tight.flux.beta.2.margin_FLUX2
\newcommand{\NSAMarginFluxTwo}{\ensuremath{5.955}}
% NSAfourSZero [cert, upper] asm_design_N2ref4.json :: res.tight.S0
\newcommand{\NSAfourSZero}{\ensuremath{7.1878\cdot10^{-3}}}
% NSAfourSOne [cert, upper] asm_design_N2ref4.json :: res.tight.S1
\newcommand{\NSAfourSOne}{\ensuremath{0.10430}}
% NSAMarginCoreSZeroV [cert, margin] asm_design_N2ref4.json :: res.tight.margins.S0v_vs_0.0153
\newcommand{\NSAMarginCoreSZeroV}{\ensuremath{2.116}}
% NSAMarginCoreSOneV [cert, margin] asm_design_N2ref4.json :: res.tight.margins.S1v_vs_0.6
\newcommand{\NSAMarginCoreSOneV}{\ensuremath{5.748}}
% NSAMarginCoreMOne [cert, margin] asm_design_N2ref4.json :: res.tight.margins.M1core
\newcommand{\NSAMarginCoreMOne}{\ensuremath{4.409}}
% NSAMarginCoreMZero [cert, margin] asm_design_N2ref4.json :: res.tight.margins.M0core
\newcommand{\NSAMarginCoreMZero}{\ensuremath{5.450}}
% NSAMarginGhat [cert, margin] asm_design_N2ref4.json :: res.tight.margins.Ghat22
\newcommand{\NSAMarginGhat}{\ensuremath{1.917}}
% NSAMarginRampMOne [cert, margin] asm_design_N2ref4.json :: res.tight.margins.ramp M1
\newcommand{\NSAMarginRampMOne}{\ensuremath{1.184}}
% NSAfourInfoZeroDWindowConst [cert, upper] asm_design_N2ref4.json :: res.tight.info.0.D_window_const
\newcommand{\NSAfourInfoZeroDWindowConst}{\ensuremath{1.7732\cdot10^{-7}}}
% NSAfourInfoOneDWindowConst [cert, upper] asm_design_N2ref4.json :: res.tight.info.1.D_window_const
\newcommand{\NSAfourInfoOneDWindowConst}{\ensuremath{9.8058\cdot10^{-5}}}
% NZDeltaRZero [cert, upper] zone3_interface.json :: DeltaR.0
\newcommand{\NZDeltaRZero}{\ensuremath{4.4986\cdot10^{-8}}}
% NZDeltaROne [cert, upper] zone3_interface.json :: DeltaR.1
\newcommand{\NZDeltaROne}{\ensuremath{6.7972\cdot10^{-6}}}
% NZDeltaRTwo [cert, upper] zone3_interface.json :: DeltaR.2
\newcommand{\NZDeltaRTwo}{\ensuremath{0.011480}}
% NZDeltaRBZero [cert, upper] zone3_interface.json :: DeltaR_b.0
\newcommand{\NZDeltaRBZero}{\ensuremath{2.66\cdot10^{-12}}}
% NZDeltaRBOne [cert, upper] zone3_interface.json :: DeltaR_b.1
\newcommand{\NZDeltaRBOne}{\ensuremath{3.26\cdot10^{-9}}}
% NZDeltaRBTwo [cert, upper] zone3_interface.json :: DeltaR_b.2
\newcommand{\NZDeltaRBTwo}{\ensuremath{4.35\cdot10^{-8}}}
% NZDeltaRDetailNSqZero [cert, upper] zone3_interface.json :: DeltaR_detail.N_sq.0
\newcommand{\NZDeltaRDetailNSqZero}{\ensuremath{8.28\cdot10^{-5}}}
% NZDeltaRDetailNSqOne [cert, upper] zone3_interface.json :: DeltaR_detail.N_sq.1
\newcommand{\NZDeltaRDetailNSqOne}{\ensuremath{0.10149}}
% NZDeltaRDetailNSqTwo [cert, upper] zone3_interface.json :: DeltaR_detail.N_sq.2
\newcommand{\NZDeltaRDetailNSqTwo}{\ensuremath{1.3567}}
% NZDeltaRDetailNSqThree [cert, upper] zone3_interface.json :: DeltaR_detail.N_sq.3
\newcommand{\NZDeltaRDetailNSqThree}{\ensuremath{25350}}
% NZDeltaRDetailPartsTwoPairsThree [cert, upper] zone3_interface.json :: DeltaR_detail.parts.2.pairs.3
\newcommand{\NZDeltaRDetailPartsTwoPairsThree}{\ensuremath{0.011298}}
% NZDeltaRDetailPartsTwoLFourFoot [cert, upper] zone3_interface.json :: DeltaR_detail.parts.2.L4_foot
\newcommand{\NZDeltaRDetailPartsTwoLFourFoot}{\ensuremath{1.6\cdot10^{-4}}}
% NZDeltaRDetailPartsTwoCubic [cert, upper] zone3_interface.json :: DeltaR_detail.parts.2.cubic
\newcommand{\NZDeltaRDetailPartsTwoCubic}{\ensuremath{9.3\cdot10^{-6}}}
% NZDeltaRDetailPartsTwoSThreeThreeFoot [cert, upper] zone3_interface.json :: DeltaR_detail.parts.2.S33_foot
\newcommand{\NZDeltaRDetailPartsTwoSThreeThreeFoot}{\ensuremath{3.7\cdot10^{-7}}}
% NZNIIIZero [cert, upper] zone3_interface.json :: n_III.0
\newcommand{\NZNIIIZero}{\ensuremath{0.071445}}
% NZNIIIOne [cert, upper] zone3_interface.json :: n_III.1
\newcommand{\NZNIIIOne}{\ensuremath{0.32707}}
% NZNIIITwo [cert, upper] zone3_interface.json :: n_III.2
\newcommand{\NZNIIITwo}{\ensuremath{3.0445}}
% NZTauIII [cert, upper] zone3_interface.json :: tau_III
\newcommand{\NZTauIII}{\ensuremath{4.0708}}
% NZoneIIIBoxZero [cert, upper] zone3_interface.json :: box_III_380.b0
\newcommand{\NZoneIIIBoxZero}{\ensuremath{0.10032}}
% NZoneIIIBoxOne [cert, upper] zone3_interface.json :: box_III_380.b1
\newcommand{\NZoneIIIBoxOne}{\ensuremath{0.41974}}
% NZoneIIIBoxZ [cert, upper] zone3_interface.json :: box_III_380.bz
\newcommand{\NZoneIIIBoxZ}{\ensuremath{0.014291}}
% NZBoxIIIThreeEightZeroMarginsBZero [cert, margin] zone3_interface.json :: box_III_380.margins.b0
\newcommand{\NZBoxIIIThreeEightZeroMarginsBZero}{\ensuremath{1.495}}
% NZoneIIIMarginThinnest [cert, margin] zone3_interface.json :: box_III_380.margins.b1
\newcommand{\NZoneIIIMarginThinnest}{\ensuremath{1.238}}
% NZBoxIIIThreeEightZeroMarginsBz [cert, margin] zone3_interface.json :: box_III_380.margins.bz
\newcommand{\NZBoxIIIThreeEightZeroMarginsBz}{\ensuremath{11.19}}
% NZBoxIIISevenZeroZeroBZero [cert, upper] zone3_interface.json :: box_III_700.b0
\newcommand{\NZBoxIIISevenZeroZeroBZero}{\ensuremath{0.22757}}
% NZBoxIIISevenZeroZeroBOne [cert, upper] zone3_interface.json :: box_III_700.b1
\newcommand{\NZBoxIIISevenZeroZeroBOne}{\ensuremath{0.54234}}
% NZUfineKTwoAtBhatRef [cert, upper] zone3_interface.json :: ufine.K2_at_Bhat_ref
\newcommand{\NZUfineKTwoAtBhatRef}{\ensuremath{0.046956}}
% NZUfineKThreeStripAtBhatRef [cert, upper] zone3_interface.json :: ufine.K3_strip_at_Bhat_ref
\newcommand{\NZUfineKThreeStripAtBhatRef}{\ensuremath{0.054140}}
% NRecoveryMinMargin [cert, margin] zone3_interface.json :: cone_family.H6.at_Bhat_ref.min_margin
\newcommand{\NRecoveryMinMargin}{\ensuremath{1.860}}
% NZConeFamilyHSixAtBhatRefConeSlopeMargin [cert, margin] zone3_interface.json :: cone_family.H6.at_Bhat_ref.cone.slope.margin
\newcommand{\NZConeFamilyHSixAtBhatRefConeSlopeMargin}{\ensuremath{2.46}}
% NZConeFamilyHSixAtBhatRefConeKTwoMargin [cert, margin] zone3_interface.json :: cone_family.H6.at_Bhat_ref.cone.K2.margin
\newcommand{\NZConeFamilyHSixAtBhatRefConeKTwoMargin}{\ensuremath{2.54}}
% NZConeFamilyHSixAtBhatRefConeKThreeMargin [cert, margin] zone3_interface.json :: cone_family.H6.at_Bhat_ref.cone.K3.margin
\newcommand{\NZConeFamilyHSixAtBhatRefConeKThreeMargin}{\ensuremath{2.51}}
% NConeTstarMargin [cert, margin] zone3_interface.json :: acceptance.tstar_margin
\newcommand{\NConeTstarMargin}{\ensuremath{2.30}}
% NUoneTFiveYaddLevelOneTwo [cert, upper] u1_ghat.json :: t5_Yadd.level1.2
\newcommand{\NUoneTFiveYaddLevelOneTwo}{\ensuremath{0.0044}}
% NUoneTFiveYaddLevelTwoTwo [cert, upper] u1_ghat.json :: t5_Yadd.level2.2
\newcommand{\NUoneTFiveYaddLevelTwoTwo}{\ensuremath{0.1081}}
% NUoneTFiveYaddLevelTwoThree [cert, upper] u1_ghat.json :: t5_Yadd.level2.3
\newcommand{\NUoneTFiveYaddLevelTwoThree}{\ensuremath{0.374}}
% NUoneWeightThetaGammaZero [cert, upper] u1_ghat.json :: weight.Theta.Gamma.0
\newcommand{\NUoneWeightThetaGammaZero}{\ensuremath{4.00000000000000000000054}}
% NRKZeroZero [cert, upper] ramp_U.json :: K.0.0
\newcommand{\NRKZeroZero}{\ensuremath{13436.03}}
% NRKOneOne [cert, upper] ramp_U.json :: K.1.1
\newcommand{\NRKOneOne}{\ensuremath{10526.19}}
% NRKTwoTwo [cert, upper] ramp_U.json :: K.2.2
\newcommand{\NRKTwoTwo}{\ensuremath{39482.30}}
% NUoneCasesPrimaryYNextZero [cert, upper] u1_ghat.json :: cases.primary.Y_next.0
\newcommand{\NUoneCasesPrimaryYNextZero}{\ensuremath{457.092}}
% NUoneCasesPrimaryYNextOne [cert, upper] u1_ghat.json :: cases.primary.Y_next.1
\newcommand{\NUoneCasesPrimaryYNextOne}{\ensuremath{356.940}}
% NUoneCasesPrimaryYNextTwo [cert, upper] u1_ghat.json :: cases.primary.Y_next.2
\newcommand{\NUoneCasesPrimaryYNextTwo}{\ensuremath{1320.220}}
% NSfourResultsMNineHTwoTightenedFUOneRowsZeroBound [cert, upper] FINAL_manifest_v33.json :: results.M9_H2_tightened_fU1.rows.0.bound
\newcommand{\NSfourResultsMNineHTwoTightenedFUOneRowsZeroBound}{\ensuremath{3.658\cdot10^{-10}}}
% NSfourResultsMNineHTwoTightenedFUOneRowsZeroMargin [cert, margin] FINAL_manifest_v33.json :: results.M9_H2_tightened_fU1.rows.0.margin
\newcommand{\NSfourResultsMNineHTwoTightenedFUOneRowsZeroMargin}{\ensuremath{27.34}}
% NSfourResultsMNineHTwoTightenedFUOneRowsOneBound [cert, upper] FINAL_manifest_v33.json :: results.M9_H2_tightened_fU1.rows.1.bound
\newcommand{\NSfourResultsMNineHTwoTightenedFUOneRowsOneBound}{\ensuremath{1.777\cdot10^{-9}}}
% NSfourResultsMNineHTwoTightenedFUOneRowsOneMargin [cert, margin] FINAL_manifest_v33.json :: results.M9_H2_tightened_fU1.rows.1.margin
\newcommand{\NSfourResultsMNineHTwoTightenedFUOneRowsOneMargin}{\ensuremath{562.8}}
% NSfourResultsMNineHTwoTightenedFUOneRowsTwoBound [cert, upper] FINAL_manifest_v33.json :: results.M9_H2_tightened_fU1.rows.2.bound
\newcommand{\NSfourResultsMNineHTwoTightenedFUOneRowsTwoBound}{\ensuremath{1.665\cdot10^{-6}}}
% NSfourResultsMNineHTwoTightenedFUOneRowsTwoMargin [cert, margin] FINAL_manifest_v33.json :: results.M9_H2_tightened_fU1.rows.2.margin
\newcommand{\NSfourResultsMNineHTwoTightenedFUOneRowsTwoMargin}{\ensuremath{600.9}}
% NHOneAlphaThr [hand, exact] hand :: 
\newcommand{\NHOneAlphaThr}{\ensuremath{1.5\cdot10^{-4}}}
% NHOneSqrtLamThr [hand, exact] hand :: 
\newcommand{\NHOneSqrtLamThr}{\ensuremath{5\cdot10^{-8}}}
% NYbarStep [hand, exact] hand :: 
\newcommand{\NYbarStep}{\ensuremath{1\cdot10^{-12}}}
% NHTwoThrZero [hand, exact] hand :: 
\newcommand{\NHTwoThrZero}{\ensuremath{1\cdot10^{-8}}}
% NHTwoThrOne [hand, exact] hand :: 
\newcommand{\NHTwoThrOne}{\ensuremath{1\cdot10^{-6}}}
% NHTwoThrTwo [hand, exact] hand :: 
\newcommand{\NHTwoThrTwo}{\ensuremath{1\cdot10^{-3}}}
% NHThreeThrZero [hand, exact] hand :: 
\newcommand{\NHThreeThrZero}{\ensuremath{1\cdot10^{-6}}}
% NHThreeThrOne [hand, exact] hand :: 
\newcommand{\NHThreeThrOne}{\ensuremath{5\cdot10^{-6}}}
% NHThreeThrTwo [hand, exact] hand :: 
\newcommand{\NHThreeThrTwo}{\ensuremath{1\cdot10^{-5}}}
% NHFourThrZero [hand, exact] hand :: 
\newcommand{\NHFourThrZero}{\ensuremath{2\cdot10^{-7}}}
% NHFourThrOne [hand, exact] hand :: 
\newcommand{\NHFourThrOne}{\ensuremath{1\cdot10^{-5}}}
% NHFourThrTwo [hand, exact] hand :: 
\newcommand{\NHFourThrTwo}{\ensuremath{1\cdot10^{-5}}}
% NHFourThrThree [hand, exact] hand :: 
\newcommand{\NHFourThrThree}{\ensuremath{1\cdot10^{-5}}}
% NHFourKfourThr [hand, exact] hand :: 
\newcommand{\NHFourKfourThr}{\ensuremath{0.06}}
% NHSixConeSlopeThr [hand, exact] hand :: 
\newcommand{\NHSixConeSlopeThr}{\ensuremath{0.08}}
% NHSixConeKtwoThr [hand, exact] hand :: 
\newcommand{\NHSixConeKtwoThr}{\ensuremath{0.03}}
% NHSixConeNthreeThr [hand, exact] hand :: 
\newcommand{\NHSixConeNthreeThr}{\ensuremath{0.004}}
% NHSixConeJetThr [hand, exact] hand :: 
\newcommand{\NHSixConeJetThr}{\ensuremath{0.12}}
% NHSixConeHFiveThr [hand, exact] hand :: 
\newcommand{\NHSixConeHFiveThr}{\ensuremath{0.0144}}
% NHSixFarHFiveThr [hand, exact] hand :: 
\newcommand{\NHSixFarHFiveThr}{\ensuremath{0.0225}}
% NLemmaOConeM [hand, exact] hand :: 
\newcommand{\NLemmaOConeM}{\ensuremath{0.12}}
% NLemmaOConeK [hand, exact] hand :: 
\newcommand{\NLemmaOConeK}{\ensuremath{0.18}}
% NLemmaOFarM [hand, exact] hand :: 
\newcommand{\NLemmaOFarM}{\ensuremath{0.15}}
% NLemmaOFarK [hand, exact] hand :: 
\newcommand{\NLemmaOFarK}{\ensuremath{0.3}}
% NCoercivityConst [hand, exact] hand :: 
\newcommand{\NCoercivityConst}{\ensuremath{0.046875}}
% NMOnez [hand, exact] hand :: 
\newcommand{\NMOnez}{\ensuremath{0.534}}
% NBetaA [hand, exact] hand :: 
\newcommand{\NBetaA}{\ensuremath{2}}
% NTOneLower [cert, lower] t1clock.json :: t1.0
\newcommand{\NTOneLower}{\ensuremath{4.5978574775}}
% NTOneUpper [cert, upper] t1clock.json :: t1.1
\newcommand{\NTOneUpper}{\ensuremath{4.5978586144}}
% NLambdaOneLower [cert, lower] t1clock.json :: lnlam.0
\newcommand{\NLambdaOneLower}{\ensuremath{6.49\cdot10^{-5}}}
% NLambdaOneUpper [cert, upper] t1clock.json :: lnlam.1
\newcommand{\NLambdaOneUpper}{\ensuremath{6.75\cdot10^{-5}}}
% NTMinusTOneUpper [cert, upper] t1clock.json :: lnlam.1; 1.5e-4.sup_ab
\newcommand{\NTMinusTOneUpper}{\ensuremath{2.1\cdot10^{-17}}}
% NAbarUpper [cert, upper] final_alpha_points_5e-8.json :: 1.5e-4.sup_ab
\newcommand{\NAbarUpper}{\ensuremath{2.19103\cdot10^{-7}}}
% NBetabarUpper [cert, upper] final_alpha_points_5e-8.json :: 1.5e-4.sup_beta
\newcommand{\NBetabarUpper}{\ensuremath{3.33811\cdot10^{-9}}}
% NBaseKmaxLower [cert, lower] t1clock.json :: base_kmax.0
\newcommand{\NBaseKmaxLower}{\ensuremath{2.279\cdot10^{4}}}
% NBaseKmaxUpper [cert, upper] t1clock.json :: base_kmax.1
\newcommand{\NBaseKmaxUpper}{\ensuremath{2.376\cdot10^{4}}}
% NHOneAlphaBase [cert, upper] summary.json :: alpha_abs_max
\newcommand{\NHOneAlphaBase}{\ensuremath{1.1739\cdot10^{-4}}}
% NHOneAlphaErr [cert, upper] t9_h1_clam_s180.json :: C_alpha; inputs.YSTAR
\newcommand{\NHOneAlphaErr}{\ensuremath{7.21\cdot10^{-13}}}
% NHOneAlphaMargin [cert, margin] summary.json :: alpha_abs_max; C_alpha; inputs.YSTAR
\newcommand{\NHOneAlphaMargin}{\ensuremath{1.27}}
% NHOneSqrtLamBase [cert, upper] t9_h1_clam_s180.json :: sqrtLam_base_Wp
\newcommand{\NHOneSqrtLamBase}{\ensuremath{2.282\cdot10^{-8}}}
% NHOneSqrtLamTotal [cert, upper] t9_h1_clam_s180.json :: sqrtLam_total_Wp
\newcommand{\NHOneSqrtLamTotal}{\ensuremath{2.29\cdot10^{-8}}}
% NHOneSqrtLamMargin [cert, margin] t9_h1_clam_s180.json :: sqrtLam_total_Wp
\newcommand{\NHOneSqrtLamMargin}{\ensuremath{2.18}}
% NHOneYbarLimit [cert, lower] t9_h1_clam_s180.json :: Ybar_limit_H1
\newcommand{\NHOneYbarLimit}{\ensuremath{1.12\cdot10^{-9}}}
% NSwitchYUpper [cert, upper] switch_s2.json :: modulated.Y_raw
\newcommand{\NSwitchYUpper}{\ensuremath{6.21\cdot10^{-13}}}
% NSwitchYMargin [cert, margin] switch_s2.json :: modulated.Y_raw
\newcommand{\NSwitchYMargin}{\ensuremath{1.6}}
% NMinCI [cert, lower] constants.json :: meta.summary.min_c_I
\newcommand{\NMinCI}{\ensuremath{0.05283}}
% NNumCertificates [cert, exact] constants.json :: meta.summary.n_cert
\newcommand{\NNumCertificates}{\ensuremath{26}}
% NMinRhoC [cert, lower] constants.json :: meta.closure.rho_c_min
\newcommand{\NMinRhoC}{\ensuremath{0.04711}}
% NHTwoBaseZero [cert, upper] FINAL_manifest_v33.json :: results.M9_H2_tightened.rows.0.base
\newcommand{\NHTwoBaseZero}{\ensuremath{8.011\cdot10^{-11}}}
% NHTwoBaseOne [cert, upper] FINAL_manifest_v33.json :: results.M9_H2_tightened.rows.1.base
\newcommand{\NHTwoBaseOne}{\ensuremath{1.336\cdot10^{-10}}}
% NHTwoBaseTwo [cert, upper] FINAL_manifest_v33.json :: results.M9_H2_tightened.rows.2.base
\newcommand{\NHTwoBaseTwo}{\ensuremath{2.936\cdot10^{-10}}}
% NHTwoErrWZero [cert, upper] FINAL_manifest_v33.json :: results.M9_H2_tightened.rows.0.with_error
\newcommand{\NHTwoErrWZero}{\ensuremath{8.014\cdot10^{-11}}}
% NHTwoMarginWZero [cert, margin] FINAL_manifest_v33.json :: results.M9_H2_tightened.rows.0.margin
\newcommand{\NHTwoMarginWZero}{\ensuremath{124.78}}
% NHTwoErrWOne [cert, upper] FINAL_manifest_v33.json :: results.M9_H2_tightened.rows.1.with_error
\newcommand{\NHTwoErrWOne}{\ensuremath{4.706\cdot10^{-10}}}
% NHTwoMarginWOne [cert, margin] FINAL_manifest_v33.json :: results.M9_H2_tightened.rows.1.margin
\newcommand{\NHTwoMarginWOne}{\ensuremath{2124.9}}
% NHTwoErrWTwo [cert, upper] FINAL_manifest_v33.json :: results.M9_H2_tightened.rows.2.with_error
\newcommand{\NHTwoErrWTwo}{\ensuremath{1.559\cdot10^{-6}}}
% NHTwoMarginWTwo [cert, margin] FINAL_manifest_v33.json :: results.M9_H2_tightened.rows.2.margin
\newcommand{\NHTwoMarginWTwo}{\ensuremath{641.53}}
% NHThreeBaseZero [cert, upper] FINAL_manifest_v33.json :: results.M9_H2f.rows.0.base
\newcommand{\NHThreeBaseZero}{\ensuremath{3.127\cdot10^{-7}}}
% NHThreeBaseOne [cert, upper] FINAL_manifest_v33.json :: results.M9_H2f.rows.1.base
\newcommand{\NHThreeBaseOne}{\ensuremath{1.004\cdot10^{-6}}}
% NHThreeBaseTwo [cert, upper] FINAL_manifest_v33.json :: results.M9_H2f.rows.2.base
\newcommand{\NHThreeBaseTwo}{\ensuremath{2.431\cdot10^{-6}}}
% NHThreeErrZero [cert, upper] FINAL_manifest_v33.json :: results.M9_H2f.rows.0.with_error
\newcommand{\NHThreeErrZero}{\ensuremath{3.1363\cdot10^{-7}}}
% NHThreeMarginZero [cert, margin] FINAL_manifest_v33.json :: results.M9_H2f.rows.0.margin
\newcommand{\NHThreeMarginZero}{\ensuremath{3.1885}}
% NHThreeErrOne [cert, upper] FINAL_manifest_v33.json :: results.M9_H2f.rows.1.with_error
\newcommand{\NHThreeErrOne}{\ensuremath{1.0034\cdot10^{-6}}}
% NHThreeMarginOne [cert, margin] FINAL_manifest_v33.json :: results.M9_H2f.rows.1.margin
\newcommand{\NHThreeMarginOne}{\ensuremath{4.9834}}
% NHThreeErrTwo [cert, upper] FINAL_manifest_v33.json :: results.M9_H2f.rows.2.with_error
\newcommand{\NHThreeErrTwo}{\ensuremath{2.4351\cdot10^{-6}}}
% NHThreeMarginTwo [cert, margin] FINAL_manifest_v33.json :: results.M9_H2f.rows.2.margin
\newcommand{\NHThreeMarginTwo}{\ensuremath{4.1067}}
% NHFourBaseZero [cert, upper] FINAL_manifest_v33.json :: results.M9_base_sums.combined.windows_max.0
\newcommand{\NHFourBaseZero}{\ensuremath{3.48\cdot10^{-8}}}
% NHFourErrZero [cert, upper] FINAL_manifest_v33.json :: results.M9_H3_H2v32.rows.win.0.with_error_arb
\newcommand{\NHFourErrZero}{\ensuremath{3.5087\cdot10^{-8}}}
% NHFourMarginZero [cert, margin] FINAL_manifest_v33.json :: results.M9_H3_H2v32.rows.win.0.margin_arb
\newcommand{\NHFourMarginZero}{\ensuremath{5.700}}
% NHFourBaseOne [cert, upper] FINAL_manifest_v33.json :: results.M9_base_sums.combined.windows_max.1
\newcommand{\NHFourBaseOne}{\ensuremath{1.27\cdot10^{-7}}}
% NHFourErrOne [cert, upper] FINAL_manifest_v33.json :: results.M9_H3_H2v32.rows.win.1.with_error_arb
\newcommand{\NHFourErrOne}{\ensuremath{1.2616\cdot10^{-7}}}
% NHFourMarginOne [cert, margin] FINAL_manifest_v33.json :: results.M9_H3_H2v32.rows.win.1.margin_arb
\newcommand{\NHFourMarginOne}{\ensuremath{79.26}}
% NHFourBaseTwo [cert, upper] FINAL_manifest_v33.json :: results.M9_base_sums.combined.windows_max.2
\newcommand{\NHFourBaseTwo}{\ensuremath{3.76\cdot10^{-7}}}
% NHFourErrTwo [cert, upper] FINAL_manifest_v33.json :: results.M9_H3_H2v32.rows.win.2.with_error_arb
\newcommand{\NHFourErrTwo}{\ensuremath{3.7724\cdot10^{-7}}}
% NHFourMarginTwo [cert, margin] FINAL_manifest_v33.json :: results.M9_H3_H2v32.rows.win.2.margin_arb
\newcommand{\NHFourMarginTwo}{\ensuremath{26.50}}
% NHFourBaseThree [cert, upper] FINAL_manifest_v33.json :: results.M9_base_sums.combined.windows_max.3
\newcommand{\NHFourBaseThree}{\ensuremath{9.71\cdot10^{-7}}}
% NHFourErrThree [cert, upper] FINAL_manifest_v33.json :: results.M9_H3_H2v32.rows.win.3.with_error_arb
\newcommand{\NHFourErrThree}{\ensuremath{1.2822\cdot10^{-6}}}
% NHFourMarginThree [cert, margin] FINAL_manifest_v33.json :: results.M9_H3_H2v32.rows.win.3.margin_arb
\newcommand{\NHFourMarginThree}{\ensuremath{7.799}}
% NHFourYbarLimit [cert, lower] FINAL_manifest_v33.json :: results.M9_H3_H2v32.binding_Ybar_limit
\newcommand{\NHFourYbarLimit}{\ensuremath{1.575\cdot10^{-11}}}
% NHFourYbarMargin [cert, margin] FINAL_manifest_v33.json :: results.M9_H3_H2v32.binding_Ybar_limit
\newcommand{\NHFourYbarMargin}{\ensuremath{15.75}}
% NHFourKfourBase [cert, upper] amp_A1_w180.json :: inputs.BASE_T1.K4
\newcommand{\NHFourKfourBase}{\ensuremath{3.1633\cdot10^{-4}}}
% NHSixConeSlope [cert, upper] FINAL_manifest_v33.json :: results.cone_fixedL.checks.slope.value
\newcommand{\NHSixConeSlope}{\ensuremath{0.03029}}
% NHSixConeSlopeMargin [cert, margin] FINAL_manifest_v33.json :: results.cone_fixedL.checks.slope.margin
\newcommand{\NHSixConeSlopeMargin}{\ensuremath{2.641}}
% NHSixConeKtwo [cert, upper] FINAL_manifest_v33.json :: results.cone_fixedL.checks.K2.value
\newcommand{\NHSixConeKtwo}{\ensuremath{0.009876}}
% NHSixConeKtwoMargin [cert, margin] FINAL_manifest_v33.json :: results.cone_fixedL.checks.K2.margin
\newcommand{\NHSixConeKtwoMargin}{\ensuremath{3.037}}
% NHSixConeNthree [cert, upper] FINAL_manifest_v33.json :: results.cone_fixedL.checks.N3.value
\newcommand{\NHSixConeNthree}{\ensuremath{1.461\cdot10^{-4}}}
% NHSixConeNthreeMargin [cert, margin] FINAL_manifest_v33.json :: results.cone_fixedL.checks.N3.margin
\newcommand{\NHSixConeNthreeMargin}{\ensuremath{27.38}}
% NHSixConeHFive [cert, upper] FINAL_manifest_v33.json :: results.H4_fixedL.total
\newcommand{\NHSixConeHFive}{\ensuremath{1.2865\cdot10^{-5}}}
% NHSixConeHFiveMargin [cert, margin] FINAL_manifest_v33.json :: results.H4_fixedL.margin
\newcommand{\NHSixConeHFiveMargin}{\ensuremath{1119}}
% NHSixConeTstarNeeded [cert, upper] FINAL_manifest_v33.json :: results.H7_fixedL.ratio_upper
\newcommand{\NHSixConeTstarNeeded}{\ensuremath{1.6000015\cdot10^{-7}}}
% NHSixConeTstarLower [cert, lower] FINAL_manifest_v33.json :: results.H7_fixedL.t_star_lower
\newcommand{\NHSixConeTstarLower}{\ensuremath{3.68334\cdot10^{-7}}}
% NHSixConeTstarMargin [cert, margin] FINAL_manifest_v33.json :: results.H7_fixedL.margin_tstar_over_ratio
\newcommand{\NHSixConeTstarMargin}{\ensuremath{2.302}}
% NHSixConeHFiveRatioPeak [cert, upper] h4loc_A1_w180.json :: windows.constants/amp_A1_w180.json.worst.ratio
\newcommand{\NHSixConeHFiveRatioPeak}{\ensuremath{0.2393}}
% NHSixConeHOneSup [cert, upper] FINAL_manifest_v33.json :: results.H1_fixedL.sup_h0_max
\newcommand{\NHSixConeHOneSup}{\ensuremath{0.2163}}
% NHSixConeHOneH [cert, upper] FINAL_manifest_v33.json :: results.H1_fixedL.H_max
\newcommand{\NHSixConeHOneH}{\ensuremath{0.2185}}
% NHSixFarHFiveRatio [cert, upper] h4loc_v32.json :: windows.constants/amp_A1_w180.json.cylinders.18.ratio
\newcommand{\NHSixFarHFiveRatio}{\ensuremath{0.1761}}
% NHSixFarHFiveMarginOld [cert, margin] h4loc_v32.json :: windows.constants/amp_A1_w180.json.cylinders.18.ratio
\newcommand{\NHSixFarHFiveMarginOld}{\ensuremath{5.68}}
% NHSixFarHFiveMargin [cert, margin] FINAL_manifest_v33.json :: results.M9_H4_far.min_margin
\newcommand{\NHSixFarHFiveMargin}{\ensuremath{15.58}}
% NHSixConeHFiveMarginOld [cert, margin] h4loc_A1_w180.json :: windows.constants/amp_A1_w180.json.worst.ratio
\newcommand{\NHSixConeHFiveMarginOld}{\ensuremath{4.18}}
% NHSixFarKtwo [cert, upper] t9_far_v32.json :: h2
\newcommand{\NHSixFarKtwo}{\ensuremath{1.797\cdot10^{-4}}}
% NHSixFarKtwoThr [hand, exact] t9_far_v32.json :: thr.h2
\newcommand{\NHSixFarKtwoThr}{\ensuremath{3.53\cdot10^{-4}}}
% NHSixFarKtwoMargin [cert, margin] t9_far_v32.json :: margins.h2
\newcommand{\NHSixFarKtwoMargin}{\ensuremath{1.964}}
% NHSixFarKthree [cert, upper] t9_far_v32.json :: h3
\newcommand{\NHSixFarKthree}{\ensuremath{2.595\cdot10^{-8}}}
% NHSixFarKthreeThr [hand, exact] t9_far_v32.json :: thr.h3
\newcommand{\NHSixFarKthreeThr}{\ensuremath{1.17\cdot10^{-5}}}
% NHSixFarKthreeMargin [cert, margin] t9_far_v32.json :: margins.h3
\newcommand{\NHSixFarKthreeMargin}{\ensuremath{450.9}}
% NHSixFarKfour [cert, upper] t9_far_v32.json :: h4
\newcommand{\NHSixFarKfour}{\ensuremath{6.964\cdot10^{-9}}}
% NHSixFarKfourThr [hand, exact] t9_far_v32.json :: thr.h4
\newcommand{\NHSixFarKfourThr}{\ensuremath{1.03\cdot10^{-7}}}
% NHSixFarKfourMargin [cert, margin] t9_far_v32.json :: margins.h4
\newcommand{\NHSixFarKfourMargin}{\ensuremath{14.79}}
% NHSixFarSlope [cert, upper] t9_far_v32.json :: slope
\newcommand{\NHSixFarSlope}{\ensuremath{0.04153}}
% NHSixFarSlopeThr [hand, exact] t9_far_v32.json :: thr.slope
\newcommand{\NHSixFarSlopeThr}{\ensuremath{0.0800}}
% NHSixFarSlopeMargin [cert, margin] t9_far_v32.json :: margins.slope
\newcommand{\NHSixFarSlopeMargin}{\ensuremath{1.926}}
% NHSixFarNthree [cert, upper] t9_far_v32.json :: L2
\newcommand{\NHSixFarNthree}{\ensuremath{5.516\cdot10^{-7}}}
% NHSixFarNthreeThr [hand, exact] t9_far_v32.json :: thr.L2
\newcommand{\NHSixFarNthreeThr}{\ensuremath{3.32\cdot10^{-6}}}
% NHSixFarNthreeMargin [cert, margin] t9_far_v32.json :: margins.L2
\newcommand{\NHSixFarNthreeMargin}{\ensuremath{6.019}}
% NFarTstarMargin [cert, margin] lemmaO_v32_consts.json :: far_v32.case.margin_factor_lower
\newcommand{\NFarTstarMargin}{\ensuremath{15.5}}
% NFarTstarLower [cert, lower] lemmaO_v32_consts.json :: far_v32.case.t_star_lower
\newcommand{\NFarTstarLower}{\ensuremath{9.5108\cdot10^{-8}}}
% NFarTstarNeeded [cert, upper] lemmaO_v32_consts.json :: far_v32.case.needed
\newcommand{\NFarTstarNeeded}{\ensuremath{6.134\cdot10^{-9}}}
% NGammaModes [cert, exact] gamma0_final.json :: M
\newcommand{\NGammaModes}{\ensuremath{3762}}
% NGammaHalfwidth [cert, upper] gamma0_final.json :: b_j_halfwidth
\newcommand{\NGammaHalfwidth}{\ensuremath{8.138\cdot10^{-24}}}
% NGammaContraction [cert, upper] gamma0_final.json :: krawczyk_contraction
\newcommand{\NGammaContraction}{\ensuremath{3.01\cdot10^{-12}}}
% NGammaArcs [cert, exact] gamma0_final.json :: embedded.arcs
\newcommand{\NGammaArcs}{\ensuremath{701}}
% NGammaMaxTurn [cert, upper] gamma0_final.json :: embedded.maxturn
\newcommand{\NGammaMaxTurn}{\ensuremath{0.04997}}
% NGammaMaxPair [cert, upper] gamma0_final.json :: embedded.maxpair
\newcommand{\NGammaMaxPair}{\ensuremath{0.09597}}
% NGammaMinSep [cert, lower] gamma0_final.json :: embedded.minsep
\newcommand{\NGammaMinSep}{\ensuremath{8.89\cdot10^{-3}}}
% NGammaHalfGap [cert, lower] gamma0_final.json :: embedded.gap_cert
\newcommand{\NGammaHalfGap}{\ensuremath{0.0248219}}
% NPhaseISlabs [cert, exact] status.json :: slab
\newcommand{\NPhaseISlabs}{\ensuremath{96}}
% NPhaseIModes [cert, exact] status.json :: last_slab.M
\newcommand{\NPhaseIModes}{\ensuremath{2800}}
% NPhaseIKmaxLower [cert, lower] status.json :: kmax_enclosure.0
\newcommand{\NPhaseIKmaxLower}{\ensuremath{29.352874545}}
% NPhaseIKmaxUpper [cert, upper] status.json :: kmax_enclosure.1
\newcommand{\NPhaseIKmaxUpper}{\ensuremath{29.352874546}}
% NPhaseITLower [cert, lower] status.json :: t
\newcommand{\NPhaseITLower}{\ensuremath{4.597850539485830}}
% NPhaseITUpper [cert, upper] status.json :: t
\newcommand{\NPhaseITUpper}{\ensuremath{4.597850539485831}}
% NPeakSqrtLam [cert, upper] rate_table_fin12f4_a0.00015.json :: sqrtL_max
\newcommand{\NPeakSqrtLam}{\ensuremath{5.308\cdot10^{-8}}}
% NTailRate [cert, lower] rate_table_fin12f4_a0.00015.json :: tail.rho
\newcommand{\NTailRate}{\ensuremath{2.693\cdot10^{-3}}}
% NFloorSqrtLam [cert, upper] rate_table_fin12f4_a0.00015.json :: tail.sqrt_LF
\newcommand{\NFloorSqrtLam}{\ensuremath{9.872\cdot10^{-10}}}
% NAlphaFloor [cert, upper] rate_table_fin12f4_a0.00015.json :: tail.alpha_F
\newcommand{\NAlphaFloor}{\ensuremath{6.867\cdot10^{-7}}}
% NSupBetaVerifier [cert, upper] final_alpha_points_5e-8.json :: 1.5e-4.sup_beta
\newcommand{\NSupBetaVerifier}{\ensuremath{3.33811\cdot10^{-9}}}
% NMOneRamp [cert, upper] final_alpha_points_5e-8.json :: 1.5e-4.M1
\newcommand{\NMOneRamp}{\ensuremath{0.4328}}
% NSITwoPrime [cert, upper] FINAL_manifest_v33.json :: results.F12_sliver.sI2_prime
\newcommand{\NSITwoPrime}{\ensuremath{0.015457}}
% NLoopGainCoarse [cert, upper] U_cert_v35.json :: loop_gain.2.total
\newcommand{\NLoopGainCoarse}{\ensuremath{0.03904}}
% NRateMarginCoarse [cert, margin] U_cert_v35.json :: rate_margin
\newcommand{\NRateMarginCoarse}{\ensuremath{28.2}}
% NSOneEpsLower [cert, lower] bracket2.json :: eps_a
\newcommand{\NSOneEpsLower}{\ensuremath{1.5346198892965\cdot10^{-4}}}
% NSOneEpsUpper [cert, upper] bracket2.json :: eps_b
\newcommand{\NSOneEpsUpper}{\ensuremath{1.5346198899299\cdot10^{-4}}}
% NSOneRowsTip [cert, exact] ext_tip_check.json :: rows
\newcommand{\NSOneRowsTip}{\ensuremath{51246}}
% NSOneTipWidth [cert, upper] ext_tip_check.json :: max_pointwise_width.th
\newcommand{\NSOneTipWidth}{\ensuremath{3.2\cdot10^{-12}}}
% NSOneExtWidth [cert, upper] ext813_check.json :: max_pointwise_width.p
\newcommand{\NSOneExtWidth}{\ensuremath{5.7\cdot10^{-10}}}
% NSOneRowsMC [cert, exact] check_outA5.json :: prof_tip_vs_MC.rows_compared
\newcommand{\NSOneRowsMC}{\ensuremath{1025}}
% NSOneRowsMB [cert, exact] check_outA5.json :: ext813_vs_MB.rows_compared
\newcommand{\NSOneRowsMB}{\ensuremath{501}}
% NSOneMinMargin [cert, margin] check_outA5.json :: min:*.min_relative_margin.*
\newcommand{\NSOneMinMargin}{\ensuremath{1.2\cdot10^{-9}}}
% NTwoSegments [cert, exact] contour_result.json :: len:segments
\newcommand{\NTwoSegments}{\ensuremath{5058}}
% NTwoInteriorBoxes [cert, exact] interior_result.json :: len:
\newcommand{\NTwoInteriorBoxes}{\ensuremath{174}}
% NTwoContourBoxes [cert, exact] eta_cache.json :: len:
\newcommand{\NTwoContourBoxes}{\ensuremath{202}}
% NTwoEtaMax [cert, upper] eta_cache.json :: max:*.eta
\newcommand{\NTwoEtaMax}{\ensuremath{0.1163}}
% NTwoLipMax [cert, upper] eta_cache.json :: max:*.Lip
\newcommand{\NTwoLipMax}{\ensuremath{0.272}}
% NTwoLambdaZero [est, upper] cert_b0.25_th0.0_tau0.03.json :: Lam
\newcommand{\NTwoLambdaZero}{\ensuremath{1.790819}}
% NTwoLambdaOne [est, upper] cert_b0.25_th0.48_tau0.03.json :: Lam
\newcommand{\NTwoLambdaOne}{\ensuremath{1.768282}}
% NSOneDTightLower [doc, lower] doc:README.md :: D ∈ [2.50080627654, 2.50080627668]
\newcommand{\NSOneDTightLower}{\ensuremath{2.50080627654}}
% NSOneDTightUpper [doc, upper] doc:README.md :: D ∈ [2.50080627654, 2.50080627668]
\newcommand{\NSOneDTightUpper}{\ensuremath{2.50080627668}}
% NSOneKTightLower [doc, lower] doc:README.md :: k(0) ∈ [1.53982063762, 1.53982063822]
\newcommand{\NSOneKTightLower}{\ensuremath{1.53982063762}}
% NSOneKTightUpper [doc, upper] doc:README.md :: k(0) ∈ [1.53982063762, 1.53982063822]
\newcommand{\NSOneKTightUpper}{\ensuremath{1.53982063822}}
% NSOneDCondLower [doc, lower] doc:README.md :: D = −p(tip) ∈ [2.500806276619695, 2.500806276620049]
\newcommand{\NSOneDCondLower}{\ensuremath{2.500806276619695}}
% NSOneDCondUpper [doc, upper] doc:README.md :: D = −p(tip) ∈ [2.500806276619695, 2.500806276620049]
\newcommand{\NSOneDCondUpper}{\ensuremath{2.500806276620049}}
% NSOneKCondLower [doc, lower] doc:README.md :: k(0) ∈ [1.539820637871836, 1.5398206378721015]
\newcommand{\NSOneKCondLower}{\ensuremath{1.539820637871836}}
% NSOneKCondUpper [doc, upper] doc:README.md :: k(0) ∈ [1.539820637871836, 1.5398206378721015]
\newcommand{\NSOneKCondUpper}{\ensuremath{1.5398206378721015}}
% NSOneTauTipLower [doc, lower] doc:README.md :: τ_tip ∈ [10.249077651120778, 10.249077651121361]
\newcommand{\NSOneTauTipLower}{\ensuremath{10.249077651120778}}
% NSOneTauTipUpper [doc, upper] doc:README.md :: τ_tip ∈ [10.249077651120778, 10.249077651121361]
\newcommand{\NSOneTauTipUpper}{\ensuremath{10.249077651121361}}
% NSOneStateUncert [doc, upper] doc:README.md :: State uncertainty at σ = 8: 4.4e−13
\newcommand{\NSOneStateUncert}{\ensuremath{4.4\cdot10^{-13}}}
% NFarBoxEprime [doc, margin] doc:STEP4_RESULT.md :: (6.4287, 4.6917, 1.9940, 1.69999995)
\newcommand{\NFarBoxEprime}{\ensuremath{6.42}}
% NFarBoxXEsecond [doc, margin] doc:STEP4_RESULT.md :: (6.4287, 4.6917, 1.9940, 1.69999995)
\newcommand{\NFarBoxXEsecond}{\ensuremath{4.69}}
% NFarBoxXsqEthird [doc, margin] doc:STEP4_RESULT.md :: (6.4287, 4.6917, 1.9940, 1.69999995)
\newcommand{\NFarBoxXsqEthird}{\ensuremath{1.99}}
% NFarBoxE [doc, margin] doc:STEP4_RESULT.md :: (6.4287, 4.6917, 1.9940, 1.69999995)
\newcommand{\NFarBoxE}{\ensuremath{1.69}}
% NHFourKfourErr [doc, upper] doc:STEP4_RESULT.md :: | 5.60e-3 against 0.06 | ×10.7 |
\newcommand{\NHFourKfourErr}{\ensuremath{5.60\cdot10^{-3}}}
% NHFourKfourMargin [doc, margin] doc:STEP4_RESULT.md :: | 5.60e-3 against 0.06 | ×10.7 |
\newcommand{\NHFourKfourMargin}{\ensuremath{10.7}}
% NCoercivityCoefficient [doc, upper] doc:step3_proof_v32.tex :: \norm{e^{\tau\mathcal L}u}_\rho\le14.8\,e^{-0.0213\tau}\norm
\newcommand{\NCoercivityCoefficient}{\ensuremath{14.8}}
% NCoercivityRate [doc, lower] doc:step3_proof_v32.tex :: \norm{e^{\tau\mathcal L}u}_\rho\le14.8\,e^{-0.0213\tau}\norm
\newcommand{\NCoercivityRate}{\ensuremath{0.0213}}
% NRegionRampStart [hand, exact] hand :: 
\newcommand{\NRegionRampStart}{\ensuremath{14}}
% NRegionInterface [hand, exact] hand :: 
\newcommand{\NRegionInterface}{\ensuremath{28}}
% NRegionShellEnd [hand, exact] hand :: 
\newcommand{\NRegionShellEnd}{\ensuremath{60}}
% NRegionSegEnd [hand, exact] hand :: 
\newcommand{\NRegionSegEnd}{\ensuremath{350}}
% NRegionAnnEnd [hand, exact] hand :: 
\newcommand{\NRegionAnnEnd}{\ensuremath{700}}
% NZoneIIILow [hand, exact] hand :: 
\newcommand{\NZoneIIILow}{\ensuremath{253}}
% NZoneIIIHigh [hand, exact] hand :: 
\newcommand{\NZoneIIIHigh}{\ensuremath{379.5}}
% NConeCentreMin [hand, exact] hand :: 
\newcommand{\NConeCentreMin}{\ensuremath{300}}
% NConeCentreMax [hand, exact] hand :: 
\newcommand{\NConeCentreMax}{\ensuremath{380}}
% NFarCentreMin [hand, exact] hand :: 
\newcommand{\NFarCentreMin}{\ensuremath{490}}
% NConeL [hand, exact] hand :: 
\newcommand{\NConeL}{\ensuremath{50}}
% NConeX [hand, exact] hand :: 
\newcommand{\NConeX}{\ensuremath{3}}
% NFarL [hand, exact] hand :: 
\newcommand{\NFarL}{\ensuremath{113}}
% NFarX [hand, exact] hand :: 
\newcommand{\NFarX}{\ensuremath{2}}
 is found.
% ======================================================================

% ---- which document are we in ----------------------------------------
\ifdefined\ifsupplement\else\newif\ifsupplement\fi

% ---- packages ---------------------------------------------------------
\usepackage[T1]{fontenc}
\usepackage[utf8]{inputenc}
\usepackage{lmodern}
\usepackage[margin=1.1in]{geometry}
\usepackage{amsmath,amssymb,amsthm,mathtools,bm}
\usepackage{booktabs,array,longtable,tabularx,multirow}
\usepackage{enumitem}
\usepackage{graphicx}
\usepackage{xcolor}
\usepackage{tikz}
\usetikzlibrary{positioning,arrows.meta,calc}
% xr-hyper must come before hyperref
\usepackage{xr-hyper}
\usepackage[colorlinks=true,linkcolor=blue!50!black,citecolor=green!35!black,urlcolor=blue!50!black]{hyperref}

\raggedbottom
\setlength{\emergencystretch}{2em}
\setlist[enumerate]{leftmargin=*,itemsep=2pt,topsep=3pt}
\setlist[itemize]{leftmargin=*,itemsep=2pt,topsep=3pt}

% ---- numbering --------------------------------------------------------
% Main paper: sections 1,2,...; theorems 2.3 (by section); equations (2.3).
% Supplement: the Parts A..K are the \section's, numbered by letter, so
% theorems are A.1, E.4, ... (do not confuse "Theorem A" of the paper, which
% has no dot, with "Theorem A.1" of Part A).
\ifsupplement
  \renewcommand{\thesection}{\Alph{section}}
\fi
\numberwithin{equation}{section}

% ---- theorem environments (numbered by section, one shared counter) ----
\theoremstyle{plain}
\newtheorem{theorem}{Theorem}[section]
\newtheorem{proposition}[theorem]{Proposition}
\newtheorem{lemma}[theorem]{Lemma}
\newtheorem{corollary}[theorem]{Corollary}
\newtheorem{maintheorem}[theorem]{Main Theorem}
% Theorems A, B, C, D of the paper: own counter, lettered, no section prefix
\newtheorem{thmlet}{Theorem}
\renewcommand{\thethmlet}{\Alph{thmlet}}
\theoremstyle{definition}
\newtheorem{definition}[theorem]{Definition}
\newtheorem{convention}[theorem]{Convention}
\newtheorem{example}[theorem]{Example}
\theoremstyle{remark}
\newtheorem{remark}[theorem]{Remark}

% A list of hypotheses labelled (H1),(H2),... : \begin{hyplist}{H} ... \end{hyplist}
% (also {O}, {L}, {B}, {K}, {W} for the families of section 5.3 of OUTLINE.md).
\newenvironment{hyplist}[1]%
  {\begin{enumerate}[label=\textup{(#1\arabic*)},leftmargin=3em,itemsep=2pt]}%
  {\end{enumerate}}

% ---- trust classes (OUTLINE 4.1, 5.1.3) -------------------------------
% Supplement: explicit tags.  Main paper: certified is the default and
% unmarked; every non-certified number is marked "(floating point)".
\ifsupplement
  \newcommand{\tcert}{\textup{[cert]}}   % arb/outward-rounded, pinned script and key
  \newcommand{\test}{\textup{[est]}}     % floating point or estimate; never in a theorem
  \newcommand{\thand}{\textup{[hand]}}   % proved by hand in the text
  \newcommand{\tlit}{\textup{[lit]}}     % cited from the literature
\else
  \newcommand{\tcert}{}
  \newcommand{\test}{\unskip\ \textup{(floating point)}}
  \newcommand{\thand}{}
  \newcommand{\tlit}{}
\fi

% ---- gaps: visible, counted, an error in a final build -----------------
% Usage: \GAP{one sentence saying what is missing}   (and an entry in ISSUES.md)
\newcounter{gapcount}
\newif\iffinalpaper   % \finalpapertrue makes every \GAP a hard error
\newcommand{\GAP}[1]{\stepcounter{gapcount}\typeout{GAP-MARKER \thegapcount}%
  \iffinalpaper\PackageError{paper}{\string\GAP\space used in a final build}{Remove all gaps.}\fi
  {\color{red}\sffamily\bfseries[GAP~\thegapcount:\ }{\color{red}\sffamily #1}{\color{red}\sffamily\bfseries]}}
\AtEndDocument{\typeout{GAP-MARKER total: \thegapcount}}

% ---- file names and JSON keys (supplement only) -------------------------
% No internal path appears in the main paper; a use there is a warning.
\newcommand{\fn}[1]{%
  \ifsupplement\else\PackageWarning{paper}{\string\fn\space used in the main paper}\fi
  \texttt{\detokenize{#1}}}
\newcommand{\key}[1]{\texttt{\detokenize{#1}}}

% ---- cross references between main paper and supplement ----------------
% Labels carry a unit prefix (section 5.3 of OUTLINE.md and DRAFTER_GUIDE.md):
%   main paper:  intro: obj: s3: s4: s5: s6: s7: s8: s9:   (kind: thm- lem- eq- tab- fig- ...)
%   supplement:  pA: pB: ... pK:
% so all labels are globally unique and the prefixes below are only a second guard.
% \Xref{label}  : reference to the OTHER document (hyperlinked via xr-hyper)
% \Lref{label}  : reference in the SAME document (plain \ref)
\ifsupplement
  \externaldocument[M-]{main}[main.pdf]
  \newcommand{\Xref}[1]{\ref{M-#1}}
\else
  \externaldocument[S-]{supplement}[supplement.pdf]
  \newcommand{\Xref}[1]{\ref{S-#1}}
\fi
\newcommand{\Lref}[1]{\ref{#1}}
% \Supp{pE}{lem-chart} -> "Lemma E.3" style text is produced by the drafter:
%   "Lemma~\Xref{pE:lem-chart} of the supplement".  \SuppPart{E} gives "Part E".
\newcommand{\SuppPart}[1]{Part~#1 of the supplement}

% ---- numbers (generated by unit U0; macros only, no literal certified numbers) ----
\IfFileExists{numbers/numbers_paper.tex}%
  {% numbers_paper.tex -- GENERATED from ledger.json by build_ledger.py; do not edit.
% Rounding: upper bounds up, lower bounds and margins down at the printed digit; margin = threshold / bound.
% Use as $\times\NName$ for margins.  kind: cert certified, est estimate, hand hypothesis constant, doc only a document quotes it.
% NUEhatZero [cert, upper] U_cert_v35.json :: supersolution.EA.0
\newcommand{\NUEhatZero}{\ensuremath{1.2018\cdot10^{-6}}}
% NUEhatOne [cert, upper] U_cert_v35.json :: supersolution.EA.1
\newcommand{\NUEhatOne}{\ensuremath{1.0441\cdot10^{-4}}}
% NUEhatTwo [cert, upper] U_cert_v35.json :: supersolution.EA.2
\newcommand{\NUEhatTwo}{\ensuremath{0.056694}}
% NUBhatZero [cert, upper] U_cert_v35.json :: supersolution.Bhat.0
\newcommand{\NUBhatZero}{\ensuremath{4.8735\cdot10^{-6}}}
% NUBhatOne [cert, upper] U_cert_v35.json :: supersolution.Bhat.1
\newcommand{\NUBhatOne}{\ensuremath{5.5933\cdot10^{-4}}}
% NUBhatTwo [cert, upper] U_cert_v35.json :: supersolution.Bhat.2
\newcommand{\NUBhatTwo}{\ensuremath{0.36675}}
% NUSupersolutionCZeroZero [cert, upper] U_cert_v35.json :: supersolution.C0.0
\newcommand{\NUSupersolutionCZeroZero}{\ensuremath{1.1340\cdot10^{-6}}}
% NUSupersolutionCZeroOne [cert, upper] U_cert_v35.json :: supersolution.C0.1
\newcommand{\NUSupersolutionCZeroOne}{\ensuremath{1.6623\cdot10^{-4}}}
% NUSupersolutionCOneZero [cert, upper] U_cert_v35.json :: supersolution.C1.0
\newcommand{\NUSupersolutionCOneZero}{\ensuremath{1.5241\cdot10^{-9}}}
% NUSupersolutionCOneOne [cert, upper] U_cert_v35.json :: supersolution.C1.1
\newcommand{\NUSupersolutionCOneOne}{\ensuremath{5.5255\cdot10^{-8}}}
% NUSupersolutionNTwoZero [cert, upper] U_cert_v35.json :: supersolution.N2.0
\newcommand{\NUSupersolutionNTwoZero}{\ensuremath{0.071472}}
% NUSupersolutionNTwoOne [cert, upper] U_cert_v35.json :: supersolution.N2.1
\newcommand{\NUSupersolutionNTwoOne}{\ensuremath{0.33012}}
% NUSupersolutionNTwoTwo [cert, upper] U_cert_v35.json :: supersolution.N2.2
\newcommand{\NUSupersolutionNTwoTwo}{\ensuremath{5.0406}}
% NUSupersolutionNTwoOverFourNTwoRefZero [cert, upper] U_cert_v35.json :: supersolution.N2_over_4N2ref.0
\newcommand{\NUSupersolutionNTwoOverFourNTwoRefZero}{\ensuremath{0.18}}
% NUSupersolutionNTwoOverFourNTwoRefOne [cert, upper] U_cert_v35.json :: supersolution.N2_over_4N2ref.1
\newcommand{\NUSupersolutionNTwoOverFourNTwoRefOne}{\ensuremath{0.42}}
% NUSupersolutionNTwoOverFourNTwoRefTwo [cert, upper] U_cert_v35.json :: supersolution.N2_over_4N2ref.2
\newcommand{\NUSupersolutionNTwoOverFourNTwoRefTwo}{\ensuremath{0.26}}
% NUSupersolutionMOneZHypCheckZonesIIISOneTwoBhat [cert, upper] U_cert_v35.json :: supersolution.M1z_hyp_check.zones_I_II.S1_2Bhat
\newcommand{\NUSupersolutionMOneZHypCheckZonesIIISOneTwoBhat}{\ensuremath{0.28585}}
% NUSupersolutionMOneZHypCheckZoneIIIBOnePlusPUOne [cert, upper] U_cert_v35.json :: supersolution.M1z_hyp_check.zone_III.b1_plus_pU1
\newcommand{\NUSupersolutionMOneZHypCheckZoneIIIBOnePlusPUOne}{\ensuremath{0.50786}}
% NUMarginLateMOne [cert, margin] U_cert_v35.json :: margins.late_M1
\newcommand{\NUMarginLateMOne}{\ensuremath{2.353}}
% NUMarginLateMZero [cert, margin] U_cert_v35.json :: margins.late_M0
\newcommand{\NUMarginLateMZero}{\ensuremath{4.892}}
% NUMarginsLateMZeroZ [cert, margin] U_cert_v35.json :: margins.late_M0z
\newcommand{\NUMarginsLateMZeroZ}{\ensuremath{12.99}}
% NUMarginPresentMOne [cert, margin] U_cert_v35.json :: margins.present_M1_min.value
\newcommand{\NUMarginPresentMOne}{\ensuremath{2.770}}
% NUMarginPresentMZero [cert, margin] U_cert_v35.json :: margins.present_M0_min.value
\newcommand{\NUMarginPresentMZero}{\ensuremath{2.225}}
% NUMarginsPresentMZeroZMinValue [cert, margin] U_cert_v35.json :: margins.present_M0z_min.value
\newcommand{\NUMarginsPresentMZeroZMinValue}{\ensuremath{14.37}}
% NULoopGainZeroTotal [cert, upper] U_cert_v35.json :: loop_gain.0.total
\newcommand{\NULoopGainZeroTotal}{\ensuremath{1.23\cdot10^{-7}}}
% NULoopGainOneTotal [cert, upper] U_cert_v35.json :: loop_gain.1.total
\newcommand{\NULoopGainOneTotal}{\ensuremath{0.037169}}
% NULoopGainTwoNeighbour [cert, upper] U_cert_v35.json :: loop_gain.2.neighbour
\newcommand{\NULoopGainTwoNeighbour}{\ensuremath{0.037018}}
% NULoopGainTwoThroughA [cert, upper] U_cert_v35.json :: loop_gain.2.through_A
\newcommand{\NULoopGainTwoThroughA}{\ensuremath{0.0020172}}
% NUfineBhatZero [cert, upper] U_cert_v35.json :: ufine.Bhat.0
\newcommand{\NUfineBhatZero}{\ensuremath{4.8723\cdot10^{-6}}}
% NUfineBhatOne [cert, upper] U_cert_v35.json :: ufine.Bhat.1
\newcommand{\NUfineBhatOne}{\ensuremath{2.6266\cdot10^{-4}}}
% NUfineBhatTwo [cert, upper] U_cert_v35.json :: ufine.Bhat.2
\newcommand{\NUfineBhatTwo}{\ensuremath{0.025842}}
% NUcdeltaU [cert, upper] U_cert_v35.json :: launch.c_dU
\newcommand{\NUcdeltaU}{\ensuremath{1.7187\cdot10^{-8}}}
% NULaunchWmin [cert, lower] U_cert_v35.json :: launch.wmin
\newcommand{\NULaunchWmin}{\ensuremath{0.73492}}
% NULaunchCAD [cert, upper] U_cert_v35.json :: launch.CA_D
\newcommand{\NULaunchCAD}{\ensuremath{2.2429}}
% NULaunchQD [cert, upper] U_cert_v35.json :: launch.qD
\newcommand{\NULaunchQD}{\ensuremath{0.018509}}
% NULaunchDeltaInZero [cert, upper] U_cert_v35.json :: launch.Delta_in.0
\newcommand{\NULaunchDeltaInZero}{\ensuremath{3.131\cdot10^{-12}}}
% NULaunchDeltaInOne [cert, upper] U_cert_v35.json :: launch.Delta_in.1
\newcommand{\NULaunchDeltaInOne}{\ensuremath{3.287\cdot10^{-12}}}
% NULaunchDeltaInTwo [cert, upper] U_cert_v35.json :: launch.Delta_in.2
\newcommand{\NULaunchDeltaInTwo}{\ensuremath{3.451\cdot10^{-12}}}
% NULaunchPhiZero [cert, upper] U_cert_v35.json :: launch.Phi.0
\newcommand{\NULaunchPhiZero}{\ensuremath{6.434\cdot10^{-9}}}
% NULaunchPhiOne [cert, upper] U_cert_v35.json :: launch.Phi.1
\newcommand{\NULaunchPhiOne}{\ensuremath{1.4217\cdot10^{-8}}}
% NULaunchPhiTwo [cert, upper] U_cert_v35.json :: launch.Phi.2
\newcommand{\NULaunchPhiTwo}{\ensuremath{3.8404\cdot10^{-8}}}
% NUGammaZero [cert, upper] U_cert_v35.json :: Gamma.0
\newcommand{\NUGammaZero}{\ensuremath{4.0000358}}
% NUGammaOne [cert, upper] U_cert_v35.json :: Gamma.1
\newcommand{\NUGammaOne}{\ensuremath{16.000322}}
% NUGammaTwo [cert, upper] U_cert_v35.json :: Gamma.2
\newcommand{\NUGammaTwo}{\ensuremath{64.00218}}
% NUKappa [cert, upper] U_cert_v35.json :: kappa
\newcommand{\NUKappa}{\ensuremath{8.0000358}}
% NUKappaTwo [cert, upper] U_cert_v35.json :: kappa2
\newcommand{\NUKappaTwo}{\ensuremath{64.00061}}
% NUGhatZeroZero [cert, upper] U_cert_v35.json :: Ghat.0.0
\newcommand{\NUGhatZeroZero}{\ensuremath{51322.6}}
% NUGhatOneOne [cert, upper] U_cert_v35.json :: Ghat.1.1
\newcommand{\NUGhatOneOne}{\ensuremath{33962.9}}
% NUGhatTwoTwo [cert, upper] U_cert_v35.json :: Ghat.2.2
\newcommand{\NUGhatTwoTwo}{\ensuremath{348341.1}}
% NUGhatOneZero [cert, upper] U_cert_v35.json :: Ghat.1.0
\newcommand{\NUGhatOneZero}{\ensuremath{4.938}}
% NUGhatTwoZero [cert, upper] U_cert_v35.json :: Ghat.2.0
\newcommand{\NUGhatTwoZero}{\ensuremath{27.984}}
% NUGhatTwoOne [cert, upper] U_cert_v35.json :: Ghat.2.1
\newcommand{\NUGhatTwoOne}{\ensuremath{458.53}}
% NUKZeroZero [cert, upper] U_cert_v35.json :: K.0.0
\newcommand{\NUKZeroZero}{\ensuremath{8034.21}}
% NUKZeroOne [cert, upper] U_cert_v35.json :: K.0.1
\newcommand{\NUKZeroOne}{\ensuremath{724.89}}
% NUKOneZero [cert, upper] U_cert_v35.json :: K.1.0
\newcommand{\NUKOneZero}{\ensuremath{0.08152}}
% NUKOneOne [cert, upper] U_cert_v35.json :: K.1.1
\newcommand{\NUKOneOne}{\ensuremath{6336.94}}
% NUKOneTwo [cert, upper] U_cert_v35.json :: K.1.2
\newcommand{\NUKOneTwo}{\ensuremath{545.59}}
% NUKOneThree [cert, upper] U_cert_v35.json :: K.1.3
\newcommand{\NUKOneThree}{\ensuremath{0.58700}}
% NUKTwoZero [cert, upper] U_cert_v35.json :: K.2.0
\newcommand{\NUKTwoZero}{\ensuremath{0.87186}}
% NUKTwoOne [cert, upper] U_cert_v35.json :: K.2.1
\newcommand{\NUKTwoOne}{\ensuremath{528.62}}
% NUKTwoTwo [cert, upper] U_cert_v35.json :: K.2.2
\newcommand{\NUKTwoTwo}{\ensuremath{23425.98}}
% NUKTwoThree [cert, upper] U_cert_v35.json :: K.2.3
\newcommand{\NUKTwoThree}{\ensuremath{555.62}}
% NUFootwOne [cert, upper] U_cert_v35.json :: footw.1
\newcommand{\NUFootwOne}{\ensuremath{2.6361}}
% NUFootwTwo [cert, upper] U_cert_v35.json :: footw.2
\newcommand{\NUFootwTwo}{\ensuremath{2663.3}}
% NUCThreeThree [cert, upper] U_cert_v35.json :: c33
\newcommand{\NUCThreeThree}{\ensuremath{108.02}}
% NUCTwoTwoThree [cert, upper] U_cert_v35.json :: c223
\newcommand{\NUCTwoTwoThree}{\ensuremath{7.3284}}
% NUCEta [cert, upper] U_cert_v35.json :: C_eta
\newcommand{\NUCEta}{\ensuremath{13.376}}
% NULMax [cert, upper] U_cert_v35.json :: L_max
\newcommand{\NULMax}{\ensuremath{3.5352}}
% NURRatio [cert, upper] U_cert_v35.json :: R_ratio
\newcommand{\NURRatio}{\ensuremath{11.593}}
% NUSweepsSigmaInCoverSup [cert, upper] U_cert_v35.json :: sweeps.Sigma_in.cover_sup
\newcommand{\NUSweepsSigmaInCoverSup}{\ensuremath{5.0892}}
% NUSweepsSigmaECoverSup [cert, upper] U_cert_v35.json :: sweeps.Sigma_E.cover_sup
\newcommand{\NUSweepsSigmaECoverSup}{\ensuremath{0.30846}}
% NUSweepsSXCoverSup [cert, upper] U_cert_v35.json :: sweeps.S_X.cover_sup
\newcommand{\NUSweepsSXCoverSup}{\ensuremath{1.5243}}
% NUSweepsSigmaXCoverSup [cert, upper] U_cert_v35.json :: sweeps.Sigma_X.cover_sup
\newcommand{\NUSweepsSigmaXCoverSup}{\ensuremath{3.0081}}
% NUSweepsSigmaRsTwoCoverSup [cert, upper] U_cert_v35.json :: sweeps.Sigma_rs.2.cover_sup
\newcommand{\NUSweepsSigmaRsTwoCoverSup}{\ensuremath{4.3350\cdot10^{-3}}}
% NUDataTermsLateDZero [cert, upper] U_cert_v35.json :: data_terms_late.D.0
\newcommand{\NUDataTermsLateDZero}{\ensuremath{1.7330\cdot10^{-6}}}
% NUDataTermsLateDOne [cert, upper] U_cert_v35.json :: data_terms_late.D.1
\newcommand{\NUDataTermsLateDOne}{\ensuremath{2.5353\cdot10^{-4}}}
% NUDataTermsLateSlopeShareCOneSigmaXOverDkZero [cert, upper] U_cert_v35.json :: data_terms_late.slope_share_C1SigmaX_over_Dk.0
\newcommand{\NUDataTermsLateSlopeShareCOneSigmaXOverDkZero}{\ensuremath{2.7\cdot10^{-3}}}
% NUDataTermsLateSlopeShareCOneSigmaXOverDkOne [cert, upper] U_cert_v35.json :: data_terms_late.slope_share_C1SigmaX_over_Dk.1
\newcommand{\NUDataTermsLateSlopeShareCOneSigmaXOverDkOne}{\ensuremath{6.6\cdot10^{-4}}}
% NURLtTwoMarginZero [cert, margin] U_cert_v35.json :: R_lt2_margin.0
\newcommand{\NURLtTwoMarginZero}{\ensuremath{4.432}}
% NURLtTwoMarginOne [cert, margin] U_cert_v35.json :: R_lt2_margin.1
\newcommand{\NURLtTwoMarginOne}{\ensuremath{6.697}}
% NUThetaTwoEight [cert, upper] U_cert_v35.json :: theta28
\newcommand{\NUThetaTwoEight}{\ensuremath{0.37011}}
% NUDFiveSix [cert, upper] U_cert_v35.json :: D56
\newcommand{\NUDFiveSix}{\ensuremath{1.0368}}
% NUDTwoEight [cert, upper] U_cert_v35.json :: D28
\newcommand{\NUDTwoEight}{\ensuremath{1.5876}}
% NUSThreeThreeFootOverBhatTwo [cert, upper] U_cert_v35.json :: S33_foot_over_Bhat2
\newcommand{\NUSThreeThreeFootOverBhatTwo}{\ensuremath{1.69\cdot10^{-3}}}
% NSAHJWKappaJ [cert, upper] A_interface.json :: HJW.kappa_J
\newcommand{\NSAHJWKappaJ}{\ensuremath{0.042351}}
% NSAGpZero [cert, lower] A_interface.json :: res.tight.gp.0
\newcommand{\NSAGpZero}{\ensuremath{0.24292}}
% NSAGpOne [cert, lower] A_interface.json :: res.tight.gp.1
\newcommand{\NSAGpOne}{\ensuremath{0.24106}}
% NSAGpTwo [cert, lower] A_interface.json :: res.tight.gp.2
\newcommand{\NSAGpTwo}{\ensuremath{0.11961}}
% NSAAminC [cert, lower] A_interface.json :: res.tight.amin_c
\newcommand{\NSAAminC}{\ensuremath{0.97623}}
% NSARhoJ [cert, upper] A_interface.json :: res.tight.rhoJ
\newcommand{\NSARhoJ}{\ensuremath{5.4995\cdot10^{-7}}}
% NSAKappaOTwoTwo [cert, upper] A_interface.json :: res.tight.kappa_o.2.2
\newcommand{\NSAKappaOTwoTwo}{\ensuremath{1.2416\cdot10^{-5}}}
% NSADsrcZero [cert, upper] A_interface.json :: res.tight.Dsrc.0
\newcommand{\NSADsrcZero}{\ensuremath{1.8092\cdot10^{-7}}}
% NSADsrcOne [cert, upper] A_interface.json :: res.tight.Dsrc.1
\newcommand{\NSADsrcOne}{\ensuremath{9.8276\cdot10^{-5}}}
% NSABZeroDWindowSlopeZero [cert, upper] A_interface.json :: res.tight.B0.D_window_slope.0
\newcommand{\NSABZeroDWindowSlopeZero}{\ensuremath{8.7399\cdot10^{-10}}}
% NSABZeroDWindowSlopeOne [cert, upper] A_interface.json :: res.tight.B0.D_window_slope.1
\newcommand{\NSABZeroDWindowSlopeOne}{\ensuremath{5.2599\cdot10^{-8}}}
% NSAFTwoZero [cert, upper] A_interface.json :: res.tight.F2.0
\newcommand{\NSAFTwoZero}{\ensuremath{5.6506\cdot10^{-13}}}
% NSAFTwoOne [cert, upper] A_interface.json :: res.tight.F2.1
\newcommand{\NSAFTwoOne}{\ensuremath{5.9320\cdot10^{-13}}}
% NSACoupOne [cert, upper] A_interface.json :: res.tight.coup.1
\newcommand{\NSACoupOne}{\ensuremath{1.1568\cdot10^{-12}}}
% NSACOneZeroZero [cert, upper] A_interface.json :: res.tight.C1_0.0
\newcommand{\NSACOneZeroZero}{\ensuremath{7.2060\cdot10^{-10}}}
% NSACOneZeroOne [cert, upper] A_interface.json :: res.tight.C1_0.1
\newcommand{\NSACOneZeroOne}{\ensuremath{4.3143\cdot10^{-8}}}
% NSABZeroSupCoreZero [cert, upper] A_interface.json :: res.tight.B0.sup_core.0
\newcommand{\NSABZeroSupCoreZero}{\ensuremath{4.23\cdot10^{-5}}}
% NSACoreFluxMarginTwoFLUXZeroOverKZero [cert, margin] A_interface.json :: res.tight.core.flux_margin.2.FLUX0_over_K0
\newcommand{\NSACoreFluxMarginTwoFLUXZeroOverKZero}{\ensuremath{3.396}}
% NSACoreFluxMarginTwoFLUXTwoOverKTwo [cert, margin] A_interface.json :: res.tight.core.flux_margin.2.FLUX2_over_K2
\newcommand{\NSACoreFluxMarginTwoFLUXTwoOverKTwo}{\ensuremath{6.012}}
% NSAMarginFluxZero [cert, margin] asm_design_N2ref4.json :: res.tight.flux.beta.2.margin_FLUX0
\newcommand{\NSAMarginFluxZero}{\ensuremath{3.365}}
% NSAMarginFluxTwo [cert, margin] asm_design_N2ref4.json :: res.tight.flux.beta.2.margin_FLUX2
\newcommand{\NSAMarginFluxTwo}{\ensuremath{5.955}}
% NSAfourSZero [cert, upper] asm_design_N2ref4.json :: res.tight.S0
\newcommand{\NSAfourSZero}{\ensuremath{7.1878\cdot10^{-3}}}
% NSAfourSOne [cert, upper] asm_design_N2ref4.json :: res.tight.S1
\newcommand{\NSAfourSOne}{\ensuremath{0.10430}}
% NSAMarginCoreSZeroV [cert, margin] asm_design_N2ref4.json :: res.tight.margins.S0v_vs_0.0153
\newcommand{\NSAMarginCoreSZeroV}{\ensuremath{2.116}}
% NSAMarginCoreSOneV [cert, margin] asm_design_N2ref4.json :: res.tight.margins.S1v_vs_0.6
\newcommand{\NSAMarginCoreSOneV}{\ensuremath{5.748}}
% NSAMarginCoreMOne [cert, margin] asm_design_N2ref4.json :: res.tight.margins.M1core
\newcommand{\NSAMarginCoreMOne}{\ensuremath{4.409}}
% NSAMarginCoreMZero [cert, margin] asm_design_N2ref4.json :: res.tight.margins.M0core
\newcommand{\NSAMarginCoreMZero}{\ensuremath{5.450}}
% NSAMarginGhat [cert, margin] asm_design_N2ref4.json :: res.tight.margins.Ghat22
\newcommand{\NSAMarginGhat}{\ensuremath{1.917}}
% NSAMarginRampMOne [cert, margin] asm_design_N2ref4.json :: res.tight.margins.ramp M1
\newcommand{\NSAMarginRampMOne}{\ensuremath{1.184}}
% NSAfourInfoZeroDWindowConst [cert, upper] asm_design_N2ref4.json :: res.tight.info.0.D_window_const
\newcommand{\NSAfourInfoZeroDWindowConst}{\ensuremath{1.7732\cdot10^{-7}}}
% NSAfourInfoOneDWindowConst [cert, upper] asm_design_N2ref4.json :: res.tight.info.1.D_window_const
\newcommand{\NSAfourInfoOneDWindowConst}{\ensuremath{9.8058\cdot10^{-5}}}
% NZDeltaRZero [cert, upper] zone3_interface.json :: DeltaR.0
\newcommand{\NZDeltaRZero}{\ensuremath{4.4986\cdot10^{-8}}}
% NZDeltaROne [cert, upper] zone3_interface.json :: DeltaR.1
\newcommand{\NZDeltaROne}{\ensuremath{6.7972\cdot10^{-6}}}
% NZDeltaRTwo [cert, upper] zone3_interface.json :: DeltaR.2
\newcommand{\NZDeltaRTwo}{\ensuremath{0.011480}}
% NZDeltaRBZero [cert, upper] zone3_interface.json :: DeltaR_b.0
\newcommand{\NZDeltaRBZero}{\ensuremath{2.66\cdot10^{-12}}}
% NZDeltaRBOne [cert, upper] zone3_interface.json :: DeltaR_b.1
\newcommand{\NZDeltaRBOne}{\ensuremath{3.26\cdot10^{-9}}}
% NZDeltaRBTwo [cert, upper] zone3_interface.json :: DeltaR_b.2
\newcommand{\NZDeltaRBTwo}{\ensuremath{4.35\cdot10^{-8}}}
% NZDeltaRDetailNSqZero [cert, upper] zone3_interface.json :: DeltaR_detail.N_sq.0
\newcommand{\NZDeltaRDetailNSqZero}{\ensuremath{8.28\cdot10^{-5}}}
% NZDeltaRDetailNSqOne [cert, upper] zone3_interface.json :: DeltaR_detail.N_sq.1
\newcommand{\NZDeltaRDetailNSqOne}{\ensuremath{0.10149}}
% NZDeltaRDetailNSqTwo [cert, upper] zone3_interface.json :: DeltaR_detail.N_sq.2
\newcommand{\NZDeltaRDetailNSqTwo}{\ensuremath{1.3567}}
% NZDeltaRDetailNSqThree [cert, upper] zone3_interface.json :: DeltaR_detail.N_sq.3
\newcommand{\NZDeltaRDetailNSqThree}{\ensuremath{25350}}
% NZDeltaRDetailPartsTwoPairsThree [cert, upper] zone3_interface.json :: DeltaR_detail.parts.2.pairs.3
\newcommand{\NZDeltaRDetailPartsTwoPairsThree}{\ensuremath{0.011298}}
% NZDeltaRDetailPartsTwoLFourFoot [cert, upper] zone3_interface.json :: DeltaR_detail.parts.2.L4_foot
\newcommand{\NZDeltaRDetailPartsTwoLFourFoot}{\ensuremath{1.6\cdot10^{-4}}}
% NZDeltaRDetailPartsTwoCubic [cert, upper] zone3_interface.json :: DeltaR_detail.parts.2.cubic
\newcommand{\NZDeltaRDetailPartsTwoCubic}{\ensuremath{9.3\cdot10^{-6}}}
% NZDeltaRDetailPartsTwoSThreeThreeFoot [cert, upper] zone3_interface.json :: DeltaR_detail.parts.2.S33_foot
\newcommand{\NZDeltaRDetailPartsTwoSThreeThreeFoot}{\ensuremath{3.7\cdot10^{-7}}}
% NZNIIIZero [cert, upper] zone3_interface.json :: n_III.0
\newcommand{\NZNIIIZero}{\ensuremath{0.071445}}
% NZNIIIOne [cert, upper] zone3_interface.json :: n_III.1
\newcommand{\NZNIIIOne}{\ensuremath{0.32707}}
% NZNIIITwo [cert, upper] zone3_interface.json :: n_III.2
\newcommand{\NZNIIITwo}{\ensuremath{3.0445}}
% NZTauIII [cert, upper] zone3_interface.json :: tau_III
\newcommand{\NZTauIII}{\ensuremath{4.0708}}
% NZoneIIIBoxZero [cert, upper] zone3_interface.json :: box_III_380.b0
\newcommand{\NZoneIIIBoxZero}{\ensuremath{0.10032}}
% NZoneIIIBoxOne [cert, upper] zone3_interface.json :: box_III_380.b1
\newcommand{\NZoneIIIBoxOne}{\ensuremath{0.41974}}
% NZoneIIIBoxZ [cert, upper] zone3_interface.json :: box_III_380.bz
\newcommand{\NZoneIIIBoxZ}{\ensuremath{0.014291}}
% NZBoxIIIThreeEightZeroMarginsBZero [cert, margin] zone3_interface.json :: box_III_380.margins.b0
\newcommand{\NZBoxIIIThreeEightZeroMarginsBZero}{\ensuremath{1.495}}
% NZoneIIIMarginThinnest [cert, margin] zone3_interface.json :: box_III_380.margins.b1
\newcommand{\NZoneIIIMarginThinnest}{\ensuremath{1.238}}
% NZBoxIIIThreeEightZeroMarginsBz [cert, margin] zone3_interface.json :: box_III_380.margins.bz
\newcommand{\NZBoxIIIThreeEightZeroMarginsBz}{\ensuremath{11.19}}
% NZBoxIIISevenZeroZeroBZero [cert, upper] zone3_interface.json :: box_III_700.b0
\newcommand{\NZBoxIIISevenZeroZeroBZero}{\ensuremath{0.22757}}
% NZBoxIIISevenZeroZeroBOne [cert, upper] zone3_interface.json :: box_III_700.b1
\newcommand{\NZBoxIIISevenZeroZeroBOne}{\ensuremath{0.54234}}
% NZUfineKTwoAtBhatRef [cert, upper] zone3_interface.json :: ufine.K2_at_Bhat_ref
\newcommand{\NZUfineKTwoAtBhatRef}{\ensuremath{0.046956}}
% NZUfineKThreeStripAtBhatRef [cert, upper] zone3_interface.json :: ufine.K3_strip_at_Bhat_ref
\newcommand{\NZUfineKThreeStripAtBhatRef}{\ensuremath{0.054140}}
% NRecoveryMinMargin [cert, margin] zone3_interface.json :: cone_family.H6.at_Bhat_ref.min_margin
\newcommand{\NRecoveryMinMargin}{\ensuremath{1.860}}
% NZConeFamilyHSixAtBhatRefConeSlopeMargin [cert, margin] zone3_interface.json :: cone_family.H6.at_Bhat_ref.cone.slope.margin
\newcommand{\NZConeFamilyHSixAtBhatRefConeSlopeMargin}{\ensuremath{2.46}}
% NZConeFamilyHSixAtBhatRefConeKTwoMargin [cert, margin] zone3_interface.json :: cone_family.H6.at_Bhat_ref.cone.K2.margin
\newcommand{\NZConeFamilyHSixAtBhatRefConeKTwoMargin}{\ensuremath{2.54}}
% NZConeFamilyHSixAtBhatRefConeKThreeMargin [cert, margin] zone3_interface.json :: cone_family.H6.at_Bhat_ref.cone.K3.margin
\newcommand{\NZConeFamilyHSixAtBhatRefConeKThreeMargin}{\ensuremath{2.51}}
% NConeTstarMargin [cert, margin] zone3_interface.json :: acceptance.tstar_margin
\newcommand{\NConeTstarMargin}{\ensuremath{2.30}}
% NUoneTFiveYaddLevelOneTwo [cert, upper] u1_ghat.json :: t5_Yadd.level1.2
\newcommand{\NUoneTFiveYaddLevelOneTwo}{\ensuremath{0.0044}}
% NUoneTFiveYaddLevelTwoTwo [cert, upper] u1_ghat.json :: t5_Yadd.level2.2
\newcommand{\NUoneTFiveYaddLevelTwoTwo}{\ensuremath{0.1081}}
% NUoneTFiveYaddLevelTwoThree [cert, upper] u1_ghat.json :: t5_Yadd.level2.3
\newcommand{\NUoneTFiveYaddLevelTwoThree}{\ensuremath{0.374}}
% NUoneWeightThetaGammaZero [cert, upper] u1_ghat.json :: weight.Theta.Gamma.0
\newcommand{\NUoneWeightThetaGammaZero}{\ensuremath{4.00000000000000000000054}}
% NRKZeroZero [cert, upper] ramp_U.json :: K.0.0
\newcommand{\NRKZeroZero}{\ensuremath{13436.03}}
% NRKOneOne [cert, upper] ramp_U.json :: K.1.1
\newcommand{\NRKOneOne}{\ensuremath{10526.19}}
% NRKTwoTwo [cert, upper] ramp_U.json :: K.2.2
\newcommand{\NRKTwoTwo}{\ensuremath{39482.30}}
% NUoneCasesPrimaryYNextZero [cert, upper] u1_ghat.json :: cases.primary.Y_next.0
\newcommand{\NUoneCasesPrimaryYNextZero}{\ensuremath{457.092}}
% NUoneCasesPrimaryYNextOne [cert, upper] u1_ghat.json :: cases.primary.Y_next.1
\newcommand{\NUoneCasesPrimaryYNextOne}{\ensuremath{356.940}}
% NUoneCasesPrimaryYNextTwo [cert, upper] u1_ghat.json :: cases.primary.Y_next.2
\newcommand{\NUoneCasesPrimaryYNextTwo}{\ensuremath{1320.220}}
% NSfourResultsMNineHTwoTightenedFUOneRowsZeroBound [cert, upper] FINAL_manifest_v33.json :: results.M9_H2_tightened_fU1.rows.0.bound
\newcommand{\NSfourResultsMNineHTwoTightenedFUOneRowsZeroBound}{\ensuremath{3.658\cdot10^{-10}}}
% NSfourResultsMNineHTwoTightenedFUOneRowsZeroMargin [cert, margin] FINAL_manifest_v33.json :: results.M9_H2_tightened_fU1.rows.0.margin
\newcommand{\NSfourResultsMNineHTwoTightenedFUOneRowsZeroMargin}{\ensuremath{27.34}}
% NSfourResultsMNineHTwoTightenedFUOneRowsOneBound [cert, upper] FINAL_manifest_v33.json :: results.M9_H2_tightened_fU1.rows.1.bound
\newcommand{\NSfourResultsMNineHTwoTightenedFUOneRowsOneBound}{\ensuremath{1.777\cdot10^{-9}}}
% NSfourResultsMNineHTwoTightenedFUOneRowsOneMargin [cert, margin] FINAL_manifest_v33.json :: results.M9_H2_tightened_fU1.rows.1.margin
\newcommand{\NSfourResultsMNineHTwoTightenedFUOneRowsOneMargin}{\ensuremath{562.8}}
% NSfourResultsMNineHTwoTightenedFUOneRowsTwoBound [cert, upper] FINAL_manifest_v33.json :: results.M9_H2_tightened_fU1.rows.2.bound
\newcommand{\NSfourResultsMNineHTwoTightenedFUOneRowsTwoBound}{\ensuremath{1.665\cdot10^{-6}}}
% NSfourResultsMNineHTwoTightenedFUOneRowsTwoMargin [cert, margin] FINAL_manifest_v33.json :: results.M9_H2_tightened_fU1.rows.2.margin
\newcommand{\NSfourResultsMNineHTwoTightenedFUOneRowsTwoMargin}{\ensuremath{600.9}}
% NHOneAlphaThr [hand, exact] hand :: 
\newcommand{\NHOneAlphaThr}{\ensuremath{1.5\cdot10^{-4}}}
% NHOneSqrtLamThr [hand, exact] hand :: 
\newcommand{\NHOneSqrtLamThr}{\ensuremath{5\cdot10^{-8}}}
% NYbarStep [hand, exact] hand :: 
\newcommand{\NYbarStep}{\ensuremath{1\cdot10^{-12}}}
% NHTwoThrZero [hand, exact] hand :: 
\newcommand{\NHTwoThrZero}{\ensuremath{1\cdot10^{-8}}}
% NHTwoThrOne [hand, exact] hand :: 
\newcommand{\NHTwoThrOne}{\ensuremath{1\cdot10^{-6}}}
% NHTwoThrTwo [hand, exact] hand :: 
\newcommand{\NHTwoThrTwo}{\ensuremath{1\cdot10^{-3}}}
% NHThreeThrZero [hand, exact] hand :: 
\newcommand{\NHThreeThrZero}{\ensuremath{1\cdot10^{-6}}}
% NHThreeThrOne [hand, exact] hand :: 
\newcommand{\NHThreeThrOne}{\ensuremath{5\cdot10^{-6}}}
% NHThreeThrTwo [hand, exact] hand :: 
\newcommand{\NHThreeThrTwo}{\ensuremath{1\cdot10^{-5}}}
% NHFourThrZero [hand, exact] hand :: 
\newcommand{\NHFourThrZero}{\ensuremath{2\cdot10^{-7}}}
% NHFourThrOne [hand, exact] hand :: 
\newcommand{\NHFourThrOne}{\ensuremath{1\cdot10^{-5}}}
% NHFourThrTwo [hand, exact] hand :: 
\newcommand{\NHFourThrTwo}{\ensuremath{1\cdot10^{-5}}}
% NHFourThrThree [hand, exact] hand :: 
\newcommand{\NHFourThrThree}{\ensuremath{1\cdot10^{-5}}}
% NHFourKfourThr [hand, exact] hand :: 
\newcommand{\NHFourKfourThr}{\ensuremath{0.06}}
% NHSixConeSlopeThr [hand, exact] hand :: 
\newcommand{\NHSixConeSlopeThr}{\ensuremath{0.08}}
% NHSixConeKtwoThr [hand, exact] hand :: 
\newcommand{\NHSixConeKtwoThr}{\ensuremath{0.03}}
% NHSixConeNthreeThr [hand, exact] hand :: 
\newcommand{\NHSixConeNthreeThr}{\ensuremath{0.004}}
% NHSixConeJetThr [hand, exact] hand :: 
\newcommand{\NHSixConeJetThr}{\ensuremath{0.12}}
% NHSixConeHFiveThr [hand, exact] hand :: 
\newcommand{\NHSixConeHFiveThr}{\ensuremath{0.0144}}
% NHSixFarHFiveThr [hand, exact] hand :: 
\newcommand{\NHSixFarHFiveThr}{\ensuremath{0.0225}}
% NLemmaOConeM [hand, exact] hand :: 
\newcommand{\NLemmaOConeM}{\ensuremath{0.12}}
% NLemmaOConeK [hand, exact] hand :: 
\newcommand{\NLemmaOConeK}{\ensuremath{0.18}}
% NLemmaOFarM [hand, exact] hand :: 
\newcommand{\NLemmaOFarM}{\ensuremath{0.15}}
% NLemmaOFarK [hand, exact] hand :: 
\newcommand{\NLemmaOFarK}{\ensuremath{0.3}}
% NCoercivityConst [hand, exact] hand :: 
\newcommand{\NCoercivityConst}{\ensuremath{0.046875}}
% NMOnez [hand, exact] hand :: 
\newcommand{\NMOnez}{\ensuremath{0.534}}
% NBetaA [hand, exact] hand :: 
\newcommand{\NBetaA}{\ensuremath{2}}
% NTOneLower [cert, lower] t1clock.json :: t1.0
\newcommand{\NTOneLower}{\ensuremath{4.5978574775}}
% NTOneUpper [cert, upper] t1clock.json :: t1.1
\newcommand{\NTOneUpper}{\ensuremath{4.5978586144}}
% NLambdaOneLower [cert, lower] t1clock.json :: lnlam.0
\newcommand{\NLambdaOneLower}{\ensuremath{6.49\cdot10^{-5}}}
% NLambdaOneUpper [cert, upper] t1clock.json :: lnlam.1
\newcommand{\NLambdaOneUpper}{\ensuremath{6.75\cdot10^{-5}}}
% NTMinusTOneUpper [cert, upper] t1clock.json :: lnlam.1; 1.5e-4.sup_ab
\newcommand{\NTMinusTOneUpper}{\ensuremath{2.1\cdot10^{-17}}}
% NAbarUpper [cert, upper] final_alpha_points_5e-8.json :: 1.5e-4.sup_ab
\newcommand{\NAbarUpper}{\ensuremath{2.19103\cdot10^{-7}}}
% NBetabarUpper [cert, upper] final_alpha_points_5e-8.json :: 1.5e-4.sup_beta
\newcommand{\NBetabarUpper}{\ensuremath{3.33811\cdot10^{-9}}}
% NBaseKmaxLower [cert, lower] t1clock.json :: base_kmax.0
\newcommand{\NBaseKmaxLower}{\ensuremath{2.279\cdot10^{4}}}
% NBaseKmaxUpper [cert, upper] t1clock.json :: base_kmax.1
\newcommand{\NBaseKmaxUpper}{\ensuremath{2.376\cdot10^{4}}}
% NHOneAlphaBase [cert, upper] summary.json :: alpha_abs_max
\newcommand{\NHOneAlphaBase}{\ensuremath{1.1739\cdot10^{-4}}}
% NHOneAlphaErr [cert, upper] t9_h1_clam_s180.json :: C_alpha; inputs.YSTAR
\newcommand{\NHOneAlphaErr}{\ensuremath{7.21\cdot10^{-13}}}
% NHOneAlphaMargin [cert, margin] summary.json :: alpha_abs_max; C_alpha; inputs.YSTAR
\newcommand{\NHOneAlphaMargin}{\ensuremath{1.27}}
% NHOneSqrtLamBase [cert, upper] t9_h1_clam_s180.json :: sqrtLam_base_Wp
\newcommand{\NHOneSqrtLamBase}{\ensuremath{2.282\cdot10^{-8}}}
% NHOneSqrtLamTotal [cert, upper] t9_h1_clam_s180.json :: sqrtLam_total_Wp
\newcommand{\NHOneSqrtLamTotal}{\ensuremath{2.29\cdot10^{-8}}}
% NHOneSqrtLamMargin [cert, margin] t9_h1_clam_s180.json :: sqrtLam_total_Wp
\newcommand{\NHOneSqrtLamMargin}{\ensuremath{2.18}}
% NHOneYbarLimit [cert, lower] t9_h1_clam_s180.json :: Ybar_limit_H1
\newcommand{\NHOneYbarLimit}{\ensuremath{1.12\cdot10^{-9}}}
% NSwitchYUpper [cert, upper] switch_s2.json :: modulated.Y_raw
\newcommand{\NSwitchYUpper}{\ensuremath{6.21\cdot10^{-13}}}
% NSwitchYMargin [cert, margin] switch_s2.json :: modulated.Y_raw
\newcommand{\NSwitchYMargin}{\ensuremath{1.6}}
% NMinCI [cert, lower] constants.json :: meta.summary.min_c_I
\newcommand{\NMinCI}{\ensuremath{0.05283}}
% NNumCertificates [cert, exact] constants.json :: meta.summary.n_cert
\newcommand{\NNumCertificates}{\ensuremath{26}}
% NMinRhoC [cert, lower] constants.json :: meta.closure.rho_c_min
\newcommand{\NMinRhoC}{\ensuremath{0.04711}}
% NHTwoBaseZero [cert, upper] FINAL_manifest_v33.json :: results.M9_H2_tightened.rows.0.base
\newcommand{\NHTwoBaseZero}{\ensuremath{8.011\cdot10^{-11}}}
% NHTwoBaseOne [cert, upper] FINAL_manifest_v33.json :: results.M9_H2_tightened.rows.1.base
\newcommand{\NHTwoBaseOne}{\ensuremath{1.336\cdot10^{-10}}}
% NHTwoBaseTwo [cert, upper] FINAL_manifest_v33.json :: results.M9_H2_tightened.rows.2.base
\newcommand{\NHTwoBaseTwo}{\ensuremath{2.936\cdot10^{-10}}}
% NHTwoErrWZero [cert, upper] FINAL_manifest_v33.json :: results.M9_H2_tightened.rows.0.with_error
\newcommand{\NHTwoErrWZero}{\ensuremath{8.014\cdot10^{-11}}}
% NHTwoMarginWZero [cert, margin] FINAL_manifest_v33.json :: results.M9_H2_tightened.rows.0.margin
\newcommand{\NHTwoMarginWZero}{\ensuremath{124.78}}
% NHTwoErrWOne [cert, upper] FINAL_manifest_v33.json :: results.M9_H2_tightened.rows.1.with_error
\newcommand{\NHTwoErrWOne}{\ensuremath{4.706\cdot10^{-10}}}
% NHTwoMarginWOne [cert, margin] FINAL_manifest_v33.json :: results.M9_H2_tightened.rows.1.margin
\newcommand{\NHTwoMarginWOne}{\ensuremath{2124.9}}
% NHTwoErrWTwo [cert, upper] FINAL_manifest_v33.json :: results.M9_H2_tightened.rows.2.with_error
\newcommand{\NHTwoErrWTwo}{\ensuremath{1.559\cdot10^{-6}}}
% NHTwoMarginWTwo [cert, margin] FINAL_manifest_v33.json :: results.M9_H2_tightened.rows.2.margin
\newcommand{\NHTwoMarginWTwo}{\ensuremath{641.53}}
% NHThreeBaseZero [cert, upper] FINAL_manifest_v33.json :: results.M9_H2f.rows.0.base
\newcommand{\NHThreeBaseZero}{\ensuremath{3.127\cdot10^{-7}}}
% NHThreeBaseOne [cert, upper] FINAL_manifest_v33.json :: results.M9_H2f.rows.1.base
\newcommand{\NHThreeBaseOne}{\ensuremath{1.004\cdot10^{-6}}}
% NHThreeBaseTwo [cert, upper] FINAL_manifest_v33.json :: results.M9_H2f.rows.2.base
\newcommand{\NHThreeBaseTwo}{\ensuremath{2.431\cdot10^{-6}}}
% NHThreeErrZero [cert, upper] FINAL_manifest_v33.json :: results.M9_H2f.rows.0.with_error
\newcommand{\NHThreeErrZero}{\ensuremath{3.1363\cdot10^{-7}}}
% NHThreeMarginZero [cert, margin] FINAL_manifest_v33.json :: results.M9_H2f.rows.0.margin
\newcommand{\NHThreeMarginZero}{\ensuremath{3.1885}}
% NHThreeErrOne [cert, upper] FINAL_manifest_v33.json :: results.M9_H2f.rows.1.with_error
\newcommand{\NHThreeErrOne}{\ensuremath{1.0034\cdot10^{-6}}}
% NHThreeMarginOne [cert, margin] FINAL_manifest_v33.json :: results.M9_H2f.rows.1.margin
\newcommand{\NHThreeMarginOne}{\ensuremath{4.9834}}
% NHThreeErrTwo [cert, upper] FINAL_manifest_v33.json :: results.M9_H2f.rows.2.with_error
\newcommand{\NHThreeErrTwo}{\ensuremath{2.4351\cdot10^{-6}}}
% NHThreeMarginTwo [cert, margin] FINAL_manifest_v33.json :: results.M9_H2f.rows.2.margin
\newcommand{\NHThreeMarginTwo}{\ensuremath{4.1067}}
% NHFourBaseZero [cert, upper] FINAL_manifest_v33.json :: results.M9_base_sums.combined.windows_max.0
\newcommand{\NHFourBaseZero}{\ensuremath{3.48\cdot10^{-8}}}
% NHFourErrZero [cert, upper] FINAL_manifest_v33.json :: results.M9_H3_H2v32.rows.win.0.with_error_arb
\newcommand{\NHFourErrZero}{\ensuremath{3.5087\cdot10^{-8}}}
% NHFourMarginZero [cert, margin] FINAL_manifest_v33.json :: results.M9_H3_H2v32.rows.win.0.margin_arb
\newcommand{\NHFourMarginZero}{\ensuremath{5.700}}
% NHFourBaseOne [cert, upper] FINAL_manifest_v33.json :: results.M9_base_sums.combined.windows_max.1
\newcommand{\NHFourBaseOne}{\ensuremath{1.27\cdot10^{-7}}}
% NHFourErrOne [cert, upper] FINAL_manifest_v33.json :: results.M9_H3_H2v32.rows.win.1.with_error_arb
\newcommand{\NHFourErrOne}{\ensuremath{1.2616\cdot10^{-7}}}
% NHFourMarginOne [cert, margin] FINAL_manifest_v33.json :: results.M9_H3_H2v32.rows.win.1.margin_arb
\newcommand{\NHFourMarginOne}{\ensuremath{79.26}}
% NHFourBaseTwo [cert, upper] FINAL_manifest_v33.json :: results.M9_base_sums.combined.windows_max.2
\newcommand{\NHFourBaseTwo}{\ensuremath{3.76\cdot10^{-7}}}
% NHFourErrTwo [cert, upper] FINAL_manifest_v33.json :: results.M9_H3_H2v32.rows.win.2.with_error_arb
\newcommand{\NHFourErrTwo}{\ensuremath{3.7724\cdot10^{-7}}}
% NHFourMarginTwo [cert, margin] FINAL_manifest_v33.json :: results.M9_H3_H2v32.rows.win.2.margin_arb
\newcommand{\NHFourMarginTwo}{\ensuremath{26.50}}
% NHFourBaseThree [cert, upper] FINAL_manifest_v33.json :: results.M9_base_sums.combined.windows_max.3
\newcommand{\NHFourBaseThree}{\ensuremath{9.71\cdot10^{-7}}}
% NHFourErrThree [cert, upper] FINAL_manifest_v33.json :: results.M9_H3_H2v32.rows.win.3.with_error_arb
\newcommand{\NHFourErrThree}{\ensuremath{1.2822\cdot10^{-6}}}
% NHFourMarginThree [cert, margin] FINAL_manifest_v33.json :: results.M9_H3_H2v32.rows.win.3.margin_arb
\newcommand{\NHFourMarginThree}{\ensuremath{7.799}}
% NHFourYbarLimit [cert, lower] FINAL_manifest_v33.json :: results.M9_H3_H2v32.binding_Ybar_limit
\newcommand{\NHFourYbarLimit}{\ensuremath{1.575\cdot10^{-11}}}
% NHFourYbarMargin [cert, margin] FINAL_manifest_v33.json :: results.M9_H3_H2v32.binding_Ybar_limit
\newcommand{\NHFourYbarMargin}{\ensuremath{15.75}}
% NHFourKfourBase [cert, upper] amp_A1_w180.json :: inputs.BASE_T1.K4
\newcommand{\NHFourKfourBase}{\ensuremath{3.1633\cdot10^{-4}}}
% NHSixConeSlope [cert, upper] FINAL_manifest_v33.json :: results.cone_fixedL.checks.slope.value
\newcommand{\NHSixConeSlope}{\ensuremath{0.03029}}
% NHSixConeSlopeMargin [cert, margin] FINAL_manifest_v33.json :: results.cone_fixedL.checks.slope.margin
\newcommand{\NHSixConeSlopeMargin}{\ensuremath{2.641}}
% NHSixConeKtwo [cert, upper] FINAL_manifest_v33.json :: results.cone_fixedL.checks.K2.value
\newcommand{\NHSixConeKtwo}{\ensuremath{0.009876}}
% NHSixConeKtwoMargin [cert, margin] FINAL_manifest_v33.json :: results.cone_fixedL.checks.K2.margin
\newcommand{\NHSixConeKtwoMargin}{\ensuremath{3.037}}
% NHSixConeNthree [cert, upper] FINAL_manifest_v33.json :: results.cone_fixedL.checks.N3.value
\newcommand{\NHSixConeNthree}{\ensuremath{1.461\cdot10^{-4}}}
% NHSixConeNthreeMargin [cert, margin] FINAL_manifest_v33.json :: results.cone_fixedL.checks.N3.margin
\newcommand{\NHSixConeNthreeMargin}{\ensuremath{27.38}}
% NHSixConeHFive [cert, upper] FINAL_manifest_v33.json :: results.H4_fixedL.total
\newcommand{\NHSixConeHFive}{\ensuremath{1.2865\cdot10^{-5}}}
% NHSixConeHFiveMargin [cert, margin] FINAL_manifest_v33.json :: results.H4_fixedL.margin
\newcommand{\NHSixConeHFiveMargin}{\ensuremath{1119}}
% NHSixConeTstarNeeded [cert, upper] FINAL_manifest_v33.json :: results.H7_fixedL.ratio_upper
\newcommand{\NHSixConeTstarNeeded}{\ensuremath{1.6000015\cdot10^{-7}}}
% NHSixConeTstarLower [cert, lower] FINAL_manifest_v33.json :: results.H7_fixedL.t_star_lower
\newcommand{\NHSixConeTstarLower}{\ensuremath{3.68334\cdot10^{-7}}}
% NHSixConeTstarMargin [cert, margin] FINAL_manifest_v33.json :: results.H7_fixedL.margin_tstar_over_ratio
\newcommand{\NHSixConeTstarMargin}{\ensuremath{2.302}}
% NHSixConeHFiveRatioPeak [cert, upper] h4loc_A1_w180.json :: windows.constants/amp_A1_w180.json.worst.ratio
\newcommand{\NHSixConeHFiveRatioPeak}{\ensuremath{0.2393}}
% NHSixConeHOneSup [cert, upper] FINAL_manifest_v33.json :: results.H1_fixedL.sup_h0_max
\newcommand{\NHSixConeHOneSup}{\ensuremath{0.2163}}
% NHSixConeHOneH [cert, upper] FINAL_manifest_v33.json :: results.H1_fixedL.H_max
\newcommand{\NHSixConeHOneH}{\ensuremath{0.2185}}
% NHSixFarHFiveRatio [cert, upper] h4loc_v32.json :: windows.constants/amp_A1_w180.json.cylinders.18.ratio
\newcommand{\NHSixFarHFiveRatio}{\ensuremath{0.1761}}
% NHSixFarHFiveMarginOld [cert, margin] h4loc_v32.json :: windows.constants/amp_A1_w180.json.cylinders.18.ratio
\newcommand{\NHSixFarHFiveMarginOld}{\ensuremath{5.68}}
% NHSixFarHFiveMargin [cert, margin] FINAL_manifest_v33.json :: results.M9_H4_far.min_margin
\newcommand{\NHSixFarHFiveMargin}{\ensuremath{15.58}}
% NHSixConeHFiveMarginOld [cert, margin] h4loc_A1_w180.json :: windows.constants/amp_A1_w180.json.worst.ratio
\newcommand{\NHSixConeHFiveMarginOld}{\ensuremath{4.18}}
% NHSixFarKtwo [cert, upper] t9_far_v32.json :: h2
\newcommand{\NHSixFarKtwo}{\ensuremath{1.797\cdot10^{-4}}}
% NHSixFarKtwoThr [hand, exact] t9_far_v32.json :: thr.h2
\newcommand{\NHSixFarKtwoThr}{\ensuremath{3.53\cdot10^{-4}}}
% NHSixFarKtwoMargin [cert, margin] t9_far_v32.json :: margins.h2
\newcommand{\NHSixFarKtwoMargin}{\ensuremath{1.964}}
% NHSixFarKthree [cert, upper] t9_far_v32.json :: h3
\newcommand{\NHSixFarKthree}{\ensuremath{2.595\cdot10^{-8}}}
% NHSixFarKthreeThr [hand, exact] t9_far_v32.json :: thr.h3
\newcommand{\NHSixFarKthreeThr}{\ensuremath{1.17\cdot10^{-5}}}
% NHSixFarKthreeMargin [cert, margin] t9_far_v32.json :: margins.h3
\newcommand{\NHSixFarKthreeMargin}{\ensuremath{450.9}}
% NHSixFarKfour [cert, upper] t9_far_v32.json :: h4
\newcommand{\NHSixFarKfour}{\ensuremath{6.964\cdot10^{-9}}}
% NHSixFarKfourThr [hand, exact] t9_far_v32.json :: thr.h4
\newcommand{\NHSixFarKfourThr}{\ensuremath{1.03\cdot10^{-7}}}
% NHSixFarKfourMargin [cert, margin] t9_far_v32.json :: margins.h4
\newcommand{\NHSixFarKfourMargin}{\ensuremath{14.79}}
% NHSixFarSlope [cert, upper] t9_far_v32.json :: slope
\newcommand{\NHSixFarSlope}{\ensuremath{0.04153}}
% NHSixFarSlopeThr [hand, exact] t9_far_v32.json :: thr.slope
\newcommand{\NHSixFarSlopeThr}{\ensuremath{0.0800}}
% NHSixFarSlopeMargin [cert, margin] t9_far_v32.json :: margins.slope
\newcommand{\NHSixFarSlopeMargin}{\ensuremath{1.926}}
% NHSixFarNthree [cert, upper] t9_far_v32.json :: L2
\newcommand{\NHSixFarNthree}{\ensuremath{5.516\cdot10^{-7}}}
% NHSixFarNthreeThr [hand, exact] t9_far_v32.json :: thr.L2
\newcommand{\NHSixFarNthreeThr}{\ensuremath{3.32\cdot10^{-6}}}
% NHSixFarNthreeMargin [cert, margin] t9_far_v32.json :: margins.L2
\newcommand{\NHSixFarNthreeMargin}{\ensuremath{6.019}}
% NFarTstarMargin [cert, margin] lemmaO_v32_consts.json :: far_v32.case.margin_factor_lower
\newcommand{\NFarTstarMargin}{\ensuremath{15.5}}
% NFarTstarLower [cert, lower] lemmaO_v32_consts.json :: far_v32.case.t_star_lower
\newcommand{\NFarTstarLower}{\ensuremath{9.5108\cdot10^{-8}}}
% NFarTstarNeeded [cert, upper] lemmaO_v32_consts.json :: far_v32.case.needed
\newcommand{\NFarTstarNeeded}{\ensuremath{6.134\cdot10^{-9}}}
% NGammaModes [cert, exact] gamma0_final.json :: M
\newcommand{\NGammaModes}{\ensuremath{3762}}
% NGammaHalfwidth [cert, upper] gamma0_final.json :: b_j_halfwidth
\newcommand{\NGammaHalfwidth}{\ensuremath{8.138\cdot10^{-24}}}
% NGammaContraction [cert, upper] gamma0_final.json :: krawczyk_contraction
\newcommand{\NGammaContraction}{\ensuremath{3.01\cdot10^{-12}}}
% NGammaArcs [cert, exact] gamma0_final.json :: embedded.arcs
\newcommand{\NGammaArcs}{\ensuremath{701}}
% NGammaMaxTurn [cert, upper] gamma0_final.json :: embedded.maxturn
\newcommand{\NGammaMaxTurn}{\ensuremath{0.04997}}
% NGammaMaxPair [cert, upper] gamma0_final.json :: embedded.maxpair
\newcommand{\NGammaMaxPair}{\ensuremath{0.09597}}
% NGammaMinSep [cert, lower] gamma0_final.json :: embedded.minsep
\newcommand{\NGammaMinSep}{\ensuremath{8.89\cdot10^{-3}}}
% NGammaHalfGap [cert, lower] gamma0_final.json :: embedded.gap_cert
\newcommand{\NGammaHalfGap}{\ensuremath{0.0248219}}
% NPhaseISlabs [cert, exact] status.json :: slab
\newcommand{\NPhaseISlabs}{\ensuremath{96}}
% NPhaseIModes [cert, exact] status.json :: last_slab.M
\newcommand{\NPhaseIModes}{\ensuremath{2800}}
% NPhaseIKmaxLower [cert, lower] status.json :: kmax_enclosure.0
\newcommand{\NPhaseIKmaxLower}{\ensuremath{29.352874545}}
% NPhaseIKmaxUpper [cert, upper] status.json :: kmax_enclosure.1
\newcommand{\NPhaseIKmaxUpper}{\ensuremath{29.352874546}}
% NPhaseITLower [cert, lower] status.json :: t
\newcommand{\NPhaseITLower}{\ensuremath{4.597850539485830}}
% NPhaseITUpper [cert, upper] status.json :: t
\newcommand{\NPhaseITUpper}{\ensuremath{4.597850539485831}}
% NPeakSqrtLam [cert, upper] rate_table_fin12f4_a0.00015.json :: sqrtL_max
\newcommand{\NPeakSqrtLam}{\ensuremath{5.308\cdot10^{-8}}}
% NTailRate [cert, lower] rate_table_fin12f4_a0.00015.json :: tail.rho
\newcommand{\NTailRate}{\ensuremath{2.693\cdot10^{-3}}}
% NFloorSqrtLam [cert, upper] rate_table_fin12f4_a0.00015.json :: tail.sqrt_LF
\newcommand{\NFloorSqrtLam}{\ensuremath{9.872\cdot10^{-10}}}
% NAlphaFloor [cert, upper] rate_table_fin12f4_a0.00015.json :: tail.alpha_F
\newcommand{\NAlphaFloor}{\ensuremath{6.867\cdot10^{-7}}}
% NSupBetaVerifier [cert, upper] final_alpha_points_5e-8.json :: 1.5e-4.sup_beta
\newcommand{\NSupBetaVerifier}{\ensuremath{3.33811\cdot10^{-9}}}
% NMOneRamp [cert, upper] final_alpha_points_5e-8.json :: 1.5e-4.M1
\newcommand{\NMOneRamp}{\ensuremath{0.4328}}
% NSITwoPrime [cert, upper] FINAL_manifest_v33.json :: results.F12_sliver.sI2_prime
\newcommand{\NSITwoPrime}{\ensuremath{0.015457}}
% NLoopGainCoarse [cert, upper] U_cert_v35.json :: loop_gain.2.total
\newcommand{\NLoopGainCoarse}{\ensuremath{0.03904}}
% NRateMarginCoarse [cert, margin] U_cert_v35.json :: rate_margin
\newcommand{\NRateMarginCoarse}{\ensuremath{28.2}}
% NSOneEpsLower [cert, lower] bracket2.json :: eps_a
\newcommand{\NSOneEpsLower}{\ensuremath{1.5346198892965\cdot10^{-4}}}
% NSOneEpsUpper [cert, upper] bracket2.json :: eps_b
\newcommand{\NSOneEpsUpper}{\ensuremath{1.5346198899299\cdot10^{-4}}}
% NSOneRowsTip [cert, exact] ext_tip_check.json :: rows
\newcommand{\NSOneRowsTip}{\ensuremath{51246}}
% NSOneTipWidth [cert, upper] ext_tip_check.json :: max_pointwise_width.th
\newcommand{\NSOneTipWidth}{\ensuremath{3.2\cdot10^{-12}}}
% NSOneExtWidth [cert, upper] ext813_check.json :: max_pointwise_width.p
\newcommand{\NSOneExtWidth}{\ensuremath{5.7\cdot10^{-10}}}
% NSOneRowsMC [cert, exact] check_outA5.json :: prof_tip_vs_MC.rows_compared
\newcommand{\NSOneRowsMC}{\ensuremath{1025}}
% NSOneRowsMB [cert, exact] check_outA5.json :: ext813_vs_MB.rows_compared
\newcommand{\NSOneRowsMB}{\ensuremath{501}}
% NSOneMinMargin [cert, margin] check_outA5.json :: min:*.min_relative_margin.*
\newcommand{\NSOneMinMargin}{\ensuremath{1.2\cdot10^{-9}}}
% NTwoSegments [cert, exact] contour_result.json :: len:segments
\newcommand{\NTwoSegments}{\ensuremath{5058}}
% NTwoInteriorBoxes [cert, exact] interior_result.json :: len:
\newcommand{\NTwoInteriorBoxes}{\ensuremath{174}}
% NTwoContourBoxes [cert, exact] eta_cache.json :: len:
\newcommand{\NTwoContourBoxes}{\ensuremath{202}}
% NTwoEtaMax [cert, upper] eta_cache.json :: max:*.eta
\newcommand{\NTwoEtaMax}{\ensuremath{0.1163}}
% NTwoLipMax [cert, upper] eta_cache.json :: max:*.Lip
\newcommand{\NTwoLipMax}{\ensuremath{0.272}}
% NTwoLambdaZero [est, upper] cert_b0.25_th0.0_tau0.03.json :: Lam
\newcommand{\NTwoLambdaZero}{\ensuremath{1.790819}}
% NTwoLambdaOne [est, upper] cert_b0.25_th0.48_tau0.03.json :: Lam
\newcommand{\NTwoLambdaOne}{\ensuremath{1.768282}}
% NSOneDTightLower [doc, lower] doc:README.md :: D ∈ [2.50080627654, 2.50080627668]
\newcommand{\NSOneDTightLower}{\ensuremath{2.50080627654}}
% NSOneDTightUpper [doc, upper] doc:README.md :: D ∈ [2.50080627654, 2.50080627668]
\newcommand{\NSOneDTightUpper}{\ensuremath{2.50080627668}}
% NSOneKTightLower [doc, lower] doc:README.md :: k(0) ∈ [1.53982063762, 1.53982063822]
\newcommand{\NSOneKTightLower}{\ensuremath{1.53982063762}}
% NSOneKTightUpper [doc, upper] doc:README.md :: k(0) ∈ [1.53982063762, 1.53982063822]
\newcommand{\NSOneKTightUpper}{\ensuremath{1.53982063822}}
% NSOneDCondLower [doc, lower] doc:README.md :: D = −p(tip) ∈ [2.500806276619695, 2.500806276620049]
\newcommand{\NSOneDCondLower}{\ensuremath{2.500806276619695}}
% NSOneDCondUpper [doc, upper] doc:README.md :: D = −p(tip) ∈ [2.500806276619695, 2.500806276620049]
\newcommand{\NSOneDCondUpper}{\ensuremath{2.500806276620049}}
% NSOneKCondLower [doc, lower] doc:README.md :: k(0) ∈ [1.539820637871836, 1.5398206378721015]
\newcommand{\NSOneKCondLower}{\ensuremath{1.539820637871836}}
% NSOneKCondUpper [doc, upper] doc:README.md :: k(0) ∈ [1.539820637871836, 1.5398206378721015]
\newcommand{\NSOneKCondUpper}{\ensuremath{1.5398206378721015}}
% NSOneTauTipLower [doc, lower] doc:README.md :: τ_tip ∈ [10.249077651120778, 10.249077651121361]
\newcommand{\NSOneTauTipLower}{\ensuremath{10.249077651120778}}
% NSOneTauTipUpper [doc, upper] doc:README.md :: τ_tip ∈ [10.249077651120778, 10.249077651121361]
\newcommand{\NSOneTauTipUpper}{\ensuremath{10.249077651121361}}
% NSOneStateUncert [doc, upper] doc:README.md :: State uncertainty at σ = 8: 4.4e−13
\newcommand{\NSOneStateUncert}{\ensuremath{4.4\cdot10^{-13}}}
% NFarBoxEprime [doc, margin] doc:STEP4_RESULT.md :: (6.4287, 4.6917, 1.9940, 1.69999995)
\newcommand{\NFarBoxEprime}{\ensuremath{6.42}}
% NFarBoxXEsecond [doc, margin] doc:STEP4_RESULT.md :: (6.4287, 4.6917, 1.9940, 1.69999995)
\newcommand{\NFarBoxXEsecond}{\ensuremath{4.69}}
% NFarBoxXsqEthird [doc, margin] doc:STEP4_RESULT.md :: (6.4287, 4.6917, 1.9940, 1.69999995)
\newcommand{\NFarBoxXsqEthird}{\ensuremath{1.99}}
% NFarBoxE [doc, margin] doc:STEP4_RESULT.md :: (6.4287, 4.6917, 1.9940, 1.69999995)
\newcommand{\NFarBoxE}{\ensuremath{1.69}}
% NHFourKfourErr [doc, upper] doc:STEP4_RESULT.md :: | 5.60e-3 against 0.06 | ×10.7 |
\newcommand{\NHFourKfourErr}{\ensuremath{5.60\cdot10^{-3}}}
% NHFourKfourMargin [doc, margin] doc:STEP4_RESULT.md :: | 5.60e-3 against 0.06 | ×10.7 |
\newcommand{\NHFourKfourMargin}{\ensuremath{10.7}}
% NCoercivityCoefficient [doc, upper] doc:step3_proof_v32.tex :: \norm{e^{\tau\mathcal L}u}_\rho\le14.8\,e^{-0.0213\tau}\norm
\newcommand{\NCoercivityCoefficient}{\ensuremath{14.8}}
% NCoercivityRate [doc, lower] doc:step3_proof_v32.tex :: \norm{e^{\tau\mathcal L}u}_\rho\le14.8\,e^{-0.0213\tau}\norm
\newcommand{\NCoercivityRate}{\ensuremath{0.0213}}
% NRegionRampStart [hand, exact] hand :: 
\newcommand{\NRegionRampStart}{\ensuremath{14}}
% NRegionInterface [hand, exact] hand :: 
\newcommand{\NRegionInterface}{\ensuremath{28}}
% NRegionShellEnd [hand, exact] hand :: 
\newcommand{\NRegionShellEnd}{\ensuremath{60}}
% NRegionSegEnd [hand, exact] hand :: 
\newcommand{\NRegionSegEnd}{\ensuremath{350}}
% NRegionAnnEnd [hand, exact] hand :: 
\newcommand{\NRegionAnnEnd}{\ensuremath{700}}
% NZoneIIILow [hand, exact] hand :: 
\newcommand{\NZoneIIILow}{\ensuremath{253}}
% NZoneIIIHigh [hand, exact] hand :: 
\newcommand{\NZoneIIIHigh}{\ensuremath{379.5}}
% NConeCentreMin [hand, exact] hand :: 
\newcommand{\NConeCentreMin}{\ensuremath{300}}
% NConeCentreMax [hand, exact] hand :: 
\newcommand{\NConeCentreMax}{\ensuremath{380}}
% NFarCentreMin [hand, exact] hand :: 
\newcommand{\NFarCentreMin}{\ensuremath{490}}
% NConeL [hand, exact] hand :: 
\newcommand{\NConeL}{\ensuremath{50}}
% NConeX [hand, exact] hand :: 
\newcommand{\NConeX}{\ensuremath{3}}
% NFarL [hand, exact] hand :: 
\newcommand{\NFarL}{\ensuremath{113}}
% NFarX [hand, exact] hand :: 
\newcommand{\NFarX}{\ensuremath{2}}
}%
  {\PackageWarning{paper}{numbers/numbers_paper.tex not found (unit U0 not finished)}}

% ======================================================================
% NOTATION MACROS (OUTLINE 5.2 and notation.md).  One macro per symbol
% whose spelling is fixed by the notation table.  Plain letters (a, b, f, ...)
% have no macro.
% ======================================================================
% basic
\newcommand{\R}{\mathbb{R}}
\newcommand{\C}{\mathbb{C}}
\newcommand{\N}{\mathbb{N}}
\newcommand{\Z}{\mathbb{Z}}
\newcommand{\eps}{\varepsilon}
\newcommand{\ip}[2]{\langle #1,#2\rangle}
\newcommand{\abs}[1]{\lvert #1\rvert}
\newcommand{\norm}[1]{\lVert #1\rVert}
\newcommand{\dd}{\mathrm d}
\newcommand{\re}{\operatorname{Re}}
\newcommand{\im}{\operatorname{Im}}
% flow and curves
\newcommand{\arclen}{\mathfrak{s}}                 % arclength of the flow curve
\newcommand{\kss}{\kappa_{\mathfrak{ss}}}          % kappa_{ss} in the flow equation
\newcommand{\Area}{\operatorname{Area}}
\newcommand{\Len}{\mathcal L}                      % length of a curve
\newcommand{\spp}{\operatorname{sp}}                 % sup quantities sp^v_0, sp^v_1, sp^f_1 (were S^v_0, S^v_1, S^f_1)
\newcommand{\rot}{\operatorname{rot}}              % turning number
\newcommand{\Tsing}{T_*}                           % singular time
\newcommand{\Tmax}{T_{\max}}                       % maximal time of existence
\newcommand{\pstar}{p_*}                           % image point of the singularity
\newcommand{\Gz}{\Gamma_0}                         % initial curve
\newcommand{\vth}{\vartheta}                       % parameter of Gamma_0
\newcommand{\dclk}{\delta_t}                       % clock defect of step 4
\newcommand{\fall}{\varpi}                         % Fourier rates varpi_n = n^4-n^2
% rescaling
\newcommand{\lamone}{\lambda_1}
\newcommand{\rate}{\varsigma}                      % 1/4 - a
\newcommand{\xia}{\xi_a}                           % (1-4\bar a)^{-1}
% shrinker and step 1
\newcommand{\Sone}{S_1}
\newcommand{\Gsh}{\Gamma_{\!*}}                    % profile curve of S_1
\newcommand{\cSet}{\mathcal S}                     % certified set of shrinkers
\newcommand{\sigtip}{\sigma_{\rm tip}}
\newcommand{\Xo}{X_0}                              % 2.2490776515
\newcommand{\Dtip}{D}
\newcommand{\ktip}{k_0}
\newcommand{\shootM}{M}                            % shooting function
\newcommand{\Ypar}{\mathcal Y}                     % step-1 parallelepiped
\newcommand{\axis}{\Sigma}                         % symmetry axis of S_1
\newcommand{\sigc}{\sigma_{\rm c}}
\newcommand{\sigrev}{\sigma^{\rm rev}}
% linear theory
\newcommand{\Hrho}{\mathcal H_\rho}
\newcommand{\muslow}{\mu_3}
\newcommand{\rre}{\mathfrak r_{\rm re}}
\newcommand{\rim}{\mathfrak r_{\rm im}}
\newcommand{\Evev}{\mathrm{Ev}_{\rm even}}
\newcommand{\Evod}{\mathrm{Ev}_{\rm odd}}
\newcommand{\gvec}{\mathbf g}                      % translation vector in step 2
\newcommand{\cons}[1]{\mathfrak f_{#1}}            % constraint functions
\newcommand{\phirho}{\varphi_\rho}
% energies
\newcommand{\Lam}{\Lambda}
\newcommand{\En}{E}
\newcommand{\EnI}{E_I}
\newcommand{\muI}{\mu_I}                           % 0.08
\newcommand{\epsE}{\varepsilon_E}                  % 1e-5
\newcommand{\epsI}{\varepsilon_I}                  % 3e-5
\newcommand{\wn}[1]{\mathsf w_{#1}}                % ||w~^{(j)}||_rho
\newcommand{\vI}{v_I}                              % zeta w'
\newcommand{\muerr}{\mu_{\rm err}}                 % step-4 error-energy weight 0.036
% far field
\newcommand{\pU}{p_U}                              % 1.0246
\newcommand{\Uone}{U_1}
\newcommand{\sm}{\operatorname{sm}}                % smoothstep t^3(10-15t+6t^2)
\newcommand{\rc}{r_c}                              % centre radius (Lemma U)
\newcommand{\rchi}{r_\chi}                         % 253
\newcommand{\Iw}{I_w}                              % hand-off windows [y_w,1.5y_w]
\newcommand{\qa}{\mathfrak q_1}                    % quasilinear coefficients of the slope equation
\newcommand{\qb}{\mathfrak q_2}
% cylinders (Lemma O)
\newcommand{\cyl}{\ell}                            % cylinder length unit
\newcommand{\MO}{M_O}
\newcommand{\KO}{K_O}
\newcommand{\Hcyl}{H_{\rm cyl}}
\newcommand{\XiO}{\Xi_O}
\newcommand{\fO}{f_O}
% step 4
\newcommand{\Ybar}{\bar Y}
\newcommand{\Lt}{\widetilde L}                     % linearisation about the step-4 base
\newcommand{\Trans}{\Pi}                           % transport integral
\newcommand{\resid}{\mathrm{res}}
\newcommand{\Resid}[1]{\mathrm{Res}_{#1}}
\newcommand{\boxtip}{\Box_{\rm tip}}
\newcommand{\boxfar}{\Box_{\rm far}}
\newcommand{\Eenv}{\mathcal E_{\rm env}}
\newcommand{\betab}{\beta_{\rm base}}
\newcommand{\tauI}{\tau^{\rm I}}                   % Phase I clock
% names of families of conditions: \Hh{3} = (H3), \Oh{4} = (O4), ...
\newcommand{\Hh}[1]{\textup{(H#1)}}
\newcommand{\Oh}[1]{\textup{(O#1)}}
\newcommand{\Lh}[1]{\textup{(L#1)}}
\newcommand{\Bh}[1]{\textup{(B#1)}}
\newcommand{\Kh}[1]{\textup{(K#1)}}
\newcommand{\Wh}[1]{\textup{(W#1)}}
\newcommand{\PhaseI}{Phase~I}
\newcommand{\PhaseII}{Phase~II}
% margin: threshold/bound (always)
\newcommand{\mgn}[1]{{\times}#1}

% ---- figures: placeholder macro so that drafts compile before F1-F7 exist ----
\newcommand{\figplaceholder}[1]{\fbox{\parbox{0.8\linewidth}{\centering\small[figure #1 to be supplied]}}}
 ...
%   paper_v1/supplement.tex  \documentclass[11pt]{amsart}
%                            \newif\ifsupplement \supplementtrue
%                            % ======================================================================
% preamble.tex -- shared preamble of the main paper and the supplement.
% Owner: coordinator.  Drafters add macros by request only (see DRAFTER_GUIDE.md).
%
% USE.  This file contains NO \documentclass.  The two master files are
%   paper_v1/main.tex        \documentclass[11pt]{amsart} % ======================================================================
% preamble.tex -- shared preamble of the main paper and the supplement.
% Owner: coordinator.  Drafters add macros by request only (see DRAFTER_GUIDE.md).
%
% USE.  This file contains NO \documentclass.  The two master files are
%   paper_v1/main.tex        \documentclass[11pt]{amsart} \input{preamble} ...
%   paper_v1/supplement.tex  \documentclass[11pt]{amsart}
%                            \newif\ifsupplement \supplementtrue
%                            \input{preamble} ...
% (the main paper does nothing special: \ifsupplement defaults to false).
% Compile from paper_v1/ so that \input{numbers/numbers_paper.tex} is found.
% ======================================================================

% ---- which document are we in ----------------------------------------
\ifdefined\ifsupplement\else\newif\ifsupplement\fi

% ---- packages ---------------------------------------------------------
\usepackage[T1]{fontenc}
\usepackage[utf8]{inputenc}
\usepackage{lmodern}
\usepackage[margin=1.1in]{geometry}
\usepackage{amsmath,amssymb,amsthm,mathtools,bm}
\usepackage{booktabs,array,longtable,tabularx,multirow}
\usepackage{enumitem}
\usepackage{graphicx}
\usepackage{xcolor}
\usepackage{tikz}
\usetikzlibrary{positioning,arrows.meta,calc}
% xr-hyper must come before hyperref
\usepackage{xr-hyper}
\usepackage[colorlinks=true,linkcolor=blue!50!black,citecolor=green!35!black,urlcolor=blue!50!black]{hyperref}

\raggedbottom
\setlength{\emergencystretch}{2em}
\setlist[enumerate]{leftmargin=*,itemsep=2pt,topsep=3pt}
\setlist[itemize]{leftmargin=*,itemsep=2pt,topsep=3pt}

% ---- numbering --------------------------------------------------------
% Main paper: sections 1,2,...; theorems 2.3 (by section); equations (2.3).
% Supplement: the Parts A..K are the \section's, numbered by letter, so
% theorems are A.1, E.4, ... (do not confuse "Theorem A" of the paper, which
% has no dot, with "Theorem A.1" of Part A).
\ifsupplement
  \renewcommand{\thesection}{\Alph{section}}
\fi
\numberwithin{equation}{section}

% ---- theorem environments (numbered by section, one shared counter) ----
\theoremstyle{plain}
\newtheorem{theorem}{Theorem}[section]
\newtheorem{proposition}[theorem]{Proposition}
\newtheorem{lemma}[theorem]{Lemma}
\newtheorem{corollary}[theorem]{Corollary}
\newtheorem{maintheorem}[theorem]{Main Theorem}
% Theorems A, B, C, D of the paper: own counter, lettered, no section prefix
\newtheorem{thmlet}{Theorem}
\renewcommand{\thethmlet}{\Alph{thmlet}}
\theoremstyle{definition}
\newtheorem{definition}[theorem]{Definition}
\newtheorem{convention}[theorem]{Convention}
\newtheorem{example}[theorem]{Example}
\theoremstyle{remark}
\newtheorem{remark}[theorem]{Remark}

% A list of hypotheses labelled (H1),(H2),... : \begin{hyplist}{H} ... \end{hyplist}
% (also {O}, {L}, {B}, {K}, {W} for the families of section 5.3 of OUTLINE.md).
\newenvironment{hyplist}[1]%
  {\begin{enumerate}[label=\textup{(#1\arabic*)},leftmargin=3em,itemsep=2pt]}%
  {\end{enumerate}}

% ---- trust classes (OUTLINE 4.1, 5.1.3) -------------------------------
% Supplement: explicit tags.  Main paper: certified is the default and
% unmarked; every non-certified number is marked "(floating point)".
\ifsupplement
  \newcommand{\tcert}{\textup{[cert]}}   % arb/outward-rounded, pinned script and key
  \newcommand{\test}{\textup{[est]}}     % floating point or estimate; never in a theorem
  \newcommand{\thand}{\textup{[hand]}}   % proved by hand in the text
  \newcommand{\tlit}{\textup{[lit]}}     % cited from the literature
\else
  \newcommand{\tcert}{}
  \newcommand{\test}{\unskip\ \textup{(floating point)}}
  \newcommand{\thand}{}
  \newcommand{\tlit}{}
\fi

% ---- gaps: visible, counted, an error in a final build -----------------
% Usage: \GAP{one sentence saying what is missing}   (and an entry in ISSUES.md)
\newcounter{gapcount}
\newif\iffinalpaper   % \finalpapertrue makes every \GAP a hard error
\newcommand{\GAP}[1]{\stepcounter{gapcount}\typeout{GAP-MARKER \thegapcount}%
  \iffinalpaper\PackageError{paper}{\string\GAP\space used in a final build}{Remove all gaps.}\fi
  {\color{red}\sffamily\bfseries[GAP~\thegapcount:\ }{\color{red}\sffamily #1}{\color{red}\sffamily\bfseries]}}
\AtEndDocument{\typeout{GAP-MARKER total: \thegapcount}}

% ---- file names and JSON keys (supplement only) -------------------------
% No internal path appears in the main paper; a use there is a warning.
\newcommand{\fn}[1]{%
  \ifsupplement\else\PackageWarning{paper}{\string\fn\space used in the main paper}\fi
  \texttt{\detokenize{#1}}}
\newcommand{\key}[1]{\texttt{\detokenize{#1}}}

% ---- cross references between main paper and supplement ----------------
% Labels carry a unit prefix (section 5.3 of OUTLINE.md and DRAFTER_GUIDE.md):
%   main paper:  intro: obj: s3: s4: s5: s6: s7: s8: s9:   (kind: thm- lem- eq- tab- fig- ...)
%   supplement:  pA: pB: ... pK:
% so all labels are globally unique and the prefixes below are only a second guard.
% \Xref{label}  : reference to the OTHER document (hyperlinked via xr-hyper)
% \Lref{label}  : reference in the SAME document (plain \ref)
\ifsupplement
  \externaldocument[M-]{main}[main.pdf]
  \newcommand{\Xref}[1]{\ref{M-#1}}
\else
  \externaldocument[S-]{supplement}[supplement.pdf]
  \newcommand{\Xref}[1]{\ref{S-#1}}
\fi
\newcommand{\Lref}[1]{\ref{#1}}
% \Supp{pE}{lem-chart} -> "Lemma E.3" style text is produced by the drafter:
%   "Lemma~\Xref{pE:lem-chart} of the supplement".  \SuppPart{E} gives "Part E".
\newcommand{\SuppPart}[1]{Part~#1 of the supplement}

% ---- numbers (generated by unit U0; macros only, no literal certified numbers) ----
\IfFileExists{numbers/numbers_paper.tex}%
  {\input{numbers/numbers_paper.tex}}%
  {\PackageWarning{paper}{numbers/numbers_paper.tex not found (unit U0 not finished)}}

% ======================================================================
% NOTATION MACROS (OUTLINE 5.2 and notation.md).  One macro per symbol
% whose spelling is fixed by the notation table.  Plain letters (a, b, f, ...)
% have no macro.
% ======================================================================
% basic
\newcommand{\R}{\mathbb{R}}
\newcommand{\C}{\mathbb{C}}
\newcommand{\N}{\mathbb{N}}
\newcommand{\Z}{\mathbb{Z}}
\newcommand{\eps}{\varepsilon}
\newcommand{\ip}[2]{\langle #1,#2\rangle}
\newcommand{\abs}[1]{\lvert #1\rvert}
\newcommand{\norm}[1]{\lVert #1\rVert}
\newcommand{\dd}{\mathrm d}
\newcommand{\re}{\operatorname{Re}}
\newcommand{\im}{\operatorname{Im}}
% flow and curves
\newcommand{\arclen}{\mathfrak{s}}                 % arclength of the flow curve
\newcommand{\kss}{\kappa_{\mathfrak{ss}}}          % kappa_{ss} in the flow equation
\newcommand{\Area}{\operatorname{Area}}
\newcommand{\Len}{\mathcal L}                      % length of a curve
\newcommand{\spp}{\operatorname{sp}}                 % sup quantities sp^v_0, sp^v_1, sp^f_1 (were S^v_0, S^v_1, S^f_1)
\newcommand{\rot}{\operatorname{rot}}              % turning number
\newcommand{\Tsing}{T_*}                           % singular time
\newcommand{\Tmax}{T_{\max}}                       % maximal time of existence
\newcommand{\pstar}{p_*}                           % image point of the singularity
\newcommand{\Gz}{\Gamma_0}                         % initial curve
\newcommand{\vth}{\vartheta}                       % parameter of Gamma_0
\newcommand{\dclk}{\delta_t}                       % clock defect of step 4
\newcommand{\fall}{\varpi}                         % Fourier rates varpi_n = n^4-n^2
% rescaling
\newcommand{\lamone}{\lambda_1}
\newcommand{\rate}{\varsigma}                      % 1/4 - a
\newcommand{\xia}{\xi_a}                           % (1-4\bar a)^{-1}
% shrinker and step 1
\newcommand{\Sone}{S_1}
\newcommand{\Gsh}{\Gamma_{\!*}}                    % profile curve of S_1
\newcommand{\cSet}{\mathcal S}                     % certified set of shrinkers
\newcommand{\sigtip}{\sigma_{\rm tip}}
\newcommand{\Xo}{X_0}                              % 2.2490776515
\newcommand{\Dtip}{D}
\newcommand{\ktip}{k_0}
\newcommand{\shootM}{M}                            % shooting function
\newcommand{\Ypar}{\mathcal Y}                     % step-1 parallelepiped
\newcommand{\axis}{\Sigma}                         % symmetry axis of S_1
\newcommand{\sigc}{\sigma_{\rm c}}
\newcommand{\sigrev}{\sigma^{\rm rev}}
% linear theory
\newcommand{\Hrho}{\mathcal H_\rho}
\newcommand{\muslow}{\mu_3}
\newcommand{\rre}{\mathfrak r_{\rm re}}
\newcommand{\rim}{\mathfrak r_{\rm im}}
\newcommand{\Evev}{\mathrm{Ev}_{\rm even}}
\newcommand{\Evod}{\mathrm{Ev}_{\rm odd}}
\newcommand{\gvec}{\mathbf g}                      % translation vector in step 2
\newcommand{\cons}[1]{\mathfrak f_{#1}}            % constraint functions
\newcommand{\phirho}{\varphi_\rho}
% energies
\newcommand{\Lam}{\Lambda}
\newcommand{\En}{E}
\newcommand{\EnI}{E_I}
\newcommand{\muI}{\mu_I}                           % 0.08
\newcommand{\epsE}{\varepsilon_E}                  % 1e-5
\newcommand{\epsI}{\varepsilon_I}                  % 3e-5
\newcommand{\wn}[1]{\mathsf w_{#1}}                % ||w~^{(j)}||_rho
\newcommand{\vI}{v_I}                              % zeta w'
\newcommand{\muerr}{\mu_{\rm err}}                 % step-4 error-energy weight 0.036
% far field
\newcommand{\pU}{p_U}                              % 1.0246
\newcommand{\Uone}{U_1}
\newcommand{\sm}{\operatorname{sm}}                % smoothstep t^3(10-15t+6t^2)
\newcommand{\rc}{r_c}                              % centre radius (Lemma U)
\newcommand{\rchi}{r_\chi}                         % 253
\newcommand{\Iw}{I_w}                              % hand-off windows [y_w,1.5y_w]
\newcommand{\qa}{\mathfrak q_1}                    % quasilinear coefficients of the slope equation
\newcommand{\qb}{\mathfrak q_2}
% cylinders (Lemma O)
\newcommand{\cyl}{\ell}                            % cylinder length unit
\newcommand{\MO}{M_O}
\newcommand{\KO}{K_O}
\newcommand{\Hcyl}{H_{\rm cyl}}
\newcommand{\XiO}{\Xi_O}
\newcommand{\fO}{f_O}
% step 4
\newcommand{\Ybar}{\bar Y}
\newcommand{\Lt}{\widetilde L}                     % linearisation about the step-4 base
\newcommand{\Trans}{\Pi}                           % transport integral
\newcommand{\resid}{\mathrm{res}}
\newcommand{\Resid}[1]{\mathrm{Res}_{#1}}
\newcommand{\boxtip}{\Box_{\rm tip}}
\newcommand{\boxfar}{\Box_{\rm far}}
\newcommand{\Eenv}{\mathcal E_{\rm env}}
\newcommand{\betab}{\beta_{\rm base}}
\newcommand{\tauI}{\tau^{\rm I}}                   % Phase I clock
% names of families of conditions: \Hh{3} = (H3), \Oh{4} = (O4), ...
\newcommand{\Hh}[1]{\textup{(H#1)}}
\newcommand{\Oh}[1]{\textup{(O#1)}}
\newcommand{\Lh}[1]{\textup{(L#1)}}
\newcommand{\Bh}[1]{\textup{(B#1)}}
\newcommand{\Kh}[1]{\textup{(K#1)}}
\newcommand{\Wh}[1]{\textup{(W#1)}}
\newcommand{\PhaseI}{Phase~I}
\newcommand{\PhaseII}{Phase~II}
% margin: threshold/bound (always)
\newcommand{\mgn}[1]{{\times}#1}

% ---- figures: placeholder macro so that drafts compile before F1-F7 exist ----
\newcommand{\figplaceholder}[1]{\fbox{\parbox{0.8\linewidth}{\centering\small[figure #1 to be supplied]}}}
 ...
%   paper_v1/supplement.tex  \documentclass[11pt]{amsart}
%                            \newif\ifsupplement \supplementtrue
%                            % ======================================================================
% preamble.tex -- shared preamble of the main paper and the supplement.
% Owner: coordinator.  Drafters add macros by request only (see DRAFTER_GUIDE.md).
%
% USE.  This file contains NO \documentclass.  The two master files are
%   paper_v1/main.tex        \documentclass[11pt]{amsart} \input{preamble} ...
%   paper_v1/supplement.tex  \documentclass[11pt]{amsart}
%                            \newif\ifsupplement \supplementtrue
%                            \input{preamble} ...
% (the main paper does nothing special: \ifsupplement defaults to false).
% Compile from paper_v1/ so that \input{numbers/numbers_paper.tex} is found.
% ======================================================================

% ---- which document are we in ----------------------------------------
\ifdefined\ifsupplement\else\newif\ifsupplement\fi

% ---- packages ---------------------------------------------------------
\usepackage[T1]{fontenc}
\usepackage[utf8]{inputenc}
\usepackage{lmodern}
\usepackage[margin=1.1in]{geometry}
\usepackage{amsmath,amssymb,amsthm,mathtools,bm}
\usepackage{booktabs,array,longtable,tabularx,multirow}
\usepackage{enumitem}
\usepackage{graphicx}
\usepackage{xcolor}
\usepackage{tikz}
\usetikzlibrary{positioning,arrows.meta,calc}
% xr-hyper must come before hyperref
\usepackage{xr-hyper}
\usepackage[colorlinks=true,linkcolor=blue!50!black,citecolor=green!35!black,urlcolor=blue!50!black]{hyperref}

\raggedbottom
\setlength{\emergencystretch}{2em}
\setlist[enumerate]{leftmargin=*,itemsep=2pt,topsep=3pt}
\setlist[itemize]{leftmargin=*,itemsep=2pt,topsep=3pt}

% ---- numbering --------------------------------------------------------
% Main paper: sections 1,2,...; theorems 2.3 (by section); equations (2.3).
% Supplement: the Parts A..K are the \section's, numbered by letter, so
% theorems are A.1, E.4, ... (do not confuse "Theorem A" of the paper, which
% has no dot, with "Theorem A.1" of Part A).
\ifsupplement
  \renewcommand{\thesection}{\Alph{section}}
\fi
\numberwithin{equation}{section}

% ---- theorem environments (numbered by section, one shared counter) ----
\theoremstyle{plain}
\newtheorem{theorem}{Theorem}[section]
\newtheorem{proposition}[theorem]{Proposition}
\newtheorem{lemma}[theorem]{Lemma}
\newtheorem{corollary}[theorem]{Corollary}
\newtheorem{maintheorem}[theorem]{Main Theorem}
% Theorems A, B, C, D of the paper: own counter, lettered, no section prefix
\newtheorem{thmlet}{Theorem}
\renewcommand{\thethmlet}{\Alph{thmlet}}
\theoremstyle{definition}
\newtheorem{definition}[theorem]{Definition}
\newtheorem{convention}[theorem]{Convention}
\newtheorem{example}[theorem]{Example}
\theoremstyle{remark}
\newtheorem{remark}[theorem]{Remark}

% A list of hypotheses labelled (H1),(H2),... : \begin{hyplist}{H} ... \end{hyplist}
% (also {O}, {L}, {B}, {K}, {W} for the families of section 5.3 of OUTLINE.md).
\newenvironment{hyplist}[1]%
  {\begin{enumerate}[label=\textup{(#1\arabic*)},leftmargin=3em,itemsep=2pt]}%
  {\end{enumerate}}

% ---- trust classes (OUTLINE 4.1, 5.1.3) -------------------------------
% Supplement: explicit tags.  Main paper: certified is the default and
% unmarked; every non-certified number is marked "(floating point)".
\ifsupplement
  \newcommand{\tcert}{\textup{[cert]}}   % arb/outward-rounded, pinned script and key
  \newcommand{\test}{\textup{[est]}}     % floating point or estimate; never in a theorem
  \newcommand{\thand}{\textup{[hand]}}   % proved by hand in the text
  \newcommand{\tlit}{\textup{[lit]}}     % cited from the literature
\else
  \newcommand{\tcert}{}
  \newcommand{\test}{\unskip\ \textup{(floating point)}}
  \newcommand{\thand}{}
  \newcommand{\tlit}{}
\fi

% ---- gaps: visible, counted, an error in a final build -----------------
% Usage: \GAP{one sentence saying what is missing}   (and an entry in ISSUES.md)
\newcounter{gapcount}
\newif\iffinalpaper   % \finalpapertrue makes every \GAP a hard error
\newcommand{\GAP}[1]{\stepcounter{gapcount}\typeout{GAP-MARKER \thegapcount}%
  \iffinalpaper\PackageError{paper}{\string\GAP\space used in a final build}{Remove all gaps.}\fi
  {\color{red}\sffamily\bfseries[GAP~\thegapcount:\ }{\color{red}\sffamily #1}{\color{red}\sffamily\bfseries]}}
\AtEndDocument{\typeout{GAP-MARKER total: \thegapcount}}

% ---- file names and JSON keys (supplement only) -------------------------
% No internal path appears in the main paper; a use there is a warning.
\newcommand{\fn}[1]{%
  \ifsupplement\else\PackageWarning{paper}{\string\fn\space used in the main paper}\fi
  \texttt{\detokenize{#1}}}
\newcommand{\key}[1]{\texttt{\detokenize{#1}}}

% ---- cross references between main paper and supplement ----------------
% Labels carry a unit prefix (section 5.3 of OUTLINE.md and DRAFTER_GUIDE.md):
%   main paper:  intro: obj: s3: s4: s5: s6: s7: s8: s9:   (kind: thm- lem- eq- tab- fig- ...)
%   supplement:  pA: pB: ... pK:
% so all labels are globally unique and the prefixes below are only a second guard.
% \Xref{label}  : reference to the OTHER document (hyperlinked via xr-hyper)
% \Lref{label}  : reference in the SAME document (plain \ref)
\ifsupplement
  \externaldocument[M-]{main}[main.pdf]
  \newcommand{\Xref}[1]{\ref{M-#1}}
\else
  \externaldocument[S-]{supplement}[supplement.pdf]
  \newcommand{\Xref}[1]{\ref{S-#1}}
\fi
\newcommand{\Lref}[1]{\ref{#1}}
% \Supp{pE}{lem-chart} -> "Lemma E.3" style text is produced by the drafter:
%   "Lemma~\Xref{pE:lem-chart} of the supplement".  \SuppPart{E} gives "Part E".
\newcommand{\SuppPart}[1]{Part~#1 of the supplement}

% ---- numbers (generated by unit U0; macros only, no literal certified numbers) ----
\IfFileExists{numbers/numbers_paper.tex}%
  {\input{numbers/numbers_paper.tex}}%
  {\PackageWarning{paper}{numbers/numbers_paper.tex not found (unit U0 not finished)}}

% ======================================================================
% NOTATION MACROS (OUTLINE 5.2 and notation.md).  One macro per symbol
% whose spelling is fixed by the notation table.  Plain letters (a, b, f, ...)
% have no macro.
% ======================================================================
% basic
\newcommand{\R}{\mathbb{R}}
\newcommand{\C}{\mathbb{C}}
\newcommand{\N}{\mathbb{N}}
\newcommand{\Z}{\mathbb{Z}}
\newcommand{\eps}{\varepsilon}
\newcommand{\ip}[2]{\langle #1,#2\rangle}
\newcommand{\abs}[1]{\lvert #1\rvert}
\newcommand{\norm}[1]{\lVert #1\rVert}
\newcommand{\dd}{\mathrm d}
\newcommand{\re}{\operatorname{Re}}
\newcommand{\im}{\operatorname{Im}}
% flow and curves
\newcommand{\arclen}{\mathfrak{s}}                 % arclength of the flow curve
\newcommand{\kss}{\kappa_{\mathfrak{ss}}}          % kappa_{ss} in the flow equation
\newcommand{\Area}{\operatorname{Area}}
\newcommand{\Len}{\mathcal L}                      % length of a curve
\newcommand{\spp}{\operatorname{sp}}                 % sup quantities sp^v_0, sp^v_1, sp^f_1 (were S^v_0, S^v_1, S^f_1)
\newcommand{\rot}{\operatorname{rot}}              % turning number
\newcommand{\Tsing}{T_*}                           % singular time
\newcommand{\Tmax}{T_{\max}}                       % maximal time of existence
\newcommand{\pstar}{p_*}                           % image point of the singularity
\newcommand{\Gz}{\Gamma_0}                         % initial curve
\newcommand{\vth}{\vartheta}                       % parameter of Gamma_0
\newcommand{\dclk}{\delta_t}                       % clock defect of step 4
\newcommand{\fall}{\varpi}                         % Fourier rates varpi_n = n^4-n^2
% rescaling
\newcommand{\lamone}{\lambda_1}
\newcommand{\rate}{\varsigma}                      % 1/4 - a
\newcommand{\xia}{\xi_a}                           % (1-4\bar a)^{-1}
% shrinker and step 1
\newcommand{\Sone}{S_1}
\newcommand{\Gsh}{\Gamma_{\!*}}                    % profile curve of S_1
\newcommand{\cSet}{\mathcal S}                     % certified set of shrinkers
\newcommand{\sigtip}{\sigma_{\rm tip}}
\newcommand{\Xo}{X_0}                              % 2.2490776515
\newcommand{\Dtip}{D}
\newcommand{\ktip}{k_0}
\newcommand{\shootM}{M}                            % shooting function
\newcommand{\Ypar}{\mathcal Y}                     % step-1 parallelepiped
\newcommand{\axis}{\Sigma}                         % symmetry axis of S_1
\newcommand{\sigc}{\sigma_{\rm c}}
\newcommand{\sigrev}{\sigma^{\rm rev}}
% linear theory
\newcommand{\Hrho}{\mathcal H_\rho}
\newcommand{\muslow}{\mu_3}
\newcommand{\rre}{\mathfrak r_{\rm re}}
\newcommand{\rim}{\mathfrak r_{\rm im}}
\newcommand{\Evev}{\mathrm{Ev}_{\rm even}}
\newcommand{\Evod}{\mathrm{Ev}_{\rm odd}}
\newcommand{\gvec}{\mathbf g}                      % translation vector in step 2
\newcommand{\cons}[1]{\mathfrak f_{#1}}            % constraint functions
\newcommand{\phirho}{\varphi_\rho}
% energies
\newcommand{\Lam}{\Lambda}
\newcommand{\En}{E}
\newcommand{\EnI}{E_I}
\newcommand{\muI}{\mu_I}                           % 0.08
\newcommand{\epsE}{\varepsilon_E}                  % 1e-5
\newcommand{\epsI}{\varepsilon_I}                  % 3e-5
\newcommand{\wn}[1]{\mathsf w_{#1}}                % ||w~^{(j)}||_rho
\newcommand{\vI}{v_I}                              % zeta w'
\newcommand{\muerr}{\mu_{\rm err}}                 % step-4 error-energy weight 0.036
% far field
\newcommand{\pU}{p_U}                              % 1.0246
\newcommand{\Uone}{U_1}
\newcommand{\sm}{\operatorname{sm}}                % smoothstep t^3(10-15t+6t^2)
\newcommand{\rc}{r_c}                              % centre radius (Lemma U)
\newcommand{\rchi}{r_\chi}                         % 253
\newcommand{\Iw}{I_w}                              % hand-off windows [y_w,1.5y_w]
\newcommand{\qa}{\mathfrak q_1}                    % quasilinear coefficients of the slope equation
\newcommand{\qb}{\mathfrak q_2}
% cylinders (Lemma O)
\newcommand{\cyl}{\ell}                            % cylinder length unit
\newcommand{\MO}{M_O}
\newcommand{\KO}{K_O}
\newcommand{\Hcyl}{H_{\rm cyl}}
\newcommand{\XiO}{\Xi_O}
\newcommand{\fO}{f_O}
% step 4
\newcommand{\Ybar}{\bar Y}
\newcommand{\Lt}{\widetilde L}                     % linearisation about the step-4 base
\newcommand{\Trans}{\Pi}                           % transport integral
\newcommand{\resid}{\mathrm{res}}
\newcommand{\Resid}[1]{\mathrm{Res}_{#1}}
\newcommand{\boxtip}{\Box_{\rm tip}}
\newcommand{\boxfar}{\Box_{\rm far}}
\newcommand{\Eenv}{\mathcal E_{\rm env}}
\newcommand{\betab}{\beta_{\rm base}}
\newcommand{\tauI}{\tau^{\rm I}}                   % Phase I clock
% names of families of conditions: \Hh{3} = (H3), \Oh{4} = (O4), ...
\newcommand{\Hh}[1]{\textup{(H#1)}}
\newcommand{\Oh}[1]{\textup{(O#1)}}
\newcommand{\Lh}[1]{\textup{(L#1)}}
\newcommand{\Bh}[1]{\textup{(B#1)}}
\newcommand{\Kh}[1]{\textup{(K#1)}}
\newcommand{\Wh}[1]{\textup{(W#1)}}
\newcommand{\PhaseI}{Phase~I}
\newcommand{\PhaseII}{Phase~II}
% margin: threshold/bound (always)
\newcommand{\mgn}[1]{{\times}#1}

% ---- figures: placeholder macro so that drafts compile before F1-F7 exist ----
\newcommand{\figplaceholder}[1]{\fbox{\parbox{0.8\linewidth}{\centering\small[figure #1 to be supplied]}}}
 ...
% (the main paper does nothing special: \ifsupplement defaults to false).
% Compile from paper_v1/ so that % numbers_paper.tex -- GENERATED from ledger.json by build_ledger.py; do not edit.
% Rounding: upper bounds up, lower bounds and margins down at the printed digit; margin = threshold / bound.
% Use as $\times\NName$ for margins.  kind: cert certified, est estimate, hand hypothesis constant, doc only a document quotes it.
% NUEhatZero [cert, upper] U_cert_v35.json :: supersolution.EA.0
\newcommand{\NUEhatZero}{\ensuremath{1.2018\cdot10^{-6}}}
% NUEhatOne [cert, upper] U_cert_v35.json :: supersolution.EA.1
\newcommand{\NUEhatOne}{\ensuremath{1.0441\cdot10^{-4}}}
% NUEhatTwo [cert, upper] U_cert_v35.json :: supersolution.EA.2
\newcommand{\NUEhatTwo}{\ensuremath{0.056694}}
% NUBhatZero [cert, upper] U_cert_v35.json :: supersolution.Bhat.0
\newcommand{\NUBhatZero}{\ensuremath{4.8735\cdot10^{-6}}}
% NUBhatOne [cert, upper] U_cert_v35.json :: supersolution.Bhat.1
\newcommand{\NUBhatOne}{\ensuremath{5.5933\cdot10^{-4}}}
% NUBhatTwo [cert, upper] U_cert_v35.json :: supersolution.Bhat.2
\newcommand{\NUBhatTwo}{\ensuremath{0.36675}}
% NUSupersolutionCZeroZero [cert, upper] U_cert_v35.json :: supersolution.C0.0
\newcommand{\NUSupersolutionCZeroZero}{\ensuremath{1.1340\cdot10^{-6}}}
% NUSupersolutionCZeroOne [cert, upper] U_cert_v35.json :: supersolution.C0.1
\newcommand{\NUSupersolutionCZeroOne}{\ensuremath{1.6623\cdot10^{-4}}}
% NUSupersolutionCOneZero [cert, upper] U_cert_v35.json :: supersolution.C1.0
\newcommand{\NUSupersolutionCOneZero}{\ensuremath{1.5241\cdot10^{-9}}}
% NUSupersolutionCOneOne [cert, upper] U_cert_v35.json :: supersolution.C1.1
\newcommand{\NUSupersolutionCOneOne}{\ensuremath{5.5255\cdot10^{-8}}}
% NUSupersolutionNTwoZero [cert, upper] U_cert_v35.json :: supersolution.N2.0
\newcommand{\NUSupersolutionNTwoZero}{\ensuremath{0.071472}}
% NUSupersolutionNTwoOne [cert, upper] U_cert_v35.json :: supersolution.N2.1
\newcommand{\NUSupersolutionNTwoOne}{\ensuremath{0.33012}}
% NUSupersolutionNTwoTwo [cert, upper] U_cert_v35.json :: supersolution.N2.2
\newcommand{\NUSupersolutionNTwoTwo}{\ensuremath{5.0406}}
% NUSupersolutionNTwoOverFourNTwoRefZero [cert, upper] U_cert_v35.json :: supersolution.N2_over_4N2ref.0
\newcommand{\NUSupersolutionNTwoOverFourNTwoRefZero}{\ensuremath{0.18}}
% NUSupersolutionNTwoOverFourNTwoRefOne [cert, upper] U_cert_v35.json :: supersolution.N2_over_4N2ref.1
\newcommand{\NUSupersolutionNTwoOverFourNTwoRefOne}{\ensuremath{0.42}}
% NUSupersolutionNTwoOverFourNTwoRefTwo [cert, upper] U_cert_v35.json :: supersolution.N2_over_4N2ref.2
\newcommand{\NUSupersolutionNTwoOverFourNTwoRefTwo}{\ensuremath{0.26}}
% NUSupersolutionMOneZHypCheckZonesIIISOneTwoBhat [cert, upper] U_cert_v35.json :: supersolution.M1z_hyp_check.zones_I_II.S1_2Bhat
\newcommand{\NUSupersolutionMOneZHypCheckZonesIIISOneTwoBhat}{\ensuremath{0.28585}}
% NUSupersolutionMOneZHypCheckZoneIIIBOnePlusPUOne [cert, upper] U_cert_v35.json :: supersolution.M1z_hyp_check.zone_III.b1_plus_pU1
\newcommand{\NUSupersolutionMOneZHypCheckZoneIIIBOnePlusPUOne}{\ensuremath{0.50786}}
% NUMarginLateMOne [cert, margin] U_cert_v35.json :: margins.late_M1
\newcommand{\NUMarginLateMOne}{\ensuremath{2.353}}
% NUMarginLateMZero [cert, margin] U_cert_v35.json :: margins.late_M0
\newcommand{\NUMarginLateMZero}{\ensuremath{4.892}}
% NUMarginsLateMZeroZ [cert, margin] U_cert_v35.json :: margins.late_M0z
\newcommand{\NUMarginsLateMZeroZ}{\ensuremath{12.99}}
% NUMarginPresentMOne [cert, margin] U_cert_v35.json :: margins.present_M1_min.value
\newcommand{\NUMarginPresentMOne}{\ensuremath{2.770}}
% NUMarginPresentMZero [cert, margin] U_cert_v35.json :: margins.present_M0_min.value
\newcommand{\NUMarginPresentMZero}{\ensuremath{2.225}}
% NUMarginsPresentMZeroZMinValue [cert, margin] U_cert_v35.json :: margins.present_M0z_min.value
\newcommand{\NUMarginsPresentMZeroZMinValue}{\ensuremath{14.37}}
% NULoopGainZeroTotal [cert, upper] U_cert_v35.json :: loop_gain.0.total
\newcommand{\NULoopGainZeroTotal}{\ensuremath{1.23\cdot10^{-7}}}
% NULoopGainOneTotal [cert, upper] U_cert_v35.json :: loop_gain.1.total
\newcommand{\NULoopGainOneTotal}{\ensuremath{0.037169}}
% NULoopGainTwoNeighbour [cert, upper] U_cert_v35.json :: loop_gain.2.neighbour
\newcommand{\NULoopGainTwoNeighbour}{\ensuremath{0.037018}}
% NULoopGainTwoThroughA [cert, upper] U_cert_v35.json :: loop_gain.2.through_A
\newcommand{\NULoopGainTwoThroughA}{\ensuremath{0.0020172}}
% NUfineBhatZero [cert, upper] U_cert_v35.json :: ufine.Bhat.0
\newcommand{\NUfineBhatZero}{\ensuremath{4.8723\cdot10^{-6}}}
% NUfineBhatOne [cert, upper] U_cert_v35.json :: ufine.Bhat.1
\newcommand{\NUfineBhatOne}{\ensuremath{2.6266\cdot10^{-4}}}
% NUfineBhatTwo [cert, upper] U_cert_v35.json :: ufine.Bhat.2
\newcommand{\NUfineBhatTwo}{\ensuremath{0.025842}}
% NUcdeltaU [cert, upper] U_cert_v35.json :: launch.c_dU
\newcommand{\NUcdeltaU}{\ensuremath{1.7187\cdot10^{-8}}}
% NULaunchWmin [cert, lower] U_cert_v35.json :: launch.wmin
\newcommand{\NULaunchWmin}{\ensuremath{0.73492}}
% NULaunchCAD [cert, upper] U_cert_v35.json :: launch.CA_D
\newcommand{\NULaunchCAD}{\ensuremath{2.2429}}
% NULaunchQD [cert, upper] U_cert_v35.json :: launch.qD
\newcommand{\NULaunchQD}{\ensuremath{0.018509}}
% NULaunchDeltaInZero [cert, upper] U_cert_v35.json :: launch.Delta_in.0
\newcommand{\NULaunchDeltaInZero}{\ensuremath{3.131\cdot10^{-12}}}
% NULaunchDeltaInOne [cert, upper] U_cert_v35.json :: launch.Delta_in.1
\newcommand{\NULaunchDeltaInOne}{\ensuremath{3.287\cdot10^{-12}}}
% NULaunchDeltaInTwo [cert, upper] U_cert_v35.json :: launch.Delta_in.2
\newcommand{\NULaunchDeltaInTwo}{\ensuremath{3.451\cdot10^{-12}}}
% NULaunchPhiZero [cert, upper] U_cert_v35.json :: launch.Phi.0
\newcommand{\NULaunchPhiZero}{\ensuremath{6.434\cdot10^{-9}}}
% NULaunchPhiOne [cert, upper] U_cert_v35.json :: launch.Phi.1
\newcommand{\NULaunchPhiOne}{\ensuremath{1.4217\cdot10^{-8}}}
% NULaunchPhiTwo [cert, upper] U_cert_v35.json :: launch.Phi.2
\newcommand{\NULaunchPhiTwo}{\ensuremath{3.8404\cdot10^{-8}}}
% NUGammaZero [cert, upper] U_cert_v35.json :: Gamma.0
\newcommand{\NUGammaZero}{\ensuremath{4.0000358}}
% NUGammaOne [cert, upper] U_cert_v35.json :: Gamma.1
\newcommand{\NUGammaOne}{\ensuremath{16.000322}}
% NUGammaTwo [cert, upper] U_cert_v35.json :: Gamma.2
\newcommand{\NUGammaTwo}{\ensuremath{64.00218}}
% NUKappa [cert, upper] U_cert_v35.json :: kappa
\newcommand{\NUKappa}{\ensuremath{8.0000358}}
% NUKappaTwo [cert, upper] U_cert_v35.json :: kappa2
\newcommand{\NUKappaTwo}{\ensuremath{64.00061}}
% NUGhatZeroZero [cert, upper] U_cert_v35.json :: Ghat.0.0
\newcommand{\NUGhatZeroZero}{\ensuremath{51322.6}}
% NUGhatOneOne [cert, upper] U_cert_v35.json :: Ghat.1.1
\newcommand{\NUGhatOneOne}{\ensuremath{33962.9}}
% NUGhatTwoTwo [cert, upper] U_cert_v35.json :: Ghat.2.2
\newcommand{\NUGhatTwoTwo}{\ensuremath{348341.1}}
% NUGhatOneZero [cert, upper] U_cert_v35.json :: Ghat.1.0
\newcommand{\NUGhatOneZero}{\ensuremath{4.938}}
% NUGhatTwoZero [cert, upper] U_cert_v35.json :: Ghat.2.0
\newcommand{\NUGhatTwoZero}{\ensuremath{27.984}}
% NUGhatTwoOne [cert, upper] U_cert_v35.json :: Ghat.2.1
\newcommand{\NUGhatTwoOne}{\ensuremath{458.53}}
% NUKZeroZero [cert, upper] U_cert_v35.json :: K.0.0
\newcommand{\NUKZeroZero}{\ensuremath{8034.21}}
% NUKZeroOne [cert, upper] U_cert_v35.json :: K.0.1
\newcommand{\NUKZeroOne}{\ensuremath{724.89}}
% NUKOneZero [cert, upper] U_cert_v35.json :: K.1.0
\newcommand{\NUKOneZero}{\ensuremath{0.08152}}
% NUKOneOne [cert, upper] U_cert_v35.json :: K.1.1
\newcommand{\NUKOneOne}{\ensuremath{6336.94}}
% NUKOneTwo [cert, upper] U_cert_v35.json :: K.1.2
\newcommand{\NUKOneTwo}{\ensuremath{545.59}}
% NUKOneThree [cert, upper] U_cert_v35.json :: K.1.3
\newcommand{\NUKOneThree}{\ensuremath{0.58700}}
% NUKTwoZero [cert, upper] U_cert_v35.json :: K.2.0
\newcommand{\NUKTwoZero}{\ensuremath{0.87186}}
% NUKTwoOne [cert, upper] U_cert_v35.json :: K.2.1
\newcommand{\NUKTwoOne}{\ensuremath{528.62}}
% NUKTwoTwo [cert, upper] U_cert_v35.json :: K.2.2
\newcommand{\NUKTwoTwo}{\ensuremath{23425.98}}
% NUKTwoThree [cert, upper] U_cert_v35.json :: K.2.3
\newcommand{\NUKTwoThree}{\ensuremath{555.62}}
% NUFootwOne [cert, upper] U_cert_v35.json :: footw.1
\newcommand{\NUFootwOne}{\ensuremath{2.6361}}
% NUFootwTwo [cert, upper] U_cert_v35.json :: footw.2
\newcommand{\NUFootwTwo}{\ensuremath{2663.3}}
% NUCThreeThree [cert, upper] U_cert_v35.json :: c33
\newcommand{\NUCThreeThree}{\ensuremath{108.02}}
% NUCTwoTwoThree [cert, upper] U_cert_v35.json :: c223
\newcommand{\NUCTwoTwoThree}{\ensuremath{7.3284}}
% NUCEta [cert, upper] U_cert_v35.json :: C_eta
\newcommand{\NUCEta}{\ensuremath{13.376}}
% NULMax [cert, upper] U_cert_v35.json :: L_max
\newcommand{\NULMax}{\ensuremath{3.5352}}
% NURRatio [cert, upper] U_cert_v35.json :: R_ratio
\newcommand{\NURRatio}{\ensuremath{11.593}}
% NUSweepsSigmaInCoverSup [cert, upper] U_cert_v35.json :: sweeps.Sigma_in.cover_sup
\newcommand{\NUSweepsSigmaInCoverSup}{\ensuremath{5.0892}}
% NUSweepsSigmaECoverSup [cert, upper] U_cert_v35.json :: sweeps.Sigma_E.cover_sup
\newcommand{\NUSweepsSigmaECoverSup}{\ensuremath{0.30846}}
% NUSweepsSXCoverSup [cert, upper] U_cert_v35.json :: sweeps.S_X.cover_sup
\newcommand{\NUSweepsSXCoverSup}{\ensuremath{1.5243}}
% NUSweepsSigmaXCoverSup [cert, upper] U_cert_v35.json :: sweeps.Sigma_X.cover_sup
\newcommand{\NUSweepsSigmaXCoverSup}{\ensuremath{3.0081}}
% NUSweepsSigmaRsTwoCoverSup [cert, upper] U_cert_v35.json :: sweeps.Sigma_rs.2.cover_sup
\newcommand{\NUSweepsSigmaRsTwoCoverSup}{\ensuremath{4.3350\cdot10^{-3}}}
% NUDataTermsLateDZero [cert, upper] U_cert_v35.json :: data_terms_late.D.0
\newcommand{\NUDataTermsLateDZero}{\ensuremath{1.7330\cdot10^{-6}}}
% NUDataTermsLateDOne [cert, upper] U_cert_v35.json :: data_terms_late.D.1
\newcommand{\NUDataTermsLateDOne}{\ensuremath{2.5353\cdot10^{-4}}}
% NUDataTermsLateSlopeShareCOneSigmaXOverDkZero [cert, upper] U_cert_v35.json :: data_terms_late.slope_share_C1SigmaX_over_Dk.0
\newcommand{\NUDataTermsLateSlopeShareCOneSigmaXOverDkZero}{\ensuremath{2.7\cdot10^{-3}}}
% NUDataTermsLateSlopeShareCOneSigmaXOverDkOne [cert, upper] U_cert_v35.json :: data_terms_late.slope_share_C1SigmaX_over_Dk.1
\newcommand{\NUDataTermsLateSlopeShareCOneSigmaXOverDkOne}{\ensuremath{6.6\cdot10^{-4}}}
% NURLtTwoMarginZero [cert, margin] U_cert_v35.json :: R_lt2_margin.0
\newcommand{\NURLtTwoMarginZero}{\ensuremath{4.432}}
% NURLtTwoMarginOne [cert, margin] U_cert_v35.json :: R_lt2_margin.1
\newcommand{\NURLtTwoMarginOne}{\ensuremath{6.697}}
% NUThetaTwoEight [cert, upper] U_cert_v35.json :: theta28
\newcommand{\NUThetaTwoEight}{\ensuremath{0.37011}}
% NUDFiveSix [cert, upper] U_cert_v35.json :: D56
\newcommand{\NUDFiveSix}{\ensuremath{1.0368}}
% NUDTwoEight [cert, upper] U_cert_v35.json :: D28
\newcommand{\NUDTwoEight}{\ensuremath{1.5876}}
% NUSThreeThreeFootOverBhatTwo [cert, upper] U_cert_v35.json :: S33_foot_over_Bhat2
\newcommand{\NUSThreeThreeFootOverBhatTwo}{\ensuremath{1.69\cdot10^{-3}}}
% NSAHJWKappaJ [cert, upper] A_interface.json :: HJW.kappa_J
\newcommand{\NSAHJWKappaJ}{\ensuremath{0.042351}}
% NSAGpZero [cert, lower] A_interface.json :: res.tight.gp.0
\newcommand{\NSAGpZero}{\ensuremath{0.24292}}
% NSAGpOne [cert, lower] A_interface.json :: res.tight.gp.1
\newcommand{\NSAGpOne}{\ensuremath{0.24106}}
% NSAGpTwo [cert, lower] A_interface.json :: res.tight.gp.2
\newcommand{\NSAGpTwo}{\ensuremath{0.11961}}
% NSAAminC [cert, lower] A_interface.json :: res.tight.amin_c
\newcommand{\NSAAminC}{\ensuremath{0.97623}}
% NSARhoJ [cert, upper] A_interface.json :: res.tight.rhoJ
\newcommand{\NSARhoJ}{\ensuremath{5.4995\cdot10^{-7}}}
% NSAKappaOTwoTwo [cert, upper] A_interface.json :: res.tight.kappa_o.2.2
\newcommand{\NSAKappaOTwoTwo}{\ensuremath{1.2416\cdot10^{-5}}}
% NSADsrcZero [cert, upper] A_interface.json :: res.tight.Dsrc.0
\newcommand{\NSADsrcZero}{\ensuremath{1.8092\cdot10^{-7}}}
% NSADsrcOne [cert, upper] A_interface.json :: res.tight.Dsrc.1
\newcommand{\NSADsrcOne}{\ensuremath{9.8276\cdot10^{-5}}}
% NSABZeroDWindowSlopeZero [cert, upper] A_interface.json :: res.tight.B0.D_window_slope.0
\newcommand{\NSABZeroDWindowSlopeZero}{\ensuremath{8.7399\cdot10^{-10}}}
% NSABZeroDWindowSlopeOne [cert, upper] A_interface.json :: res.tight.B0.D_window_slope.1
\newcommand{\NSABZeroDWindowSlopeOne}{\ensuremath{5.2599\cdot10^{-8}}}
% NSAFTwoZero [cert, upper] A_interface.json :: res.tight.F2.0
\newcommand{\NSAFTwoZero}{\ensuremath{5.6506\cdot10^{-13}}}
% NSAFTwoOne [cert, upper] A_interface.json :: res.tight.F2.1
\newcommand{\NSAFTwoOne}{\ensuremath{5.9320\cdot10^{-13}}}
% NSACoupOne [cert, upper] A_interface.json :: res.tight.coup.1
\newcommand{\NSACoupOne}{\ensuremath{1.1568\cdot10^{-12}}}
% NSACOneZeroZero [cert, upper] A_interface.json :: res.tight.C1_0.0
\newcommand{\NSACOneZeroZero}{\ensuremath{7.2060\cdot10^{-10}}}
% NSACOneZeroOne [cert, upper] A_interface.json :: res.tight.C1_0.1
\newcommand{\NSACOneZeroOne}{\ensuremath{4.3143\cdot10^{-8}}}
% NSABZeroSupCoreZero [cert, upper] A_interface.json :: res.tight.B0.sup_core.0
\newcommand{\NSABZeroSupCoreZero}{\ensuremath{4.23\cdot10^{-5}}}
% NSACoreFluxMarginTwoFLUXZeroOverKZero [cert, margin] A_interface.json :: res.tight.core.flux_margin.2.FLUX0_over_K0
\newcommand{\NSACoreFluxMarginTwoFLUXZeroOverKZero}{\ensuremath{3.396}}
% NSACoreFluxMarginTwoFLUXTwoOverKTwo [cert, margin] A_interface.json :: res.tight.core.flux_margin.2.FLUX2_over_K2
\newcommand{\NSACoreFluxMarginTwoFLUXTwoOverKTwo}{\ensuremath{6.012}}
% NSAMarginFluxZero [cert, margin] asm_design_N2ref4.json :: res.tight.flux.beta.2.margin_FLUX0
\newcommand{\NSAMarginFluxZero}{\ensuremath{3.365}}
% NSAMarginFluxTwo [cert, margin] asm_design_N2ref4.json :: res.tight.flux.beta.2.margin_FLUX2
\newcommand{\NSAMarginFluxTwo}{\ensuremath{5.955}}
% NSAfourSZero [cert, upper] asm_design_N2ref4.json :: res.tight.S0
\newcommand{\NSAfourSZero}{\ensuremath{7.1878\cdot10^{-3}}}
% NSAfourSOne [cert, upper] asm_design_N2ref4.json :: res.tight.S1
\newcommand{\NSAfourSOne}{\ensuremath{0.10430}}
% NSAMarginCoreSZeroV [cert, margin] asm_design_N2ref4.json :: res.tight.margins.S0v_vs_0.0153
\newcommand{\NSAMarginCoreSZeroV}{\ensuremath{2.116}}
% NSAMarginCoreSOneV [cert, margin] asm_design_N2ref4.json :: res.tight.margins.S1v_vs_0.6
\newcommand{\NSAMarginCoreSOneV}{\ensuremath{5.748}}
% NSAMarginCoreMOne [cert, margin] asm_design_N2ref4.json :: res.tight.margins.M1core
\newcommand{\NSAMarginCoreMOne}{\ensuremath{4.409}}
% NSAMarginCoreMZero [cert, margin] asm_design_N2ref4.json :: res.tight.margins.M0core
\newcommand{\NSAMarginCoreMZero}{\ensuremath{5.450}}
% NSAMarginGhat [cert, margin] asm_design_N2ref4.json :: res.tight.margins.Ghat22
\newcommand{\NSAMarginGhat}{\ensuremath{1.917}}
% NSAMarginRampMOne [cert, margin] asm_design_N2ref4.json :: res.tight.margins.ramp M1
\newcommand{\NSAMarginRampMOne}{\ensuremath{1.184}}
% NSAfourInfoZeroDWindowConst [cert, upper] asm_design_N2ref4.json :: res.tight.info.0.D_window_const
\newcommand{\NSAfourInfoZeroDWindowConst}{\ensuremath{1.7732\cdot10^{-7}}}
% NSAfourInfoOneDWindowConst [cert, upper] asm_design_N2ref4.json :: res.tight.info.1.D_window_const
\newcommand{\NSAfourInfoOneDWindowConst}{\ensuremath{9.8058\cdot10^{-5}}}
% NZDeltaRZero [cert, upper] zone3_interface.json :: DeltaR.0
\newcommand{\NZDeltaRZero}{\ensuremath{4.4986\cdot10^{-8}}}
% NZDeltaROne [cert, upper] zone3_interface.json :: DeltaR.1
\newcommand{\NZDeltaROne}{\ensuremath{6.7972\cdot10^{-6}}}
% NZDeltaRTwo [cert, upper] zone3_interface.json :: DeltaR.2
\newcommand{\NZDeltaRTwo}{\ensuremath{0.011480}}
% NZDeltaRBZero [cert, upper] zone3_interface.json :: DeltaR_b.0
\newcommand{\NZDeltaRBZero}{\ensuremath{2.66\cdot10^{-12}}}
% NZDeltaRBOne [cert, upper] zone3_interface.json :: DeltaR_b.1
\newcommand{\NZDeltaRBOne}{\ensuremath{3.26\cdot10^{-9}}}
% NZDeltaRBTwo [cert, upper] zone3_interface.json :: DeltaR_b.2
\newcommand{\NZDeltaRBTwo}{\ensuremath{4.35\cdot10^{-8}}}
% NZDeltaRDetailNSqZero [cert, upper] zone3_interface.json :: DeltaR_detail.N_sq.0
\newcommand{\NZDeltaRDetailNSqZero}{\ensuremath{8.28\cdot10^{-5}}}
% NZDeltaRDetailNSqOne [cert, upper] zone3_interface.json :: DeltaR_detail.N_sq.1
\newcommand{\NZDeltaRDetailNSqOne}{\ensuremath{0.10149}}
% NZDeltaRDetailNSqTwo [cert, upper] zone3_interface.json :: DeltaR_detail.N_sq.2
\newcommand{\NZDeltaRDetailNSqTwo}{\ensuremath{1.3567}}
% NZDeltaRDetailNSqThree [cert, upper] zone3_interface.json :: DeltaR_detail.N_sq.3
\newcommand{\NZDeltaRDetailNSqThree}{\ensuremath{25350}}
% NZDeltaRDetailPartsTwoPairsThree [cert, upper] zone3_interface.json :: DeltaR_detail.parts.2.pairs.3
\newcommand{\NZDeltaRDetailPartsTwoPairsThree}{\ensuremath{0.011298}}
% NZDeltaRDetailPartsTwoLFourFoot [cert, upper] zone3_interface.json :: DeltaR_detail.parts.2.L4_foot
\newcommand{\NZDeltaRDetailPartsTwoLFourFoot}{\ensuremath{1.6\cdot10^{-4}}}
% NZDeltaRDetailPartsTwoCubic [cert, upper] zone3_interface.json :: DeltaR_detail.parts.2.cubic
\newcommand{\NZDeltaRDetailPartsTwoCubic}{\ensuremath{9.3\cdot10^{-6}}}
% NZDeltaRDetailPartsTwoSThreeThreeFoot [cert, upper] zone3_interface.json :: DeltaR_detail.parts.2.S33_foot
\newcommand{\NZDeltaRDetailPartsTwoSThreeThreeFoot}{\ensuremath{3.7\cdot10^{-7}}}
% NZNIIIZero [cert, upper] zone3_interface.json :: n_III.0
\newcommand{\NZNIIIZero}{\ensuremath{0.071445}}
% NZNIIIOne [cert, upper] zone3_interface.json :: n_III.1
\newcommand{\NZNIIIOne}{\ensuremath{0.32707}}
% NZNIIITwo [cert, upper] zone3_interface.json :: n_III.2
\newcommand{\NZNIIITwo}{\ensuremath{3.0445}}
% NZTauIII [cert, upper] zone3_interface.json :: tau_III
\newcommand{\NZTauIII}{\ensuremath{4.0708}}
% NZoneIIIBoxZero [cert, upper] zone3_interface.json :: box_III_380.b0
\newcommand{\NZoneIIIBoxZero}{\ensuremath{0.10032}}
% NZoneIIIBoxOne [cert, upper] zone3_interface.json :: box_III_380.b1
\newcommand{\NZoneIIIBoxOne}{\ensuremath{0.41974}}
% NZoneIIIBoxZ [cert, upper] zone3_interface.json :: box_III_380.bz
\newcommand{\NZoneIIIBoxZ}{\ensuremath{0.014291}}
% NZBoxIIIThreeEightZeroMarginsBZero [cert, margin] zone3_interface.json :: box_III_380.margins.b0
\newcommand{\NZBoxIIIThreeEightZeroMarginsBZero}{\ensuremath{1.495}}
% NZoneIIIMarginThinnest [cert, margin] zone3_interface.json :: box_III_380.margins.b1
\newcommand{\NZoneIIIMarginThinnest}{\ensuremath{1.238}}
% NZBoxIIIThreeEightZeroMarginsBz [cert, margin] zone3_interface.json :: box_III_380.margins.bz
\newcommand{\NZBoxIIIThreeEightZeroMarginsBz}{\ensuremath{11.19}}
% NZBoxIIISevenZeroZeroBZero [cert, upper] zone3_interface.json :: box_III_700.b0
\newcommand{\NZBoxIIISevenZeroZeroBZero}{\ensuremath{0.22757}}
% NZBoxIIISevenZeroZeroBOne [cert, upper] zone3_interface.json :: box_III_700.b1
\newcommand{\NZBoxIIISevenZeroZeroBOne}{\ensuremath{0.54234}}
% NZUfineKTwoAtBhatRef [cert, upper] zone3_interface.json :: ufine.K2_at_Bhat_ref
\newcommand{\NZUfineKTwoAtBhatRef}{\ensuremath{0.046956}}
% NZUfineKThreeStripAtBhatRef [cert, upper] zone3_interface.json :: ufine.K3_strip_at_Bhat_ref
\newcommand{\NZUfineKThreeStripAtBhatRef}{\ensuremath{0.054140}}
% NRecoveryMinMargin [cert, margin] zone3_interface.json :: cone_family.H6.at_Bhat_ref.min_margin
\newcommand{\NRecoveryMinMargin}{\ensuremath{1.860}}
% NZConeFamilyHSixAtBhatRefConeSlopeMargin [cert, margin] zone3_interface.json :: cone_family.H6.at_Bhat_ref.cone.slope.margin
\newcommand{\NZConeFamilyHSixAtBhatRefConeSlopeMargin}{\ensuremath{2.46}}
% NZConeFamilyHSixAtBhatRefConeKTwoMargin [cert, margin] zone3_interface.json :: cone_family.H6.at_Bhat_ref.cone.K2.margin
\newcommand{\NZConeFamilyHSixAtBhatRefConeKTwoMargin}{\ensuremath{2.54}}
% NZConeFamilyHSixAtBhatRefConeKThreeMargin [cert, margin] zone3_interface.json :: cone_family.H6.at_Bhat_ref.cone.K3.margin
\newcommand{\NZConeFamilyHSixAtBhatRefConeKThreeMargin}{\ensuremath{2.51}}
% NConeTstarMargin [cert, margin] zone3_interface.json :: acceptance.tstar_margin
\newcommand{\NConeTstarMargin}{\ensuremath{2.30}}
% NUoneTFiveYaddLevelOneTwo [cert, upper] u1_ghat.json :: t5_Yadd.level1.2
\newcommand{\NUoneTFiveYaddLevelOneTwo}{\ensuremath{0.0044}}
% NUoneTFiveYaddLevelTwoTwo [cert, upper] u1_ghat.json :: t5_Yadd.level2.2
\newcommand{\NUoneTFiveYaddLevelTwoTwo}{\ensuremath{0.1081}}
% NUoneTFiveYaddLevelTwoThree [cert, upper] u1_ghat.json :: t5_Yadd.level2.3
\newcommand{\NUoneTFiveYaddLevelTwoThree}{\ensuremath{0.374}}
% NUoneWeightThetaGammaZero [cert, upper] u1_ghat.json :: weight.Theta.Gamma.0
\newcommand{\NUoneWeightThetaGammaZero}{\ensuremath{4.00000000000000000000054}}
% NRKZeroZero [cert, upper] ramp_U.json :: K.0.0
\newcommand{\NRKZeroZero}{\ensuremath{13436.03}}
% NRKOneOne [cert, upper] ramp_U.json :: K.1.1
\newcommand{\NRKOneOne}{\ensuremath{10526.19}}
% NRKTwoTwo [cert, upper] ramp_U.json :: K.2.2
\newcommand{\NRKTwoTwo}{\ensuremath{39482.30}}
% NUoneCasesPrimaryYNextZero [cert, upper] u1_ghat.json :: cases.primary.Y_next.0
\newcommand{\NUoneCasesPrimaryYNextZero}{\ensuremath{457.092}}
% NUoneCasesPrimaryYNextOne [cert, upper] u1_ghat.json :: cases.primary.Y_next.1
\newcommand{\NUoneCasesPrimaryYNextOne}{\ensuremath{356.940}}
% NUoneCasesPrimaryYNextTwo [cert, upper] u1_ghat.json :: cases.primary.Y_next.2
\newcommand{\NUoneCasesPrimaryYNextTwo}{\ensuremath{1320.220}}
% NSfourResultsMNineHTwoTightenedFUOneRowsZeroBound [cert, upper] FINAL_manifest_v33.json :: results.M9_H2_tightened_fU1.rows.0.bound
\newcommand{\NSfourResultsMNineHTwoTightenedFUOneRowsZeroBound}{\ensuremath{3.658\cdot10^{-10}}}
% NSfourResultsMNineHTwoTightenedFUOneRowsZeroMargin [cert, margin] FINAL_manifest_v33.json :: results.M9_H2_tightened_fU1.rows.0.margin
\newcommand{\NSfourResultsMNineHTwoTightenedFUOneRowsZeroMargin}{\ensuremath{27.34}}
% NSfourResultsMNineHTwoTightenedFUOneRowsOneBound [cert, upper] FINAL_manifest_v33.json :: results.M9_H2_tightened_fU1.rows.1.bound
\newcommand{\NSfourResultsMNineHTwoTightenedFUOneRowsOneBound}{\ensuremath{1.777\cdot10^{-9}}}
% NSfourResultsMNineHTwoTightenedFUOneRowsOneMargin [cert, margin] FINAL_manifest_v33.json :: results.M9_H2_tightened_fU1.rows.1.margin
\newcommand{\NSfourResultsMNineHTwoTightenedFUOneRowsOneMargin}{\ensuremath{562.8}}
% NSfourResultsMNineHTwoTightenedFUOneRowsTwoBound [cert, upper] FINAL_manifest_v33.json :: results.M9_H2_tightened_fU1.rows.2.bound
\newcommand{\NSfourResultsMNineHTwoTightenedFUOneRowsTwoBound}{\ensuremath{1.665\cdot10^{-6}}}
% NSfourResultsMNineHTwoTightenedFUOneRowsTwoMargin [cert, margin] FINAL_manifest_v33.json :: results.M9_H2_tightened_fU1.rows.2.margin
\newcommand{\NSfourResultsMNineHTwoTightenedFUOneRowsTwoMargin}{\ensuremath{600.9}}
% NHOneAlphaThr [hand, exact] hand :: 
\newcommand{\NHOneAlphaThr}{\ensuremath{1.5\cdot10^{-4}}}
% NHOneSqrtLamThr [hand, exact] hand :: 
\newcommand{\NHOneSqrtLamThr}{\ensuremath{5\cdot10^{-8}}}
% NYbarStep [hand, exact] hand :: 
\newcommand{\NYbarStep}{\ensuremath{1\cdot10^{-12}}}
% NHTwoThrZero [hand, exact] hand :: 
\newcommand{\NHTwoThrZero}{\ensuremath{1\cdot10^{-8}}}
% NHTwoThrOne [hand, exact] hand :: 
\newcommand{\NHTwoThrOne}{\ensuremath{1\cdot10^{-6}}}
% NHTwoThrTwo [hand, exact] hand :: 
\newcommand{\NHTwoThrTwo}{\ensuremath{1\cdot10^{-3}}}
% NHThreeThrZero [hand, exact] hand :: 
\newcommand{\NHThreeThrZero}{\ensuremath{1\cdot10^{-6}}}
% NHThreeThrOne [hand, exact] hand :: 
\newcommand{\NHThreeThrOne}{\ensuremath{5\cdot10^{-6}}}
% NHThreeThrTwo [hand, exact] hand :: 
\newcommand{\NHThreeThrTwo}{\ensuremath{1\cdot10^{-5}}}
% NHFourThrZero [hand, exact] hand :: 
\newcommand{\NHFourThrZero}{\ensuremath{2\cdot10^{-7}}}
% NHFourThrOne [hand, exact] hand :: 
\newcommand{\NHFourThrOne}{\ensuremath{1\cdot10^{-5}}}
% NHFourThrTwo [hand, exact] hand :: 
\newcommand{\NHFourThrTwo}{\ensuremath{1\cdot10^{-5}}}
% NHFourThrThree [hand, exact] hand :: 
\newcommand{\NHFourThrThree}{\ensuremath{1\cdot10^{-5}}}
% NHFourKfourThr [hand, exact] hand :: 
\newcommand{\NHFourKfourThr}{\ensuremath{0.06}}
% NHSixConeSlopeThr [hand, exact] hand :: 
\newcommand{\NHSixConeSlopeThr}{\ensuremath{0.08}}
% NHSixConeKtwoThr [hand, exact] hand :: 
\newcommand{\NHSixConeKtwoThr}{\ensuremath{0.03}}
% NHSixConeNthreeThr [hand, exact] hand :: 
\newcommand{\NHSixConeNthreeThr}{\ensuremath{0.004}}
% NHSixConeJetThr [hand, exact] hand :: 
\newcommand{\NHSixConeJetThr}{\ensuremath{0.12}}
% NHSixConeHFiveThr [hand, exact] hand :: 
\newcommand{\NHSixConeHFiveThr}{\ensuremath{0.0144}}
% NHSixFarHFiveThr [hand, exact] hand :: 
\newcommand{\NHSixFarHFiveThr}{\ensuremath{0.0225}}
% NLemmaOConeM [hand, exact] hand :: 
\newcommand{\NLemmaOConeM}{\ensuremath{0.12}}
% NLemmaOConeK [hand, exact] hand :: 
\newcommand{\NLemmaOConeK}{\ensuremath{0.18}}
% NLemmaOFarM [hand, exact] hand :: 
\newcommand{\NLemmaOFarM}{\ensuremath{0.15}}
% NLemmaOFarK [hand, exact] hand :: 
\newcommand{\NLemmaOFarK}{\ensuremath{0.3}}
% NCoercivityConst [hand, exact] hand :: 
\newcommand{\NCoercivityConst}{\ensuremath{0.046875}}
% NMOnez [hand, exact] hand :: 
\newcommand{\NMOnez}{\ensuremath{0.534}}
% NBetaA [hand, exact] hand :: 
\newcommand{\NBetaA}{\ensuremath{2}}
% NTOneLower [cert, lower] t1clock.json :: t1.0
\newcommand{\NTOneLower}{\ensuremath{4.5978574775}}
% NTOneUpper [cert, upper] t1clock.json :: t1.1
\newcommand{\NTOneUpper}{\ensuremath{4.5978586144}}
% NLambdaOneLower [cert, lower] t1clock.json :: lnlam.0
\newcommand{\NLambdaOneLower}{\ensuremath{6.49\cdot10^{-5}}}
% NLambdaOneUpper [cert, upper] t1clock.json :: lnlam.1
\newcommand{\NLambdaOneUpper}{\ensuremath{6.75\cdot10^{-5}}}
% NTMinusTOneUpper [cert, upper] t1clock.json :: lnlam.1; 1.5e-4.sup_ab
\newcommand{\NTMinusTOneUpper}{\ensuremath{2.1\cdot10^{-17}}}
% NAbarUpper [cert, upper] final_alpha_points_5e-8.json :: 1.5e-4.sup_ab
\newcommand{\NAbarUpper}{\ensuremath{2.19103\cdot10^{-7}}}
% NBetabarUpper [cert, upper] final_alpha_points_5e-8.json :: 1.5e-4.sup_beta
\newcommand{\NBetabarUpper}{\ensuremath{3.33811\cdot10^{-9}}}
% NBaseKmaxLower [cert, lower] t1clock.json :: base_kmax.0
\newcommand{\NBaseKmaxLower}{\ensuremath{2.279\cdot10^{4}}}
% NBaseKmaxUpper [cert, upper] t1clock.json :: base_kmax.1
\newcommand{\NBaseKmaxUpper}{\ensuremath{2.376\cdot10^{4}}}
% NHOneAlphaBase [cert, upper] summary.json :: alpha_abs_max
\newcommand{\NHOneAlphaBase}{\ensuremath{1.1739\cdot10^{-4}}}
% NHOneAlphaErr [cert, upper] t9_h1_clam_s180.json :: C_alpha; inputs.YSTAR
\newcommand{\NHOneAlphaErr}{\ensuremath{7.21\cdot10^{-13}}}
% NHOneAlphaMargin [cert, margin] summary.json :: alpha_abs_max; C_alpha; inputs.YSTAR
\newcommand{\NHOneAlphaMargin}{\ensuremath{1.27}}
% NHOneSqrtLamBase [cert, upper] t9_h1_clam_s180.json :: sqrtLam_base_Wp
\newcommand{\NHOneSqrtLamBase}{\ensuremath{2.282\cdot10^{-8}}}
% NHOneSqrtLamTotal [cert, upper] t9_h1_clam_s180.json :: sqrtLam_total_Wp
\newcommand{\NHOneSqrtLamTotal}{\ensuremath{2.29\cdot10^{-8}}}
% NHOneSqrtLamMargin [cert, margin] t9_h1_clam_s180.json :: sqrtLam_total_Wp
\newcommand{\NHOneSqrtLamMargin}{\ensuremath{2.18}}
% NHOneYbarLimit [cert, lower] t9_h1_clam_s180.json :: Ybar_limit_H1
\newcommand{\NHOneYbarLimit}{\ensuremath{1.12\cdot10^{-9}}}
% NSwitchYUpper [cert, upper] switch_s2.json :: modulated.Y_raw
\newcommand{\NSwitchYUpper}{\ensuremath{6.21\cdot10^{-13}}}
% NSwitchYMargin [cert, margin] switch_s2.json :: modulated.Y_raw
\newcommand{\NSwitchYMargin}{\ensuremath{1.6}}
% NMinCI [cert, lower] constants.json :: meta.summary.min_c_I
\newcommand{\NMinCI}{\ensuremath{0.05283}}
% NNumCertificates [cert, exact] constants.json :: meta.summary.n_cert
\newcommand{\NNumCertificates}{\ensuremath{26}}
% NMinRhoC [cert, lower] constants.json :: meta.closure.rho_c_min
\newcommand{\NMinRhoC}{\ensuremath{0.04711}}
% NHTwoBaseZero [cert, upper] FINAL_manifest_v33.json :: results.M9_H2_tightened.rows.0.base
\newcommand{\NHTwoBaseZero}{\ensuremath{8.011\cdot10^{-11}}}
% NHTwoBaseOne [cert, upper] FINAL_manifest_v33.json :: results.M9_H2_tightened.rows.1.base
\newcommand{\NHTwoBaseOne}{\ensuremath{1.336\cdot10^{-10}}}
% NHTwoBaseTwo [cert, upper] FINAL_manifest_v33.json :: results.M9_H2_tightened.rows.2.base
\newcommand{\NHTwoBaseTwo}{\ensuremath{2.936\cdot10^{-10}}}
% NHTwoErrWZero [cert, upper] FINAL_manifest_v33.json :: results.M9_H2_tightened.rows.0.with_error
\newcommand{\NHTwoErrWZero}{\ensuremath{8.014\cdot10^{-11}}}
% NHTwoMarginWZero [cert, margin] FINAL_manifest_v33.json :: results.M9_H2_tightened.rows.0.margin
\newcommand{\NHTwoMarginWZero}{\ensuremath{124.78}}
% NHTwoErrWOne [cert, upper] FINAL_manifest_v33.json :: results.M9_H2_tightened.rows.1.with_error
\newcommand{\NHTwoErrWOne}{\ensuremath{4.706\cdot10^{-10}}}
% NHTwoMarginWOne [cert, margin] FINAL_manifest_v33.json :: results.M9_H2_tightened.rows.1.margin
\newcommand{\NHTwoMarginWOne}{\ensuremath{2124.9}}
% NHTwoErrWTwo [cert, upper] FINAL_manifest_v33.json :: results.M9_H2_tightened.rows.2.with_error
\newcommand{\NHTwoErrWTwo}{\ensuremath{1.559\cdot10^{-6}}}
% NHTwoMarginWTwo [cert, margin] FINAL_manifest_v33.json :: results.M9_H2_tightened.rows.2.margin
\newcommand{\NHTwoMarginWTwo}{\ensuremath{641.53}}
% NHThreeBaseZero [cert, upper] FINAL_manifest_v33.json :: results.M9_H2f.rows.0.base
\newcommand{\NHThreeBaseZero}{\ensuremath{3.127\cdot10^{-7}}}
% NHThreeBaseOne [cert, upper] FINAL_manifest_v33.json :: results.M9_H2f.rows.1.base
\newcommand{\NHThreeBaseOne}{\ensuremath{1.004\cdot10^{-6}}}
% NHThreeBaseTwo [cert, upper] FINAL_manifest_v33.json :: results.M9_H2f.rows.2.base
\newcommand{\NHThreeBaseTwo}{\ensuremath{2.431\cdot10^{-6}}}
% NHThreeErrZero [cert, upper] FINAL_manifest_v33.json :: results.M9_H2f.rows.0.with_error
\newcommand{\NHThreeErrZero}{\ensuremath{3.1363\cdot10^{-7}}}
% NHThreeMarginZero [cert, margin] FINAL_manifest_v33.json :: results.M9_H2f.rows.0.margin
\newcommand{\NHThreeMarginZero}{\ensuremath{3.1885}}
% NHThreeErrOne [cert, upper] FINAL_manifest_v33.json :: results.M9_H2f.rows.1.with_error
\newcommand{\NHThreeErrOne}{\ensuremath{1.0034\cdot10^{-6}}}
% NHThreeMarginOne [cert, margin] FINAL_manifest_v33.json :: results.M9_H2f.rows.1.margin
\newcommand{\NHThreeMarginOne}{\ensuremath{4.9834}}
% NHThreeErrTwo [cert, upper] FINAL_manifest_v33.json :: results.M9_H2f.rows.2.with_error
\newcommand{\NHThreeErrTwo}{\ensuremath{2.4351\cdot10^{-6}}}
% NHThreeMarginTwo [cert, margin] FINAL_manifest_v33.json :: results.M9_H2f.rows.2.margin
\newcommand{\NHThreeMarginTwo}{\ensuremath{4.1067}}
% NHFourBaseZero [cert, upper] FINAL_manifest_v33.json :: results.M9_base_sums.combined.windows_max.0
\newcommand{\NHFourBaseZero}{\ensuremath{3.48\cdot10^{-8}}}
% NHFourErrZero [cert, upper] FINAL_manifest_v33.json :: results.M9_H3_H2v32.rows.win.0.with_error_arb
\newcommand{\NHFourErrZero}{\ensuremath{3.5087\cdot10^{-8}}}
% NHFourMarginZero [cert, margin] FINAL_manifest_v33.json :: results.M9_H3_H2v32.rows.win.0.margin_arb
\newcommand{\NHFourMarginZero}{\ensuremath{5.700}}
% NHFourBaseOne [cert, upper] FINAL_manifest_v33.json :: results.M9_base_sums.combined.windows_max.1
\newcommand{\NHFourBaseOne}{\ensuremath{1.27\cdot10^{-7}}}
% NHFourErrOne [cert, upper] FINAL_manifest_v33.json :: results.M9_H3_H2v32.rows.win.1.with_error_arb
\newcommand{\NHFourErrOne}{\ensuremath{1.2616\cdot10^{-7}}}
% NHFourMarginOne [cert, margin] FINAL_manifest_v33.json :: results.M9_H3_H2v32.rows.win.1.margin_arb
\newcommand{\NHFourMarginOne}{\ensuremath{79.26}}
% NHFourBaseTwo [cert, upper] FINAL_manifest_v33.json :: results.M9_base_sums.combined.windows_max.2
\newcommand{\NHFourBaseTwo}{\ensuremath{3.76\cdot10^{-7}}}
% NHFourErrTwo [cert, upper] FINAL_manifest_v33.json :: results.M9_H3_H2v32.rows.win.2.with_error_arb
\newcommand{\NHFourErrTwo}{\ensuremath{3.7724\cdot10^{-7}}}
% NHFourMarginTwo [cert, margin] FINAL_manifest_v33.json :: results.M9_H3_H2v32.rows.win.2.margin_arb
\newcommand{\NHFourMarginTwo}{\ensuremath{26.50}}
% NHFourBaseThree [cert, upper] FINAL_manifest_v33.json :: results.M9_base_sums.combined.windows_max.3
\newcommand{\NHFourBaseThree}{\ensuremath{9.71\cdot10^{-7}}}
% NHFourErrThree [cert, upper] FINAL_manifest_v33.json :: results.M9_H3_H2v32.rows.win.3.with_error_arb
\newcommand{\NHFourErrThree}{\ensuremath{1.2822\cdot10^{-6}}}
% NHFourMarginThree [cert, margin] FINAL_manifest_v33.json :: results.M9_H3_H2v32.rows.win.3.margin_arb
\newcommand{\NHFourMarginThree}{\ensuremath{7.799}}
% NHFourYbarLimit [cert, lower] FINAL_manifest_v33.json :: results.M9_H3_H2v32.binding_Ybar_limit
\newcommand{\NHFourYbarLimit}{\ensuremath{1.575\cdot10^{-11}}}
% NHFourYbarMargin [cert, margin] FINAL_manifest_v33.json :: results.M9_H3_H2v32.binding_Ybar_limit
\newcommand{\NHFourYbarMargin}{\ensuremath{15.75}}
% NHFourKfourBase [cert, upper] amp_A1_w180.json :: inputs.BASE_T1.K4
\newcommand{\NHFourKfourBase}{\ensuremath{3.1633\cdot10^{-4}}}
% NHSixConeSlope [cert, upper] FINAL_manifest_v33.json :: results.cone_fixedL.checks.slope.value
\newcommand{\NHSixConeSlope}{\ensuremath{0.03029}}
% NHSixConeSlopeMargin [cert, margin] FINAL_manifest_v33.json :: results.cone_fixedL.checks.slope.margin
\newcommand{\NHSixConeSlopeMargin}{\ensuremath{2.641}}
% NHSixConeKtwo [cert, upper] FINAL_manifest_v33.json :: results.cone_fixedL.checks.K2.value
\newcommand{\NHSixConeKtwo}{\ensuremath{0.009876}}
% NHSixConeKtwoMargin [cert, margin] FINAL_manifest_v33.json :: results.cone_fixedL.checks.K2.margin
\newcommand{\NHSixConeKtwoMargin}{\ensuremath{3.037}}
% NHSixConeNthree [cert, upper] FINAL_manifest_v33.json :: results.cone_fixedL.checks.N3.value
\newcommand{\NHSixConeNthree}{\ensuremath{1.461\cdot10^{-4}}}
% NHSixConeNthreeMargin [cert, margin] FINAL_manifest_v33.json :: results.cone_fixedL.checks.N3.margin
\newcommand{\NHSixConeNthreeMargin}{\ensuremath{27.38}}
% NHSixConeHFive [cert, upper] FINAL_manifest_v33.json :: results.H4_fixedL.total
\newcommand{\NHSixConeHFive}{\ensuremath{1.2865\cdot10^{-5}}}
% NHSixConeHFiveMargin [cert, margin] FINAL_manifest_v33.json :: results.H4_fixedL.margin
\newcommand{\NHSixConeHFiveMargin}{\ensuremath{1119}}
% NHSixConeTstarNeeded [cert, upper] FINAL_manifest_v33.json :: results.H7_fixedL.ratio_upper
\newcommand{\NHSixConeTstarNeeded}{\ensuremath{1.6000015\cdot10^{-7}}}
% NHSixConeTstarLower [cert, lower] FINAL_manifest_v33.json :: results.H7_fixedL.t_star_lower
\newcommand{\NHSixConeTstarLower}{\ensuremath{3.68334\cdot10^{-7}}}
% NHSixConeTstarMargin [cert, margin] FINAL_manifest_v33.json :: results.H7_fixedL.margin_tstar_over_ratio
\newcommand{\NHSixConeTstarMargin}{\ensuremath{2.302}}
% NHSixConeHFiveRatioPeak [cert, upper] h4loc_A1_w180.json :: windows.constants/amp_A1_w180.json.worst.ratio
\newcommand{\NHSixConeHFiveRatioPeak}{\ensuremath{0.2393}}
% NHSixConeHOneSup [cert, upper] FINAL_manifest_v33.json :: results.H1_fixedL.sup_h0_max
\newcommand{\NHSixConeHOneSup}{\ensuremath{0.2163}}
% NHSixConeHOneH [cert, upper] FINAL_manifest_v33.json :: results.H1_fixedL.H_max
\newcommand{\NHSixConeHOneH}{\ensuremath{0.2185}}
% NHSixFarHFiveRatio [cert, upper] h4loc_v32.json :: windows.constants/amp_A1_w180.json.cylinders.18.ratio
\newcommand{\NHSixFarHFiveRatio}{\ensuremath{0.1761}}
% NHSixFarHFiveMarginOld [cert, margin] h4loc_v32.json :: windows.constants/amp_A1_w180.json.cylinders.18.ratio
\newcommand{\NHSixFarHFiveMarginOld}{\ensuremath{5.68}}
% NHSixFarHFiveMargin [cert, margin] FINAL_manifest_v33.json :: results.M9_H4_far.min_margin
\newcommand{\NHSixFarHFiveMargin}{\ensuremath{15.58}}
% NHSixConeHFiveMarginOld [cert, margin] h4loc_A1_w180.json :: windows.constants/amp_A1_w180.json.worst.ratio
\newcommand{\NHSixConeHFiveMarginOld}{\ensuremath{4.18}}
% NHSixFarKtwo [cert, upper] t9_far_v32.json :: h2
\newcommand{\NHSixFarKtwo}{\ensuremath{1.797\cdot10^{-4}}}
% NHSixFarKtwoThr [hand, exact] t9_far_v32.json :: thr.h2
\newcommand{\NHSixFarKtwoThr}{\ensuremath{3.53\cdot10^{-4}}}
% NHSixFarKtwoMargin [cert, margin] t9_far_v32.json :: margins.h2
\newcommand{\NHSixFarKtwoMargin}{\ensuremath{1.964}}
% NHSixFarKthree [cert, upper] t9_far_v32.json :: h3
\newcommand{\NHSixFarKthree}{\ensuremath{2.595\cdot10^{-8}}}
% NHSixFarKthreeThr [hand, exact] t9_far_v32.json :: thr.h3
\newcommand{\NHSixFarKthreeThr}{\ensuremath{1.17\cdot10^{-5}}}
% NHSixFarKthreeMargin [cert, margin] t9_far_v32.json :: margins.h3
\newcommand{\NHSixFarKthreeMargin}{\ensuremath{450.9}}
% NHSixFarKfour [cert, upper] t9_far_v32.json :: h4
\newcommand{\NHSixFarKfour}{\ensuremath{6.964\cdot10^{-9}}}
% NHSixFarKfourThr [hand, exact] t9_far_v32.json :: thr.h4
\newcommand{\NHSixFarKfourThr}{\ensuremath{1.03\cdot10^{-7}}}
% NHSixFarKfourMargin [cert, margin] t9_far_v32.json :: margins.h4
\newcommand{\NHSixFarKfourMargin}{\ensuremath{14.79}}
% NHSixFarSlope [cert, upper] t9_far_v32.json :: slope
\newcommand{\NHSixFarSlope}{\ensuremath{0.04153}}
% NHSixFarSlopeThr [hand, exact] t9_far_v32.json :: thr.slope
\newcommand{\NHSixFarSlopeThr}{\ensuremath{0.0800}}
% NHSixFarSlopeMargin [cert, margin] t9_far_v32.json :: margins.slope
\newcommand{\NHSixFarSlopeMargin}{\ensuremath{1.926}}
% NHSixFarNthree [cert, upper] t9_far_v32.json :: L2
\newcommand{\NHSixFarNthree}{\ensuremath{5.516\cdot10^{-7}}}
% NHSixFarNthreeThr [hand, exact] t9_far_v32.json :: thr.L2
\newcommand{\NHSixFarNthreeThr}{\ensuremath{3.32\cdot10^{-6}}}
% NHSixFarNthreeMargin [cert, margin] t9_far_v32.json :: margins.L2
\newcommand{\NHSixFarNthreeMargin}{\ensuremath{6.019}}
% NFarTstarMargin [cert, margin] lemmaO_v32_consts.json :: far_v32.case.margin_factor_lower
\newcommand{\NFarTstarMargin}{\ensuremath{15.5}}
% NFarTstarLower [cert, lower] lemmaO_v32_consts.json :: far_v32.case.t_star_lower
\newcommand{\NFarTstarLower}{\ensuremath{9.5108\cdot10^{-8}}}
% NFarTstarNeeded [cert, upper] lemmaO_v32_consts.json :: far_v32.case.needed
\newcommand{\NFarTstarNeeded}{\ensuremath{6.134\cdot10^{-9}}}
% NGammaModes [cert, exact] gamma0_final.json :: M
\newcommand{\NGammaModes}{\ensuremath{3762}}
% NGammaHalfwidth [cert, upper] gamma0_final.json :: b_j_halfwidth
\newcommand{\NGammaHalfwidth}{\ensuremath{8.138\cdot10^{-24}}}
% NGammaContraction [cert, upper] gamma0_final.json :: krawczyk_contraction
\newcommand{\NGammaContraction}{\ensuremath{3.01\cdot10^{-12}}}
% NGammaArcs [cert, exact] gamma0_final.json :: embedded.arcs
\newcommand{\NGammaArcs}{\ensuremath{701}}
% NGammaMaxTurn [cert, upper] gamma0_final.json :: embedded.maxturn
\newcommand{\NGammaMaxTurn}{\ensuremath{0.04997}}
% NGammaMaxPair [cert, upper] gamma0_final.json :: embedded.maxpair
\newcommand{\NGammaMaxPair}{\ensuremath{0.09597}}
% NGammaMinSep [cert, lower] gamma0_final.json :: embedded.minsep
\newcommand{\NGammaMinSep}{\ensuremath{8.89\cdot10^{-3}}}
% NGammaHalfGap [cert, lower] gamma0_final.json :: embedded.gap_cert
\newcommand{\NGammaHalfGap}{\ensuremath{0.0248219}}
% NPhaseISlabs [cert, exact] status.json :: slab
\newcommand{\NPhaseISlabs}{\ensuremath{96}}
% NPhaseIModes [cert, exact] status.json :: last_slab.M
\newcommand{\NPhaseIModes}{\ensuremath{2800}}
% NPhaseIKmaxLower [cert, lower] status.json :: kmax_enclosure.0
\newcommand{\NPhaseIKmaxLower}{\ensuremath{29.352874545}}
% NPhaseIKmaxUpper [cert, upper] status.json :: kmax_enclosure.1
\newcommand{\NPhaseIKmaxUpper}{\ensuremath{29.352874546}}
% NPhaseITLower [cert, lower] status.json :: t
\newcommand{\NPhaseITLower}{\ensuremath{4.597850539485830}}
% NPhaseITUpper [cert, upper] status.json :: t
\newcommand{\NPhaseITUpper}{\ensuremath{4.597850539485831}}
% NPeakSqrtLam [cert, upper] rate_table_fin12f4_a0.00015.json :: sqrtL_max
\newcommand{\NPeakSqrtLam}{\ensuremath{5.308\cdot10^{-8}}}
% NTailRate [cert, lower] rate_table_fin12f4_a0.00015.json :: tail.rho
\newcommand{\NTailRate}{\ensuremath{2.693\cdot10^{-3}}}
% NFloorSqrtLam [cert, upper] rate_table_fin12f4_a0.00015.json :: tail.sqrt_LF
\newcommand{\NFloorSqrtLam}{\ensuremath{9.872\cdot10^{-10}}}
% NAlphaFloor [cert, upper] rate_table_fin12f4_a0.00015.json :: tail.alpha_F
\newcommand{\NAlphaFloor}{\ensuremath{6.867\cdot10^{-7}}}
% NSupBetaVerifier [cert, upper] final_alpha_points_5e-8.json :: 1.5e-4.sup_beta
\newcommand{\NSupBetaVerifier}{\ensuremath{3.33811\cdot10^{-9}}}
% NMOneRamp [cert, upper] final_alpha_points_5e-8.json :: 1.5e-4.M1
\newcommand{\NMOneRamp}{\ensuremath{0.4328}}
% NSITwoPrime [cert, upper] FINAL_manifest_v33.json :: results.F12_sliver.sI2_prime
\newcommand{\NSITwoPrime}{\ensuremath{0.015457}}
% NLoopGainCoarse [cert, upper] U_cert_v35.json :: loop_gain.2.total
\newcommand{\NLoopGainCoarse}{\ensuremath{0.03904}}
% NRateMarginCoarse [cert, margin] U_cert_v35.json :: rate_margin
\newcommand{\NRateMarginCoarse}{\ensuremath{28.2}}
% NSOneEpsLower [cert, lower] bracket2.json :: eps_a
\newcommand{\NSOneEpsLower}{\ensuremath{1.5346198892965\cdot10^{-4}}}
% NSOneEpsUpper [cert, upper] bracket2.json :: eps_b
\newcommand{\NSOneEpsUpper}{\ensuremath{1.5346198899299\cdot10^{-4}}}
% NSOneRowsTip [cert, exact] ext_tip_check.json :: rows
\newcommand{\NSOneRowsTip}{\ensuremath{51246}}
% NSOneTipWidth [cert, upper] ext_tip_check.json :: max_pointwise_width.th
\newcommand{\NSOneTipWidth}{\ensuremath{3.2\cdot10^{-12}}}
% NSOneExtWidth [cert, upper] ext813_check.json :: max_pointwise_width.p
\newcommand{\NSOneExtWidth}{\ensuremath{5.7\cdot10^{-10}}}
% NSOneRowsMC [cert, exact] check_outA5.json :: prof_tip_vs_MC.rows_compared
\newcommand{\NSOneRowsMC}{\ensuremath{1025}}
% NSOneRowsMB [cert, exact] check_outA5.json :: ext813_vs_MB.rows_compared
\newcommand{\NSOneRowsMB}{\ensuremath{501}}
% NSOneMinMargin [cert, margin] check_outA5.json :: min:*.min_relative_margin.*
\newcommand{\NSOneMinMargin}{\ensuremath{1.2\cdot10^{-9}}}
% NTwoSegments [cert, exact] contour_result.json :: len:segments
\newcommand{\NTwoSegments}{\ensuremath{5058}}
% NTwoInteriorBoxes [cert, exact] interior_result.json :: len:
\newcommand{\NTwoInteriorBoxes}{\ensuremath{174}}
% NTwoContourBoxes [cert, exact] eta_cache.json :: len:
\newcommand{\NTwoContourBoxes}{\ensuremath{202}}
% NTwoEtaMax [cert, upper] eta_cache.json :: max:*.eta
\newcommand{\NTwoEtaMax}{\ensuremath{0.1163}}
% NTwoLipMax [cert, upper] eta_cache.json :: max:*.Lip
\newcommand{\NTwoLipMax}{\ensuremath{0.272}}
% NTwoLambdaZero [est, upper] cert_b0.25_th0.0_tau0.03.json :: Lam
\newcommand{\NTwoLambdaZero}{\ensuremath{1.790819}}
% NTwoLambdaOne [est, upper] cert_b0.25_th0.48_tau0.03.json :: Lam
\newcommand{\NTwoLambdaOne}{\ensuremath{1.768282}}
% NSOneDTightLower [doc, lower] doc:README.md :: D ∈ [2.50080627654, 2.50080627668]
\newcommand{\NSOneDTightLower}{\ensuremath{2.50080627654}}
% NSOneDTightUpper [doc, upper] doc:README.md :: D ∈ [2.50080627654, 2.50080627668]
\newcommand{\NSOneDTightUpper}{\ensuremath{2.50080627668}}
% NSOneKTightLower [doc, lower] doc:README.md :: k(0) ∈ [1.53982063762, 1.53982063822]
\newcommand{\NSOneKTightLower}{\ensuremath{1.53982063762}}
% NSOneKTightUpper [doc, upper] doc:README.md :: k(0) ∈ [1.53982063762, 1.53982063822]
\newcommand{\NSOneKTightUpper}{\ensuremath{1.53982063822}}
% NSOneDCondLower [doc, lower] doc:README.md :: D = −p(tip) ∈ [2.500806276619695, 2.500806276620049]
\newcommand{\NSOneDCondLower}{\ensuremath{2.500806276619695}}
% NSOneDCondUpper [doc, upper] doc:README.md :: D = −p(tip) ∈ [2.500806276619695, 2.500806276620049]
\newcommand{\NSOneDCondUpper}{\ensuremath{2.500806276620049}}
% NSOneKCondLower [doc, lower] doc:README.md :: k(0) ∈ [1.539820637871836, 1.5398206378721015]
\newcommand{\NSOneKCondLower}{\ensuremath{1.539820637871836}}
% NSOneKCondUpper [doc, upper] doc:README.md :: k(0) ∈ [1.539820637871836, 1.5398206378721015]
\newcommand{\NSOneKCondUpper}{\ensuremath{1.5398206378721015}}
% NSOneTauTipLower [doc, lower] doc:README.md :: τ_tip ∈ [10.249077651120778, 10.249077651121361]
\newcommand{\NSOneTauTipLower}{\ensuremath{10.249077651120778}}
% NSOneTauTipUpper [doc, upper] doc:README.md :: τ_tip ∈ [10.249077651120778, 10.249077651121361]
\newcommand{\NSOneTauTipUpper}{\ensuremath{10.249077651121361}}
% NSOneStateUncert [doc, upper] doc:README.md :: State uncertainty at σ = 8: 4.4e−13
\newcommand{\NSOneStateUncert}{\ensuremath{4.4\cdot10^{-13}}}
% NFarBoxEprime [doc, margin] doc:STEP4_RESULT.md :: (6.4287, 4.6917, 1.9940, 1.69999995)
\newcommand{\NFarBoxEprime}{\ensuremath{6.42}}
% NFarBoxXEsecond [doc, margin] doc:STEP4_RESULT.md :: (6.4287, 4.6917, 1.9940, 1.69999995)
\newcommand{\NFarBoxXEsecond}{\ensuremath{4.69}}
% NFarBoxXsqEthird [doc, margin] doc:STEP4_RESULT.md :: (6.4287, 4.6917, 1.9940, 1.69999995)
\newcommand{\NFarBoxXsqEthird}{\ensuremath{1.99}}
% NFarBoxE [doc, margin] doc:STEP4_RESULT.md :: (6.4287, 4.6917, 1.9940, 1.69999995)
\newcommand{\NFarBoxE}{\ensuremath{1.69}}
% NHFourKfourErr [doc, upper] doc:STEP4_RESULT.md :: | 5.60e-3 against 0.06 | ×10.7 |
\newcommand{\NHFourKfourErr}{\ensuremath{5.60\cdot10^{-3}}}
% NHFourKfourMargin [doc, margin] doc:STEP4_RESULT.md :: | 5.60e-3 against 0.06 | ×10.7 |
\newcommand{\NHFourKfourMargin}{\ensuremath{10.7}}
% NCoercivityCoefficient [doc, upper] doc:step3_proof_v32.tex :: \norm{e^{\tau\mathcal L}u}_\rho\le14.8\,e^{-0.0213\tau}\norm
\newcommand{\NCoercivityCoefficient}{\ensuremath{14.8}}
% NCoercivityRate [doc, lower] doc:step3_proof_v32.tex :: \norm{e^{\tau\mathcal L}u}_\rho\le14.8\,e^{-0.0213\tau}\norm
\newcommand{\NCoercivityRate}{\ensuremath{0.0213}}
% NRegionRampStart [hand, exact] hand :: 
\newcommand{\NRegionRampStart}{\ensuremath{14}}
% NRegionInterface [hand, exact] hand :: 
\newcommand{\NRegionInterface}{\ensuremath{28}}
% NRegionShellEnd [hand, exact] hand :: 
\newcommand{\NRegionShellEnd}{\ensuremath{60}}
% NRegionSegEnd [hand, exact] hand :: 
\newcommand{\NRegionSegEnd}{\ensuremath{350}}
% NRegionAnnEnd [hand, exact] hand :: 
\newcommand{\NRegionAnnEnd}{\ensuremath{700}}
% NZoneIIILow [hand, exact] hand :: 
\newcommand{\NZoneIIILow}{\ensuremath{253}}
% NZoneIIIHigh [hand, exact] hand :: 
\newcommand{\NZoneIIIHigh}{\ensuremath{379.5}}
% NConeCentreMin [hand, exact] hand :: 
\newcommand{\NConeCentreMin}{\ensuremath{300}}
% NConeCentreMax [hand, exact] hand :: 
\newcommand{\NConeCentreMax}{\ensuremath{380}}
% NFarCentreMin [hand, exact] hand :: 
\newcommand{\NFarCentreMin}{\ensuremath{490}}
% NConeL [hand, exact] hand :: 
\newcommand{\NConeL}{\ensuremath{50}}
% NConeX [hand, exact] hand :: 
\newcommand{\NConeX}{\ensuremath{3}}
% NFarL [hand, exact] hand :: 
\newcommand{\NFarL}{\ensuremath{113}}
% NFarX [hand, exact] hand :: 
\newcommand{\NFarX}{\ensuremath{2}}
 is found.
% ======================================================================

% ---- which document are we in ----------------------------------------
\ifdefined\ifsupplement\else\newif\ifsupplement\fi

% ---- packages ---------------------------------------------------------
\usepackage[T1]{fontenc}
\usepackage[utf8]{inputenc}
\usepackage{lmodern}
\usepackage[margin=1.1in]{geometry}
\usepackage{amsmath,amssymb,amsthm,mathtools,bm}
\usepackage{booktabs,array,longtable,tabularx,multirow}
\usepackage{enumitem}
\usepackage{graphicx}
\usepackage{xcolor}
\usepackage{tikz}
\usetikzlibrary{positioning,arrows.meta,calc}
% xr-hyper must come before hyperref
\usepackage{xr-hyper}
\usepackage[colorlinks=true,linkcolor=blue!50!black,citecolor=green!35!black,urlcolor=blue!50!black]{hyperref}

\raggedbottom
\setlength{\emergencystretch}{2em}
\setlist[enumerate]{leftmargin=*,itemsep=2pt,topsep=3pt}
\setlist[itemize]{leftmargin=*,itemsep=2pt,topsep=3pt}

% ---- numbering --------------------------------------------------------
% Main paper: sections 1,2,...; theorems 2.3 (by section); equations (2.3).
% Supplement: the Parts A..K are the \section's, numbered by letter, so
% theorems are A.1, E.4, ... (do not confuse "Theorem A" of the paper, which
% has no dot, with "Theorem A.1" of Part A).
\ifsupplement
  \renewcommand{\thesection}{\Alph{section}}
\fi
\numberwithin{equation}{section}

% ---- theorem environments (numbered by section, one shared counter) ----
\theoremstyle{plain}
\newtheorem{theorem}{Theorem}[section]
\newtheorem{proposition}[theorem]{Proposition}
\newtheorem{lemma}[theorem]{Lemma}
\newtheorem{corollary}[theorem]{Corollary}
\newtheorem{maintheorem}[theorem]{Main Theorem}
% Theorems A, B, C, D of the paper: own counter, lettered, no section prefix
\newtheorem{thmlet}{Theorem}
\renewcommand{\thethmlet}{\Alph{thmlet}}
\theoremstyle{definition}
\newtheorem{definition}[theorem]{Definition}
\newtheorem{convention}[theorem]{Convention}
\newtheorem{example}[theorem]{Example}
\theoremstyle{remark}
\newtheorem{remark}[theorem]{Remark}

% A list of hypotheses labelled (H1),(H2),... : \begin{hyplist}{H} ... \end{hyplist}
% (also {O}, {L}, {B}, {K}, {W} for the families of section 5.3 of OUTLINE.md).
\newenvironment{hyplist}[1]%
  {\begin{enumerate}[label=\textup{(#1\arabic*)},leftmargin=3em,itemsep=2pt]}%
  {\end{enumerate}}

% ---- trust classes (OUTLINE 4.1, 5.1.3) -------------------------------
% Supplement: explicit tags.  Main paper: certified is the default and
% unmarked; every non-certified number is marked "(floating point)".
\ifsupplement
  \newcommand{\tcert}{\textup{[cert]}}   % arb/outward-rounded, pinned script and key
  \newcommand{\test}{\textup{[est]}}     % floating point or estimate; never in a theorem
  \newcommand{\thand}{\textup{[hand]}}   % proved by hand in the text
  \newcommand{\tlit}{\textup{[lit]}}     % cited from the literature
\else
  \newcommand{\tcert}{}
  \newcommand{\test}{\unskip\ \textup{(floating point)}}
  \newcommand{\thand}{}
  \newcommand{\tlit}{}
\fi

% ---- gaps: visible, counted, an error in a final build -----------------
% Usage: \GAP{one sentence saying what is missing}   (and an entry in ISSUES.md)
\newcounter{gapcount}
\newif\iffinalpaper   % \finalpapertrue makes every \GAP a hard error
\newcommand{\GAP}[1]{\stepcounter{gapcount}\typeout{GAP-MARKER \thegapcount}%
  \iffinalpaper\PackageError{paper}{\string\GAP\space used in a final build}{Remove all gaps.}\fi
  {\color{red}\sffamily\bfseries[GAP~\thegapcount:\ }{\color{red}\sffamily #1}{\color{red}\sffamily\bfseries]}}
\AtEndDocument{\typeout{GAP-MARKER total: \thegapcount}}

% ---- file names and JSON keys (supplement only) -------------------------
% No internal path appears in the main paper; a use there is a warning.
\newcommand{\fn}[1]{%
  \ifsupplement\else\PackageWarning{paper}{\string\fn\space used in the main paper}\fi
  \texttt{\detokenize{#1}}}
\newcommand{\key}[1]{\texttt{\detokenize{#1}}}

% ---- cross references between main paper and supplement ----------------
% Labels carry a unit prefix (section 5.3 of OUTLINE.md and DRAFTER_GUIDE.md):
%   main paper:  intro: obj: s3: s4: s5: s6: s7: s8: s9:   (kind: thm- lem- eq- tab- fig- ...)
%   supplement:  pA: pB: ... pK:
% so all labels are globally unique and the prefixes below are only a second guard.
% \Xref{label}  : reference to the OTHER document (hyperlinked via xr-hyper)
% \Lref{label}  : reference in the SAME document (plain \ref)
\ifsupplement
  \externaldocument[M-]{main}[main.pdf]
  \newcommand{\Xref}[1]{\ref{M-#1}}
\else
  \externaldocument[S-]{supplement}[supplement.pdf]
  \newcommand{\Xref}[1]{\ref{S-#1}}
\fi
\newcommand{\Lref}[1]{\ref{#1}}
% \Supp{pE}{lem-chart} -> "Lemma E.3" style text is produced by the drafter:
%   "Lemma~\Xref{pE:lem-chart} of the supplement".  \SuppPart{E} gives "Part E".
\newcommand{\SuppPart}[1]{Part~#1 of the supplement}

% ---- numbers (generated by unit U0; macros only, no literal certified numbers) ----
\IfFileExists{numbers/numbers_paper.tex}%
  {% numbers_paper.tex -- GENERATED from ledger.json by build_ledger.py; do not edit.
% Rounding: upper bounds up, lower bounds and margins down at the printed digit; margin = threshold / bound.
% Use as $\times\NName$ for margins.  kind: cert certified, est estimate, hand hypothesis constant, doc only a document quotes it.
% NUEhatZero [cert, upper] U_cert_v35.json :: supersolution.EA.0
\newcommand{\NUEhatZero}{\ensuremath{1.2018\cdot10^{-6}}}
% NUEhatOne [cert, upper] U_cert_v35.json :: supersolution.EA.1
\newcommand{\NUEhatOne}{\ensuremath{1.0441\cdot10^{-4}}}
% NUEhatTwo [cert, upper] U_cert_v35.json :: supersolution.EA.2
\newcommand{\NUEhatTwo}{\ensuremath{0.056694}}
% NUBhatZero [cert, upper] U_cert_v35.json :: supersolution.Bhat.0
\newcommand{\NUBhatZero}{\ensuremath{4.8735\cdot10^{-6}}}
% NUBhatOne [cert, upper] U_cert_v35.json :: supersolution.Bhat.1
\newcommand{\NUBhatOne}{\ensuremath{5.5933\cdot10^{-4}}}
% NUBhatTwo [cert, upper] U_cert_v35.json :: supersolution.Bhat.2
\newcommand{\NUBhatTwo}{\ensuremath{0.36675}}
% NUSupersolutionCZeroZero [cert, upper] U_cert_v35.json :: supersolution.C0.0
\newcommand{\NUSupersolutionCZeroZero}{\ensuremath{1.1340\cdot10^{-6}}}
% NUSupersolutionCZeroOne [cert, upper] U_cert_v35.json :: supersolution.C0.1
\newcommand{\NUSupersolutionCZeroOne}{\ensuremath{1.6623\cdot10^{-4}}}
% NUSupersolutionCOneZero [cert, upper] U_cert_v35.json :: supersolution.C1.0
\newcommand{\NUSupersolutionCOneZero}{\ensuremath{1.5241\cdot10^{-9}}}
% NUSupersolutionCOneOne [cert, upper] U_cert_v35.json :: supersolution.C1.1
\newcommand{\NUSupersolutionCOneOne}{\ensuremath{5.5255\cdot10^{-8}}}
% NUSupersolutionNTwoZero [cert, upper] U_cert_v35.json :: supersolution.N2.0
\newcommand{\NUSupersolutionNTwoZero}{\ensuremath{0.071472}}
% NUSupersolutionNTwoOne [cert, upper] U_cert_v35.json :: supersolution.N2.1
\newcommand{\NUSupersolutionNTwoOne}{\ensuremath{0.33012}}
% NUSupersolutionNTwoTwo [cert, upper] U_cert_v35.json :: supersolution.N2.2
\newcommand{\NUSupersolutionNTwoTwo}{\ensuremath{5.0406}}
% NUSupersolutionNTwoOverFourNTwoRefZero [cert, upper] U_cert_v35.json :: supersolution.N2_over_4N2ref.0
\newcommand{\NUSupersolutionNTwoOverFourNTwoRefZero}{\ensuremath{0.18}}
% NUSupersolutionNTwoOverFourNTwoRefOne [cert, upper] U_cert_v35.json :: supersolution.N2_over_4N2ref.1
\newcommand{\NUSupersolutionNTwoOverFourNTwoRefOne}{\ensuremath{0.42}}
% NUSupersolutionNTwoOverFourNTwoRefTwo [cert, upper] U_cert_v35.json :: supersolution.N2_over_4N2ref.2
\newcommand{\NUSupersolutionNTwoOverFourNTwoRefTwo}{\ensuremath{0.26}}
% NUSupersolutionMOneZHypCheckZonesIIISOneTwoBhat [cert, upper] U_cert_v35.json :: supersolution.M1z_hyp_check.zones_I_II.S1_2Bhat
\newcommand{\NUSupersolutionMOneZHypCheckZonesIIISOneTwoBhat}{\ensuremath{0.28585}}
% NUSupersolutionMOneZHypCheckZoneIIIBOnePlusPUOne [cert, upper] U_cert_v35.json :: supersolution.M1z_hyp_check.zone_III.b1_plus_pU1
\newcommand{\NUSupersolutionMOneZHypCheckZoneIIIBOnePlusPUOne}{\ensuremath{0.50786}}
% NUMarginLateMOne [cert, margin] U_cert_v35.json :: margins.late_M1
\newcommand{\NUMarginLateMOne}{\ensuremath{2.353}}
% NUMarginLateMZero [cert, margin] U_cert_v35.json :: margins.late_M0
\newcommand{\NUMarginLateMZero}{\ensuremath{4.892}}
% NUMarginsLateMZeroZ [cert, margin] U_cert_v35.json :: margins.late_M0z
\newcommand{\NUMarginsLateMZeroZ}{\ensuremath{12.99}}
% NUMarginPresentMOne [cert, margin] U_cert_v35.json :: margins.present_M1_min.value
\newcommand{\NUMarginPresentMOne}{\ensuremath{2.770}}
% NUMarginPresentMZero [cert, margin] U_cert_v35.json :: margins.present_M0_min.value
\newcommand{\NUMarginPresentMZero}{\ensuremath{2.225}}
% NUMarginsPresentMZeroZMinValue [cert, margin] U_cert_v35.json :: margins.present_M0z_min.value
\newcommand{\NUMarginsPresentMZeroZMinValue}{\ensuremath{14.37}}
% NULoopGainZeroTotal [cert, upper] U_cert_v35.json :: loop_gain.0.total
\newcommand{\NULoopGainZeroTotal}{\ensuremath{1.23\cdot10^{-7}}}
% NULoopGainOneTotal [cert, upper] U_cert_v35.json :: loop_gain.1.total
\newcommand{\NULoopGainOneTotal}{\ensuremath{0.037169}}
% NULoopGainTwoNeighbour [cert, upper] U_cert_v35.json :: loop_gain.2.neighbour
\newcommand{\NULoopGainTwoNeighbour}{\ensuremath{0.037018}}
% NULoopGainTwoThroughA [cert, upper] U_cert_v35.json :: loop_gain.2.through_A
\newcommand{\NULoopGainTwoThroughA}{\ensuremath{0.0020172}}
% NUfineBhatZero [cert, upper] U_cert_v35.json :: ufine.Bhat.0
\newcommand{\NUfineBhatZero}{\ensuremath{4.8723\cdot10^{-6}}}
% NUfineBhatOne [cert, upper] U_cert_v35.json :: ufine.Bhat.1
\newcommand{\NUfineBhatOne}{\ensuremath{2.6266\cdot10^{-4}}}
% NUfineBhatTwo [cert, upper] U_cert_v35.json :: ufine.Bhat.2
\newcommand{\NUfineBhatTwo}{\ensuremath{0.025842}}
% NUcdeltaU [cert, upper] U_cert_v35.json :: launch.c_dU
\newcommand{\NUcdeltaU}{\ensuremath{1.7187\cdot10^{-8}}}
% NULaunchWmin [cert, lower] U_cert_v35.json :: launch.wmin
\newcommand{\NULaunchWmin}{\ensuremath{0.73492}}
% NULaunchCAD [cert, upper] U_cert_v35.json :: launch.CA_D
\newcommand{\NULaunchCAD}{\ensuremath{2.2429}}
% NULaunchQD [cert, upper] U_cert_v35.json :: launch.qD
\newcommand{\NULaunchQD}{\ensuremath{0.018509}}
% NULaunchDeltaInZero [cert, upper] U_cert_v35.json :: launch.Delta_in.0
\newcommand{\NULaunchDeltaInZero}{\ensuremath{3.131\cdot10^{-12}}}
% NULaunchDeltaInOne [cert, upper] U_cert_v35.json :: launch.Delta_in.1
\newcommand{\NULaunchDeltaInOne}{\ensuremath{3.287\cdot10^{-12}}}
% NULaunchDeltaInTwo [cert, upper] U_cert_v35.json :: launch.Delta_in.2
\newcommand{\NULaunchDeltaInTwo}{\ensuremath{3.451\cdot10^{-12}}}
% NULaunchPhiZero [cert, upper] U_cert_v35.json :: launch.Phi.0
\newcommand{\NULaunchPhiZero}{\ensuremath{6.434\cdot10^{-9}}}
% NULaunchPhiOne [cert, upper] U_cert_v35.json :: launch.Phi.1
\newcommand{\NULaunchPhiOne}{\ensuremath{1.4217\cdot10^{-8}}}
% NULaunchPhiTwo [cert, upper] U_cert_v35.json :: launch.Phi.2
\newcommand{\NULaunchPhiTwo}{\ensuremath{3.8404\cdot10^{-8}}}
% NUGammaZero [cert, upper] U_cert_v35.json :: Gamma.0
\newcommand{\NUGammaZero}{\ensuremath{4.0000358}}
% NUGammaOne [cert, upper] U_cert_v35.json :: Gamma.1
\newcommand{\NUGammaOne}{\ensuremath{16.000322}}
% NUGammaTwo [cert, upper] U_cert_v35.json :: Gamma.2
\newcommand{\NUGammaTwo}{\ensuremath{64.00218}}
% NUKappa [cert, upper] U_cert_v35.json :: kappa
\newcommand{\NUKappa}{\ensuremath{8.0000358}}
% NUKappaTwo [cert, upper] U_cert_v35.json :: kappa2
\newcommand{\NUKappaTwo}{\ensuremath{64.00061}}
% NUGhatZeroZero [cert, upper] U_cert_v35.json :: Ghat.0.0
\newcommand{\NUGhatZeroZero}{\ensuremath{51322.6}}
% NUGhatOneOne [cert, upper] U_cert_v35.json :: Ghat.1.1
\newcommand{\NUGhatOneOne}{\ensuremath{33962.9}}
% NUGhatTwoTwo [cert, upper] U_cert_v35.json :: Ghat.2.2
\newcommand{\NUGhatTwoTwo}{\ensuremath{348341.1}}
% NUGhatOneZero [cert, upper] U_cert_v35.json :: Ghat.1.0
\newcommand{\NUGhatOneZero}{\ensuremath{4.938}}
% NUGhatTwoZero [cert, upper] U_cert_v35.json :: Ghat.2.0
\newcommand{\NUGhatTwoZero}{\ensuremath{27.984}}
% NUGhatTwoOne [cert, upper] U_cert_v35.json :: Ghat.2.1
\newcommand{\NUGhatTwoOne}{\ensuremath{458.53}}
% NUKZeroZero [cert, upper] U_cert_v35.json :: K.0.0
\newcommand{\NUKZeroZero}{\ensuremath{8034.21}}
% NUKZeroOne [cert, upper] U_cert_v35.json :: K.0.1
\newcommand{\NUKZeroOne}{\ensuremath{724.89}}
% NUKOneZero [cert, upper] U_cert_v35.json :: K.1.0
\newcommand{\NUKOneZero}{\ensuremath{0.08152}}
% NUKOneOne [cert, upper] U_cert_v35.json :: K.1.1
\newcommand{\NUKOneOne}{\ensuremath{6336.94}}
% NUKOneTwo [cert, upper] U_cert_v35.json :: K.1.2
\newcommand{\NUKOneTwo}{\ensuremath{545.59}}
% NUKOneThree [cert, upper] U_cert_v35.json :: K.1.3
\newcommand{\NUKOneThree}{\ensuremath{0.58700}}
% NUKTwoZero [cert, upper] U_cert_v35.json :: K.2.0
\newcommand{\NUKTwoZero}{\ensuremath{0.87186}}
% NUKTwoOne [cert, upper] U_cert_v35.json :: K.2.1
\newcommand{\NUKTwoOne}{\ensuremath{528.62}}
% NUKTwoTwo [cert, upper] U_cert_v35.json :: K.2.2
\newcommand{\NUKTwoTwo}{\ensuremath{23425.98}}
% NUKTwoThree [cert, upper] U_cert_v35.json :: K.2.3
\newcommand{\NUKTwoThree}{\ensuremath{555.62}}
% NUFootwOne [cert, upper] U_cert_v35.json :: footw.1
\newcommand{\NUFootwOne}{\ensuremath{2.6361}}
% NUFootwTwo [cert, upper] U_cert_v35.json :: footw.2
\newcommand{\NUFootwTwo}{\ensuremath{2663.3}}
% NUCThreeThree [cert, upper] U_cert_v35.json :: c33
\newcommand{\NUCThreeThree}{\ensuremath{108.02}}
% NUCTwoTwoThree [cert, upper] U_cert_v35.json :: c223
\newcommand{\NUCTwoTwoThree}{\ensuremath{7.3284}}
% NUCEta [cert, upper] U_cert_v35.json :: C_eta
\newcommand{\NUCEta}{\ensuremath{13.376}}
% NULMax [cert, upper] U_cert_v35.json :: L_max
\newcommand{\NULMax}{\ensuremath{3.5352}}
% NURRatio [cert, upper] U_cert_v35.json :: R_ratio
\newcommand{\NURRatio}{\ensuremath{11.593}}
% NUSweepsSigmaInCoverSup [cert, upper] U_cert_v35.json :: sweeps.Sigma_in.cover_sup
\newcommand{\NUSweepsSigmaInCoverSup}{\ensuremath{5.0892}}
% NUSweepsSigmaECoverSup [cert, upper] U_cert_v35.json :: sweeps.Sigma_E.cover_sup
\newcommand{\NUSweepsSigmaECoverSup}{\ensuremath{0.30846}}
% NUSweepsSXCoverSup [cert, upper] U_cert_v35.json :: sweeps.S_X.cover_sup
\newcommand{\NUSweepsSXCoverSup}{\ensuremath{1.5243}}
% NUSweepsSigmaXCoverSup [cert, upper] U_cert_v35.json :: sweeps.Sigma_X.cover_sup
\newcommand{\NUSweepsSigmaXCoverSup}{\ensuremath{3.0081}}
% NUSweepsSigmaRsTwoCoverSup [cert, upper] U_cert_v35.json :: sweeps.Sigma_rs.2.cover_sup
\newcommand{\NUSweepsSigmaRsTwoCoverSup}{\ensuremath{4.3350\cdot10^{-3}}}
% NUDataTermsLateDZero [cert, upper] U_cert_v35.json :: data_terms_late.D.0
\newcommand{\NUDataTermsLateDZero}{\ensuremath{1.7330\cdot10^{-6}}}
% NUDataTermsLateDOne [cert, upper] U_cert_v35.json :: data_terms_late.D.1
\newcommand{\NUDataTermsLateDOne}{\ensuremath{2.5353\cdot10^{-4}}}
% NUDataTermsLateSlopeShareCOneSigmaXOverDkZero [cert, upper] U_cert_v35.json :: data_terms_late.slope_share_C1SigmaX_over_Dk.0
\newcommand{\NUDataTermsLateSlopeShareCOneSigmaXOverDkZero}{\ensuremath{2.7\cdot10^{-3}}}
% NUDataTermsLateSlopeShareCOneSigmaXOverDkOne [cert, upper] U_cert_v35.json :: data_terms_late.slope_share_C1SigmaX_over_Dk.1
\newcommand{\NUDataTermsLateSlopeShareCOneSigmaXOverDkOne}{\ensuremath{6.6\cdot10^{-4}}}
% NURLtTwoMarginZero [cert, margin] U_cert_v35.json :: R_lt2_margin.0
\newcommand{\NURLtTwoMarginZero}{\ensuremath{4.432}}
% NURLtTwoMarginOne [cert, margin] U_cert_v35.json :: R_lt2_margin.1
\newcommand{\NURLtTwoMarginOne}{\ensuremath{6.697}}
% NUThetaTwoEight [cert, upper] U_cert_v35.json :: theta28
\newcommand{\NUThetaTwoEight}{\ensuremath{0.37011}}
% NUDFiveSix [cert, upper] U_cert_v35.json :: D56
\newcommand{\NUDFiveSix}{\ensuremath{1.0368}}
% NUDTwoEight [cert, upper] U_cert_v35.json :: D28
\newcommand{\NUDTwoEight}{\ensuremath{1.5876}}
% NUSThreeThreeFootOverBhatTwo [cert, upper] U_cert_v35.json :: S33_foot_over_Bhat2
\newcommand{\NUSThreeThreeFootOverBhatTwo}{\ensuremath{1.69\cdot10^{-3}}}
% NSAHJWKappaJ [cert, upper] A_interface.json :: HJW.kappa_J
\newcommand{\NSAHJWKappaJ}{\ensuremath{0.042351}}
% NSAGpZero [cert, lower] A_interface.json :: res.tight.gp.0
\newcommand{\NSAGpZero}{\ensuremath{0.24292}}
% NSAGpOne [cert, lower] A_interface.json :: res.tight.gp.1
\newcommand{\NSAGpOne}{\ensuremath{0.24106}}
% NSAGpTwo [cert, lower] A_interface.json :: res.tight.gp.2
\newcommand{\NSAGpTwo}{\ensuremath{0.11961}}
% NSAAminC [cert, lower] A_interface.json :: res.tight.amin_c
\newcommand{\NSAAminC}{\ensuremath{0.97623}}
% NSARhoJ [cert, upper] A_interface.json :: res.tight.rhoJ
\newcommand{\NSARhoJ}{\ensuremath{5.4995\cdot10^{-7}}}
% NSAKappaOTwoTwo [cert, upper] A_interface.json :: res.tight.kappa_o.2.2
\newcommand{\NSAKappaOTwoTwo}{\ensuremath{1.2416\cdot10^{-5}}}
% NSADsrcZero [cert, upper] A_interface.json :: res.tight.Dsrc.0
\newcommand{\NSADsrcZero}{\ensuremath{1.8092\cdot10^{-7}}}
% NSADsrcOne [cert, upper] A_interface.json :: res.tight.Dsrc.1
\newcommand{\NSADsrcOne}{\ensuremath{9.8276\cdot10^{-5}}}
% NSABZeroDWindowSlopeZero [cert, upper] A_interface.json :: res.tight.B0.D_window_slope.0
\newcommand{\NSABZeroDWindowSlopeZero}{\ensuremath{8.7399\cdot10^{-10}}}
% NSABZeroDWindowSlopeOne [cert, upper] A_interface.json :: res.tight.B0.D_window_slope.1
\newcommand{\NSABZeroDWindowSlopeOne}{\ensuremath{5.2599\cdot10^{-8}}}
% NSAFTwoZero [cert, upper] A_interface.json :: res.tight.F2.0
\newcommand{\NSAFTwoZero}{\ensuremath{5.6506\cdot10^{-13}}}
% NSAFTwoOne [cert, upper] A_interface.json :: res.tight.F2.1
\newcommand{\NSAFTwoOne}{\ensuremath{5.9320\cdot10^{-13}}}
% NSACoupOne [cert, upper] A_interface.json :: res.tight.coup.1
\newcommand{\NSACoupOne}{\ensuremath{1.1568\cdot10^{-12}}}
% NSACOneZeroZero [cert, upper] A_interface.json :: res.tight.C1_0.0
\newcommand{\NSACOneZeroZero}{\ensuremath{7.2060\cdot10^{-10}}}
% NSACOneZeroOne [cert, upper] A_interface.json :: res.tight.C1_0.1
\newcommand{\NSACOneZeroOne}{\ensuremath{4.3143\cdot10^{-8}}}
% NSABZeroSupCoreZero [cert, upper] A_interface.json :: res.tight.B0.sup_core.0
\newcommand{\NSABZeroSupCoreZero}{\ensuremath{4.23\cdot10^{-5}}}
% NSACoreFluxMarginTwoFLUXZeroOverKZero [cert, margin] A_interface.json :: res.tight.core.flux_margin.2.FLUX0_over_K0
\newcommand{\NSACoreFluxMarginTwoFLUXZeroOverKZero}{\ensuremath{3.396}}
% NSACoreFluxMarginTwoFLUXTwoOverKTwo [cert, margin] A_interface.json :: res.tight.core.flux_margin.2.FLUX2_over_K2
\newcommand{\NSACoreFluxMarginTwoFLUXTwoOverKTwo}{\ensuremath{6.012}}
% NSAMarginFluxZero [cert, margin] asm_design_N2ref4.json :: res.tight.flux.beta.2.margin_FLUX0
\newcommand{\NSAMarginFluxZero}{\ensuremath{3.365}}
% NSAMarginFluxTwo [cert, margin] asm_design_N2ref4.json :: res.tight.flux.beta.2.margin_FLUX2
\newcommand{\NSAMarginFluxTwo}{\ensuremath{5.955}}
% NSAfourSZero [cert, upper] asm_design_N2ref4.json :: res.tight.S0
\newcommand{\NSAfourSZero}{\ensuremath{7.1878\cdot10^{-3}}}
% NSAfourSOne [cert, upper] asm_design_N2ref4.json :: res.tight.S1
\newcommand{\NSAfourSOne}{\ensuremath{0.10430}}
% NSAMarginCoreSZeroV [cert, margin] asm_design_N2ref4.json :: res.tight.margins.S0v_vs_0.0153
\newcommand{\NSAMarginCoreSZeroV}{\ensuremath{2.116}}
% NSAMarginCoreSOneV [cert, margin] asm_design_N2ref4.json :: res.tight.margins.S1v_vs_0.6
\newcommand{\NSAMarginCoreSOneV}{\ensuremath{5.748}}
% NSAMarginCoreMOne [cert, margin] asm_design_N2ref4.json :: res.tight.margins.M1core
\newcommand{\NSAMarginCoreMOne}{\ensuremath{4.409}}
% NSAMarginCoreMZero [cert, margin] asm_design_N2ref4.json :: res.tight.margins.M0core
\newcommand{\NSAMarginCoreMZero}{\ensuremath{5.450}}
% NSAMarginGhat [cert, margin] asm_design_N2ref4.json :: res.tight.margins.Ghat22
\newcommand{\NSAMarginGhat}{\ensuremath{1.917}}
% NSAMarginRampMOne [cert, margin] asm_design_N2ref4.json :: res.tight.margins.ramp M1
\newcommand{\NSAMarginRampMOne}{\ensuremath{1.184}}
% NSAfourInfoZeroDWindowConst [cert, upper] asm_design_N2ref4.json :: res.tight.info.0.D_window_const
\newcommand{\NSAfourInfoZeroDWindowConst}{\ensuremath{1.7732\cdot10^{-7}}}
% NSAfourInfoOneDWindowConst [cert, upper] asm_design_N2ref4.json :: res.tight.info.1.D_window_const
\newcommand{\NSAfourInfoOneDWindowConst}{\ensuremath{9.8058\cdot10^{-5}}}
% NZDeltaRZero [cert, upper] zone3_interface.json :: DeltaR.0
\newcommand{\NZDeltaRZero}{\ensuremath{4.4986\cdot10^{-8}}}
% NZDeltaROne [cert, upper] zone3_interface.json :: DeltaR.1
\newcommand{\NZDeltaROne}{\ensuremath{6.7972\cdot10^{-6}}}
% NZDeltaRTwo [cert, upper] zone3_interface.json :: DeltaR.2
\newcommand{\NZDeltaRTwo}{\ensuremath{0.011480}}
% NZDeltaRBZero [cert, upper] zone3_interface.json :: DeltaR_b.0
\newcommand{\NZDeltaRBZero}{\ensuremath{2.66\cdot10^{-12}}}
% NZDeltaRBOne [cert, upper] zone3_interface.json :: DeltaR_b.1
\newcommand{\NZDeltaRBOne}{\ensuremath{3.26\cdot10^{-9}}}
% NZDeltaRBTwo [cert, upper] zone3_interface.json :: DeltaR_b.2
\newcommand{\NZDeltaRBTwo}{\ensuremath{4.35\cdot10^{-8}}}
% NZDeltaRDetailNSqZero [cert, upper] zone3_interface.json :: DeltaR_detail.N_sq.0
\newcommand{\NZDeltaRDetailNSqZero}{\ensuremath{8.28\cdot10^{-5}}}
% NZDeltaRDetailNSqOne [cert, upper] zone3_interface.json :: DeltaR_detail.N_sq.1
\newcommand{\NZDeltaRDetailNSqOne}{\ensuremath{0.10149}}
% NZDeltaRDetailNSqTwo [cert, upper] zone3_interface.json :: DeltaR_detail.N_sq.2
\newcommand{\NZDeltaRDetailNSqTwo}{\ensuremath{1.3567}}
% NZDeltaRDetailNSqThree [cert, upper] zone3_interface.json :: DeltaR_detail.N_sq.3
\newcommand{\NZDeltaRDetailNSqThree}{\ensuremath{25350}}
% NZDeltaRDetailPartsTwoPairsThree [cert, upper] zone3_interface.json :: DeltaR_detail.parts.2.pairs.3
\newcommand{\NZDeltaRDetailPartsTwoPairsThree}{\ensuremath{0.011298}}
% NZDeltaRDetailPartsTwoLFourFoot [cert, upper] zone3_interface.json :: DeltaR_detail.parts.2.L4_foot
\newcommand{\NZDeltaRDetailPartsTwoLFourFoot}{\ensuremath{1.6\cdot10^{-4}}}
% NZDeltaRDetailPartsTwoCubic [cert, upper] zone3_interface.json :: DeltaR_detail.parts.2.cubic
\newcommand{\NZDeltaRDetailPartsTwoCubic}{\ensuremath{9.3\cdot10^{-6}}}
% NZDeltaRDetailPartsTwoSThreeThreeFoot [cert, upper] zone3_interface.json :: DeltaR_detail.parts.2.S33_foot
\newcommand{\NZDeltaRDetailPartsTwoSThreeThreeFoot}{\ensuremath{3.7\cdot10^{-7}}}
% NZNIIIZero [cert, upper] zone3_interface.json :: n_III.0
\newcommand{\NZNIIIZero}{\ensuremath{0.071445}}
% NZNIIIOne [cert, upper] zone3_interface.json :: n_III.1
\newcommand{\NZNIIIOne}{\ensuremath{0.32707}}
% NZNIIITwo [cert, upper] zone3_interface.json :: n_III.2
\newcommand{\NZNIIITwo}{\ensuremath{3.0445}}
% NZTauIII [cert, upper] zone3_interface.json :: tau_III
\newcommand{\NZTauIII}{\ensuremath{4.0708}}
% NZoneIIIBoxZero [cert, upper] zone3_interface.json :: box_III_380.b0
\newcommand{\NZoneIIIBoxZero}{\ensuremath{0.10032}}
% NZoneIIIBoxOne [cert, upper] zone3_interface.json :: box_III_380.b1
\newcommand{\NZoneIIIBoxOne}{\ensuremath{0.41974}}
% NZoneIIIBoxZ [cert, upper] zone3_interface.json :: box_III_380.bz
\newcommand{\NZoneIIIBoxZ}{\ensuremath{0.014291}}
% NZBoxIIIThreeEightZeroMarginsBZero [cert, margin] zone3_interface.json :: box_III_380.margins.b0
\newcommand{\NZBoxIIIThreeEightZeroMarginsBZero}{\ensuremath{1.495}}
% NZoneIIIMarginThinnest [cert, margin] zone3_interface.json :: box_III_380.margins.b1
\newcommand{\NZoneIIIMarginThinnest}{\ensuremath{1.238}}
% NZBoxIIIThreeEightZeroMarginsBz [cert, margin] zone3_interface.json :: box_III_380.margins.bz
\newcommand{\NZBoxIIIThreeEightZeroMarginsBz}{\ensuremath{11.19}}
% NZBoxIIISevenZeroZeroBZero [cert, upper] zone3_interface.json :: box_III_700.b0
\newcommand{\NZBoxIIISevenZeroZeroBZero}{\ensuremath{0.22757}}
% NZBoxIIISevenZeroZeroBOne [cert, upper] zone3_interface.json :: box_III_700.b1
\newcommand{\NZBoxIIISevenZeroZeroBOne}{\ensuremath{0.54234}}
% NZUfineKTwoAtBhatRef [cert, upper] zone3_interface.json :: ufine.K2_at_Bhat_ref
\newcommand{\NZUfineKTwoAtBhatRef}{\ensuremath{0.046956}}
% NZUfineKThreeStripAtBhatRef [cert, upper] zone3_interface.json :: ufine.K3_strip_at_Bhat_ref
\newcommand{\NZUfineKThreeStripAtBhatRef}{\ensuremath{0.054140}}
% NRecoveryMinMargin [cert, margin] zone3_interface.json :: cone_family.H6.at_Bhat_ref.min_margin
\newcommand{\NRecoveryMinMargin}{\ensuremath{1.860}}
% NZConeFamilyHSixAtBhatRefConeSlopeMargin [cert, margin] zone3_interface.json :: cone_family.H6.at_Bhat_ref.cone.slope.margin
\newcommand{\NZConeFamilyHSixAtBhatRefConeSlopeMargin}{\ensuremath{2.46}}
% NZConeFamilyHSixAtBhatRefConeKTwoMargin [cert, margin] zone3_interface.json :: cone_family.H6.at_Bhat_ref.cone.K2.margin
\newcommand{\NZConeFamilyHSixAtBhatRefConeKTwoMargin}{\ensuremath{2.54}}
% NZConeFamilyHSixAtBhatRefConeKThreeMargin [cert, margin] zone3_interface.json :: cone_family.H6.at_Bhat_ref.cone.K3.margin
\newcommand{\NZConeFamilyHSixAtBhatRefConeKThreeMargin}{\ensuremath{2.51}}
% NConeTstarMargin [cert, margin] zone3_interface.json :: acceptance.tstar_margin
\newcommand{\NConeTstarMargin}{\ensuremath{2.30}}
% NUoneTFiveYaddLevelOneTwo [cert, upper] u1_ghat.json :: t5_Yadd.level1.2
\newcommand{\NUoneTFiveYaddLevelOneTwo}{\ensuremath{0.0044}}
% NUoneTFiveYaddLevelTwoTwo [cert, upper] u1_ghat.json :: t5_Yadd.level2.2
\newcommand{\NUoneTFiveYaddLevelTwoTwo}{\ensuremath{0.1081}}
% NUoneTFiveYaddLevelTwoThree [cert, upper] u1_ghat.json :: t5_Yadd.level2.3
\newcommand{\NUoneTFiveYaddLevelTwoThree}{\ensuremath{0.374}}
% NUoneWeightThetaGammaZero [cert, upper] u1_ghat.json :: weight.Theta.Gamma.0
\newcommand{\NUoneWeightThetaGammaZero}{\ensuremath{4.00000000000000000000054}}
% NRKZeroZero [cert, upper] ramp_U.json :: K.0.0
\newcommand{\NRKZeroZero}{\ensuremath{13436.03}}
% NRKOneOne [cert, upper] ramp_U.json :: K.1.1
\newcommand{\NRKOneOne}{\ensuremath{10526.19}}
% NRKTwoTwo [cert, upper] ramp_U.json :: K.2.2
\newcommand{\NRKTwoTwo}{\ensuremath{39482.30}}
% NUoneCasesPrimaryYNextZero [cert, upper] u1_ghat.json :: cases.primary.Y_next.0
\newcommand{\NUoneCasesPrimaryYNextZero}{\ensuremath{457.092}}
% NUoneCasesPrimaryYNextOne [cert, upper] u1_ghat.json :: cases.primary.Y_next.1
\newcommand{\NUoneCasesPrimaryYNextOne}{\ensuremath{356.940}}
% NUoneCasesPrimaryYNextTwo [cert, upper] u1_ghat.json :: cases.primary.Y_next.2
\newcommand{\NUoneCasesPrimaryYNextTwo}{\ensuremath{1320.220}}
% NSfourResultsMNineHTwoTightenedFUOneRowsZeroBound [cert, upper] FINAL_manifest_v33.json :: results.M9_H2_tightened_fU1.rows.0.bound
\newcommand{\NSfourResultsMNineHTwoTightenedFUOneRowsZeroBound}{\ensuremath{3.658\cdot10^{-10}}}
% NSfourResultsMNineHTwoTightenedFUOneRowsZeroMargin [cert, margin] FINAL_manifest_v33.json :: results.M9_H2_tightened_fU1.rows.0.margin
\newcommand{\NSfourResultsMNineHTwoTightenedFUOneRowsZeroMargin}{\ensuremath{27.34}}
% NSfourResultsMNineHTwoTightenedFUOneRowsOneBound [cert, upper] FINAL_manifest_v33.json :: results.M9_H2_tightened_fU1.rows.1.bound
\newcommand{\NSfourResultsMNineHTwoTightenedFUOneRowsOneBound}{\ensuremath{1.777\cdot10^{-9}}}
% NSfourResultsMNineHTwoTightenedFUOneRowsOneMargin [cert, margin] FINAL_manifest_v33.json :: results.M9_H2_tightened_fU1.rows.1.margin
\newcommand{\NSfourResultsMNineHTwoTightenedFUOneRowsOneMargin}{\ensuremath{562.8}}
% NSfourResultsMNineHTwoTightenedFUOneRowsTwoBound [cert, upper] FINAL_manifest_v33.json :: results.M9_H2_tightened_fU1.rows.2.bound
\newcommand{\NSfourResultsMNineHTwoTightenedFUOneRowsTwoBound}{\ensuremath{1.665\cdot10^{-6}}}
% NSfourResultsMNineHTwoTightenedFUOneRowsTwoMargin [cert, margin] FINAL_manifest_v33.json :: results.M9_H2_tightened_fU1.rows.2.margin
\newcommand{\NSfourResultsMNineHTwoTightenedFUOneRowsTwoMargin}{\ensuremath{600.9}}
% NHOneAlphaThr [hand, exact] hand :: 
\newcommand{\NHOneAlphaThr}{\ensuremath{1.5\cdot10^{-4}}}
% NHOneSqrtLamThr [hand, exact] hand :: 
\newcommand{\NHOneSqrtLamThr}{\ensuremath{5\cdot10^{-8}}}
% NYbarStep [hand, exact] hand :: 
\newcommand{\NYbarStep}{\ensuremath{1\cdot10^{-12}}}
% NHTwoThrZero [hand, exact] hand :: 
\newcommand{\NHTwoThrZero}{\ensuremath{1\cdot10^{-8}}}
% NHTwoThrOne [hand, exact] hand :: 
\newcommand{\NHTwoThrOne}{\ensuremath{1\cdot10^{-6}}}
% NHTwoThrTwo [hand, exact] hand :: 
\newcommand{\NHTwoThrTwo}{\ensuremath{1\cdot10^{-3}}}
% NHThreeThrZero [hand, exact] hand :: 
\newcommand{\NHThreeThrZero}{\ensuremath{1\cdot10^{-6}}}
% NHThreeThrOne [hand, exact] hand :: 
\newcommand{\NHThreeThrOne}{\ensuremath{5\cdot10^{-6}}}
% NHThreeThrTwo [hand, exact] hand :: 
\newcommand{\NHThreeThrTwo}{\ensuremath{1\cdot10^{-5}}}
% NHFourThrZero [hand, exact] hand :: 
\newcommand{\NHFourThrZero}{\ensuremath{2\cdot10^{-7}}}
% NHFourThrOne [hand, exact] hand :: 
\newcommand{\NHFourThrOne}{\ensuremath{1\cdot10^{-5}}}
% NHFourThrTwo [hand, exact] hand :: 
\newcommand{\NHFourThrTwo}{\ensuremath{1\cdot10^{-5}}}
% NHFourThrThree [hand, exact] hand :: 
\newcommand{\NHFourThrThree}{\ensuremath{1\cdot10^{-5}}}
% NHFourKfourThr [hand, exact] hand :: 
\newcommand{\NHFourKfourThr}{\ensuremath{0.06}}
% NHSixConeSlopeThr [hand, exact] hand :: 
\newcommand{\NHSixConeSlopeThr}{\ensuremath{0.08}}
% NHSixConeKtwoThr [hand, exact] hand :: 
\newcommand{\NHSixConeKtwoThr}{\ensuremath{0.03}}
% NHSixConeNthreeThr [hand, exact] hand :: 
\newcommand{\NHSixConeNthreeThr}{\ensuremath{0.004}}
% NHSixConeJetThr [hand, exact] hand :: 
\newcommand{\NHSixConeJetThr}{\ensuremath{0.12}}
% NHSixConeHFiveThr [hand, exact] hand :: 
\newcommand{\NHSixConeHFiveThr}{\ensuremath{0.0144}}
% NHSixFarHFiveThr [hand, exact] hand :: 
\newcommand{\NHSixFarHFiveThr}{\ensuremath{0.0225}}
% NLemmaOConeM [hand, exact] hand :: 
\newcommand{\NLemmaOConeM}{\ensuremath{0.12}}
% NLemmaOConeK [hand, exact] hand :: 
\newcommand{\NLemmaOConeK}{\ensuremath{0.18}}
% NLemmaOFarM [hand, exact] hand :: 
\newcommand{\NLemmaOFarM}{\ensuremath{0.15}}
% NLemmaOFarK [hand, exact] hand :: 
\newcommand{\NLemmaOFarK}{\ensuremath{0.3}}
% NCoercivityConst [hand, exact] hand :: 
\newcommand{\NCoercivityConst}{\ensuremath{0.046875}}
% NMOnez [hand, exact] hand :: 
\newcommand{\NMOnez}{\ensuremath{0.534}}
% NBetaA [hand, exact] hand :: 
\newcommand{\NBetaA}{\ensuremath{2}}
% NTOneLower [cert, lower] t1clock.json :: t1.0
\newcommand{\NTOneLower}{\ensuremath{4.5978574775}}
% NTOneUpper [cert, upper] t1clock.json :: t1.1
\newcommand{\NTOneUpper}{\ensuremath{4.5978586144}}
% NLambdaOneLower [cert, lower] t1clock.json :: lnlam.0
\newcommand{\NLambdaOneLower}{\ensuremath{6.49\cdot10^{-5}}}
% NLambdaOneUpper [cert, upper] t1clock.json :: lnlam.1
\newcommand{\NLambdaOneUpper}{\ensuremath{6.75\cdot10^{-5}}}
% NTMinusTOneUpper [cert, upper] t1clock.json :: lnlam.1; 1.5e-4.sup_ab
\newcommand{\NTMinusTOneUpper}{\ensuremath{2.1\cdot10^{-17}}}
% NAbarUpper [cert, upper] final_alpha_points_5e-8.json :: 1.5e-4.sup_ab
\newcommand{\NAbarUpper}{\ensuremath{2.19103\cdot10^{-7}}}
% NBetabarUpper [cert, upper] final_alpha_points_5e-8.json :: 1.5e-4.sup_beta
\newcommand{\NBetabarUpper}{\ensuremath{3.33811\cdot10^{-9}}}
% NBaseKmaxLower [cert, lower] t1clock.json :: base_kmax.0
\newcommand{\NBaseKmaxLower}{\ensuremath{2.279\cdot10^{4}}}
% NBaseKmaxUpper [cert, upper] t1clock.json :: base_kmax.1
\newcommand{\NBaseKmaxUpper}{\ensuremath{2.376\cdot10^{4}}}
% NHOneAlphaBase [cert, upper] summary.json :: alpha_abs_max
\newcommand{\NHOneAlphaBase}{\ensuremath{1.1739\cdot10^{-4}}}
% NHOneAlphaErr [cert, upper] t9_h1_clam_s180.json :: C_alpha; inputs.YSTAR
\newcommand{\NHOneAlphaErr}{\ensuremath{7.21\cdot10^{-13}}}
% NHOneAlphaMargin [cert, margin] summary.json :: alpha_abs_max; C_alpha; inputs.YSTAR
\newcommand{\NHOneAlphaMargin}{\ensuremath{1.27}}
% NHOneSqrtLamBase [cert, upper] t9_h1_clam_s180.json :: sqrtLam_base_Wp
\newcommand{\NHOneSqrtLamBase}{\ensuremath{2.282\cdot10^{-8}}}
% NHOneSqrtLamTotal [cert, upper] t9_h1_clam_s180.json :: sqrtLam_total_Wp
\newcommand{\NHOneSqrtLamTotal}{\ensuremath{2.29\cdot10^{-8}}}
% NHOneSqrtLamMargin [cert, margin] t9_h1_clam_s180.json :: sqrtLam_total_Wp
\newcommand{\NHOneSqrtLamMargin}{\ensuremath{2.18}}
% NHOneYbarLimit [cert, lower] t9_h1_clam_s180.json :: Ybar_limit_H1
\newcommand{\NHOneYbarLimit}{\ensuremath{1.12\cdot10^{-9}}}
% NSwitchYUpper [cert, upper] switch_s2.json :: modulated.Y_raw
\newcommand{\NSwitchYUpper}{\ensuremath{6.21\cdot10^{-13}}}
% NSwitchYMargin [cert, margin] switch_s2.json :: modulated.Y_raw
\newcommand{\NSwitchYMargin}{\ensuremath{1.6}}
% NMinCI [cert, lower] constants.json :: meta.summary.min_c_I
\newcommand{\NMinCI}{\ensuremath{0.05283}}
% NNumCertificates [cert, exact] constants.json :: meta.summary.n_cert
\newcommand{\NNumCertificates}{\ensuremath{26}}
% NMinRhoC [cert, lower] constants.json :: meta.closure.rho_c_min
\newcommand{\NMinRhoC}{\ensuremath{0.04711}}
% NHTwoBaseZero [cert, upper] FINAL_manifest_v33.json :: results.M9_H2_tightened.rows.0.base
\newcommand{\NHTwoBaseZero}{\ensuremath{8.011\cdot10^{-11}}}
% NHTwoBaseOne [cert, upper] FINAL_manifest_v33.json :: results.M9_H2_tightened.rows.1.base
\newcommand{\NHTwoBaseOne}{\ensuremath{1.336\cdot10^{-10}}}
% NHTwoBaseTwo [cert, upper] FINAL_manifest_v33.json :: results.M9_H2_tightened.rows.2.base
\newcommand{\NHTwoBaseTwo}{\ensuremath{2.936\cdot10^{-10}}}
% NHTwoErrWZero [cert, upper] FINAL_manifest_v33.json :: results.M9_H2_tightened.rows.0.with_error
\newcommand{\NHTwoErrWZero}{\ensuremath{8.014\cdot10^{-11}}}
% NHTwoMarginWZero [cert, margin] FINAL_manifest_v33.json :: results.M9_H2_tightened.rows.0.margin
\newcommand{\NHTwoMarginWZero}{\ensuremath{124.78}}
% NHTwoErrWOne [cert, upper] FINAL_manifest_v33.json :: results.M9_H2_tightened.rows.1.with_error
\newcommand{\NHTwoErrWOne}{\ensuremath{4.706\cdot10^{-10}}}
% NHTwoMarginWOne [cert, margin] FINAL_manifest_v33.json :: results.M9_H2_tightened.rows.1.margin
\newcommand{\NHTwoMarginWOne}{\ensuremath{2124.9}}
% NHTwoErrWTwo [cert, upper] FINAL_manifest_v33.json :: results.M9_H2_tightened.rows.2.with_error
\newcommand{\NHTwoErrWTwo}{\ensuremath{1.559\cdot10^{-6}}}
% NHTwoMarginWTwo [cert, margin] FINAL_manifest_v33.json :: results.M9_H2_tightened.rows.2.margin
\newcommand{\NHTwoMarginWTwo}{\ensuremath{641.53}}
% NHThreeBaseZero [cert, upper] FINAL_manifest_v33.json :: results.M9_H2f.rows.0.base
\newcommand{\NHThreeBaseZero}{\ensuremath{3.127\cdot10^{-7}}}
% NHThreeBaseOne [cert, upper] FINAL_manifest_v33.json :: results.M9_H2f.rows.1.base
\newcommand{\NHThreeBaseOne}{\ensuremath{1.004\cdot10^{-6}}}
% NHThreeBaseTwo [cert, upper] FINAL_manifest_v33.json :: results.M9_H2f.rows.2.base
\newcommand{\NHThreeBaseTwo}{\ensuremath{2.431\cdot10^{-6}}}
% NHThreeErrZero [cert, upper] FINAL_manifest_v33.json :: results.M9_H2f.rows.0.with_error
\newcommand{\NHThreeErrZero}{\ensuremath{3.1363\cdot10^{-7}}}
% NHThreeMarginZero [cert, margin] FINAL_manifest_v33.json :: results.M9_H2f.rows.0.margin
\newcommand{\NHThreeMarginZero}{\ensuremath{3.1885}}
% NHThreeErrOne [cert, upper] FINAL_manifest_v33.json :: results.M9_H2f.rows.1.with_error
\newcommand{\NHThreeErrOne}{\ensuremath{1.0034\cdot10^{-6}}}
% NHThreeMarginOne [cert, margin] FINAL_manifest_v33.json :: results.M9_H2f.rows.1.margin
\newcommand{\NHThreeMarginOne}{\ensuremath{4.9834}}
% NHThreeErrTwo [cert, upper] FINAL_manifest_v33.json :: results.M9_H2f.rows.2.with_error
\newcommand{\NHThreeErrTwo}{\ensuremath{2.4351\cdot10^{-6}}}
% NHThreeMarginTwo [cert, margin] FINAL_manifest_v33.json :: results.M9_H2f.rows.2.margin
\newcommand{\NHThreeMarginTwo}{\ensuremath{4.1067}}
% NHFourBaseZero [cert, upper] FINAL_manifest_v33.json :: results.M9_base_sums.combined.windows_max.0
\newcommand{\NHFourBaseZero}{\ensuremath{3.48\cdot10^{-8}}}
% NHFourErrZero [cert, upper] FINAL_manifest_v33.json :: results.M9_H3_H2v32.rows.win.0.with_error_arb
\newcommand{\NHFourErrZero}{\ensuremath{3.5087\cdot10^{-8}}}
% NHFourMarginZero [cert, margin] FINAL_manifest_v33.json :: results.M9_H3_H2v32.rows.win.0.margin_arb
\newcommand{\NHFourMarginZero}{\ensuremath{5.700}}
% NHFourBaseOne [cert, upper] FINAL_manifest_v33.json :: results.M9_base_sums.combined.windows_max.1
\newcommand{\NHFourBaseOne}{\ensuremath{1.27\cdot10^{-7}}}
% NHFourErrOne [cert, upper] FINAL_manifest_v33.json :: results.M9_H3_H2v32.rows.win.1.with_error_arb
\newcommand{\NHFourErrOne}{\ensuremath{1.2616\cdot10^{-7}}}
% NHFourMarginOne [cert, margin] FINAL_manifest_v33.json :: results.M9_H3_H2v32.rows.win.1.margin_arb
\newcommand{\NHFourMarginOne}{\ensuremath{79.26}}
% NHFourBaseTwo [cert, upper] FINAL_manifest_v33.json :: results.M9_base_sums.combined.windows_max.2
\newcommand{\NHFourBaseTwo}{\ensuremath{3.76\cdot10^{-7}}}
% NHFourErrTwo [cert, upper] FINAL_manifest_v33.json :: results.M9_H3_H2v32.rows.win.2.with_error_arb
\newcommand{\NHFourErrTwo}{\ensuremath{3.7724\cdot10^{-7}}}
% NHFourMarginTwo [cert, margin] FINAL_manifest_v33.json :: results.M9_H3_H2v32.rows.win.2.margin_arb
\newcommand{\NHFourMarginTwo}{\ensuremath{26.50}}
% NHFourBaseThree [cert, upper] FINAL_manifest_v33.json :: results.M9_base_sums.combined.windows_max.3
\newcommand{\NHFourBaseThree}{\ensuremath{9.71\cdot10^{-7}}}
% NHFourErrThree [cert, upper] FINAL_manifest_v33.json :: results.M9_H3_H2v32.rows.win.3.with_error_arb
\newcommand{\NHFourErrThree}{\ensuremath{1.2822\cdot10^{-6}}}
% NHFourMarginThree [cert, margin] FINAL_manifest_v33.json :: results.M9_H3_H2v32.rows.win.3.margin_arb
\newcommand{\NHFourMarginThree}{\ensuremath{7.799}}
% NHFourYbarLimit [cert, lower] FINAL_manifest_v33.json :: results.M9_H3_H2v32.binding_Ybar_limit
\newcommand{\NHFourYbarLimit}{\ensuremath{1.575\cdot10^{-11}}}
% NHFourYbarMargin [cert, margin] FINAL_manifest_v33.json :: results.M9_H3_H2v32.binding_Ybar_limit
\newcommand{\NHFourYbarMargin}{\ensuremath{15.75}}
% NHFourKfourBase [cert, upper] amp_A1_w180.json :: inputs.BASE_T1.K4
\newcommand{\NHFourKfourBase}{\ensuremath{3.1633\cdot10^{-4}}}
% NHSixConeSlope [cert, upper] FINAL_manifest_v33.json :: results.cone_fixedL.checks.slope.value
\newcommand{\NHSixConeSlope}{\ensuremath{0.03029}}
% NHSixConeSlopeMargin [cert, margin] FINAL_manifest_v33.json :: results.cone_fixedL.checks.slope.margin
\newcommand{\NHSixConeSlopeMargin}{\ensuremath{2.641}}
% NHSixConeKtwo [cert, upper] FINAL_manifest_v33.json :: results.cone_fixedL.checks.K2.value
\newcommand{\NHSixConeKtwo}{\ensuremath{0.009876}}
% NHSixConeKtwoMargin [cert, margin] FINAL_manifest_v33.json :: results.cone_fixedL.checks.K2.margin
\newcommand{\NHSixConeKtwoMargin}{\ensuremath{3.037}}
% NHSixConeNthree [cert, upper] FINAL_manifest_v33.json :: results.cone_fixedL.checks.N3.value
\newcommand{\NHSixConeNthree}{\ensuremath{1.461\cdot10^{-4}}}
% NHSixConeNthreeMargin [cert, margin] FINAL_manifest_v33.json :: results.cone_fixedL.checks.N3.margin
\newcommand{\NHSixConeNthreeMargin}{\ensuremath{27.38}}
% NHSixConeHFive [cert, upper] FINAL_manifest_v33.json :: results.H4_fixedL.total
\newcommand{\NHSixConeHFive}{\ensuremath{1.2865\cdot10^{-5}}}
% NHSixConeHFiveMargin [cert, margin] FINAL_manifest_v33.json :: results.H4_fixedL.margin
\newcommand{\NHSixConeHFiveMargin}{\ensuremath{1119}}
% NHSixConeTstarNeeded [cert, upper] FINAL_manifest_v33.json :: results.H7_fixedL.ratio_upper
\newcommand{\NHSixConeTstarNeeded}{\ensuremath{1.6000015\cdot10^{-7}}}
% NHSixConeTstarLower [cert, lower] FINAL_manifest_v33.json :: results.H7_fixedL.t_star_lower
\newcommand{\NHSixConeTstarLower}{\ensuremath{3.68334\cdot10^{-7}}}
% NHSixConeTstarMargin [cert, margin] FINAL_manifest_v33.json :: results.H7_fixedL.margin_tstar_over_ratio
\newcommand{\NHSixConeTstarMargin}{\ensuremath{2.302}}
% NHSixConeHFiveRatioPeak [cert, upper] h4loc_A1_w180.json :: windows.constants/amp_A1_w180.json.worst.ratio
\newcommand{\NHSixConeHFiveRatioPeak}{\ensuremath{0.2393}}
% NHSixConeHOneSup [cert, upper] FINAL_manifest_v33.json :: results.H1_fixedL.sup_h0_max
\newcommand{\NHSixConeHOneSup}{\ensuremath{0.2163}}
% NHSixConeHOneH [cert, upper] FINAL_manifest_v33.json :: results.H1_fixedL.H_max
\newcommand{\NHSixConeHOneH}{\ensuremath{0.2185}}
% NHSixFarHFiveRatio [cert, upper] h4loc_v32.json :: windows.constants/amp_A1_w180.json.cylinders.18.ratio
\newcommand{\NHSixFarHFiveRatio}{\ensuremath{0.1761}}
% NHSixFarHFiveMarginOld [cert, margin] h4loc_v32.json :: windows.constants/amp_A1_w180.json.cylinders.18.ratio
\newcommand{\NHSixFarHFiveMarginOld}{\ensuremath{5.68}}
% NHSixFarHFiveMargin [cert, margin] FINAL_manifest_v33.json :: results.M9_H4_far.min_margin
\newcommand{\NHSixFarHFiveMargin}{\ensuremath{15.58}}
% NHSixConeHFiveMarginOld [cert, margin] h4loc_A1_w180.json :: windows.constants/amp_A1_w180.json.worst.ratio
\newcommand{\NHSixConeHFiveMarginOld}{\ensuremath{4.18}}
% NHSixFarKtwo [cert, upper] t9_far_v32.json :: h2
\newcommand{\NHSixFarKtwo}{\ensuremath{1.797\cdot10^{-4}}}
% NHSixFarKtwoThr [hand, exact] t9_far_v32.json :: thr.h2
\newcommand{\NHSixFarKtwoThr}{\ensuremath{3.53\cdot10^{-4}}}
% NHSixFarKtwoMargin [cert, margin] t9_far_v32.json :: margins.h2
\newcommand{\NHSixFarKtwoMargin}{\ensuremath{1.964}}
% NHSixFarKthree [cert, upper] t9_far_v32.json :: h3
\newcommand{\NHSixFarKthree}{\ensuremath{2.595\cdot10^{-8}}}
% NHSixFarKthreeThr [hand, exact] t9_far_v32.json :: thr.h3
\newcommand{\NHSixFarKthreeThr}{\ensuremath{1.17\cdot10^{-5}}}
% NHSixFarKthreeMargin [cert, margin] t9_far_v32.json :: margins.h3
\newcommand{\NHSixFarKthreeMargin}{\ensuremath{450.9}}
% NHSixFarKfour [cert, upper] t9_far_v32.json :: h4
\newcommand{\NHSixFarKfour}{\ensuremath{6.964\cdot10^{-9}}}
% NHSixFarKfourThr [hand, exact] t9_far_v32.json :: thr.h4
\newcommand{\NHSixFarKfourThr}{\ensuremath{1.03\cdot10^{-7}}}
% NHSixFarKfourMargin [cert, margin] t9_far_v32.json :: margins.h4
\newcommand{\NHSixFarKfourMargin}{\ensuremath{14.79}}
% NHSixFarSlope [cert, upper] t9_far_v32.json :: slope
\newcommand{\NHSixFarSlope}{\ensuremath{0.04153}}
% NHSixFarSlopeThr [hand, exact] t9_far_v32.json :: thr.slope
\newcommand{\NHSixFarSlopeThr}{\ensuremath{0.0800}}
% NHSixFarSlopeMargin [cert, margin] t9_far_v32.json :: margins.slope
\newcommand{\NHSixFarSlopeMargin}{\ensuremath{1.926}}
% NHSixFarNthree [cert, upper] t9_far_v32.json :: L2
\newcommand{\NHSixFarNthree}{\ensuremath{5.516\cdot10^{-7}}}
% NHSixFarNthreeThr [hand, exact] t9_far_v32.json :: thr.L2
\newcommand{\NHSixFarNthreeThr}{\ensuremath{3.32\cdot10^{-6}}}
% NHSixFarNthreeMargin [cert, margin] t9_far_v32.json :: margins.L2
\newcommand{\NHSixFarNthreeMargin}{\ensuremath{6.019}}
% NFarTstarMargin [cert, margin] lemmaO_v32_consts.json :: far_v32.case.margin_factor_lower
\newcommand{\NFarTstarMargin}{\ensuremath{15.5}}
% NFarTstarLower [cert, lower] lemmaO_v32_consts.json :: far_v32.case.t_star_lower
\newcommand{\NFarTstarLower}{\ensuremath{9.5108\cdot10^{-8}}}
% NFarTstarNeeded [cert, upper] lemmaO_v32_consts.json :: far_v32.case.needed
\newcommand{\NFarTstarNeeded}{\ensuremath{6.134\cdot10^{-9}}}
% NGammaModes [cert, exact] gamma0_final.json :: M
\newcommand{\NGammaModes}{\ensuremath{3762}}
% NGammaHalfwidth [cert, upper] gamma0_final.json :: b_j_halfwidth
\newcommand{\NGammaHalfwidth}{\ensuremath{8.138\cdot10^{-24}}}
% NGammaContraction [cert, upper] gamma0_final.json :: krawczyk_contraction
\newcommand{\NGammaContraction}{\ensuremath{3.01\cdot10^{-12}}}
% NGammaArcs [cert, exact] gamma0_final.json :: embedded.arcs
\newcommand{\NGammaArcs}{\ensuremath{701}}
% NGammaMaxTurn [cert, upper] gamma0_final.json :: embedded.maxturn
\newcommand{\NGammaMaxTurn}{\ensuremath{0.04997}}
% NGammaMaxPair [cert, upper] gamma0_final.json :: embedded.maxpair
\newcommand{\NGammaMaxPair}{\ensuremath{0.09597}}
% NGammaMinSep [cert, lower] gamma0_final.json :: embedded.minsep
\newcommand{\NGammaMinSep}{\ensuremath{8.89\cdot10^{-3}}}
% NGammaHalfGap [cert, lower] gamma0_final.json :: embedded.gap_cert
\newcommand{\NGammaHalfGap}{\ensuremath{0.0248219}}
% NPhaseISlabs [cert, exact] status.json :: slab
\newcommand{\NPhaseISlabs}{\ensuremath{96}}
% NPhaseIModes [cert, exact] status.json :: last_slab.M
\newcommand{\NPhaseIModes}{\ensuremath{2800}}
% NPhaseIKmaxLower [cert, lower] status.json :: kmax_enclosure.0
\newcommand{\NPhaseIKmaxLower}{\ensuremath{29.352874545}}
% NPhaseIKmaxUpper [cert, upper] status.json :: kmax_enclosure.1
\newcommand{\NPhaseIKmaxUpper}{\ensuremath{29.352874546}}
% NPhaseITLower [cert, lower] status.json :: t
\newcommand{\NPhaseITLower}{\ensuremath{4.597850539485830}}
% NPhaseITUpper [cert, upper] status.json :: t
\newcommand{\NPhaseITUpper}{\ensuremath{4.597850539485831}}
% NPeakSqrtLam [cert, upper] rate_table_fin12f4_a0.00015.json :: sqrtL_max
\newcommand{\NPeakSqrtLam}{\ensuremath{5.308\cdot10^{-8}}}
% NTailRate [cert, lower] rate_table_fin12f4_a0.00015.json :: tail.rho
\newcommand{\NTailRate}{\ensuremath{2.693\cdot10^{-3}}}
% NFloorSqrtLam [cert, upper] rate_table_fin12f4_a0.00015.json :: tail.sqrt_LF
\newcommand{\NFloorSqrtLam}{\ensuremath{9.872\cdot10^{-10}}}
% NAlphaFloor [cert, upper] rate_table_fin12f4_a0.00015.json :: tail.alpha_F
\newcommand{\NAlphaFloor}{\ensuremath{6.867\cdot10^{-7}}}
% NSupBetaVerifier [cert, upper] final_alpha_points_5e-8.json :: 1.5e-4.sup_beta
\newcommand{\NSupBetaVerifier}{\ensuremath{3.33811\cdot10^{-9}}}
% NMOneRamp [cert, upper] final_alpha_points_5e-8.json :: 1.5e-4.M1
\newcommand{\NMOneRamp}{\ensuremath{0.4328}}
% NSITwoPrime [cert, upper] FINAL_manifest_v33.json :: results.F12_sliver.sI2_prime
\newcommand{\NSITwoPrime}{\ensuremath{0.015457}}
% NLoopGainCoarse [cert, upper] U_cert_v35.json :: loop_gain.2.total
\newcommand{\NLoopGainCoarse}{\ensuremath{0.03904}}
% NRateMarginCoarse [cert, margin] U_cert_v35.json :: rate_margin
\newcommand{\NRateMarginCoarse}{\ensuremath{28.2}}
% NSOneEpsLower [cert, lower] bracket2.json :: eps_a
\newcommand{\NSOneEpsLower}{\ensuremath{1.5346198892965\cdot10^{-4}}}
% NSOneEpsUpper [cert, upper] bracket2.json :: eps_b
\newcommand{\NSOneEpsUpper}{\ensuremath{1.5346198899299\cdot10^{-4}}}
% NSOneRowsTip [cert, exact] ext_tip_check.json :: rows
\newcommand{\NSOneRowsTip}{\ensuremath{51246}}
% NSOneTipWidth [cert, upper] ext_tip_check.json :: max_pointwise_width.th
\newcommand{\NSOneTipWidth}{\ensuremath{3.2\cdot10^{-12}}}
% NSOneExtWidth [cert, upper] ext813_check.json :: max_pointwise_width.p
\newcommand{\NSOneExtWidth}{\ensuremath{5.7\cdot10^{-10}}}
% NSOneRowsMC [cert, exact] check_outA5.json :: prof_tip_vs_MC.rows_compared
\newcommand{\NSOneRowsMC}{\ensuremath{1025}}
% NSOneRowsMB [cert, exact] check_outA5.json :: ext813_vs_MB.rows_compared
\newcommand{\NSOneRowsMB}{\ensuremath{501}}
% NSOneMinMargin [cert, margin] check_outA5.json :: min:*.min_relative_margin.*
\newcommand{\NSOneMinMargin}{\ensuremath{1.2\cdot10^{-9}}}
% NTwoSegments [cert, exact] contour_result.json :: len:segments
\newcommand{\NTwoSegments}{\ensuremath{5058}}
% NTwoInteriorBoxes [cert, exact] interior_result.json :: len:
\newcommand{\NTwoInteriorBoxes}{\ensuremath{174}}
% NTwoContourBoxes [cert, exact] eta_cache.json :: len:
\newcommand{\NTwoContourBoxes}{\ensuremath{202}}
% NTwoEtaMax [cert, upper] eta_cache.json :: max:*.eta
\newcommand{\NTwoEtaMax}{\ensuremath{0.1163}}
% NTwoLipMax [cert, upper] eta_cache.json :: max:*.Lip
\newcommand{\NTwoLipMax}{\ensuremath{0.272}}
% NTwoLambdaZero [est, upper] cert_b0.25_th0.0_tau0.03.json :: Lam
\newcommand{\NTwoLambdaZero}{\ensuremath{1.790819}}
% NTwoLambdaOne [est, upper] cert_b0.25_th0.48_tau0.03.json :: Lam
\newcommand{\NTwoLambdaOne}{\ensuremath{1.768282}}
% NSOneDTightLower [doc, lower] doc:README.md :: D ∈ [2.50080627654, 2.50080627668]
\newcommand{\NSOneDTightLower}{\ensuremath{2.50080627654}}
% NSOneDTightUpper [doc, upper] doc:README.md :: D ∈ [2.50080627654, 2.50080627668]
\newcommand{\NSOneDTightUpper}{\ensuremath{2.50080627668}}
% NSOneKTightLower [doc, lower] doc:README.md :: k(0) ∈ [1.53982063762, 1.53982063822]
\newcommand{\NSOneKTightLower}{\ensuremath{1.53982063762}}
% NSOneKTightUpper [doc, upper] doc:README.md :: k(0) ∈ [1.53982063762, 1.53982063822]
\newcommand{\NSOneKTightUpper}{\ensuremath{1.53982063822}}
% NSOneDCondLower [doc, lower] doc:README.md :: D = −p(tip) ∈ [2.500806276619695, 2.500806276620049]
\newcommand{\NSOneDCondLower}{\ensuremath{2.500806276619695}}
% NSOneDCondUpper [doc, upper] doc:README.md :: D = −p(tip) ∈ [2.500806276619695, 2.500806276620049]
\newcommand{\NSOneDCondUpper}{\ensuremath{2.500806276620049}}
% NSOneKCondLower [doc, lower] doc:README.md :: k(0) ∈ [1.539820637871836, 1.5398206378721015]
\newcommand{\NSOneKCondLower}{\ensuremath{1.539820637871836}}
% NSOneKCondUpper [doc, upper] doc:README.md :: k(0) ∈ [1.539820637871836, 1.5398206378721015]
\newcommand{\NSOneKCondUpper}{\ensuremath{1.5398206378721015}}
% NSOneTauTipLower [doc, lower] doc:README.md :: τ_tip ∈ [10.249077651120778, 10.249077651121361]
\newcommand{\NSOneTauTipLower}{\ensuremath{10.249077651120778}}
% NSOneTauTipUpper [doc, upper] doc:README.md :: τ_tip ∈ [10.249077651120778, 10.249077651121361]
\newcommand{\NSOneTauTipUpper}{\ensuremath{10.249077651121361}}
% NSOneStateUncert [doc, upper] doc:README.md :: State uncertainty at σ = 8: 4.4e−13
\newcommand{\NSOneStateUncert}{\ensuremath{4.4\cdot10^{-13}}}
% NFarBoxEprime [doc, margin] doc:STEP4_RESULT.md :: (6.4287, 4.6917, 1.9940, 1.69999995)
\newcommand{\NFarBoxEprime}{\ensuremath{6.42}}
% NFarBoxXEsecond [doc, margin] doc:STEP4_RESULT.md :: (6.4287, 4.6917, 1.9940, 1.69999995)
\newcommand{\NFarBoxXEsecond}{\ensuremath{4.69}}
% NFarBoxXsqEthird [doc, margin] doc:STEP4_RESULT.md :: (6.4287, 4.6917, 1.9940, 1.69999995)
\newcommand{\NFarBoxXsqEthird}{\ensuremath{1.99}}
% NFarBoxE [doc, margin] doc:STEP4_RESULT.md :: (6.4287, 4.6917, 1.9940, 1.69999995)
\newcommand{\NFarBoxE}{\ensuremath{1.69}}
% NHFourKfourErr [doc, upper] doc:STEP4_RESULT.md :: | 5.60e-3 against 0.06 | ×10.7 |
\newcommand{\NHFourKfourErr}{\ensuremath{5.60\cdot10^{-3}}}
% NHFourKfourMargin [doc, margin] doc:STEP4_RESULT.md :: | 5.60e-3 against 0.06 | ×10.7 |
\newcommand{\NHFourKfourMargin}{\ensuremath{10.7}}
% NCoercivityCoefficient [doc, upper] doc:step3_proof_v32.tex :: \norm{e^{\tau\mathcal L}u}_\rho\le14.8\,e^{-0.0213\tau}\norm
\newcommand{\NCoercivityCoefficient}{\ensuremath{14.8}}
% NCoercivityRate [doc, lower] doc:step3_proof_v32.tex :: \norm{e^{\tau\mathcal L}u}_\rho\le14.8\,e^{-0.0213\tau}\norm
\newcommand{\NCoercivityRate}{\ensuremath{0.0213}}
% NRegionRampStart [hand, exact] hand :: 
\newcommand{\NRegionRampStart}{\ensuremath{14}}
% NRegionInterface [hand, exact] hand :: 
\newcommand{\NRegionInterface}{\ensuremath{28}}
% NRegionShellEnd [hand, exact] hand :: 
\newcommand{\NRegionShellEnd}{\ensuremath{60}}
% NRegionSegEnd [hand, exact] hand :: 
\newcommand{\NRegionSegEnd}{\ensuremath{350}}
% NRegionAnnEnd [hand, exact] hand :: 
\newcommand{\NRegionAnnEnd}{\ensuremath{700}}
% NZoneIIILow [hand, exact] hand :: 
\newcommand{\NZoneIIILow}{\ensuremath{253}}
% NZoneIIIHigh [hand, exact] hand :: 
\newcommand{\NZoneIIIHigh}{\ensuremath{379.5}}
% NConeCentreMin [hand, exact] hand :: 
\newcommand{\NConeCentreMin}{\ensuremath{300}}
% NConeCentreMax [hand, exact] hand :: 
\newcommand{\NConeCentreMax}{\ensuremath{380}}
% NFarCentreMin [hand, exact] hand :: 
\newcommand{\NFarCentreMin}{\ensuremath{490}}
% NConeL [hand, exact] hand :: 
\newcommand{\NConeL}{\ensuremath{50}}
% NConeX [hand, exact] hand :: 
\newcommand{\NConeX}{\ensuremath{3}}
% NFarL [hand, exact] hand :: 
\newcommand{\NFarL}{\ensuremath{113}}
% NFarX [hand, exact] hand :: 
\newcommand{\NFarX}{\ensuremath{2}}
}%
  {\PackageWarning{paper}{numbers/numbers_paper.tex not found (unit U0 not finished)}}

% ======================================================================
% NOTATION MACROS (OUTLINE 5.2 and notation.md).  One macro per symbol
% whose spelling is fixed by the notation table.  Plain letters (a, b, f, ...)
% have no macro.
% ======================================================================
% basic
\newcommand{\R}{\mathbb{R}}
\newcommand{\C}{\mathbb{C}}
\newcommand{\N}{\mathbb{N}}
\newcommand{\Z}{\mathbb{Z}}
\newcommand{\eps}{\varepsilon}
\newcommand{\ip}[2]{\langle #1,#2\rangle}
\newcommand{\abs}[1]{\lvert #1\rvert}
\newcommand{\norm}[1]{\lVert #1\rVert}
\newcommand{\dd}{\mathrm d}
\newcommand{\re}{\operatorname{Re}}
\newcommand{\im}{\operatorname{Im}}
% flow and curves
\newcommand{\arclen}{\mathfrak{s}}                 % arclength of the flow curve
\newcommand{\kss}{\kappa_{\mathfrak{ss}}}          % kappa_{ss} in the flow equation
\newcommand{\Area}{\operatorname{Area}}
\newcommand{\Len}{\mathcal L}                      % length of a curve
\newcommand{\spp}{\operatorname{sp}}                 % sup quantities sp^v_0, sp^v_1, sp^f_1 (were S^v_0, S^v_1, S^f_1)
\newcommand{\rot}{\operatorname{rot}}              % turning number
\newcommand{\Tsing}{T_*}                           % singular time
\newcommand{\Tmax}{T_{\max}}                       % maximal time of existence
\newcommand{\pstar}{p_*}                           % image point of the singularity
\newcommand{\Gz}{\Gamma_0}                         % initial curve
\newcommand{\vth}{\vartheta}                       % parameter of Gamma_0
\newcommand{\dclk}{\delta_t}                       % clock defect of step 4
\newcommand{\fall}{\varpi}                         % Fourier rates varpi_n = n^4-n^2
% rescaling
\newcommand{\lamone}{\lambda_1}
\newcommand{\rate}{\varsigma}                      % 1/4 - a
\newcommand{\xia}{\xi_a}                           % (1-4\bar a)^{-1}
% shrinker and step 1
\newcommand{\Sone}{S_1}
\newcommand{\Gsh}{\Gamma_{\!*}}                    % profile curve of S_1
\newcommand{\cSet}{\mathcal S}                     % certified set of shrinkers
\newcommand{\sigtip}{\sigma_{\rm tip}}
\newcommand{\Xo}{X_0}                              % 2.2490776515
\newcommand{\Dtip}{D}
\newcommand{\ktip}{k_0}
\newcommand{\shootM}{M}                            % shooting function
\newcommand{\Ypar}{\mathcal Y}                     % step-1 parallelepiped
\newcommand{\axis}{\Sigma}                         % symmetry axis of S_1
\newcommand{\sigc}{\sigma_{\rm c}}
\newcommand{\sigrev}{\sigma^{\rm rev}}
% linear theory
\newcommand{\Hrho}{\mathcal H_\rho}
\newcommand{\muslow}{\mu_3}
\newcommand{\rre}{\mathfrak r_{\rm re}}
\newcommand{\rim}{\mathfrak r_{\rm im}}
\newcommand{\Evev}{\mathrm{Ev}_{\rm even}}
\newcommand{\Evod}{\mathrm{Ev}_{\rm odd}}
\newcommand{\gvec}{\mathbf g}                      % translation vector in step 2
\newcommand{\cons}[1]{\mathfrak f_{#1}}            % constraint functions
\newcommand{\phirho}{\varphi_\rho}
% energies
\newcommand{\Lam}{\Lambda}
\newcommand{\En}{E}
\newcommand{\EnI}{E_I}
\newcommand{\muI}{\mu_I}                           % 0.08
\newcommand{\epsE}{\varepsilon_E}                  % 1e-5
\newcommand{\epsI}{\varepsilon_I}                  % 3e-5
\newcommand{\wn}[1]{\mathsf w_{#1}}                % ||w~^{(j)}||_rho
\newcommand{\vI}{v_I}                              % zeta w'
\newcommand{\muerr}{\mu_{\rm err}}                 % step-4 error-energy weight 0.036
% far field
\newcommand{\pU}{p_U}                              % 1.0246
\newcommand{\Uone}{U_1}
\newcommand{\sm}{\operatorname{sm}}                % smoothstep t^3(10-15t+6t^2)
\newcommand{\rc}{r_c}                              % centre radius (Lemma U)
\newcommand{\rchi}{r_\chi}                         % 253
\newcommand{\Iw}{I_w}                              % hand-off windows [y_w,1.5y_w]
\newcommand{\qa}{\mathfrak q_1}                    % quasilinear coefficients of the slope equation
\newcommand{\qb}{\mathfrak q_2}
% cylinders (Lemma O)
\newcommand{\cyl}{\ell}                            % cylinder length unit
\newcommand{\MO}{M_O}
\newcommand{\KO}{K_O}
\newcommand{\Hcyl}{H_{\rm cyl}}
\newcommand{\XiO}{\Xi_O}
\newcommand{\fO}{f_O}
% step 4
\newcommand{\Ybar}{\bar Y}
\newcommand{\Lt}{\widetilde L}                     % linearisation about the step-4 base
\newcommand{\Trans}{\Pi}                           % transport integral
\newcommand{\resid}{\mathrm{res}}
\newcommand{\Resid}[1]{\mathrm{Res}_{#1}}
\newcommand{\boxtip}{\Box_{\rm tip}}
\newcommand{\boxfar}{\Box_{\rm far}}
\newcommand{\Eenv}{\mathcal E_{\rm env}}
\newcommand{\betab}{\beta_{\rm base}}
\newcommand{\tauI}{\tau^{\rm I}}                   % Phase I clock
% names of families of conditions: \Hh{3} = (H3), \Oh{4} = (O4), ...
\newcommand{\Hh}[1]{\textup{(H#1)}}
\newcommand{\Oh}[1]{\textup{(O#1)}}
\newcommand{\Lh}[1]{\textup{(L#1)}}
\newcommand{\Bh}[1]{\textup{(B#1)}}
\newcommand{\Kh}[1]{\textup{(K#1)}}
\newcommand{\Wh}[1]{\textup{(W#1)}}
\newcommand{\PhaseI}{Phase~I}
\newcommand{\PhaseII}{Phase~II}
% margin: threshold/bound (always)
\newcommand{\mgn}[1]{{\times}#1}

% ---- figures: placeholder macro so that drafts compile before F1-F7 exist ----
\newcommand{\figplaceholder}[1]{\fbox{\parbox{0.8\linewidth}{\centering\small[figure #1 to be supplied]}}}
 ...
% (the main paper does nothing special: \ifsupplement defaults to false).
% Compile from paper_v1/ so that % numbers_paper.tex -- GENERATED from ledger.json by build_ledger.py; do not edit.
% Rounding: upper bounds up, lower bounds and margins down at the printed digit; margin = threshold / bound.
% Use as $\times\NName$ for margins.  kind: cert certified, est estimate, hand hypothesis constant, doc only a document quotes it.
% NUEhatZero [cert, upper] U_cert_v35.json :: supersolution.EA.0
\newcommand{\NUEhatZero}{\ensuremath{1.2018\cdot10^{-6}}}
% NUEhatOne [cert, upper] U_cert_v35.json :: supersolution.EA.1
\newcommand{\NUEhatOne}{\ensuremath{1.0441\cdot10^{-4}}}
% NUEhatTwo [cert, upper] U_cert_v35.json :: supersolution.EA.2
\newcommand{\NUEhatTwo}{\ensuremath{0.056694}}
% NUBhatZero [cert, upper] U_cert_v35.json :: supersolution.Bhat.0
\newcommand{\NUBhatZero}{\ensuremath{4.8735\cdot10^{-6}}}
% NUBhatOne [cert, upper] U_cert_v35.json :: supersolution.Bhat.1
\newcommand{\NUBhatOne}{\ensuremath{5.5933\cdot10^{-4}}}
% NUBhatTwo [cert, upper] U_cert_v35.json :: supersolution.Bhat.2
\newcommand{\NUBhatTwo}{\ensuremath{0.36675}}
% NUSupersolutionCZeroZero [cert, upper] U_cert_v35.json :: supersolution.C0.0
\newcommand{\NUSupersolutionCZeroZero}{\ensuremath{1.1340\cdot10^{-6}}}
% NUSupersolutionCZeroOne [cert, upper] U_cert_v35.json :: supersolution.C0.1
\newcommand{\NUSupersolutionCZeroOne}{\ensuremath{1.6623\cdot10^{-4}}}
% NUSupersolutionCOneZero [cert, upper] U_cert_v35.json :: supersolution.C1.0
\newcommand{\NUSupersolutionCOneZero}{\ensuremath{1.5241\cdot10^{-9}}}
% NUSupersolutionCOneOne [cert, upper] U_cert_v35.json :: supersolution.C1.1
\newcommand{\NUSupersolutionCOneOne}{\ensuremath{5.5255\cdot10^{-8}}}
% NUSupersolutionNTwoZero [cert, upper] U_cert_v35.json :: supersolution.N2.0
\newcommand{\NUSupersolutionNTwoZero}{\ensuremath{0.071472}}
% NUSupersolutionNTwoOne [cert, upper] U_cert_v35.json :: supersolution.N2.1
\newcommand{\NUSupersolutionNTwoOne}{\ensuremath{0.33012}}
% NUSupersolutionNTwoTwo [cert, upper] U_cert_v35.json :: supersolution.N2.2
\newcommand{\NUSupersolutionNTwoTwo}{\ensuremath{5.0406}}
% NUSupersolutionNTwoOverFourNTwoRefZero [cert, upper] U_cert_v35.json :: supersolution.N2_over_4N2ref.0
\newcommand{\NUSupersolutionNTwoOverFourNTwoRefZero}{\ensuremath{0.18}}
% NUSupersolutionNTwoOverFourNTwoRefOne [cert, upper] U_cert_v35.json :: supersolution.N2_over_4N2ref.1
\newcommand{\NUSupersolutionNTwoOverFourNTwoRefOne}{\ensuremath{0.42}}
% NUSupersolutionNTwoOverFourNTwoRefTwo [cert, upper] U_cert_v35.json :: supersolution.N2_over_4N2ref.2
\newcommand{\NUSupersolutionNTwoOverFourNTwoRefTwo}{\ensuremath{0.26}}
% NUSupersolutionMOneZHypCheckZonesIIISOneTwoBhat [cert, upper] U_cert_v35.json :: supersolution.M1z_hyp_check.zones_I_II.S1_2Bhat
\newcommand{\NUSupersolutionMOneZHypCheckZonesIIISOneTwoBhat}{\ensuremath{0.28585}}
% NUSupersolutionMOneZHypCheckZoneIIIBOnePlusPUOne [cert, upper] U_cert_v35.json :: supersolution.M1z_hyp_check.zone_III.b1_plus_pU1
\newcommand{\NUSupersolutionMOneZHypCheckZoneIIIBOnePlusPUOne}{\ensuremath{0.50786}}
% NUMarginLateMOne [cert, margin] U_cert_v35.json :: margins.late_M1
\newcommand{\NUMarginLateMOne}{\ensuremath{2.353}}
% NUMarginLateMZero [cert, margin] U_cert_v35.json :: margins.late_M0
\newcommand{\NUMarginLateMZero}{\ensuremath{4.892}}
% NUMarginsLateMZeroZ [cert, margin] U_cert_v35.json :: margins.late_M0z
\newcommand{\NUMarginsLateMZeroZ}{\ensuremath{12.99}}
% NUMarginPresentMOne [cert, margin] U_cert_v35.json :: margins.present_M1_min.value
\newcommand{\NUMarginPresentMOne}{\ensuremath{2.770}}
% NUMarginPresentMZero [cert, margin] U_cert_v35.json :: margins.present_M0_min.value
\newcommand{\NUMarginPresentMZero}{\ensuremath{2.225}}
% NUMarginsPresentMZeroZMinValue [cert, margin] U_cert_v35.json :: margins.present_M0z_min.value
\newcommand{\NUMarginsPresentMZeroZMinValue}{\ensuremath{14.37}}
% NULoopGainZeroTotal [cert, upper] U_cert_v35.json :: loop_gain.0.total
\newcommand{\NULoopGainZeroTotal}{\ensuremath{1.23\cdot10^{-7}}}
% NULoopGainOneTotal [cert, upper] U_cert_v35.json :: loop_gain.1.total
\newcommand{\NULoopGainOneTotal}{\ensuremath{0.037169}}
% NULoopGainTwoNeighbour [cert, upper] U_cert_v35.json :: loop_gain.2.neighbour
\newcommand{\NULoopGainTwoNeighbour}{\ensuremath{0.037018}}
% NULoopGainTwoThroughA [cert, upper] U_cert_v35.json :: loop_gain.2.through_A
\newcommand{\NULoopGainTwoThroughA}{\ensuremath{0.0020172}}
% NUfineBhatZero [cert, upper] U_cert_v35.json :: ufine.Bhat.0
\newcommand{\NUfineBhatZero}{\ensuremath{4.8723\cdot10^{-6}}}
% NUfineBhatOne [cert, upper] U_cert_v35.json :: ufine.Bhat.1
\newcommand{\NUfineBhatOne}{\ensuremath{2.6266\cdot10^{-4}}}
% NUfineBhatTwo [cert, upper] U_cert_v35.json :: ufine.Bhat.2
\newcommand{\NUfineBhatTwo}{\ensuremath{0.025842}}
% NUcdeltaU [cert, upper] U_cert_v35.json :: launch.c_dU
\newcommand{\NUcdeltaU}{\ensuremath{1.7187\cdot10^{-8}}}
% NULaunchWmin [cert, lower] U_cert_v35.json :: launch.wmin
\newcommand{\NULaunchWmin}{\ensuremath{0.73492}}
% NULaunchCAD [cert, upper] U_cert_v35.json :: launch.CA_D
\newcommand{\NULaunchCAD}{\ensuremath{2.2429}}
% NULaunchQD [cert, upper] U_cert_v35.json :: launch.qD
\newcommand{\NULaunchQD}{\ensuremath{0.018509}}
% NULaunchDeltaInZero [cert, upper] U_cert_v35.json :: launch.Delta_in.0
\newcommand{\NULaunchDeltaInZero}{\ensuremath{3.131\cdot10^{-12}}}
% NULaunchDeltaInOne [cert, upper] U_cert_v35.json :: launch.Delta_in.1
\newcommand{\NULaunchDeltaInOne}{\ensuremath{3.287\cdot10^{-12}}}
% NULaunchDeltaInTwo [cert, upper] U_cert_v35.json :: launch.Delta_in.2
\newcommand{\NULaunchDeltaInTwo}{\ensuremath{3.451\cdot10^{-12}}}
% NULaunchPhiZero [cert, upper] U_cert_v35.json :: launch.Phi.0
\newcommand{\NULaunchPhiZero}{\ensuremath{6.434\cdot10^{-9}}}
% NULaunchPhiOne [cert, upper] U_cert_v35.json :: launch.Phi.1
\newcommand{\NULaunchPhiOne}{\ensuremath{1.4217\cdot10^{-8}}}
% NULaunchPhiTwo [cert, upper] U_cert_v35.json :: launch.Phi.2
\newcommand{\NULaunchPhiTwo}{\ensuremath{3.8404\cdot10^{-8}}}
% NUGammaZero [cert, upper] U_cert_v35.json :: Gamma.0
\newcommand{\NUGammaZero}{\ensuremath{4.0000358}}
% NUGammaOne [cert, upper] U_cert_v35.json :: Gamma.1
\newcommand{\NUGammaOne}{\ensuremath{16.000322}}
% NUGammaTwo [cert, upper] U_cert_v35.json :: Gamma.2
\newcommand{\NUGammaTwo}{\ensuremath{64.00218}}
% NUKappa [cert, upper] U_cert_v35.json :: kappa
\newcommand{\NUKappa}{\ensuremath{8.0000358}}
% NUKappaTwo [cert, upper] U_cert_v35.json :: kappa2
\newcommand{\NUKappaTwo}{\ensuremath{64.00061}}
% NUGhatZeroZero [cert, upper] U_cert_v35.json :: Ghat.0.0
\newcommand{\NUGhatZeroZero}{\ensuremath{51322.6}}
% NUGhatOneOne [cert, upper] U_cert_v35.json :: Ghat.1.1
\newcommand{\NUGhatOneOne}{\ensuremath{33962.9}}
% NUGhatTwoTwo [cert, upper] U_cert_v35.json :: Ghat.2.2
\newcommand{\NUGhatTwoTwo}{\ensuremath{348341.1}}
% NUGhatOneZero [cert, upper] U_cert_v35.json :: Ghat.1.0
\newcommand{\NUGhatOneZero}{\ensuremath{4.938}}
% NUGhatTwoZero [cert, upper] U_cert_v35.json :: Ghat.2.0
\newcommand{\NUGhatTwoZero}{\ensuremath{27.984}}
% NUGhatTwoOne [cert, upper] U_cert_v35.json :: Ghat.2.1
\newcommand{\NUGhatTwoOne}{\ensuremath{458.53}}
% NUKZeroZero [cert, upper] U_cert_v35.json :: K.0.0
\newcommand{\NUKZeroZero}{\ensuremath{8034.21}}
% NUKZeroOne [cert, upper] U_cert_v35.json :: K.0.1
\newcommand{\NUKZeroOne}{\ensuremath{724.89}}
% NUKOneZero [cert, upper] U_cert_v35.json :: K.1.0
\newcommand{\NUKOneZero}{\ensuremath{0.08152}}
% NUKOneOne [cert, upper] U_cert_v35.json :: K.1.1
\newcommand{\NUKOneOne}{\ensuremath{6336.94}}
% NUKOneTwo [cert, upper] U_cert_v35.json :: K.1.2
\newcommand{\NUKOneTwo}{\ensuremath{545.59}}
% NUKOneThree [cert, upper] U_cert_v35.json :: K.1.3
\newcommand{\NUKOneThree}{\ensuremath{0.58700}}
% NUKTwoZero [cert, upper] U_cert_v35.json :: K.2.0
\newcommand{\NUKTwoZero}{\ensuremath{0.87186}}
% NUKTwoOne [cert, upper] U_cert_v35.json :: K.2.1
\newcommand{\NUKTwoOne}{\ensuremath{528.62}}
% NUKTwoTwo [cert, upper] U_cert_v35.json :: K.2.2
\newcommand{\NUKTwoTwo}{\ensuremath{23425.98}}
% NUKTwoThree [cert, upper] U_cert_v35.json :: K.2.3
\newcommand{\NUKTwoThree}{\ensuremath{555.62}}
% NUFootwOne [cert, upper] U_cert_v35.json :: footw.1
\newcommand{\NUFootwOne}{\ensuremath{2.6361}}
% NUFootwTwo [cert, upper] U_cert_v35.json :: footw.2
\newcommand{\NUFootwTwo}{\ensuremath{2663.3}}
% NUCThreeThree [cert, upper] U_cert_v35.json :: c33
\newcommand{\NUCThreeThree}{\ensuremath{108.02}}
% NUCTwoTwoThree [cert, upper] U_cert_v35.json :: c223
\newcommand{\NUCTwoTwoThree}{\ensuremath{7.3284}}
% NUCEta [cert, upper] U_cert_v35.json :: C_eta
\newcommand{\NUCEta}{\ensuremath{13.376}}
% NULMax [cert, upper] U_cert_v35.json :: L_max
\newcommand{\NULMax}{\ensuremath{3.5352}}
% NURRatio [cert, upper] U_cert_v35.json :: R_ratio
\newcommand{\NURRatio}{\ensuremath{11.593}}
% NUSweepsSigmaInCoverSup [cert, upper] U_cert_v35.json :: sweeps.Sigma_in.cover_sup
\newcommand{\NUSweepsSigmaInCoverSup}{\ensuremath{5.0892}}
% NUSweepsSigmaECoverSup [cert, upper] U_cert_v35.json :: sweeps.Sigma_E.cover_sup
\newcommand{\NUSweepsSigmaECoverSup}{\ensuremath{0.30846}}
% NUSweepsSXCoverSup [cert, upper] U_cert_v35.json :: sweeps.S_X.cover_sup
\newcommand{\NUSweepsSXCoverSup}{\ensuremath{1.5243}}
% NUSweepsSigmaXCoverSup [cert, upper] U_cert_v35.json :: sweeps.Sigma_X.cover_sup
\newcommand{\NUSweepsSigmaXCoverSup}{\ensuremath{3.0081}}
% NUSweepsSigmaRsTwoCoverSup [cert, upper] U_cert_v35.json :: sweeps.Sigma_rs.2.cover_sup
\newcommand{\NUSweepsSigmaRsTwoCoverSup}{\ensuremath{4.3350\cdot10^{-3}}}
% NUDataTermsLateDZero [cert, upper] U_cert_v35.json :: data_terms_late.D.0
\newcommand{\NUDataTermsLateDZero}{\ensuremath{1.7330\cdot10^{-6}}}
% NUDataTermsLateDOne [cert, upper] U_cert_v35.json :: data_terms_late.D.1
\newcommand{\NUDataTermsLateDOne}{\ensuremath{2.5353\cdot10^{-4}}}
% NUDataTermsLateSlopeShareCOneSigmaXOverDkZero [cert, upper] U_cert_v35.json :: data_terms_late.slope_share_C1SigmaX_over_Dk.0
\newcommand{\NUDataTermsLateSlopeShareCOneSigmaXOverDkZero}{\ensuremath{2.7\cdot10^{-3}}}
% NUDataTermsLateSlopeShareCOneSigmaXOverDkOne [cert, upper] U_cert_v35.json :: data_terms_late.slope_share_C1SigmaX_over_Dk.1
\newcommand{\NUDataTermsLateSlopeShareCOneSigmaXOverDkOne}{\ensuremath{6.6\cdot10^{-4}}}
% NURLtTwoMarginZero [cert, margin] U_cert_v35.json :: R_lt2_margin.0
\newcommand{\NURLtTwoMarginZero}{\ensuremath{4.432}}
% NURLtTwoMarginOne [cert, margin] U_cert_v35.json :: R_lt2_margin.1
\newcommand{\NURLtTwoMarginOne}{\ensuremath{6.697}}
% NUThetaTwoEight [cert, upper] U_cert_v35.json :: theta28
\newcommand{\NUThetaTwoEight}{\ensuremath{0.37011}}
% NUDFiveSix [cert, upper] U_cert_v35.json :: D56
\newcommand{\NUDFiveSix}{\ensuremath{1.0368}}
% NUDTwoEight [cert, upper] U_cert_v35.json :: D28
\newcommand{\NUDTwoEight}{\ensuremath{1.5876}}
% NUSThreeThreeFootOverBhatTwo [cert, upper] U_cert_v35.json :: S33_foot_over_Bhat2
\newcommand{\NUSThreeThreeFootOverBhatTwo}{\ensuremath{1.69\cdot10^{-3}}}
% NSAHJWKappaJ [cert, upper] A_interface.json :: HJW.kappa_J
\newcommand{\NSAHJWKappaJ}{\ensuremath{0.042351}}
% NSAGpZero [cert, lower] A_interface.json :: res.tight.gp.0
\newcommand{\NSAGpZero}{\ensuremath{0.24292}}
% NSAGpOne [cert, lower] A_interface.json :: res.tight.gp.1
\newcommand{\NSAGpOne}{\ensuremath{0.24106}}
% NSAGpTwo [cert, lower] A_interface.json :: res.tight.gp.2
\newcommand{\NSAGpTwo}{\ensuremath{0.11961}}
% NSAAminC [cert, lower] A_interface.json :: res.tight.amin_c
\newcommand{\NSAAminC}{\ensuremath{0.97623}}
% NSARhoJ [cert, upper] A_interface.json :: res.tight.rhoJ
\newcommand{\NSARhoJ}{\ensuremath{5.4995\cdot10^{-7}}}
% NSAKappaOTwoTwo [cert, upper] A_interface.json :: res.tight.kappa_o.2.2
\newcommand{\NSAKappaOTwoTwo}{\ensuremath{1.2416\cdot10^{-5}}}
% NSADsrcZero [cert, upper] A_interface.json :: res.tight.Dsrc.0
\newcommand{\NSADsrcZero}{\ensuremath{1.8092\cdot10^{-7}}}
% NSADsrcOne [cert, upper] A_interface.json :: res.tight.Dsrc.1
\newcommand{\NSADsrcOne}{\ensuremath{9.8276\cdot10^{-5}}}
% NSABZeroDWindowSlopeZero [cert, upper] A_interface.json :: res.tight.B0.D_window_slope.0
\newcommand{\NSABZeroDWindowSlopeZero}{\ensuremath{8.7399\cdot10^{-10}}}
% NSABZeroDWindowSlopeOne [cert, upper] A_interface.json :: res.tight.B0.D_window_slope.1
\newcommand{\NSABZeroDWindowSlopeOne}{\ensuremath{5.2599\cdot10^{-8}}}
% NSAFTwoZero [cert, upper] A_interface.json :: res.tight.F2.0
\newcommand{\NSAFTwoZero}{\ensuremath{5.6506\cdot10^{-13}}}
% NSAFTwoOne [cert, upper] A_interface.json :: res.tight.F2.1
\newcommand{\NSAFTwoOne}{\ensuremath{5.9320\cdot10^{-13}}}
% NSACoupOne [cert, upper] A_interface.json :: res.tight.coup.1
\newcommand{\NSACoupOne}{\ensuremath{1.1568\cdot10^{-12}}}
% NSACOneZeroZero [cert, upper] A_interface.json :: res.tight.C1_0.0
\newcommand{\NSACOneZeroZero}{\ensuremath{7.2060\cdot10^{-10}}}
% NSACOneZeroOne [cert, upper] A_interface.json :: res.tight.C1_0.1
\newcommand{\NSACOneZeroOne}{\ensuremath{4.3143\cdot10^{-8}}}
% NSABZeroSupCoreZero [cert, upper] A_interface.json :: res.tight.B0.sup_core.0
\newcommand{\NSABZeroSupCoreZero}{\ensuremath{4.23\cdot10^{-5}}}
% NSACoreFluxMarginTwoFLUXZeroOverKZero [cert, margin] A_interface.json :: res.tight.core.flux_margin.2.FLUX0_over_K0
\newcommand{\NSACoreFluxMarginTwoFLUXZeroOverKZero}{\ensuremath{3.396}}
% NSACoreFluxMarginTwoFLUXTwoOverKTwo [cert, margin] A_interface.json :: res.tight.core.flux_margin.2.FLUX2_over_K2
\newcommand{\NSACoreFluxMarginTwoFLUXTwoOverKTwo}{\ensuremath{6.012}}
% NSAMarginFluxZero [cert, margin] asm_design_N2ref4.json :: res.tight.flux.beta.2.margin_FLUX0
\newcommand{\NSAMarginFluxZero}{\ensuremath{3.365}}
% NSAMarginFluxTwo [cert, margin] asm_design_N2ref4.json :: res.tight.flux.beta.2.margin_FLUX2
\newcommand{\NSAMarginFluxTwo}{\ensuremath{5.955}}
% NSAfourSZero [cert, upper] asm_design_N2ref4.json :: res.tight.S0
\newcommand{\NSAfourSZero}{\ensuremath{7.1878\cdot10^{-3}}}
% NSAfourSOne [cert, upper] asm_design_N2ref4.json :: res.tight.S1
\newcommand{\NSAfourSOne}{\ensuremath{0.10430}}
% NSAMarginCoreSZeroV [cert, margin] asm_design_N2ref4.json :: res.tight.margins.S0v_vs_0.0153
\newcommand{\NSAMarginCoreSZeroV}{\ensuremath{2.116}}
% NSAMarginCoreSOneV [cert, margin] asm_design_N2ref4.json :: res.tight.margins.S1v_vs_0.6
\newcommand{\NSAMarginCoreSOneV}{\ensuremath{5.748}}
% NSAMarginCoreMOne [cert, margin] asm_design_N2ref4.json :: res.tight.margins.M1core
\newcommand{\NSAMarginCoreMOne}{\ensuremath{4.409}}
% NSAMarginCoreMZero [cert, margin] asm_design_N2ref4.json :: res.tight.margins.M0core
\newcommand{\NSAMarginCoreMZero}{\ensuremath{5.450}}
% NSAMarginGhat [cert, margin] asm_design_N2ref4.json :: res.tight.margins.Ghat22
\newcommand{\NSAMarginGhat}{\ensuremath{1.917}}
% NSAMarginRampMOne [cert, margin] asm_design_N2ref4.json :: res.tight.margins.ramp M1
\newcommand{\NSAMarginRampMOne}{\ensuremath{1.184}}
% NSAfourInfoZeroDWindowConst [cert, upper] asm_design_N2ref4.json :: res.tight.info.0.D_window_const
\newcommand{\NSAfourInfoZeroDWindowConst}{\ensuremath{1.7732\cdot10^{-7}}}
% NSAfourInfoOneDWindowConst [cert, upper] asm_design_N2ref4.json :: res.tight.info.1.D_window_const
\newcommand{\NSAfourInfoOneDWindowConst}{\ensuremath{9.8058\cdot10^{-5}}}
% NZDeltaRZero [cert, upper] zone3_interface.json :: DeltaR.0
\newcommand{\NZDeltaRZero}{\ensuremath{4.4986\cdot10^{-8}}}
% NZDeltaROne [cert, upper] zone3_interface.json :: DeltaR.1
\newcommand{\NZDeltaROne}{\ensuremath{6.7972\cdot10^{-6}}}
% NZDeltaRTwo [cert, upper] zone3_interface.json :: DeltaR.2
\newcommand{\NZDeltaRTwo}{\ensuremath{0.011480}}
% NZDeltaRBZero [cert, upper] zone3_interface.json :: DeltaR_b.0
\newcommand{\NZDeltaRBZero}{\ensuremath{2.66\cdot10^{-12}}}
% NZDeltaRBOne [cert, upper] zone3_interface.json :: DeltaR_b.1
\newcommand{\NZDeltaRBOne}{\ensuremath{3.26\cdot10^{-9}}}
% NZDeltaRBTwo [cert, upper] zone3_interface.json :: DeltaR_b.2
\newcommand{\NZDeltaRBTwo}{\ensuremath{4.35\cdot10^{-8}}}
% NZDeltaRDetailNSqZero [cert, upper] zone3_interface.json :: DeltaR_detail.N_sq.0
\newcommand{\NZDeltaRDetailNSqZero}{\ensuremath{8.28\cdot10^{-5}}}
% NZDeltaRDetailNSqOne [cert, upper] zone3_interface.json :: DeltaR_detail.N_sq.1
\newcommand{\NZDeltaRDetailNSqOne}{\ensuremath{0.10149}}
% NZDeltaRDetailNSqTwo [cert, upper] zone3_interface.json :: DeltaR_detail.N_sq.2
\newcommand{\NZDeltaRDetailNSqTwo}{\ensuremath{1.3567}}
% NZDeltaRDetailNSqThree [cert, upper] zone3_interface.json :: DeltaR_detail.N_sq.3
\newcommand{\NZDeltaRDetailNSqThree}{\ensuremath{25350}}
% NZDeltaRDetailPartsTwoPairsThree [cert, upper] zone3_interface.json :: DeltaR_detail.parts.2.pairs.3
\newcommand{\NZDeltaRDetailPartsTwoPairsThree}{\ensuremath{0.011298}}
% NZDeltaRDetailPartsTwoLFourFoot [cert, upper] zone3_interface.json :: DeltaR_detail.parts.2.L4_foot
\newcommand{\NZDeltaRDetailPartsTwoLFourFoot}{\ensuremath{1.6\cdot10^{-4}}}
% NZDeltaRDetailPartsTwoCubic [cert, upper] zone3_interface.json :: DeltaR_detail.parts.2.cubic
\newcommand{\NZDeltaRDetailPartsTwoCubic}{\ensuremath{9.3\cdot10^{-6}}}
% NZDeltaRDetailPartsTwoSThreeThreeFoot [cert, upper] zone3_interface.json :: DeltaR_detail.parts.2.S33_foot
\newcommand{\NZDeltaRDetailPartsTwoSThreeThreeFoot}{\ensuremath{3.7\cdot10^{-7}}}
% NZNIIIZero [cert, upper] zone3_interface.json :: n_III.0
\newcommand{\NZNIIIZero}{\ensuremath{0.071445}}
% NZNIIIOne [cert, upper] zone3_interface.json :: n_III.1
\newcommand{\NZNIIIOne}{\ensuremath{0.32707}}
% NZNIIITwo [cert, upper] zone3_interface.json :: n_III.2
\newcommand{\NZNIIITwo}{\ensuremath{3.0445}}
% NZTauIII [cert, upper] zone3_interface.json :: tau_III
\newcommand{\NZTauIII}{\ensuremath{4.0708}}
% NZoneIIIBoxZero [cert, upper] zone3_interface.json :: box_III_380.b0
\newcommand{\NZoneIIIBoxZero}{\ensuremath{0.10032}}
% NZoneIIIBoxOne [cert, upper] zone3_interface.json :: box_III_380.b1
\newcommand{\NZoneIIIBoxOne}{\ensuremath{0.41974}}
% NZoneIIIBoxZ [cert, upper] zone3_interface.json :: box_III_380.bz
\newcommand{\NZoneIIIBoxZ}{\ensuremath{0.014291}}
% NZBoxIIIThreeEightZeroMarginsBZero [cert, margin] zone3_interface.json :: box_III_380.margins.b0
\newcommand{\NZBoxIIIThreeEightZeroMarginsBZero}{\ensuremath{1.495}}
% NZoneIIIMarginThinnest [cert, margin] zone3_interface.json :: box_III_380.margins.b1
\newcommand{\NZoneIIIMarginThinnest}{\ensuremath{1.238}}
% NZBoxIIIThreeEightZeroMarginsBz [cert, margin] zone3_interface.json :: box_III_380.margins.bz
\newcommand{\NZBoxIIIThreeEightZeroMarginsBz}{\ensuremath{11.19}}
% NZBoxIIISevenZeroZeroBZero [cert, upper] zone3_interface.json :: box_III_700.b0
\newcommand{\NZBoxIIISevenZeroZeroBZero}{\ensuremath{0.22757}}
% NZBoxIIISevenZeroZeroBOne [cert, upper] zone3_interface.json :: box_III_700.b1
\newcommand{\NZBoxIIISevenZeroZeroBOne}{\ensuremath{0.54234}}
% NZUfineKTwoAtBhatRef [cert, upper] zone3_interface.json :: ufine.K2_at_Bhat_ref
\newcommand{\NZUfineKTwoAtBhatRef}{\ensuremath{0.046956}}
% NZUfineKThreeStripAtBhatRef [cert, upper] zone3_interface.json :: ufine.K3_strip_at_Bhat_ref
\newcommand{\NZUfineKThreeStripAtBhatRef}{\ensuremath{0.054140}}
% NRecoveryMinMargin [cert, margin] zone3_interface.json :: cone_family.H6.at_Bhat_ref.min_margin
\newcommand{\NRecoveryMinMargin}{\ensuremath{1.860}}
% NZConeFamilyHSixAtBhatRefConeSlopeMargin [cert, margin] zone3_interface.json :: cone_family.H6.at_Bhat_ref.cone.slope.margin
\newcommand{\NZConeFamilyHSixAtBhatRefConeSlopeMargin}{\ensuremath{2.46}}
% NZConeFamilyHSixAtBhatRefConeKTwoMargin [cert, margin] zone3_interface.json :: cone_family.H6.at_Bhat_ref.cone.K2.margin
\newcommand{\NZConeFamilyHSixAtBhatRefConeKTwoMargin}{\ensuremath{2.54}}
% NZConeFamilyHSixAtBhatRefConeKThreeMargin [cert, margin] zone3_interface.json :: cone_family.H6.at_Bhat_ref.cone.K3.margin
\newcommand{\NZConeFamilyHSixAtBhatRefConeKThreeMargin}{\ensuremath{2.51}}
% NConeTstarMargin [cert, margin] zone3_interface.json :: acceptance.tstar_margin
\newcommand{\NConeTstarMargin}{\ensuremath{2.30}}
% NUoneTFiveYaddLevelOneTwo [cert, upper] u1_ghat.json :: t5_Yadd.level1.2
\newcommand{\NUoneTFiveYaddLevelOneTwo}{\ensuremath{0.0044}}
% NUoneTFiveYaddLevelTwoTwo [cert, upper] u1_ghat.json :: t5_Yadd.level2.2
\newcommand{\NUoneTFiveYaddLevelTwoTwo}{\ensuremath{0.1081}}
% NUoneTFiveYaddLevelTwoThree [cert, upper] u1_ghat.json :: t5_Yadd.level2.3
\newcommand{\NUoneTFiveYaddLevelTwoThree}{\ensuremath{0.374}}
% NUoneWeightThetaGammaZero [cert, upper] u1_ghat.json :: weight.Theta.Gamma.0
\newcommand{\NUoneWeightThetaGammaZero}{\ensuremath{4.00000000000000000000054}}
% NRKZeroZero [cert, upper] ramp_U.json :: K.0.0
\newcommand{\NRKZeroZero}{\ensuremath{13436.03}}
% NRKOneOne [cert, upper] ramp_U.json :: K.1.1
\newcommand{\NRKOneOne}{\ensuremath{10526.19}}
% NRKTwoTwo [cert, upper] ramp_U.json :: K.2.2
\newcommand{\NRKTwoTwo}{\ensuremath{39482.30}}
% NUoneCasesPrimaryYNextZero [cert, upper] u1_ghat.json :: cases.primary.Y_next.0
\newcommand{\NUoneCasesPrimaryYNextZero}{\ensuremath{457.092}}
% NUoneCasesPrimaryYNextOne [cert, upper] u1_ghat.json :: cases.primary.Y_next.1
\newcommand{\NUoneCasesPrimaryYNextOne}{\ensuremath{356.940}}
% NUoneCasesPrimaryYNextTwo [cert, upper] u1_ghat.json :: cases.primary.Y_next.2
\newcommand{\NUoneCasesPrimaryYNextTwo}{\ensuremath{1320.220}}
% NSfourResultsMNineHTwoTightenedFUOneRowsZeroBound [cert, upper] FINAL_manifest_v33.json :: results.M9_H2_tightened_fU1.rows.0.bound
\newcommand{\NSfourResultsMNineHTwoTightenedFUOneRowsZeroBound}{\ensuremath{3.658\cdot10^{-10}}}
% NSfourResultsMNineHTwoTightenedFUOneRowsZeroMargin [cert, margin] FINAL_manifest_v33.json :: results.M9_H2_tightened_fU1.rows.0.margin
\newcommand{\NSfourResultsMNineHTwoTightenedFUOneRowsZeroMargin}{\ensuremath{27.34}}
% NSfourResultsMNineHTwoTightenedFUOneRowsOneBound [cert, upper] FINAL_manifest_v33.json :: results.M9_H2_tightened_fU1.rows.1.bound
\newcommand{\NSfourResultsMNineHTwoTightenedFUOneRowsOneBound}{\ensuremath{1.777\cdot10^{-9}}}
% NSfourResultsMNineHTwoTightenedFUOneRowsOneMargin [cert, margin] FINAL_manifest_v33.json :: results.M9_H2_tightened_fU1.rows.1.margin
\newcommand{\NSfourResultsMNineHTwoTightenedFUOneRowsOneMargin}{\ensuremath{562.8}}
% NSfourResultsMNineHTwoTightenedFUOneRowsTwoBound [cert, upper] FINAL_manifest_v33.json :: results.M9_H2_tightened_fU1.rows.2.bound
\newcommand{\NSfourResultsMNineHTwoTightenedFUOneRowsTwoBound}{\ensuremath{1.665\cdot10^{-6}}}
% NSfourResultsMNineHTwoTightenedFUOneRowsTwoMargin [cert, margin] FINAL_manifest_v33.json :: results.M9_H2_tightened_fU1.rows.2.margin
\newcommand{\NSfourResultsMNineHTwoTightenedFUOneRowsTwoMargin}{\ensuremath{600.9}}
% NHOneAlphaThr [hand, exact] hand :: 
\newcommand{\NHOneAlphaThr}{\ensuremath{1.5\cdot10^{-4}}}
% NHOneSqrtLamThr [hand, exact] hand :: 
\newcommand{\NHOneSqrtLamThr}{\ensuremath{5\cdot10^{-8}}}
% NYbarStep [hand, exact] hand :: 
\newcommand{\NYbarStep}{\ensuremath{1\cdot10^{-12}}}
% NHTwoThrZero [hand, exact] hand :: 
\newcommand{\NHTwoThrZero}{\ensuremath{1\cdot10^{-8}}}
% NHTwoThrOne [hand, exact] hand :: 
\newcommand{\NHTwoThrOne}{\ensuremath{1\cdot10^{-6}}}
% NHTwoThrTwo [hand, exact] hand :: 
\newcommand{\NHTwoThrTwo}{\ensuremath{1\cdot10^{-3}}}
% NHThreeThrZero [hand, exact] hand :: 
\newcommand{\NHThreeThrZero}{\ensuremath{1\cdot10^{-6}}}
% NHThreeThrOne [hand, exact] hand :: 
\newcommand{\NHThreeThrOne}{\ensuremath{5\cdot10^{-6}}}
% NHThreeThrTwo [hand, exact] hand :: 
\newcommand{\NHThreeThrTwo}{\ensuremath{1\cdot10^{-5}}}
% NHFourThrZero [hand, exact] hand :: 
\newcommand{\NHFourThrZero}{\ensuremath{2\cdot10^{-7}}}
% NHFourThrOne [hand, exact] hand :: 
\newcommand{\NHFourThrOne}{\ensuremath{1\cdot10^{-5}}}
% NHFourThrTwo [hand, exact] hand :: 
\newcommand{\NHFourThrTwo}{\ensuremath{1\cdot10^{-5}}}
% NHFourThrThree [hand, exact] hand :: 
\newcommand{\NHFourThrThree}{\ensuremath{1\cdot10^{-5}}}
% NHFourKfourThr [hand, exact] hand :: 
\newcommand{\NHFourKfourThr}{\ensuremath{0.06}}
% NHSixConeSlopeThr [hand, exact] hand :: 
\newcommand{\NHSixConeSlopeThr}{\ensuremath{0.08}}
% NHSixConeKtwoThr [hand, exact] hand :: 
\newcommand{\NHSixConeKtwoThr}{\ensuremath{0.03}}
% NHSixConeNthreeThr [hand, exact] hand :: 
\newcommand{\NHSixConeNthreeThr}{\ensuremath{0.004}}
% NHSixConeJetThr [hand, exact] hand :: 
\newcommand{\NHSixConeJetThr}{\ensuremath{0.12}}
% NHSixConeHFiveThr [hand, exact] hand :: 
\newcommand{\NHSixConeHFiveThr}{\ensuremath{0.0144}}
% NHSixFarHFiveThr [hand, exact] hand :: 
\newcommand{\NHSixFarHFiveThr}{\ensuremath{0.0225}}
% NLemmaOConeM [hand, exact] hand :: 
\newcommand{\NLemmaOConeM}{\ensuremath{0.12}}
% NLemmaOConeK [hand, exact] hand :: 
\newcommand{\NLemmaOConeK}{\ensuremath{0.18}}
% NLemmaOFarM [hand, exact] hand :: 
\newcommand{\NLemmaOFarM}{\ensuremath{0.15}}
% NLemmaOFarK [hand, exact] hand :: 
\newcommand{\NLemmaOFarK}{\ensuremath{0.3}}
% NCoercivityConst [hand, exact] hand :: 
\newcommand{\NCoercivityConst}{\ensuremath{0.046875}}
% NMOnez [hand, exact] hand :: 
\newcommand{\NMOnez}{\ensuremath{0.534}}
% NBetaA [hand, exact] hand :: 
\newcommand{\NBetaA}{\ensuremath{2}}
% NTOneLower [cert, lower] t1clock.json :: t1.0
\newcommand{\NTOneLower}{\ensuremath{4.5978574775}}
% NTOneUpper [cert, upper] t1clock.json :: t1.1
\newcommand{\NTOneUpper}{\ensuremath{4.5978586144}}
% NLambdaOneLower [cert, lower] t1clock.json :: lnlam.0
\newcommand{\NLambdaOneLower}{\ensuremath{6.49\cdot10^{-5}}}
% NLambdaOneUpper [cert, upper] t1clock.json :: lnlam.1
\newcommand{\NLambdaOneUpper}{\ensuremath{6.75\cdot10^{-5}}}
% NTMinusTOneUpper [cert, upper] t1clock.json :: lnlam.1; 1.5e-4.sup_ab
\newcommand{\NTMinusTOneUpper}{\ensuremath{2.1\cdot10^{-17}}}
% NAbarUpper [cert, upper] final_alpha_points_5e-8.json :: 1.5e-4.sup_ab
\newcommand{\NAbarUpper}{\ensuremath{2.19103\cdot10^{-7}}}
% NBetabarUpper [cert, upper] final_alpha_points_5e-8.json :: 1.5e-4.sup_beta
\newcommand{\NBetabarUpper}{\ensuremath{3.33811\cdot10^{-9}}}
% NBaseKmaxLower [cert, lower] t1clock.json :: base_kmax.0
\newcommand{\NBaseKmaxLower}{\ensuremath{2.279\cdot10^{4}}}
% NBaseKmaxUpper [cert, upper] t1clock.json :: base_kmax.1
\newcommand{\NBaseKmaxUpper}{\ensuremath{2.376\cdot10^{4}}}
% NHOneAlphaBase [cert, upper] summary.json :: alpha_abs_max
\newcommand{\NHOneAlphaBase}{\ensuremath{1.1739\cdot10^{-4}}}
% NHOneAlphaErr [cert, upper] t9_h1_clam_s180.json :: C_alpha; inputs.YSTAR
\newcommand{\NHOneAlphaErr}{\ensuremath{7.21\cdot10^{-13}}}
% NHOneAlphaMargin [cert, margin] summary.json :: alpha_abs_max; C_alpha; inputs.YSTAR
\newcommand{\NHOneAlphaMargin}{\ensuremath{1.27}}
% NHOneSqrtLamBase [cert, upper] t9_h1_clam_s180.json :: sqrtLam_base_Wp
\newcommand{\NHOneSqrtLamBase}{\ensuremath{2.282\cdot10^{-8}}}
% NHOneSqrtLamTotal [cert, upper] t9_h1_clam_s180.json :: sqrtLam_total_Wp
\newcommand{\NHOneSqrtLamTotal}{\ensuremath{2.29\cdot10^{-8}}}
% NHOneSqrtLamMargin [cert, margin] t9_h1_clam_s180.json :: sqrtLam_total_Wp
\newcommand{\NHOneSqrtLamMargin}{\ensuremath{2.18}}
% NHOneYbarLimit [cert, lower] t9_h1_clam_s180.json :: Ybar_limit_H1
\newcommand{\NHOneYbarLimit}{\ensuremath{1.12\cdot10^{-9}}}
% NSwitchYUpper [cert, upper] switch_s2.json :: modulated.Y_raw
\newcommand{\NSwitchYUpper}{\ensuremath{6.21\cdot10^{-13}}}
% NSwitchYMargin [cert, margin] switch_s2.json :: modulated.Y_raw
\newcommand{\NSwitchYMargin}{\ensuremath{1.6}}
% NMinCI [cert, lower] constants.json :: meta.summary.min_c_I
\newcommand{\NMinCI}{\ensuremath{0.05283}}
% NNumCertificates [cert, exact] constants.json :: meta.summary.n_cert
\newcommand{\NNumCertificates}{\ensuremath{26}}
% NMinRhoC [cert, lower] constants.json :: meta.closure.rho_c_min
\newcommand{\NMinRhoC}{\ensuremath{0.04711}}
% NHTwoBaseZero [cert, upper] FINAL_manifest_v33.json :: results.M9_H2_tightened.rows.0.base
\newcommand{\NHTwoBaseZero}{\ensuremath{8.011\cdot10^{-11}}}
% NHTwoBaseOne [cert, upper] FINAL_manifest_v33.json :: results.M9_H2_tightened.rows.1.base
\newcommand{\NHTwoBaseOne}{\ensuremath{1.336\cdot10^{-10}}}
% NHTwoBaseTwo [cert, upper] FINAL_manifest_v33.json :: results.M9_H2_tightened.rows.2.base
\newcommand{\NHTwoBaseTwo}{\ensuremath{2.936\cdot10^{-10}}}
% NHTwoErrWZero [cert, upper] FINAL_manifest_v33.json :: results.M9_H2_tightened.rows.0.with_error
\newcommand{\NHTwoErrWZero}{\ensuremath{8.014\cdot10^{-11}}}
% NHTwoMarginWZero [cert, margin] FINAL_manifest_v33.json :: results.M9_H2_tightened.rows.0.margin
\newcommand{\NHTwoMarginWZero}{\ensuremath{124.78}}
% NHTwoErrWOne [cert, upper] FINAL_manifest_v33.json :: results.M9_H2_tightened.rows.1.with_error
\newcommand{\NHTwoErrWOne}{\ensuremath{4.706\cdot10^{-10}}}
% NHTwoMarginWOne [cert, margin] FINAL_manifest_v33.json :: results.M9_H2_tightened.rows.1.margin
\newcommand{\NHTwoMarginWOne}{\ensuremath{2124.9}}
% NHTwoErrWTwo [cert, upper] FINAL_manifest_v33.json :: results.M9_H2_tightened.rows.2.with_error
\newcommand{\NHTwoErrWTwo}{\ensuremath{1.559\cdot10^{-6}}}
% NHTwoMarginWTwo [cert, margin] FINAL_manifest_v33.json :: results.M9_H2_tightened.rows.2.margin
\newcommand{\NHTwoMarginWTwo}{\ensuremath{641.53}}
% NHThreeBaseZero [cert, upper] FINAL_manifest_v33.json :: results.M9_H2f.rows.0.base
\newcommand{\NHThreeBaseZero}{\ensuremath{3.127\cdot10^{-7}}}
% NHThreeBaseOne [cert, upper] FINAL_manifest_v33.json :: results.M9_H2f.rows.1.base
\newcommand{\NHThreeBaseOne}{\ensuremath{1.004\cdot10^{-6}}}
% NHThreeBaseTwo [cert, upper] FINAL_manifest_v33.json :: results.M9_H2f.rows.2.base
\newcommand{\NHThreeBaseTwo}{\ensuremath{2.431\cdot10^{-6}}}
% NHThreeErrZero [cert, upper] FINAL_manifest_v33.json :: results.M9_H2f.rows.0.with_error
\newcommand{\NHThreeErrZero}{\ensuremath{3.1363\cdot10^{-7}}}
% NHThreeMarginZero [cert, margin] FINAL_manifest_v33.json :: results.M9_H2f.rows.0.margin
\newcommand{\NHThreeMarginZero}{\ensuremath{3.1885}}
% NHThreeErrOne [cert, upper] FINAL_manifest_v33.json :: results.M9_H2f.rows.1.with_error
\newcommand{\NHThreeErrOne}{\ensuremath{1.0034\cdot10^{-6}}}
% NHThreeMarginOne [cert, margin] FINAL_manifest_v33.json :: results.M9_H2f.rows.1.margin
\newcommand{\NHThreeMarginOne}{\ensuremath{4.9834}}
% NHThreeErrTwo [cert, upper] FINAL_manifest_v33.json :: results.M9_H2f.rows.2.with_error
\newcommand{\NHThreeErrTwo}{\ensuremath{2.4351\cdot10^{-6}}}
% NHThreeMarginTwo [cert, margin] FINAL_manifest_v33.json :: results.M9_H2f.rows.2.margin
\newcommand{\NHThreeMarginTwo}{\ensuremath{4.1067}}
% NHFourBaseZero [cert, upper] FINAL_manifest_v33.json :: results.M9_base_sums.combined.windows_max.0
\newcommand{\NHFourBaseZero}{\ensuremath{3.48\cdot10^{-8}}}
% NHFourErrZero [cert, upper] FINAL_manifest_v33.json :: results.M9_H3_H2v32.rows.win.0.with_error_arb
\newcommand{\NHFourErrZero}{\ensuremath{3.5087\cdot10^{-8}}}
% NHFourMarginZero [cert, margin] FINAL_manifest_v33.json :: results.M9_H3_H2v32.rows.win.0.margin_arb
\newcommand{\NHFourMarginZero}{\ensuremath{5.700}}
% NHFourBaseOne [cert, upper] FINAL_manifest_v33.json :: results.M9_base_sums.combined.windows_max.1
\newcommand{\NHFourBaseOne}{\ensuremath{1.27\cdot10^{-7}}}
% NHFourErrOne [cert, upper] FINAL_manifest_v33.json :: results.M9_H3_H2v32.rows.win.1.with_error_arb
\newcommand{\NHFourErrOne}{\ensuremath{1.2616\cdot10^{-7}}}
% NHFourMarginOne [cert, margin] FINAL_manifest_v33.json :: results.M9_H3_H2v32.rows.win.1.margin_arb
\newcommand{\NHFourMarginOne}{\ensuremath{79.26}}
% NHFourBaseTwo [cert, upper] FINAL_manifest_v33.json :: results.M9_base_sums.combined.windows_max.2
\newcommand{\NHFourBaseTwo}{\ensuremath{3.76\cdot10^{-7}}}
% NHFourErrTwo [cert, upper] FINAL_manifest_v33.json :: results.M9_H3_H2v32.rows.win.2.with_error_arb
\newcommand{\NHFourErrTwo}{\ensuremath{3.7724\cdot10^{-7}}}
% NHFourMarginTwo [cert, margin] FINAL_manifest_v33.json :: results.M9_H3_H2v32.rows.win.2.margin_arb
\newcommand{\NHFourMarginTwo}{\ensuremath{26.50}}
% NHFourBaseThree [cert, upper] FINAL_manifest_v33.json :: results.M9_base_sums.combined.windows_max.3
\newcommand{\NHFourBaseThree}{\ensuremath{9.71\cdot10^{-7}}}
% NHFourErrThree [cert, upper] FINAL_manifest_v33.json :: results.M9_H3_H2v32.rows.win.3.with_error_arb
\newcommand{\NHFourErrThree}{\ensuremath{1.2822\cdot10^{-6}}}
% NHFourMarginThree [cert, margin] FINAL_manifest_v33.json :: results.M9_H3_H2v32.rows.win.3.margin_arb
\newcommand{\NHFourMarginThree}{\ensuremath{7.799}}
% NHFourYbarLimit [cert, lower] FINAL_manifest_v33.json :: results.M9_H3_H2v32.binding_Ybar_limit
\newcommand{\NHFourYbarLimit}{\ensuremath{1.575\cdot10^{-11}}}
% NHFourYbarMargin [cert, margin] FINAL_manifest_v33.json :: results.M9_H3_H2v32.binding_Ybar_limit
\newcommand{\NHFourYbarMargin}{\ensuremath{15.75}}
% NHFourKfourBase [cert, upper] amp_A1_w180.json :: inputs.BASE_T1.K4
\newcommand{\NHFourKfourBase}{\ensuremath{3.1633\cdot10^{-4}}}
% NHSixConeSlope [cert, upper] FINAL_manifest_v33.json :: results.cone_fixedL.checks.slope.value
\newcommand{\NHSixConeSlope}{\ensuremath{0.03029}}
% NHSixConeSlopeMargin [cert, margin] FINAL_manifest_v33.json :: results.cone_fixedL.checks.slope.margin
\newcommand{\NHSixConeSlopeMargin}{\ensuremath{2.641}}
% NHSixConeKtwo [cert, upper] FINAL_manifest_v33.json :: results.cone_fixedL.checks.K2.value
\newcommand{\NHSixConeKtwo}{\ensuremath{0.009876}}
% NHSixConeKtwoMargin [cert, margin] FINAL_manifest_v33.json :: results.cone_fixedL.checks.K2.margin
\newcommand{\NHSixConeKtwoMargin}{\ensuremath{3.037}}
% NHSixConeNthree [cert, upper] FINAL_manifest_v33.json :: results.cone_fixedL.checks.N3.value
\newcommand{\NHSixConeNthree}{\ensuremath{1.461\cdot10^{-4}}}
% NHSixConeNthreeMargin [cert, margin] FINAL_manifest_v33.json :: results.cone_fixedL.checks.N3.margin
\newcommand{\NHSixConeNthreeMargin}{\ensuremath{27.38}}
% NHSixConeHFive [cert, upper] FINAL_manifest_v33.json :: results.H4_fixedL.total
\newcommand{\NHSixConeHFive}{\ensuremath{1.2865\cdot10^{-5}}}
% NHSixConeHFiveMargin [cert, margin] FINAL_manifest_v33.json :: results.H4_fixedL.margin
\newcommand{\NHSixConeHFiveMargin}{\ensuremath{1119}}
% NHSixConeTstarNeeded [cert, upper] FINAL_manifest_v33.json :: results.H7_fixedL.ratio_upper
\newcommand{\NHSixConeTstarNeeded}{\ensuremath{1.6000015\cdot10^{-7}}}
% NHSixConeTstarLower [cert, lower] FINAL_manifest_v33.json :: results.H7_fixedL.t_star_lower
\newcommand{\NHSixConeTstarLower}{\ensuremath{3.68334\cdot10^{-7}}}
% NHSixConeTstarMargin [cert, margin] FINAL_manifest_v33.json :: results.H7_fixedL.margin_tstar_over_ratio
\newcommand{\NHSixConeTstarMargin}{\ensuremath{2.302}}
% NHSixConeHFiveRatioPeak [cert, upper] h4loc_A1_w180.json :: windows.constants/amp_A1_w180.json.worst.ratio
\newcommand{\NHSixConeHFiveRatioPeak}{\ensuremath{0.2393}}
% NHSixConeHOneSup [cert, upper] FINAL_manifest_v33.json :: results.H1_fixedL.sup_h0_max
\newcommand{\NHSixConeHOneSup}{\ensuremath{0.2163}}
% NHSixConeHOneH [cert, upper] FINAL_manifest_v33.json :: results.H1_fixedL.H_max
\newcommand{\NHSixConeHOneH}{\ensuremath{0.2185}}
% NHSixFarHFiveRatio [cert, upper] h4loc_v32.json :: windows.constants/amp_A1_w180.json.cylinders.18.ratio
\newcommand{\NHSixFarHFiveRatio}{\ensuremath{0.1761}}
% NHSixFarHFiveMarginOld [cert, margin] h4loc_v32.json :: windows.constants/amp_A1_w180.json.cylinders.18.ratio
\newcommand{\NHSixFarHFiveMarginOld}{\ensuremath{5.68}}
% NHSixFarHFiveMargin [cert, margin] FINAL_manifest_v33.json :: results.M9_H4_far.min_margin
\newcommand{\NHSixFarHFiveMargin}{\ensuremath{15.58}}
% NHSixConeHFiveMarginOld [cert, margin] h4loc_A1_w180.json :: windows.constants/amp_A1_w180.json.worst.ratio
\newcommand{\NHSixConeHFiveMarginOld}{\ensuremath{4.18}}
% NHSixFarKtwo [cert, upper] t9_far_v32.json :: h2
\newcommand{\NHSixFarKtwo}{\ensuremath{1.797\cdot10^{-4}}}
% NHSixFarKtwoThr [hand, exact] t9_far_v32.json :: thr.h2
\newcommand{\NHSixFarKtwoThr}{\ensuremath{3.53\cdot10^{-4}}}
% NHSixFarKtwoMargin [cert, margin] t9_far_v32.json :: margins.h2
\newcommand{\NHSixFarKtwoMargin}{\ensuremath{1.964}}
% NHSixFarKthree [cert, upper] t9_far_v32.json :: h3
\newcommand{\NHSixFarKthree}{\ensuremath{2.595\cdot10^{-8}}}
% NHSixFarKthreeThr [hand, exact] t9_far_v32.json :: thr.h3
\newcommand{\NHSixFarKthreeThr}{\ensuremath{1.17\cdot10^{-5}}}
% NHSixFarKthreeMargin [cert, margin] t9_far_v32.json :: margins.h3
\newcommand{\NHSixFarKthreeMargin}{\ensuremath{450.9}}
% NHSixFarKfour [cert, upper] t9_far_v32.json :: h4
\newcommand{\NHSixFarKfour}{\ensuremath{6.964\cdot10^{-9}}}
% NHSixFarKfourThr [hand, exact] t9_far_v32.json :: thr.h4
\newcommand{\NHSixFarKfourThr}{\ensuremath{1.03\cdot10^{-7}}}
% NHSixFarKfourMargin [cert, margin] t9_far_v32.json :: margins.h4
\newcommand{\NHSixFarKfourMargin}{\ensuremath{14.79}}
% NHSixFarSlope [cert, upper] t9_far_v32.json :: slope
\newcommand{\NHSixFarSlope}{\ensuremath{0.04153}}
% NHSixFarSlopeThr [hand, exact] t9_far_v32.json :: thr.slope
\newcommand{\NHSixFarSlopeThr}{\ensuremath{0.0800}}
% NHSixFarSlopeMargin [cert, margin] t9_far_v32.json :: margins.slope
\newcommand{\NHSixFarSlopeMargin}{\ensuremath{1.926}}
% NHSixFarNthree [cert, upper] t9_far_v32.json :: L2
\newcommand{\NHSixFarNthree}{\ensuremath{5.516\cdot10^{-7}}}
% NHSixFarNthreeThr [hand, exact] t9_far_v32.json :: thr.L2
\newcommand{\NHSixFarNthreeThr}{\ensuremath{3.32\cdot10^{-6}}}
% NHSixFarNthreeMargin [cert, margin] t9_far_v32.json :: margins.L2
\newcommand{\NHSixFarNthreeMargin}{\ensuremath{6.019}}
% NFarTstarMargin [cert, margin] lemmaO_v32_consts.json :: far_v32.case.margin_factor_lower
\newcommand{\NFarTstarMargin}{\ensuremath{15.5}}
% NFarTstarLower [cert, lower] lemmaO_v32_consts.json :: far_v32.case.t_star_lower
\newcommand{\NFarTstarLower}{\ensuremath{9.5108\cdot10^{-8}}}
% NFarTstarNeeded [cert, upper] lemmaO_v32_consts.json :: far_v32.case.needed
\newcommand{\NFarTstarNeeded}{\ensuremath{6.134\cdot10^{-9}}}
% NGammaModes [cert, exact] gamma0_final.json :: M
\newcommand{\NGammaModes}{\ensuremath{3762}}
% NGammaHalfwidth [cert, upper] gamma0_final.json :: b_j_halfwidth
\newcommand{\NGammaHalfwidth}{\ensuremath{8.138\cdot10^{-24}}}
% NGammaContraction [cert, upper] gamma0_final.json :: krawczyk_contraction
\newcommand{\NGammaContraction}{\ensuremath{3.01\cdot10^{-12}}}
% NGammaArcs [cert, exact] gamma0_final.json :: embedded.arcs
\newcommand{\NGammaArcs}{\ensuremath{701}}
% NGammaMaxTurn [cert, upper] gamma0_final.json :: embedded.maxturn
\newcommand{\NGammaMaxTurn}{\ensuremath{0.04997}}
% NGammaMaxPair [cert, upper] gamma0_final.json :: embedded.maxpair
\newcommand{\NGammaMaxPair}{\ensuremath{0.09597}}
% NGammaMinSep [cert, lower] gamma0_final.json :: embedded.minsep
\newcommand{\NGammaMinSep}{\ensuremath{8.89\cdot10^{-3}}}
% NGammaHalfGap [cert, lower] gamma0_final.json :: embedded.gap_cert
\newcommand{\NGammaHalfGap}{\ensuremath{0.0248219}}
% NPhaseISlabs [cert, exact] status.json :: slab
\newcommand{\NPhaseISlabs}{\ensuremath{96}}
% NPhaseIModes [cert, exact] status.json :: last_slab.M
\newcommand{\NPhaseIModes}{\ensuremath{2800}}
% NPhaseIKmaxLower [cert, lower] status.json :: kmax_enclosure.0
\newcommand{\NPhaseIKmaxLower}{\ensuremath{29.352874545}}
% NPhaseIKmaxUpper [cert, upper] status.json :: kmax_enclosure.1
\newcommand{\NPhaseIKmaxUpper}{\ensuremath{29.352874546}}
% NPhaseITLower [cert, lower] status.json :: t
\newcommand{\NPhaseITLower}{\ensuremath{4.597850539485830}}
% NPhaseITUpper [cert, upper] status.json :: t
\newcommand{\NPhaseITUpper}{\ensuremath{4.597850539485831}}
% NPeakSqrtLam [cert, upper] rate_table_fin12f4_a0.00015.json :: sqrtL_max
\newcommand{\NPeakSqrtLam}{\ensuremath{5.308\cdot10^{-8}}}
% NTailRate [cert, lower] rate_table_fin12f4_a0.00015.json :: tail.rho
\newcommand{\NTailRate}{\ensuremath{2.693\cdot10^{-3}}}
% NFloorSqrtLam [cert, upper] rate_table_fin12f4_a0.00015.json :: tail.sqrt_LF
\newcommand{\NFloorSqrtLam}{\ensuremath{9.872\cdot10^{-10}}}
% NAlphaFloor [cert, upper] rate_table_fin12f4_a0.00015.json :: tail.alpha_F
\newcommand{\NAlphaFloor}{\ensuremath{6.867\cdot10^{-7}}}
% NSupBetaVerifier [cert, upper] final_alpha_points_5e-8.json :: 1.5e-4.sup_beta
\newcommand{\NSupBetaVerifier}{\ensuremath{3.33811\cdot10^{-9}}}
% NMOneRamp [cert, upper] final_alpha_points_5e-8.json :: 1.5e-4.M1
\newcommand{\NMOneRamp}{\ensuremath{0.4328}}
% NSITwoPrime [cert, upper] FINAL_manifest_v33.json :: results.F12_sliver.sI2_prime
\newcommand{\NSITwoPrime}{\ensuremath{0.015457}}
% NLoopGainCoarse [cert, upper] U_cert_v35.json :: loop_gain.2.total
\newcommand{\NLoopGainCoarse}{\ensuremath{0.03904}}
% NRateMarginCoarse [cert, margin] U_cert_v35.json :: rate_margin
\newcommand{\NRateMarginCoarse}{\ensuremath{28.2}}
% NSOneEpsLower [cert, lower] bracket2.json :: eps_a
\newcommand{\NSOneEpsLower}{\ensuremath{1.5346198892965\cdot10^{-4}}}
% NSOneEpsUpper [cert, upper] bracket2.json :: eps_b
\newcommand{\NSOneEpsUpper}{\ensuremath{1.5346198899299\cdot10^{-4}}}
% NSOneRowsTip [cert, exact] ext_tip_check.json :: rows
\newcommand{\NSOneRowsTip}{\ensuremath{51246}}
% NSOneTipWidth [cert, upper] ext_tip_check.json :: max_pointwise_width.th
\newcommand{\NSOneTipWidth}{\ensuremath{3.2\cdot10^{-12}}}
% NSOneExtWidth [cert, upper] ext813_check.json :: max_pointwise_width.p
\newcommand{\NSOneExtWidth}{\ensuremath{5.7\cdot10^{-10}}}
% NSOneRowsMC [cert, exact] check_outA5.json :: prof_tip_vs_MC.rows_compared
\newcommand{\NSOneRowsMC}{\ensuremath{1025}}
% NSOneRowsMB [cert, exact] check_outA5.json :: ext813_vs_MB.rows_compared
\newcommand{\NSOneRowsMB}{\ensuremath{501}}
% NSOneMinMargin [cert, margin] check_outA5.json :: min:*.min_relative_margin.*
\newcommand{\NSOneMinMargin}{\ensuremath{1.2\cdot10^{-9}}}
% NTwoSegments [cert, exact] contour_result.json :: len:segments
\newcommand{\NTwoSegments}{\ensuremath{5058}}
% NTwoInteriorBoxes [cert, exact] interior_result.json :: len:
\newcommand{\NTwoInteriorBoxes}{\ensuremath{174}}
% NTwoContourBoxes [cert, exact] eta_cache.json :: len:
\newcommand{\NTwoContourBoxes}{\ensuremath{202}}
% NTwoEtaMax [cert, upper] eta_cache.json :: max:*.eta
\newcommand{\NTwoEtaMax}{\ensuremath{0.1163}}
% NTwoLipMax [cert, upper] eta_cache.json :: max:*.Lip
\newcommand{\NTwoLipMax}{\ensuremath{0.272}}
% NTwoLambdaZero [est, upper] cert_b0.25_th0.0_tau0.03.json :: Lam
\newcommand{\NTwoLambdaZero}{\ensuremath{1.790819}}
% NTwoLambdaOne [est, upper] cert_b0.25_th0.48_tau0.03.json :: Lam
\newcommand{\NTwoLambdaOne}{\ensuremath{1.768282}}
% NSOneDTightLower [doc, lower] doc:README.md :: D ∈ [2.50080627654, 2.50080627668]
\newcommand{\NSOneDTightLower}{\ensuremath{2.50080627654}}
% NSOneDTightUpper [doc, upper] doc:README.md :: D ∈ [2.50080627654, 2.50080627668]
\newcommand{\NSOneDTightUpper}{\ensuremath{2.50080627668}}
% NSOneKTightLower [doc, lower] doc:README.md :: k(0) ∈ [1.53982063762, 1.53982063822]
\newcommand{\NSOneKTightLower}{\ensuremath{1.53982063762}}
% NSOneKTightUpper [doc, upper] doc:README.md :: k(0) ∈ [1.53982063762, 1.53982063822]
\newcommand{\NSOneKTightUpper}{\ensuremath{1.53982063822}}
% NSOneDCondLower [doc, lower] doc:README.md :: D = −p(tip) ∈ [2.500806276619695, 2.500806276620049]
\newcommand{\NSOneDCondLower}{\ensuremath{2.500806276619695}}
% NSOneDCondUpper [doc, upper] doc:README.md :: D = −p(tip) ∈ [2.500806276619695, 2.500806276620049]
\newcommand{\NSOneDCondUpper}{\ensuremath{2.500806276620049}}
% NSOneKCondLower [doc, lower] doc:README.md :: k(0) ∈ [1.539820637871836, 1.5398206378721015]
\newcommand{\NSOneKCondLower}{\ensuremath{1.539820637871836}}
% NSOneKCondUpper [doc, upper] doc:README.md :: k(0) ∈ [1.539820637871836, 1.5398206378721015]
\newcommand{\NSOneKCondUpper}{\ensuremath{1.5398206378721015}}
% NSOneTauTipLower [doc, lower] doc:README.md :: τ_tip ∈ [10.249077651120778, 10.249077651121361]
\newcommand{\NSOneTauTipLower}{\ensuremath{10.249077651120778}}
% NSOneTauTipUpper [doc, upper] doc:README.md :: τ_tip ∈ [10.249077651120778, 10.249077651121361]
\newcommand{\NSOneTauTipUpper}{\ensuremath{10.249077651121361}}
% NSOneStateUncert [doc, upper] doc:README.md :: State uncertainty at σ = 8: 4.4e−13
\newcommand{\NSOneStateUncert}{\ensuremath{4.4\cdot10^{-13}}}
% NFarBoxEprime [doc, margin] doc:STEP4_RESULT.md :: (6.4287, 4.6917, 1.9940, 1.69999995)
\newcommand{\NFarBoxEprime}{\ensuremath{6.42}}
% NFarBoxXEsecond [doc, margin] doc:STEP4_RESULT.md :: (6.4287, 4.6917, 1.9940, 1.69999995)
\newcommand{\NFarBoxXEsecond}{\ensuremath{4.69}}
% NFarBoxXsqEthird [doc, margin] doc:STEP4_RESULT.md :: (6.4287, 4.6917, 1.9940, 1.69999995)
\newcommand{\NFarBoxXsqEthird}{\ensuremath{1.99}}
% NFarBoxE [doc, margin] doc:STEP4_RESULT.md :: (6.4287, 4.6917, 1.9940, 1.69999995)
\newcommand{\NFarBoxE}{\ensuremath{1.69}}
% NHFourKfourErr [doc, upper] doc:STEP4_RESULT.md :: | 5.60e-3 against 0.06 | ×10.7 |
\newcommand{\NHFourKfourErr}{\ensuremath{5.60\cdot10^{-3}}}
% NHFourKfourMargin [doc, margin] doc:STEP4_RESULT.md :: | 5.60e-3 against 0.06 | ×10.7 |
\newcommand{\NHFourKfourMargin}{\ensuremath{10.7}}
% NCoercivityCoefficient [doc, upper] doc:step3_proof_v32.tex :: \norm{e^{\tau\mathcal L}u}_\rho\le14.8\,e^{-0.0213\tau}\norm
\newcommand{\NCoercivityCoefficient}{\ensuremath{14.8}}
% NCoercivityRate [doc, lower] doc:step3_proof_v32.tex :: \norm{e^{\tau\mathcal L}u}_\rho\le14.8\,e^{-0.0213\tau}\norm
\newcommand{\NCoercivityRate}{\ensuremath{0.0213}}
% NRegionRampStart [hand, exact] hand :: 
\newcommand{\NRegionRampStart}{\ensuremath{14}}
% NRegionInterface [hand, exact] hand :: 
\newcommand{\NRegionInterface}{\ensuremath{28}}
% NRegionShellEnd [hand, exact] hand :: 
\newcommand{\NRegionShellEnd}{\ensuremath{60}}
% NRegionSegEnd [hand, exact] hand :: 
\newcommand{\NRegionSegEnd}{\ensuremath{350}}
% NRegionAnnEnd [hand, exact] hand :: 
\newcommand{\NRegionAnnEnd}{\ensuremath{700}}
% NZoneIIILow [hand, exact] hand :: 
\newcommand{\NZoneIIILow}{\ensuremath{253}}
% NZoneIIIHigh [hand, exact] hand :: 
\newcommand{\NZoneIIIHigh}{\ensuremath{379.5}}
% NConeCentreMin [hand, exact] hand :: 
\newcommand{\NConeCentreMin}{\ensuremath{300}}
% NConeCentreMax [hand, exact] hand :: 
\newcommand{\NConeCentreMax}{\ensuremath{380}}
% NFarCentreMin [hand, exact] hand :: 
\newcommand{\NFarCentreMin}{\ensuremath{490}}
% NConeL [hand, exact] hand :: 
\newcommand{\NConeL}{\ensuremath{50}}
% NConeX [hand, exact] hand :: 
\newcommand{\NConeX}{\ensuremath{3}}
% NFarL [hand, exact] hand :: 
\newcommand{\NFarL}{\ensuremath{113}}
% NFarX [hand, exact] hand :: 
\newcommand{\NFarX}{\ensuremath{2}}
 is found.
% ======================================================================

% ---- which document are we in ----------------------------------------
\ifdefined\ifsupplement\else\newif\ifsupplement\fi

% ---- packages ---------------------------------------------------------
\usepackage[T1]{fontenc}
\usepackage[utf8]{inputenc}
\usepackage{lmodern}
\usepackage[margin=1.1in]{geometry}
\usepackage{amsmath,amssymb,amsthm,mathtools,bm}
\usepackage{booktabs,array,longtable,tabularx,multirow}
\usepackage{enumitem}
\usepackage{graphicx}
\usepackage{xcolor}
\usepackage{tikz}
\usetikzlibrary{positioning,arrows.meta,calc}
% xr-hyper must come before hyperref
\usepackage{xr-hyper}
\usepackage[colorlinks=true,linkcolor=blue!50!black,citecolor=green!35!black,urlcolor=blue!50!black]{hyperref}

\raggedbottom
\setlength{\emergencystretch}{2em}
\setlist[enumerate]{leftmargin=*,itemsep=2pt,topsep=3pt}
\setlist[itemize]{leftmargin=*,itemsep=2pt,topsep=3pt}

% ---- numbering --------------------------------------------------------
% Main paper: sections 1,2,...; theorems 2.3 (by section); equations (2.3).
% Supplement: the Parts A..K are the \section's, numbered by letter, so
% theorems are A.1, E.4, ... (do not confuse "Theorem A" of the paper, which
% has no dot, with "Theorem A.1" of Part A).
\ifsupplement
  \renewcommand{\thesection}{\Alph{section}}
\fi
\numberwithin{equation}{section}

% ---- theorem environments (numbered by section, one shared counter) ----
\theoremstyle{plain}
\newtheorem{theorem}{Theorem}[section]
\newtheorem{proposition}[theorem]{Proposition}
\newtheorem{lemma}[theorem]{Lemma}
\newtheorem{corollary}[theorem]{Corollary}
\newtheorem{maintheorem}[theorem]{Main Theorem}
% Theorems A, B, C, D of the paper: own counter, lettered, no section prefix
\newtheorem{thmlet}{Theorem}
\renewcommand{\thethmlet}{\Alph{thmlet}}
\theoremstyle{definition}
\newtheorem{definition}[theorem]{Definition}
\newtheorem{convention}[theorem]{Convention}
\newtheorem{example}[theorem]{Example}
\theoremstyle{remark}
\newtheorem{remark}[theorem]{Remark}

% A list of hypotheses labelled (H1),(H2),... : \begin{hyplist}{H} ... \end{hyplist}
% (also {O}, {L}, {B}, {K}, {W} for the families of section 5.3 of OUTLINE.md).
\newenvironment{hyplist}[1]%
  {\begin{enumerate}[label=\textup{(#1\arabic*)},leftmargin=3em,itemsep=2pt]}%
  {\end{enumerate}}

% ---- trust classes (OUTLINE 4.1, 5.1.3) -------------------------------
% Supplement: explicit tags.  Main paper: certified is the default and
% unmarked; every non-certified number is marked "(floating point)".
\ifsupplement
  \newcommand{\tcert}{\textup{[cert]}}   % arb/outward-rounded, pinned script and key
  \newcommand{\test}{\textup{[est]}}     % floating point or estimate; never in a theorem
  \newcommand{\thand}{\textup{[hand]}}   % proved by hand in the text
  \newcommand{\tlit}{\textup{[lit]}}     % cited from the literature
\else
  \newcommand{\tcert}{}
  \newcommand{\test}{\unskip\ \textup{(floating point)}}
  \newcommand{\thand}{}
  \newcommand{\tlit}{}
\fi

% ---- gaps: visible, counted, an error in a final build -----------------
% Usage: \GAP{one sentence saying what is missing}   (and an entry in ISSUES.md)
\newcounter{gapcount}
\newif\iffinalpaper   % \finalpapertrue makes every \GAP a hard error
\newcommand{\GAP}[1]{\stepcounter{gapcount}\typeout{GAP-MARKER \thegapcount}%
  \iffinalpaper\PackageError{paper}{\string\GAP\space used in a final build}{Remove all gaps.}\fi
  {\color{red}\sffamily\bfseries[GAP~\thegapcount:\ }{\color{red}\sffamily #1}{\color{red}\sffamily\bfseries]}}
\AtEndDocument{\typeout{GAP-MARKER total: \thegapcount}}

% ---- file names and JSON keys (supplement only) -------------------------
% No internal path appears in the main paper; a use there is a warning.
\newcommand{\fn}[1]{%
  \ifsupplement\else\PackageWarning{paper}{\string\fn\space used in the main paper}\fi
  \texttt{\detokenize{#1}}}
\newcommand{\key}[1]{\texttt{\detokenize{#1}}}

% ---- cross references between main paper and supplement ----------------
% Labels carry a unit prefix (section 5.3 of OUTLINE.md and DRAFTER_GUIDE.md):
%   main paper:  intro: obj: s3: s4: s5: s6: s7: s8: s9:   (kind: thm- lem- eq- tab- fig- ...)
%   supplement:  pA: pB: ... pK:
% so all labels are globally unique and the prefixes below are only a second guard.
% \Xref{label}  : reference to the OTHER document (hyperlinked via xr-hyper)
% \Lref{label}  : reference in the SAME document (plain \ref)
\ifsupplement
  \externaldocument[M-]{main}[main.pdf]
  \newcommand{\Xref}[1]{\ref{M-#1}}
\else
  \externaldocument[S-]{supplement}[supplement.pdf]
  \newcommand{\Xref}[1]{\ref{S-#1}}
\fi
\newcommand{\Lref}[1]{\ref{#1}}
% \Supp{pE}{lem-chart} -> "Lemma E.3" style text is produced by the drafter:
%   "Lemma~\Xref{pE:lem-chart} of the supplement".  \SuppPart{E} gives "Part E".
\newcommand{\SuppPart}[1]{Part~#1 of the supplement}

% ---- numbers (generated by unit U0; macros only, no literal certified numbers) ----
\IfFileExists{numbers/numbers_paper.tex}%
  {% numbers_paper.tex -- GENERATED from ledger.json by build_ledger.py; do not edit.
% Rounding: upper bounds up, lower bounds and margins down at the printed digit; margin = threshold / bound.
% Use as $\times\NName$ for margins.  kind: cert certified, est estimate, hand hypothesis constant, doc only a document quotes it.
% NUEhatZero [cert, upper] U_cert_v35.json :: supersolution.EA.0
\newcommand{\NUEhatZero}{\ensuremath{1.2018\cdot10^{-6}}}
% NUEhatOne [cert, upper] U_cert_v35.json :: supersolution.EA.1
\newcommand{\NUEhatOne}{\ensuremath{1.0441\cdot10^{-4}}}
% NUEhatTwo [cert, upper] U_cert_v35.json :: supersolution.EA.2
\newcommand{\NUEhatTwo}{\ensuremath{0.056694}}
% NUBhatZero [cert, upper] U_cert_v35.json :: supersolution.Bhat.0
\newcommand{\NUBhatZero}{\ensuremath{4.8735\cdot10^{-6}}}
% NUBhatOne [cert, upper] U_cert_v35.json :: supersolution.Bhat.1
\newcommand{\NUBhatOne}{\ensuremath{5.5933\cdot10^{-4}}}
% NUBhatTwo [cert, upper] U_cert_v35.json :: supersolution.Bhat.2
\newcommand{\NUBhatTwo}{\ensuremath{0.36675}}
% NUSupersolutionCZeroZero [cert, upper] U_cert_v35.json :: supersolution.C0.0
\newcommand{\NUSupersolutionCZeroZero}{\ensuremath{1.1340\cdot10^{-6}}}
% NUSupersolutionCZeroOne [cert, upper] U_cert_v35.json :: supersolution.C0.1
\newcommand{\NUSupersolutionCZeroOne}{\ensuremath{1.6623\cdot10^{-4}}}
% NUSupersolutionCOneZero [cert, upper] U_cert_v35.json :: supersolution.C1.0
\newcommand{\NUSupersolutionCOneZero}{\ensuremath{1.5241\cdot10^{-9}}}
% NUSupersolutionCOneOne [cert, upper] U_cert_v35.json :: supersolution.C1.1
\newcommand{\NUSupersolutionCOneOne}{\ensuremath{5.5255\cdot10^{-8}}}
% NUSupersolutionNTwoZero [cert, upper] U_cert_v35.json :: supersolution.N2.0
\newcommand{\NUSupersolutionNTwoZero}{\ensuremath{0.071472}}
% NUSupersolutionNTwoOne [cert, upper] U_cert_v35.json :: supersolution.N2.1
\newcommand{\NUSupersolutionNTwoOne}{\ensuremath{0.33012}}
% NUSupersolutionNTwoTwo [cert, upper] U_cert_v35.json :: supersolution.N2.2
\newcommand{\NUSupersolutionNTwoTwo}{\ensuremath{5.0406}}
% NUSupersolutionNTwoOverFourNTwoRefZero [cert, upper] U_cert_v35.json :: supersolution.N2_over_4N2ref.0
\newcommand{\NUSupersolutionNTwoOverFourNTwoRefZero}{\ensuremath{0.18}}
% NUSupersolutionNTwoOverFourNTwoRefOne [cert, upper] U_cert_v35.json :: supersolution.N2_over_4N2ref.1
\newcommand{\NUSupersolutionNTwoOverFourNTwoRefOne}{\ensuremath{0.42}}
% NUSupersolutionNTwoOverFourNTwoRefTwo [cert, upper] U_cert_v35.json :: supersolution.N2_over_4N2ref.2
\newcommand{\NUSupersolutionNTwoOverFourNTwoRefTwo}{\ensuremath{0.26}}
% NUSupersolutionMOneZHypCheckZonesIIISOneTwoBhat [cert, upper] U_cert_v35.json :: supersolution.M1z_hyp_check.zones_I_II.S1_2Bhat
\newcommand{\NUSupersolutionMOneZHypCheckZonesIIISOneTwoBhat}{\ensuremath{0.28585}}
% NUSupersolutionMOneZHypCheckZoneIIIBOnePlusPUOne [cert, upper] U_cert_v35.json :: supersolution.M1z_hyp_check.zone_III.b1_plus_pU1
\newcommand{\NUSupersolutionMOneZHypCheckZoneIIIBOnePlusPUOne}{\ensuremath{0.50786}}
% NUMarginLateMOne [cert, margin] U_cert_v35.json :: margins.late_M1
\newcommand{\NUMarginLateMOne}{\ensuremath{2.353}}
% NUMarginLateMZero [cert, margin] U_cert_v35.json :: margins.late_M0
\newcommand{\NUMarginLateMZero}{\ensuremath{4.892}}
% NUMarginsLateMZeroZ [cert, margin] U_cert_v35.json :: margins.late_M0z
\newcommand{\NUMarginsLateMZeroZ}{\ensuremath{12.99}}
% NUMarginPresentMOne [cert, margin] U_cert_v35.json :: margins.present_M1_min.value
\newcommand{\NUMarginPresentMOne}{\ensuremath{2.770}}
% NUMarginPresentMZero [cert, margin] U_cert_v35.json :: margins.present_M0_min.value
\newcommand{\NUMarginPresentMZero}{\ensuremath{2.225}}
% NUMarginsPresentMZeroZMinValue [cert, margin] U_cert_v35.json :: margins.present_M0z_min.value
\newcommand{\NUMarginsPresentMZeroZMinValue}{\ensuremath{14.37}}
% NULoopGainZeroTotal [cert, upper] U_cert_v35.json :: loop_gain.0.total
\newcommand{\NULoopGainZeroTotal}{\ensuremath{1.23\cdot10^{-7}}}
% NULoopGainOneTotal [cert, upper] U_cert_v35.json :: loop_gain.1.total
\newcommand{\NULoopGainOneTotal}{\ensuremath{0.037169}}
% NULoopGainTwoNeighbour [cert, upper] U_cert_v35.json :: loop_gain.2.neighbour
\newcommand{\NULoopGainTwoNeighbour}{\ensuremath{0.037018}}
% NULoopGainTwoThroughA [cert, upper] U_cert_v35.json :: loop_gain.2.through_A
\newcommand{\NULoopGainTwoThroughA}{\ensuremath{0.0020172}}
% NUfineBhatZero [cert, upper] U_cert_v35.json :: ufine.Bhat.0
\newcommand{\NUfineBhatZero}{\ensuremath{4.8723\cdot10^{-6}}}
% NUfineBhatOne [cert, upper] U_cert_v35.json :: ufine.Bhat.1
\newcommand{\NUfineBhatOne}{\ensuremath{2.6266\cdot10^{-4}}}
% NUfineBhatTwo [cert, upper] U_cert_v35.json :: ufine.Bhat.2
\newcommand{\NUfineBhatTwo}{\ensuremath{0.025842}}
% NUcdeltaU [cert, upper] U_cert_v35.json :: launch.c_dU
\newcommand{\NUcdeltaU}{\ensuremath{1.7187\cdot10^{-8}}}
% NULaunchWmin [cert, lower] U_cert_v35.json :: launch.wmin
\newcommand{\NULaunchWmin}{\ensuremath{0.73492}}
% NULaunchCAD [cert, upper] U_cert_v35.json :: launch.CA_D
\newcommand{\NULaunchCAD}{\ensuremath{2.2429}}
% NULaunchQD [cert, upper] U_cert_v35.json :: launch.qD
\newcommand{\NULaunchQD}{\ensuremath{0.018509}}
% NULaunchDeltaInZero [cert, upper] U_cert_v35.json :: launch.Delta_in.0
\newcommand{\NULaunchDeltaInZero}{\ensuremath{3.131\cdot10^{-12}}}
% NULaunchDeltaInOne [cert, upper] U_cert_v35.json :: launch.Delta_in.1
\newcommand{\NULaunchDeltaInOne}{\ensuremath{3.287\cdot10^{-12}}}
% NULaunchDeltaInTwo [cert, upper] U_cert_v35.json :: launch.Delta_in.2
\newcommand{\NULaunchDeltaInTwo}{\ensuremath{3.451\cdot10^{-12}}}
% NULaunchPhiZero [cert, upper] U_cert_v35.json :: launch.Phi.0
\newcommand{\NULaunchPhiZero}{\ensuremath{6.434\cdot10^{-9}}}
% NULaunchPhiOne [cert, upper] U_cert_v35.json :: launch.Phi.1
\newcommand{\NULaunchPhiOne}{\ensuremath{1.4217\cdot10^{-8}}}
% NULaunchPhiTwo [cert, upper] U_cert_v35.json :: launch.Phi.2
\newcommand{\NULaunchPhiTwo}{\ensuremath{3.8404\cdot10^{-8}}}
% NUGammaZero [cert, upper] U_cert_v35.json :: Gamma.0
\newcommand{\NUGammaZero}{\ensuremath{4.0000358}}
% NUGammaOne [cert, upper] U_cert_v35.json :: Gamma.1
\newcommand{\NUGammaOne}{\ensuremath{16.000322}}
% NUGammaTwo [cert, upper] U_cert_v35.json :: Gamma.2
\newcommand{\NUGammaTwo}{\ensuremath{64.00218}}
% NUKappa [cert, upper] U_cert_v35.json :: kappa
\newcommand{\NUKappa}{\ensuremath{8.0000358}}
% NUKappaTwo [cert, upper] U_cert_v35.json :: kappa2
\newcommand{\NUKappaTwo}{\ensuremath{64.00061}}
% NUGhatZeroZero [cert, upper] U_cert_v35.json :: Ghat.0.0
\newcommand{\NUGhatZeroZero}{\ensuremath{51322.6}}
% NUGhatOneOne [cert, upper] U_cert_v35.json :: Ghat.1.1
\newcommand{\NUGhatOneOne}{\ensuremath{33962.9}}
% NUGhatTwoTwo [cert, upper] U_cert_v35.json :: Ghat.2.2
\newcommand{\NUGhatTwoTwo}{\ensuremath{348341.1}}
% NUGhatOneZero [cert, upper] U_cert_v35.json :: Ghat.1.0
\newcommand{\NUGhatOneZero}{\ensuremath{4.938}}
% NUGhatTwoZero [cert, upper] U_cert_v35.json :: Ghat.2.0
\newcommand{\NUGhatTwoZero}{\ensuremath{27.984}}
% NUGhatTwoOne [cert, upper] U_cert_v35.json :: Ghat.2.1
\newcommand{\NUGhatTwoOne}{\ensuremath{458.53}}
% NUKZeroZero [cert, upper] U_cert_v35.json :: K.0.0
\newcommand{\NUKZeroZero}{\ensuremath{8034.21}}
% NUKZeroOne [cert, upper] U_cert_v35.json :: K.0.1
\newcommand{\NUKZeroOne}{\ensuremath{724.89}}
% NUKOneZero [cert, upper] U_cert_v35.json :: K.1.0
\newcommand{\NUKOneZero}{\ensuremath{0.08152}}
% NUKOneOne [cert, upper] U_cert_v35.json :: K.1.1
\newcommand{\NUKOneOne}{\ensuremath{6336.94}}
% NUKOneTwo [cert, upper] U_cert_v35.json :: K.1.2
\newcommand{\NUKOneTwo}{\ensuremath{545.59}}
% NUKOneThree [cert, upper] U_cert_v35.json :: K.1.3
\newcommand{\NUKOneThree}{\ensuremath{0.58700}}
% NUKTwoZero [cert, upper] U_cert_v35.json :: K.2.0
\newcommand{\NUKTwoZero}{\ensuremath{0.87186}}
% NUKTwoOne [cert, upper] U_cert_v35.json :: K.2.1
\newcommand{\NUKTwoOne}{\ensuremath{528.62}}
% NUKTwoTwo [cert, upper] U_cert_v35.json :: K.2.2
\newcommand{\NUKTwoTwo}{\ensuremath{23425.98}}
% NUKTwoThree [cert, upper] U_cert_v35.json :: K.2.3
\newcommand{\NUKTwoThree}{\ensuremath{555.62}}
% NUFootwOne [cert, upper] U_cert_v35.json :: footw.1
\newcommand{\NUFootwOne}{\ensuremath{2.6361}}
% NUFootwTwo [cert, upper] U_cert_v35.json :: footw.2
\newcommand{\NUFootwTwo}{\ensuremath{2663.3}}
% NUCThreeThree [cert, upper] U_cert_v35.json :: c33
\newcommand{\NUCThreeThree}{\ensuremath{108.02}}
% NUCTwoTwoThree [cert, upper] U_cert_v35.json :: c223
\newcommand{\NUCTwoTwoThree}{\ensuremath{7.3284}}
% NUCEta [cert, upper] U_cert_v35.json :: C_eta
\newcommand{\NUCEta}{\ensuremath{13.376}}
% NULMax [cert, upper] U_cert_v35.json :: L_max
\newcommand{\NULMax}{\ensuremath{3.5352}}
% NURRatio [cert, upper] U_cert_v35.json :: R_ratio
\newcommand{\NURRatio}{\ensuremath{11.593}}
% NUSweepsSigmaInCoverSup [cert, upper] U_cert_v35.json :: sweeps.Sigma_in.cover_sup
\newcommand{\NUSweepsSigmaInCoverSup}{\ensuremath{5.0892}}
% NUSweepsSigmaECoverSup [cert, upper] U_cert_v35.json :: sweeps.Sigma_E.cover_sup
\newcommand{\NUSweepsSigmaECoverSup}{\ensuremath{0.30846}}
% NUSweepsSXCoverSup [cert, upper] U_cert_v35.json :: sweeps.S_X.cover_sup
\newcommand{\NUSweepsSXCoverSup}{\ensuremath{1.5243}}
% NUSweepsSigmaXCoverSup [cert, upper] U_cert_v35.json :: sweeps.Sigma_X.cover_sup
\newcommand{\NUSweepsSigmaXCoverSup}{\ensuremath{3.0081}}
% NUSweepsSigmaRsTwoCoverSup [cert, upper] U_cert_v35.json :: sweeps.Sigma_rs.2.cover_sup
\newcommand{\NUSweepsSigmaRsTwoCoverSup}{\ensuremath{4.3350\cdot10^{-3}}}
% NUDataTermsLateDZero [cert, upper] U_cert_v35.json :: data_terms_late.D.0
\newcommand{\NUDataTermsLateDZero}{\ensuremath{1.7330\cdot10^{-6}}}
% NUDataTermsLateDOne [cert, upper] U_cert_v35.json :: data_terms_late.D.1
\newcommand{\NUDataTermsLateDOne}{\ensuremath{2.5353\cdot10^{-4}}}
% NUDataTermsLateSlopeShareCOneSigmaXOverDkZero [cert, upper] U_cert_v35.json :: data_terms_late.slope_share_C1SigmaX_over_Dk.0
\newcommand{\NUDataTermsLateSlopeShareCOneSigmaXOverDkZero}{\ensuremath{2.7\cdot10^{-3}}}
% NUDataTermsLateSlopeShareCOneSigmaXOverDkOne [cert, upper] U_cert_v35.json :: data_terms_late.slope_share_C1SigmaX_over_Dk.1
\newcommand{\NUDataTermsLateSlopeShareCOneSigmaXOverDkOne}{\ensuremath{6.6\cdot10^{-4}}}
% NURLtTwoMarginZero [cert, margin] U_cert_v35.json :: R_lt2_margin.0
\newcommand{\NURLtTwoMarginZero}{\ensuremath{4.432}}
% NURLtTwoMarginOne [cert, margin] U_cert_v35.json :: R_lt2_margin.1
\newcommand{\NURLtTwoMarginOne}{\ensuremath{6.697}}
% NUThetaTwoEight [cert, upper] U_cert_v35.json :: theta28
\newcommand{\NUThetaTwoEight}{\ensuremath{0.37011}}
% NUDFiveSix [cert, upper] U_cert_v35.json :: D56
\newcommand{\NUDFiveSix}{\ensuremath{1.0368}}
% NUDTwoEight [cert, upper] U_cert_v35.json :: D28
\newcommand{\NUDTwoEight}{\ensuremath{1.5876}}
% NUSThreeThreeFootOverBhatTwo [cert, upper] U_cert_v35.json :: S33_foot_over_Bhat2
\newcommand{\NUSThreeThreeFootOverBhatTwo}{\ensuremath{1.69\cdot10^{-3}}}
% NSAHJWKappaJ [cert, upper] A_interface.json :: HJW.kappa_J
\newcommand{\NSAHJWKappaJ}{\ensuremath{0.042351}}
% NSAGpZero [cert, lower] A_interface.json :: res.tight.gp.0
\newcommand{\NSAGpZero}{\ensuremath{0.24292}}
% NSAGpOne [cert, lower] A_interface.json :: res.tight.gp.1
\newcommand{\NSAGpOne}{\ensuremath{0.24106}}
% NSAGpTwo [cert, lower] A_interface.json :: res.tight.gp.2
\newcommand{\NSAGpTwo}{\ensuremath{0.11961}}
% NSAAminC [cert, lower] A_interface.json :: res.tight.amin_c
\newcommand{\NSAAminC}{\ensuremath{0.97623}}
% NSARhoJ [cert, upper] A_interface.json :: res.tight.rhoJ
\newcommand{\NSARhoJ}{\ensuremath{5.4995\cdot10^{-7}}}
% NSAKappaOTwoTwo [cert, upper] A_interface.json :: res.tight.kappa_o.2.2
\newcommand{\NSAKappaOTwoTwo}{\ensuremath{1.2416\cdot10^{-5}}}
% NSADsrcZero [cert, upper] A_interface.json :: res.tight.Dsrc.0
\newcommand{\NSADsrcZero}{\ensuremath{1.8092\cdot10^{-7}}}
% NSADsrcOne [cert, upper] A_interface.json :: res.tight.Dsrc.1
\newcommand{\NSADsrcOne}{\ensuremath{9.8276\cdot10^{-5}}}
% NSABZeroDWindowSlopeZero [cert, upper] A_interface.json :: res.tight.B0.D_window_slope.0
\newcommand{\NSABZeroDWindowSlopeZero}{\ensuremath{8.7399\cdot10^{-10}}}
% NSABZeroDWindowSlopeOne [cert, upper] A_interface.json :: res.tight.B0.D_window_slope.1
\newcommand{\NSABZeroDWindowSlopeOne}{\ensuremath{5.2599\cdot10^{-8}}}
% NSAFTwoZero [cert, upper] A_interface.json :: res.tight.F2.0
\newcommand{\NSAFTwoZero}{\ensuremath{5.6506\cdot10^{-13}}}
% NSAFTwoOne [cert, upper] A_interface.json :: res.tight.F2.1
\newcommand{\NSAFTwoOne}{\ensuremath{5.9320\cdot10^{-13}}}
% NSACoupOne [cert, upper] A_interface.json :: res.tight.coup.1
\newcommand{\NSACoupOne}{\ensuremath{1.1568\cdot10^{-12}}}
% NSACOneZeroZero [cert, upper] A_interface.json :: res.tight.C1_0.0
\newcommand{\NSACOneZeroZero}{\ensuremath{7.2060\cdot10^{-10}}}
% NSACOneZeroOne [cert, upper] A_interface.json :: res.tight.C1_0.1
\newcommand{\NSACOneZeroOne}{\ensuremath{4.3143\cdot10^{-8}}}
% NSABZeroSupCoreZero [cert, upper] A_interface.json :: res.tight.B0.sup_core.0
\newcommand{\NSABZeroSupCoreZero}{\ensuremath{4.23\cdot10^{-5}}}
% NSACoreFluxMarginTwoFLUXZeroOverKZero [cert, margin] A_interface.json :: res.tight.core.flux_margin.2.FLUX0_over_K0
\newcommand{\NSACoreFluxMarginTwoFLUXZeroOverKZero}{\ensuremath{3.396}}
% NSACoreFluxMarginTwoFLUXTwoOverKTwo [cert, margin] A_interface.json :: res.tight.core.flux_margin.2.FLUX2_over_K2
\newcommand{\NSACoreFluxMarginTwoFLUXTwoOverKTwo}{\ensuremath{6.012}}
% NSAMarginFluxZero [cert, margin] asm_design_N2ref4.json :: res.tight.flux.beta.2.margin_FLUX0
\newcommand{\NSAMarginFluxZero}{\ensuremath{3.365}}
% NSAMarginFluxTwo [cert, margin] asm_design_N2ref4.json :: res.tight.flux.beta.2.margin_FLUX2
\newcommand{\NSAMarginFluxTwo}{\ensuremath{5.955}}
% NSAfourSZero [cert, upper] asm_design_N2ref4.json :: res.tight.S0
\newcommand{\NSAfourSZero}{\ensuremath{7.1878\cdot10^{-3}}}
% NSAfourSOne [cert, upper] asm_design_N2ref4.json :: res.tight.S1
\newcommand{\NSAfourSOne}{\ensuremath{0.10430}}
% NSAMarginCoreSZeroV [cert, margin] asm_design_N2ref4.json :: res.tight.margins.S0v_vs_0.0153
\newcommand{\NSAMarginCoreSZeroV}{\ensuremath{2.116}}
% NSAMarginCoreSOneV [cert, margin] asm_design_N2ref4.json :: res.tight.margins.S1v_vs_0.6
\newcommand{\NSAMarginCoreSOneV}{\ensuremath{5.748}}
% NSAMarginCoreMOne [cert, margin] asm_design_N2ref4.json :: res.tight.margins.M1core
\newcommand{\NSAMarginCoreMOne}{\ensuremath{4.409}}
% NSAMarginCoreMZero [cert, margin] asm_design_N2ref4.json :: res.tight.margins.M0core
\newcommand{\NSAMarginCoreMZero}{\ensuremath{5.450}}
% NSAMarginGhat [cert, margin] asm_design_N2ref4.json :: res.tight.margins.Ghat22
\newcommand{\NSAMarginGhat}{\ensuremath{1.917}}
% NSAMarginRampMOne [cert, margin] asm_design_N2ref4.json :: res.tight.margins.ramp M1
\newcommand{\NSAMarginRampMOne}{\ensuremath{1.184}}
% NSAfourInfoZeroDWindowConst [cert, upper] asm_design_N2ref4.json :: res.tight.info.0.D_window_const
\newcommand{\NSAfourInfoZeroDWindowConst}{\ensuremath{1.7732\cdot10^{-7}}}
% NSAfourInfoOneDWindowConst [cert, upper] asm_design_N2ref4.json :: res.tight.info.1.D_window_const
\newcommand{\NSAfourInfoOneDWindowConst}{\ensuremath{9.8058\cdot10^{-5}}}
% NZDeltaRZero [cert, upper] zone3_interface.json :: DeltaR.0
\newcommand{\NZDeltaRZero}{\ensuremath{4.4986\cdot10^{-8}}}
% NZDeltaROne [cert, upper] zone3_interface.json :: DeltaR.1
\newcommand{\NZDeltaROne}{\ensuremath{6.7972\cdot10^{-6}}}
% NZDeltaRTwo [cert, upper] zone3_interface.json :: DeltaR.2
\newcommand{\NZDeltaRTwo}{\ensuremath{0.011480}}
% NZDeltaRBZero [cert, upper] zone3_interface.json :: DeltaR_b.0
\newcommand{\NZDeltaRBZero}{\ensuremath{2.66\cdot10^{-12}}}
% NZDeltaRBOne [cert, upper] zone3_interface.json :: DeltaR_b.1
\newcommand{\NZDeltaRBOne}{\ensuremath{3.26\cdot10^{-9}}}
% NZDeltaRBTwo [cert, upper] zone3_interface.json :: DeltaR_b.2
\newcommand{\NZDeltaRBTwo}{\ensuremath{4.35\cdot10^{-8}}}
% NZDeltaRDetailNSqZero [cert, upper] zone3_interface.json :: DeltaR_detail.N_sq.0
\newcommand{\NZDeltaRDetailNSqZero}{\ensuremath{8.28\cdot10^{-5}}}
% NZDeltaRDetailNSqOne [cert, upper] zone3_interface.json :: DeltaR_detail.N_sq.1
\newcommand{\NZDeltaRDetailNSqOne}{\ensuremath{0.10149}}
% NZDeltaRDetailNSqTwo [cert, upper] zone3_interface.json :: DeltaR_detail.N_sq.2
\newcommand{\NZDeltaRDetailNSqTwo}{\ensuremath{1.3567}}
% NZDeltaRDetailNSqThree [cert, upper] zone3_interface.json :: DeltaR_detail.N_sq.3
\newcommand{\NZDeltaRDetailNSqThree}{\ensuremath{25350}}
% NZDeltaRDetailPartsTwoPairsThree [cert, upper] zone3_interface.json :: DeltaR_detail.parts.2.pairs.3
\newcommand{\NZDeltaRDetailPartsTwoPairsThree}{\ensuremath{0.011298}}
% NZDeltaRDetailPartsTwoLFourFoot [cert, upper] zone3_interface.json :: DeltaR_detail.parts.2.L4_foot
\newcommand{\NZDeltaRDetailPartsTwoLFourFoot}{\ensuremath{1.6\cdot10^{-4}}}
% NZDeltaRDetailPartsTwoCubic [cert, upper] zone3_interface.json :: DeltaR_detail.parts.2.cubic
\newcommand{\NZDeltaRDetailPartsTwoCubic}{\ensuremath{9.3\cdot10^{-6}}}
% NZDeltaRDetailPartsTwoSThreeThreeFoot [cert, upper] zone3_interface.json :: DeltaR_detail.parts.2.S33_foot
\newcommand{\NZDeltaRDetailPartsTwoSThreeThreeFoot}{\ensuremath{3.7\cdot10^{-7}}}
% NZNIIIZero [cert, upper] zone3_interface.json :: n_III.0
\newcommand{\NZNIIIZero}{\ensuremath{0.071445}}
% NZNIIIOne [cert, upper] zone3_interface.json :: n_III.1
\newcommand{\NZNIIIOne}{\ensuremath{0.32707}}
% NZNIIITwo [cert, upper] zone3_interface.json :: n_III.2
\newcommand{\NZNIIITwo}{\ensuremath{3.0445}}
% NZTauIII [cert, upper] zone3_interface.json :: tau_III
\newcommand{\NZTauIII}{\ensuremath{4.0708}}
% NZoneIIIBoxZero [cert, upper] zone3_interface.json :: box_III_380.b0
\newcommand{\NZoneIIIBoxZero}{\ensuremath{0.10032}}
% NZoneIIIBoxOne [cert, upper] zone3_interface.json :: box_III_380.b1
\newcommand{\NZoneIIIBoxOne}{\ensuremath{0.41974}}
% NZoneIIIBoxZ [cert, upper] zone3_interface.json :: box_III_380.bz
\newcommand{\NZoneIIIBoxZ}{\ensuremath{0.014291}}
% NZBoxIIIThreeEightZeroMarginsBZero [cert, margin] zone3_interface.json :: box_III_380.margins.b0
\newcommand{\NZBoxIIIThreeEightZeroMarginsBZero}{\ensuremath{1.495}}
% NZoneIIIMarginThinnest [cert, margin] zone3_interface.json :: box_III_380.margins.b1
\newcommand{\NZoneIIIMarginThinnest}{\ensuremath{1.238}}
% NZBoxIIIThreeEightZeroMarginsBz [cert, margin] zone3_interface.json :: box_III_380.margins.bz
\newcommand{\NZBoxIIIThreeEightZeroMarginsBz}{\ensuremath{11.19}}
% NZBoxIIISevenZeroZeroBZero [cert, upper] zone3_interface.json :: box_III_700.b0
\newcommand{\NZBoxIIISevenZeroZeroBZero}{\ensuremath{0.22757}}
% NZBoxIIISevenZeroZeroBOne [cert, upper] zone3_interface.json :: box_III_700.b1
\newcommand{\NZBoxIIISevenZeroZeroBOne}{\ensuremath{0.54234}}
% NZUfineKTwoAtBhatRef [cert, upper] zone3_interface.json :: ufine.K2_at_Bhat_ref
\newcommand{\NZUfineKTwoAtBhatRef}{\ensuremath{0.046956}}
% NZUfineKThreeStripAtBhatRef [cert, upper] zone3_interface.json :: ufine.K3_strip_at_Bhat_ref
\newcommand{\NZUfineKThreeStripAtBhatRef}{\ensuremath{0.054140}}
% NRecoveryMinMargin [cert, margin] zone3_interface.json :: cone_family.H6.at_Bhat_ref.min_margin
\newcommand{\NRecoveryMinMargin}{\ensuremath{1.860}}
% NZConeFamilyHSixAtBhatRefConeSlopeMargin [cert, margin] zone3_interface.json :: cone_family.H6.at_Bhat_ref.cone.slope.margin
\newcommand{\NZConeFamilyHSixAtBhatRefConeSlopeMargin}{\ensuremath{2.46}}
% NZConeFamilyHSixAtBhatRefConeKTwoMargin [cert, margin] zone3_interface.json :: cone_family.H6.at_Bhat_ref.cone.K2.margin
\newcommand{\NZConeFamilyHSixAtBhatRefConeKTwoMargin}{\ensuremath{2.54}}
% NZConeFamilyHSixAtBhatRefConeKThreeMargin [cert, margin] zone3_interface.json :: cone_family.H6.at_Bhat_ref.cone.K3.margin
\newcommand{\NZConeFamilyHSixAtBhatRefConeKThreeMargin}{\ensuremath{2.51}}
% NConeTstarMargin [cert, margin] zone3_interface.json :: acceptance.tstar_margin
\newcommand{\NConeTstarMargin}{\ensuremath{2.30}}
% NUoneTFiveYaddLevelOneTwo [cert, upper] u1_ghat.json :: t5_Yadd.level1.2
\newcommand{\NUoneTFiveYaddLevelOneTwo}{\ensuremath{0.0044}}
% NUoneTFiveYaddLevelTwoTwo [cert, upper] u1_ghat.json :: t5_Yadd.level2.2
\newcommand{\NUoneTFiveYaddLevelTwoTwo}{\ensuremath{0.1081}}
% NUoneTFiveYaddLevelTwoThree [cert, upper] u1_ghat.json :: t5_Yadd.level2.3
\newcommand{\NUoneTFiveYaddLevelTwoThree}{\ensuremath{0.374}}
% NUoneWeightThetaGammaZero [cert, upper] u1_ghat.json :: weight.Theta.Gamma.0
\newcommand{\NUoneWeightThetaGammaZero}{\ensuremath{4.00000000000000000000054}}
% NRKZeroZero [cert, upper] ramp_U.json :: K.0.0
\newcommand{\NRKZeroZero}{\ensuremath{13436.03}}
% NRKOneOne [cert, upper] ramp_U.json :: K.1.1
\newcommand{\NRKOneOne}{\ensuremath{10526.19}}
% NRKTwoTwo [cert, upper] ramp_U.json :: K.2.2
\newcommand{\NRKTwoTwo}{\ensuremath{39482.30}}
% NUoneCasesPrimaryYNextZero [cert, upper] u1_ghat.json :: cases.primary.Y_next.0
\newcommand{\NUoneCasesPrimaryYNextZero}{\ensuremath{457.092}}
% NUoneCasesPrimaryYNextOne [cert, upper] u1_ghat.json :: cases.primary.Y_next.1
\newcommand{\NUoneCasesPrimaryYNextOne}{\ensuremath{356.940}}
% NUoneCasesPrimaryYNextTwo [cert, upper] u1_ghat.json :: cases.primary.Y_next.2
\newcommand{\NUoneCasesPrimaryYNextTwo}{\ensuremath{1320.220}}
% NSfourResultsMNineHTwoTightenedFUOneRowsZeroBound [cert, upper] FINAL_manifest_v33.json :: results.M9_H2_tightened_fU1.rows.0.bound
\newcommand{\NSfourResultsMNineHTwoTightenedFUOneRowsZeroBound}{\ensuremath{3.658\cdot10^{-10}}}
% NSfourResultsMNineHTwoTightenedFUOneRowsZeroMargin [cert, margin] FINAL_manifest_v33.json :: results.M9_H2_tightened_fU1.rows.0.margin
\newcommand{\NSfourResultsMNineHTwoTightenedFUOneRowsZeroMargin}{\ensuremath{27.34}}
% NSfourResultsMNineHTwoTightenedFUOneRowsOneBound [cert, upper] FINAL_manifest_v33.json :: results.M9_H2_tightened_fU1.rows.1.bound
\newcommand{\NSfourResultsMNineHTwoTightenedFUOneRowsOneBound}{\ensuremath{1.777\cdot10^{-9}}}
% NSfourResultsMNineHTwoTightenedFUOneRowsOneMargin [cert, margin] FINAL_manifest_v33.json :: results.M9_H2_tightened_fU1.rows.1.margin
\newcommand{\NSfourResultsMNineHTwoTightenedFUOneRowsOneMargin}{\ensuremath{562.8}}
% NSfourResultsMNineHTwoTightenedFUOneRowsTwoBound [cert, upper] FINAL_manifest_v33.json :: results.M9_H2_tightened_fU1.rows.2.bound
\newcommand{\NSfourResultsMNineHTwoTightenedFUOneRowsTwoBound}{\ensuremath{1.665\cdot10^{-6}}}
% NSfourResultsMNineHTwoTightenedFUOneRowsTwoMargin [cert, margin] FINAL_manifest_v33.json :: results.M9_H2_tightened_fU1.rows.2.margin
\newcommand{\NSfourResultsMNineHTwoTightenedFUOneRowsTwoMargin}{\ensuremath{600.9}}
% NHOneAlphaThr [hand, exact] hand :: 
\newcommand{\NHOneAlphaThr}{\ensuremath{1.5\cdot10^{-4}}}
% NHOneSqrtLamThr [hand, exact] hand :: 
\newcommand{\NHOneSqrtLamThr}{\ensuremath{5\cdot10^{-8}}}
% NYbarStep [hand, exact] hand :: 
\newcommand{\NYbarStep}{\ensuremath{1\cdot10^{-12}}}
% NHTwoThrZero [hand, exact] hand :: 
\newcommand{\NHTwoThrZero}{\ensuremath{1\cdot10^{-8}}}
% NHTwoThrOne [hand, exact] hand :: 
\newcommand{\NHTwoThrOne}{\ensuremath{1\cdot10^{-6}}}
% NHTwoThrTwo [hand, exact] hand :: 
\newcommand{\NHTwoThrTwo}{\ensuremath{1\cdot10^{-3}}}
% NHThreeThrZero [hand, exact] hand :: 
\newcommand{\NHThreeThrZero}{\ensuremath{1\cdot10^{-6}}}
% NHThreeThrOne [hand, exact] hand :: 
\newcommand{\NHThreeThrOne}{\ensuremath{5\cdot10^{-6}}}
% NHThreeThrTwo [hand, exact] hand :: 
\newcommand{\NHThreeThrTwo}{\ensuremath{1\cdot10^{-5}}}
% NHFourThrZero [hand, exact] hand :: 
\newcommand{\NHFourThrZero}{\ensuremath{2\cdot10^{-7}}}
% NHFourThrOne [hand, exact] hand :: 
\newcommand{\NHFourThrOne}{\ensuremath{1\cdot10^{-5}}}
% NHFourThrTwo [hand, exact] hand :: 
\newcommand{\NHFourThrTwo}{\ensuremath{1\cdot10^{-5}}}
% NHFourThrThree [hand, exact] hand :: 
\newcommand{\NHFourThrThree}{\ensuremath{1\cdot10^{-5}}}
% NHFourKfourThr [hand, exact] hand :: 
\newcommand{\NHFourKfourThr}{\ensuremath{0.06}}
% NHSixConeSlopeThr [hand, exact] hand :: 
\newcommand{\NHSixConeSlopeThr}{\ensuremath{0.08}}
% NHSixConeKtwoThr [hand, exact] hand :: 
\newcommand{\NHSixConeKtwoThr}{\ensuremath{0.03}}
% NHSixConeNthreeThr [hand, exact] hand :: 
\newcommand{\NHSixConeNthreeThr}{\ensuremath{0.004}}
% NHSixConeJetThr [hand, exact] hand :: 
\newcommand{\NHSixConeJetThr}{\ensuremath{0.12}}
% NHSixConeHFiveThr [hand, exact] hand :: 
\newcommand{\NHSixConeHFiveThr}{\ensuremath{0.0144}}
% NHSixFarHFiveThr [hand, exact] hand :: 
\newcommand{\NHSixFarHFiveThr}{\ensuremath{0.0225}}
% NLemmaOConeM [hand, exact] hand :: 
\newcommand{\NLemmaOConeM}{\ensuremath{0.12}}
% NLemmaOConeK [hand, exact] hand :: 
\newcommand{\NLemmaOConeK}{\ensuremath{0.18}}
% NLemmaOFarM [hand, exact] hand :: 
\newcommand{\NLemmaOFarM}{\ensuremath{0.15}}
% NLemmaOFarK [hand, exact] hand :: 
\newcommand{\NLemmaOFarK}{\ensuremath{0.3}}
% NCoercivityConst [hand, exact] hand :: 
\newcommand{\NCoercivityConst}{\ensuremath{0.046875}}
% NMOnez [hand, exact] hand :: 
\newcommand{\NMOnez}{\ensuremath{0.534}}
% NBetaA [hand, exact] hand :: 
\newcommand{\NBetaA}{\ensuremath{2}}
% NTOneLower [cert, lower] t1clock.json :: t1.0
\newcommand{\NTOneLower}{\ensuremath{4.5978574775}}
% NTOneUpper [cert, upper] t1clock.json :: t1.1
\newcommand{\NTOneUpper}{\ensuremath{4.5978586144}}
% NLambdaOneLower [cert, lower] t1clock.json :: lnlam.0
\newcommand{\NLambdaOneLower}{\ensuremath{6.49\cdot10^{-5}}}
% NLambdaOneUpper [cert, upper] t1clock.json :: lnlam.1
\newcommand{\NLambdaOneUpper}{\ensuremath{6.75\cdot10^{-5}}}
% NTMinusTOneUpper [cert, upper] t1clock.json :: lnlam.1; 1.5e-4.sup_ab
\newcommand{\NTMinusTOneUpper}{\ensuremath{2.1\cdot10^{-17}}}
% NAbarUpper [cert, upper] final_alpha_points_5e-8.json :: 1.5e-4.sup_ab
\newcommand{\NAbarUpper}{\ensuremath{2.19103\cdot10^{-7}}}
% NBetabarUpper [cert, upper] final_alpha_points_5e-8.json :: 1.5e-4.sup_beta
\newcommand{\NBetabarUpper}{\ensuremath{3.33811\cdot10^{-9}}}
% NBaseKmaxLower [cert, lower] t1clock.json :: base_kmax.0
\newcommand{\NBaseKmaxLower}{\ensuremath{2.279\cdot10^{4}}}
% NBaseKmaxUpper [cert, upper] t1clock.json :: base_kmax.1
\newcommand{\NBaseKmaxUpper}{\ensuremath{2.376\cdot10^{4}}}
% NHOneAlphaBase [cert, upper] summary.json :: alpha_abs_max
\newcommand{\NHOneAlphaBase}{\ensuremath{1.1739\cdot10^{-4}}}
% NHOneAlphaErr [cert, upper] t9_h1_clam_s180.json :: C_alpha; inputs.YSTAR
\newcommand{\NHOneAlphaErr}{\ensuremath{7.21\cdot10^{-13}}}
% NHOneAlphaMargin [cert, margin] summary.json :: alpha_abs_max; C_alpha; inputs.YSTAR
\newcommand{\NHOneAlphaMargin}{\ensuremath{1.27}}
% NHOneSqrtLamBase [cert, upper] t9_h1_clam_s180.json :: sqrtLam_base_Wp
\newcommand{\NHOneSqrtLamBase}{\ensuremath{2.282\cdot10^{-8}}}
% NHOneSqrtLamTotal [cert, upper] t9_h1_clam_s180.json :: sqrtLam_total_Wp
\newcommand{\NHOneSqrtLamTotal}{\ensuremath{2.29\cdot10^{-8}}}
% NHOneSqrtLamMargin [cert, margin] t9_h1_clam_s180.json :: sqrtLam_total_Wp
\newcommand{\NHOneSqrtLamMargin}{\ensuremath{2.18}}
% NHOneYbarLimit [cert, lower] t9_h1_clam_s180.json :: Ybar_limit_H1
\newcommand{\NHOneYbarLimit}{\ensuremath{1.12\cdot10^{-9}}}
% NSwitchYUpper [cert, upper] switch_s2.json :: modulated.Y_raw
\newcommand{\NSwitchYUpper}{\ensuremath{6.21\cdot10^{-13}}}
% NSwitchYMargin [cert, margin] switch_s2.json :: modulated.Y_raw
\newcommand{\NSwitchYMargin}{\ensuremath{1.6}}
% NMinCI [cert, lower] constants.json :: meta.summary.min_c_I
\newcommand{\NMinCI}{\ensuremath{0.05283}}
% NNumCertificates [cert, exact] constants.json :: meta.summary.n_cert
\newcommand{\NNumCertificates}{\ensuremath{26}}
% NMinRhoC [cert, lower] constants.json :: meta.closure.rho_c_min
\newcommand{\NMinRhoC}{\ensuremath{0.04711}}
% NHTwoBaseZero [cert, upper] FINAL_manifest_v33.json :: results.M9_H2_tightened.rows.0.base
\newcommand{\NHTwoBaseZero}{\ensuremath{8.011\cdot10^{-11}}}
% NHTwoBaseOne [cert, upper] FINAL_manifest_v33.json :: results.M9_H2_tightened.rows.1.base
\newcommand{\NHTwoBaseOne}{\ensuremath{1.336\cdot10^{-10}}}
% NHTwoBaseTwo [cert, upper] FINAL_manifest_v33.json :: results.M9_H2_tightened.rows.2.base
\newcommand{\NHTwoBaseTwo}{\ensuremath{2.936\cdot10^{-10}}}
% NHTwoErrWZero [cert, upper] FINAL_manifest_v33.json :: results.M9_H2_tightened.rows.0.with_error
\newcommand{\NHTwoErrWZero}{\ensuremath{8.014\cdot10^{-11}}}
% NHTwoMarginWZero [cert, margin] FINAL_manifest_v33.json :: results.M9_H2_tightened.rows.0.margin
\newcommand{\NHTwoMarginWZero}{\ensuremath{124.78}}
% NHTwoErrWOne [cert, upper] FINAL_manifest_v33.json :: results.M9_H2_tightened.rows.1.with_error
\newcommand{\NHTwoErrWOne}{\ensuremath{4.706\cdot10^{-10}}}
% NHTwoMarginWOne [cert, margin] FINAL_manifest_v33.json :: results.M9_H2_tightened.rows.1.margin
\newcommand{\NHTwoMarginWOne}{\ensuremath{2124.9}}
% NHTwoErrWTwo [cert, upper] FINAL_manifest_v33.json :: results.M9_H2_tightened.rows.2.with_error
\newcommand{\NHTwoErrWTwo}{\ensuremath{1.559\cdot10^{-6}}}
% NHTwoMarginWTwo [cert, margin] FINAL_manifest_v33.json :: results.M9_H2_tightened.rows.2.margin
\newcommand{\NHTwoMarginWTwo}{\ensuremath{641.53}}
% NHThreeBaseZero [cert, upper] FINAL_manifest_v33.json :: results.M9_H2f.rows.0.base
\newcommand{\NHThreeBaseZero}{\ensuremath{3.127\cdot10^{-7}}}
% NHThreeBaseOne [cert, upper] FINAL_manifest_v33.json :: results.M9_H2f.rows.1.base
\newcommand{\NHThreeBaseOne}{\ensuremath{1.004\cdot10^{-6}}}
% NHThreeBaseTwo [cert, upper] FINAL_manifest_v33.json :: results.M9_H2f.rows.2.base
\newcommand{\NHThreeBaseTwo}{\ensuremath{2.431\cdot10^{-6}}}
% NHThreeErrZero [cert, upper] FINAL_manifest_v33.json :: results.M9_H2f.rows.0.with_error
\newcommand{\NHThreeErrZero}{\ensuremath{3.1363\cdot10^{-7}}}
% NHThreeMarginZero [cert, margin] FINAL_manifest_v33.json :: results.M9_H2f.rows.0.margin
\newcommand{\NHThreeMarginZero}{\ensuremath{3.1885}}
% NHThreeErrOne [cert, upper] FINAL_manifest_v33.json :: results.M9_H2f.rows.1.with_error
\newcommand{\NHThreeErrOne}{\ensuremath{1.0034\cdot10^{-6}}}
% NHThreeMarginOne [cert, margin] FINAL_manifest_v33.json :: results.M9_H2f.rows.1.margin
\newcommand{\NHThreeMarginOne}{\ensuremath{4.9834}}
% NHThreeErrTwo [cert, upper] FINAL_manifest_v33.json :: results.M9_H2f.rows.2.with_error
\newcommand{\NHThreeErrTwo}{\ensuremath{2.4351\cdot10^{-6}}}
% NHThreeMarginTwo [cert, margin] FINAL_manifest_v33.json :: results.M9_H2f.rows.2.margin
\newcommand{\NHThreeMarginTwo}{\ensuremath{4.1067}}
% NHFourBaseZero [cert, upper] FINAL_manifest_v33.json :: results.M9_base_sums.combined.windows_max.0
\newcommand{\NHFourBaseZero}{\ensuremath{3.48\cdot10^{-8}}}
% NHFourErrZero [cert, upper] FINAL_manifest_v33.json :: results.M9_H3_H2v32.rows.win.0.with_error_arb
\newcommand{\NHFourErrZero}{\ensuremath{3.5087\cdot10^{-8}}}
% NHFourMarginZero [cert, margin] FINAL_manifest_v33.json :: results.M9_H3_H2v32.rows.win.0.margin_arb
\newcommand{\NHFourMarginZero}{\ensuremath{5.700}}
% NHFourBaseOne [cert, upper] FINAL_manifest_v33.json :: results.M9_base_sums.combined.windows_max.1
\newcommand{\NHFourBaseOne}{\ensuremath{1.27\cdot10^{-7}}}
% NHFourErrOne [cert, upper] FINAL_manifest_v33.json :: results.M9_H3_H2v32.rows.win.1.with_error_arb
\newcommand{\NHFourErrOne}{\ensuremath{1.2616\cdot10^{-7}}}
% NHFourMarginOne [cert, margin] FINAL_manifest_v33.json :: results.M9_H3_H2v32.rows.win.1.margin_arb
\newcommand{\NHFourMarginOne}{\ensuremath{79.26}}
% NHFourBaseTwo [cert, upper] FINAL_manifest_v33.json :: results.M9_base_sums.combined.windows_max.2
\newcommand{\NHFourBaseTwo}{\ensuremath{3.76\cdot10^{-7}}}
% NHFourErrTwo [cert, upper] FINAL_manifest_v33.json :: results.M9_H3_H2v32.rows.win.2.with_error_arb
\newcommand{\NHFourErrTwo}{\ensuremath{3.7724\cdot10^{-7}}}
% NHFourMarginTwo [cert, margin] FINAL_manifest_v33.json :: results.M9_H3_H2v32.rows.win.2.margin_arb
\newcommand{\NHFourMarginTwo}{\ensuremath{26.50}}
% NHFourBaseThree [cert, upper] FINAL_manifest_v33.json :: results.M9_base_sums.combined.windows_max.3
\newcommand{\NHFourBaseThree}{\ensuremath{9.71\cdot10^{-7}}}
% NHFourErrThree [cert, upper] FINAL_manifest_v33.json :: results.M9_H3_H2v32.rows.win.3.with_error_arb
\newcommand{\NHFourErrThree}{\ensuremath{1.2822\cdot10^{-6}}}
% NHFourMarginThree [cert, margin] FINAL_manifest_v33.json :: results.M9_H3_H2v32.rows.win.3.margin_arb
\newcommand{\NHFourMarginThree}{\ensuremath{7.799}}
% NHFourYbarLimit [cert, lower] FINAL_manifest_v33.json :: results.M9_H3_H2v32.binding_Ybar_limit
\newcommand{\NHFourYbarLimit}{\ensuremath{1.575\cdot10^{-11}}}
% NHFourYbarMargin [cert, margin] FINAL_manifest_v33.json :: results.M9_H3_H2v32.binding_Ybar_limit
\newcommand{\NHFourYbarMargin}{\ensuremath{15.75}}
% NHFourKfourBase [cert, upper] amp_A1_w180.json :: inputs.BASE_T1.K4
\newcommand{\NHFourKfourBase}{\ensuremath{3.1633\cdot10^{-4}}}
% NHSixConeSlope [cert, upper] FINAL_manifest_v33.json :: results.cone_fixedL.checks.slope.value
\newcommand{\NHSixConeSlope}{\ensuremath{0.03029}}
% NHSixConeSlopeMargin [cert, margin] FINAL_manifest_v33.json :: results.cone_fixedL.checks.slope.margin
\newcommand{\NHSixConeSlopeMargin}{\ensuremath{2.641}}
% NHSixConeKtwo [cert, upper] FINAL_manifest_v33.json :: results.cone_fixedL.checks.K2.value
\newcommand{\NHSixConeKtwo}{\ensuremath{0.009876}}
% NHSixConeKtwoMargin [cert, margin] FINAL_manifest_v33.json :: results.cone_fixedL.checks.K2.margin
\newcommand{\NHSixConeKtwoMargin}{\ensuremath{3.037}}
% NHSixConeNthree [cert, upper] FINAL_manifest_v33.json :: results.cone_fixedL.checks.N3.value
\newcommand{\NHSixConeNthree}{\ensuremath{1.461\cdot10^{-4}}}
% NHSixConeNthreeMargin [cert, margin] FINAL_manifest_v33.json :: results.cone_fixedL.checks.N3.margin
\newcommand{\NHSixConeNthreeMargin}{\ensuremath{27.38}}
% NHSixConeHFive [cert, upper] FINAL_manifest_v33.json :: results.H4_fixedL.total
\newcommand{\NHSixConeHFive}{\ensuremath{1.2865\cdot10^{-5}}}
% NHSixConeHFiveMargin [cert, margin] FINAL_manifest_v33.json :: results.H4_fixedL.margin
\newcommand{\NHSixConeHFiveMargin}{\ensuremath{1119}}
% NHSixConeTstarNeeded [cert, upper] FINAL_manifest_v33.json :: results.H7_fixedL.ratio_upper
\newcommand{\NHSixConeTstarNeeded}{\ensuremath{1.6000015\cdot10^{-7}}}
% NHSixConeTstarLower [cert, lower] FINAL_manifest_v33.json :: results.H7_fixedL.t_star_lower
\newcommand{\NHSixConeTstarLower}{\ensuremath{3.68334\cdot10^{-7}}}
% NHSixConeTstarMargin [cert, margin] FINAL_manifest_v33.json :: results.H7_fixedL.margin_tstar_over_ratio
\newcommand{\NHSixConeTstarMargin}{\ensuremath{2.302}}
% NHSixConeHFiveRatioPeak [cert, upper] h4loc_A1_w180.json :: windows.constants/amp_A1_w180.json.worst.ratio
\newcommand{\NHSixConeHFiveRatioPeak}{\ensuremath{0.2393}}
% NHSixConeHOneSup [cert, upper] FINAL_manifest_v33.json :: results.H1_fixedL.sup_h0_max
\newcommand{\NHSixConeHOneSup}{\ensuremath{0.2163}}
% NHSixConeHOneH [cert, upper] FINAL_manifest_v33.json :: results.H1_fixedL.H_max
\newcommand{\NHSixConeHOneH}{\ensuremath{0.2185}}
% NHSixFarHFiveRatio [cert, upper] h4loc_v32.json :: windows.constants/amp_A1_w180.json.cylinders.18.ratio
\newcommand{\NHSixFarHFiveRatio}{\ensuremath{0.1761}}
% NHSixFarHFiveMarginOld [cert, margin] h4loc_v32.json :: windows.constants/amp_A1_w180.json.cylinders.18.ratio
\newcommand{\NHSixFarHFiveMarginOld}{\ensuremath{5.68}}
% NHSixFarHFiveMargin [cert, margin] FINAL_manifest_v33.json :: results.M9_H4_far.min_margin
\newcommand{\NHSixFarHFiveMargin}{\ensuremath{15.58}}
% NHSixConeHFiveMarginOld [cert, margin] h4loc_A1_w180.json :: windows.constants/amp_A1_w180.json.worst.ratio
\newcommand{\NHSixConeHFiveMarginOld}{\ensuremath{4.18}}
% NHSixFarKtwo [cert, upper] t9_far_v32.json :: h2
\newcommand{\NHSixFarKtwo}{\ensuremath{1.797\cdot10^{-4}}}
% NHSixFarKtwoThr [hand, exact] t9_far_v32.json :: thr.h2
\newcommand{\NHSixFarKtwoThr}{\ensuremath{3.53\cdot10^{-4}}}
% NHSixFarKtwoMargin [cert, margin] t9_far_v32.json :: margins.h2
\newcommand{\NHSixFarKtwoMargin}{\ensuremath{1.964}}
% NHSixFarKthree [cert, upper] t9_far_v32.json :: h3
\newcommand{\NHSixFarKthree}{\ensuremath{2.595\cdot10^{-8}}}
% NHSixFarKthreeThr [hand, exact] t9_far_v32.json :: thr.h3
\newcommand{\NHSixFarKthreeThr}{\ensuremath{1.17\cdot10^{-5}}}
% NHSixFarKthreeMargin [cert, margin] t9_far_v32.json :: margins.h3
\newcommand{\NHSixFarKthreeMargin}{\ensuremath{450.9}}
% NHSixFarKfour [cert, upper] t9_far_v32.json :: h4
\newcommand{\NHSixFarKfour}{\ensuremath{6.964\cdot10^{-9}}}
% NHSixFarKfourThr [hand, exact] t9_far_v32.json :: thr.h4
\newcommand{\NHSixFarKfourThr}{\ensuremath{1.03\cdot10^{-7}}}
% NHSixFarKfourMargin [cert, margin] t9_far_v32.json :: margins.h4
\newcommand{\NHSixFarKfourMargin}{\ensuremath{14.79}}
% NHSixFarSlope [cert, upper] t9_far_v32.json :: slope
\newcommand{\NHSixFarSlope}{\ensuremath{0.04153}}
% NHSixFarSlopeThr [hand, exact] t9_far_v32.json :: thr.slope
\newcommand{\NHSixFarSlopeThr}{\ensuremath{0.0800}}
% NHSixFarSlopeMargin [cert, margin] t9_far_v32.json :: margins.slope
\newcommand{\NHSixFarSlopeMargin}{\ensuremath{1.926}}
% NHSixFarNthree [cert, upper] t9_far_v32.json :: L2
\newcommand{\NHSixFarNthree}{\ensuremath{5.516\cdot10^{-7}}}
% NHSixFarNthreeThr [hand, exact] t9_far_v32.json :: thr.L2
\newcommand{\NHSixFarNthreeThr}{\ensuremath{3.32\cdot10^{-6}}}
% NHSixFarNthreeMargin [cert, margin] t9_far_v32.json :: margins.L2
\newcommand{\NHSixFarNthreeMargin}{\ensuremath{6.019}}
% NFarTstarMargin [cert, margin] lemmaO_v32_consts.json :: far_v32.case.margin_factor_lower
\newcommand{\NFarTstarMargin}{\ensuremath{15.5}}
% NFarTstarLower [cert, lower] lemmaO_v32_consts.json :: far_v32.case.t_star_lower
\newcommand{\NFarTstarLower}{\ensuremath{9.5108\cdot10^{-8}}}
% NFarTstarNeeded [cert, upper] lemmaO_v32_consts.json :: far_v32.case.needed
\newcommand{\NFarTstarNeeded}{\ensuremath{6.134\cdot10^{-9}}}
% NGammaModes [cert, exact] gamma0_final.json :: M
\newcommand{\NGammaModes}{\ensuremath{3762}}
% NGammaHalfwidth [cert, upper] gamma0_final.json :: b_j_halfwidth
\newcommand{\NGammaHalfwidth}{\ensuremath{8.138\cdot10^{-24}}}
% NGammaContraction [cert, upper] gamma0_final.json :: krawczyk_contraction
\newcommand{\NGammaContraction}{\ensuremath{3.01\cdot10^{-12}}}
% NGammaArcs [cert, exact] gamma0_final.json :: embedded.arcs
\newcommand{\NGammaArcs}{\ensuremath{701}}
% NGammaMaxTurn [cert, upper] gamma0_final.json :: embedded.maxturn
\newcommand{\NGammaMaxTurn}{\ensuremath{0.04997}}
% NGammaMaxPair [cert, upper] gamma0_final.json :: embedded.maxpair
\newcommand{\NGammaMaxPair}{\ensuremath{0.09597}}
% NGammaMinSep [cert, lower] gamma0_final.json :: embedded.minsep
\newcommand{\NGammaMinSep}{\ensuremath{8.89\cdot10^{-3}}}
% NGammaHalfGap [cert, lower] gamma0_final.json :: embedded.gap_cert
\newcommand{\NGammaHalfGap}{\ensuremath{0.0248219}}
% NPhaseISlabs [cert, exact] status.json :: slab
\newcommand{\NPhaseISlabs}{\ensuremath{96}}
% NPhaseIModes [cert, exact] status.json :: last_slab.M
\newcommand{\NPhaseIModes}{\ensuremath{2800}}
% NPhaseIKmaxLower [cert, lower] status.json :: kmax_enclosure.0
\newcommand{\NPhaseIKmaxLower}{\ensuremath{29.352874545}}
% NPhaseIKmaxUpper [cert, upper] status.json :: kmax_enclosure.1
\newcommand{\NPhaseIKmaxUpper}{\ensuremath{29.352874546}}
% NPhaseITLower [cert, lower] status.json :: t
\newcommand{\NPhaseITLower}{\ensuremath{4.597850539485830}}
% NPhaseITUpper [cert, upper] status.json :: t
\newcommand{\NPhaseITUpper}{\ensuremath{4.597850539485831}}
% NPeakSqrtLam [cert, upper] rate_table_fin12f4_a0.00015.json :: sqrtL_max
\newcommand{\NPeakSqrtLam}{\ensuremath{5.308\cdot10^{-8}}}
% NTailRate [cert, lower] rate_table_fin12f4_a0.00015.json :: tail.rho
\newcommand{\NTailRate}{\ensuremath{2.693\cdot10^{-3}}}
% NFloorSqrtLam [cert, upper] rate_table_fin12f4_a0.00015.json :: tail.sqrt_LF
\newcommand{\NFloorSqrtLam}{\ensuremath{9.872\cdot10^{-10}}}
% NAlphaFloor [cert, upper] rate_table_fin12f4_a0.00015.json :: tail.alpha_F
\newcommand{\NAlphaFloor}{\ensuremath{6.867\cdot10^{-7}}}
% NSupBetaVerifier [cert, upper] final_alpha_points_5e-8.json :: 1.5e-4.sup_beta
\newcommand{\NSupBetaVerifier}{\ensuremath{3.33811\cdot10^{-9}}}
% NMOneRamp [cert, upper] final_alpha_points_5e-8.json :: 1.5e-4.M1
\newcommand{\NMOneRamp}{\ensuremath{0.4328}}
% NSITwoPrime [cert, upper] FINAL_manifest_v33.json :: results.F12_sliver.sI2_prime
\newcommand{\NSITwoPrime}{\ensuremath{0.015457}}
% NLoopGainCoarse [cert, upper] U_cert_v35.json :: loop_gain.2.total
\newcommand{\NLoopGainCoarse}{\ensuremath{0.03904}}
% NRateMarginCoarse [cert, margin] U_cert_v35.json :: rate_margin
\newcommand{\NRateMarginCoarse}{\ensuremath{28.2}}
% NSOneEpsLower [cert, lower] bracket2.json :: eps_a
\newcommand{\NSOneEpsLower}{\ensuremath{1.5346198892965\cdot10^{-4}}}
% NSOneEpsUpper [cert, upper] bracket2.json :: eps_b
\newcommand{\NSOneEpsUpper}{\ensuremath{1.5346198899299\cdot10^{-4}}}
% NSOneRowsTip [cert, exact] ext_tip_check.json :: rows
\newcommand{\NSOneRowsTip}{\ensuremath{51246}}
% NSOneTipWidth [cert, upper] ext_tip_check.json :: max_pointwise_width.th
\newcommand{\NSOneTipWidth}{\ensuremath{3.2\cdot10^{-12}}}
% NSOneExtWidth [cert, upper] ext813_check.json :: max_pointwise_width.p
\newcommand{\NSOneExtWidth}{\ensuremath{5.7\cdot10^{-10}}}
% NSOneRowsMC [cert, exact] check_outA5.json :: prof_tip_vs_MC.rows_compared
\newcommand{\NSOneRowsMC}{\ensuremath{1025}}
% NSOneRowsMB [cert, exact] check_outA5.json :: ext813_vs_MB.rows_compared
\newcommand{\NSOneRowsMB}{\ensuremath{501}}
% NSOneMinMargin [cert, margin] check_outA5.json :: min:*.min_relative_margin.*
\newcommand{\NSOneMinMargin}{\ensuremath{1.2\cdot10^{-9}}}
% NTwoSegments [cert, exact] contour_result.json :: len:segments
\newcommand{\NTwoSegments}{\ensuremath{5058}}
% NTwoInteriorBoxes [cert, exact] interior_result.json :: len:
\newcommand{\NTwoInteriorBoxes}{\ensuremath{174}}
% NTwoContourBoxes [cert, exact] eta_cache.json :: len:
\newcommand{\NTwoContourBoxes}{\ensuremath{202}}
% NTwoEtaMax [cert, upper] eta_cache.json :: max:*.eta
\newcommand{\NTwoEtaMax}{\ensuremath{0.1163}}
% NTwoLipMax [cert, upper] eta_cache.json :: max:*.Lip
\newcommand{\NTwoLipMax}{\ensuremath{0.272}}
% NTwoLambdaZero [est, upper] cert_b0.25_th0.0_tau0.03.json :: Lam
\newcommand{\NTwoLambdaZero}{\ensuremath{1.790819}}
% NTwoLambdaOne [est, upper] cert_b0.25_th0.48_tau0.03.json :: Lam
\newcommand{\NTwoLambdaOne}{\ensuremath{1.768282}}
% NSOneDTightLower [doc, lower] doc:README.md :: D ∈ [2.50080627654, 2.50080627668]
\newcommand{\NSOneDTightLower}{\ensuremath{2.50080627654}}
% NSOneDTightUpper [doc, upper] doc:README.md :: D ∈ [2.50080627654, 2.50080627668]
\newcommand{\NSOneDTightUpper}{\ensuremath{2.50080627668}}
% NSOneKTightLower [doc, lower] doc:README.md :: k(0) ∈ [1.53982063762, 1.53982063822]
\newcommand{\NSOneKTightLower}{\ensuremath{1.53982063762}}
% NSOneKTightUpper [doc, upper] doc:README.md :: k(0) ∈ [1.53982063762, 1.53982063822]
\newcommand{\NSOneKTightUpper}{\ensuremath{1.53982063822}}
% NSOneDCondLower [doc, lower] doc:README.md :: D = −p(tip) ∈ [2.500806276619695, 2.500806276620049]
\newcommand{\NSOneDCondLower}{\ensuremath{2.500806276619695}}
% NSOneDCondUpper [doc, upper] doc:README.md :: D = −p(tip) ∈ [2.500806276619695, 2.500806276620049]
\newcommand{\NSOneDCondUpper}{\ensuremath{2.500806276620049}}
% NSOneKCondLower [doc, lower] doc:README.md :: k(0) ∈ [1.539820637871836, 1.5398206378721015]
\newcommand{\NSOneKCondLower}{\ensuremath{1.539820637871836}}
% NSOneKCondUpper [doc, upper] doc:README.md :: k(0) ∈ [1.539820637871836, 1.5398206378721015]
\newcommand{\NSOneKCondUpper}{\ensuremath{1.5398206378721015}}
% NSOneTauTipLower [doc, lower] doc:README.md :: τ_tip ∈ [10.249077651120778, 10.249077651121361]
\newcommand{\NSOneTauTipLower}{\ensuremath{10.249077651120778}}
% NSOneTauTipUpper [doc, upper] doc:README.md :: τ_tip ∈ [10.249077651120778, 10.249077651121361]
\newcommand{\NSOneTauTipUpper}{\ensuremath{10.249077651121361}}
% NSOneStateUncert [doc, upper] doc:README.md :: State uncertainty at σ = 8: 4.4e−13
\newcommand{\NSOneStateUncert}{\ensuremath{4.4\cdot10^{-13}}}
% NFarBoxEprime [doc, margin] doc:STEP4_RESULT.md :: (6.4287, 4.6917, 1.9940, 1.69999995)
\newcommand{\NFarBoxEprime}{\ensuremath{6.42}}
% NFarBoxXEsecond [doc, margin] doc:STEP4_RESULT.md :: (6.4287, 4.6917, 1.9940, 1.69999995)
\newcommand{\NFarBoxXEsecond}{\ensuremath{4.69}}
% NFarBoxXsqEthird [doc, margin] doc:STEP4_RESULT.md :: (6.4287, 4.6917, 1.9940, 1.69999995)
\newcommand{\NFarBoxXsqEthird}{\ensuremath{1.99}}
% NFarBoxE [doc, margin] doc:STEP4_RESULT.md :: (6.4287, 4.6917, 1.9940, 1.69999995)
\newcommand{\NFarBoxE}{\ensuremath{1.69}}
% NHFourKfourErr [doc, upper] doc:STEP4_RESULT.md :: | 5.60e-3 against 0.06 | ×10.7 |
\newcommand{\NHFourKfourErr}{\ensuremath{5.60\cdot10^{-3}}}
% NHFourKfourMargin [doc, margin] doc:STEP4_RESULT.md :: | 5.60e-3 against 0.06 | ×10.7 |
\newcommand{\NHFourKfourMargin}{\ensuremath{10.7}}
% NCoercivityCoefficient [doc, upper] doc:step3_proof_v32.tex :: \norm{e^{\tau\mathcal L}u}_\rho\le14.8\,e^{-0.0213\tau}\norm
\newcommand{\NCoercivityCoefficient}{\ensuremath{14.8}}
% NCoercivityRate [doc, lower] doc:step3_proof_v32.tex :: \norm{e^{\tau\mathcal L}u}_\rho\le14.8\,e^{-0.0213\tau}\norm
\newcommand{\NCoercivityRate}{\ensuremath{0.0213}}
% NRegionRampStart [hand, exact] hand :: 
\newcommand{\NRegionRampStart}{\ensuremath{14}}
% NRegionInterface [hand, exact] hand :: 
\newcommand{\NRegionInterface}{\ensuremath{28}}
% NRegionShellEnd [hand, exact] hand :: 
\newcommand{\NRegionShellEnd}{\ensuremath{60}}
% NRegionSegEnd [hand, exact] hand :: 
\newcommand{\NRegionSegEnd}{\ensuremath{350}}
% NRegionAnnEnd [hand, exact] hand :: 
\newcommand{\NRegionAnnEnd}{\ensuremath{700}}
% NZoneIIILow [hand, exact] hand :: 
\newcommand{\NZoneIIILow}{\ensuremath{253}}
% NZoneIIIHigh [hand, exact] hand :: 
\newcommand{\NZoneIIIHigh}{\ensuremath{379.5}}
% NConeCentreMin [hand, exact] hand :: 
\newcommand{\NConeCentreMin}{\ensuremath{300}}
% NConeCentreMax [hand, exact] hand :: 
\newcommand{\NConeCentreMax}{\ensuremath{380}}
% NFarCentreMin [hand, exact] hand :: 
\newcommand{\NFarCentreMin}{\ensuremath{490}}
% NConeL [hand, exact] hand :: 
\newcommand{\NConeL}{\ensuremath{50}}
% NConeX [hand, exact] hand :: 
\newcommand{\NConeX}{\ensuremath{3}}
% NFarL [hand, exact] hand :: 
\newcommand{\NFarL}{\ensuremath{113}}
% NFarX [hand, exact] hand :: 
\newcommand{\NFarX}{\ensuremath{2}}
}%
  {\PackageWarning{paper}{numbers/numbers_paper.tex not found (unit U0 not finished)}}

% ======================================================================
% NOTATION MACROS (OUTLINE 5.2 and notation.md).  One macro per symbol
% whose spelling is fixed by the notation table.  Plain letters (a, b, f, ...)
% have no macro.
% ======================================================================
% basic
\newcommand{\R}{\mathbb{R}}
\newcommand{\C}{\mathbb{C}}
\newcommand{\N}{\mathbb{N}}
\newcommand{\Z}{\mathbb{Z}}
\newcommand{\eps}{\varepsilon}
\newcommand{\ip}[2]{\langle #1,#2\rangle}
\newcommand{\abs}[1]{\lvert #1\rvert}
\newcommand{\norm}[1]{\lVert #1\rVert}
\newcommand{\dd}{\mathrm d}
\newcommand{\re}{\operatorname{Re}}
\newcommand{\im}{\operatorname{Im}}
% flow and curves
\newcommand{\arclen}{\mathfrak{s}}                 % arclength of the flow curve
\newcommand{\kss}{\kappa_{\mathfrak{ss}}}          % kappa_{ss} in the flow equation
\newcommand{\Area}{\operatorname{Area}}
\newcommand{\Len}{\mathcal L}                      % length of a curve
\newcommand{\spp}{\operatorname{sp}}                 % sup quantities sp^v_0, sp^v_1, sp^f_1 (were S^v_0, S^v_1, S^f_1)
\newcommand{\rot}{\operatorname{rot}}              % turning number
\newcommand{\Tsing}{T_*}                           % singular time
\newcommand{\Tmax}{T_{\max}}                       % maximal time of existence
\newcommand{\pstar}{p_*}                           % image point of the singularity
\newcommand{\Gz}{\Gamma_0}                         % initial curve
\newcommand{\vth}{\vartheta}                       % parameter of Gamma_0
\newcommand{\dclk}{\delta_t}                       % clock defect of step 4
\newcommand{\fall}{\varpi}                         % Fourier rates varpi_n = n^4-n^2
% rescaling
\newcommand{\lamone}{\lambda_1}
\newcommand{\rate}{\varsigma}                      % 1/4 - a
\newcommand{\xia}{\xi_a}                           % (1-4\bar a)^{-1}
% shrinker and step 1
\newcommand{\Sone}{S_1}
\newcommand{\Gsh}{\Gamma_{\!*}}                    % profile curve of S_1
\newcommand{\cSet}{\mathcal S}                     % certified set of shrinkers
\newcommand{\sigtip}{\sigma_{\rm tip}}
\newcommand{\Xo}{X_0}                              % 2.2490776515
\newcommand{\Dtip}{D}
\newcommand{\ktip}{k_0}
\newcommand{\shootM}{M}                            % shooting function
\newcommand{\Ypar}{\mathcal Y}                     % step-1 parallelepiped
\newcommand{\axis}{\Sigma}                         % symmetry axis of S_1
\newcommand{\sigc}{\sigma_{\rm c}}
\newcommand{\sigrev}{\sigma^{\rm rev}}
% linear theory
\newcommand{\Hrho}{\mathcal H_\rho}
\newcommand{\muslow}{\mu_3}
\newcommand{\rre}{\mathfrak r_{\rm re}}
\newcommand{\rim}{\mathfrak r_{\rm im}}
\newcommand{\Evev}{\mathrm{Ev}_{\rm even}}
\newcommand{\Evod}{\mathrm{Ev}_{\rm odd}}
\newcommand{\gvec}{\mathbf g}                      % translation vector in step 2
\newcommand{\cons}[1]{\mathfrak f_{#1}}            % constraint functions
\newcommand{\phirho}{\varphi_\rho}
% energies
\newcommand{\Lam}{\Lambda}
\newcommand{\En}{E}
\newcommand{\EnI}{E_I}
\newcommand{\muI}{\mu_I}                           % 0.08
\newcommand{\epsE}{\varepsilon_E}                  % 1e-5
\newcommand{\epsI}{\varepsilon_I}                  % 3e-5
\newcommand{\wn}[1]{\mathsf w_{#1}}                % ||w~^{(j)}||_rho
\newcommand{\vI}{v_I}                              % zeta w'
\newcommand{\muerr}{\mu_{\rm err}}                 % step-4 error-energy weight 0.036
% far field
\newcommand{\pU}{p_U}                              % 1.0246
\newcommand{\Uone}{U_1}
\newcommand{\sm}{\operatorname{sm}}                % smoothstep t^3(10-15t+6t^2)
\newcommand{\rc}{r_c}                              % centre radius (Lemma U)
\newcommand{\rchi}{r_\chi}                         % 253
\newcommand{\Iw}{I_w}                              % hand-off windows [y_w,1.5y_w]
\newcommand{\qa}{\mathfrak q_1}                    % quasilinear coefficients of the slope equation
\newcommand{\qb}{\mathfrak q_2}
% cylinders (Lemma O)
\newcommand{\cyl}{\ell}                            % cylinder length unit
\newcommand{\MO}{M_O}
\newcommand{\KO}{K_O}
\newcommand{\Hcyl}{H_{\rm cyl}}
\newcommand{\XiO}{\Xi_O}
\newcommand{\fO}{f_O}
% step 4
\newcommand{\Ybar}{\bar Y}
\newcommand{\Lt}{\widetilde L}                     % linearisation about the step-4 base
\newcommand{\Trans}{\Pi}                           % transport integral
\newcommand{\resid}{\mathrm{res}}
\newcommand{\Resid}[1]{\mathrm{Res}_{#1}}
\newcommand{\boxtip}{\Box_{\rm tip}}
\newcommand{\boxfar}{\Box_{\rm far}}
\newcommand{\Eenv}{\mathcal E_{\rm env}}
\newcommand{\betab}{\beta_{\rm base}}
\newcommand{\tauI}{\tau^{\rm I}}                   % Phase I clock
% names of families of conditions: \Hh{3} = (H3), \Oh{4} = (O4), ...
\newcommand{\Hh}[1]{\textup{(H#1)}}
\newcommand{\Oh}[1]{\textup{(O#1)}}
\newcommand{\Lh}[1]{\textup{(L#1)}}
\newcommand{\Bh}[1]{\textup{(B#1)}}
\newcommand{\Kh}[1]{\textup{(K#1)}}
\newcommand{\Wh}[1]{\textup{(W#1)}}
\newcommand{\PhaseI}{Phase~I}
\newcommand{\PhaseII}{Phase~II}
% margin: threshold/bound (always)
\newcommand{\mgn}[1]{{\times}#1}

% ---- figures: placeholder macro so that drafts compile before F1-F7 exist ----
\newcommand{\figplaceholder}[1]{\fbox{\parbox{0.8\linewidth}{\centering\small[figure #1 to be supplied]}}}
 ...
% (the main paper does nothing special: \ifsupplement defaults to false).
% Compile from paper_v1/ so that % numbers_paper.tex -- GENERATED from ledger.json by build_ledger.py; do not edit.
% Rounding: upper bounds up, lower bounds and margins down at the printed digit; margin = threshold / bound.
% Use as $\times\NName$ for margins.  kind: cert certified, est estimate, hand hypothesis constant, doc only a document quotes it.
% NUEhatZero [cert, upper] U_cert_v35.json :: supersolution.EA.0
\newcommand{\NUEhatZero}{\ensuremath{1.2018\cdot10^{-6}}}
% NUEhatOne [cert, upper] U_cert_v35.json :: supersolution.EA.1
\newcommand{\NUEhatOne}{\ensuremath{1.0441\cdot10^{-4}}}
% NUEhatTwo [cert, upper] U_cert_v35.json :: supersolution.EA.2
\newcommand{\NUEhatTwo}{\ensuremath{0.056694}}
% NUBhatZero [cert, upper] U_cert_v35.json :: supersolution.Bhat.0
\newcommand{\NUBhatZero}{\ensuremath{4.8735\cdot10^{-6}}}
% NUBhatOne [cert, upper] U_cert_v35.json :: supersolution.Bhat.1
\newcommand{\NUBhatOne}{\ensuremath{5.5933\cdot10^{-4}}}
% NUBhatTwo [cert, upper] U_cert_v35.json :: supersolution.Bhat.2
\newcommand{\NUBhatTwo}{\ensuremath{0.36675}}
% NUSupersolutionCZeroZero [cert, upper] U_cert_v35.json :: supersolution.C0.0
\newcommand{\NUSupersolutionCZeroZero}{\ensuremath{1.1340\cdot10^{-6}}}
% NUSupersolutionCZeroOne [cert, upper] U_cert_v35.json :: supersolution.C0.1
\newcommand{\NUSupersolutionCZeroOne}{\ensuremath{1.6623\cdot10^{-4}}}
% NUSupersolutionCOneZero [cert, upper] U_cert_v35.json :: supersolution.C1.0
\newcommand{\NUSupersolutionCOneZero}{\ensuremath{1.5241\cdot10^{-9}}}
% NUSupersolutionCOneOne [cert, upper] U_cert_v35.json :: supersolution.C1.1
\newcommand{\NUSupersolutionCOneOne}{\ensuremath{5.5255\cdot10^{-8}}}
% NUSupersolutionNTwoZero [cert, upper] U_cert_v35.json :: supersolution.N2.0
\newcommand{\NUSupersolutionNTwoZero}{\ensuremath{0.071472}}
% NUSupersolutionNTwoOne [cert, upper] U_cert_v35.json :: supersolution.N2.1
\newcommand{\NUSupersolutionNTwoOne}{\ensuremath{0.33012}}
% NUSupersolutionNTwoTwo [cert, upper] U_cert_v35.json :: supersolution.N2.2
\newcommand{\NUSupersolutionNTwoTwo}{\ensuremath{5.0406}}
% NUSupersolutionNTwoOverFourNTwoRefZero [cert, upper] U_cert_v35.json :: supersolution.N2_over_4N2ref.0
\newcommand{\NUSupersolutionNTwoOverFourNTwoRefZero}{\ensuremath{0.18}}
% NUSupersolutionNTwoOverFourNTwoRefOne [cert, upper] U_cert_v35.json :: supersolution.N2_over_4N2ref.1
\newcommand{\NUSupersolutionNTwoOverFourNTwoRefOne}{\ensuremath{0.42}}
% NUSupersolutionNTwoOverFourNTwoRefTwo [cert, upper] U_cert_v35.json :: supersolution.N2_over_4N2ref.2
\newcommand{\NUSupersolutionNTwoOverFourNTwoRefTwo}{\ensuremath{0.26}}
% NUSupersolutionMOneZHypCheckZonesIIISOneTwoBhat [cert, upper] U_cert_v35.json :: supersolution.M1z_hyp_check.zones_I_II.S1_2Bhat
\newcommand{\NUSupersolutionMOneZHypCheckZonesIIISOneTwoBhat}{\ensuremath{0.28585}}
% NUSupersolutionMOneZHypCheckZoneIIIBOnePlusPUOne [cert, upper] U_cert_v35.json :: supersolution.M1z_hyp_check.zone_III.b1_plus_pU1
\newcommand{\NUSupersolutionMOneZHypCheckZoneIIIBOnePlusPUOne}{\ensuremath{0.50786}}
% NUMarginLateMOne [cert, margin] U_cert_v35.json :: margins.late_M1
\newcommand{\NUMarginLateMOne}{\ensuremath{2.353}}
% NUMarginLateMZero [cert, margin] U_cert_v35.json :: margins.late_M0
\newcommand{\NUMarginLateMZero}{\ensuremath{4.892}}
% NUMarginsLateMZeroZ [cert, margin] U_cert_v35.json :: margins.late_M0z
\newcommand{\NUMarginsLateMZeroZ}{\ensuremath{12.99}}
% NUMarginPresentMOne [cert, margin] U_cert_v35.json :: margins.present_M1_min.value
\newcommand{\NUMarginPresentMOne}{\ensuremath{2.770}}
% NUMarginPresentMZero [cert, margin] U_cert_v35.json :: margins.present_M0_min.value
\newcommand{\NUMarginPresentMZero}{\ensuremath{2.225}}
% NUMarginsPresentMZeroZMinValue [cert, margin] U_cert_v35.json :: margins.present_M0z_min.value
\newcommand{\NUMarginsPresentMZeroZMinValue}{\ensuremath{14.37}}
% NULoopGainZeroTotal [cert, upper] U_cert_v35.json :: loop_gain.0.total
\newcommand{\NULoopGainZeroTotal}{\ensuremath{1.23\cdot10^{-7}}}
% NULoopGainOneTotal [cert, upper] U_cert_v35.json :: loop_gain.1.total
\newcommand{\NULoopGainOneTotal}{\ensuremath{0.037169}}
% NULoopGainTwoNeighbour [cert, upper] U_cert_v35.json :: loop_gain.2.neighbour
\newcommand{\NULoopGainTwoNeighbour}{\ensuremath{0.037018}}
% NULoopGainTwoThroughA [cert, upper] U_cert_v35.json :: loop_gain.2.through_A
\newcommand{\NULoopGainTwoThroughA}{\ensuremath{0.0020172}}
% NUfineBhatZero [cert, upper] U_cert_v35.json :: ufine.Bhat.0
\newcommand{\NUfineBhatZero}{\ensuremath{4.8723\cdot10^{-6}}}
% NUfineBhatOne [cert, upper] U_cert_v35.json :: ufine.Bhat.1
\newcommand{\NUfineBhatOne}{\ensuremath{2.6266\cdot10^{-4}}}
% NUfineBhatTwo [cert, upper] U_cert_v35.json :: ufine.Bhat.2
\newcommand{\NUfineBhatTwo}{\ensuremath{0.025842}}
% NUcdeltaU [cert, upper] U_cert_v35.json :: launch.c_dU
\newcommand{\NUcdeltaU}{\ensuremath{1.7187\cdot10^{-8}}}
% NULaunchWmin [cert, lower] U_cert_v35.json :: launch.wmin
\newcommand{\NULaunchWmin}{\ensuremath{0.73492}}
% NULaunchCAD [cert, upper] U_cert_v35.json :: launch.CA_D
\newcommand{\NULaunchCAD}{\ensuremath{2.2429}}
% NULaunchQD [cert, upper] U_cert_v35.json :: launch.qD
\newcommand{\NULaunchQD}{\ensuremath{0.018509}}
% NULaunchDeltaInZero [cert, upper] U_cert_v35.json :: launch.Delta_in.0
\newcommand{\NULaunchDeltaInZero}{\ensuremath{3.131\cdot10^{-12}}}
% NULaunchDeltaInOne [cert, upper] U_cert_v35.json :: launch.Delta_in.1
\newcommand{\NULaunchDeltaInOne}{\ensuremath{3.287\cdot10^{-12}}}
% NULaunchDeltaInTwo [cert, upper] U_cert_v35.json :: launch.Delta_in.2
\newcommand{\NULaunchDeltaInTwo}{\ensuremath{3.451\cdot10^{-12}}}
% NULaunchPhiZero [cert, upper] U_cert_v35.json :: launch.Phi.0
\newcommand{\NULaunchPhiZero}{\ensuremath{6.434\cdot10^{-9}}}
% NULaunchPhiOne [cert, upper] U_cert_v35.json :: launch.Phi.1
\newcommand{\NULaunchPhiOne}{\ensuremath{1.4217\cdot10^{-8}}}
% NULaunchPhiTwo [cert, upper] U_cert_v35.json :: launch.Phi.2
\newcommand{\NULaunchPhiTwo}{\ensuremath{3.8404\cdot10^{-8}}}
% NUGammaZero [cert, upper] U_cert_v35.json :: Gamma.0
\newcommand{\NUGammaZero}{\ensuremath{4.0000358}}
% NUGammaOne [cert, upper] U_cert_v35.json :: Gamma.1
\newcommand{\NUGammaOne}{\ensuremath{16.000322}}
% NUGammaTwo [cert, upper] U_cert_v35.json :: Gamma.2
\newcommand{\NUGammaTwo}{\ensuremath{64.00218}}
% NUKappa [cert, upper] U_cert_v35.json :: kappa
\newcommand{\NUKappa}{\ensuremath{8.0000358}}
% NUKappaTwo [cert, upper] U_cert_v35.json :: kappa2
\newcommand{\NUKappaTwo}{\ensuremath{64.00061}}
% NUGhatZeroZero [cert, upper] U_cert_v35.json :: Ghat.0.0
\newcommand{\NUGhatZeroZero}{\ensuremath{51322.6}}
% NUGhatOneOne [cert, upper] U_cert_v35.json :: Ghat.1.1
\newcommand{\NUGhatOneOne}{\ensuremath{33962.9}}
% NUGhatTwoTwo [cert, upper] U_cert_v35.json :: Ghat.2.2
\newcommand{\NUGhatTwoTwo}{\ensuremath{348341.1}}
% NUGhatOneZero [cert, upper] U_cert_v35.json :: Ghat.1.0
\newcommand{\NUGhatOneZero}{\ensuremath{4.938}}
% NUGhatTwoZero [cert, upper] U_cert_v35.json :: Ghat.2.0
\newcommand{\NUGhatTwoZero}{\ensuremath{27.984}}
% NUGhatTwoOne [cert, upper] U_cert_v35.json :: Ghat.2.1
\newcommand{\NUGhatTwoOne}{\ensuremath{458.53}}
% NUKZeroZero [cert, upper] U_cert_v35.json :: K.0.0
\newcommand{\NUKZeroZero}{\ensuremath{8034.21}}
% NUKZeroOne [cert, upper] U_cert_v35.json :: K.0.1
\newcommand{\NUKZeroOne}{\ensuremath{724.89}}
% NUKOneZero [cert, upper] U_cert_v35.json :: K.1.0
\newcommand{\NUKOneZero}{\ensuremath{0.08152}}
% NUKOneOne [cert, upper] U_cert_v35.json :: K.1.1
\newcommand{\NUKOneOne}{\ensuremath{6336.94}}
% NUKOneTwo [cert, upper] U_cert_v35.json :: K.1.2
\newcommand{\NUKOneTwo}{\ensuremath{545.59}}
% NUKOneThree [cert, upper] U_cert_v35.json :: K.1.3
\newcommand{\NUKOneThree}{\ensuremath{0.58700}}
% NUKTwoZero [cert, upper] U_cert_v35.json :: K.2.0
\newcommand{\NUKTwoZero}{\ensuremath{0.87186}}
% NUKTwoOne [cert, upper] U_cert_v35.json :: K.2.1
\newcommand{\NUKTwoOne}{\ensuremath{528.62}}
% NUKTwoTwo [cert, upper] U_cert_v35.json :: K.2.2
\newcommand{\NUKTwoTwo}{\ensuremath{23425.98}}
% NUKTwoThree [cert, upper] U_cert_v35.json :: K.2.3
\newcommand{\NUKTwoThree}{\ensuremath{555.62}}
% NUFootwOne [cert, upper] U_cert_v35.json :: footw.1
\newcommand{\NUFootwOne}{\ensuremath{2.6361}}
% NUFootwTwo [cert, upper] U_cert_v35.json :: footw.2
\newcommand{\NUFootwTwo}{\ensuremath{2663.3}}
% NUCThreeThree [cert, upper] U_cert_v35.json :: c33
\newcommand{\NUCThreeThree}{\ensuremath{108.02}}
% NUCTwoTwoThree [cert, upper] U_cert_v35.json :: c223
\newcommand{\NUCTwoTwoThree}{\ensuremath{7.3284}}
% NUCEta [cert, upper] U_cert_v35.json :: C_eta
\newcommand{\NUCEta}{\ensuremath{13.376}}
% NULMax [cert, upper] U_cert_v35.json :: L_max
\newcommand{\NULMax}{\ensuremath{3.5352}}
% NURRatio [cert, upper] U_cert_v35.json :: R_ratio
\newcommand{\NURRatio}{\ensuremath{11.593}}
% NUSweepsSigmaInCoverSup [cert, upper] U_cert_v35.json :: sweeps.Sigma_in.cover_sup
\newcommand{\NUSweepsSigmaInCoverSup}{\ensuremath{5.0892}}
% NUSweepsSigmaECoverSup [cert, upper] U_cert_v35.json :: sweeps.Sigma_E.cover_sup
\newcommand{\NUSweepsSigmaECoverSup}{\ensuremath{0.30846}}
% NUSweepsSXCoverSup [cert, upper] U_cert_v35.json :: sweeps.S_X.cover_sup
\newcommand{\NUSweepsSXCoverSup}{\ensuremath{1.5243}}
% NUSweepsSigmaXCoverSup [cert, upper] U_cert_v35.json :: sweeps.Sigma_X.cover_sup
\newcommand{\NUSweepsSigmaXCoverSup}{\ensuremath{3.0081}}
% NUSweepsSigmaRsTwoCoverSup [cert, upper] U_cert_v35.json :: sweeps.Sigma_rs.2.cover_sup
\newcommand{\NUSweepsSigmaRsTwoCoverSup}{\ensuremath{4.3350\cdot10^{-3}}}
% NUDataTermsLateDZero [cert, upper] U_cert_v35.json :: data_terms_late.D.0
\newcommand{\NUDataTermsLateDZero}{\ensuremath{1.7330\cdot10^{-6}}}
% NUDataTermsLateDOne [cert, upper] U_cert_v35.json :: data_terms_late.D.1
\newcommand{\NUDataTermsLateDOne}{\ensuremath{2.5353\cdot10^{-4}}}
% NUDataTermsLateSlopeShareCOneSigmaXOverDkZero [cert, upper] U_cert_v35.json :: data_terms_late.slope_share_C1SigmaX_over_Dk.0
\newcommand{\NUDataTermsLateSlopeShareCOneSigmaXOverDkZero}{\ensuremath{2.7\cdot10^{-3}}}
% NUDataTermsLateSlopeShareCOneSigmaXOverDkOne [cert, upper] U_cert_v35.json :: data_terms_late.slope_share_C1SigmaX_over_Dk.1
\newcommand{\NUDataTermsLateSlopeShareCOneSigmaXOverDkOne}{\ensuremath{6.6\cdot10^{-4}}}
% NURLtTwoMarginZero [cert, margin] U_cert_v35.json :: R_lt2_margin.0
\newcommand{\NURLtTwoMarginZero}{\ensuremath{4.432}}
% NURLtTwoMarginOne [cert, margin] U_cert_v35.json :: R_lt2_margin.1
\newcommand{\NURLtTwoMarginOne}{\ensuremath{6.697}}
% NUThetaTwoEight [cert, upper] U_cert_v35.json :: theta28
\newcommand{\NUThetaTwoEight}{\ensuremath{0.37011}}
% NUDFiveSix [cert, upper] U_cert_v35.json :: D56
\newcommand{\NUDFiveSix}{\ensuremath{1.0368}}
% NUDTwoEight [cert, upper] U_cert_v35.json :: D28
\newcommand{\NUDTwoEight}{\ensuremath{1.5876}}
% NUSThreeThreeFootOverBhatTwo [cert, upper] U_cert_v35.json :: S33_foot_over_Bhat2
\newcommand{\NUSThreeThreeFootOverBhatTwo}{\ensuremath{1.69\cdot10^{-3}}}
% NSAHJWKappaJ [cert, upper] A_interface.json :: HJW.kappa_J
\newcommand{\NSAHJWKappaJ}{\ensuremath{0.042351}}
% NSAGpZero [cert, lower] A_interface.json :: res.tight.gp.0
\newcommand{\NSAGpZero}{\ensuremath{0.24292}}
% NSAGpOne [cert, lower] A_interface.json :: res.tight.gp.1
\newcommand{\NSAGpOne}{\ensuremath{0.24106}}
% NSAGpTwo [cert, lower] A_interface.json :: res.tight.gp.2
\newcommand{\NSAGpTwo}{\ensuremath{0.11961}}
% NSAAminC [cert, lower] A_interface.json :: res.tight.amin_c
\newcommand{\NSAAminC}{\ensuremath{0.97623}}
% NSARhoJ [cert, upper] A_interface.json :: res.tight.rhoJ
\newcommand{\NSARhoJ}{\ensuremath{5.4995\cdot10^{-7}}}
% NSAKappaOTwoTwo [cert, upper] A_interface.json :: res.tight.kappa_o.2.2
\newcommand{\NSAKappaOTwoTwo}{\ensuremath{1.2416\cdot10^{-5}}}
% NSADsrcZero [cert, upper] A_interface.json :: res.tight.Dsrc.0
\newcommand{\NSADsrcZero}{\ensuremath{1.8092\cdot10^{-7}}}
% NSADsrcOne [cert, upper] A_interface.json :: res.tight.Dsrc.1
\newcommand{\NSADsrcOne}{\ensuremath{9.8276\cdot10^{-5}}}
% NSABZeroDWindowSlopeZero [cert, upper] A_interface.json :: res.tight.B0.D_window_slope.0
\newcommand{\NSABZeroDWindowSlopeZero}{\ensuremath{8.7399\cdot10^{-10}}}
% NSABZeroDWindowSlopeOne [cert, upper] A_interface.json :: res.tight.B0.D_window_slope.1
\newcommand{\NSABZeroDWindowSlopeOne}{\ensuremath{5.2599\cdot10^{-8}}}
% NSAFTwoZero [cert, upper] A_interface.json :: res.tight.F2.0
\newcommand{\NSAFTwoZero}{\ensuremath{5.6506\cdot10^{-13}}}
% NSAFTwoOne [cert, upper] A_interface.json :: res.tight.F2.1
\newcommand{\NSAFTwoOne}{\ensuremath{5.9320\cdot10^{-13}}}
% NSACoupOne [cert, upper] A_interface.json :: res.tight.coup.1
\newcommand{\NSACoupOne}{\ensuremath{1.1568\cdot10^{-12}}}
% NSACOneZeroZero [cert, upper] A_interface.json :: res.tight.C1_0.0
\newcommand{\NSACOneZeroZero}{\ensuremath{7.2060\cdot10^{-10}}}
% NSACOneZeroOne [cert, upper] A_interface.json :: res.tight.C1_0.1
\newcommand{\NSACOneZeroOne}{\ensuremath{4.3143\cdot10^{-8}}}
% NSABZeroSupCoreZero [cert, upper] A_interface.json :: res.tight.B0.sup_core.0
\newcommand{\NSABZeroSupCoreZero}{\ensuremath{4.23\cdot10^{-5}}}
% NSACoreFluxMarginTwoFLUXZeroOverKZero [cert, margin] A_interface.json :: res.tight.core.flux_margin.2.FLUX0_over_K0
\newcommand{\NSACoreFluxMarginTwoFLUXZeroOverKZero}{\ensuremath{3.396}}
% NSACoreFluxMarginTwoFLUXTwoOverKTwo [cert, margin] A_interface.json :: res.tight.core.flux_margin.2.FLUX2_over_K2
\newcommand{\NSACoreFluxMarginTwoFLUXTwoOverKTwo}{\ensuremath{6.012}}
% NSAMarginFluxZero [cert, margin] asm_design_N2ref4.json :: res.tight.flux.beta.2.margin_FLUX0
\newcommand{\NSAMarginFluxZero}{\ensuremath{3.365}}
% NSAMarginFluxTwo [cert, margin] asm_design_N2ref4.json :: res.tight.flux.beta.2.margin_FLUX2
\newcommand{\NSAMarginFluxTwo}{\ensuremath{5.955}}
% NSAfourSZero [cert, upper] asm_design_N2ref4.json :: res.tight.S0
\newcommand{\NSAfourSZero}{\ensuremath{7.1878\cdot10^{-3}}}
% NSAfourSOne [cert, upper] asm_design_N2ref4.json :: res.tight.S1
\newcommand{\NSAfourSOne}{\ensuremath{0.10430}}
% NSAMarginCoreSZeroV [cert, margin] asm_design_N2ref4.json :: res.tight.margins.S0v_vs_0.0153
\newcommand{\NSAMarginCoreSZeroV}{\ensuremath{2.116}}
% NSAMarginCoreSOneV [cert, margin] asm_design_N2ref4.json :: res.tight.margins.S1v_vs_0.6
\newcommand{\NSAMarginCoreSOneV}{\ensuremath{5.748}}
% NSAMarginCoreMOne [cert, margin] asm_design_N2ref4.json :: res.tight.margins.M1core
\newcommand{\NSAMarginCoreMOne}{\ensuremath{4.409}}
% NSAMarginCoreMZero [cert, margin] asm_design_N2ref4.json :: res.tight.margins.M0core
\newcommand{\NSAMarginCoreMZero}{\ensuremath{5.450}}
% NSAMarginGhat [cert, margin] asm_design_N2ref4.json :: res.tight.margins.Ghat22
\newcommand{\NSAMarginGhat}{\ensuremath{1.917}}
% NSAMarginRampMOne [cert, margin] asm_design_N2ref4.json :: res.tight.margins.ramp M1
\newcommand{\NSAMarginRampMOne}{\ensuremath{1.184}}
% NSAfourInfoZeroDWindowConst [cert, upper] asm_design_N2ref4.json :: res.tight.info.0.D_window_const
\newcommand{\NSAfourInfoZeroDWindowConst}{\ensuremath{1.7732\cdot10^{-7}}}
% NSAfourInfoOneDWindowConst [cert, upper] asm_design_N2ref4.json :: res.tight.info.1.D_window_const
\newcommand{\NSAfourInfoOneDWindowConst}{\ensuremath{9.8058\cdot10^{-5}}}
% NZDeltaRZero [cert, upper] zone3_interface.json :: DeltaR.0
\newcommand{\NZDeltaRZero}{\ensuremath{4.4986\cdot10^{-8}}}
% NZDeltaROne [cert, upper] zone3_interface.json :: DeltaR.1
\newcommand{\NZDeltaROne}{\ensuremath{6.7972\cdot10^{-6}}}
% NZDeltaRTwo [cert, upper] zone3_interface.json :: DeltaR.2
\newcommand{\NZDeltaRTwo}{\ensuremath{0.011480}}
% NZDeltaRBZero [cert, upper] zone3_interface.json :: DeltaR_b.0
\newcommand{\NZDeltaRBZero}{\ensuremath{2.66\cdot10^{-12}}}
% NZDeltaRBOne [cert, upper] zone3_interface.json :: DeltaR_b.1
\newcommand{\NZDeltaRBOne}{\ensuremath{3.26\cdot10^{-9}}}
% NZDeltaRBTwo [cert, upper] zone3_interface.json :: DeltaR_b.2
\newcommand{\NZDeltaRBTwo}{\ensuremath{4.35\cdot10^{-8}}}
% NZDeltaRDetailNSqZero [cert, upper] zone3_interface.json :: DeltaR_detail.N_sq.0
\newcommand{\NZDeltaRDetailNSqZero}{\ensuremath{8.28\cdot10^{-5}}}
% NZDeltaRDetailNSqOne [cert, upper] zone3_interface.json :: DeltaR_detail.N_sq.1
\newcommand{\NZDeltaRDetailNSqOne}{\ensuremath{0.10149}}
% NZDeltaRDetailNSqTwo [cert, upper] zone3_interface.json :: DeltaR_detail.N_sq.2
\newcommand{\NZDeltaRDetailNSqTwo}{\ensuremath{1.3567}}
% NZDeltaRDetailNSqThree [cert, upper] zone3_interface.json :: DeltaR_detail.N_sq.3
\newcommand{\NZDeltaRDetailNSqThree}{\ensuremath{25350}}
% NZDeltaRDetailPartsTwoPairsThree [cert, upper] zone3_interface.json :: DeltaR_detail.parts.2.pairs.3
\newcommand{\NZDeltaRDetailPartsTwoPairsThree}{\ensuremath{0.011298}}
% NZDeltaRDetailPartsTwoLFourFoot [cert, upper] zone3_interface.json :: DeltaR_detail.parts.2.L4_foot
\newcommand{\NZDeltaRDetailPartsTwoLFourFoot}{\ensuremath{1.6\cdot10^{-4}}}
% NZDeltaRDetailPartsTwoCubic [cert, upper] zone3_interface.json :: DeltaR_detail.parts.2.cubic
\newcommand{\NZDeltaRDetailPartsTwoCubic}{\ensuremath{9.3\cdot10^{-6}}}
% NZDeltaRDetailPartsTwoSThreeThreeFoot [cert, upper] zone3_interface.json :: DeltaR_detail.parts.2.S33_foot
\newcommand{\NZDeltaRDetailPartsTwoSThreeThreeFoot}{\ensuremath{3.7\cdot10^{-7}}}
% NZNIIIZero [cert, upper] zone3_interface.json :: n_III.0
\newcommand{\NZNIIIZero}{\ensuremath{0.071445}}
% NZNIIIOne [cert, upper] zone3_interface.json :: n_III.1
\newcommand{\NZNIIIOne}{\ensuremath{0.32707}}
% NZNIIITwo [cert, upper] zone3_interface.json :: n_III.2
\newcommand{\NZNIIITwo}{\ensuremath{3.0445}}
% NZTauIII [cert, upper] zone3_interface.json :: tau_III
\newcommand{\NZTauIII}{\ensuremath{4.0708}}
% NZoneIIIBoxZero [cert, upper] zone3_interface.json :: box_III_380.b0
\newcommand{\NZoneIIIBoxZero}{\ensuremath{0.10032}}
% NZoneIIIBoxOne [cert, upper] zone3_interface.json :: box_III_380.b1
\newcommand{\NZoneIIIBoxOne}{\ensuremath{0.41974}}
% NZoneIIIBoxZ [cert, upper] zone3_interface.json :: box_III_380.bz
\newcommand{\NZoneIIIBoxZ}{\ensuremath{0.014291}}
% NZBoxIIIThreeEightZeroMarginsBZero [cert, margin] zone3_interface.json :: box_III_380.margins.b0
\newcommand{\NZBoxIIIThreeEightZeroMarginsBZero}{\ensuremath{1.495}}
% NZoneIIIMarginThinnest [cert, margin] zone3_interface.json :: box_III_380.margins.b1
\newcommand{\NZoneIIIMarginThinnest}{\ensuremath{1.238}}
% NZBoxIIIThreeEightZeroMarginsBz [cert, margin] zone3_interface.json :: box_III_380.margins.bz
\newcommand{\NZBoxIIIThreeEightZeroMarginsBz}{\ensuremath{11.19}}
% NZBoxIIISevenZeroZeroBZero [cert, upper] zone3_interface.json :: box_III_700.b0
\newcommand{\NZBoxIIISevenZeroZeroBZero}{\ensuremath{0.22757}}
% NZBoxIIISevenZeroZeroBOne [cert, upper] zone3_interface.json :: box_III_700.b1
\newcommand{\NZBoxIIISevenZeroZeroBOne}{\ensuremath{0.54234}}
% NZUfineKTwoAtBhatRef [cert, upper] zone3_interface.json :: ufine.K2_at_Bhat_ref
\newcommand{\NZUfineKTwoAtBhatRef}{\ensuremath{0.046956}}
% NZUfineKThreeStripAtBhatRef [cert, upper] zone3_interface.json :: ufine.K3_strip_at_Bhat_ref
\newcommand{\NZUfineKThreeStripAtBhatRef}{\ensuremath{0.054140}}
% NRecoveryMinMargin [cert, margin] zone3_interface.json :: cone_family.H6.at_Bhat_ref.min_margin
\newcommand{\NRecoveryMinMargin}{\ensuremath{1.860}}
% NZConeFamilyHSixAtBhatRefConeSlopeMargin [cert, margin] zone3_interface.json :: cone_family.H6.at_Bhat_ref.cone.slope.margin
\newcommand{\NZConeFamilyHSixAtBhatRefConeSlopeMargin}{\ensuremath{2.46}}
% NZConeFamilyHSixAtBhatRefConeKTwoMargin [cert, margin] zone3_interface.json :: cone_family.H6.at_Bhat_ref.cone.K2.margin
\newcommand{\NZConeFamilyHSixAtBhatRefConeKTwoMargin}{\ensuremath{2.54}}
% NZConeFamilyHSixAtBhatRefConeKThreeMargin [cert, margin] zone3_interface.json :: cone_family.H6.at_Bhat_ref.cone.K3.margin
\newcommand{\NZConeFamilyHSixAtBhatRefConeKThreeMargin}{\ensuremath{2.51}}
% NConeTstarMargin [cert, margin] zone3_interface.json :: acceptance.tstar_margin
\newcommand{\NConeTstarMargin}{\ensuremath{2.30}}
% NUoneTFiveYaddLevelOneTwo [cert, upper] u1_ghat.json :: t5_Yadd.level1.2
\newcommand{\NUoneTFiveYaddLevelOneTwo}{\ensuremath{0.0044}}
% NUoneTFiveYaddLevelTwoTwo [cert, upper] u1_ghat.json :: t5_Yadd.level2.2
\newcommand{\NUoneTFiveYaddLevelTwoTwo}{\ensuremath{0.1081}}
% NUoneTFiveYaddLevelTwoThree [cert, upper] u1_ghat.json :: t5_Yadd.level2.3
\newcommand{\NUoneTFiveYaddLevelTwoThree}{\ensuremath{0.374}}
% NUoneWeightThetaGammaZero [cert, upper] u1_ghat.json :: weight.Theta.Gamma.0
\newcommand{\NUoneWeightThetaGammaZero}{\ensuremath{4.00000000000000000000054}}
% NRKZeroZero [cert, upper] ramp_U.json :: K.0.0
\newcommand{\NRKZeroZero}{\ensuremath{13436.03}}
% NRKOneOne [cert, upper] ramp_U.json :: K.1.1
\newcommand{\NRKOneOne}{\ensuremath{10526.19}}
% NRKTwoTwo [cert, upper] ramp_U.json :: K.2.2
\newcommand{\NRKTwoTwo}{\ensuremath{39482.30}}
% NUoneCasesPrimaryYNextZero [cert, upper] u1_ghat.json :: cases.primary.Y_next.0
\newcommand{\NUoneCasesPrimaryYNextZero}{\ensuremath{457.092}}
% NUoneCasesPrimaryYNextOne [cert, upper] u1_ghat.json :: cases.primary.Y_next.1
\newcommand{\NUoneCasesPrimaryYNextOne}{\ensuremath{356.940}}
% NUoneCasesPrimaryYNextTwo [cert, upper] u1_ghat.json :: cases.primary.Y_next.2
\newcommand{\NUoneCasesPrimaryYNextTwo}{\ensuremath{1320.220}}
% NSfourResultsMNineHTwoTightenedFUOneRowsZeroBound [cert, upper] FINAL_manifest_v33.json :: results.M9_H2_tightened_fU1.rows.0.bound
\newcommand{\NSfourResultsMNineHTwoTightenedFUOneRowsZeroBound}{\ensuremath{3.658\cdot10^{-10}}}
% NSfourResultsMNineHTwoTightenedFUOneRowsZeroMargin [cert, margin] FINAL_manifest_v33.json :: results.M9_H2_tightened_fU1.rows.0.margin
\newcommand{\NSfourResultsMNineHTwoTightenedFUOneRowsZeroMargin}{\ensuremath{27.34}}
% NSfourResultsMNineHTwoTightenedFUOneRowsOneBound [cert, upper] FINAL_manifest_v33.json :: results.M9_H2_tightened_fU1.rows.1.bound
\newcommand{\NSfourResultsMNineHTwoTightenedFUOneRowsOneBound}{\ensuremath{1.777\cdot10^{-9}}}
% NSfourResultsMNineHTwoTightenedFUOneRowsOneMargin [cert, margin] FINAL_manifest_v33.json :: results.M9_H2_tightened_fU1.rows.1.margin
\newcommand{\NSfourResultsMNineHTwoTightenedFUOneRowsOneMargin}{\ensuremath{562.8}}
% NSfourResultsMNineHTwoTightenedFUOneRowsTwoBound [cert, upper] FINAL_manifest_v33.json :: results.M9_H2_tightened_fU1.rows.2.bound
\newcommand{\NSfourResultsMNineHTwoTightenedFUOneRowsTwoBound}{\ensuremath{1.665\cdot10^{-6}}}
% NSfourResultsMNineHTwoTightenedFUOneRowsTwoMargin [cert, margin] FINAL_manifest_v33.json :: results.M9_H2_tightened_fU1.rows.2.margin
\newcommand{\NSfourResultsMNineHTwoTightenedFUOneRowsTwoMargin}{\ensuremath{600.9}}
% NHOneAlphaThr [hand, exact] hand :: 
\newcommand{\NHOneAlphaThr}{\ensuremath{1.5\cdot10^{-4}}}
% NHOneSqrtLamThr [hand, exact] hand :: 
\newcommand{\NHOneSqrtLamThr}{\ensuremath{5\cdot10^{-8}}}
% NYbarStep [hand, exact] hand :: 
\newcommand{\NYbarStep}{\ensuremath{1\cdot10^{-12}}}
% NHTwoThrZero [hand, exact] hand :: 
\newcommand{\NHTwoThrZero}{\ensuremath{1\cdot10^{-8}}}
% NHTwoThrOne [hand, exact] hand :: 
\newcommand{\NHTwoThrOne}{\ensuremath{1\cdot10^{-6}}}
% NHTwoThrTwo [hand, exact] hand :: 
\newcommand{\NHTwoThrTwo}{\ensuremath{1\cdot10^{-3}}}
% NHThreeThrZero [hand, exact] hand :: 
\newcommand{\NHThreeThrZero}{\ensuremath{1\cdot10^{-6}}}
% NHThreeThrOne [hand, exact] hand :: 
\newcommand{\NHThreeThrOne}{\ensuremath{5\cdot10^{-6}}}
% NHThreeThrTwo [hand, exact] hand :: 
\newcommand{\NHThreeThrTwo}{\ensuremath{1\cdot10^{-5}}}
% NHFourThrZero [hand, exact] hand :: 
\newcommand{\NHFourThrZero}{\ensuremath{2\cdot10^{-7}}}
% NHFourThrOne [hand, exact] hand :: 
\newcommand{\NHFourThrOne}{\ensuremath{1\cdot10^{-5}}}
% NHFourThrTwo [hand, exact] hand :: 
\newcommand{\NHFourThrTwo}{\ensuremath{1\cdot10^{-5}}}
% NHFourThrThree [hand, exact] hand :: 
\newcommand{\NHFourThrThree}{\ensuremath{1\cdot10^{-5}}}
% NHFourKfourThr [hand, exact] hand :: 
\newcommand{\NHFourKfourThr}{\ensuremath{0.06}}
% NHSixConeSlopeThr [hand, exact] hand :: 
\newcommand{\NHSixConeSlopeThr}{\ensuremath{0.08}}
% NHSixConeKtwoThr [hand, exact] hand :: 
\newcommand{\NHSixConeKtwoThr}{\ensuremath{0.03}}
% NHSixConeNthreeThr [hand, exact] hand :: 
\newcommand{\NHSixConeNthreeThr}{\ensuremath{0.004}}
% NHSixConeJetThr [hand, exact] hand :: 
\newcommand{\NHSixConeJetThr}{\ensuremath{0.12}}
% NHSixConeHFiveThr [hand, exact] hand :: 
\newcommand{\NHSixConeHFiveThr}{\ensuremath{0.0144}}
% NHSixFarHFiveThr [hand, exact] hand :: 
\newcommand{\NHSixFarHFiveThr}{\ensuremath{0.0225}}
% NLemmaOConeM [hand, exact] hand :: 
\newcommand{\NLemmaOConeM}{\ensuremath{0.12}}
% NLemmaOConeK [hand, exact] hand :: 
\newcommand{\NLemmaOConeK}{\ensuremath{0.18}}
% NLemmaOFarM [hand, exact] hand :: 
\newcommand{\NLemmaOFarM}{\ensuremath{0.15}}
% NLemmaOFarK [hand, exact] hand :: 
\newcommand{\NLemmaOFarK}{\ensuremath{0.3}}
% NCoercivityConst [hand, exact] hand :: 
\newcommand{\NCoercivityConst}{\ensuremath{0.046875}}
% NMOnez [hand, exact] hand :: 
\newcommand{\NMOnez}{\ensuremath{0.534}}
% NBetaA [hand, exact] hand :: 
\newcommand{\NBetaA}{\ensuremath{2}}
% NTOneLower [cert, lower] t1clock.json :: t1.0
\newcommand{\NTOneLower}{\ensuremath{4.5978574775}}
% NTOneUpper [cert, upper] t1clock.json :: t1.1
\newcommand{\NTOneUpper}{\ensuremath{4.5978586144}}
% NLambdaOneLower [cert, lower] t1clock.json :: lnlam.0
\newcommand{\NLambdaOneLower}{\ensuremath{6.49\cdot10^{-5}}}
% NLambdaOneUpper [cert, upper] t1clock.json :: lnlam.1
\newcommand{\NLambdaOneUpper}{\ensuremath{6.75\cdot10^{-5}}}
% NTMinusTOneUpper [cert, upper] t1clock.json :: lnlam.1; 1.5e-4.sup_ab
\newcommand{\NTMinusTOneUpper}{\ensuremath{2.1\cdot10^{-17}}}
% NAbarUpper [cert, upper] final_alpha_points_5e-8.json :: 1.5e-4.sup_ab
\newcommand{\NAbarUpper}{\ensuremath{2.19103\cdot10^{-7}}}
% NBetabarUpper [cert, upper] final_alpha_points_5e-8.json :: 1.5e-4.sup_beta
\newcommand{\NBetabarUpper}{\ensuremath{3.33811\cdot10^{-9}}}
% NBaseKmaxLower [cert, lower] t1clock.json :: base_kmax.0
\newcommand{\NBaseKmaxLower}{\ensuremath{2.279\cdot10^{4}}}
% NBaseKmaxUpper [cert, upper] t1clock.json :: base_kmax.1
\newcommand{\NBaseKmaxUpper}{\ensuremath{2.376\cdot10^{4}}}
% NHOneAlphaBase [cert, upper] summary.json :: alpha_abs_max
\newcommand{\NHOneAlphaBase}{\ensuremath{1.1739\cdot10^{-4}}}
% NHOneAlphaErr [cert, upper] t9_h1_clam_s180.json :: C_alpha; inputs.YSTAR
\newcommand{\NHOneAlphaErr}{\ensuremath{7.21\cdot10^{-13}}}
% NHOneAlphaMargin [cert, margin] summary.json :: alpha_abs_max; C_alpha; inputs.YSTAR
\newcommand{\NHOneAlphaMargin}{\ensuremath{1.27}}
% NHOneSqrtLamBase [cert, upper] t9_h1_clam_s180.json :: sqrtLam_base_Wp
\newcommand{\NHOneSqrtLamBase}{\ensuremath{2.282\cdot10^{-8}}}
% NHOneSqrtLamTotal [cert, upper] t9_h1_clam_s180.json :: sqrtLam_total_Wp
\newcommand{\NHOneSqrtLamTotal}{\ensuremath{2.29\cdot10^{-8}}}
% NHOneSqrtLamMargin [cert, margin] t9_h1_clam_s180.json :: sqrtLam_total_Wp
\newcommand{\NHOneSqrtLamMargin}{\ensuremath{2.18}}
% NHOneYbarLimit [cert, lower] t9_h1_clam_s180.json :: Ybar_limit_H1
\newcommand{\NHOneYbarLimit}{\ensuremath{1.12\cdot10^{-9}}}
% NSwitchYUpper [cert, upper] switch_s2.json :: modulated.Y_raw
\newcommand{\NSwitchYUpper}{\ensuremath{6.21\cdot10^{-13}}}
% NSwitchYMargin [cert, margin] switch_s2.json :: modulated.Y_raw
\newcommand{\NSwitchYMargin}{\ensuremath{1.6}}
% NMinCI [cert, lower] constants.json :: meta.summary.min_c_I
\newcommand{\NMinCI}{\ensuremath{0.05283}}
% NNumCertificates [cert, exact] constants.json :: meta.summary.n_cert
\newcommand{\NNumCertificates}{\ensuremath{26}}
% NMinRhoC [cert, lower] constants.json :: meta.closure.rho_c_min
\newcommand{\NMinRhoC}{\ensuremath{0.04711}}
% NHTwoBaseZero [cert, upper] FINAL_manifest_v33.json :: results.M9_H2_tightened.rows.0.base
\newcommand{\NHTwoBaseZero}{\ensuremath{8.011\cdot10^{-11}}}
% NHTwoBaseOne [cert, upper] FINAL_manifest_v33.json :: results.M9_H2_tightened.rows.1.base
\newcommand{\NHTwoBaseOne}{\ensuremath{1.336\cdot10^{-10}}}
% NHTwoBaseTwo [cert, upper] FINAL_manifest_v33.json :: results.M9_H2_tightened.rows.2.base
\newcommand{\NHTwoBaseTwo}{\ensuremath{2.936\cdot10^{-10}}}
% NHTwoErrWZero [cert, upper] FINAL_manifest_v33.json :: results.M9_H2_tightened.rows.0.with_error
\newcommand{\NHTwoErrWZero}{\ensuremath{8.014\cdot10^{-11}}}
% NHTwoMarginWZero [cert, margin] FINAL_manifest_v33.json :: results.M9_H2_tightened.rows.0.margin
\newcommand{\NHTwoMarginWZero}{\ensuremath{124.78}}
% NHTwoErrWOne [cert, upper] FINAL_manifest_v33.json :: results.M9_H2_tightened.rows.1.with_error
\newcommand{\NHTwoErrWOne}{\ensuremath{4.706\cdot10^{-10}}}
% NHTwoMarginWOne [cert, margin] FINAL_manifest_v33.json :: results.M9_H2_tightened.rows.1.margin
\newcommand{\NHTwoMarginWOne}{\ensuremath{2124.9}}
% NHTwoErrWTwo [cert, upper] FINAL_manifest_v33.json :: results.M9_H2_tightened.rows.2.with_error
\newcommand{\NHTwoErrWTwo}{\ensuremath{1.559\cdot10^{-6}}}
% NHTwoMarginWTwo [cert, margin] FINAL_manifest_v33.json :: results.M9_H2_tightened.rows.2.margin
\newcommand{\NHTwoMarginWTwo}{\ensuremath{641.53}}
% NHThreeBaseZero [cert, upper] FINAL_manifest_v33.json :: results.M9_H2f.rows.0.base
\newcommand{\NHThreeBaseZero}{\ensuremath{3.127\cdot10^{-7}}}
% NHThreeBaseOne [cert, upper] FINAL_manifest_v33.json :: results.M9_H2f.rows.1.base
\newcommand{\NHThreeBaseOne}{\ensuremath{1.004\cdot10^{-6}}}
% NHThreeBaseTwo [cert, upper] FINAL_manifest_v33.json :: results.M9_H2f.rows.2.base
\newcommand{\NHThreeBaseTwo}{\ensuremath{2.431\cdot10^{-6}}}
% NHThreeErrZero [cert, upper] FINAL_manifest_v33.json :: results.M9_H2f.rows.0.with_error
\newcommand{\NHThreeErrZero}{\ensuremath{3.1363\cdot10^{-7}}}
% NHThreeMarginZero [cert, margin] FINAL_manifest_v33.json :: results.M9_H2f.rows.0.margin
\newcommand{\NHThreeMarginZero}{\ensuremath{3.1885}}
% NHThreeErrOne [cert, upper] FINAL_manifest_v33.json :: results.M9_H2f.rows.1.with_error
\newcommand{\NHThreeErrOne}{\ensuremath{1.0034\cdot10^{-6}}}
% NHThreeMarginOne [cert, margin] FINAL_manifest_v33.json :: results.M9_H2f.rows.1.margin
\newcommand{\NHThreeMarginOne}{\ensuremath{4.9834}}
% NHThreeErrTwo [cert, upper] FINAL_manifest_v33.json :: results.M9_H2f.rows.2.with_error
\newcommand{\NHThreeErrTwo}{\ensuremath{2.4351\cdot10^{-6}}}
% NHThreeMarginTwo [cert, margin] FINAL_manifest_v33.json :: results.M9_H2f.rows.2.margin
\newcommand{\NHThreeMarginTwo}{\ensuremath{4.1067}}
% NHFourBaseZero [cert, upper] FINAL_manifest_v33.json :: results.M9_base_sums.combined.windows_max.0
\newcommand{\NHFourBaseZero}{\ensuremath{3.48\cdot10^{-8}}}
% NHFourErrZero [cert, upper] FINAL_manifest_v33.json :: results.M9_H3_H2v32.rows.win.0.with_error_arb
\newcommand{\NHFourErrZero}{\ensuremath{3.5087\cdot10^{-8}}}
% NHFourMarginZero [cert, margin] FINAL_manifest_v33.json :: results.M9_H3_H2v32.rows.win.0.margin_arb
\newcommand{\NHFourMarginZero}{\ensuremath{5.700}}
% NHFourBaseOne [cert, upper] FINAL_manifest_v33.json :: results.M9_base_sums.combined.windows_max.1
\newcommand{\NHFourBaseOne}{\ensuremath{1.27\cdot10^{-7}}}
% NHFourErrOne [cert, upper] FINAL_manifest_v33.json :: results.M9_H3_H2v32.rows.win.1.with_error_arb
\newcommand{\NHFourErrOne}{\ensuremath{1.2616\cdot10^{-7}}}
% NHFourMarginOne [cert, margin] FINAL_manifest_v33.json :: results.M9_H3_H2v32.rows.win.1.margin_arb
\newcommand{\NHFourMarginOne}{\ensuremath{79.26}}
% NHFourBaseTwo [cert, upper] FINAL_manifest_v33.json :: results.M9_base_sums.combined.windows_max.2
\newcommand{\NHFourBaseTwo}{\ensuremath{3.76\cdot10^{-7}}}
% NHFourErrTwo [cert, upper] FINAL_manifest_v33.json :: results.M9_H3_H2v32.rows.win.2.with_error_arb
\newcommand{\NHFourErrTwo}{\ensuremath{3.7724\cdot10^{-7}}}
% NHFourMarginTwo [cert, margin] FINAL_manifest_v33.json :: results.M9_H3_H2v32.rows.win.2.margin_arb
\newcommand{\NHFourMarginTwo}{\ensuremath{26.50}}
% NHFourBaseThree [cert, upper] FINAL_manifest_v33.json :: results.M9_base_sums.combined.windows_max.3
\newcommand{\NHFourBaseThree}{\ensuremath{9.71\cdot10^{-7}}}
% NHFourErrThree [cert, upper] FINAL_manifest_v33.json :: results.M9_H3_H2v32.rows.win.3.with_error_arb
\newcommand{\NHFourErrThree}{\ensuremath{1.2822\cdot10^{-6}}}
% NHFourMarginThree [cert, margin] FINAL_manifest_v33.json :: results.M9_H3_H2v32.rows.win.3.margin_arb
\newcommand{\NHFourMarginThree}{\ensuremath{7.799}}
% NHFourYbarLimit [cert, lower] FINAL_manifest_v33.json :: results.M9_H3_H2v32.binding_Ybar_limit
\newcommand{\NHFourYbarLimit}{\ensuremath{1.575\cdot10^{-11}}}
% NHFourYbarMargin [cert, margin] FINAL_manifest_v33.json :: results.M9_H3_H2v32.binding_Ybar_limit
\newcommand{\NHFourYbarMargin}{\ensuremath{15.75}}
% NHFourKfourBase [cert, upper] amp_A1_w180.json :: inputs.BASE_T1.K4
\newcommand{\NHFourKfourBase}{\ensuremath{3.1633\cdot10^{-4}}}
% NHSixConeSlope [cert, upper] FINAL_manifest_v33.json :: results.cone_fixedL.checks.slope.value
\newcommand{\NHSixConeSlope}{\ensuremath{0.03029}}
% NHSixConeSlopeMargin [cert, margin] FINAL_manifest_v33.json :: results.cone_fixedL.checks.slope.margin
\newcommand{\NHSixConeSlopeMargin}{\ensuremath{2.641}}
% NHSixConeKtwo [cert, upper] FINAL_manifest_v33.json :: results.cone_fixedL.checks.K2.value
\newcommand{\NHSixConeKtwo}{\ensuremath{0.009876}}
% NHSixConeKtwoMargin [cert, margin] FINAL_manifest_v33.json :: results.cone_fixedL.checks.K2.margin
\newcommand{\NHSixConeKtwoMargin}{\ensuremath{3.037}}
% NHSixConeNthree [cert, upper] FINAL_manifest_v33.json :: results.cone_fixedL.checks.N3.value
\newcommand{\NHSixConeNthree}{\ensuremath{1.461\cdot10^{-4}}}
% NHSixConeNthreeMargin [cert, margin] FINAL_manifest_v33.json :: results.cone_fixedL.checks.N3.margin
\newcommand{\NHSixConeNthreeMargin}{\ensuremath{27.38}}
% NHSixConeHFive [cert, upper] FINAL_manifest_v33.json :: results.H4_fixedL.total
\newcommand{\NHSixConeHFive}{\ensuremath{1.2865\cdot10^{-5}}}
% NHSixConeHFiveMargin [cert, margin] FINAL_manifest_v33.json :: results.H4_fixedL.margin
\newcommand{\NHSixConeHFiveMargin}{\ensuremath{1119}}
% NHSixConeTstarNeeded [cert, upper] FINAL_manifest_v33.json :: results.H7_fixedL.ratio_upper
\newcommand{\NHSixConeTstarNeeded}{\ensuremath{1.6000015\cdot10^{-7}}}
% NHSixConeTstarLower [cert, lower] FINAL_manifest_v33.json :: results.H7_fixedL.t_star_lower
\newcommand{\NHSixConeTstarLower}{\ensuremath{3.68334\cdot10^{-7}}}
% NHSixConeTstarMargin [cert, margin] FINAL_manifest_v33.json :: results.H7_fixedL.margin_tstar_over_ratio
\newcommand{\NHSixConeTstarMargin}{\ensuremath{2.302}}
% NHSixConeHFiveRatioPeak [cert, upper] h4loc_A1_w180.json :: windows.constants/amp_A1_w180.json.worst.ratio
\newcommand{\NHSixConeHFiveRatioPeak}{\ensuremath{0.2393}}
% NHSixConeHOneSup [cert, upper] FINAL_manifest_v33.json :: results.H1_fixedL.sup_h0_max
\newcommand{\NHSixConeHOneSup}{\ensuremath{0.2163}}
% NHSixConeHOneH [cert, upper] FINAL_manifest_v33.json :: results.H1_fixedL.H_max
\newcommand{\NHSixConeHOneH}{\ensuremath{0.2185}}
% NHSixFarHFiveRatio [cert, upper] h4loc_v32.json :: windows.constants/amp_A1_w180.json.cylinders.18.ratio
\newcommand{\NHSixFarHFiveRatio}{\ensuremath{0.1761}}
% NHSixFarHFiveMarginOld [cert, margin] h4loc_v32.json :: windows.constants/amp_A1_w180.json.cylinders.18.ratio
\newcommand{\NHSixFarHFiveMarginOld}{\ensuremath{5.68}}
% NHSixFarHFiveMargin [cert, margin] FINAL_manifest_v33.json :: results.M9_H4_far.min_margin
\newcommand{\NHSixFarHFiveMargin}{\ensuremath{15.58}}
% NHSixConeHFiveMarginOld [cert, margin] h4loc_A1_w180.json :: windows.constants/amp_A1_w180.json.worst.ratio
\newcommand{\NHSixConeHFiveMarginOld}{\ensuremath{4.18}}
% NHSixFarKtwo [cert, upper] t9_far_v32.json :: h2
\newcommand{\NHSixFarKtwo}{\ensuremath{1.797\cdot10^{-4}}}
% NHSixFarKtwoThr [hand, exact] t9_far_v32.json :: thr.h2
\newcommand{\NHSixFarKtwoThr}{\ensuremath{3.53\cdot10^{-4}}}
% NHSixFarKtwoMargin [cert, margin] t9_far_v32.json :: margins.h2
\newcommand{\NHSixFarKtwoMargin}{\ensuremath{1.964}}
% NHSixFarKthree [cert, upper] t9_far_v32.json :: h3
\newcommand{\NHSixFarKthree}{\ensuremath{2.595\cdot10^{-8}}}
% NHSixFarKthreeThr [hand, exact] t9_far_v32.json :: thr.h3
\newcommand{\NHSixFarKthreeThr}{\ensuremath{1.17\cdot10^{-5}}}
% NHSixFarKthreeMargin [cert, margin] t9_far_v32.json :: margins.h3
\newcommand{\NHSixFarKthreeMargin}{\ensuremath{450.9}}
% NHSixFarKfour [cert, upper] t9_far_v32.json :: h4
\newcommand{\NHSixFarKfour}{\ensuremath{6.964\cdot10^{-9}}}
% NHSixFarKfourThr [hand, exact] t9_far_v32.json :: thr.h4
\newcommand{\NHSixFarKfourThr}{\ensuremath{1.03\cdot10^{-7}}}
% NHSixFarKfourMargin [cert, margin] t9_far_v32.json :: margins.h4
\newcommand{\NHSixFarKfourMargin}{\ensuremath{14.79}}
% NHSixFarSlope [cert, upper] t9_far_v32.json :: slope
\newcommand{\NHSixFarSlope}{\ensuremath{0.04153}}
% NHSixFarSlopeThr [hand, exact] t9_far_v32.json :: thr.slope
\newcommand{\NHSixFarSlopeThr}{\ensuremath{0.0800}}
% NHSixFarSlopeMargin [cert, margin] t9_far_v32.json :: margins.slope
\newcommand{\NHSixFarSlopeMargin}{\ensuremath{1.926}}
% NHSixFarNthree [cert, upper] t9_far_v32.json :: L2
\newcommand{\NHSixFarNthree}{\ensuremath{5.516\cdot10^{-7}}}
% NHSixFarNthreeThr [hand, exact] t9_far_v32.json :: thr.L2
\newcommand{\NHSixFarNthreeThr}{\ensuremath{3.32\cdot10^{-6}}}
% NHSixFarNthreeMargin [cert, margin] t9_far_v32.json :: margins.L2
\newcommand{\NHSixFarNthreeMargin}{\ensuremath{6.019}}
% NFarTstarMargin [cert, margin] lemmaO_v32_consts.json :: far_v32.case.margin_factor_lower
\newcommand{\NFarTstarMargin}{\ensuremath{15.5}}
% NFarTstarLower [cert, lower] lemmaO_v32_consts.json :: far_v32.case.t_star_lower
\newcommand{\NFarTstarLower}{\ensuremath{9.5108\cdot10^{-8}}}
% NFarTstarNeeded [cert, upper] lemmaO_v32_consts.json :: far_v32.case.needed
\newcommand{\NFarTstarNeeded}{\ensuremath{6.134\cdot10^{-9}}}
% NGammaModes [cert, exact] gamma0_final.json :: M
\newcommand{\NGammaModes}{\ensuremath{3762}}
% NGammaHalfwidth [cert, upper] gamma0_final.json :: b_j_halfwidth
\newcommand{\NGammaHalfwidth}{\ensuremath{8.138\cdot10^{-24}}}
% NGammaContraction [cert, upper] gamma0_final.json :: krawczyk_contraction
\newcommand{\NGammaContraction}{\ensuremath{3.01\cdot10^{-12}}}
% NGammaArcs [cert, exact] gamma0_final.json :: embedded.arcs
\newcommand{\NGammaArcs}{\ensuremath{701}}
% NGammaMaxTurn [cert, upper] gamma0_final.json :: embedded.maxturn
\newcommand{\NGammaMaxTurn}{\ensuremath{0.04997}}
% NGammaMaxPair [cert, upper] gamma0_final.json :: embedded.maxpair
\newcommand{\NGammaMaxPair}{\ensuremath{0.09597}}
% NGammaMinSep [cert, lower] gamma0_final.json :: embedded.minsep
\newcommand{\NGammaMinSep}{\ensuremath{8.89\cdot10^{-3}}}
% NGammaHalfGap [cert, lower] gamma0_final.json :: embedded.gap_cert
\newcommand{\NGammaHalfGap}{\ensuremath{0.0248219}}
% NPhaseISlabs [cert, exact] status.json :: slab
\newcommand{\NPhaseISlabs}{\ensuremath{96}}
% NPhaseIModes [cert, exact] status.json :: last_slab.M
\newcommand{\NPhaseIModes}{\ensuremath{2800}}
% NPhaseIKmaxLower [cert, lower] status.json :: kmax_enclosure.0
\newcommand{\NPhaseIKmaxLower}{\ensuremath{29.352874545}}
% NPhaseIKmaxUpper [cert, upper] status.json :: kmax_enclosure.1
\newcommand{\NPhaseIKmaxUpper}{\ensuremath{29.352874546}}
% NPhaseITLower [cert, lower] status.json :: t
\newcommand{\NPhaseITLower}{\ensuremath{4.597850539485830}}
% NPhaseITUpper [cert, upper] status.json :: t
\newcommand{\NPhaseITUpper}{\ensuremath{4.597850539485831}}
% NPeakSqrtLam [cert, upper] rate_table_fin12f4_a0.00015.json :: sqrtL_max
\newcommand{\NPeakSqrtLam}{\ensuremath{5.308\cdot10^{-8}}}
% NTailRate [cert, lower] rate_table_fin12f4_a0.00015.json :: tail.rho
\newcommand{\NTailRate}{\ensuremath{2.693\cdot10^{-3}}}
% NFloorSqrtLam [cert, upper] rate_table_fin12f4_a0.00015.json :: tail.sqrt_LF
\newcommand{\NFloorSqrtLam}{\ensuremath{9.872\cdot10^{-10}}}
% NAlphaFloor [cert, upper] rate_table_fin12f4_a0.00015.json :: tail.alpha_F
\newcommand{\NAlphaFloor}{\ensuremath{6.867\cdot10^{-7}}}
% NSupBetaVerifier [cert, upper] final_alpha_points_5e-8.json :: 1.5e-4.sup_beta
\newcommand{\NSupBetaVerifier}{\ensuremath{3.33811\cdot10^{-9}}}
% NMOneRamp [cert, upper] final_alpha_points_5e-8.json :: 1.5e-4.M1
\newcommand{\NMOneRamp}{\ensuremath{0.4328}}
% NSITwoPrime [cert, upper] FINAL_manifest_v33.json :: results.F12_sliver.sI2_prime
\newcommand{\NSITwoPrime}{\ensuremath{0.015457}}
% NLoopGainCoarse [cert, upper] U_cert_v35.json :: loop_gain.2.total
\newcommand{\NLoopGainCoarse}{\ensuremath{0.03904}}
% NRateMarginCoarse [cert, margin] U_cert_v35.json :: rate_margin
\newcommand{\NRateMarginCoarse}{\ensuremath{28.2}}
% NSOneEpsLower [cert, lower] bracket2.json :: eps_a
\newcommand{\NSOneEpsLower}{\ensuremath{1.5346198892965\cdot10^{-4}}}
% NSOneEpsUpper [cert, upper] bracket2.json :: eps_b
\newcommand{\NSOneEpsUpper}{\ensuremath{1.5346198899299\cdot10^{-4}}}
% NSOneRowsTip [cert, exact] ext_tip_check.json :: rows
\newcommand{\NSOneRowsTip}{\ensuremath{51246}}
% NSOneTipWidth [cert, upper] ext_tip_check.json :: max_pointwise_width.th
\newcommand{\NSOneTipWidth}{\ensuremath{3.2\cdot10^{-12}}}
% NSOneExtWidth [cert, upper] ext813_check.json :: max_pointwise_width.p
\newcommand{\NSOneExtWidth}{\ensuremath{5.7\cdot10^{-10}}}
% NSOneRowsMC [cert, exact] check_outA5.json :: prof_tip_vs_MC.rows_compared
\newcommand{\NSOneRowsMC}{\ensuremath{1025}}
% NSOneRowsMB [cert, exact] check_outA5.json :: ext813_vs_MB.rows_compared
\newcommand{\NSOneRowsMB}{\ensuremath{501}}
% NSOneMinMargin [cert, margin] check_outA5.json :: min:*.min_relative_margin.*
\newcommand{\NSOneMinMargin}{\ensuremath{1.2\cdot10^{-9}}}
% NTwoSegments [cert, exact] contour_result.json :: len:segments
\newcommand{\NTwoSegments}{\ensuremath{5058}}
% NTwoInteriorBoxes [cert, exact] interior_result.json :: len:
\newcommand{\NTwoInteriorBoxes}{\ensuremath{174}}
% NTwoContourBoxes [cert, exact] eta_cache.json :: len:
\newcommand{\NTwoContourBoxes}{\ensuremath{202}}
% NTwoEtaMax [cert, upper] eta_cache.json :: max:*.eta
\newcommand{\NTwoEtaMax}{\ensuremath{0.1163}}
% NTwoLipMax [cert, upper] eta_cache.json :: max:*.Lip
\newcommand{\NTwoLipMax}{\ensuremath{0.272}}
% NTwoLambdaZero [est, upper] cert_b0.25_th0.0_tau0.03.json :: Lam
\newcommand{\NTwoLambdaZero}{\ensuremath{1.790819}}
% NTwoLambdaOne [est, upper] cert_b0.25_th0.48_tau0.03.json :: Lam
\newcommand{\NTwoLambdaOne}{\ensuremath{1.768282}}
% NSOneDTightLower [doc, lower] doc:README.md :: D ∈ [2.50080627654, 2.50080627668]
\newcommand{\NSOneDTightLower}{\ensuremath{2.50080627654}}
% NSOneDTightUpper [doc, upper] doc:README.md :: D ∈ [2.50080627654, 2.50080627668]
\newcommand{\NSOneDTightUpper}{\ensuremath{2.50080627668}}
% NSOneKTightLower [doc, lower] doc:README.md :: k(0) ∈ [1.53982063762, 1.53982063822]
\newcommand{\NSOneKTightLower}{\ensuremath{1.53982063762}}
% NSOneKTightUpper [doc, upper] doc:README.md :: k(0) ∈ [1.53982063762, 1.53982063822]
\newcommand{\NSOneKTightUpper}{\ensuremath{1.53982063822}}
% NSOneDCondLower [doc, lower] doc:README.md :: D = −p(tip) ∈ [2.500806276619695, 2.500806276620049]
\newcommand{\NSOneDCondLower}{\ensuremath{2.500806276619695}}
% NSOneDCondUpper [doc, upper] doc:README.md :: D = −p(tip) ∈ [2.500806276619695, 2.500806276620049]
\newcommand{\NSOneDCondUpper}{\ensuremath{2.500806276620049}}
% NSOneKCondLower [doc, lower] doc:README.md :: k(0) ∈ [1.539820637871836, 1.5398206378721015]
\newcommand{\NSOneKCondLower}{\ensuremath{1.539820637871836}}
% NSOneKCondUpper [doc, upper] doc:README.md :: k(0) ∈ [1.539820637871836, 1.5398206378721015]
\newcommand{\NSOneKCondUpper}{\ensuremath{1.5398206378721015}}
% NSOneTauTipLower [doc, lower] doc:README.md :: τ_tip ∈ [10.249077651120778, 10.249077651121361]
\newcommand{\NSOneTauTipLower}{\ensuremath{10.249077651120778}}
% NSOneTauTipUpper [doc, upper] doc:README.md :: τ_tip ∈ [10.249077651120778, 10.249077651121361]
\newcommand{\NSOneTauTipUpper}{\ensuremath{10.249077651121361}}
% NSOneStateUncert [doc, upper] doc:README.md :: State uncertainty at σ = 8: 4.4e−13
\newcommand{\NSOneStateUncert}{\ensuremath{4.4\cdot10^{-13}}}
% NFarBoxEprime [doc, margin] doc:STEP4_RESULT.md :: (6.4287, 4.6917, 1.9940, 1.69999995)
\newcommand{\NFarBoxEprime}{\ensuremath{6.42}}
% NFarBoxXEsecond [doc, margin] doc:STEP4_RESULT.md :: (6.4287, 4.6917, 1.9940, 1.69999995)
\newcommand{\NFarBoxXEsecond}{\ensuremath{4.69}}
% NFarBoxXsqEthird [doc, margin] doc:STEP4_RESULT.md :: (6.4287, 4.6917, 1.9940, 1.69999995)
\newcommand{\NFarBoxXsqEthird}{\ensuremath{1.99}}
% NFarBoxE [doc, margin] doc:STEP4_RESULT.md :: (6.4287, 4.6917, 1.9940, 1.69999995)
\newcommand{\NFarBoxE}{\ensuremath{1.69}}
% NHFourKfourErr [doc, upper] doc:STEP4_RESULT.md :: | 5.60e-3 against 0.06 | ×10.7 |
\newcommand{\NHFourKfourErr}{\ensuremath{5.60\cdot10^{-3}}}
% NHFourKfourMargin [doc, margin] doc:STEP4_RESULT.md :: | 5.60e-3 against 0.06 | ×10.7 |
\newcommand{\NHFourKfourMargin}{\ensuremath{10.7}}
% NCoercivityCoefficient [doc, upper] doc:step3_proof_v32.tex :: \norm{e^{\tau\mathcal L}u}_\rho\le14.8\,e^{-0.0213\tau}\norm
\newcommand{\NCoercivityCoefficient}{\ensuremath{14.8}}
% NCoercivityRate [doc, lower] doc:step3_proof_v32.tex :: \norm{e^{\tau\mathcal L}u}_\rho\le14.8\,e^{-0.0213\tau}\norm
\newcommand{\NCoercivityRate}{\ensuremath{0.0213}}
% NRegionRampStart [hand, exact] hand :: 
\newcommand{\NRegionRampStart}{\ensuremath{14}}
% NRegionInterface [hand, exact] hand :: 
\newcommand{\NRegionInterface}{\ensuremath{28}}
% NRegionShellEnd [hand, exact] hand :: 
\newcommand{\NRegionShellEnd}{\ensuremath{60}}
% NRegionSegEnd [hand, exact] hand :: 
\newcommand{\NRegionSegEnd}{\ensuremath{350}}
% NRegionAnnEnd [hand, exact] hand :: 
\newcommand{\NRegionAnnEnd}{\ensuremath{700}}
% NZoneIIILow [hand, exact] hand :: 
\newcommand{\NZoneIIILow}{\ensuremath{253}}
% NZoneIIIHigh [hand, exact] hand :: 
\newcommand{\NZoneIIIHigh}{\ensuremath{379.5}}
% NConeCentreMin [hand, exact] hand :: 
\newcommand{\NConeCentreMin}{\ensuremath{300}}
% NConeCentreMax [hand, exact] hand :: 
\newcommand{\NConeCentreMax}{\ensuremath{380}}
% NFarCentreMin [hand, exact] hand :: 
\newcommand{\NFarCentreMin}{\ensuremath{490}}
% NConeL [hand, exact] hand :: 
\newcommand{\NConeL}{\ensuremath{50}}
% NConeX [hand, exact] hand :: 
\newcommand{\NConeX}{\ensuremath{3}}
% NFarL [hand, exact] hand :: 
\newcommand{\NFarL}{\ensuremath{113}}
% NFarX [hand, exact] hand :: 
\newcommand{\NFarX}{\ensuremath{2}}
 is found.
% ======================================================================

% ---- which document are we in ----------------------------------------
\ifdefined\ifsupplement\else\newif\ifsupplement\fi

% ---- packages ---------------------------------------------------------
\usepackage[T1]{fontenc}
\usepackage[utf8]{inputenc}
\usepackage{lmodern}
\usepackage[margin=1.1in]{geometry}
\usepackage{amsmath,amssymb,amsthm,mathtools,bm}
\usepackage{booktabs,array,longtable,tabularx,multirow}
\usepackage{enumitem}
\usepackage{graphicx}
\usepackage{xcolor}
\usepackage{tikz}
\usetikzlibrary{positioning,arrows.meta,calc}
% xr-hyper must come before hyperref
\usepackage{xr-hyper}
\usepackage[colorlinks=true,linkcolor=blue!50!black,citecolor=green!35!black,urlcolor=blue!50!black]{hyperref}

\raggedbottom
\setlength{\emergencystretch}{2em}
\setlist[enumerate]{leftmargin=*,itemsep=2pt,topsep=3pt}
\setlist[itemize]{leftmargin=*,itemsep=2pt,topsep=3pt}

% ---- numbering --------------------------------------------------------
% Main paper: sections 1,2,...; theorems 2.3 (by section); equations (2.3).
% Supplement: the Parts A..K are the \section's, numbered by letter, so
% theorems are A.1, E.4, ... (do not confuse "Theorem A" of the paper, which
% has no dot, with "Theorem A.1" of Part A).
\ifsupplement
  \renewcommand{\thesection}{\Alph{section}}
\fi
\numberwithin{equation}{section}

% ---- theorem environments (numbered by section, one shared counter) ----
\theoremstyle{plain}
\newtheorem{theorem}{Theorem}[section]
\newtheorem{proposition}[theorem]{Proposition}
\newtheorem{lemma}[theorem]{Lemma}
\newtheorem{corollary}[theorem]{Corollary}
\newtheorem{maintheorem}[theorem]{Main Theorem}
% Theorems A, B, C, D of the paper: own counter, lettered, no section prefix
\newtheorem{thmlet}{Theorem}
\renewcommand{\thethmlet}{\Alph{thmlet}}
\theoremstyle{definition}
\newtheorem{definition}[theorem]{Definition}
\newtheorem{convention}[theorem]{Convention}
\newtheorem{example}[theorem]{Example}
\theoremstyle{remark}
\newtheorem{remark}[theorem]{Remark}

% A list of hypotheses labelled (H1),(H2),... : \begin{hyplist}{H} ... \end{hyplist}
% (also {O}, {L}, {B}, {K}, {W} for the families of section 5.3 of OUTLINE.md).
\newenvironment{hyplist}[1]%
  {\begin{enumerate}[label=\textup{(#1\arabic*)},leftmargin=3em,itemsep=2pt]}%
  {\end{enumerate}}

% ---- trust classes (OUTLINE 4.1, 5.1.3) -------------------------------
% Supplement: explicit tags.  Main paper: certified is the default and
% unmarked; every non-certified number is marked "(floating point)".
\ifsupplement
  \newcommand{\tcert}{\textup{[cert]}}   % arb/outward-rounded, pinned script and key
  \newcommand{\test}{\textup{[est]}}     % floating point or estimate; never in a theorem
  \newcommand{\thand}{\textup{[hand]}}   % proved by hand in the text
  \newcommand{\tlit}{\textup{[lit]}}     % cited from the literature
\else
  \newcommand{\tcert}{}
  \newcommand{\test}{\unskip\ \textup{(floating point)}}
  \newcommand{\thand}{}
  \newcommand{\tlit}{}
\fi

% ---- gaps: visible, counted, an error in a final build -----------------
% Usage: \GAP{one sentence saying what is missing}   (and an entry in ISSUES.md)
\newcounter{gapcount}
\newif\iffinalpaper   % \finalpapertrue makes every \GAP a hard error
\newcommand{\GAP}[1]{\stepcounter{gapcount}\typeout{GAP-MARKER \thegapcount}%
  \iffinalpaper\PackageError{paper}{\string\GAP\space used in a final build}{Remove all gaps.}\fi
  {\color{red}\sffamily\bfseries[GAP~\thegapcount:\ }{\color{red}\sffamily #1}{\color{red}\sffamily\bfseries]}}
\AtEndDocument{\typeout{GAP-MARKER total: \thegapcount}}

% ---- file names and JSON keys (supplement only) -------------------------
% No internal path appears in the main paper; a use there is a warning.
\newcommand{\fn}[1]{%
  \ifsupplement\else\PackageWarning{paper}{\string\fn\space used in the main paper}\fi
  \texttt{\detokenize{#1}}}
\newcommand{\key}[1]{\texttt{\detokenize{#1}}}

% ---- cross references between main paper and supplement ----------------
% Labels carry a unit prefix (section 5.3 of OUTLINE.md and DRAFTER_GUIDE.md):
%   main paper:  intro: obj: s3: s4: s5: s6: s7: s8: s9:   (kind: thm- lem- eq- tab- fig- ...)
%   supplement:  pA: pB: ... pK:
% so all labels are globally unique and the prefixes below are only a second guard.
% \Xref{label}  : reference to the OTHER document (hyperlinked via xr-hyper)
% \Lref{label}  : reference in the SAME document (plain \ref)
\ifsupplement
  \externaldocument[M-]{main}[main.pdf]
  \newcommand{\Xref}[1]{\ref{M-#1}}
\else
  \externaldocument[S-]{supplement}[supplement.pdf]
  \newcommand{\Xref}[1]{\ref{S-#1}}
\fi
\newcommand{\Lref}[1]{\ref{#1}}
% \Supp{pE}{lem-chart} -> "Lemma E.3" style text is produced by the drafter:
%   "Lemma~\Xref{pE:lem-chart} of the supplement".  \SuppPart{E} gives "Part E".
\newcommand{\SuppPart}[1]{Part~#1 of the supplement}

% ---- numbers (generated by unit U0; macros only, no literal certified numbers) ----
\IfFileExists{numbers/numbers_paper.tex}%
  {% numbers_paper.tex -- GENERATED from ledger.json by build_ledger.py; do not edit.
% Rounding: upper bounds up, lower bounds and margins down at the printed digit; margin = threshold / bound.
% Use as $\times\NName$ for margins.  kind: cert certified, est estimate, hand hypothesis constant, doc only a document quotes it.
% NUEhatZero [cert, upper] U_cert_v35.json :: supersolution.EA.0
\newcommand{\NUEhatZero}{\ensuremath{1.2018\cdot10^{-6}}}
% NUEhatOne [cert, upper] U_cert_v35.json :: supersolution.EA.1
\newcommand{\NUEhatOne}{\ensuremath{1.0441\cdot10^{-4}}}
% NUEhatTwo [cert, upper] U_cert_v35.json :: supersolution.EA.2
\newcommand{\NUEhatTwo}{\ensuremath{0.056694}}
% NUBhatZero [cert, upper] U_cert_v35.json :: supersolution.Bhat.0
\newcommand{\NUBhatZero}{\ensuremath{4.8735\cdot10^{-6}}}
% NUBhatOne [cert, upper] U_cert_v35.json :: supersolution.Bhat.1
\newcommand{\NUBhatOne}{\ensuremath{5.5933\cdot10^{-4}}}
% NUBhatTwo [cert, upper] U_cert_v35.json :: supersolution.Bhat.2
\newcommand{\NUBhatTwo}{\ensuremath{0.36675}}
% NUSupersolutionCZeroZero [cert, upper] U_cert_v35.json :: supersolution.C0.0
\newcommand{\NUSupersolutionCZeroZero}{\ensuremath{1.1340\cdot10^{-6}}}
% NUSupersolutionCZeroOne [cert, upper] U_cert_v35.json :: supersolution.C0.1
\newcommand{\NUSupersolutionCZeroOne}{\ensuremath{1.6623\cdot10^{-4}}}
% NUSupersolutionCOneZero [cert, upper] U_cert_v35.json :: supersolution.C1.0
\newcommand{\NUSupersolutionCOneZero}{\ensuremath{1.5241\cdot10^{-9}}}
% NUSupersolutionCOneOne [cert, upper] U_cert_v35.json :: supersolution.C1.1
\newcommand{\NUSupersolutionCOneOne}{\ensuremath{5.5255\cdot10^{-8}}}
% NUSupersolutionNTwoZero [cert, upper] U_cert_v35.json :: supersolution.N2.0
\newcommand{\NUSupersolutionNTwoZero}{\ensuremath{0.071472}}
% NUSupersolutionNTwoOne [cert, upper] U_cert_v35.json :: supersolution.N2.1
\newcommand{\NUSupersolutionNTwoOne}{\ensuremath{0.33012}}
% NUSupersolutionNTwoTwo [cert, upper] U_cert_v35.json :: supersolution.N2.2
\newcommand{\NUSupersolutionNTwoTwo}{\ensuremath{5.0406}}
% NUSupersolutionNTwoOverFourNTwoRefZero [cert, upper] U_cert_v35.json :: supersolution.N2_over_4N2ref.0
\newcommand{\NUSupersolutionNTwoOverFourNTwoRefZero}{\ensuremath{0.18}}
% NUSupersolutionNTwoOverFourNTwoRefOne [cert, upper] U_cert_v35.json :: supersolution.N2_over_4N2ref.1
\newcommand{\NUSupersolutionNTwoOverFourNTwoRefOne}{\ensuremath{0.42}}
% NUSupersolutionNTwoOverFourNTwoRefTwo [cert, upper] U_cert_v35.json :: supersolution.N2_over_4N2ref.2
\newcommand{\NUSupersolutionNTwoOverFourNTwoRefTwo}{\ensuremath{0.26}}
% NUSupersolutionMOneZHypCheckZonesIIISOneTwoBhat [cert, upper] U_cert_v35.json :: supersolution.M1z_hyp_check.zones_I_II.S1_2Bhat
\newcommand{\NUSupersolutionMOneZHypCheckZonesIIISOneTwoBhat}{\ensuremath{0.28585}}
% NUSupersolutionMOneZHypCheckZoneIIIBOnePlusPUOne [cert, upper] U_cert_v35.json :: supersolution.M1z_hyp_check.zone_III.b1_plus_pU1
\newcommand{\NUSupersolutionMOneZHypCheckZoneIIIBOnePlusPUOne}{\ensuremath{0.50786}}
% NUMarginLateMOne [cert, margin] U_cert_v35.json :: margins.late_M1
\newcommand{\NUMarginLateMOne}{\ensuremath{2.353}}
% NUMarginLateMZero [cert, margin] U_cert_v35.json :: margins.late_M0
\newcommand{\NUMarginLateMZero}{\ensuremath{4.892}}
% NUMarginsLateMZeroZ [cert, margin] U_cert_v35.json :: margins.late_M0z
\newcommand{\NUMarginsLateMZeroZ}{\ensuremath{12.99}}
% NUMarginPresentMOne [cert, margin] U_cert_v35.json :: margins.present_M1_min.value
\newcommand{\NUMarginPresentMOne}{\ensuremath{2.770}}
% NUMarginPresentMZero [cert, margin] U_cert_v35.json :: margins.present_M0_min.value
\newcommand{\NUMarginPresentMZero}{\ensuremath{2.225}}
% NUMarginsPresentMZeroZMinValue [cert, margin] U_cert_v35.json :: margins.present_M0z_min.value
\newcommand{\NUMarginsPresentMZeroZMinValue}{\ensuremath{14.37}}
% NULoopGainZeroTotal [cert, upper] U_cert_v35.json :: loop_gain.0.total
\newcommand{\NULoopGainZeroTotal}{\ensuremath{1.23\cdot10^{-7}}}
% NULoopGainOneTotal [cert, upper] U_cert_v35.json :: loop_gain.1.total
\newcommand{\NULoopGainOneTotal}{\ensuremath{0.037169}}
% NULoopGainTwoNeighbour [cert, upper] U_cert_v35.json :: loop_gain.2.neighbour
\newcommand{\NULoopGainTwoNeighbour}{\ensuremath{0.037018}}
% NULoopGainTwoThroughA [cert, upper] U_cert_v35.json :: loop_gain.2.through_A
\newcommand{\NULoopGainTwoThroughA}{\ensuremath{0.0020172}}
% NUfineBhatZero [cert, upper] U_cert_v35.json :: ufine.Bhat.0
\newcommand{\NUfineBhatZero}{\ensuremath{4.8723\cdot10^{-6}}}
% NUfineBhatOne [cert, upper] U_cert_v35.json :: ufine.Bhat.1
\newcommand{\NUfineBhatOne}{\ensuremath{2.6266\cdot10^{-4}}}
% NUfineBhatTwo [cert, upper] U_cert_v35.json :: ufine.Bhat.2
\newcommand{\NUfineBhatTwo}{\ensuremath{0.025842}}
% NUcdeltaU [cert, upper] U_cert_v35.json :: launch.c_dU
\newcommand{\NUcdeltaU}{\ensuremath{1.7187\cdot10^{-8}}}
% NULaunchWmin [cert, lower] U_cert_v35.json :: launch.wmin
\newcommand{\NULaunchWmin}{\ensuremath{0.73492}}
% NULaunchCAD [cert, upper] U_cert_v35.json :: launch.CA_D
\newcommand{\NULaunchCAD}{\ensuremath{2.2429}}
% NULaunchQD [cert, upper] U_cert_v35.json :: launch.qD
\newcommand{\NULaunchQD}{\ensuremath{0.018509}}
% NULaunchDeltaInZero [cert, upper] U_cert_v35.json :: launch.Delta_in.0
\newcommand{\NULaunchDeltaInZero}{\ensuremath{3.131\cdot10^{-12}}}
% NULaunchDeltaInOne [cert, upper] U_cert_v35.json :: launch.Delta_in.1
\newcommand{\NULaunchDeltaInOne}{\ensuremath{3.287\cdot10^{-12}}}
% NULaunchDeltaInTwo [cert, upper] U_cert_v35.json :: launch.Delta_in.2
\newcommand{\NULaunchDeltaInTwo}{\ensuremath{3.451\cdot10^{-12}}}
% NULaunchPhiZero [cert, upper] U_cert_v35.json :: launch.Phi.0
\newcommand{\NULaunchPhiZero}{\ensuremath{6.434\cdot10^{-9}}}
% NULaunchPhiOne [cert, upper] U_cert_v35.json :: launch.Phi.1
\newcommand{\NULaunchPhiOne}{\ensuremath{1.4217\cdot10^{-8}}}
% NULaunchPhiTwo [cert, upper] U_cert_v35.json :: launch.Phi.2
\newcommand{\NULaunchPhiTwo}{\ensuremath{3.8404\cdot10^{-8}}}
% NUGammaZero [cert, upper] U_cert_v35.json :: Gamma.0
\newcommand{\NUGammaZero}{\ensuremath{4.0000358}}
% NUGammaOne [cert, upper] U_cert_v35.json :: Gamma.1
\newcommand{\NUGammaOne}{\ensuremath{16.000322}}
% NUGammaTwo [cert, upper] U_cert_v35.json :: Gamma.2
\newcommand{\NUGammaTwo}{\ensuremath{64.00218}}
% NUKappa [cert, upper] U_cert_v35.json :: kappa
\newcommand{\NUKappa}{\ensuremath{8.0000358}}
% NUKappaTwo [cert, upper] U_cert_v35.json :: kappa2
\newcommand{\NUKappaTwo}{\ensuremath{64.00061}}
% NUGhatZeroZero [cert, upper] U_cert_v35.json :: Ghat.0.0
\newcommand{\NUGhatZeroZero}{\ensuremath{51322.6}}
% NUGhatOneOne [cert, upper] U_cert_v35.json :: Ghat.1.1
\newcommand{\NUGhatOneOne}{\ensuremath{33962.9}}
% NUGhatTwoTwo [cert, upper] U_cert_v35.json :: Ghat.2.2
\newcommand{\NUGhatTwoTwo}{\ensuremath{348341.1}}
% NUGhatOneZero [cert, upper] U_cert_v35.json :: Ghat.1.0
\newcommand{\NUGhatOneZero}{\ensuremath{4.938}}
% NUGhatTwoZero [cert, upper] U_cert_v35.json :: Ghat.2.0
\newcommand{\NUGhatTwoZero}{\ensuremath{27.984}}
% NUGhatTwoOne [cert, upper] U_cert_v35.json :: Ghat.2.1
\newcommand{\NUGhatTwoOne}{\ensuremath{458.53}}
% NUKZeroZero [cert, upper] U_cert_v35.json :: K.0.0
\newcommand{\NUKZeroZero}{\ensuremath{8034.21}}
% NUKZeroOne [cert, upper] U_cert_v35.json :: K.0.1
\newcommand{\NUKZeroOne}{\ensuremath{724.89}}
% NUKOneZero [cert, upper] U_cert_v35.json :: K.1.0
\newcommand{\NUKOneZero}{\ensuremath{0.08152}}
% NUKOneOne [cert, upper] U_cert_v35.json :: K.1.1
\newcommand{\NUKOneOne}{\ensuremath{6336.94}}
% NUKOneTwo [cert, upper] U_cert_v35.json :: K.1.2
\newcommand{\NUKOneTwo}{\ensuremath{545.59}}
% NUKOneThree [cert, upper] U_cert_v35.json :: K.1.3
\newcommand{\NUKOneThree}{\ensuremath{0.58700}}
% NUKTwoZero [cert, upper] U_cert_v35.json :: K.2.0
\newcommand{\NUKTwoZero}{\ensuremath{0.87186}}
% NUKTwoOne [cert, upper] U_cert_v35.json :: K.2.1
\newcommand{\NUKTwoOne}{\ensuremath{528.62}}
% NUKTwoTwo [cert, upper] U_cert_v35.json :: K.2.2
\newcommand{\NUKTwoTwo}{\ensuremath{23425.98}}
% NUKTwoThree [cert, upper] U_cert_v35.json :: K.2.3
\newcommand{\NUKTwoThree}{\ensuremath{555.62}}
% NUFootwOne [cert, upper] U_cert_v35.json :: footw.1
\newcommand{\NUFootwOne}{\ensuremath{2.6361}}
% NUFootwTwo [cert, upper] U_cert_v35.json :: footw.2
\newcommand{\NUFootwTwo}{\ensuremath{2663.3}}
% NUCThreeThree [cert, upper] U_cert_v35.json :: c33
\newcommand{\NUCThreeThree}{\ensuremath{108.02}}
% NUCTwoTwoThree [cert, upper] U_cert_v35.json :: c223
\newcommand{\NUCTwoTwoThree}{\ensuremath{7.3284}}
% NUCEta [cert, upper] U_cert_v35.json :: C_eta
\newcommand{\NUCEta}{\ensuremath{13.376}}
% NULMax [cert, upper] U_cert_v35.json :: L_max
\newcommand{\NULMax}{\ensuremath{3.5352}}
% NURRatio [cert, upper] U_cert_v35.json :: R_ratio
\newcommand{\NURRatio}{\ensuremath{11.593}}
% NUSweepsSigmaInCoverSup [cert, upper] U_cert_v35.json :: sweeps.Sigma_in.cover_sup
\newcommand{\NUSweepsSigmaInCoverSup}{\ensuremath{5.0892}}
% NUSweepsSigmaECoverSup [cert, upper] U_cert_v35.json :: sweeps.Sigma_E.cover_sup
\newcommand{\NUSweepsSigmaECoverSup}{\ensuremath{0.30846}}
% NUSweepsSXCoverSup [cert, upper] U_cert_v35.json :: sweeps.S_X.cover_sup
\newcommand{\NUSweepsSXCoverSup}{\ensuremath{1.5243}}
% NUSweepsSigmaXCoverSup [cert, upper] U_cert_v35.json :: sweeps.Sigma_X.cover_sup
\newcommand{\NUSweepsSigmaXCoverSup}{\ensuremath{3.0081}}
% NUSweepsSigmaRsTwoCoverSup [cert, upper] U_cert_v35.json :: sweeps.Sigma_rs.2.cover_sup
\newcommand{\NUSweepsSigmaRsTwoCoverSup}{\ensuremath{4.3350\cdot10^{-3}}}
% NUDataTermsLateDZero [cert, upper] U_cert_v35.json :: data_terms_late.D.0
\newcommand{\NUDataTermsLateDZero}{\ensuremath{1.7330\cdot10^{-6}}}
% NUDataTermsLateDOne [cert, upper] U_cert_v35.json :: data_terms_late.D.1
\newcommand{\NUDataTermsLateDOne}{\ensuremath{2.5353\cdot10^{-4}}}
% NUDataTermsLateSlopeShareCOneSigmaXOverDkZero [cert, upper] U_cert_v35.json :: data_terms_late.slope_share_C1SigmaX_over_Dk.0
\newcommand{\NUDataTermsLateSlopeShareCOneSigmaXOverDkZero}{\ensuremath{2.7\cdot10^{-3}}}
% NUDataTermsLateSlopeShareCOneSigmaXOverDkOne [cert, upper] U_cert_v35.json :: data_terms_late.slope_share_C1SigmaX_over_Dk.1
\newcommand{\NUDataTermsLateSlopeShareCOneSigmaXOverDkOne}{\ensuremath{6.6\cdot10^{-4}}}
% NURLtTwoMarginZero [cert, margin] U_cert_v35.json :: R_lt2_margin.0
\newcommand{\NURLtTwoMarginZero}{\ensuremath{4.432}}
% NURLtTwoMarginOne [cert, margin] U_cert_v35.json :: R_lt2_margin.1
\newcommand{\NURLtTwoMarginOne}{\ensuremath{6.697}}
% NUThetaTwoEight [cert, upper] U_cert_v35.json :: theta28
\newcommand{\NUThetaTwoEight}{\ensuremath{0.37011}}
% NUDFiveSix [cert, upper] U_cert_v35.json :: D56
\newcommand{\NUDFiveSix}{\ensuremath{1.0368}}
% NUDTwoEight [cert, upper] U_cert_v35.json :: D28
\newcommand{\NUDTwoEight}{\ensuremath{1.5876}}
% NUSThreeThreeFootOverBhatTwo [cert, upper] U_cert_v35.json :: S33_foot_over_Bhat2
\newcommand{\NUSThreeThreeFootOverBhatTwo}{\ensuremath{1.69\cdot10^{-3}}}
% NSAHJWKappaJ [cert, upper] A_interface.json :: HJW.kappa_J
\newcommand{\NSAHJWKappaJ}{\ensuremath{0.042351}}
% NSAGpZero [cert, lower] A_interface.json :: res.tight.gp.0
\newcommand{\NSAGpZero}{\ensuremath{0.24292}}
% NSAGpOne [cert, lower] A_interface.json :: res.tight.gp.1
\newcommand{\NSAGpOne}{\ensuremath{0.24106}}
% NSAGpTwo [cert, lower] A_interface.json :: res.tight.gp.2
\newcommand{\NSAGpTwo}{\ensuremath{0.11961}}
% NSAAminC [cert, lower] A_interface.json :: res.tight.amin_c
\newcommand{\NSAAminC}{\ensuremath{0.97623}}
% NSARhoJ [cert, upper] A_interface.json :: res.tight.rhoJ
\newcommand{\NSARhoJ}{\ensuremath{5.4995\cdot10^{-7}}}
% NSAKappaOTwoTwo [cert, upper] A_interface.json :: res.tight.kappa_o.2.2
\newcommand{\NSAKappaOTwoTwo}{\ensuremath{1.2416\cdot10^{-5}}}
% NSADsrcZero [cert, upper] A_interface.json :: res.tight.Dsrc.0
\newcommand{\NSADsrcZero}{\ensuremath{1.8092\cdot10^{-7}}}
% NSADsrcOne [cert, upper] A_interface.json :: res.tight.Dsrc.1
\newcommand{\NSADsrcOne}{\ensuremath{9.8276\cdot10^{-5}}}
% NSABZeroDWindowSlopeZero [cert, upper] A_interface.json :: res.tight.B0.D_window_slope.0
\newcommand{\NSABZeroDWindowSlopeZero}{\ensuremath{8.7399\cdot10^{-10}}}
% NSABZeroDWindowSlopeOne [cert, upper] A_interface.json :: res.tight.B0.D_window_slope.1
\newcommand{\NSABZeroDWindowSlopeOne}{\ensuremath{5.2599\cdot10^{-8}}}
% NSAFTwoZero [cert, upper] A_interface.json :: res.tight.F2.0
\newcommand{\NSAFTwoZero}{\ensuremath{5.6506\cdot10^{-13}}}
% NSAFTwoOne [cert, upper] A_interface.json :: res.tight.F2.1
\newcommand{\NSAFTwoOne}{\ensuremath{5.9320\cdot10^{-13}}}
% NSACoupOne [cert, upper] A_interface.json :: res.tight.coup.1
\newcommand{\NSACoupOne}{\ensuremath{1.1568\cdot10^{-12}}}
% NSACOneZeroZero [cert, upper] A_interface.json :: res.tight.C1_0.0
\newcommand{\NSACOneZeroZero}{\ensuremath{7.2060\cdot10^{-10}}}
% NSACOneZeroOne [cert, upper] A_interface.json :: res.tight.C1_0.1
\newcommand{\NSACOneZeroOne}{\ensuremath{4.3143\cdot10^{-8}}}
% NSABZeroSupCoreZero [cert, upper] A_interface.json :: res.tight.B0.sup_core.0
\newcommand{\NSABZeroSupCoreZero}{\ensuremath{4.23\cdot10^{-5}}}
% NSACoreFluxMarginTwoFLUXZeroOverKZero [cert, margin] A_interface.json :: res.tight.core.flux_margin.2.FLUX0_over_K0
\newcommand{\NSACoreFluxMarginTwoFLUXZeroOverKZero}{\ensuremath{3.396}}
% NSACoreFluxMarginTwoFLUXTwoOverKTwo [cert, margin] A_interface.json :: res.tight.core.flux_margin.2.FLUX2_over_K2
\newcommand{\NSACoreFluxMarginTwoFLUXTwoOverKTwo}{\ensuremath{6.012}}
% NSAMarginFluxZero [cert, margin] asm_design_N2ref4.json :: res.tight.flux.beta.2.margin_FLUX0
\newcommand{\NSAMarginFluxZero}{\ensuremath{3.365}}
% NSAMarginFluxTwo [cert, margin] asm_design_N2ref4.json :: res.tight.flux.beta.2.margin_FLUX2
\newcommand{\NSAMarginFluxTwo}{\ensuremath{5.955}}
% NSAfourSZero [cert, upper] asm_design_N2ref4.json :: res.tight.S0
\newcommand{\NSAfourSZero}{\ensuremath{7.1878\cdot10^{-3}}}
% NSAfourSOne [cert, upper] asm_design_N2ref4.json :: res.tight.S1
\newcommand{\NSAfourSOne}{\ensuremath{0.10430}}
% NSAMarginCoreSZeroV [cert, margin] asm_design_N2ref4.json :: res.tight.margins.S0v_vs_0.0153
\newcommand{\NSAMarginCoreSZeroV}{\ensuremath{2.116}}
% NSAMarginCoreSOneV [cert, margin] asm_design_N2ref4.json :: res.tight.margins.S1v_vs_0.6
\newcommand{\NSAMarginCoreSOneV}{\ensuremath{5.748}}
% NSAMarginCoreMOne [cert, margin] asm_design_N2ref4.json :: res.tight.margins.M1core
\newcommand{\NSAMarginCoreMOne}{\ensuremath{4.409}}
% NSAMarginCoreMZero [cert, margin] asm_design_N2ref4.json :: res.tight.margins.M0core
\newcommand{\NSAMarginCoreMZero}{\ensuremath{5.450}}
% NSAMarginGhat [cert, margin] asm_design_N2ref4.json :: res.tight.margins.Ghat22
\newcommand{\NSAMarginGhat}{\ensuremath{1.917}}
% NSAMarginRampMOne [cert, margin] asm_design_N2ref4.json :: res.tight.margins.ramp M1
\newcommand{\NSAMarginRampMOne}{\ensuremath{1.184}}
% NSAfourInfoZeroDWindowConst [cert, upper] asm_design_N2ref4.json :: res.tight.info.0.D_window_const
\newcommand{\NSAfourInfoZeroDWindowConst}{\ensuremath{1.7732\cdot10^{-7}}}
% NSAfourInfoOneDWindowConst [cert, upper] asm_design_N2ref4.json :: res.tight.info.1.D_window_const
\newcommand{\NSAfourInfoOneDWindowConst}{\ensuremath{9.8058\cdot10^{-5}}}
% NZDeltaRZero [cert, upper] zone3_interface.json :: DeltaR.0
\newcommand{\NZDeltaRZero}{\ensuremath{4.4986\cdot10^{-8}}}
% NZDeltaROne [cert, upper] zone3_interface.json :: DeltaR.1
\newcommand{\NZDeltaROne}{\ensuremath{6.7972\cdot10^{-6}}}
% NZDeltaRTwo [cert, upper] zone3_interface.json :: DeltaR.2
\newcommand{\NZDeltaRTwo}{\ensuremath{0.011480}}
% NZDeltaRBZero [cert, upper] zone3_interface.json :: DeltaR_b.0
\newcommand{\NZDeltaRBZero}{\ensuremath{2.66\cdot10^{-12}}}
% NZDeltaRBOne [cert, upper] zone3_interface.json :: DeltaR_b.1
\newcommand{\NZDeltaRBOne}{\ensuremath{3.26\cdot10^{-9}}}
% NZDeltaRBTwo [cert, upper] zone3_interface.json :: DeltaR_b.2
\newcommand{\NZDeltaRBTwo}{\ensuremath{4.35\cdot10^{-8}}}
% NZDeltaRDetailNSqZero [cert, upper] zone3_interface.json :: DeltaR_detail.N_sq.0
\newcommand{\NZDeltaRDetailNSqZero}{\ensuremath{8.28\cdot10^{-5}}}
% NZDeltaRDetailNSqOne [cert, upper] zone3_interface.json :: DeltaR_detail.N_sq.1
\newcommand{\NZDeltaRDetailNSqOne}{\ensuremath{0.10149}}
% NZDeltaRDetailNSqTwo [cert, upper] zone3_interface.json :: DeltaR_detail.N_sq.2
\newcommand{\NZDeltaRDetailNSqTwo}{\ensuremath{1.3567}}
% NZDeltaRDetailNSqThree [cert, upper] zone3_interface.json :: DeltaR_detail.N_sq.3
\newcommand{\NZDeltaRDetailNSqThree}{\ensuremath{25350}}
% NZDeltaRDetailPartsTwoPairsThree [cert, upper] zone3_interface.json :: DeltaR_detail.parts.2.pairs.3
\newcommand{\NZDeltaRDetailPartsTwoPairsThree}{\ensuremath{0.011298}}
% NZDeltaRDetailPartsTwoLFourFoot [cert, upper] zone3_interface.json :: DeltaR_detail.parts.2.L4_foot
\newcommand{\NZDeltaRDetailPartsTwoLFourFoot}{\ensuremath{1.6\cdot10^{-4}}}
% NZDeltaRDetailPartsTwoCubic [cert, upper] zone3_interface.json :: DeltaR_detail.parts.2.cubic
\newcommand{\NZDeltaRDetailPartsTwoCubic}{\ensuremath{9.3\cdot10^{-6}}}
% NZDeltaRDetailPartsTwoSThreeThreeFoot [cert, upper] zone3_interface.json :: DeltaR_detail.parts.2.S33_foot
\newcommand{\NZDeltaRDetailPartsTwoSThreeThreeFoot}{\ensuremath{3.7\cdot10^{-7}}}
% NZNIIIZero [cert, upper] zone3_interface.json :: n_III.0
\newcommand{\NZNIIIZero}{\ensuremath{0.071445}}
% NZNIIIOne [cert, upper] zone3_interface.json :: n_III.1
\newcommand{\NZNIIIOne}{\ensuremath{0.32707}}
% NZNIIITwo [cert, upper] zone3_interface.json :: n_III.2
\newcommand{\NZNIIITwo}{\ensuremath{3.0445}}
% NZTauIII [cert, upper] zone3_interface.json :: tau_III
\newcommand{\NZTauIII}{\ensuremath{4.0708}}
% NZoneIIIBoxZero [cert, upper] zone3_interface.json :: box_III_380.b0
\newcommand{\NZoneIIIBoxZero}{\ensuremath{0.10032}}
% NZoneIIIBoxOne [cert, upper] zone3_interface.json :: box_III_380.b1
\newcommand{\NZoneIIIBoxOne}{\ensuremath{0.41974}}
% NZoneIIIBoxZ [cert, upper] zone3_interface.json :: box_III_380.bz
\newcommand{\NZoneIIIBoxZ}{\ensuremath{0.014291}}
% NZBoxIIIThreeEightZeroMarginsBZero [cert, margin] zone3_interface.json :: box_III_380.margins.b0
\newcommand{\NZBoxIIIThreeEightZeroMarginsBZero}{\ensuremath{1.495}}
% NZoneIIIMarginThinnest [cert, margin] zone3_interface.json :: box_III_380.margins.b1
\newcommand{\NZoneIIIMarginThinnest}{\ensuremath{1.238}}
% NZBoxIIIThreeEightZeroMarginsBz [cert, margin] zone3_interface.json :: box_III_380.margins.bz
\newcommand{\NZBoxIIIThreeEightZeroMarginsBz}{\ensuremath{11.19}}
% NZBoxIIISevenZeroZeroBZero [cert, upper] zone3_interface.json :: box_III_700.b0
\newcommand{\NZBoxIIISevenZeroZeroBZero}{\ensuremath{0.22757}}
% NZBoxIIISevenZeroZeroBOne [cert, upper] zone3_interface.json :: box_III_700.b1
\newcommand{\NZBoxIIISevenZeroZeroBOne}{\ensuremath{0.54234}}
% NZUfineKTwoAtBhatRef [cert, upper] zone3_interface.json :: ufine.K2_at_Bhat_ref
\newcommand{\NZUfineKTwoAtBhatRef}{\ensuremath{0.046956}}
% NZUfineKThreeStripAtBhatRef [cert, upper] zone3_interface.json :: ufine.K3_strip_at_Bhat_ref
\newcommand{\NZUfineKThreeStripAtBhatRef}{\ensuremath{0.054140}}
% NRecoveryMinMargin [cert, margin] zone3_interface.json :: cone_family.H6.at_Bhat_ref.min_margin
\newcommand{\NRecoveryMinMargin}{\ensuremath{1.860}}
% NZConeFamilyHSixAtBhatRefConeSlopeMargin [cert, margin] zone3_interface.json :: cone_family.H6.at_Bhat_ref.cone.slope.margin
\newcommand{\NZConeFamilyHSixAtBhatRefConeSlopeMargin}{\ensuremath{2.46}}
% NZConeFamilyHSixAtBhatRefConeKTwoMargin [cert, margin] zone3_interface.json :: cone_family.H6.at_Bhat_ref.cone.K2.margin
\newcommand{\NZConeFamilyHSixAtBhatRefConeKTwoMargin}{\ensuremath{2.54}}
% NZConeFamilyHSixAtBhatRefConeKThreeMargin [cert, margin] zone3_interface.json :: cone_family.H6.at_Bhat_ref.cone.K3.margin
\newcommand{\NZConeFamilyHSixAtBhatRefConeKThreeMargin}{\ensuremath{2.51}}
% NConeTstarMargin [cert, margin] zone3_interface.json :: acceptance.tstar_margin
\newcommand{\NConeTstarMargin}{\ensuremath{2.30}}
% NUoneTFiveYaddLevelOneTwo [cert, upper] u1_ghat.json :: t5_Yadd.level1.2
\newcommand{\NUoneTFiveYaddLevelOneTwo}{\ensuremath{0.0044}}
% NUoneTFiveYaddLevelTwoTwo [cert, upper] u1_ghat.json :: t5_Yadd.level2.2
\newcommand{\NUoneTFiveYaddLevelTwoTwo}{\ensuremath{0.1081}}
% NUoneTFiveYaddLevelTwoThree [cert, upper] u1_ghat.json :: t5_Yadd.level2.3
\newcommand{\NUoneTFiveYaddLevelTwoThree}{\ensuremath{0.374}}
% NUoneWeightThetaGammaZero [cert, upper] u1_ghat.json :: weight.Theta.Gamma.0
\newcommand{\NUoneWeightThetaGammaZero}{\ensuremath{4.00000000000000000000054}}
% NRKZeroZero [cert, upper] ramp_U.json :: K.0.0
\newcommand{\NRKZeroZero}{\ensuremath{13436.03}}
% NRKOneOne [cert, upper] ramp_U.json :: K.1.1
\newcommand{\NRKOneOne}{\ensuremath{10526.19}}
% NRKTwoTwo [cert, upper] ramp_U.json :: K.2.2
\newcommand{\NRKTwoTwo}{\ensuremath{39482.30}}
% NUoneCasesPrimaryYNextZero [cert, upper] u1_ghat.json :: cases.primary.Y_next.0
\newcommand{\NUoneCasesPrimaryYNextZero}{\ensuremath{457.092}}
% NUoneCasesPrimaryYNextOne [cert, upper] u1_ghat.json :: cases.primary.Y_next.1
\newcommand{\NUoneCasesPrimaryYNextOne}{\ensuremath{356.940}}
% NUoneCasesPrimaryYNextTwo [cert, upper] u1_ghat.json :: cases.primary.Y_next.2
\newcommand{\NUoneCasesPrimaryYNextTwo}{\ensuremath{1320.220}}
% NSfourResultsMNineHTwoTightenedFUOneRowsZeroBound [cert, upper] FINAL_manifest_v33.json :: results.M9_H2_tightened_fU1.rows.0.bound
\newcommand{\NSfourResultsMNineHTwoTightenedFUOneRowsZeroBound}{\ensuremath{3.658\cdot10^{-10}}}
% NSfourResultsMNineHTwoTightenedFUOneRowsZeroMargin [cert, margin] FINAL_manifest_v33.json :: results.M9_H2_tightened_fU1.rows.0.margin
\newcommand{\NSfourResultsMNineHTwoTightenedFUOneRowsZeroMargin}{\ensuremath{27.34}}
% NSfourResultsMNineHTwoTightenedFUOneRowsOneBound [cert, upper] FINAL_manifest_v33.json :: results.M9_H2_tightened_fU1.rows.1.bound
\newcommand{\NSfourResultsMNineHTwoTightenedFUOneRowsOneBound}{\ensuremath{1.777\cdot10^{-9}}}
% NSfourResultsMNineHTwoTightenedFUOneRowsOneMargin [cert, margin] FINAL_manifest_v33.json :: results.M9_H2_tightened_fU1.rows.1.margin
\newcommand{\NSfourResultsMNineHTwoTightenedFUOneRowsOneMargin}{\ensuremath{562.8}}
% NSfourResultsMNineHTwoTightenedFUOneRowsTwoBound [cert, upper] FINAL_manifest_v33.json :: results.M9_H2_tightened_fU1.rows.2.bound
\newcommand{\NSfourResultsMNineHTwoTightenedFUOneRowsTwoBound}{\ensuremath{1.665\cdot10^{-6}}}
% NSfourResultsMNineHTwoTightenedFUOneRowsTwoMargin [cert, margin] FINAL_manifest_v33.json :: results.M9_H2_tightened_fU1.rows.2.margin
\newcommand{\NSfourResultsMNineHTwoTightenedFUOneRowsTwoMargin}{\ensuremath{600.9}}
% NHOneAlphaThr [hand, exact] hand :: 
\newcommand{\NHOneAlphaThr}{\ensuremath{1.5\cdot10^{-4}}}
% NHOneSqrtLamThr [hand, exact] hand :: 
\newcommand{\NHOneSqrtLamThr}{\ensuremath{5\cdot10^{-8}}}
% NYbarStep [hand, exact] hand :: 
\newcommand{\NYbarStep}{\ensuremath{1\cdot10^{-12}}}
% NHTwoThrZero [hand, exact] hand :: 
\newcommand{\NHTwoThrZero}{\ensuremath{1\cdot10^{-8}}}
% NHTwoThrOne [hand, exact] hand :: 
\newcommand{\NHTwoThrOne}{\ensuremath{1\cdot10^{-6}}}
% NHTwoThrTwo [hand, exact] hand :: 
\newcommand{\NHTwoThrTwo}{\ensuremath{1\cdot10^{-3}}}
% NHThreeThrZero [hand, exact] hand :: 
\newcommand{\NHThreeThrZero}{\ensuremath{1\cdot10^{-6}}}
% NHThreeThrOne [hand, exact] hand :: 
\newcommand{\NHThreeThrOne}{\ensuremath{5\cdot10^{-6}}}
% NHThreeThrTwo [hand, exact] hand :: 
\newcommand{\NHThreeThrTwo}{\ensuremath{1\cdot10^{-5}}}
% NHFourThrZero [hand, exact] hand :: 
\newcommand{\NHFourThrZero}{\ensuremath{2\cdot10^{-7}}}
% NHFourThrOne [hand, exact] hand :: 
\newcommand{\NHFourThrOne}{\ensuremath{1\cdot10^{-5}}}
% NHFourThrTwo [hand, exact] hand :: 
\newcommand{\NHFourThrTwo}{\ensuremath{1\cdot10^{-5}}}
% NHFourThrThree [hand, exact] hand :: 
\newcommand{\NHFourThrThree}{\ensuremath{1\cdot10^{-5}}}
% NHFourKfourThr [hand, exact] hand :: 
\newcommand{\NHFourKfourThr}{\ensuremath{0.06}}
% NHSixConeSlopeThr [hand, exact] hand :: 
\newcommand{\NHSixConeSlopeThr}{\ensuremath{0.08}}
% NHSixConeKtwoThr [hand, exact] hand :: 
\newcommand{\NHSixConeKtwoThr}{\ensuremath{0.03}}
% NHSixConeNthreeThr [hand, exact] hand :: 
\newcommand{\NHSixConeNthreeThr}{\ensuremath{0.004}}
% NHSixConeJetThr [hand, exact] hand :: 
\newcommand{\NHSixConeJetThr}{\ensuremath{0.12}}
% NHSixConeHFiveThr [hand, exact] hand :: 
\newcommand{\NHSixConeHFiveThr}{\ensuremath{0.0144}}
% NHSixFarHFiveThr [hand, exact] hand :: 
\newcommand{\NHSixFarHFiveThr}{\ensuremath{0.0225}}
% NLemmaOConeM [hand, exact] hand :: 
\newcommand{\NLemmaOConeM}{\ensuremath{0.12}}
% NLemmaOConeK [hand, exact] hand :: 
\newcommand{\NLemmaOConeK}{\ensuremath{0.18}}
% NLemmaOFarM [hand, exact] hand :: 
\newcommand{\NLemmaOFarM}{\ensuremath{0.15}}
% NLemmaOFarK [hand, exact] hand :: 
\newcommand{\NLemmaOFarK}{\ensuremath{0.3}}
% NCoercivityConst [hand, exact] hand :: 
\newcommand{\NCoercivityConst}{\ensuremath{0.046875}}
% NMOnez [hand, exact] hand :: 
\newcommand{\NMOnez}{\ensuremath{0.534}}
% NBetaA [hand, exact] hand :: 
\newcommand{\NBetaA}{\ensuremath{2}}
% NTOneLower [cert, lower] t1clock.json :: t1.0
\newcommand{\NTOneLower}{\ensuremath{4.5978574775}}
% NTOneUpper [cert, upper] t1clock.json :: t1.1
\newcommand{\NTOneUpper}{\ensuremath{4.5978586144}}
% NLambdaOneLower [cert, lower] t1clock.json :: lnlam.0
\newcommand{\NLambdaOneLower}{\ensuremath{6.49\cdot10^{-5}}}
% NLambdaOneUpper [cert, upper] t1clock.json :: lnlam.1
\newcommand{\NLambdaOneUpper}{\ensuremath{6.75\cdot10^{-5}}}
% NTMinusTOneUpper [cert, upper] t1clock.json :: lnlam.1; 1.5e-4.sup_ab
\newcommand{\NTMinusTOneUpper}{\ensuremath{2.1\cdot10^{-17}}}
% NAbarUpper [cert, upper] final_alpha_points_5e-8.json :: 1.5e-4.sup_ab
\newcommand{\NAbarUpper}{\ensuremath{2.19103\cdot10^{-7}}}
% NBetabarUpper [cert, upper] final_alpha_points_5e-8.json :: 1.5e-4.sup_beta
\newcommand{\NBetabarUpper}{\ensuremath{3.33811\cdot10^{-9}}}
% NBaseKmaxLower [cert, lower] t1clock.json :: base_kmax.0
\newcommand{\NBaseKmaxLower}{\ensuremath{2.279\cdot10^{4}}}
% NBaseKmaxUpper [cert, upper] t1clock.json :: base_kmax.1
\newcommand{\NBaseKmaxUpper}{\ensuremath{2.376\cdot10^{4}}}
% NHOneAlphaBase [cert, upper] summary.json :: alpha_abs_max
\newcommand{\NHOneAlphaBase}{\ensuremath{1.1739\cdot10^{-4}}}
% NHOneAlphaErr [cert, upper] t9_h1_clam_s180.json :: C_alpha; inputs.YSTAR
\newcommand{\NHOneAlphaErr}{\ensuremath{7.21\cdot10^{-13}}}
% NHOneAlphaMargin [cert, margin] summary.json :: alpha_abs_max; C_alpha; inputs.YSTAR
\newcommand{\NHOneAlphaMargin}{\ensuremath{1.27}}
% NHOneSqrtLamBase [cert, upper] t9_h1_clam_s180.json :: sqrtLam_base_Wp
\newcommand{\NHOneSqrtLamBase}{\ensuremath{2.282\cdot10^{-8}}}
% NHOneSqrtLamTotal [cert, upper] t9_h1_clam_s180.json :: sqrtLam_total_Wp
\newcommand{\NHOneSqrtLamTotal}{\ensuremath{2.29\cdot10^{-8}}}
% NHOneSqrtLamMargin [cert, margin] t9_h1_clam_s180.json :: sqrtLam_total_Wp
\newcommand{\NHOneSqrtLamMargin}{\ensuremath{2.18}}
% NHOneYbarLimit [cert, lower] t9_h1_clam_s180.json :: Ybar_limit_H1
\newcommand{\NHOneYbarLimit}{\ensuremath{1.12\cdot10^{-9}}}
% NSwitchYUpper [cert, upper] switch_s2.json :: modulated.Y_raw
\newcommand{\NSwitchYUpper}{\ensuremath{6.21\cdot10^{-13}}}
% NSwitchYMargin [cert, margin] switch_s2.json :: modulated.Y_raw
\newcommand{\NSwitchYMargin}{\ensuremath{1.6}}
% NMinCI [cert, lower] constants.json :: meta.summary.min_c_I
\newcommand{\NMinCI}{\ensuremath{0.05283}}
% NNumCertificates [cert, exact] constants.json :: meta.summary.n_cert
\newcommand{\NNumCertificates}{\ensuremath{26}}
% NMinRhoC [cert, lower] constants.json :: meta.closure.rho_c_min
\newcommand{\NMinRhoC}{\ensuremath{0.04711}}
% NHTwoBaseZero [cert, upper] FINAL_manifest_v33.json :: results.M9_H2_tightened.rows.0.base
\newcommand{\NHTwoBaseZero}{\ensuremath{8.011\cdot10^{-11}}}
% NHTwoBaseOne [cert, upper] FINAL_manifest_v33.json :: results.M9_H2_tightened.rows.1.base
\newcommand{\NHTwoBaseOne}{\ensuremath{1.336\cdot10^{-10}}}
% NHTwoBaseTwo [cert, upper] FINAL_manifest_v33.json :: results.M9_H2_tightened.rows.2.base
\newcommand{\NHTwoBaseTwo}{\ensuremath{2.936\cdot10^{-10}}}
% NHTwoErrWZero [cert, upper] FINAL_manifest_v33.json :: results.M9_H2_tightened.rows.0.with_error
\newcommand{\NHTwoErrWZero}{\ensuremath{8.014\cdot10^{-11}}}
% NHTwoMarginWZero [cert, margin] FINAL_manifest_v33.json :: results.M9_H2_tightened.rows.0.margin
\newcommand{\NHTwoMarginWZero}{\ensuremath{124.78}}
% NHTwoErrWOne [cert, upper] FINAL_manifest_v33.json :: results.M9_H2_tightened.rows.1.with_error
\newcommand{\NHTwoErrWOne}{\ensuremath{4.706\cdot10^{-10}}}
% NHTwoMarginWOne [cert, margin] FINAL_manifest_v33.json :: results.M9_H2_tightened.rows.1.margin
\newcommand{\NHTwoMarginWOne}{\ensuremath{2124.9}}
% NHTwoErrWTwo [cert, upper] FINAL_manifest_v33.json :: results.M9_H2_tightened.rows.2.with_error
\newcommand{\NHTwoErrWTwo}{\ensuremath{1.559\cdot10^{-6}}}
% NHTwoMarginWTwo [cert, margin] FINAL_manifest_v33.json :: results.M9_H2_tightened.rows.2.margin
\newcommand{\NHTwoMarginWTwo}{\ensuremath{641.53}}
% NHThreeBaseZero [cert, upper] FINAL_manifest_v33.json :: results.M9_H2f.rows.0.base
\newcommand{\NHThreeBaseZero}{\ensuremath{3.127\cdot10^{-7}}}
% NHThreeBaseOne [cert, upper] FINAL_manifest_v33.json :: results.M9_H2f.rows.1.base
\newcommand{\NHThreeBaseOne}{\ensuremath{1.004\cdot10^{-6}}}
% NHThreeBaseTwo [cert, upper] FINAL_manifest_v33.json :: results.M9_H2f.rows.2.base
\newcommand{\NHThreeBaseTwo}{\ensuremath{2.431\cdot10^{-6}}}
% NHThreeErrZero [cert, upper] FINAL_manifest_v33.json :: results.M9_H2f.rows.0.with_error
\newcommand{\NHThreeErrZero}{\ensuremath{3.1363\cdot10^{-7}}}
% NHThreeMarginZero [cert, margin] FINAL_manifest_v33.json :: results.M9_H2f.rows.0.margin
\newcommand{\NHThreeMarginZero}{\ensuremath{3.1885}}
% NHThreeErrOne [cert, upper] FINAL_manifest_v33.json :: results.M9_H2f.rows.1.with_error
\newcommand{\NHThreeErrOne}{\ensuremath{1.0034\cdot10^{-6}}}
% NHThreeMarginOne [cert, margin] FINAL_manifest_v33.json :: results.M9_H2f.rows.1.margin
\newcommand{\NHThreeMarginOne}{\ensuremath{4.9834}}
% NHThreeErrTwo [cert, upper] FINAL_manifest_v33.json :: results.M9_H2f.rows.2.with_error
\newcommand{\NHThreeErrTwo}{\ensuremath{2.4351\cdot10^{-6}}}
% NHThreeMarginTwo [cert, margin] FINAL_manifest_v33.json :: results.M9_H2f.rows.2.margin
\newcommand{\NHThreeMarginTwo}{\ensuremath{4.1067}}
% NHFourBaseZero [cert, upper] FINAL_manifest_v33.json :: results.M9_base_sums.combined.windows_max.0
\newcommand{\NHFourBaseZero}{\ensuremath{3.48\cdot10^{-8}}}
% NHFourErrZero [cert, upper] FINAL_manifest_v33.json :: results.M9_H3_H2v32.rows.win.0.with_error_arb
\newcommand{\NHFourErrZero}{\ensuremath{3.5087\cdot10^{-8}}}
% NHFourMarginZero [cert, margin] FINAL_manifest_v33.json :: results.M9_H3_H2v32.rows.win.0.margin_arb
\newcommand{\NHFourMarginZero}{\ensuremath{5.700}}
% NHFourBaseOne [cert, upper] FINAL_manifest_v33.json :: results.M9_base_sums.combined.windows_max.1
\newcommand{\NHFourBaseOne}{\ensuremath{1.27\cdot10^{-7}}}
% NHFourErrOne [cert, upper] FINAL_manifest_v33.json :: results.M9_H3_H2v32.rows.win.1.with_error_arb
\newcommand{\NHFourErrOne}{\ensuremath{1.2616\cdot10^{-7}}}
% NHFourMarginOne [cert, margin] FINAL_manifest_v33.json :: results.M9_H3_H2v32.rows.win.1.margin_arb
\newcommand{\NHFourMarginOne}{\ensuremath{79.26}}
% NHFourBaseTwo [cert, upper] FINAL_manifest_v33.json :: results.M9_base_sums.combined.windows_max.2
\newcommand{\NHFourBaseTwo}{\ensuremath{3.76\cdot10^{-7}}}
% NHFourErrTwo [cert, upper] FINAL_manifest_v33.json :: results.M9_H3_H2v32.rows.win.2.with_error_arb
\newcommand{\NHFourErrTwo}{\ensuremath{3.7724\cdot10^{-7}}}
% NHFourMarginTwo [cert, margin] FINAL_manifest_v33.json :: results.M9_H3_H2v32.rows.win.2.margin_arb
\newcommand{\NHFourMarginTwo}{\ensuremath{26.50}}
% NHFourBaseThree [cert, upper] FINAL_manifest_v33.json :: results.M9_base_sums.combined.windows_max.3
\newcommand{\NHFourBaseThree}{\ensuremath{9.71\cdot10^{-7}}}
% NHFourErrThree [cert, upper] FINAL_manifest_v33.json :: results.M9_H3_H2v32.rows.win.3.with_error_arb
\newcommand{\NHFourErrThree}{\ensuremath{1.2822\cdot10^{-6}}}
% NHFourMarginThree [cert, margin] FINAL_manifest_v33.json :: results.M9_H3_H2v32.rows.win.3.margin_arb
\newcommand{\NHFourMarginThree}{\ensuremath{7.799}}
% NHFourYbarLimit [cert, lower] FINAL_manifest_v33.json :: results.M9_H3_H2v32.binding_Ybar_limit
\newcommand{\NHFourYbarLimit}{\ensuremath{1.575\cdot10^{-11}}}
% NHFourYbarMargin [cert, margin] FINAL_manifest_v33.json :: results.M9_H3_H2v32.binding_Ybar_limit
\newcommand{\NHFourYbarMargin}{\ensuremath{15.75}}
% NHFourKfourBase [cert, upper] amp_A1_w180.json :: inputs.BASE_T1.K4
\newcommand{\NHFourKfourBase}{\ensuremath{3.1633\cdot10^{-4}}}
% NHSixConeSlope [cert, upper] FINAL_manifest_v33.json :: results.cone_fixedL.checks.slope.value
\newcommand{\NHSixConeSlope}{\ensuremath{0.03029}}
% NHSixConeSlopeMargin [cert, margin] FINAL_manifest_v33.json :: results.cone_fixedL.checks.slope.margin
\newcommand{\NHSixConeSlopeMargin}{\ensuremath{2.641}}
% NHSixConeKtwo [cert, upper] FINAL_manifest_v33.json :: results.cone_fixedL.checks.K2.value
\newcommand{\NHSixConeKtwo}{\ensuremath{0.009876}}
% NHSixConeKtwoMargin [cert, margin] FINAL_manifest_v33.json :: results.cone_fixedL.checks.K2.margin
\newcommand{\NHSixConeKtwoMargin}{\ensuremath{3.037}}
% NHSixConeNthree [cert, upper] FINAL_manifest_v33.json :: results.cone_fixedL.checks.N3.value
\newcommand{\NHSixConeNthree}{\ensuremath{1.461\cdot10^{-4}}}
% NHSixConeNthreeMargin [cert, margin] FINAL_manifest_v33.json :: results.cone_fixedL.checks.N3.margin
\newcommand{\NHSixConeNthreeMargin}{\ensuremath{27.38}}
% NHSixConeHFive [cert, upper] FINAL_manifest_v33.json :: results.H4_fixedL.total
\newcommand{\NHSixConeHFive}{\ensuremath{1.2865\cdot10^{-5}}}
% NHSixConeHFiveMargin [cert, margin] FINAL_manifest_v33.json :: results.H4_fixedL.margin
\newcommand{\NHSixConeHFiveMargin}{\ensuremath{1119}}
% NHSixConeTstarNeeded [cert, upper] FINAL_manifest_v33.json :: results.H7_fixedL.ratio_upper
\newcommand{\NHSixConeTstarNeeded}{\ensuremath{1.6000015\cdot10^{-7}}}
% NHSixConeTstarLower [cert, lower] FINAL_manifest_v33.json :: results.H7_fixedL.t_star_lower
\newcommand{\NHSixConeTstarLower}{\ensuremath{3.68334\cdot10^{-7}}}
% NHSixConeTstarMargin [cert, margin] FINAL_manifest_v33.json :: results.H7_fixedL.margin_tstar_over_ratio
\newcommand{\NHSixConeTstarMargin}{\ensuremath{2.302}}
% NHSixConeHFiveRatioPeak [cert, upper] h4loc_A1_w180.json :: windows.constants/amp_A1_w180.json.worst.ratio
\newcommand{\NHSixConeHFiveRatioPeak}{\ensuremath{0.2393}}
% NHSixConeHOneSup [cert, upper] FINAL_manifest_v33.json :: results.H1_fixedL.sup_h0_max
\newcommand{\NHSixConeHOneSup}{\ensuremath{0.2163}}
% NHSixConeHOneH [cert, upper] FINAL_manifest_v33.json :: results.H1_fixedL.H_max
\newcommand{\NHSixConeHOneH}{\ensuremath{0.2185}}
% NHSixFarHFiveRatio [cert, upper] h4loc_v32.json :: windows.constants/amp_A1_w180.json.cylinders.18.ratio
\newcommand{\NHSixFarHFiveRatio}{\ensuremath{0.1761}}
% NHSixFarHFiveMarginOld [cert, margin] h4loc_v32.json :: windows.constants/amp_A1_w180.json.cylinders.18.ratio
\newcommand{\NHSixFarHFiveMarginOld}{\ensuremath{5.68}}
% NHSixFarHFiveMargin [cert, margin] FINAL_manifest_v33.json :: results.M9_H4_far.min_margin
\newcommand{\NHSixFarHFiveMargin}{\ensuremath{15.58}}
% NHSixConeHFiveMarginOld [cert, margin] h4loc_A1_w180.json :: windows.constants/amp_A1_w180.json.worst.ratio
\newcommand{\NHSixConeHFiveMarginOld}{\ensuremath{4.18}}
% NHSixFarKtwo [cert, upper] t9_far_v32.json :: h2
\newcommand{\NHSixFarKtwo}{\ensuremath{1.797\cdot10^{-4}}}
% NHSixFarKtwoThr [hand, exact] t9_far_v32.json :: thr.h2
\newcommand{\NHSixFarKtwoThr}{\ensuremath{3.53\cdot10^{-4}}}
% NHSixFarKtwoMargin [cert, margin] t9_far_v32.json :: margins.h2
\newcommand{\NHSixFarKtwoMargin}{\ensuremath{1.964}}
% NHSixFarKthree [cert, upper] t9_far_v32.json :: h3
\newcommand{\NHSixFarKthree}{\ensuremath{2.595\cdot10^{-8}}}
% NHSixFarKthreeThr [hand, exact] t9_far_v32.json :: thr.h3
\newcommand{\NHSixFarKthreeThr}{\ensuremath{1.17\cdot10^{-5}}}
% NHSixFarKthreeMargin [cert, margin] t9_far_v32.json :: margins.h3
\newcommand{\NHSixFarKthreeMargin}{\ensuremath{450.9}}
% NHSixFarKfour [cert, upper] t9_far_v32.json :: h4
\newcommand{\NHSixFarKfour}{\ensuremath{6.964\cdot10^{-9}}}
% NHSixFarKfourThr [hand, exact] t9_far_v32.json :: thr.h4
\newcommand{\NHSixFarKfourThr}{\ensuremath{1.03\cdot10^{-7}}}
% NHSixFarKfourMargin [cert, margin] t9_far_v32.json :: margins.h4
\newcommand{\NHSixFarKfourMargin}{\ensuremath{14.79}}
% NHSixFarSlope [cert, upper] t9_far_v32.json :: slope
\newcommand{\NHSixFarSlope}{\ensuremath{0.04153}}
% NHSixFarSlopeThr [hand, exact] t9_far_v32.json :: thr.slope
\newcommand{\NHSixFarSlopeThr}{\ensuremath{0.0800}}
% NHSixFarSlopeMargin [cert, margin] t9_far_v32.json :: margins.slope
\newcommand{\NHSixFarSlopeMargin}{\ensuremath{1.926}}
% NHSixFarNthree [cert, upper] t9_far_v32.json :: L2
\newcommand{\NHSixFarNthree}{\ensuremath{5.516\cdot10^{-7}}}
% NHSixFarNthreeThr [hand, exact] t9_far_v32.json :: thr.L2
\newcommand{\NHSixFarNthreeThr}{\ensuremath{3.32\cdot10^{-6}}}
% NHSixFarNthreeMargin [cert, margin] t9_far_v32.json :: margins.L2
\newcommand{\NHSixFarNthreeMargin}{\ensuremath{6.019}}
% NFarTstarMargin [cert, margin] lemmaO_v32_consts.json :: far_v32.case.margin_factor_lower
\newcommand{\NFarTstarMargin}{\ensuremath{15.5}}
% NFarTstarLower [cert, lower] lemmaO_v32_consts.json :: far_v32.case.t_star_lower
\newcommand{\NFarTstarLower}{\ensuremath{9.5108\cdot10^{-8}}}
% NFarTstarNeeded [cert, upper] lemmaO_v32_consts.json :: far_v32.case.needed
\newcommand{\NFarTstarNeeded}{\ensuremath{6.134\cdot10^{-9}}}
% NGammaModes [cert, exact] gamma0_final.json :: M
\newcommand{\NGammaModes}{\ensuremath{3762}}
% NGammaHalfwidth [cert, upper] gamma0_final.json :: b_j_halfwidth
\newcommand{\NGammaHalfwidth}{\ensuremath{8.138\cdot10^{-24}}}
% NGammaContraction [cert, upper] gamma0_final.json :: krawczyk_contraction
\newcommand{\NGammaContraction}{\ensuremath{3.01\cdot10^{-12}}}
% NGammaArcs [cert, exact] gamma0_final.json :: embedded.arcs
\newcommand{\NGammaArcs}{\ensuremath{701}}
% NGammaMaxTurn [cert, upper] gamma0_final.json :: embedded.maxturn
\newcommand{\NGammaMaxTurn}{\ensuremath{0.04997}}
% NGammaMaxPair [cert, upper] gamma0_final.json :: embedded.maxpair
\newcommand{\NGammaMaxPair}{\ensuremath{0.09597}}
% NGammaMinSep [cert, lower] gamma0_final.json :: embedded.minsep
\newcommand{\NGammaMinSep}{\ensuremath{8.89\cdot10^{-3}}}
% NGammaHalfGap [cert, lower] gamma0_final.json :: embedded.gap_cert
\newcommand{\NGammaHalfGap}{\ensuremath{0.0248219}}
% NPhaseISlabs [cert, exact] status.json :: slab
\newcommand{\NPhaseISlabs}{\ensuremath{96}}
% NPhaseIModes [cert, exact] status.json :: last_slab.M
\newcommand{\NPhaseIModes}{\ensuremath{2800}}
% NPhaseIKmaxLower [cert, lower] status.json :: kmax_enclosure.0
\newcommand{\NPhaseIKmaxLower}{\ensuremath{29.352874545}}
% NPhaseIKmaxUpper [cert, upper] status.json :: kmax_enclosure.1
\newcommand{\NPhaseIKmaxUpper}{\ensuremath{29.352874546}}
% NPhaseITLower [cert, lower] status.json :: t
\newcommand{\NPhaseITLower}{\ensuremath{4.597850539485830}}
% NPhaseITUpper [cert, upper] status.json :: t
\newcommand{\NPhaseITUpper}{\ensuremath{4.597850539485831}}
% NPeakSqrtLam [cert, upper] rate_table_fin12f4_a0.00015.json :: sqrtL_max
\newcommand{\NPeakSqrtLam}{\ensuremath{5.308\cdot10^{-8}}}
% NTailRate [cert, lower] rate_table_fin12f4_a0.00015.json :: tail.rho
\newcommand{\NTailRate}{\ensuremath{2.693\cdot10^{-3}}}
% NFloorSqrtLam [cert, upper] rate_table_fin12f4_a0.00015.json :: tail.sqrt_LF
\newcommand{\NFloorSqrtLam}{\ensuremath{9.872\cdot10^{-10}}}
% NAlphaFloor [cert, upper] rate_table_fin12f4_a0.00015.json :: tail.alpha_F
\newcommand{\NAlphaFloor}{\ensuremath{6.867\cdot10^{-7}}}
% NSupBetaVerifier [cert, upper] final_alpha_points_5e-8.json :: 1.5e-4.sup_beta
\newcommand{\NSupBetaVerifier}{\ensuremath{3.33811\cdot10^{-9}}}
% NMOneRamp [cert, upper] final_alpha_points_5e-8.json :: 1.5e-4.M1
\newcommand{\NMOneRamp}{\ensuremath{0.4328}}
% NSITwoPrime [cert, upper] FINAL_manifest_v33.json :: results.F12_sliver.sI2_prime
\newcommand{\NSITwoPrime}{\ensuremath{0.015457}}
% NLoopGainCoarse [cert, upper] U_cert_v35.json :: loop_gain.2.total
\newcommand{\NLoopGainCoarse}{\ensuremath{0.03904}}
% NRateMarginCoarse [cert, margin] U_cert_v35.json :: rate_margin
\newcommand{\NRateMarginCoarse}{\ensuremath{28.2}}
% NSOneEpsLower [cert, lower] bracket2.json :: eps_a
\newcommand{\NSOneEpsLower}{\ensuremath{1.5346198892965\cdot10^{-4}}}
% NSOneEpsUpper [cert, upper] bracket2.json :: eps_b
\newcommand{\NSOneEpsUpper}{\ensuremath{1.5346198899299\cdot10^{-4}}}
% NSOneRowsTip [cert, exact] ext_tip_check.json :: rows
\newcommand{\NSOneRowsTip}{\ensuremath{51246}}
% NSOneTipWidth [cert, upper] ext_tip_check.json :: max_pointwise_width.th
\newcommand{\NSOneTipWidth}{\ensuremath{3.2\cdot10^{-12}}}
% NSOneExtWidth [cert, upper] ext813_check.json :: max_pointwise_width.p
\newcommand{\NSOneExtWidth}{\ensuremath{5.7\cdot10^{-10}}}
% NSOneRowsMC [cert, exact] check_outA5.json :: prof_tip_vs_MC.rows_compared
\newcommand{\NSOneRowsMC}{\ensuremath{1025}}
% NSOneRowsMB [cert, exact] check_outA5.json :: ext813_vs_MB.rows_compared
\newcommand{\NSOneRowsMB}{\ensuremath{501}}
% NSOneMinMargin [cert, margin] check_outA5.json :: min:*.min_relative_margin.*
\newcommand{\NSOneMinMargin}{\ensuremath{1.2\cdot10^{-9}}}
% NTwoSegments [cert, exact] contour_result.json :: len:segments
\newcommand{\NTwoSegments}{\ensuremath{5058}}
% NTwoInteriorBoxes [cert, exact] interior_result.json :: len:
\newcommand{\NTwoInteriorBoxes}{\ensuremath{174}}
% NTwoContourBoxes [cert, exact] eta_cache.json :: len:
\newcommand{\NTwoContourBoxes}{\ensuremath{202}}
% NTwoEtaMax [cert, upper] eta_cache.json :: max:*.eta
\newcommand{\NTwoEtaMax}{\ensuremath{0.1163}}
% NTwoLipMax [cert, upper] eta_cache.json :: max:*.Lip
\newcommand{\NTwoLipMax}{\ensuremath{0.272}}
% NTwoLambdaZero [est, upper] cert_b0.25_th0.0_tau0.03.json :: Lam
\newcommand{\NTwoLambdaZero}{\ensuremath{1.790819}}
% NTwoLambdaOne [est, upper] cert_b0.25_th0.48_tau0.03.json :: Lam
\newcommand{\NTwoLambdaOne}{\ensuremath{1.768282}}
% NSOneDTightLower [doc, lower] doc:README.md :: D ∈ [2.50080627654, 2.50080627668]
\newcommand{\NSOneDTightLower}{\ensuremath{2.50080627654}}
% NSOneDTightUpper [doc, upper] doc:README.md :: D ∈ [2.50080627654, 2.50080627668]
\newcommand{\NSOneDTightUpper}{\ensuremath{2.50080627668}}
% NSOneKTightLower [doc, lower] doc:README.md :: k(0) ∈ [1.53982063762, 1.53982063822]
\newcommand{\NSOneKTightLower}{\ensuremath{1.53982063762}}
% NSOneKTightUpper [doc, upper] doc:README.md :: k(0) ∈ [1.53982063762, 1.53982063822]
\newcommand{\NSOneKTightUpper}{\ensuremath{1.53982063822}}
% NSOneDCondLower [doc, lower] doc:README.md :: D = −p(tip) ∈ [2.500806276619695, 2.500806276620049]
\newcommand{\NSOneDCondLower}{\ensuremath{2.500806276619695}}
% NSOneDCondUpper [doc, upper] doc:README.md :: D = −p(tip) ∈ [2.500806276619695, 2.500806276620049]
\newcommand{\NSOneDCondUpper}{\ensuremath{2.500806276620049}}
% NSOneKCondLower [doc, lower] doc:README.md :: k(0) ∈ [1.539820637871836, 1.5398206378721015]
\newcommand{\NSOneKCondLower}{\ensuremath{1.539820637871836}}
% NSOneKCondUpper [doc, upper] doc:README.md :: k(0) ∈ [1.539820637871836, 1.5398206378721015]
\newcommand{\NSOneKCondUpper}{\ensuremath{1.5398206378721015}}
% NSOneTauTipLower [doc, lower] doc:README.md :: τ_tip ∈ [10.249077651120778, 10.249077651121361]
\newcommand{\NSOneTauTipLower}{\ensuremath{10.249077651120778}}
% NSOneTauTipUpper [doc, upper] doc:README.md :: τ_tip ∈ [10.249077651120778, 10.249077651121361]
\newcommand{\NSOneTauTipUpper}{\ensuremath{10.249077651121361}}
% NSOneStateUncert [doc, upper] doc:README.md :: State uncertainty at σ = 8: 4.4e−13
\newcommand{\NSOneStateUncert}{\ensuremath{4.4\cdot10^{-13}}}
% NFarBoxEprime [doc, margin] doc:STEP4_RESULT.md :: (6.4287, 4.6917, 1.9940, 1.69999995)
\newcommand{\NFarBoxEprime}{\ensuremath{6.42}}
% NFarBoxXEsecond [doc, margin] doc:STEP4_RESULT.md :: (6.4287, 4.6917, 1.9940, 1.69999995)
\newcommand{\NFarBoxXEsecond}{\ensuremath{4.69}}
% NFarBoxXsqEthird [doc, margin] doc:STEP4_RESULT.md :: (6.4287, 4.6917, 1.9940, 1.69999995)
\newcommand{\NFarBoxXsqEthird}{\ensuremath{1.99}}
% NFarBoxE [doc, margin] doc:STEP4_RESULT.md :: (6.4287, 4.6917, 1.9940, 1.69999995)
\newcommand{\NFarBoxE}{\ensuremath{1.69}}
% NHFourKfourErr [doc, upper] doc:STEP4_RESULT.md :: | 5.60e-3 against 0.06 | ×10.7 |
\newcommand{\NHFourKfourErr}{\ensuremath{5.60\cdot10^{-3}}}
% NHFourKfourMargin [doc, margin] doc:STEP4_RESULT.md :: | 5.60e-3 against 0.06 | ×10.7 |
\newcommand{\NHFourKfourMargin}{\ensuremath{10.7}}
% NCoercivityCoefficient [doc, upper] doc:step3_proof_v32.tex :: \norm{e^{\tau\mathcal L}u}_\rho\le14.8\,e^{-0.0213\tau}\norm
\newcommand{\NCoercivityCoefficient}{\ensuremath{14.8}}
% NCoercivityRate [doc, lower] doc:step3_proof_v32.tex :: \norm{e^{\tau\mathcal L}u}_\rho\le14.8\,e^{-0.0213\tau}\norm
\newcommand{\NCoercivityRate}{\ensuremath{0.0213}}
% NRegionRampStart [hand, exact] hand :: 
\newcommand{\NRegionRampStart}{\ensuremath{14}}
% NRegionInterface [hand, exact] hand :: 
\newcommand{\NRegionInterface}{\ensuremath{28}}
% NRegionShellEnd [hand, exact] hand :: 
\newcommand{\NRegionShellEnd}{\ensuremath{60}}
% NRegionSegEnd [hand, exact] hand :: 
\newcommand{\NRegionSegEnd}{\ensuremath{350}}
% NRegionAnnEnd [hand, exact] hand :: 
\newcommand{\NRegionAnnEnd}{\ensuremath{700}}
% NZoneIIILow [hand, exact] hand :: 
\newcommand{\NZoneIIILow}{\ensuremath{253}}
% NZoneIIIHigh [hand, exact] hand :: 
\newcommand{\NZoneIIIHigh}{\ensuremath{379.5}}
% NConeCentreMin [hand, exact] hand :: 
\newcommand{\NConeCentreMin}{\ensuremath{300}}
% NConeCentreMax [hand, exact] hand :: 
\newcommand{\NConeCentreMax}{\ensuremath{380}}
% NFarCentreMin [hand, exact] hand :: 
\newcommand{\NFarCentreMin}{\ensuremath{490}}
% NConeL [hand, exact] hand :: 
\newcommand{\NConeL}{\ensuremath{50}}
% NConeX [hand, exact] hand :: 
\newcommand{\NConeX}{\ensuremath{3}}
% NFarL [hand, exact] hand :: 
\newcommand{\NFarL}{\ensuremath{113}}
% NFarX [hand, exact] hand :: 
\newcommand{\NFarX}{\ensuremath{2}}
}%
  {\PackageWarning{paper}{numbers/numbers_paper.tex not found (unit U0 not finished)}}

% ======================================================================
% NOTATION MACROS (OUTLINE 5.2 and notation.md).  One macro per symbol
% whose spelling is fixed by the notation table.  Plain letters (a, b, f, ...)
% have no macro.
% ======================================================================
% basic
\newcommand{\R}{\mathbb{R}}
\newcommand{\C}{\mathbb{C}}
\newcommand{\N}{\mathbb{N}}
\newcommand{\Z}{\mathbb{Z}}
\newcommand{\eps}{\varepsilon}
\newcommand{\ip}[2]{\langle #1,#2\rangle}
\newcommand{\abs}[1]{\lvert #1\rvert}
\newcommand{\norm}[1]{\lVert #1\rVert}
\newcommand{\dd}{\mathrm d}
\newcommand{\re}{\operatorname{Re}}
\newcommand{\im}{\operatorname{Im}}
% flow and curves
\newcommand{\arclen}{\mathfrak{s}}                 % arclength of the flow curve
\newcommand{\kss}{\kappa_{\mathfrak{ss}}}          % kappa_{ss} in the flow equation
\newcommand{\Area}{\operatorname{Area}}
\newcommand{\Len}{\mathcal L}                      % length of a curve
\newcommand{\spp}{\operatorname{sp}}                 % sup quantities sp^v_0, sp^v_1, sp^f_1 (were S^v_0, S^v_1, S^f_1)
\newcommand{\rot}{\operatorname{rot}}              % turning number
\newcommand{\Tsing}{T_*}                           % singular time
\newcommand{\Tmax}{T_{\max}}                       % maximal time of existence
\newcommand{\pstar}{p_*}                           % image point of the singularity
\newcommand{\Gz}{\Gamma_0}                         % initial curve
\newcommand{\vth}{\vartheta}                       % parameter of Gamma_0
\newcommand{\dclk}{\delta_t}                       % clock defect of step 4
\newcommand{\fall}{\varpi}                         % Fourier rates varpi_n = n^4-n^2
% rescaling
\newcommand{\lamone}{\lambda_1}
\newcommand{\rate}{\varsigma}                      % 1/4 - a
\newcommand{\xia}{\xi_a}                           % (1-4\bar a)^{-1}
% shrinker and step 1
\newcommand{\Sone}{S_1}
\newcommand{\Gsh}{\Gamma_{\!*}}                    % profile curve of S_1
\newcommand{\cSet}{\mathcal S}                     % certified set of shrinkers
\newcommand{\sigtip}{\sigma_{\rm tip}}
\newcommand{\Xo}{X_0}                              % 2.2490776515
\newcommand{\Dtip}{D}
\newcommand{\ktip}{k_0}
\newcommand{\shootM}{M}                            % shooting function
\newcommand{\Ypar}{\mathcal Y}                     % step-1 parallelepiped
\newcommand{\axis}{\Sigma}                         % symmetry axis of S_1
\newcommand{\sigc}{\sigma_{\rm c}}
\newcommand{\sigrev}{\sigma^{\rm rev}}
% linear theory
\newcommand{\Hrho}{\mathcal H_\rho}
\newcommand{\muslow}{\mu_3}
\newcommand{\rre}{\mathfrak r_{\rm re}}
\newcommand{\rim}{\mathfrak r_{\rm im}}
\newcommand{\Evev}{\mathrm{Ev}_{\rm even}}
\newcommand{\Evod}{\mathrm{Ev}_{\rm odd}}
\newcommand{\gvec}{\mathbf g}                      % translation vector in step 2
\newcommand{\cons}[1]{\mathfrak f_{#1}}            % constraint functions
\newcommand{\phirho}{\varphi_\rho}
% energies
\newcommand{\Lam}{\Lambda}
\newcommand{\En}{E}
\newcommand{\EnI}{E_I}
\newcommand{\muI}{\mu_I}                           % 0.08
\newcommand{\epsE}{\varepsilon_E}                  % 1e-5
\newcommand{\epsI}{\varepsilon_I}                  % 3e-5
\newcommand{\wn}[1]{\mathsf w_{#1}}                % ||w~^{(j)}||_rho
\newcommand{\vI}{v_I}                              % zeta w'
\newcommand{\muerr}{\mu_{\rm err}}                 % step-4 error-energy weight 0.036
% far field
\newcommand{\pU}{p_U}                              % 1.0246
\newcommand{\Uone}{U_1}
\newcommand{\sm}{\operatorname{sm}}                % smoothstep t^3(10-15t+6t^2)
\newcommand{\rc}{r_c}                              % centre radius (Lemma U)
\newcommand{\rchi}{r_\chi}                         % 253
\newcommand{\Iw}{I_w}                              % hand-off windows [y_w,1.5y_w]
\newcommand{\qa}{\mathfrak q_1}                    % quasilinear coefficients of the slope equation
\newcommand{\qb}{\mathfrak q_2}
% cylinders (Lemma O)
\newcommand{\cyl}{\ell}                            % cylinder length unit
\newcommand{\MO}{M_O}
\newcommand{\KO}{K_O}
\newcommand{\Hcyl}{H_{\rm cyl}}
\newcommand{\XiO}{\Xi_O}
\newcommand{\fO}{f_O}
% step 4
\newcommand{\Ybar}{\bar Y}
\newcommand{\Lt}{\widetilde L}                     % linearisation about the step-4 base
\newcommand{\Trans}{\Pi}                           % transport integral
\newcommand{\resid}{\mathrm{res}}
\newcommand{\Resid}[1]{\mathrm{Res}_{#1}}
\newcommand{\boxtip}{\Box_{\rm tip}}
\newcommand{\boxfar}{\Box_{\rm far}}
\newcommand{\Eenv}{\mathcal E_{\rm env}}
\newcommand{\betab}{\beta_{\rm base}}
\newcommand{\tauI}{\tau^{\rm I}}                   % Phase I clock
% names of families of conditions: \Hh{3} = (H3), \Oh{4} = (O4), ...
\newcommand{\Hh}[1]{\textup{(H#1)}}
\newcommand{\Oh}[1]{\textup{(O#1)}}
\newcommand{\Lh}[1]{\textup{(L#1)}}
\newcommand{\Bh}[1]{\textup{(B#1)}}
\newcommand{\Kh}[1]{\textup{(K#1)}}
\newcommand{\Wh}[1]{\textup{(W#1)}}
\newcommand{\PhaseI}{Phase~I}
\newcommand{\PhaseII}{Phase~II}
% margin: threshold/bound (always)
\newcommand{\mgn}[1]{{\times}#1}

% ---- figures: placeholder macro so that drafts compile before F1-F7 exist ----
\newcommand{\figplaceholder}[1]{\fbox{\parbox{0.8\linewidth}{\centering\small[figure #1 to be supplied]}}}
